\\documentclass{article}

% ---------------------------------------------------------------------------
% 中译本改动（xelatex + 中文支持），详见 VERIFY.md
% xeCJK 必须在 hyperref 之前加载
% ---------------------------------------------------------------------------
\usepackage{xeCJK}
\setCJKmainfont{Songti SC}[ItalicFont={Kaiti SC}]
\usepackage[utf8]{inputenc} %
\usepackage[T1]{fontenc}    %
\usepackage{lmodern}
\usepackage{hyperref}       %
\usepackage{url}            %
\usepackage{booktabs}       %
\usepackage{amsfonts}       %
\usepackage{nicefrac}       %
% \usepackage{microtype}      % 中译本：XeLaTeX 下 microtype 与 Type1 TS1+ptm 冲突
% 触发 "Cannot use XeTeXglyph with ptmr8c" + Missing number 致命错误，故停用（详见 VERIFY.md）

\usepackage{amsmath}
\usepackage{amsthm}
\usepackage{amssymb}
\newcommand\hmmax{0}
\newcommand\bmmax{0}
\usepackage{bm}
\usepackage[title]{appendix}
\usepackage{courier}
\usepackage{enumitem}
\usepackage{graphicx}
\usepackage{algorithmic}
\usepackage[caption=false,font=footnotesize]{subfig}
\usepackage{float} 
\usepackage[margin=1in]{geometry}
\usepackage{authblk}
\usepackage{nicefrac}
\usepackage{mathtools}
\usepackage{color}
\usepackage{xspace} %
\usepackage{times}
\usepackage{stmaryrd}



\newcommand{\llangle}{\langle\!\langle}
\newcommand{\rrangle}{\rangle\!\rangle}

\usepackage{todonotes}

\newcommand{\guillaume}[1]{\todo[color=blue!40,inline,caption={}]{{\it Guillaume:~}#1}}
\newcommand{\Beheshteh}[1]{\todo[color=gray!40,inline,caption={}]{{\it Beheshteh:~}#1}}

\newcommand{\toaddress}[1]{\todo[color=yellow!10,inline,caption={}]{{\it Here we need to:~}#1}}
\newcommand{\dcite}[1]{{\small [\citetitle{#1} (\citeauthor{#1}, \citeyear{#1}) \cite{#1}]}} %
\newcommand{\alert}[1]{{\color{red}{\fontencoding{U}\fontfamily{futs}\selectfont\char 66\relax}
#1}}
\newcommand{\TODO}[1]{{\color{orange}TODO:
#1}}
\newcommand{\Sreg}{\widetilde{\Sigmab}}
\newcommand{\hWb}{\widehat{\Wb}}

\newcommand{\fTT}[1]{f_{\mathrm{TT}(#1)}}
\newcommand{\fCP}[1]{f_{\mathrm{CP}(#1)}}


\newtheorem*{theorem*}{定理}
\newtheorem*{lemma*}{引理}
\newcommand{\bigomega}{\Omega}



\RequirePackage{latexsym}
\RequirePackage{amsthm}
\RequirePackage{amssymb}
\RequirePackage{bm}
\RequirePackage[mathscr]{euscript}

\newcommand{\D}{\text{d}}
\newcommand{\ex}{\mathbb{E}}
\newcommand{\Ex}[1]{\mbox{}\mathbb{E}\left[#1\right]}
\newcommand{\Var}[1]{\mbox{}\textup{Var}\left(#1\right)}
\newcommand{\Cov}[1]{\mbox{}\textup{Cov}\left(#1\right)}
\newcommand{\Tr}[1]{\mbox{}\textup{Tr}\left[#1\right]}
\newcommand{\Trsq}[1]{\mbox{}\textup{tr}^2\left(#1\right)}
\newcommand{\pr}[1]{\mbox{}\mathbb{P}\left(#1\right)}
\newcommand{\kl}[2]{\mbox{}\mathbb{KL}\left(#1\,||\,#2\right)}
\newcommand{\tAb}{\widetilde{\Ab}}
\newcommand{\tp}{\widetilde{p}}
\newcommand{\ts}{\mathsf{T}}
\newcommand{\ind}[1]{\mathds{1}\{#1\}}
\newcommand{\nnz}[1]{\mathsf{nnz}(#1)}
\newcommand{\Diag}[1]{\diag \left\{#1\right\}} %
\newcommand{\Pinv}[1]{\mbox{}\left(#1\right)^{\dagger}} %
\newcommand{\inv}[1]{\mbox{}\left(#1\right)^{-1}} %
\newcommand{\RR}[2]{\mathbb{R}^{#1 \times #2}} %
\newcommand{\R}[1]{\mathbb{R}^{#1}} %
\newcommand{\Sqrt}[1]{#1^{\nicefrac{1}{2}}} %
\newcommand{\defeq}{\triangleq}
\newcommand{\dist}{\mathrm{dist}}
\newcommand{\disteq}{\stackrel{d}{=}}
\newcommand{\deriv}[2]{\frac{\partial #1}{\partial #2}}

\newcommand{\vectorize}{\mathrm{vec}}
\newcommand{\bigo}{\mathcal{O}}


\newcommand{\St}{\bm{\mathcal{S}}}
\newcommand{\At}{\bm{\mathcal{A}}}
\newcommand{\Bt}{\bm{\mathcal{B}}}
\newcommand{\Xt}{\bm{\mathcal{X}}}
\newcommand{\Mt}{\bm{\mathcal{M}}}
\newcommand{\Gt}{\bm{\mathcal{G}}}
\newcommand{\Tt}{\bm{\mathcal{T}}}
\newcommand{\It}{\bm{\mathcal{I}}}
\newcommand{\Yt}{\bm{\mathcal{Y}}}
\newcommand{\Qt}{\bm{\mathcal{Q}}}
\newcommand{\Rt}{\bm{\mathcal{R}}}
\newcommand{\Ht}{\bm{\mathcal{H}}}
\newcommand{\Ut}{\bm{\mathcal{U}}}
\newcommand{\Wt}{\bm{\mathcal{W}}}
\newcommand{\Ct}{\bm{\mathcal{C}}}
\newcommand{\CP}[1]{\llbracket #1 \rrbracket}
\newcommand{\Tucker}[1]{\llbracket #1 \rrbracket}
\newcommand{\T}{\Tt}
\newcommand{\Rbb}{\RR}




\newcommand{\ab}{\mathbf{a}}
\newcommand{\bb}{\mathbf{b}}
\newcommand{\cbb}{\mathbf{c}}
\newcommand{\db}{\mathbf{d}}
\newcommand{\eb}{\mathbf{e}}
\newcommand{\fb}{\mathbf{f}}
\newcommand{\gb}{\mathbf{g}}
\newcommand{\hb}{\mathbf{h}}
\newcommand{\ib}{\mathbf{i}}
\newcommand{\jb}{\mathbf{j}}
\newcommand{\kb}{\mathbf{k}}
\newcommand{\lb}{\mathbf{l}}
\newcommand{\mb}{\mathbf{m}}
\newcommand{\nbb}{\mathbf{n}}
\newcommand{\ob}{\mathbf{o}}
\newcommand{\pb}{\mathbf{p}}
\newcommand{\qb}{\mathbf{q}}
\newcommand{\rb}{\mathbf{r}}
\newcommand{\sbb}{\mathbf{s}}
\newcommand{\tb}{\mathbf{t}}
\newcommand{\ub}{\mathbf{u}}
\newcommand{\vb}{\mathbf{v}}
\newcommand{\wb}{\mathbf{w}}
\newcommand{\xb}{\mathbf{x}}
\newcommand{\yb}{\mathbf{y}}
\newcommand{\zb}{\mathbf{z}}

\newcommand{\abtil}{\tilde{\ab}}
\newcommand{\bbtil}{\tilde{\bb}}
\newcommand{\cbtil}{\tilde{\cbb}}
\newcommand{\dbtil}{\tilde{\db}}
\newcommand{\ebtil}{\tilde{\eb}}
\newcommand{\fbtil}{\tilde{\fb}}
\newcommand{\gbtil}{\tilde{\gb}}
\newcommand{\hbtil}{\tilde{\hb}}
\newcommand{\ibtil}{\tilde{\ib}}
\newcommand{\jbtil}{\tilde{\jb}}
\newcommand{\kbtil}{\tilde{\kb}}
\newcommand{\lbtil}{\tilde{\lb}}
\newcommand{\mbtil}{\tilde{\mb}}
\newcommand{\nbtil}{\tilde{\nbb}}
\newcommand{\obtil}{\tilde{\ob}}
\newcommand{\pbtil}{\tilde{\pb}}
\newcommand{\qbtil}{\tilde{\qb}}
\newcommand{\rbtil}{\tilde{\rb}}
\newcommand{\sbtil}{\tilde{\sbb}}
\newcommand{\tbtil}{\tilde{\tb}}
\newcommand{\ubtil}{\tilde{\ub}}
\newcommand{\vbtil}{\tilde{\vb}}
\newcommand{\wbtil}{\tilde{\wb}}
\newcommand{\xbtil}{\tilde{\xb}}
\newcommand{\ybtil}{\tilde{\yb}}
\newcommand{\zbtil}{\tilde{\zb}}

\newcommand{\atil}{\tilde{a}}
\newcommand{\btil}{\tilde{b}}
\newcommand{\ctil}{\tilde{c}}
\newcommand{\dtil}{\tilde{d}}
\newcommand{\etil}{\tilde{e}}
\newcommand{\ftil}{\tilde{f}}
\newcommand{\gtil}{\tilde{g}}
\newcommand{\htil}{\tilde{h}}
\newcommand{\itil}{\tilde{i}}
\newcommand{\jtil}{\tilde{j}}
\newcommand{\ktil}{\tilde{k}}
\newcommand{\ltil}{\tilde{l}}
\newcommand{\mtil}{\tilde{m}}
\newcommand{\ntil}{\tilde{n}}
\newcommand{\otil}{\tilde{o}}
\newcommand{\ptil}{\tilde{p}}
\newcommand{\qtil}{\tilde{q}}
\newcommand{\rtil}{\tilde{r}}
\newcommand{\stil}{\tilde{s}}
\newcommand{\ttil}{\tilde{t}}
\newcommand{\util}{\tilde{u}}
\newcommand{\vtil}{\tilde{v}}
\newcommand{\wtil}{\tilde{w}}
\newcommand{\xtil}{\tilde{x}}
\newcommand{\ytil}{\tilde{y}}
\newcommand{\ztil}{\tilde{z}}

\newcommand{\ahat}{\hat{a}}
\newcommand{\bhat}{\hat{b}}
\newcommand{\chat}{\hat{c}}
\newcommand{\dhat}{\hat{d}}
\newcommand{\ehat}{\hat{e}}
\newcommand{\fhat}{\hat{f}}
\newcommand{\ghat}{\hat{g}}
\newcommand{\hhat}{\hat{h}}
\newcommand{\ihat}{\hat{i}}
\newcommand{\jhat}{\hat{j}}
\newcommand{\khat}{\hat{k}}
\newcommand{\lhat}{\hat{l}}
\newcommand{\mhat}{\hat{m}}
\newcommand{\nhat}{\hat{n}}
\newcommand{\ohat}{\hat{o}}
\newcommand{\phat}{\hat{p}}
\newcommand{\qhat}{\hat{q}}
\newcommand{\rhat}{\hat{r}}
\newcommand{\shat}{\hat{s}}
\newcommand{\that}{\hat{t}}
\newcommand{\uhat}{\hat{u}}
\newcommand{\vhat}{\hat{v}}
\newcommand{\what}{\hat{w}}
\newcommand{\xhat}{\hat{x}}
\newcommand{\yhat}{\hat{y}}
\newcommand{\zhat}{\hat{z}}

\newcommand{\abhat}{\hat{\ab}}
\newcommand{\bbhat}{\hat{\bb}}
\newcommand{\cbhat}{\hat{\cbb}}
\newcommand{\dbhat}{\hat{\db}}
\newcommand{\ebhat}{\hat{\eb}}
\newcommand{\fbhat}{\hat{\fb}}
\newcommand{\gbhat}{\hat{\gb}}
\newcommand{\hbhat}{\hat{\hb}}
\newcommand{\ibhat}{\hat{\ib}}
\newcommand{\jbhat}{\hat{\jb}}
\newcommand{\kbhat}{\hat{\kb}}
\newcommand{\lbhat}{\hat{\lb}}
\newcommand{\mbhat}{\hat{\mb}}
\newcommand{\nbhat}{\hat{\nb}}
\newcommand{\obhat}{\hat{\ob}}
\newcommand{\pbhat}{\hat{\pb}}
\newcommand{\qbhat}{\hat{\qb}}
\newcommand{\rbhat}{\hat{\rb}}
\newcommand{\sbhat}{\hat{\sb}}
\newcommand{\tbhat}{\hat{\tb}}
\newcommand{\ubhat}{\hat{\ub}}
\newcommand{\vbhat}{\hat{\vb}}
\newcommand{\wbhat}{\hat{\wb}}
\newcommand{\xbhat}{\hat{\xb}}
\newcommand{\ybhat}{\hat{\yb}}
\newcommand{\zbhat}{\hat{\zb}}

\newcommand{\Atil}{\tilde{A}}
\newcommand{\Btil}{\tilde{B}}
\newcommand{\Ctil}{\tilde{C}}
\newcommand{\Dtil}{\tilde{D}}
\newcommand{\Etil}{\tilde{E}}
\newcommand{\Ftil}{\tilde{F}}
\newcommand{\Gtil}{\tilde{G}}
\newcommand{\Htil}{\tilde{H}}
\newcommand{\Itil}{\tilde{I}}
\newcommand{\Jtil}{\tilde{J}}
\newcommand{\Ktil}{\tilde{K}}
\newcommand{\Ltil}{\tilde{L}}
\newcommand{\Mtil}{\tilde{M}}
\newcommand{\Ntil}{\tilde{N}}
\newcommand{\Otil}{\tilde{O}}
\newcommand{\Ptil}{\tilde{P}}
\newcommand{\Qtil}{\tilde{Q}}
\newcommand{\Rtil}{\tilde{R}}
\newcommand{\Stil}{\tilde{S}}
\newcommand{\Ttil}{\tilde{T}}
\newcommand{\Util}{\tilde{U}}
\newcommand{\Vtil}{\tilde{V}}
\newcommand{\Wtil}{\tilde{W}}
\newcommand{\Xtil}{\tilde{X}}
\newcommand{\Ytil}{\tilde{Y}}
\newcommand{\Ztil}{\tilde{Z}}

\newcommand{\abar}{\bar{a}}
\newcommand{\bbar}{\bar{b}}
\newcommand{\cbar}{\bar{c}}
\newcommand{\dbar}{\bar{d}}
\newcommand{\ebar}{\bar{e}}
\newcommand{\fbar}{\bar{f}}
\newcommand{\gbar}{\bar{g}}
\newcommand{\hbr}{\bar{h}}
\newcommand{\ibar}{\bar{i}}
\newcommand{\jbar}{\bar{j}}
\newcommand{\kbar}{\bar{k}}
\newcommand{\lbar}{\bar{l}}
\newcommand{\mbar}{\bar{m}}
\newcommand{\nbar}{\bar{n}}
\newcommand{\pbar}{\bar{p}}
\newcommand{\qbar}{\bar{q}}
\newcommand{\rbar}{\bar{r}}
\newcommand{\sbar}{\bar{s}}
\newcommand{\tbar}{\bar{t}}
\newcommand{\ubar}{\bar{u}}
\newcommand{\vbar}{\bar{v}}
\newcommand{\wbar}{\bar{w}}
\newcommand{\xbar}{\bar{x}}
\newcommand{\ybar}{\bar{y}}
\newcommand{\zbar}{\bar{z}}

\newcommand{\abbar}{\bar{\ab}}
\newcommand{\bbbar}{\bar{\bb}}
\newcommand{\cbbar}{\bar{\cb}}
\newcommand{\dbbar}{\bar{\db}}
\newcommand{\ebbar}{\bar{\eb}}
\newcommand{\fbbar}{\bar{\fb}}
\newcommand{\gbbar}{\bar{\gb}}
\newcommand{\hbbar}{\bar{\hb}}
\newcommand{\ibbar}{\bar{\ib}}
\newcommand{\jbbar}{\bar{\jb}}
\newcommand{\kbbar}{\bar{\kb}}
\newcommand{\lbbar}{\bar{\lb}}
\newcommand{\mbbar}{\bar{\mb}}
\newcommand{\nbbar}{\bar{\nbb}}
\newcommand{\obbar}{\bar{\ob}}
\newcommand{\pbbar}{\bar{\pb}}
\newcommand{\qbbar}{\bar{\qb}}
\newcommand{\rbbar}{\bar{\rb}}
\newcommand{\sbbar}{\bar{\sbb}}
\newcommand{\tbbar}{\bar{\tb}}
\newcommand{\ubbar}{\bar{\ub}}
\newcommand{\vbbar}{\bar{\vb}}
\newcommand{\wbbar}{\bar{\wb}}
\newcommand{\xbbar}{\bar{\xb}}
\newcommand{\ybbar}{\bar{\yb}}
\newcommand{\zbbar}{\bar{\zb}}

\newcommand{\Ab}{\mathbf{A}}
\newcommand{\Bb}{\mathbf{B}}
\newcommand{\Cb}{\mathbf{C}}
\newcommand{\Db}{\mathbf{D}}
\newcommand{\Eb}{\mathbf{E}}
\newcommand{\Fb}{\mathbf{F}}
\newcommand{\Gb}{\mathbf{G}}
\newcommand{\Hb}{\mathbf{H}}
\newcommand{\Ib}{\mathbf{I}}
\newcommand{\Jb}{\mathbf{J}}
\newcommand{\Kb}{\mathbf{K}}
\newcommand{\Lb}{\mathbf{L}}
\newcommand{\Mb}{\mathbf{M}}
\newcommand{\Nb}{\mathbf{N}}
\newcommand{\Ob}{\mathbf{O}}
\newcommand{\Pb}{\mathbf{P}}
\newcommand{\Qb}{\mathbf{Q}}
\newcommand{\Rb}{\mathbf{R}}
\newcommand{\Sb}{\mathbf{S}}
\newcommand{\Tb}{\mathbf{T}}
\newcommand{\Ub}{\mathbf{U}}
\newcommand{\Vb}{\mathbf{V}}
\newcommand{\Wb}{\mathbf{W}}
\newcommand{\Xb}{\mathbf{X}}
\newcommand{\Yb}{\mathbf{Y}}
\newcommand{\Zb}{\mathbf{Z}}

\newcommand{\Abtil}{\widetilde{\Ab}}
\newcommand{\Bbtil}{\tilde{\Bb}}
\newcommand{\Cbtil}{\tilde{\Cb}}
\newcommand{\Dbtil}{\tilde{\Db}}
\newcommand{\Ebtil}{\tilde{\Eb}}
\newcommand{\Fbtil}{\tilde{\Fb}}
\newcommand{\Gbtil}{\tilde{\Gb}}
\newcommand{\Hbtil}{\tilde{\Hb}}
\newcommand{\Ibtil}{\tilde{\Ib}}
\newcommand{\Jbtil}{\tilde{\Jb}}
\newcommand{\Kbtil}{\tilde{\Kb}}
\newcommand{\Lbtil}{\tilde{\Lb}}
\newcommand{\Mbtil}{\tilde{\Mb}}
\newcommand{\Nbtil}{\tilde{\Nb}}
\newcommand{\Obtil}{\tilde{\Ob}}
\newcommand{\Pbtil}{\tilde{\Pb}}
\newcommand{\Qbtil}{\tilde{\Qb}}
\newcommand{\Rbtil}{\tilde{\Rb}}
\newcommand{\Sbtil}{\tilde{\Sb}}
\newcommand{\Tbtil}{\tilde{\Tb}}
\newcommand{\Ubtil}{\tilde{\Ub}}
\newcommand{\Vbtil}{\tilde{\Vb}}
\newcommand{\Wbtil}{\tilde{\Wb}}
\newcommand{\Xbtil}{\tilde{\Xb}}
\newcommand{\Ybtil}{\tilde{\Yb}}
\newcommand{\Zbtil}{\tilde{\Zb}}

\newcommand{\Abar}{\bar{A}}
\newcommand{\Bbar}{\bar{B}}
\newcommand{\Cbar}{\bar{C}}
\newcommand{\Dbar}{\bar{D}}
\newcommand{\Ebar}{\bar{E}}
\newcommand{\Fbar}{\bar{F}}
\newcommand{\Gbar}{\bar{G}}
\newcommand{\Hbar}{\bar{H}}
\newcommand{\Ibar}{\bar{I}}
\newcommand{\Jbar}{\bar{J}}
\newcommand{\Kbar}{\bar{K}}
\newcommand{\Lbar}{\bar{L}}
\newcommand{\Mbar}{\bar{M}}
\newcommand{\Nbar}{\bar{N}}
\newcommand{\Obar}{\bar{O}}
\newcommand{\Pbar}{\bar{P}}
\newcommand{\Qbar}{\bar{Q}}
\newcommand{\Rbar}{\bar{R}}
\newcommand{\Sbar}{\bar{S}}
\newcommand{\Tbar}{\bar{T}}
\newcommand{\Ubar}{\bar{U}}
\newcommand{\Vbar}{\bar{V}}
\newcommand{\Wbar}{\bar{W}}
\newcommand{\Xbar}{\bar{X}}
\newcommand{\Ybar}{\bar{Y}}
\newcommand{\Zbar}{\bar{Z}}

\newcommand{\Abbar}{\bar{\Ab}}
\newcommand{\Bbbar}{\bar{\Bb}}
\newcommand{\Cbbar}{\bar{\Cb}}
\newcommand{\Dbbar}{\bar{\Db}}
\newcommand{\Ebbar}{\bar{\Eb}}
\newcommand{\Fbbar}{\bar{\Fb}}
\newcommand{\Gbbar}{\bar{\Gb}}
\newcommand{\Hbbar}{\bar{\Hb}}
\newcommand{\Ibbar}{\bar{\Ib}}
\newcommand{\Jbbar}{\bar{\Jb}}
\newcommand{\Kbbar}{\bar{\Kb}}
\newcommand{\Lbbar}{\bar{\Lb}}
\newcommand{\Mbbar}{\bar{\Mb}}
\newcommand{\Nbbar}{\bar{\Nb}}
\newcommand{\Obbar}{\bar{\Ob}}
\newcommand{\Pbbar}{\bar{\Pb}}
\newcommand{\Qbbar}{\bar{\Qb}}
\newcommand{\Rbbar}{\bar{\Rb}}
\newcommand{\Sbbar}{\bar{\Sb}}
\newcommand{\Tbbar}{\bar{\Tb}}
\newcommand{\Ubbar}{\bar{\Ub}}
\newcommand{\Vbbar}{\bar{\Vb}}
\newcommand{\Wbbar}{\bar{\Wb}}
\newcommand{\Xbbar}{\bar{\Xb}}
\newcommand{\Ybbar}{\bar{\Yb}}
\newcommand{\Zbbar}{\bar{\Zb}}

\newcommand{\Ahat}{\hat{A}}
\newcommand{\Bhat}{\hat{B}}
\newcommand{\Chat}{\hat{C}}
\newcommand{\Dhat}{\hat{D}}
\newcommand{\Ehat}{\hat{E}}
\newcommand{\Fhat}{\hat{F}}
\newcommand{\Ghat}{\hat{G}}
\newcommand{\Hhat}{\hat{H}}
\newcommand{\Ihat}{\hat{I}}
\newcommand{\Jhat}{\hat{J}}
\newcommand{\Khat}{\hat{K}}
\newcommand{\Lhat}{\hat{L}}
\newcommand{\Mhat}{\hat{M}}
\newcommand{\Nhat}{\hat{N}}
\newcommand{\Ohat}{\hat{O}}
\newcommand{\Phat}{\hat{P}}
\newcommand{\Qhat}{\hat{Q}}
\newcommand{\Rhat}{\hat{R}}
\newcommand{\Shat}{\hat{S}}
\newcommand{\That}{\hat{T}}
\newcommand{\Uhat}{\hat{U}}
\newcommand{\Vhat}{\hat{V}}
\newcommand{\What}{\hat{W}}
\newcommand{\Xhat}{\hat{X}}
\newcommand{\Yhat}{\hat{Y}}
\newcommand{\Zhat}{\hat{Z}}

\newcommand{\Abhat}{\hat{\Ab}}
\newcommand{\Bbhat}{\hat{\Bb}}
\newcommand{\Cbhat}{\hat{\Cb}}
\newcommand{\Dbhat}{\hat{\Db}}
\newcommand{\Ebhat}{\hat{\Eb}}
\newcommand{\Fbhat}{\hat{\Fb}}
\newcommand{\Gbhat}{\hat{\Gb}}
\newcommand{\Hbhat}{\hat{\Hb}}
\newcommand{\Ibhat}{\hat{\Ib}}
\newcommand{\Jbhat}{\hat{\Jb}}
\newcommand{\Kbhat}{\hat{\Kb}}
\newcommand{\Lbhat}{\hat{\Lb}}
\newcommand{\Mbhat}{\hat{\Mb}}
\newcommand{\Nbhat}{\hat{\Nb}}
\newcommand{\Obhat}{\hat{\Ob}}
\newcommand{\Pbhat}{\hat{\Pb}}
\newcommand{\Qbhat}{\hat{\Qb}}
\newcommand{\Rbhat}{\hat{\Rb}}
\newcommand{\Sbhat}{\hat{\Sb}}
\newcommand{\Tbhat}{\hat{\Tb}}
\newcommand{\Ubhat}{\hat{\Ub}}
\newcommand{\Vbhat}{\hat{\Vb}}
\newcommand{\Wbhat}{\hat{\Wb}}
\newcommand{\Xbhat}{\hat{\Xb}}
\newcommand{\Ybhat}{\hat{\Yb}}
\newcommand{\Zbhat}{\hat{\Zb}}

\newcommand{\Acal}{\mathcal{A}}
\newcommand{\Bcal}{\mathcal{B}}
\newcommand{\Ccal}{\mathcal{C}}
\newcommand{\Dcal}{\mathcal{D}}
\newcommand{\Ecal}{\mathcal{E}}
\newcommand{\Fcal}{\mathcal{F}}
\newcommand{\Gcal}{\mathcal{G}}
\newcommand{\Hcal}{\mathcal{H}}
\newcommand{\Ical}{\mathcal{I}}
\newcommand{\Jcal}{\mathcal{J}}
\newcommand{\Kcal}{\mathcal{K}}
\newcommand{\Lcal}{\mathcal{L}}
\newcommand{\Mcal}{\mathcal{M}}
\newcommand{\Ncal}{\mathcal{N}}
\newcommand{\Ocal}{\mathcal{O}}
\newcommand{\Pcal}{\mathcal{P}}
\newcommand{\Qcal}{\mathcal{Q}}
\newcommand{\Rcal}{\mathcal{R}}
\newcommand{\Scal}{\mathcal{S}}
\newcommand{\Tcal}{\mathcal{T}}
\newcommand{\Ucal}{\mathcal{U}}
\newcommand{\Vcal}{\mathcal{V}}
\newcommand{\Wcal}{\mathcal{W}}
\newcommand{\Xcal}{\mathcal{X}}
\newcommand{\Ycal}{\mathcal{Y}}
\newcommand{\Zcal}{\mathcal{Z}}

\newcommand{\Ascr}{\mathscr{A}}
\newcommand{\Bscr}{\mathscr{B}}
\newcommand{\Cscr}{\mathscr{C}}
\newcommand{\Dscr}{\mathscr{D}}
\newcommand{\Escr}{\mathscr{E}}
\newcommand{\Fscr}{\mathscr{F}}
\newcommand{\Gscr}{\mathscr{G}}
\newcommand{\Hscr}{\mathscr{H}}
\newcommand{\Iscr}{\mathscr{I}}
\newcommand{\Jscr}{\mathscr{J}}
\newcommand{\Kscr}{\mathscr{K}}
\newcommand{\Lscr}{\mathscr{L}}
\newcommand{\Mscr}{\mathscr{M}}
\newcommand{\Nscr}{\mathscr{N}}
\newcommand{\Oscr}{\mathscr{O}}
\newcommand{\Pscr}{\mathscr{P}}
\newcommand{\Qscr}{\mathscr{Q}}
\newcommand{\Rscr}{\mathscr{R}}
\newcommand{\Sscr}{\mathscr{S}}
\newcommand{\Tscr}{\mathscr{T}}
\newcommand{\Uscr}{\mathscr{U}}
\newcommand{\Vscr}{\mathscr{V}}
\newcommand{\Wscr}{\mathscr{W}}
\newcommand{\Xscr}{\mathscr{X}}
\newcommand{\Yscr}{\mathscr{Y}}
\newcommand{\Zscr}{\mathscr{Z}}

\newcommand{\Ascrb}{\boldsymbol{\Ascr}}
\newcommand{\Bscrb}{\boldsymbol{\Bscr}}
\newcommand{\Cscrb}{\boldsymbol{\Cscr}}
\newcommand{\Dscrb}{\boldsymbol{\Dscr}}
\newcommand{\Escrb}{\boldsymbol{\Escr}}
\newcommand{\Fscrb}{\boldsymbol{\Fscr}}
\newcommand{\Gscrb}{\boldsymbol{\Gscr}}
\newcommand{\Hscrb}{\boldsymbol{\Hscr}}
\newcommand{\Iscrb}{\boldsymbol{\Iscr}}
\newcommand{\Jscrb}{\boldsymbol{\Jscr}}
\newcommand{\Kscrb}{\boldsymbol{\Kscr}}
\newcommand{\Lscrb}{\boldsymbol{\Lscr}}
\newcommand{\Mscrb}{\boldsymbol{\Mscr}}
\newcommand{\Nscrb}{\boldsymbol{\Nscr}}
\newcommand{\Oscrb}{\boldsymbol{\Oscr}}
\newcommand{\Pscrb}{\boldsymbol{\Pscr}}
\newcommand{\Qscrb}{\boldsymbol{\Qscr}}
\newcommand{\Rscrb}{\boldsymbol{\Rscr}}
\newcommand{\Sscrb}{\boldsymbol{\Sscr}}
\newcommand{\Tscrb}{\boldsymbol{\Tscr}}
\newcommand{\Uscrb}{\boldsymbol{\Uscr}}
\newcommand{\Vscrb}{\boldsymbol{\Vscr}}
\newcommand{\Wscrb}{\boldsymbol{\Wscr}}
\newcommand{\Xscrb}{\boldsymbol{\Xscr}}
\newcommand{\Xscrbh}{\hat{\Xscrb}}
\newcommand{\Yscrb}{\boldsymbol{\Yscr}}
\newcommand{\Zscrb}{\boldsymbol{\Zscr}}

\newcommand{\Atilde}{\widetilde{A}}
\newcommand{\Btilde}{\widetilde{B}}
\newcommand{\Ctilde}{\widetilde{C}}
\newcommand{\Dtilde}{\widetilde{D}}
\newcommand{\Etilde}{\widetilde{E}}
\newcommand{\Ftilde}{\widetilde{F}}
\newcommand{\Gtilde}{\widetilde{G}}
\newcommand{\Htilde}{\widetilde{H}}
\newcommand{\Itilde}{\widetilde{I}}
\newcommand{\Jtilde}{\widetilde{J}}
\newcommand{\Ktilde}{\widetilde{K}}
\newcommand{\Ltilde}{\widetilde{L}}
\newcommand{\Mtilde}{\widetilde{M}}
\newcommand{\Ntilde}{\widetilde{N}}
\newcommand{\Otilde}{\widetilde{O}}
\newcommand{\Ptilde}{\widetilde{P}}
\newcommand{\Qtilde}{\widetilde{Q}}
\newcommand{\Rtilde}{\widetilde{R}}
\newcommand{\Stilde}{\widetilde{S}}
\newcommand{\Ttilde}{\widetilde{T}}
\newcommand{\Utilde}{\widetilde{U}}
\newcommand{\Vtilde}{\widetilde{V}}
\newcommand{\Wtilde}{\widetilde{W}}
\newcommand{\Xtilde}{\widetilde{X}}
\newcommand{\Ytilde}{\widetilde{Y}}
\newcommand{\Ztilde}{\widetilde{Z}}

\renewcommand{\vec}[1]{\mathbf{\boldsymbol{#1}}}

\newcommand{\avec}{\vec{a}}
\newcommand{\bvec}{\vec{b}}
\newcommand{\cvec}{\vec{c}}
\newcommand{\dvec}{\vec{d}}
\newcommand{\evec}{\vec{e}}
\newcommand{\fvec}{\vec{f}}
\newcommand{\gvec}{\vec{g}}
\newcommand{\hvec}{\vec{h}}
\newcommand{\ivec}{\vec{i}}
\newcommand{\jvec}{\vec{j}}
\newcommand{\kvec}{\vec{k}}
\newcommand{\lvec}{\vec{l}}
\newcommand{\mvec}{\vec{m}}
\newcommand{\nvec}{\vec{n}}
\newcommand{\ovec}{\vec{o}}
\newcommand{\pvec}{\vec{p}}
\newcommand{\qvec}{\vec{q}}
\newcommand{\rvec}{\vec{r}}
\newcommand{\svec}{\vec{s}}
\newcommand{\tvec}{\vec{t}}
\newcommand{\uvec}{\vec{u}}
\newcommand{\vvec}{\vec{v}}
\newcommand{\wvec}{\vec{w}}
\newcommand{\xvec}{\vec{x}}
\newcommand{\yvec}{\vec{y}}
\newcommand{\zvec}{\vec{z}}

\newcommand{\va}{\vec{a}}
\newcommand{\vecb}{\vec{b}}
\newcommand{\vc}{\vec{c}}
\newcommand{\vd}{\vec{d}}
\newcommand{\veee}{\vec{e}}
\newcommand{\vf}{\vec{f}}
\newcommand{\vg}{\vec{g}}
\newcommand{\vh}{\vec{h}}
\newcommand{\vi}{\vec{i}}
\newcommand{\vo}{\vec{o}}
\newcommand{\vp}{\vec{p}}
\newcommand{\vq}{\vec{q}}
\newcommand{\vr}{\vec{r}}
\newcommand{\vs}{\vec{s}}
\newcommand{\vt}{\vec{t}}
\newcommand{\vu}{\vec{u}}
\newcommand{\vv}{\vec{v}}
\newcommand{\vx}{\vec{x}}
\newcommand{\vw}{\vec{w}}
\newcommand{\vy}{\vec{y}}
\newcommand{\vz}{\vec{z}}

\newcommand{\Avec}{\vec{A}}
\newcommand{\Bvec}{\vec{B}}
\newcommand{\Cvec}{\vec{C}}
\newcommand{\Dvec}{\vec{D}}
\newcommand{\Evec}{\vec{E}}
\newcommand{\Fvec}{\vec{F}}
\newcommand{\Gvec}{\vec{G}}
\newcommand{\Hvec}{\vec{H}}
\newcommand{\Ivec}{\vec{I}}
\newcommand{\Jvec}{\vec{J}}
\newcommand{\Kvec}{\vec{K}}
\newcommand{\Lvec}{\vec{L}}
\newcommand{\Mvec}{\vec{M}}
\newcommand{\Nvec}{\vec{N}}
\newcommand{\Ovec}{\vec{O}}
\newcommand{\Pvec}{\vec{P}}
\newcommand{\Qvec}{\vec{Q}}
\newcommand{\Rvec}{\vec{R}}
\newcommand{\Svec}{\vec{S}}
\newcommand{\Tvec}{\vec{T}}
\newcommand{\Uvec}{\vec{U}}
\newcommand{\Vvec}{\vec{V}}
\newcommand{\Wvec}{\vec{W}}
\newcommand{\Xvec}{\vec{X}}
\newcommand{\Yvec}{\vec{Y}}
\newcommand{\Zvec}{\vec{Z}}

\newcommand{\Amat}{\Ab}
\newcommand{\Bmat}{\Bb}
\newcommand{\Cmat}{\Cb}
\newcommand{\Dmat}{\Db}
\newcommand{\Emat}{\Eb}
\newcommand{\Fmat}{\Fb}
\newcommand{\Gmat}{\Gb}
\newcommand{\Hmat}{\Hb}
\newcommand{\Imat}{\Ib}
\newcommand{\Lmat}{\Lb}

\newcommand{\Vmat}{\Vb}
\newcommand{\Wmat}{\Wb}
\newcommand{\Xmat}{\Xb}
\newcommand{\Ymat}{\Yb}
\newcommand{\Zmat}{\Zb}


\newcommand{\yvecbar}{\bar{\vec{y}}}
\newcommand{\wvecbar}{\bar{\vec{w}}}
\newcommand{\xvecbar}{\bar{\vec{x}}}
\newcommand{\yvectil}{\tilde{\vec{y}}}
\newcommand{\yvechat}{\hat{\vec{y}}}

\newcommand{\vol}{\text{vol}}




\ifx\BlackBox\undefined
\newcommand{\BlackBox}{\rule{1.5ex}{1.5ex}}  %
\fi

\ifx\QED\undefined
\def\QED{~\rule[-1pt]{5pt}{5pt}\par\medskip}
\fi

\ifx\proof\undefined
\newenvironment{proof}{\par\noindent{\bf Proof\ }}{\hfill\BlackBox\\[2mm]}
\fi

\ifx\theorem\undefined
\newtheorem{theorem}{定理}
\fi

\ifx\example\undefined
\newtheorem{example}{例}
\fi

\ifx\property\undefined
\newtheorem{properties}{性质}
\fi

\ifx\lemma\undefined
\newtheorem{lemma}[theorem]{引理}
\fi

\ifx\proposition\undefined
\newtheorem{proposition}[theorem]{命题}
\fi

\ifx\remark\undefined
\newtheorem{remark}[theorem]{注}
\fi

\ifx\corollary\undefined
\newtheorem{corollary}[theorem]{推论}
\fi

\ifx\definition\undefined
\newtheorem{definition}{定义}
\fi

\ifx\step\undefined
\newtheorem{step}{步骤}
\fi


\ifx\conjecture\undefined
\newtheorem{conjecture}[theorem]{猜想}
\fi

\ifx\axiom\undefined
\newtheorem{axiom}[theorem]{公理}
\fi

\ifx\claim\undefined
\newtheorem{claim}[theorem]{断言}
\fi

\ifx\assumption\undefined
\newtheorem{assumption}[theorem]{假设}
\fi





\newcommand{\CC}{\mathbb{C}} %
\newcommand{\EE}{\mathbb{E}} %
\newcommand{\FF}{\mathbb{F}} %
\newcommand{\HH}{\mathbb{H}} %
\newcommand{\II}{\mathbb{I}} %
\newcommand{\KK}{\mathbb{K}} %
\newcommand{\MM}{\mathbb{M}} %
\newcommand{\NN}{\mathbb{N}} %
\newcommand{\PP}{\mathbb{P}} %
\newcommand{\QQ}{\mathbb{Q}} %
\renewcommand{\RR}{\mathbb{R}} %
\newcommand{\RRbar}{\overline{\RR}} %
\newcommand{\VV}{\mathbb{V}} %
\newcommand{\XX}{\mathbb{X}}
\newcommand{\YY}{\mathbb{Y}}
\newcommand{\ZZ}{\mathbb{Z}} %


\newcommand*{\mini}{\mathop{\mathrm{minimize}}}
\newcommand*{\maxi}{\mathop{\mathrm{maximize}}}
\newcommand*{\argmin}{\mathop{\mathrm{argmin}}}
\newcommand*{\argmax}{\mathop{\mathrm{argmax}}}
\newcommand*{\arginf}{\mathop{\mathrm{arginf}}}
\newcommand*{\argsup}{\mathop{\mathrm{argsup}}}
\newcommand{\sgn}{\mathop{\mathrm{sign}}}
\newcommand{\tr}{\mathop{\mathrm{tr}}}
\newcommand{\diag}{\mathop{\mathrm{diag}}}
\newcommand{\vect}{\mathop{\mathrm{vec}}}
\newcommand{\traj}{\mathop{\mathrm{Traj}}}
\newcommand{\cov}{\mathrm{Cov}}
\newcommand{\conv}{\mathrm{conv}}
\newcommand{\const}{\mathrm{constant}}
\newcommand{\dom}{\mathop{\mathrm{dom}}}
\newcommand{\ri}{\mathop{\mathrm{ri}}}
\newcommand{\cl}{\mathop{\mathrm{cl}}}
\newcommand{\intr}{\mathop{\mathrm{int}}}
\newcommand{\bd}{\mathop{\mathrm{bd}}}
\newcommand{\emp}{\mathop{\mathrm{emp}}}
\newcommand{\core}{\mathop{\mathrm{core}}}
\newcommand{\co}{\mathop{\mathrm{co}}}


\newcommand{\eq}[1]{(\ref{#1})}
\newcommand{\mymatrix}[2]{\left[\begin{array}{#1} #2 \end{array}\right]}
\newcommand{\mychoose}[2]{\left(\begin{array}{c} #1 \\ #2 \end{array}\right)}
\newcommand{\mydet}[1]{\det\left[ #1 \right]}
\newcommand{\myspan}[1]{\mathrm{span}\cbr{#1}}
\newcommand{\smallfrac}[2]{{\textstyle \frac{#1}{#2}}}
\newcommand{\pwrt}[1]{\frac{\partial}{\partial #1}}
\newcommand{\ppwrt}[1]{\frac{\partial^2}{(\partial #1)^2}}
\newcommand{\aleq}{\preccurlyeq}
\newcommand{\ageq}{\succcurlyeq}

\newcommand{\ea}{\emph{et al.} }
\newcommand{\eg}{e.g.\xspace}
\newcommand{\ie}{i.e.\xspace}
\newcommand{\iid}{\emph{iid}}
\newcommand{\wrt}{\emph{w.r.t.}\ }
\newcommand{\cf}{\emph{cf.}\ }


\newcommand{\rbr}[1]{\left(#1\right)}
\newcommand{\sbr}[1]{\left[#1\right]}
\newcommand{\cbr}[1]{\left\{#1\right\}}
\newcommand{\nbr}[1]{\left\|#1\right\|}
\newcommand{\abr}[1]{\left|#1\right|}
\newcommand{\abs}[1]{\left|#1\right|}
\newcommand{\floor}[1]{\left\lfloor #1 \right\rfloor}
\newcommand{\ceil}[1]{\left\lceil #1 \right\rceil}
\newcommand{\inner}[2]{\left\langle #1,#2 \right\rangle}
\newcommand{\norm}[1]{\left\lVert#1\right\rVert}
\newcommand{\ccc}[1]{|\!|\!|#1|\!|\!|}
\newcommand{\sembrack}[1]{[\![#1]\!]}


\newcommand{\one}{\mathbf{1}}  %
\newcommand{\zero}{\mathbf{0}} %
\newcommand{\half}{\frac{1}{2}}
\newcommand{\sqrttwo}{\sqrt{2}}
\newcommand{\invsqrttwo}{\frac{1}{\sqrt{2}}}
\newcommand{\rmd}{\ensuremath{\mathrm{d}}}


\newcommand{\sigmab}{\bm{\sigma}}
\newcommand{\Sigmab}{\mathbf{\Sigma}}
\newcommand{\Thetab}{\mathbf{\Theta}}

\newcommand{\val}{\vec{\alpha}}
\newcommand{\vbeta}{\vec{\beta}}
\newcommand{\vde}{\vec{\delta}}
\newcommand{\vga}{\vec{\gamma}}
\newcommand{\vsig}{\vec{\sigma}}
\newcommand{\vth}{\vec{\theta}}
\newcommand{\veta}{\vec{\eta}}
\newcommand{\vnu}{\vec{\nu}}
\newcommand{\vpsi}{\vec{\psi}}
\newcommand{\vpi}{\vec{\pi}}
\newcommand{\vphi}{\vec{\phi}}
\newcommand{\vxi}{\vec{\xi}}
\newcommand{\vze}{\vec{\zeta}}

\newcommand{\valbar}{\bar{\val}}
\newcommand{\valhat}{\hat{\val}}
\newcommand{\valtil}{\tilde{\val}}
\newcommand{\vthtil}{\tilde{\vth}}
\newcommand{\vthhat}{\hat{\vth}}
\newcommand{\vthbar}{\bar{\vth}}
\newcommand{\vbetatil}{\tilde{\vbeta}}
\newcommand{\vbetahat}{\hat{\vbeta}}

\newcommand{\alphavec}{\val}
\newcommand{\alphavecbar}{\bar{\val}}
\newcommand{\alphavechat}{\hat{\val}}
\newcommand{\alphavectil}{\tilde{\val}}
\newcommand{\betavec}{\vec{\beta}}
\newcommand{\gammavec}{\vec{\gamma}}
\newcommand{\deltavec}{\vec{\delta}}
\newcommand{\etavec}{\vec{\eta}}
\newcommand{\phivec}{\vec{\phi}}
\newcommand{\psivec}{\vec{\psi}}
\newcommand{\thetavec}{\vec{\theta}}
\newcommand{\muvec}{\vec{\mu}}
\newcommand{\xivec}{\vec{\xi}}
\newcommand{\chivec}{\vec{\chi}}
\newcommand{\lambdavec}{\vec{\lambda}}

\newcommand{\alphab}{\boldsymbol{\alpha}}
\newcommand{\betab}{\boldsymbol{\beta}}
\newcommand{\gammab}{\boldsymbol{\gamma}}
\newcommand{\thetab}{\boldsymbol{\theta}}
\newcommand{\mub}{\boldsymbol{\mu}}
\newcommand{\xib}{\boldsymbol{\xi}}
\newcommand{\Deltab}{\boldsymbol{\Delta}}
\newcommand{\Pib}{\boldsymbol{\Pi}}
\newcommand{\etab}{\boldsymbol{\eta}}
\newcommand{\taub}{\boldsymbol{\tau}}
\newcommand{\lambdab}{\boldsymbol{\lambda}}
\newcommand{\pitil}{\tilde{\pi}}
\newcommand{\elltil}{\tilde{\ell}}
\newcommand{\rhob}{\boldsymbol{\rho}}
\newcommand{\Omegab}{\boldsymbol{\Omega}}
\newcommand{\Phib}{\boldsymbol{\Phi}}

\newcommand{\delhat}{\hat{\delta}}
\newcommand{\phihat}{\hat{\phi}}
\newcommand{\delbar}{\bar{\delta}}
\newcommand{\alphahat}{\hat{\alpha}}
\newcommand{\alphatil}{\tilde{\alpha}}
\newcommand{\mubar}{\bar{\mu}}
\newcommand{\thetahat}{\hat{\theta}}
\newcommand{\thetabar}{\bar{\theta}}
\newcommand{\betahat}{\hat{\beta}}

\newcommand{\xibar}{\bar{\xi}}
\newcommand{\lambdavectil}{\tilde{\vec{\lambda}}}
\newcommand{\thetatil}{\tilde{\theta}}
\newcommand{\betatil}{\tilde{\beta}}


\newcommand{\BMRM}{{\sf BMRM}}
\newcommand{\bmrm}{{\sf BMRM}}
\newcommand{\cpm}{{\sf CPM}}
\newcommand{\liblinear}{{\sf liblinear}}
\newcommand{\lsbmrm}{{\sf ls-bmrm}}
\newcommand{\qpbmrm}{{\sf qp-bmrm}}
\newcommand{\pegasos}{{\sf pegasos}}
\newcommand{\pegan}{{\sf pegasos-$n$}}
\newcommand{\pegaone}{{\sf pegasos-1}}
\newcommand{\svmstruct}{{\sf SVM-Struct}}
\newcommand{\svmperf}{{\sf SVM-Perf}}
\newcommand{\nest}{{\sf pragam}}
\newcommand{\nestb}{{\sf pragam-b}}
\newcommand{\mcn}{{\sf M$^3$N}}
\newcommand{\expgrad}{{\sf ExpGrad}}
\newcommand{\smo}{{\sf SMO}}
\newcommand{\lasvm}{{\sf LaSVM}}
\newcommand{\sms}{{\sf SMS}}


\newcommand{\prbep}{{\sf PRBEP}}
\newcommand{\rocarea}{{\sf ROCArea}}

\newcommand{\arow}[2]{#1_{#2\cdot}}
\newcommand{\acol}[2]{#1_{\cdot#2}}
\newcommand{\expunder}[1]{\mathop{\EE}\limits_{#1}}
\newcommand{\mean}[1]{\mathop{\EE}\sbr{{#1}}}
\newcommand{\vars}[2]{\var_{{#1}}\sbr{{#2}}}
\newcommand{\expec}[2]{\mathop{\EE}_{{#1}}\sbr{{#2}}}

\newcommand{\twoco}[1]{\multicolumn{2}{c|}{#1}}
\def\ci{\perp\!\!\!\perp}

\newcommand{\ve}{\varepsilon}

\newcommand{\gstar}{g^{\star}}
\newcommand{\fstar}{f^{\star}}
\newcommand{\hstar}{h^{\star}}
\newcommand{\Astar}{A^{\star}}
\newcommand{\Kstar}{K^{\star}}

\newcommand{\var}{\mathrm{Var}}
\newcommand{\grad}{{\nabla}}
\newcommand{\gradtil}{{\tilde{\grad}}}
\newcommand{\MED}{{\text{MED}}}
\newcommand{\where}{{\quad \text{where} \quad}}
\newcommand{\lcg}{{\textit{l.c.g}}}
\newcommand{\cp}{\mathrm{cp}}
\newcommand{\err}{{\sf err}}
\newcommand{\wvechat}{\hat{\vec{w}}}

\newcommand{\ydiff}[2]{\delta(#1, #2)}
\newcommand{\ykdiff}[2]{\delta(#1, #2)^{\kappa}}
\newcommand{\myfdiff}[2]{\Delta f(#1, #2)}


\newcommand{\Aone}{\textbf{A1}}
\newcommand{\Atwo}{\textbf{A2}}

\newcommand{\kron}{\otimes}
\newcommand{\krao}{\odot}
\newcommand{\bigkrao}{\bigodot}
\newcommand{\hadam}{*}
\newcommand{\outprod}{\circ}
\newcommand{\bigoutprod}{\bigcirc}
\newcommand{\dirsum}{\oplus}

\newcommand{\cprank}{\mathrm{rank}_{\mathrm{CP}}}
\newcommand{\rank}{\mathrm{rank}}



\newcommand\textequal[1]{\stackrel{\mathclap{\scalebox{.75}{\normalfont\mbox{#1}}}}{=}}


\renewcommand{\vec}[1]{\ensuremath{\mathbf{#1}}}
\newcommand{\vecs}[1]{\ensuremath{\mathbf{\boldsymbol{#1}}}}
\newcommand{\mat}[1]{\ensuremath{\mathbf{#1}}}
\newcommand{\mats}[1]{\ensuremath{\mathbf{\boldsymbol{#1}}}}
\newcommand{\ten}[1]{\mat{\ensuremath{\boldsymbol{\mathcal{#1}}}}}

\newcommand{\ttm}[1]{\times_{#1}}
\newcommand{\ttv}[1]{\bullet_{#1}}

\newcommand{\TN}{\textsc{TN}}
\newcommand{\isomorphic}{\simeq}

\newcommand{\tenmat}[2]{\mat{#1}_{(#2)}}
\newcommand{\tenmatpar}[2]{(#1)_{(#2)}}

\DeclareMathOperator{\elm}{\bm{\mathcal{T}}_{\boldsymbol{-}}}
\newcommand{\dbold}{\mathbf{d}}
\newcommand{\dtuple}{(d_1, \ldots, d_N)}
\newcommand{\dprod}{d_1 \times \cdots \times d_N}
\newcommand{\varprod}[2]{#1_1 \times \cdots \times #1_{#2}}
\newcommand{\idxn}{i_1, \ldots, i_N}
\newcommand{\varidx}[2]{#1_1, \ldots ,#1_{#2}}

\newcommand{\supidx}[2]{{#1}^{(#2)}}

\newcommand{\edgelabel}[1]{\scalebox{0.7}{\textcolor{gray}{#1}}}
\DeclareMathOperator*{\TT}{\operatorname{TT}}

\newcommand{\dimcolor}{\textcolor{gray}}
\newcommand{\indcolor}{\textcolor{blue}}





\DeclareMathOperator*{\mpo}{\operatorname{MPO}}
\newcommand{\ttrad}{\mathrm{TT}_\text{Rad}}
\newcommand{\ttnormal}{\mathrm{TT}_\mathcal{N}}
\newcommand{\cpnormal}{\mathrm{CP}_\mathcal{N}}
\usepackage{xcolor}
\hypersetup{
    colorlinks,
    linkcolor={red!50!black},
    citecolor={blue!50!black},
    urlcolor={blue!80!black}
}


% ---- 中译本：其余固定名称中文化（置于文件末尾以覆盖宏包默认值）----
\renewcommand{\proofname}{证明}
\renewcommand{\contentsname}{目录}
\renewcommand{\figurename}{图}
\renewcommand{\tablename}{表}
\renewcommand{\refname}{参考文献}
\renewcommand{\appendixname}{附录}\section*{记号}


    \begin{tabular}{ll}
        $\At$ & 一般张量，包含 $n \geq 0$ 个模 \\
        $\At\in \RR^{d_1 \times d_2 \times d_3}$ & 形状为 $(d_1, d_2, d_3)$ 的三阶张量 $\At$ \\
        $\At_{i,j,k}$ & 三阶张量 $\At$ 的第 $(i,j,k)$ 个元素 \\
        $\At_{i,:,k}$ & 三阶张量 $\ten{A}$ 的模 2 纤维，等价于向量 $\vb\in \RR^{d_2}$ \\
        $\At_{(2)}$ & （三阶）张量 $\At$ 的模 2 展开，等价于矩阵 $\Mb\in \RR^{d_2 \times d_1 d_3}$ \\
        $\vectorize(\At)$ & 把 $\At$ 的模 $N$ 纤维逐一拼接所得的向量，等价于向量 $\ab\in\RR^{d_1d_2d_3}$\\
        $\At\outprod \Bt$ & 张量 $\At$ 与 $\Bt$ 的张量积（即外积） \\
         $\inner{\At}{\Bt}$ & 张量 $\At$ 与 $\Bt$ 的内积\\
         $\norm{\At}_F$  &  张量 $\At$ 的 Frobenius 范数\\
         $\At\kron\Bt$ & 张量 $\At$ 与 $\Bt$ 的 Kronecker 积\\
         \\
         $\Ab$ & 矩阵；等价地，秩 2 张量或二阶张量\\
         $\Ab_{ij}$ & 矩阵 $\Ab$ 的第 $ij$ 个元素\\
         $\Ab^{-1}$ & 矩阵 $\Ab$ 的逆矩阵\\
         $\Ab^\ts$ & 矩阵 $\Ab$ 的转置矩阵\\
         $\Ib$ & 单位矩阵\\
         \\
         $\Ab\krao\Bb$ & 矩阵 $\Ab$ 与 $\Bb$ 的 Khatri-Rao 积\\
         $\Ab\hadam\Bb$ & 矩阵 $\Ab$ 与 $\Bb$ 的 Hadamard 积 \\ 
         \\
         $\ab$ & 列向量；秩 1 张量\\
         $\ab_{i}$ & 向量 $\ab$ 的第 $i$ 个元素 \\
         $a$ & 标量；秩 0 张量\\
         
         
\end{tabular}
 \newpage
\section*{张量网络记号} 
    \begin{table}[H]
    \begin{tabular}{ll}
       $N$ 阶张量 $\At$ &
    $
    \begin{tikzpicture}[baseline=-0.5ex]
    \tikzset{tensor/.style = {minimum size = 0.5cm,shape = circle,thick,draw=black,inner sep = 0pt}, edge/.style = {thick,line width=.4mm},every loop/.style={}}
    \def\x{6}
    \node[tensor,fill=green!50!red!50!white] (T) at (\x,0) {$\At$};
    \draw[edge] (T) -- (\x-0.6,0)  node [midway,above] {\scalebox{0.7}{\dimcolor{$d_1$}}};
    \draw[edge] (T) -- (\x+0.6,0) node [midway,above] {\scalebox{0.7}{\dimcolor{$d_N$}}};
    \draw[edge] (T) -- (\x-0.5,-0.3) node [midway,below] {\scalebox{0.7}{\dimcolor{$d_2$}}};
    \draw[dotted] (T) -- (\x,-0.6);
    \draw[dotted] (T) -- (\x,-0.6);
    \draw[dotted] (T) -- (\x+0.5,-0.5);
    \draw[dotted] (T) -- (\x+0.3,-0.6);
    \end{tikzpicture}\in \RR^{d_1\times \cdots \times d_N} $\\
    三阶张量 $\At $ & $  
    \begin{tikzpicture}[baseline=-0.5ex]
    \tikzset{tensor/.style = {minimum size = 0.5cm,shape = circle,thick,draw=black,inner sep = 0pt}, edge/.style = {thick,line width=.4mm},every loop/.style={}}
    \def\y{0}
    \def\x{0}
    \node[tensor,fill=green!50!red!50!white] (A) at (\x,\y){\scalebox{0.85}{$\At$}};
    \draw[edge] (\x+0.25,\y) -- (\x+0.6,\y)  node [midway,above] {\scalebox{0.7}{\dimcolor{$d_2$}}};
    \draw[edge] (\x-0.25,\y) -- (\x-0.6,\y)  node [midway,above] {\scalebox{0.7}{\dimcolor{$d_1$}}};
    \draw[edge] (\x,\y-0.25) -- (\x,\y-0.7)  node [midway,left] {\scalebox{0.7}{\dimcolor{$d_3$}}};
    \end{tikzpicture}\in \RR^{d_1 \times d_2 \times d_3}$\\
    张量元素 $\At_{i_1,i_2,i_3}$ &
    $\begin{tikzpicture}[baseline=-0.5ex]
    \tikzset{tensor/.style = {minimum size = 0.5cm,shape = circle,thick,draw=black,inner sep = 0pt}, edge/.style = {thick,line width=.4mm},every loop/.style={}}
    \def\y{0}
    \def\x{0}
    \node[tensor,fill=green!50!red!50!white] (A) at (\x,\y){\scalebox{0.85}{$\At$}};
    \draw[edge] (\x+0.25,\y) -- (\x+0.6,\y)  node [very near end, right] {\scalebox{0.7}{\indcolor{$i_2$}}};
    \draw[edge] (\x-0.25,\y) -- (\x-0.6,\y)  node [very near end, left] {\scalebox{0.7}{\indcolor{$i_1$}}};
    \draw[edge] (\x,\y-0.25) -- (\x,\y-0.7)  node [very near end, below] {\scalebox{0.7}{\indcolor{$i_3$}}};
    \end{tikzpicture}\in \RR$ \\
    $\At\in\RR^{d_1\times d_2 \times d_3}$ 的模 2 矩阵化 & $ \Ab_{(2)}=\begin{tikzpicture}[baseline=\baseline]
    \tikzset{tensor/.style = {minimum size = 0.5cm,shape = circle,thick,draw=black,inner sep = 0pt}, edge/.style = {thick,line width=.4mm},every loop/.style={}}
    \def\y{0}
    \def\x{0}
    \node[tensor,fill=green!50!red!50!white] (A) at (\x,\y){\scalebox{0.85}{$\At$}};
    \draw[edge] (\x+0.2,\y+0.15) -- (\x+0.5,\y+0.15)-- (\x+0.5,\y+0.05) -- (\x+1,\y+0.05) node [midway,above]
    {\scalebox{0.7}{\dimcolor{$d_1$}}};
    \draw[edge] (\x+0.2,\y-0.15) -- (\x+0.5,\y-0.15)-- (\x+0.5,\y-0.05)-- (\x+1,\y-0.05) node [midway,below]
    {\scalebox{0.7}{\dimcolor{$d_3$}}};
    \draw[edge] (A) -- (\x-0.7,\y)node [midway,above]
    {\scalebox{0.7}{\dimcolor{$d_2$}}};
    \end{tikzpicture} \in \RR^{d_2\times d_1d_3}$ \\ $\At\in\RR^{d_1\times d_2\times d_3}$ 的向量化 & $\vectorize(\At)=
    \begin{tikzpicture}[baseline=\baseline]
    \tikzset{tensor/.style = {minimum size = 0.5cm,shape = circle,thick,draw=black,inner sep = 0pt}, edge/.style = {thick,line width=.4mm},every loop/.style={}}
    \def\y{0}
    \def\x{0}
    \node[tensor,fill=green!50!red!50!white] (A) at (\x,\y){\scalebox{0.85}{$\At$}};
    \draw[edge] (\x+0.15,\y+0.2) -- (\x+0.5,\y+0.2) -- (\x+0.5,\y+0.1) -- (\x+1,\y+0.1) node [midway,above]
    {\scalebox{0.7}{\dimcolor{$d_1$}}};
     \draw[edge] (A) -- (\x+1,\y)node [right]
    {\scalebox{0.7}{\dimcolor{$d_2$}}};

    \draw[edge] (\x+0.15,\y-0.2) -- (\x+0.5,\y-0.2)-- (\x+0.5,\y-0.1)-- (\x+1,\y-0.1) node [midway,below]
    {\scalebox{0.7}{\dimcolor{$d_3$}}};
    \end{tikzpicture}\in\RR^{d_1d_2d_3}$\\
    两个张量的外积 & $ \At \outprod \Bb =  \begin{tikzpicture}[baseline=-0.5ex]
    \tikzset{tensor/.style = {minimum size = 0.5cm,shape = circle,thick,draw=black,inner sep = 0pt}, edge/.style = {thick,line width=.4mm},every loop/.style={}}
    \def\x{6}
    \def\y{0}
    \node[tensor,fill=green!50!red!50!white] (A) at (\x,\y) {$\At$};
    \draw[edge] (A) -- (\x-0.6,\y) node [midway,above] {\scalebox{0.7}{\dimcolor{$d_1$}}};
    \draw[edge] (A) -- (\x+0.6,\y) node [midway,above] {\scalebox{0.7}{\dimcolor{$d_3$}}};
    \draw[edge] (A) -- (\x,\y-0.6) node [midway,left] {\scalebox{0.7}{\dimcolor{$d_2$}}};
    \end{tikzpicture}
    \hspace{0.3cm}
    \begin{tikzpicture}[baseline=-0.5ex]
    \tikzset{tensor/.style = {minimum size = 0.5cm,shape = circle,thick,draw=black,inner sep = 0pt}, edge/.style = {thick,line width=.4mm},every loop/.style={}}
    \def\x{6}
    \def\y{0}
    \node[tensor,fill=green!70!red!20!white] (B) at (\x+1,\y) {$\Bb$};
    \draw[edge] (B) -- (\x+0.3,0) node [midway,above] {\scalebox{0.7}{\dimcolor{$d_4$}}};
    \draw[edge] (B) -- (\x+1.6,0)node [midway,above] {\scalebox{0.7}{\dimcolor{$d_5$}}};
    \end{tikzpicture} \in \RR^{d_1\times d_2 \times d_3 \times d_4\times d_5}$ \\
    \\
    两个张量的内积 &$\inner{\At}{\Bt} = \begin{tikzpicture}[baseline=\baseline]
    \tikzset{tensor/.style = {minimum size = 0.5cm,shape = circle,thick,draw=black,inner sep = 0pt}, edge/.style = {thick,line width=.4mm},every loop/.style={}}
    \def\y{0}
    \def\x{0}
    \node[tensor,fill=green!50!red!50!white] (St) at (\x,\y){\scalebox{0.85}{$\At$}};
    \node[tensor,fill=green!50!red!50!white] (Tt) at (\x+1,\y){\scalebox{0.85}{$\Bt$}};
    \draw[edge] (St) -- (Tt);
    \draw[edge] (St) to [out=45,in=135] (Tt);
    \draw[edge] (St) to [out=-45,in=-135] (Tt);
    \node[draw=none] () at (\x+0.5,\y+0.5) {\scalebox{0.7}{\dimcolor{$d_1$}}};
    \node[draw=none] () at (\x+0.5,\y+0.15) {\scalebox{0.7}{\dimcolor{$d_2$}}};
    \node[draw=none] () at (\x+0.5,\y-0.5) {\scalebox{0.7}{\dimcolor{$d_3$}}};
\end{tikzpicture} \in\RR $\\
    \\
    张量的 Frobenius 范数 & $\norm{\At}_F^2 = \begin{tikzpicture}[baseline=\baseline]
    \tikzset{tensor/.style = {minimum size = 0.5cm,shape = circle,thick,draw=black,inner sep = 0pt}, edge/.style = {thick,line width=.4mm},every loop/.style={}}
    \def\y{0}
    \def\x{0}
    \node[tensor,fill=green!50!red!50!white] (St) at (\x,\y){\scalebox{0.85}{$\At$}};
    \node[tensor,fill=green!50!red!50!white] (Tt) at (\x+1,\y){\scalebox{0.85}{$\At$}};
    \draw[edge] (St) -- (Tt);
    \draw[edge] (St) to [out=45,in=135] (Tt);
    \draw[edge] (St) to [out=-45,in=-135] (Tt);
    \node[draw=none] () at (\x+0.5,\y+0.5) {\scalebox{0.7}{\dimcolor{$d_1$}}};
    \node[draw=none] () at (\x+0.5,\y+0.15) {\scalebox{0.7}{\dimcolor{$d_2$}}};
    \node[draw=none] () at (\x+0.5,\y-0.5) {\scalebox{0.7}{\dimcolor{$d_3$}}};
\end{tikzpicture}\in\RR$\\
    \\
    两个张量的 Kronecker 积 & $\At\kron\Bt=   \begin{tikzpicture}[baseline=-0.5ex]
    \tikzset{tensor/.style = {minimum width = 2.2cm,minimum height = 0.5cm,shape = rounded rectangle,thick,draw=black,inner sep = 0pt}, edge/.style = {thick,line width=.4mm},every loop/.style={}}
    \def\x{0.5}
    \def\y{0.25}
    \draw[edge] (\x-0.85,\y-0.3) -- (\x-0.85,\y-0.6) -- (\x-0.65,\y-0.6) -- (\x-0.65,\y-0.8);
    \draw[edge] (\x-0.35,\y-0.3) -- (\x-0.35,\y-0.6) -- (\x-0.55,\y-0.6) -- (\x-0.55,\y-0.8)node[below]{\scalebox{0.7}{\dimcolor{$d_1d_4$}}};
    \draw[edge] (\x+0.85,\y-0.3) -- (\x+0.85,\y-0.6) -- (\x+0.65,\y-0.6) -- (\x+0.65,\y-0.8);
    \draw[edge] (\x+0.35,\y-0.3) -- (\x+0.35,\y-0.6) -- (\x+0.55,\y-0.6) -- (\x+0.55,\y-0.8)node[below]{\scalebox{0.7}{\dimcolor{$d_3d_6$}}};
    \draw[edge](\x-0.25,\y-0.3)--(\x-0.25,\y-0.6)-- (\x-0.05,\y-0.6)--(\x-0.05,\y-0.8);
    \draw[edge](\x+0.25,\y-0.3)--(\x+0.25,\y-0.6)-- (\x+0.05,\y-0.6)--(\x+0.05,\y-0.8)node[below]{\scalebox{0.7}{\dimcolor{$d_2d_5$}}};
    \node[tensor,fill=green!30!white] (At) at (\x,\y-0.2){$\At\kron\Bt$};
    \end{tikzpicture}\in\RR^{d_1d_4\times d_2d_5\times d_3d_6}$\\

    二阶张量~（矩阵）$\Ab$ & $\begin{tikzpicture}[baseline=-0.5ex]
    \tikzset{tensor/.style = {minimum size = 0.5cm,shape = circle,thick,draw=black,inner sep = 0pt}, edge/.style = {thick,line width=.4mm},every loop/.style={}}
    \def\x{0}
    \def\y{0}
    \node[tensor,fill=red!50!white] (A) at (\x+1,\y){\scalebox{0.85}{$\Ab$}};
    \draw[edge] (\x+1.25,\y) -- (\x+1.6,\y)  node [midway,above] {\scalebox{0.7}{\dimcolor{$d_2$}}};
    \draw[edge] (\x+0.75,\y) -- (\x+0.4,\y)  node [midway,above] {\scalebox{0.7}{\dimcolor{$d_1$}}};
    \end{tikzpicture}\in\RR^{d_1\times d_2}$\\
    \\
    单位矩阵 &$\Ib_d =  \begin{tikzpicture}[baseline=\baseline]
    \tikzset{tensor/.style = {minimum size = 0.5cm,shape = circle,thick,draw=black,inner sep = 0pt}, edge/.style   = {thick,line width=.4mm},every loop/.style={}}
    \def\y{0}
    \def\x{0}
    \draw[edge](\x,\y)--(\x+1.5,\y);
    \node[draw=none] () at (\x+0.75,\y+0.15) {\scalebox{0.7}{\dimcolor{$d$}}};
    \end{tikzpicture}\in\RR^{d\times d}$\\  
    \\
    Kronecker delta 记号 &     $\begin{tikzpicture}[baseline=\baseline]
    \tikzset{tensor/.style = {minimum size = 0.5cm,shape = circle,thick,draw=black,inner sep = 0pt}, edge/.style   = {thick,line width=.4mm},every loop/.style={}}
    \def\y{0}
    \def\x{0}
    \draw  node[fill = black,circle,inner sep=0pt,minimum size=4pt] (m) at (\x,\y) {};
    \draw  node[fill = black,circle,inner sep=0pt,minimum size=4pt] (m) at (\x,\y) {};
    \draw[edge] (m) -- (\x-0.6,0) node [very near end,left] {\scalebox{0.7}{\indcolor{$i$}}};
    \draw[edge] (m) -- (\x+0.6,0) node [very near end,right]
    {\scalebox{0.7}{\indcolor{$j$}}};
    \node[draw=none] at (\x+1,\y) {\scalebox{1}{\textcolor{black}{$=$}}};;
    \draw[edge] (\x+1.3,\y) -- (\x+2.5,\y) node [very near start,left=0.05cm] {\scalebox{0.7}{\indcolor{$i$}}} node [very near end,right=0.05cm] {\scalebox{0.7}{\indcolor{$j$}}};
    \end{tikzpicture} = \delta_{ij} = \begin{cases}
      1 & \text{if } i=j\\
      0 & \text{otherwise},
    \end{cases}
    =
    \Ib_{ij}$\\
    
    Khatri-Rao 积  &  
    $\Ab\krao\Bb=  \begin{tikzpicture}[baseline=\baseline]
    \tikzset{tensor/.style = {minimum size = 0.5cm,shape = circle,thick,draw=black,inner sep = 0pt}, edge/.style   = {thick,line width=.4mm},every loop/.style={}}
    \def\y{0}
    \def\x{0}
    \node[tensor,fill=red!50!white] (A) at (\x,\y){$\Ab$};
    \node[tensor,fill=blue!50!white] (B) at (\x+1,\y){$\Bb$};
    \draw[edge] (\x,\y+0.25) -- (\x,\y+0.5) -- (\x+0.5,\y+0.7);
    \draw[edge] (\x+1,\y+0.25) -- (\x+1,\y+0.5) -- (\x+0.5,\y+0.7);
    \draw[edge] (\x+0.5,\y+0.7) -- (\x+0.5,\y+1.2)node [midway,right] {\scalebox{0.7}{\dimcolor{$R$}}};
    \draw  node[fill = black,circle,inner sep=0pt,minimum size=4pt] (m) at (\x+0.5,\y+0.7) {};
    \draw[edge] (\x,\y-0.25) -- (\x,\y-0.7);
    \draw[edge] (\x+1,\y-0.25) -- (\x+1,\y-0.7);
    \draw[edge] (\x+1,\y-0.7) -- (\x+0.5,\y-0.7) -- (\x+0.5,\y-1)node [below] {\scalebox{0.7}{\dimcolor{$d_1d_2$}}};
    \draw[edge] (\x+1,\y-0.25) -- (\x+1,\y-0.7);
    \draw[edge] (\x,\y-0.7) -- (\x+0.4,\y-0.7) -- (\x+0.4,\y-1);
    \end{tikzpicture}\in\RR^{d_1d_2\times R}$\\
    Hadamard 积 & $\Ab\hadam\Bb =  \begin{tikzpicture}[baseline=\baseline]
    \tikzset{tensor/.style = {minimum size = 0.5cm,shape = circle,thick,draw=black,inner sep = 0pt}, edge/.style   = {thick,line width=.4mm},every loop/.style={}}
    \def\y{0}
    \def\x{0}
    \node[tensor,fill=red!50!white] (A) at (\x,\y){$\Ab$};
    \node[tensor,fill=blue!50!white] (B) at (\x+1,\y){$\Bb$};
    \draw[edge] (\x,\y+0.25) -- (\x,\y+0.5) -- (\x+0.5,\y+0.7);
    \draw[edge] (\x+1,\y+0.25) -- (\x+1,\y+0.5) -- (\x+0.5,\y+0.7);
    \draw[edge] (\x+0.5,\y+0.7) -- (\x+0.5,\y+1)node [midway,right] {\scalebox{0.7}{\dimcolor{$d_1$}}};
    \draw  node[fill = black,circle,inner sep=0pt,minimum size=4pt] (m) at (\x+0.5,\y+0.7) {};
    \draw[edge] (\x,\y-0.25) -- (\x,\y-0.5) -- (\x+0.5,\y-0.7);
    \draw[edge] (\x+1,\y-0.25) -- (\x+1,\y-0.5) -- (\x+0.5,\y-0.7);
    \draw[edge] (\x+0.5,\y-0.7) -- (\x+0.5,\y-1)node [midway,right] {\scalebox{0.7}{\dimcolor{$d_2$}}};
    \draw  node[fill = black,circle,inner sep=0pt,minimum size=4pt] (m) at (\x+0.5,\y-0.7) {};
    \end{tikzpicture}\in\RR^{d_1\times d_2}$ \\ 
    一阶张量~（向量）$\ab$
    &     $\begin{tikzpicture}[baseline=-0.5ex]
    \tikzset{tensor/.style = {minimum size = 0.5cm,shape = circle,thick,draw=black,inner sep = 0pt}, edge/.style = {thick,line width=.4mm},every loop/.style={}}
    \def\y{0}
    \def\x{0}
    \node[tensor,fill=blue!40!red!60!white] (a) at (\x,\y){\scalebox{0.85}{$\ab$}};
    \draw[edge] (\x+0.25,\y) -- (\x+0.6,\y)  node [midway,above] {\scalebox{0.7}{\dimcolor{$d$}}};
\end{tikzpicture}\in\RR^d$\\
\\
    零阶张量~（标量）$a$ &
    $\begin{tikzpicture}[baseline=-0.5ex]
    \tikzset{tensor/.style = {minimum size = 0.5cm,shape = circle,thick,draw=black,inner sep = 0pt}, edge/.style = {thick,line width=.4mm},every loop/.style={}}
    \def\y{0}
    \def\x{0}
    \node[tensor,fill=blue!20!white] (a) at (\x,\y){\scalebox{0.85}{$a$}};
    \end{tikzpicture}\in\RR$
    
    \end{tabular}
 \end{table}

\newpage
\section*{引言}
高维数据在科学与工程的众多领域中自然出现，涵盖机器学习、信号处理、计算物理与统计学。这类数据常以张量表示，亦即矩阵向高维的推广。张量虽为多模态结构提供了自然的表示，但随着张量阶数增大，直接操作它很快变得困难：参数数目呈指数增长，涉及大量指标的代数表达式也难以理解与实现。

张量网络（tensor network, TN）为应对这些挑战提供了有效的框架。
由~\cite{penrose1971applications} 最早提出、又在量子物理界得到充分发展的张量网络图形语言，把张量收缩编码为图中的边；这一形式体系不仅减少了记号负担，还揭示出被指标记法掩盖的张量运算结构性质。尽管高维张量在现代机器学习与数值分析中居于核心地位，张量网络图在量子计算之外仍未得到充分利用，部分原因是缺少一部自成体系、面向广大技术读者的数学参考书。
本文稿的目标是提供一部关于张量网络及其在张量代数中用法的自洽指南。我们以图示记法讲述张量的主要运算，如收缩、乘积与重排，并说明经典张量分解与相关计算如何在这一框架下自然表达。此外，我们展示张量网络如何简化梯度的推导以及高维概率分布的处理。

贯穿全文，我们说明图示方法并不只是简化记号；它对经典恒等式、秩的界与梯度公式给出了真正更短、更清晰的证明，而这些结果若改用指标运算则需繁复的推导。 



\subsection*{结构安排}

本文稿共分五章，各章以前一章为基础，逐步建立起 
关于张量网络的完整而自洽的论述。

\medskip
\noindent\textbf{第 1 章：张量网络基础。}
我们从第一性原理出发引入张量网络的图形语言。张量表示为带标注腿
（leg）的节点，收缩则编码为共享的边。我们
在这一框架内建立核心运算：内积、外积、迹与 Frobenius 范数，并说明熟知的矩阵恒等式如何自然推广到
高阶张量。本章以复制张量（超边）与矩阵分解的图形
表示收尾，包括 QR 分解与奇异值
分解（SVD）。

\medskip
\noindent\textbf{第 2 章：张量上的运算。}
我们以形式化与图形化两种方式讨论张量的重排运算、纤维、切片、矩阵化与向量化。随后建立主要的张量乘积：模 $n$ 
乘积、Kronecker 积、Khatri-Rao 积与 Hadamard 积，每种乘积都配有其关键恒等式的图示证明。本章以一个一般性定理收尾：借助张量网络图中的割为张量矩阵化的秩给出上界，从而把若干经典秩不等式统一于同一框架之下。

\medskip
\noindent\textbf{第 3 章：张量分解。}
我们介绍高维张量的三种主要分解模型。CANDECOMP/PARAFAC（CP）分解把张量写成若干秩一项之和；Tucker 分解借助压缩的核心张量与因子矩阵将此推广，
而张量列（TT）分解给出链式结构的因子分解，其参数个数随张量阶线性增长。对每种模型，我们给出数学表述、相应的张量网络图、秩的刻画与计算算法，包括用于 Tucker 分解的高阶 SVD（HOSVD），以及用于 TT 的 TT-SVD 与交替最小二乘（ALS）算法。我们进一步讨论 TT 格式下的高效代数运算，并简要概述更先进的架构，包括张量环分解、投影纠缠对态（PEPS）分解与层次 Tucker 分解。

\medskip
\noindent\textbf{第 4 章：计算张量网络的梯度。}
我们为张量网络函数的求导建立一套图示演算。从梯度与雅可比矩阵的经典定义出发，我们证明：张量网络关于任一核心张量的雅可比矩阵，
只需把该核心从图中删去即得。由此得到对经典求导恒等式的透明、纯图示层面的证明，
包括二次型与迹泛函的梯度，并自然推广到同一核心出现多次的网络。我们以两个应用示范这一框架：CP 分解的
基于梯度的优化，以及经由以矩阵乘积算符（MPO）格式参数化的
权重矩阵的反向传播。

\medskip
\noindent\textbf{第 5 章：概率分布与随机张量网络。}
我们探讨张量网络在概率与统计中的
两个相互关联的应用。首先，我们说明离散随机变量的
联合概率分布如何编码为张量，以及边缘分布与条件分布如何化为简单的
图示运算。随后我们讨论张量列格式如何为高维分布提供
高效而结构化的参数化，其中包括源自量子物理的
Born 机器式表示。其次，我们计算元素取自标准正态
分布的随机张量网络的期望。借助 Isserlis 定理与图示推理，我们为
形如
$\EE[\|\Ab\Bb\|_F^2]$ 与
$\EE[(\Ab^\ts \Ab)^{\kron 2}]$ 这类量导出闭式表达式，其简洁程度是常规基于指标的方法难以达到的。

\bigskip
\subsection*{读者对象}

本文稿可作为一部自成体系、面向研究者的参考书，适合
处理高维数据、希望熟练使用张量网络的研究人员与
从业者。所述内容易于入门：只要具备本科第一门线性代数课程层面的
实用知识，包括矩阵乘法、QR 分解与 SVD，
就足以阅读正文的大部分内容。

行文的组织方式同时服务于多个群体。对
\textit{数学家}而言，本文稿给出严格的定义、秩的刻画与
以线性代数理论为基础的算法分析。对
\textit{机器学习与数据科学的研究者}而言，它提供一套实用的
工具，用以参数化高维权重矩阵、经由结构化的张量分解计算梯度，
并以紧凑的形式建模复杂分布。对
\textit{物理学家与量子计算研究者}而言，它把张量网络的图示
约定与应用数学的语言沟通起来。 

本文稿中的所有张量网络图均以 \LaTeX\ 配合 \texttt{TikZ} 宏
排版，图示记法逐步引入，因此不假定读者
此前接触过张量网络图。



\section{张量网络基础}\label{chapter:basics}

随着阶数升高，张量的表示与运算都变得更加复杂。
张量网络（tensor network, TN）为处理这类高阶对象提供了高效的框架，使表示与分析同时得到简化。
张量网络的图示记法（graphical notation）让复杂张量运算的可视化与化简变得直观~\citep{orus2014practical,biamonte2017tensor}。
本节用张量网络的语言介绍张量的基本概念与常见的张量运算。

\subsection{什么是张量网络？}
顾名思义，张量网络就是把张量彼此连接而成的网络。更准确地说，张量网络是一类图：顶点表示\emph{张量}，边表示\emph{已收缩}的指标（contracted indices），即共享的模。张量网络的图示记法最早由~\cite{penrose1971applications}提出。张量用带腿（leg，也称边 edge）的图形表示，例如，
$
\begin{tikzpicture}[baseline=-0.5ex]
    \tikzset{tensor/.style = {minimum size = 0.5cm,shape = circle,thick,draw=black,inner sep = 0pt}, edge/.style = {thick,line width=.4mm},every loop/.style={}}
    \def\x{6}
    \node[tensor,fill=green!50!red!50!white] (T) at (\x,0) {$\Tt$};
    \draw[edge] (T) -- (\x-0.6,0);
    \draw[edge] (T) -- (\x+0.6,0);
    \draw[edge] (T) -- (\x-0.5,-0.3);
    \draw[dotted] (T) -- (\x,-0.6);
    \draw[dotted] (T) -- (\x+0.5,-0.3);
    \draw[dotted] (T) -- (\x+0.3,-0.5);
    \end{tikzpicture}.
$
我们用矩形、三角形、圆形等不同形状以及各种颜色来表示张量。形状与颜色主要起视觉提示的作用。
张量的\emph{阶}（order）指它的维度个数。例如，$N$ 阶张量 $\Tt\in\RR^{d_1\times\cdots\times d_N}$ 是由标量~$\Tt_{i_1,\dots,i_N}$ 构成的多维数组，由 $N$ 个指标 $i_1\in [d_1], i_2\in[d_2],\cdots, i_N \in [d_N]$ 索引。每个轴称为张量的一个\emph{模}（mode）：$N$ 阶张量有 $N$ 个模。
 
 张量把向量与矩阵自然推广为高维数组。阶数越高，这类数组的表示就越困难。\emph{张量网络}为表示和运算这些高阶对象提供了简单而直观的方式。借助张量网络的图示记法，可以更轻松地表达复杂的张量运算~\citep{biamonte2017tensor,orus2014practical}。


\paragraph{张量网络的节点}\label{def:tn-nodes}
在张量网络中，每个节点（node，或称顶点）表示一个张量，腿~（与之关联的边）的数目对应张量的阶：标量是没有边的节点，向量节点有一条边，矩阵节点有两条边，依此类推。例如，
\begin{center}
    \input{TN_figures/TN_nodes}
\end{center}
分别表示标量、向量、矩阵、三阶张量与四阶张量。
在本手册中，
\begin{itemize}
    \item 张量用带颜色的图形表示，称为节点，颜色与形状并无特定含义~（除非另有说明）。
    \item 标量用小写字母表示~（如 $a$），向量用粗小写字母表示~（如 $\vb$），矩阵用粗大写字母表示（如 $\Mb$），高阶张量则用
    粗花体大写字母表示~（如 $\Tt$）。
    \item 张量每个模的维数用一个灰色字母标在相应腿的上方，例如，
$
    \begin{tikzpicture}[baseline=-0.5ex]
    \tikzset{tensor/.style = {minimum size = 0.5cm,shape = circle,thick,draw=black,inner sep = 0pt}, edge/.style = {thick,line width=.4mm},every loop/.style={}}
    \def\y{0}
    \def\x{0}
    \node[tensor,fill=green!50!red!50!white] (A) at (\x,\y){\scalebox{0.85}{$\At$}};
    \draw[edge] (\x+0.25,\y) -- (\x+0.6,\y)  node [midway,above] {\scalebox{0.7}{\dimcolor{$d_2$}}};
    \draw[edge] (\x-0.25,\y) -- (\x-0.6,\y)  node [midway,above] {\scalebox{0.7}{\dimcolor{$d_1$}}};
    \draw[edge] (\x,\y-0.25) -- (\x,\y-0.7)  node [midway,left] {\scalebox{0.7}{\dimcolor{$d_3$}}};
    \end{tikzpicture}
    \ \ \ \in \RR^{d_1\times d_2\times d_3},
$
是尺寸为 $d_1\times d_2\times d_3$ 的三阶张量。
    
$
\At_{i,j,k} = 
    \begin{tikzpicture}[baseline=-0.5ex]
    \tikzset{tensor/.style = {minimum size = 0.5cm,shape = circle,thick,draw=black,inner sep = 0pt}, edge/.style = {thick,line width=.4mm},every loop/.style={}}
    \def\y{0}
    \def\x{0}
    \node[tensor,fill=green!50!red!50!white] (A) at (\x,\y){\scalebox{0.85}{$\At$}};
    \draw[edge] (\x+0.25,\y) -- (\x+0.6,\y)  node [pos=1.3] {\scalebox{0.7}{\indcolor{$j$}}};
    \draw[edge] (\x-0.25,\y) -- (\x-0.6,\y)  node [pos=1.3] {\scalebox{0.7}{\indcolor{$i$}}};
    \draw[edge] (\x,\y-0.25) -- (\x,\y-0.7)  node [pos=1.3] {\scalebox{0.7}{\indcolor{$k$}}};
    \end{tikzpicture}
$    
是三阶张量 $\At$ 的第 $(i,j,k)$ 个元素。
\end{itemize}
注意，大矩阵有时也可借助重排（reshaping）被视为高阶张量，例如，
\begin{align*}
    &\Tb
    =
    \begin{tikzpicture}[baseline=\baseline]
    \tikzset{tensor/.style = {minimum size = 0.5cm,shape = circle,thick,draw=black,inner sep = 0pt}, edge/.style = {thick,line width=.4mm},every loop/.style={}}
    \def\x{0}
    \def\y{0}
    \node[tensor,fill=green!50!blue!40!white] (Tb) at (\x,\y){$\Tb$};
    \draw[edge] (Tb) -- (\x+1.15,\y) node [midway,above] {\scalebox{0.7}{\dimcolor{$J_1J_2J_3J_4$}}};
    \draw[edge] (Tb) -- (\x-1.15,\y) node [midway,above] {\scalebox{0.7}{\dimcolor{$I_1I_2I_3I_4$}}};
    \end{tikzpicture}\in\RR^{I_1I_2I_3I_4\times J_1J_2J_3J_4}
    \equiv
    \Tt = 
    \begin{tikzpicture}[baseline=-0.5ex]
    \tikzset{tensor/.style = {minimum width = 1.8cm,minimum height = 0.4cm,shape = rounded rectangle,thick,draw=black,inner sep = 0pt}, edge/.style = {thick,line width=.4mm},every loop/.style={}}
    \def\x{0}
    \def\y{0}
    \draw[edge] (\x-0.2,\y-0.15) -- (\x-0.2,\y-0.6)node [midway,left=-0.05cm] {\scalebox{0.7}{\dimcolor{$I_2$}}};
    \draw[edge] (\x-0.6,\y-0.15) -- (\x-0.6,\y-0.6)node [midway,left=-0.05cm] {\scalebox{0.7}{\dimcolor{$I_1$}}};
    \draw[edge] (\x+0.2,\y-0.15) -- (\x+0.2,\y-0.6)node [midway,left=-0.05cm] {\scalebox{0.7}{\dimcolor{$I_3$}}};
    \draw[edge] (\x+0.6,\y-0.15) -- (\x+0.6,\y-0.6)node [midway,left=-0.05cm] {\scalebox{0.7}{\dimcolor{$I_4$}}};
    \draw[edge] (\x-0.2,\y+0.15) -- (\x-0.2,\y+0.6)node [midway,left=-0.05cm] {\scalebox{0.7}{\dimcolor{$J_2$}}};
    \draw[edge] (\x-0.6,\y+0.15) -- (\x-0.6,\y+0.6)node [midway,left=-0.05cm] {\scalebox{0.7}{\dimcolor{$J_1$}}};
    \draw[edge] (\x+0.2,\y+0.15) -- (\x+0.2,\y+0.6)node [midway,left=-0.05cm] {\scalebox{0.7}{\dimcolor{$J_3$}}};
    \draw[edge] (\x+0.6,\y+0.15) -- (\x+0.6,\y+0.6)node [midway,left=-0.05cm] {\scalebox{0.7}{\dimcolor{$J_4$}}};
    \node[tensor,fill=yellow] (Bt) at (\x,\y){$\Tt$};
    \end{tikzpicture}
    \in\RR^{I_1\times I_2\times I_3\times I_4\times J_1 \times J_2\times J_3\times J_4}.
\end{align*}
\paragraph{张量网络的边}\label{subsec:contraction} 在张量网络图中，腿分为两类：已收缩的腿~（连接两个张量的腿）与未收缩的腿，后者也称自由腿~（free legs），一端悬空~（即不与任何其他张量相连的腿）。如前所述，未收缩的腿对应自由指标（free index）：自由腿的数目给出张量的阶——标量没有自由腿，向量有一条，矩阵有两条，高阶张量则有三条或更多。已收缩的腿表示收缩（contraction）：
张量可沿尺寸相同的腿相连，这表示对所连接的模求和~（收缩）。
求和与收缩在本文中通用。最常见的收缩运算为\textit{矩阵乘法}。对两个矩阵 $\Ab\in\RR^{d_1\times R}, \Bb\in\RR^{R\times d_2}$，其矩阵乘积相当于 $\Ab$ 的第二模与 $\Bb$ 的第一模之间的收缩：

\begin{align}\label{eq:matrix-product}
(\Ab\Bb)_{ij} 
=
\begin{tikzpicture}[baseline=\baseline]
   \tikzset{tensor/.style = {minimum size = 0.5cm,shape = circle,thick,draw=black,inner sep = 0pt}, edge/.style = {thick,line width=.4mm},every loop/.style={}}
    \def\y{0}
    \def\x{0}
    \node[tensor,fill=red!50!white] (A) at (\x,\y){\scalebox{0.85}{$\Ab$}};
    \node[tensor,fill=blue!50!white] (B) at (\x+1,\y){\scalebox{0.85}{$\Bb$}};
    \draw[edge] (A) -- (B) node [above,midway] {\scalebox{0.7}{\dimcolor{$R$}}};
    \draw[edge] (\x-0.25,\y) -- (\x-0.75,\y)  node [pos=1.3] {\scalebox{0.7} {\indcolor{$i$}}};
    \draw[edge] (\x+1.25,\y) -- (\x+1.75,\y)  node [pos=1.3] {\scalebox{0.7}{\indcolor{$j$}}};
    \end{tikzpicture}
= \sum_{r=1}^R\Ab_{ir}\Bb_{rj},\ \ \ \text{for}\ \ i\in[d_1],j\in[d_2]\ .
\end{align}
由于该等式对一切指标 $i$ 与 $j$ 都成立，可以直接简写为
\begin{align}\label{eq:matrix-product_noidx}
\Ab\Bb 
=
\begin{tikzpicture}[baseline=\baseline]
   \tikzset{tensor/.style = {minimum size = 0.5cm,shape = circle,thick,draw=black,inner sep = 0pt}, edge/.style = {thick,line width=.4mm},every loop/.style={}}
    \def\y{0}
    \def\x{0}
    \node[tensor,fill=red!50!white] (A) at (\x,\y){\scalebox{0.85}{$\Ab$}};
    \node[tensor,fill=blue!50!white] (B) at (\x+1,\y){\scalebox{0.85}{$\Bb$}};
    \draw[edge] (A) -- (B) node [above,midway] {\scalebox{0.7}{\dimcolor{$R$}}};
    \draw[edge] (\x-0.25,\y) -- (\x-0.75,\y)  node [midway, above] {\scalebox{0.7}{\dimcolor{$d_1$}}};
    \draw[edge] (\x+1.25,\y) -- (\x+1.75,\y)  node [midway,above] {\scalebox{0.7}{\dimcolor{$d_2$}}};
    \end{tikzpicture}
 \ \in \mathbb R^{d_1\times d_2}.
\end{align}

 在上图中可以看到，尺寸同为 $R$ 的两条腿相互连接。这与矩阵乘法相一致：矩阵乘法就是对维度相同的指标求和。最后，由于所得的图有两条自由腿，它表示一个矩阵。这里的最终对象是一个矩阵，而矩阵可以用单个节点表示，这说明在张量网络中节点可以被合并：
$$
\Ab\Bb =\Mb \Leftrightarrow 
\begin{tikzpicture}[baseline=\baseline]
   \tikzset{tensor/.style = {minimum size = 0.5cm,shape = circle,thick,draw=black,inner sep = 0pt}, edge/.style = {thick,line width=.4mm},every loop/.style={}}
    \def\y{0}
    \def\x{0}
    \node[tensor,fill=red!50!white] (A) at (\x,\y){\scalebox{0.85}{$\Ab$}};
    \node[tensor,fill=blue!50!white] (B) at (\x+1,\y){\scalebox{0.85}{$\Bb$}};
    \draw[edge] (A) -- (B) node [above,midway] {\scalebox{0.7}{\dimcolor{$R$}}};
    \draw[edge] (\x-0.25,\y) -- (\x-0.75,\y)  node [midway, above] {\scalebox{0.7}{\dimcolor{$d_1$}}};
    \draw[edge] (\x+1.25,\y) -- (\x+1.75,\y)  node [midway,above] {\scalebox{0.7}{\dimcolor{$d_2$}}};
    \node[draw=none] () at (\x+2,\y) {\scalebox{1}{\textcolor{black}{$=$}}};
    \node[tensor,fill=blue!30!red!20!white] (M) at (\x+3,\y){\scalebox{0.85}{$\Mb$}};
    \draw[edge] (M) -- (\x+3.75,\y)  node [midway,above] {\scalebox{0.7}{\dimcolor{$d_2$}}};
    \draw[edge] (M) -- (\x+2.25,\y)  node [midway,above] {\scalebox{0.7}{\dimcolor{$d_1$}}};
    \end{tikzpicture}.
$$
矩阵之间的收缩可以直接推广到矩阵与向量、张量与矩阵、张量与矩阵及向量的收缩等情形。下面给出若干张量网络例子及其对应的数学表达式：
\begin{align}\label{eq:matvec}
\Ab\in\RR^{d_1\times R},\vb\in\RR^{R} \ : \ \
\begin{tikzpicture}[baseline=\baseline]
   \tikzset{tensor/.style = {minimum size = 0.5cm,shape = circle,thick,draw=black,inner sep = 0pt}, edge/.style = {thick,line width=.4mm},every loop/.style={}}
    \def\y{0}
    \def\x{0}
    \node[tensor,fill=red!50!white] (A) at (\x,\y){\scalebox{0.85}{$\Ab$}};
    \node[tensor,fill=blue!10!white] (v) at (\x+1,\y){\scalebox{0.85}{$\vb$}};
    \draw[edge] (A) -- (v) node [above,midway] {\scalebox{0.7}{\dimcolor{$R$}}};
    \draw[edge] (\x-0.25,\y) -- (\x-0.75,\y)  node [pos=1.3]{\scalebox{0.7}{\indcolor{$i$}}};
    \end{tikzpicture}
=\sum_{r=1}^R\Ab_{ir}\vb_r,
\end{align}
\begin{align}\label{eq:tenmat}
\At\in\RR^{d_1\times R\times d_2},\Ab\in\RR^{R\times d_3}\ :\ \
\begin{tikzpicture}[baseline = \baseline]
    \tikzset{tensor/.style = {minimum size = 0.5cm,shape = circle,thick,draw=black,inner sep = 0pt}, edge/.style = {thick,line width=.4mm},every loop/.style={}}
    \def\y{0}
    \def\x{0}
    \node[tensor,fill=green!50!red!50!white] (At) at (\x,\y){\scalebox{0.85}{$\At$}};
    \node[tensor,fill=red!50!white] (A) at (\x+1,\y){\scalebox{0.85}{$\Ab$}};
    \draw[edge] (At) -- (A) node [above,midway] {\scalebox{0.7}{\dimcolor{$R$}}};
    \draw[edge] (\x-0.25,\y) -- (\x-0.75,\y)  node [pos=1.3] {\scalebox{0.7}{\indcolor{$i$}}};
    \draw[edge] (\x,\y-0.25) -- (\x,\y-0.75)  node [pos=1.3] {\scalebox{0.7}{\indcolor{$j$}}};
    \draw[edge] (\x+1.25,\y) -- (\x+1.75,\y)  node [pos=1.3] {\scalebox{0.7}{\indcolor{$k$}}};
\end{tikzpicture} 
= \sum_{r=1}^R\At_{ijr}\Ab_{rk},
\end{align}
\begin{align}\label{eq:tenmatvec}
\At\in\RR^{d_1\times R\times d_2},\Ab\in\RR^{R\times d_3}, \vb\in\RR^{d_3}&:
   \begin{tikzpicture}[baseline = \baseline]
    \tikzset{tensor/.style = {minimum size = 0.5cm,shape = circle,thick,draw=black,inner sep = 0pt}, edge/.style = {thick,line width=.4mm},every loop/.style={}}
    \def\y{0}
    \def\x{0}
    \node[tensor,fill=green!50!red!50!white] (At) at (\x,\y){\scalebox{0.85}{$\At$}};
    \node[tensor,fill=red!50!white] (A) at (\x+1,\y){\scalebox{0.85}{$\Ab$}};
    \node[tensor,fill=blue!40!red!60!white] (a) at (\x+2,\y){\scalebox{0.85}{$\vb$}};
    \draw[edge] (At) -- (A) node [above,midway] {\scalebox{0.7}{\dimcolor{$R$}}};
    \draw[edge] (A) -- (a) node [above,midway] {\scalebox{0.7}{\dimcolor{$d_3$}}};
    \draw[edge] (\x-0.25,\y) -- (\x-0.75,\y)  node [pos=1.3] {\scalebox{0.7}{\indcolor{$i$}}};
    \draw[edge] (\x,\y-0.25) -- (\x,\y-0.75)  node [pos=1.3] {\scalebox{0.7}{\indcolor{$j$}}};
    \end{tikzpicture} 
= \sum_{r=1}^R\sum_{k=1}^{d_3}\At_{ijr}\Ab_{rk}\ab_{k}.
\end{align}
由上述各图可见，两个节点之间的边表示求和，同时也表明两个张量在相应的模上共享尺寸相同的维度。
如前所述，张量网络中自由腿的数目对应其所表示张量的阶。对前面的例子，我们有

$$
\begin{tikzpicture}[baseline=\baseline]
   \tikzset{tensor/.style = {minimum size = 0.5cm,shape = circle,thick,draw=black,inner sep = 0pt}, edge/.style = {thick,line width=.4mm},every loop/.style={}}
    \def\y{0}
    \def\x{0}
    \node[tensor,fill=red!50!white] (A) at (\x,\y){\scalebox{0.85}{$\Ab$}};
    \node[tensor,fill=blue!10!white] (v) at (\x+1,\y){\scalebox{0.85}{$\vb$}};
    \draw[edge] (A) -- (v) node [above,midway] {\scalebox{0.7}{\dimcolor{$R$}}};
    \draw[edge] (\x-0.25,\y) -- (\x-0.75,\y)  node [midway, above] {\scalebox{0.7}{\dimcolor{$d_1$}}};
    \end{tikzpicture} 
   \in \mathbb R^{d_1},\ \ 
\begin{tikzpicture}[baseline = \baseline]
    \tikzset{tensor/.style = {minimum size = 0.5cm,shape = circle,thick,draw=black,inner sep = 0pt}, edge/.style = {thick,line width=.4mm},every loop/.style={}}
    \def\y{0}
    \def\x{0}
    \node[tensor,fill=green!50!red!50!white] (At) at (\x,\y){\scalebox{0.85}{$\At$}};
    \node[tensor,fill=red!50!white] (A) at (\x+1,\y){\scalebox{0.85}{$\Ab$}};
    \draw[edge] (At) -- (A) node [above,midway] {\scalebox{0.7}{\dimcolor{$R$}}};
    \draw[edge] (\x-0.25,\y) -- (\x-0.75,\y)  node [midway,above] {\scalebox{0.7}{\dimcolor{$d_1$}}};
    \draw[edge] (\x,\y-0.25) -- (\x,\y-0.75)  node [midway,left] {\scalebox{0.7}{\dimcolor{$d_2$}}};
    \draw[edge] (\x+1.25,\y) -- (\x+1.75,\y)  node [midway,above] {\scalebox{0.7}{\dimcolor{$d_3$}}};
\end{tikzpicture} 
\in \mathbb R^{d_1\times d_2 \times d_3},\ \ \text{and}\ \ 
    \begin{tikzpicture}[baseline = \baseline]
    \tikzset{tensor/.style = {minimum size = 0.5cm,shape = circle,thick,draw=black,inner sep = 0pt}, edge/.style = {thick,line width=.4mm},every loop/.style={}}
    \def\y{0}
    \def\x{0}
    \node[tensor,fill=green!50!red!50!white] (At) at (\x,\y){\scalebox{0.85}{$\At$}};
    \node[tensor,fill=red!50!white] (A) at (\x+1,\y){\scalebox{0.85}{$\Ab$}};
    \node[tensor,fill=blue!40!red!60!white] (a) at (\x+2,\y){\scalebox{0.85}{$\vb$}};
    \draw[edge] (At) -- (A) node [above,midway] {\scalebox{0.7}{\dimcolor{$R$}}};
    \draw[edge] (A) -- (a) node [above,midway] {\scalebox{0.7}{\dimcolor{$d_3$}}};
    \draw[edge] (\x-0.25,\y) -- (\x-0.75,\y)  node [midway,above] {\scalebox{0.7}{\dimcolor{$d_1$}}};
    \draw[edge] (\x,\y-0.25) -- (\x,\y-0.75)  node [midway,left] {\scalebox{0.7}{\dimcolor{$d_2$}}};
    \end{tikzpicture} 
\in \mathbb R^{d_1\times d_2}.
$$
 
\paragraph{从数学表达式到张量网络，再回到数学表达式} 张量网络用起来很简单，因为表示腿和节点并没有严格的规则。
在张量网络图中，节点与边可以在平面内任意摆放，重要的只是图的结构。
例如，把矩阵乘法翻译成张量网络图时，关键是保证相连腿的维度一致。例如对 $\Ab\in\RR^{d_1\times R}, \Bb\in\RR^{R\times d_2}$，下面各图表示的都是\emph{同一个}矩阵乘法：
\begin{align*} 
\Ab\Bb
=
\begin{tikzpicture}[baseline=\baseline]
   \tikzset{tensor/.style = {minimum size = 0.5cm,shape = circle,thick,draw=black,inner sep = 0pt}, edge/.style = {thick,line width=.4mm},every loop/.style={}}
    \def\y{0}
    \def\x{0}
    \node[tensor,fill=red!50!white] (A) at (\x,\y){\scalebox{0.85}{$\Ab$}};
    \node[tensor,fill=blue!50!white] (B) at (\x+1,\y){\scalebox{0.85}{$\Bb$}};
    \draw[edge] (A) -- (B) node [midway,above] {\scalebox{0.7}{\dimcolor{$R$}}};;
    \draw[edge] (\x-0.25,\y) -- (\x-0.75,\y) node [midway,above] {\scalebox{0.7}{\dimcolor{$d_1$}}};;
    \draw[edge] (\x+1.25,\y) -- (\x+1.75,\y) node [midway,above] {\scalebox{0.7}{\dimcolor{$d_2$}}};;
    \node[draw=none] () at (\x+1.95,\y) {$=$};
    \node[tensor,fill=red!50!white] (A1) at (\x+2.85,\y){\scalebox{0.85}{$\Ab$}};
    \node[tensor,fill=blue!50!white] (B1) at (\x+3.85,\y){\scalebox{0.85}{$\Bb$}};
    \draw[edge] (A1) -- (B1) node [midway,above] {\scalebox{0.7}{\dimcolor{$R$}}};;
    \draw[edge] (A1) to [out=60,in=60, distance=0.5cm] (\x+2.2,\y);
    \draw[edge] (B1) to [out=60,in=60, distance=0.5cm] (\x+4.4,\y-0.5);
    \node[draw=none] () at (\x+2.4,\y+0.65) {\scalebox{0.7}{\dimcolor{$d_1$}}};
    \node[draw=none] () at (\x+4,\y+0.65) {\scalebox{0.7}{\dimcolor{$d_2$}}};
    \node[draw=none] () at (\x+4.7,\y) {$=$};
    \node[tensor,fill=blue!50!white] (A2) at (\x+5.3,\y){\scalebox{0.85}{$\Bb$}};
    \node[tensor,fill=red!50!white] (B2) at (\x+6.3,\y){\scalebox{0.85}{$\Ab$}};
    \draw[edge] (A2) -- (B2) node [midway,above] {\scalebox{0.7}{\dimcolor{$R$}}};;
    \draw[edge] (A2) -- (\x+5.3,\y+0.75) node [midway,left] {\scalebox{0.7}{\dimcolor{$d_2$}}};
     \draw[edge] (B2) -- (\x+6.3,\y-0.75) node [midway,right] {\scalebox{0.7}{\dimcolor{$d_1$}}};   
     \node[draw=none] () at (\x+7.1,\y) {$=$};
    \node[tensor,fill=red!50!white] (A3) at (\x+7.6,\y+0.5){\scalebox{0.85}{$\Ab$}};
    \node[tensor,fill=blue!50!white] (B3) at (\x+7.6,\y-0.5){\scalebox{0.85}{$\Bb$}};
    \draw[edge] (A3) -- (B3) node [midway,right] {\scalebox{0.7}{\dimcolor{$R$}}};
    \draw[edge] (A3) -- (\x+7.6,\y+1.25) node [midway,left] {\scalebox{0.7}{\dimcolor{$d_1$}}};
    \draw[edge] (B3) -- (\x+7.6,\y-1.25) node [midway,left] {\scalebox{0.7}{\dimcolor{$d_2$}}};
    \end{tikzpicture}.
\end{align*}
注意，由于 $\Ab\in\RR^{d_1\times R}$ 与 $\Bb\in\RR^{R\times d_2}$ 的尺寸限制，二者之间只有一种收缩方式；因此，即使省去边上标注的维数，上面的图也不会有歧义。 

反之，把张量网络图翻译成数学表达式也可能得到多种解释。例如下面这个矩阵乘法图
$
\begin{tikzpicture}[baseline=\baseline]
   \tikzset{tensor/.style = {minimum size = 0.5cm,shape = circle,thick,draw=black,inner sep = 0pt}, edge/.style = {thick,line width=.4mm},every loop/.style={}}
    \def\y{0}
    \def\x{0}
    \node[tensor,fill=red!50!white] (A) at (\x,\y){\scalebox{0.85}{$\Ab$}};
    \node[tensor,fill=blue!50!white] (B) at (\x+1,\y){\scalebox{0.85}{$\Bb$}};
    \draw[edge] (A) -- (B) node [midway,above] {\scalebox{0.7}{\dimcolor{$R$}}};
    \draw[edge] (\x-0.25,\y) -- (\x-0.75,\y) node [midway,above] {\scalebox{0.7}{\dimcolor{$d_1$}}};
    \draw[edge] (\x+1.25,\y) -- (\x+1.75,\y) node [midway,above] {\scalebox{0.7}{\dimcolor{$d_2$}}};
\end{tikzpicture}
$
依据我们为 $\Ab$、$\Bb$ 与结果矩阵选定的形状不同，它有不同的解读：
\begin{itemize}
    \item 若 $\Ab\in\RR^{d_1\times R}, \Bb\in\RR^{R\times d_2}$ 且结果的大小为 $d_1 \times d_2$，则该图表示乘积 $\Ab\Bb$，
    \item 若 $\Ab\in\RR^{R\times d_1}, \Bb\in\RR^{R\times d_2}$ 且结果的大小为 $d_1 \times d_2$，则该图表示乘积 $\Ab^\ts\Bb$，
    \item 若 $\Ab\in\RR^{R\times d_1}, \Bb\in\RR^{R\times d_2}$ 且结果的大小为 $d_2 \times d_1$，则该图表示乘积 $\Bb^\ts\Ab$，
    \item ...
\end{itemize}
因此，把张量网络图翻译成数学表达式时要多加留心。但这也正是张量网络如此有用之处！当我们操作图示而非数学表达式时，就不必在意例如转置这样的问题。

\paragraph{收缩同一节点的两条腿。 } 如前所述，任意两条维度相同的腿都可以连接起来，构成张量网络中的一条边，即使它们属于同一个节点。举例来说，考虑五阶张量 $\Tt\in\mathbb R^{d_1\times d_2 \times d_3 \times d_2 \times d_4}$。它的第 2 个与第 4 个模维度相同，因此可以把二者收缩在一起。所得的张量网络有三条悬空腿，因而表示一个三阶张量： 
\begin{align*}
&\begin{tikzpicture}[baseline=-0.5ex]
    \tikzset{tensor/.style = {minimum width = 1.8cm,minimum height = 0.4cm,shape = rounded rectangle,thick,draw=black,inner sep = 0pt}, edge/.style = {thick,line width=.4mm},every loop/.style={}}
    \def\x{0}
    \def\y{0}
    \draw[edge] (\x-0.4,\y-0.15) -- (\x-0.4,\y-0.5)node [near end,left=-0.05cm] {\scalebox{0.7}{\dimcolor{$d_2$}}};
    \draw[edge] (\x-0.8,\y-0.15) -- (\x-0.8,\y-0.5)node [near end,left=-0.05cm] {\scalebox{0.7}{\dimcolor{$d_1$}}};
    \draw[edge] (\x,\y-0.15) -- (\x,\y-0.5)node [near end,left=-0.05cm] {\scalebox{0.7}{\dimcolor{$d_3$}}};
    \draw[edge] (\x+0.4,\y-0.15) -- (\x+0.4,\y-0.5)node [near end,left=-0.05cm] {\scalebox{0.7}{\dimcolor{$d_2$}}};
    \draw[edge] (\x+0.8,\y-0.15) -- (\x+0.8,\y-0.5)node [near end,left=-0.05cm] {\scalebox{0.7}{\dimcolor{$d_4$}}};
    \node[tensor,fill=green!30!white] (Bt) at (\x,\y){$\Tt$};
    \end{tikzpicture}\in \mathbb R^{d_1\times d_2\times d_3 \times d_2 \times d_4},\ \ \ \ 
  \begin{tikzpicture}[baseline=-0.5ex]
    \tikzset{tensor/.style = {minimum width = 1.8cm,minimum height = 0.4cm,shape = rounded rectangle,thick,draw=black,inner sep = 0pt}, edge/.style = {thick,line width=.4mm},every loop/.style={}}
    \def\x{0}
    \def\y{0}
    \draw[edge] (\x-0.4,\y-0.15) -- (\x-0.4,\y-0.5);
    \draw[edge] (\x-0.8,\y-0.15) -- (\x-0.8,\y-0.5)node [pos=1.3] {\scalebox{0.7}{\indcolor{$i_1$}}};
    \draw[edge] (\x,\y-0.15) -- (\x,\y-0.5)node [pos=1.3] {\scalebox{0.7}{\indcolor{$i_3$}}};
    \draw[edge] (\x+0.4,\y-0.15) -- (\x+0.4,\y-0.5);
    \draw[edge] (\x-0.4,\y-0.5) to [out=270,in=270] (\x+0.4,\y-0.5) node [midway,below=0.7cm] {\scalebox{0.7}{\dimcolor{$d_2$}}};
    \draw[edge] (\x+0.8,\y-0.15) -- (\x+0.8,\y-0.5)node[pos=1.3] {\scalebox{0.7}{\indcolor{$i_4$}}};
    \node[tensor,fill=green!30!white] (Bt) at (\x,\y){$\Tt$};
    \end{tikzpicture} = \sum_{i_2=1}^{d_2} \Tt_{i_1,i_2,i_3,i_2,i_4},   %
    \text{ and }\ \ \
    \begin{tikzpicture}[baseline=-0.5ex]
    \tikzset{tensor/.style = {minimum width = 1.8cm,minimum height = 0.4cm,shape = rounded rectangle,thick,draw=black,inner sep = 0pt}, edge/.style = {thick,line width=.4mm},every loop/.style={}}
    \def\x{0}
    \def\y{0}
    \draw[edge] (\x-0.4,\y-0.15) -- (\x-0.4,\y-0.5);
    \draw[edge] (\x-0.8,\y-0.15) -- (\x-0.8,\y-0.5)node [midway,left=-0.05cm] {\scalebox{0.7}{\dimcolor{$d_1$}}};
    \draw[edge] (\x,\y-0.15) -- (\x,\y-0.5)node [midway,left=-0.05cm] {\scalebox{0.7}{\dimcolor{$d_3$}}};
    \draw[edge] (\x+0.4,\y-0.15) -- (\x+0.4,\y-0.5);
    \draw[edge] (\x-0.4,\y-0.5) to [out=270,in=270] (\x+0.4,\y-0.5) node [midway,below=0.7cm] {\scalebox{0.7}{\dimcolor{$d_2$}}};
    \draw[edge] (\x+0.8,\y-0.15) -- (\x+0.8,\y-0.5)node [midway,left=-0.05cm] {\scalebox{0.7}{\dimcolor{$d_4$}}};
    \node[tensor,fill=green!30!white] (Bt) at (\x,\y){$\Tt$};
    \end{tikzpicture} \in \mathbb R^{d_1\times d_3 \times d_4}.
\end{align*}
收缩一个 $N$ 阶张量的两条腿后，所得张量的阶数为 $N-2$。这与后文将要介绍的\emph{偏迹}（partial trace）概念密切相关。
\subsection{内积、外积、迹与范数}\label{sec:inner-products}
本节借助张量网络说明，内积（inner product）、外积（outer product）、迹（trace）与范数（norm）这些经典概念如何推广到高阶张量。

\paragraph{内积} 与向量之间经典的内积~（欧几里得内积）类似，两个 $N$ 阶张量 $\St,\Tt\in\RR^{d_1\times\dots\times d_N}$ 的内积是它们对应元素乘积之和：
$$ \langle \St,\Tt \rangle = \sum_{i_1=1}^{d_1}\dots\sum_{i_N=1}^{d_N}\St_{i_1\dots i_N}\Tt_{i_1\dots i_N}.$$ 在张量网络图中，对所有维度求和就是把两个张量的所有腿都连接起来。这样得到的张量网络没有自由腿，表示一个标量。例如，对向量 $\ab\in\RR^{d\times 1}, \bb\in\RR^{d\times 1}$ 有
\begin{align}\label{eq:inner-product}
\langle\ab,\bb\rangle = \sum_{i=1}^d\ab_i\bb_i
=
    \begin{tikzpicture}[baseline=\baseline]
   \tikzset{tensor/.style = {minimum size = 0.5cm,shape = circle,thick,draw=black,inner sep = 0pt}, edge/.style = {thick,line width=.4mm},every loop/.style={}}
    \def\y{0}
    \def\x{0}
    \node[tensor,fill=red!30!white] (a) at (\x,\y){\scalebox{0.85}{$\ab$}};
    \node[tensor,fill=blue!30!white] (b) at (\x+1,\y){\scalebox{0.85}{$\bb$}};
    \draw[edge] (A) -- (B) node [midway,above]{\scalebox{0.7}{\dimcolor{$d$}}};
\end{tikzpicture},
\end{align}
对两个三阶张量 $\St,\Tt\in\RR^{d_1\times d_2\times d_3}$ 则有
\begin{align*}
\langle\St,\Tt\rangle = \sum_{i_1=1}^{d_1}\sum_{i_2=1}^{d_2}\sum_{i_3=1}^{d_3}\St_{i_1i_2i_3}\Tt_{i_1i_2i_3} =
\begin{tikzpicture}[baseline=\baseline]
    \tikzset{tensor/.style = {minimum size = 0.5cm,shape = circle,thick,draw=black,inner sep = 0pt}, edge/.style = {thick,line width=.4mm},every loop/.style={}}
    \def\y{0}
    \def\x{0}
    \node[tensor,fill=green!50!red!50!white] (St) at (\x,\y){\scalebox{0.85}{$\St$}};
    \node[tensor,fill=green!50!red!50!white] (Tt) at (\x+1,\y){\scalebox{0.85}{$\Tt$}};
    \draw[edge] (St) -- (Tt);
    \draw[edge] (St) to [out=45,in=135] (Tt);
    \draw[edge] (St) to [out=-45,in=-135] (Tt);
    \node[draw=none] () at (\x+0.5,\y+0.5) {\scalebox{0.7}{\dimcolor{$d_1$}}};
    \node[draw=none] () at (\x+0.5,\y+0.15) {\scalebox{0.7}{\dimcolor{$d_2$}}};
    \node[draw=none] () at (\x+0.5,\y-0.5) {\scalebox{0.7}{\dimcolor{$d_3$}}};
\end{tikzpicture}  
.
\end{align*}
本书只讨论实值张量，但张量网络这套形式语言稍加改造即可用于复值张量。
例如，若这两个张量取复值，就需对第二个张量的元素取共轭，这在张量网络中可以用节点 $\Tt^*$~（逐元素共轭）来代替 $\Tt$ 表示。
\paragraph{外积}\label{subsection:outer-product}
回忆一下，两个向量 $\ub\in\RR^{m},\vb\in\RR^n$ 的外积是 $m\times n$ 的秩一（rank-one）矩阵 $\ub\vb^\top$。
这一运算可推广到任意多个张量。例如，$N$ 个向量 $\ab_1\in\RR^{d_1},\cdots,\ab_N\in\RR^{d_N}$ 的外积是一个 $N$ 阶张量，其定义为
\begin{align*}
    (\ab_1 \outprod \ab_2 \outprod \cdots \outprod \ab_N)_{i_1,i_2,\cdots,i_n}
    &= \\
    &\hspace{-1cm}(\ab_1)_{i_1} (\ab_2)_{i_2}  \cdots (\ab_N)_{i_N}\ \text{for all } i_1 \in [d_1],\cdots, i_N\in[d_N].
\end{align*}
这样的张量称为\textbf{秩一}张量。 %
下图展示了 $N$ 个向量的外积： 
$$
\ab_1\outprod\ab_2 \outprod\cdots\outprod \ab_N=
\begin{tikzpicture}[baseline=\baseline]
    \tikzset{tensor/.style = {minimum size = 0.5cm,shape = circle,thick,draw=black,inner sep = 0pt}, edge/.style   = {thick,line width=.4mm}}
    \def\y{0}
     \def\x{0}
     \node[tensor,fill=blue!10!white] (a1) at (\x,\y){$\ab_1$};
    \node[tensor,fill=blue!10!white](a2) at (\x+1,\y){$\ab_2$};
  
    \node[tensor,fill=blue!10!white](an) at (\x+3,\y){$\ab_N$};
    \draw[edge] (a1) -- (\x,\y-0.7) node [midway, left] {\scalebox{0.7}{\dimcolor{$d_1$}}};
    \draw[edge] (a2) -- (\x+1,\y-0.7) node [midway, left] {\scalebox{0.7}{\dimcolor{$d_2$}}};
    \draw[edge] (an) -- (\x+3,\y-0.7)  node [midway, left] {\scalebox{0.7}{\dimcolor{$d_N$}}};
    \node (eq) at (\x+2,\y) {$\cdots$};
\end{tikzpicture}
\in\RR^{d_1\times\cdots\times d_N}.
$$
特别地，当 $N=2$ 时有  $\ab_1 \outprod \ab_2 = \ab_1\ab_2^\ts\in\RR^{d_1\times d_2}$。从外积的张量网络图中可以看到，节点之间没有共享边，这正对应于外积中不发生收缩（求和）。
外积的概念可以自然地推广到高阶张量。
设 $\At\in\RR^{m_1\times\dots\times m_p}$、$\Bt\in\RR^{n_1\times\dots\times n_q}$。$\At$ 与 $\Bt$ 的外积是一个 $p+q$ 阶张量，定义为 
$$ (\At\outprod\Bt)_{i_1,i_2,\cdots,i_p,j_1,j_2,\cdots,j_q} = \At_{i_1,i_2,\cdots,i_p} 
\Bt_{j_1,j_2,\cdots,j_q}. $$
与向量的情形一样，外积的张量网络只需把相应的节点并排摆放、不添加任何边即可得到：
$$
\At\outprod\Bt = 
   \begin{tikzpicture}[baseline=-0.5ex]
    \tikzset{tensor/.style = {minimum size = 0.5cm,shape = circle,thick,draw=black,inner sep = 0pt}, edge/.style = {thick,line width=.4mm},every loop/.style={}}
    \def\x{6}
    \node[tensor,fill=green!50!red!50!white] (A) at (\x,0) {$\At$};
    \draw[edge] (A) -- (\x-0.75,0) node [midway,above] {\scalebox{0.7}{\dimcolor{$m_1$}}};
    \draw[edge] (A) -- (\x+0.75,0) node [midway,above] {\scalebox{0.7}{\dimcolor{$m_p$}}};;
    \draw[edge] (A) -- (\x-0.75,-0.3) node [midway, below] {\scalebox{0.7}{\dimcolor{$m_2$}}};;
    \draw[dotted] (A) -- (\x,-0.6);
    \draw[dotted] (A) -- (\x+0.5,-0.5);
    \draw[dotted] (A) -- (\x+0.3,-0.6);
    \end{tikzpicture}
\hspace{0.3cm}
\begin{tikzpicture}[baseline=-0.5ex]
    \tikzset{tensor/.style = {minimum size = 0.5cm,shape = circle,thick,draw=black,inner sep = 0pt}, edge/.style = {thick,line width=.4mm},every loop/.style={}}
    \def\x{6}
    \def\y{0}
    \node[tensor,fill=green!70!red!20!white] (B) at (\x+1,\y) {$\Bt$};
    \draw[edge] (B) -- (\x+0.25,0) node [midway,above] {\scalebox{0.7}{\dimcolor{$n_1$}}};
    \draw[edge] (B) -- (\x+1.75,0) node [midway,above] {\scalebox{0.7}{\dimcolor{$n_q$}}};;
    \draw[edge] (B) -- (\x+0.25,-0.3) node [midway, below] {\scalebox{0.7}{\dimcolor{$n_2$}}};
    \draw[dotted] (B) -- (\x+1,-0.6);
    \draw[dotted] (B) -- (\x+1.5,-0.5);
    \draw[dotted] (B) -- (\x+1.3,-0.6);
\end{tikzpicture}\in\RR^{m_1\times\dots\times m_p\times n_1\times\dots\times n_q}.
$$
更一般地，对任意多个、阶数任意的张量，其外积的张量网络图只需把它们并排摆放即可得到。例如对  $\Tt\in\RR^{d_1\times d_2\times d_3 \times d_4}$、$\Ab\in\RR^{m\times n}$ 与 $\vb\in\RR^{p}$
$$
\Ab\outprod \Tt\outprod \vb
=
\begin{tikzpicture}[baseline=-0.5ex]
    \tikzset{tensor/.style = {minimum size = 0.5cm,shape = circle,thick,draw=black,inner sep = 0pt}, edge/.style = {thick,line width=.4mm},every loop/.style={}}
    \def\x{0}
    \def\y{0}
    \node[tensor,fill=blue!20!red!40!white] (A) at (\x,\y) {$\Ab$};
    \draw[edge] (A) -- (\x+0.75,\y)node [midway, above] {\scalebox{0.7}{\dimcolor{$n$}}};
    \draw[edge] (A) -- (\x-0.75,\y)node [midway, above] {\scalebox{0.7}{\dimcolor{$m$}}};
    \node[tensor,fill=red!20!blue!30!white] (T) at (\x+1.5,\y) {$\Tt$};
    \node[tensor,fill=red!70!blue!10!white] (v) at (\x+2.5,\y) {$\vb$};
    \draw[edge](T) -- (\x+2,\y-0.5);
    \node[draw=none] () at (\x+2,\y-0.25) {\scalebox{0.7}{\dimcolor{$d_3$}}};
    \draw[edge](T) -- (\x+2,\y+0.5);
    \node[draw=none] () at (\x+2,\y+0.25) {\scalebox{0.7}{\dimcolor{$d_2$}}};
    \draw[edge](T) -- (\x+1,\y-0.5);
    \node[draw=none] () at (\x+1,\y-0.25) {\scalebox{0.7}{\dimcolor{$d_4$}}};
    \draw[edge](T) -- (\x+1,\y+0.5);
    \node[draw=none] () at (\x+1,\y+0.25) {\scalebox{0.7}{\dimcolor{$d_1$}}};
    \draw[edge](v) -- (\x+3.25,\y)node [midway, above] {\scalebox{0.7}{\dimcolor{$p$}}};
\end{tikzpicture}\in\RR^{m\times n \times d_1\times d_2\times d_3 \times d_4\times p},
$$
其逐元素定义为 $(\Ab\outprod\Tt\outprod \vb)_{i_1i_2j_1j_2j_3j_4k} = \Ab_{i_1i_2}\Tt_{j_1\cdots j_4}\vb_{k}$，其中对 $l\in[4]$ 有 $i_1\in[m],i_2\in[n],j_l\in[d_l]$，且 $k\in[p]$。

下面的命题给出一个简洁的图示证明：外积的内积等于内积的乘积。
\begin{proposition}
    设 $\Xt,\Tt\in\RR^{d_1\times d_2\times\cdots\times d_N}$，其中 $\Xt = \ab_1\outprod\ab_2\outprod\cdots\outprod\ab_N$、$\Tt = \tb_1\outprod\tb_2\outprod\cdots\outprod\tb_N$，则 $\langle\Xt,\Tt\rangle = \prod_{n=1}^N\langle\ab_n,\tb_n\rangle$。
\begin{proof}
\begin{align*}
\langle\Xt,\Tt\rangle
&=
\left\langle\begin{tikzpicture}[baseline=\baseline]
    \tikzset{tensor/.style = {minimum size = 0.5cm,shape = circle,thick,draw=black,inner sep = 0pt}, edge/.style   = {thick,line width=.4mm}}
    \def\y{0}
     \def\x{0}
\node[tensor,fill=blue!10!white] (a1) at (\x,\y){$\ab_1$};
\node[tensor,fill=blue!10!white](a2) at (\x+1,\y){$\ab_2$};
  \node[tensor,fill=blue!10!white](an) at (\x+3,\y){$\ab_N$};
    \draw[edge] (a1) -- (\x,\y-0.7) node [midway, left] {\scalebox{0.7}{\dimcolor{$d_1$}}};
    \draw[edge] (a2) -- (\x+1,\y-0.7) node [midway, left] {\scalebox{0.7}{\dimcolor{$d_2$}}};
    \draw[edge] (an) -- (\x+3,\y-0.7)  node [midway, left] {\scalebox{0.7}{\dimcolor{$d_N$}}};
    \node (eq) at (\x+2,\y) {$\cdots$};
\end{tikzpicture} \ \scalebox{1.5}{,}
\begin{tikzpicture}[baseline=\baseline]
    \tikzset{tensor/.style = {minimum size = 0.5cm,shape = circle,thick,draw=black,inner sep = 0pt}, edge/.style   = {thick,line width=.4mm}}
    \def\y{0}
     \def\x{0}
\node[tensor,fill=red!10!white] (t1) at (\x,\y){$\tb_1$};
\node[tensor,fill=red!10!white](t2) at (\x+1,\y){$\tb_2$};
  \node[tensor,fill=red!10!white](tn) at (\x+3,\y){$\tb_N$};
    \draw[edge] (t1) -- (\x,\y-0.7) node [midway, left] {\scalebox{0.7}{\dimcolor{$d_1$}}};
    \draw[edge] (t2) -- (\x+1,\y-0.7) node [midway, left] {\scalebox{0.7}{\dimcolor{$d_2$}}};
    \draw[edge] (tn) -- (\x+3,\y-0.7)  node [midway, left] {\scalebox{0.7}{\dimcolor{$d_N$}}};
    \node (eq) at (\x+2,\y) {$\cdots$};
\end{tikzpicture} 
\right\rangle\\
&=
\begin{tikzpicture}[baseline=\baseline]
    \tikzset{tensor/.style = {minimum size = 0.5cm,shape = circle,thick,draw=black,inner sep = 0pt}, edge/.style   = {thick,line width=.4mm}}
    \def\y{0}
     \def\x{0}
\node[tensor,fill=blue!10!white] (a1) at (\x,\y+0.5){$\ab_1$};
\node[tensor,fill=blue!10!white](a2) at (\x+1,\y+0.5){$\ab_2$};
  \node[tensor,fill=blue!10!white](an) at (\x+3,\y+0.5){$\ab_N$};
  \node[tensor,fill=red!10!white] (t1) at (\x,\y-0.5){$\tb_1$};
\node[tensor,fill=red!10!white](t2) at (\x+1,\y-0.5){$\tb_2$};
  \node[tensor,fill=red!10!white](tn) at (\x+3,\y-0.5){$\tb_N$};
    \draw[edge] (a1) -- (t1) node [midway, left] {\scalebox{0.7}{\dimcolor{$d_1$}}};
    \draw[edge] (a2) -- (t2) node [midway, left] {\scalebox{0.7}{\dimcolor{$d_2$}}};
    \draw[edge] (an) -- (tn)  node [midway, left] {\scalebox{0.7}{\dimcolor{$d_N$}}};
    \node (eq) at (\x+2,\y) {$\cdots$};
\end{tikzpicture}
=
\prod_{n=1}^N\langle\ab_n,\tb_n\rangle.
\end{align*}
\end{proof}
\end{proposition}
\paragraph{维度为 1 的边等于没有边}
需要注意，在张量网络图中，大小为 1 的边等价于没有这条边。
例如在矩阵乘法中，若收缩边的大小为 1，它就等价于向量的外积：
\begin{align*}
    \Ab\in\RR^{p\times 1},\ \ \ \Bb\in\RR^{1\times q}, \ \ \ 
(\Ab\Bb)_{ij} 
&=
\begin{tikzpicture}[baseline=\baseline]
   \tikzset{tensor/.style = {minimum size = 0.5cm,shape = circle,thick,draw=black,inner sep = 0pt}, edge/.style = {thick,line width=.4mm},every loop/.style={}}
    \def\y{0}
    \def\x{0}
    \node[tensor,fill=red!50!white] (A) at (\x,\y){\scalebox{0.85}{$\Ab$}};
    \node[tensor,fill=blue!50!white] (B) at (\x+1,\y){\scalebox{0.85}{$\Bb$}};
    \draw[edge] (A) -- (B) node [above,midway] {\scalebox{0.7}{\dimcolor{$1$}}};
    \draw[edge] (\x-0.25,\y) -- (\x-0.75,\y)  node [pos=1.3] {\scalebox{0.7}{\indcolor{$i$}}};
    \draw[edge] (\x+1.25,\y) -- (\x+1.75,\y)  node [pos=1.3] {\scalebox{0.7}{\indcolor{$j$}}};
\end{tikzpicture}\\
&=
\sum_{k=1}^{1}\Ab_{i1}\Bb_{1j}
=
\begin{tikzpicture}[baseline=\baseline]
   \tikzset{tensor/.style = {minimum size = 0.5cm,shape = circle,thick,draw=black,inner sep = 0pt}, edge/.style = {thick,line width=.4mm},every loop/.style={}}
    \def\y{0}
    \def\x{0}
    \node[tensor,fill=red!50!white] (A) at (\x,\y){\scalebox{0.85}{$\ab$}};
    \node[tensor,fill=blue!50!white] (B) at (\x+1,\y){\scalebox{0.85}{$\bb$}};
    \draw[edge] (\x-0.25,\y) -- (\x-0.75,\y)  node[pos=1.3] {\scalebox{0.7}{\indcolor{$i$}}};
    \draw[edge] (\x+1.25,\y) -- (\x+1.75,\y)  node [pos=1.3] {\scalebox{0.7}{\indcolor{$j$}}};
\end{tikzpicture}\\
&=
(\ab\outprod\bb)_{ij},
\end{align*}
其中 $\ab$ 是 $\Ab$ 的唯一一列，$\bb$ 是 $\Bb$ 的唯一一行。
\paragraph{迹} 迹运算是张量收缩作用于方阵时的特例。在张量网络中，迹自然地由一条连接矩阵两条腿的环路~（自环边）表示：
\begin{align*}
\Ab\in\RR^{d \times d},\hspace{0.5cm}
\tr(\Ab) = \sum_{i=1}^d\Ab_{ii} = \hspace*{-0.5cm}
\begin{tikzpicture}[baseline=\baseline]
   \tikzset{tensor/.style = {minimum size = 0.5cm,shape = circle,thick,draw=black,inner sep = 0pt}, edge/.style = {thick,line width=.4mm},every loop/.style={}}
    \def\y{0}
    \def\x{0}
    \node[tensor,fill=red!50!white] (A) at (\x,\y){\scalebox{0.85}{$\Ab$}};
    \draw[edge] (A) to [out=20,in=160, distance=1cm] (A);
    \node[draw=none] () at (\x,\y+0.5) {\scalebox{0.7}{\dimcolor{$d$}}};
    \end{tikzpicture}\hspace*{-0.5cm}.
\end{align*}
图中没有自由腿，这与迹是标量这一事实相吻合。张量网络图为迹在循环置换下的不变性提供了一个极为简单的证明。先从两个矩阵乘积的迹入手，设 $\Ab \in \RR^{d\times R}$、$\Bb \in \RR^{R\times d}$。  如前所述，画张量网络图时重要的只是图的结构。由此即可看出
\begin{align*}
\tr(\Ab\Bb)= 
\hspace{-0.5cm}
\begin{tikzpicture}[baseline=\baseline]
   \tikzset{tensor/.style = {minimum size = 0.5cm,shape = circle,thick,draw=black,inner sep = 0pt}, edge/.style = {thick,line width=.4mm},every loop/.style={}}
    \def\y{0}
    \def\x{0}
    \node[tensor,fill=red!50!white] (A) at (\x,\y){\scalebox{0.85}{$\Ab$}};
    \node[tensor,fill=blue!50!white] (B) at (\x+1,\y){\scalebox{0.85}{$\Bb$}};
    \draw[edge] (A) -- (B) node [midway,above]{\scalebox{0.7}{\dimcolor{$R$}}};
    \draw[edge] (A) to [out=155,in=25, distance=1cm] (B);
    \node[draw=none] () at (\x+0.5,\y+0.65) {\scalebox{0.7}{\dimcolor{$d$}}};
    \end{tikzpicture}
\hspace{-0.5cm} = 
\begin{tikzpicture}[baseline=\baseline]
   \tikzset{tensor/.style = {minimum size = 0.5cm,shape = circle,thick,draw=black,inner sep = 0pt}, edge/.style = {thick,line width=.4mm},every loop/.style={}}
    \def\y{0}
    \def\x{0}
    \node[tensor,fill=red!50!white] (A) at (\x,\y-0.5){\scalebox{0.85}{$\Ab$}};
    \node[tensor,fill=blue!50!white] (B) at (\x,\y+0.5){\scalebox{0.85}{$\Bb$}};
    \draw[edge] (A) to [out=180,in=180, distance=0.5cm] (B);
    \draw[edge] (A) to [out=0,in=0, distance=0.5cm] (B);
    \node[draw=none] () at (\x+0.8,\y) {\scalebox{0.7}{\dimcolor{$R$}}};
    \node[draw=none] () at (\x-0.8,\y) {\scalebox{0.7}{\dimcolor{$d$}}};
    \end{tikzpicture}
 =\hspace{-0.5cm}
\begin{tikzpicture}[baseline=\baseline]
   \tikzset{tensor/.style = {minimum size = 0.5cm,shape = circle,thick,draw=black,inner sep = 0pt}, edge/.style = {thick,line width=.4mm},every loop/.style={}}
    \def\y{0}
    \def\x{0}
    \node[tensor,fill=blue!50!white] (B) at (\x,\y){\scalebox{0.85}{$\Bb$}};
    \node[tensor,fill=red!50!white] (A) at (\x+1,\y){\scalebox{0.85}{$\Ab$}};
    \draw[edge] (A) -- (B)node [midway,above]{\scalebox{0.7}{\dimcolor{$R$}}};;
    \draw[edge] (B) to [out=155,in=25, distance=1cm] (A);
    \node[draw=none] () at (\x+0.5,\y+0.65) {\scalebox{0.7}{\dimcolor{$d$}}};
    \end{tikzpicture}
\hspace{-0.4cm} =
\tr(\Bb\Ab),
\end{align*}
对三个矩阵 $\Ab\in\RR^{m\times n}, \Bb\in\RR^{n\times p},  \Cb\in\RR^{p\times m}$ 也有类似结论
\begin{align*}
\tr(\Ab\Bb\Cb) = 
\hspace{-0.7cm}
\begin{tikzpicture}[baseline=\baseline]
   \tikzset{tensor/.style = {minimum size = 0.5cm,shape = circle,thick,draw=black,inner sep = 0pt}, edge/.style = {thick,line width=.4mm},every loop/.style={}}
    \def\y{0}
    \def\x{0}
    \node[tensor,fill=red!50!white] (A) at (\x,\y){\scalebox{0.85}{$\Ab$}};
    \node[tensor,fill=blue!50!white] (B) at (\x+1,\y){\scalebox{0.85}{$\Bb$}};
    \node[tensor,fill=red!30!white] (C) at (\x+2,\y){\scalebox{0.85}{$\Cb$}};
    \draw[edge] (A) -- (B) node [midway,above]{\scalebox{0.7}{\dimcolor{$n$}}};
    \draw[edge] (B) -- (C) node [midway,above]{\scalebox{0.7}{\dimcolor{$p$}}};;
    \draw[edge] (A) to [out=155,in=25, distance=1cm] (C);
    \node[draw=none] () at (\x+1,\y+0.65) {\scalebox{0.7}{\dimcolor{$m$}}};
    \end{tikzpicture}
\hspace{-0.7cm}  = 
\begin{tikzpicture}[baseline=\baseline]
   \tikzset{tensor/.style = {minimum size = 0.5cm,shape = circle,thick,draw=black,inner sep = 0pt}, edge/.style = {thick,line width=.4mm},every loop/.style={}}
    \def\y{0}
    \def\x{0}
    \node[tensor,fill=red!50!white] (A) at (\x,\y+0.5){\scalebox{0.85}{$\Ab$}};
    \node[tensor,fill=blue!50!white] (B) at (\x-0.5,\y-0.3){\scalebox{0.85}{$\Bb$}};
    \node[tensor,fill=red!30!white] (C) at (\x+0.5,\y-0.3){\scalebox{0.85}{$\Cb$}};
    \draw[edge] (A) to [out=210,in=90] (B);
    \draw[edge] (A) to [out=-30,in=90] (C);
    \draw[edge] (B) to [out=300,in=240] (C);
    \node[draw=none] () at (\x+0.55,\y+0.3) {\scalebox{0.7}{\dimcolor{$m$}}};
    \node[draw=none] () at (\x-0.55,\y+0.3) {\scalebox{0.7}{\dimcolor{$n$}}};
    \node[draw=none] () at (\x,\y-0.85) {\scalebox{0.7}{\dimcolor{$p$}}};
    \end{tikzpicture}
= 
\hspace{-0.7cm}
\begin{tikzpicture}[baseline=\baseline]
   \tikzset{tensor/.style = {minimum size = 0.5cm,shape = circle,thick,draw=black,inner sep = 0pt}, edge/.style = {thick,line width=.4mm},every loop/.style={}}
    \def\y{0}
    \def\x{0}
    \node[tensor,fill=red!30!white] (C) at (\x,\y){\scalebox{0.85}{$\Cb$}};
    \node[tensor,fill=red!50!white] (A) at (\x+1,\y){\scalebox{0.85}{$\Ab$}};
    \node[tensor,fill=blue!50!white] (B) at (\x+2,\y){\scalebox{0.85}{$\Bb$}};
    \draw[edge] (C) -- (A) node [midway,above]{\scalebox{0.7}{\dimcolor{$m$}}};
    \draw[edge] (A) -- (B) node [midway,above]{\scalebox{0.7}{\dimcolor{$n$}}};
    \draw[edge] (C) to [out=155,in=25, distance=1cm] (B);
    \node[draw=none] () at (\x+1,\y+0.65) {\scalebox{0.7}{\dimcolor{$p$}}};
    \end{tikzpicture}
\hspace{-0.7cm} = \tr(\Cb\Ab\Bb).
\end{align*}
需要注意的是，张量网络图之间的这些等式是平凡的：我们只是挪动了节点的位置，而没有改变底层图的结构。因此，这只是对同一张图的两种解读，却得到了 $\tr(\Ab\Bb\Cb)$ 与 $\tr(\Cb\Ab\Bb)$ 之间并不那么平凡的等式。

\paragraph{偏迹} \emph{偏迹}（partial trace）是一类常见运算，尤其在多体量子系统的建模中：对整个系统的密度矩阵施加偏迹，便得到系统某一部分的密度矩阵。不必深究量子系统的具体细节，只需考虑一个维度为 $d_1d_2\times d_1d_2$ 的方阵，我们把它解释为一个四阶张量 \begin{tikzpicture}[baseline=-0.5ex]
    \tikzset{tensor/.style = {minimum width = 0.8cm,minimum height = 0.4cm,shape = rounded rectangle,thick,draw=black,inner sep = 0pt}, edge/.style = {thick,line width=.4mm},every loop/.style={}}
    \def\x{0}
    \def\y{0}
    \draw[edge] (\x-0.2,\y-0.15) -- (\x-0.2,\y-0.5)node [midway,left=-0.05cm] {\scalebox{0.7}{\dimcolor{$d_1$}}};
    \draw[edge] (\x+0.2,\y-0.15) -- (\x+0.2,\y-0.5)node [midway,right=-0.05cm] {\scalebox{0.7}{\dimcolor{$d_2$}}};
    \draw[edge] (\x-0.2,\y+0.15) -- (\x-0.2,\y+0.5)node [midway,left=-0.05cm] {\scalebox{0.7}{\dimcolor{$d_1$}}};
    \draw[edge] (\x+0.2,\y+0.15) -- (\x+0.2,\y+0.5)node [midway,right=-0.05cm] {\scalebox{0.7}{\dimcolor{$d_2$}}};
    \node[tensor,fill=yellow] (Bt) at (\x,\y){$\Tt$};
    \end{tikzpicture}。
$\Tt$ 关于子空间 $\mathbb R^{d_1}$ 的偏迹是一个 $d_2\times d_2$ 矩阵，把对应的两条腿相连，即对 $d_1$ 这一维``求迹''而得到： \begin{tikzpicture}[baseline=-0.5ex]
    \tikzset{tensor/.style = {minimum width = 0.6cm,minimum height = 0.4cm,shape = rounded rectangle,thick,draw=black,inner sep = 0pt}, edge/.style = {thick,line width=.4mm},every loop/.style={}}
    \def\x{0}
    \def\y{0}
    \draw[edge] (\x-0.2,\y-0.15) -- (\x-0.2,\y-0.4);
    \draw[edge] (\x+0.2,\y-0.15) -- (\x+0.2,\y-0.5)node [midway,right=-0.05cm] {\scalebox{0.7}{\dimcolor{$d_2$}}};
    \draw[edge] (\x-0.2,\y+0.15) -- (\x-0.2,\y+0.4);
    \draw[edge] (\x+0.2,\y+0.15) -- (\x+0.2,\y+0.5)node [midway,right=-0.05cm] {\scalebox{0.7}{\dimcolor{$d_2$}}};
    \draw[edge] (\x-0.2,\y-0.39) to [out=200,in=160] (\x-0.2,\y+0.39) node [midway,left=0.35cm] {\scalebox{0.7}{\dimcolor{$d_1$}}};
    \node[tensor,fill=yellow] (Bt) at (\x,\y){$\Tt$};
    \end{tikzpicture}。类似地，$\Tt$ 关于子空间 $\mathbb R^{d_2}$ 的偏迹是 $d_1\times d_1$ 矩阵
\begin{tikzpicture}[baseline=-0.5ex]
    \tikzset{tensor/.style = {minimum width = 0.6cm,minimum height = 0.4cm,shape = rounded rectangle,thick,draw=black,inner sep = 0pt}, edge/.style = {thick,line width=.4mm},every loop/.style={}}
    \def\x{0}
    \def\y{0}
    \draw[edge] (\x-0.2,\y-0.15) -- (\x-0.2,\y-0.5)node [midway,left=-0.05cm] {\scalebox{0.7}{\dimcolor{$d_1$}}};
    \draw[edge] (\x+0.2,\y-0.15) -- (\x+0.2,\y-0.4);
    \draw[edge] (\x-0.2,\y+0.15) -- (\x-0.2,\y+0.5)node [midway,left=-0.05cm] {\scalebox{0.7}{\dimcolor{$d_1$}}};
    \draw[edge] (\x+0.2,\y+0.15) -- (\x+0.2,\y+0.4);
    \draw[edge] (\x+0.2,\y-0.39) to [out=330,in=30] (\x+0.2,\y+0.39) node [midway,right=0.3cm] {\scalebox{0.7}{\dimcolor{$d_2$}}};
    \node[tensor,fill=yellow] (Bt) at (\x,\y){$\Tt$};
\end{tikzpicture}。

\paragraph{张量的 Frobenius 范数} 张量 $\At\in\RR^{d_1\times\cdots\times d_N}$ 的 Frobenius 范数是其各元素
平方之和的平方根~（即与自身内积的平方根）。在张量网络中，三阶张量 $\At\in\RR^{d_1\times d_2\times d_3}$ 的范数表示为
\begin{align*}
\norm{\At}_F^2 = \langle\At,\At\rangle = 
\sum_{i_1=1}^{d_1}\sum_{i_2=1}^{d_2}\sum_{i_3=1}^{d_3}\At_{i_1i_2i_3}^2
=
\begin{tikzpicture}[baseline=\baseline]
    \tikzset{tensor/.style = {minimum size = 0.5cm,shape = circle,thick,draw=black,inner sep = 0pt}, edge/.style = {thick,line width=.4mm},every loop/.style={}}
    \def\y{0}
    \def\x{0}
    \node[tensor,fill=green!50!red!50!white] (At) at (\x,\y){\scalebox{0.85}{$\At$}};
    \node[tensor,fill=green!50!red!50!white] (At1) at (\x+1,\y){\scalebox{0.85}{$\At$}};
    \draw[edge] (At) -- (At1);
    \draw[edge] (At) to [out=45,in=135] (At1);
    \draw[edge] (At) to [out=-45,in=-135] (At1);
    \node[draw=none] () at (\x+0.5,\y+0.6) {\scalebox{0.7}{\dimcolor{$d_1$}}};
    \node[draw=none] () at (\x+0.5,\y+0.15) {\scalebox{0.7}{\dimcolor{$d_2$}}};
    \node[draw=none] () at (\x+0.5,\y-0.5) {\scalebox{0.7}{\dimcolor{$d_3$}}};
\end{tikzpicture} .
\end{align*}
它是矩阵 Frobenius 范数最直接、自然的推广：
\begin{align*}
 \norm{\Ab}_F^2 = \tr(\Ab\Ab^\ts) = \tr(\Ab^\ts\Ab) =  \hspace*{-0.5cm}
\begin{tikzpicture}[baseline=\baseline]
   \tikzset{tensor/.style = {minimum size = 0.5cm,shape = circle,thick,draw=black,inner sep = 0pt}, edge/.style = {thick,line width=.4mm},every loop/.style={}}
    \def\y{0}
    \def\x{0}
    \node[tensor,fill=red!50!white] (A) at (\x,\y){\scalebox{0.85}{$\Ab$}};
    \node[tensor,fill=red!50!white] (B) at (\x+1,\y){\scalebox{0.85}{$\Ab$}};
    \draw[edge] (A) -- (B) node [midway,above]{\scalebox{0.7}{\dimcolor{$n$}}};
    \draw[edge] (A) to [out=155,in=25, distance=1cm] (B);
    \node[draw=none] () at (\x+0.5,\y+0.65) {\scalebox{0.7}{\dimcolor{$m$}}};
    \end{tikzpicture} \hspace*{-0.5cm}.
\end{align*}


\paragraph{外积的 Frobenius 范数} 借助张量网络，几乎可以显然地验证：外积的 Frobenius 范数等于各因子 Frobenius 范数的乘积。例如，取 $\Ab\in\RR^{m\times n}$ 与 $\Tt\in\RR^{d_1\times d_2\times d_3}$，有
$$ \begin{tikzpicture}[baseline=-0.5ex]
    \tikzset{tensor/.style = {minimum width = 1.8cm,minimum height = 0.4cm,shape = rounded rectangle,thick,draw=black,inner sep = 0pt}, edge/.style = {thick,line width=.4mm},every loop/.style={}}
    \def\x{0}
    \def\y{0}
    \draw[edge] (\x-0.7,\y-0.15) -- (\x-0.7,\y-0.5)node [midway,left=-0.05cm] {\scalebox{0.7}{\dimcolor{$m$}}};
    \draw[edge] (\x-0.35,\y-0.15) -- (\x-0.35,\y-0.5)node [midway,left=-0.05cm] {\scalebox{0.7}{\dimcolor{$n$}}};
    \draw[edge] (\x,\y-0.15) -- (\x,\y-0.5)node [midway,left=-0.1cm] {\scalebox{0.7}{\dimcolor{$d_1$}}};
    \draw[edge] (\x+0.35,\y-0.15) -- (\x+0.35,\y-0.5)node [midway,left=-0.1cm] {\scalebox{0.7}{\dimcolor{$d_2$}}};
    \draw[edge] (\x+0.7,\y-0.15) -- (\x+0.7,\y-0.5)node [midway,left=-0.1cm] {\scalebox{0.7}{\dimcolor{$d_3$}}};
    \node[tensor,fill=yellow] (Bt) at (\x,\y){$\Ab \outprod \Tt$};
    \end{tikzpicture}
    = 
    \begin{tikzpicture}[baseline=\baseline]
   \tikzset{tensor/.style = {minimum width = 0.8cm,minimum height = 0.4cm,shape = rounded rectangle,thick,draw=black,inner sep = 0pt}, edge/.style = {thick,line width=.4mm},every loop/.style={}}
    \def\y{0}
    \def\x{0}
    \draw[edge] (\x+0.2,\y-0.15) -- (\x+0.2,\y-0.5) node [midway,left=-0.05cm]{\scalebox{0.7}{\dimcolor{$n$}}};
    \draw[edge] (\x-0.2,\y-0.15) -- (\x-0.2,\y-0.5) node [midway,left=-0.05cm]{\scalebox{0.7}{\dimcolor{$m$}}};
    \node[tensor,fill=red!50!white] (A) at (\x,\y){{$\Ab$}};

   
    \def\y{0}
    \def\x{1}
    \draw[edge] (\x-0.3,\y-0.15) -- (\x-0.3,\y-0.5) node [midway,left=-0.1cm]{\scalebox{0.7}{\dimcolor{$d_1$}}};
    \draw[edge] (\x+0.05,\y-0.15) -- (\x+0.05,\y-0.5) node [midway,left=-0.1cm]{\scalebox{0.7}{\dimcolor{$d_2$}}};
    \draw[edge] (\x+0.3,\y-0.15) -- (\x+0.3,\y-0.5) node [midway,right=-0.05cm]{\scalebox{0.7}{\dimcolor{$d_3$}}};
    \node[tensor,fill=red!30!blue!20!white] (T) at (\x,\y){{$\Tt$}};
\end{tikzpicture}
    $$
于是
\begin{align*}
\norm{\Ab\outprod\Tt}_F^2
=
\norm{
\begin{tikzpicture}[baseline=\baseline]
\tikzset{tensor/.style = {minimum width = 1.8cm,minimum height = 0.4cm,shape = rounded rectangle,thick,draw=black,inner sep = 0pt}, edge/.style = {thick,line width=.4mm},every loop/.style={}}
    \def\x{0}
    \def\y{0}
    \draw[edge] (\x-0.7,\y-0.15) -- (\x-0.7,\y-0.5)node [midway,left=-0.05cm] {\scalebox{0.7}{\dimcolor{$m$}}};
    \draw[edge] (\x-0.35,\y-0.15) -- (\x-0.35,\y-0.5)node [midway,left=-0.05cm] {\scalebox{0.7}{\dimcolor{$n$}}};
    \draw[edge] (\x,\y-0.15) -- (\x,\y-0.5)node [midway,left=-0.1cm] {\scalebox{0.7}{\dimcolor{$d_1$}}};
    \draw[edge] (\x+0.35,\y-0.15) -- (\x+0.35,\y-0.5)node [midway,left=-0.1cm] {\scalebox{0.7}{\dimcolor{$d_2$}}};
    \draw[edge] (\x+0.7,\y-0.15) -- (\x+0.7,\y-0.5)node [midway,left=-0.1cm] {\scalebox{0.7}{\dimcolor{$d_3$}}};
    \node[tensor,fill=yellow] (Bt) at (\x,\y){$\Ab \outprod \Tt$};
\end{tikzpicture}
}_F^2
=
\begin{tikzpicture}[baseline=\baseline]
\tikzset{tensor/.style = {minimum width = 1.8cm,minimum height = 0.4cm,shape = rounded rectangle,thick,draw=black,inner sep = 0pt}, edge/.style = {thick,line width=.4mm},every loop/.style={}}
    \def\x{0}
    \def\y{0.35}
    \draw[edge] (\x-0.7,\y-0.15) -- (\x-0.7,\y-0.5)node [midway,left=-0.05cm] {\scalebox{0.7}{\dimcolor{$m$}}};
    \draw[edge] (\x-0.35,\y-0.15) -- (\x-0.35,\y-0.5)node [midway,left=-0.05cm] {\scalebox{0.7}{\dimcolor{$n$}}};
    \draw[edge] (\x,\y-0.15) -- (\x,\y-0.5)node [midway,left=-0.1cm] {\scalebox{0.7}{\dimcolor{$d_1$}}};
    \draw[edge] (\x+0.35,\y-0.15) -- (\x+0.35,\y-0.5)node [midway,left=-0.1cm] {\scalebox{0.7}{\dimcolor{$d_2$}}};
    \draw[edge] (\x+0.7,\y-0.15) -- (\x+0.7,\y-0.5)node [midway,left=-0.1cm] {\scalebox{0.7}{\dimcolor{$d_3$}}};
    \node[tensor,fill=yellow] (Bt) at (\x,\y){$\Ab \outprod \Tt$};
    \node[tensor,fill=yellow] (Bt) at (\x,\y-0.7){$\Ab \outprod \Tt$};
\end{tikzpicture}
=
\begin{tikzpicture}[baseline=\baseline]
   \tikzset{tensor/.style = {minimum width = 0.9cm,minimum height = 0.4cm,shape = rounded rectangle,thick,draw=black,inner sep = 0pt}, edge/.style = {thick,line width=.4mm},every loop/.style={}}
    \def\y{0.35}
    \def\x{0}
    \draw[edge] (\x+0.2,\y-0.15) -- (\x+0.2,\y-0.5) node [midway,left=-0.05cm]{\scalebox{0.7}{\dimcolor{$n$}}};
    \draw[edge] (\x-0.2,\y-0.15) -- (\x-0.2,\y-0.5) node [midway,left=-0.05cm]{\scalebox{0.7}{\dimcolor{$m$}}};
    \node[tensor,fill=red!50!white] (A) at (\x,\y){{$\Ab$}};
    \node[tensor,fill=red!50!white] (A) at (\x,\y-0.7){{$\Ab$}};
    \def\x{1}
    \draw[edge] (\x-0.3,\y-0.15) -- (\x-0.3,\y-0.5) node [midway,left=-0.1cm]{\scalebox{0.7}{\dimcolor{$d_1$}}};
    \draw[edge] (\x+0.05,\y-0.15) -- (\x+0.05,\y-0.5) node [midway,left=-0.1cm]{\scalebox{0.7}{\dimcolor{$d_2$}}};
    \draw[edge] (\x+0.3,\y-0.15) -- (\x+0.3,\y-0.5) node [midway,right=-0.1cm]{\scalebox{0.7}{\dimcolor{$d_3$}}};
    \node[tensor,fill=red!30!blue!20!white] (T) at (\x,\y){{$\Tt$}};
    \node[tensor,fill=red!30!blue!20!white] (T) at (\x,\y-0.7){{$\Tt$}};
\end{tikzpicture}
=\norm{\Ab}_F^2 \norm{\Tt}_F^2
.
\end{align*}


\paragraph{例} 本节以最后一个例子收尾，说明张量网络的图示记法在表示复杂的张量运算时多么有用：
\begin{align}\label{eq:einsum}
\begin{tikzpicture}[baseline = \baseline]
    \tikzset{tensor/.style = {minimum size = 0.5cm,shape = circle,thick,draw=black,inner sep = 0pt}, edge/.style = 
    {thick,line width=.4mm},every loop/.style={}}
    \def\y{-0.5}
    \def\x{0}
    \node[tensor,fill=green!50!red!50!white] (At) at (\x,\y){\scalebox{0.85}{$\At$}};
    \node[tensor,fill=red!20!white] (S) at (\x-1,\y){\scalebox{0.85}{$\St$}};
    \node[tensor,fill=red!50!white] (Tt) at (\x,\y+1){\scalebox{0.85}{$\Tt$}};
    \node[tensor,fill=blue!40!red!60!white] (A) at (\x+1,\y){\scalebox{0.85}{$\Ab$}};
    \draw[edge] (Tt) -- (At) node [midway,left] {\scalebox{0.7}{\dimcolor{$R$}}};
    \draw[edge] (Tt) -- (A) node [above,midway] {\scalebox{0.7}{\dimcolor{$R$}}};
    \draw[edge] (At) -- (A)  node [midway,above] {\scalebox{0.7}{\dimcolor{$R$}}};
    \draw[edge] (Tt) -- (\x,\y+1.75) node [pos=1.3] {\scalebox{0.7}{\indcolor{$k$}}};
    \draw[edge] (At) -- (S)  node [midway,above] {\scalebox{0.7}{\dimcolor{$R$}}};
    \draw[edge] (S) -- (Tt)  node [midway,above] {\scalebox{0.7}{\dimcolor{$R$}}};
    \draw[edge] (Tt) -- (\x-0.75,\y+1)  node [pos=1.3] {\scalebox{0.7}{\indcolor{$j$}}};
    \draw[edge] (Tt) -- (\x+0.75,\y+1)  node [pos=1.3] {\scalebox{0.7}{\indcolor{$l$}}};
    \draw[edge] (S) -- (\x-1.2,\y+0.75)  node [pos=1.3] {\scalebox{0.7}{\indcolor{$i$}}};
    \end{tikzpicture}
&=
\sum_{r_1=1}^R\sum_{r_2=1}^R\sum_{r_3=1}^R\sum_{r_4=1}^R\sum_{r_5=1}^R\St_{ir_1 r_2}\At_{r_2r_3r_5}\Tt_{iklr_1r_3r_4}\Ab_{r_4r_5}\nonumber\\
&=
    \begin{tikzpicture}[baseline = \baseline]
    \tikzset{tensor/.style = {minimum size = 0.5cm,shape = circle,thick,draw=black,inner sep = 0pt}, edge/.style = 
    {thick,line width=.4mm},every loop/.style={}}
    \def\y{0}
    \def\x{0}
    \node[tensor,fill=blue!50!red!30!yellow!50!white] (Ht) at (\x,\y){\scalebox{0.85}{$\Ht$}};
    \draw[edge](Ht) -- (\x,\y+0.7) node [pos=1.3] {\scalebox{0.7}{\indcolor{$j$}}};
    \draw[edge](Ht) -- (\x,\y-0.7) node [pos=1.3] {\scalebox{0.7}{\indcolor{$l$}}};
    \draw[edge](Ht) -- (\x-0.7,\y) node [pos=1.3] {\scalebox{0.7}{\indcolor{$i$}}};
    \draw[edge](Ht) -- (\x+0.7,\y) node[pos=1.3] {\scalebox{0.7}{\indcolor{$k$}}};
    \end{tikzpicture}.
\end{align}
\subsection{复制张量~（超边）}\label{sec:copy-tensor}


有时需要表示由多于两个张量共享的同时收缩。例如，把三个向量 $\ab\in\RR^d, \bb\in\RR^d,$ 与 $\cbb\in\RR^d$ 收缩成一个标量：$\sum_{i=1}^d\ab_i\bb_i\cbb_i$。
这一运算在张量网络中可用一种特殊的张量——\emph{复制张量}（copy tensor）——来刻画。复制张量等价于克罗内克 $\delta$，也就是对角元全为 1、其余元素全为 0 的超对角张量，在张量网络图中画作一个黑点。例如，复制张量
\begin{tikzpicture}[baseline=\baseline]
   \tikzset{tensor/.style = {minimum size = 0.5cm,shape = circle,thick,draw=black,inner sep = 0pt}, edge/.style = {thick,line width=.4mm},every loop/.style={}}
    \def\y{0}
    \def\x{0}
     \draw  node[fill = black,circle,inner sep=0pt,minimum size=4pt] (m) at (\x,\y) {};
     \draw  node[draw=none] (a) at (\x-0.7,\y) {};
     \draw  node[draw=none] (b) at (\x+0.7,\y) {};
     \draw  node[draw=none] (c) at (\x,\y-0.7) {};
    \draw[edge] (a) -- (m);
    \draw[edge] (m) -- (b);
    \draw[edge] (m) -- (c);
\end{tikzpicture}
其逐元素定义为
$$
 \left(\begin{tikzpicture}[baseline=\baseline]
   \tikzset{tensor/.style = {minimum size = 0.5cm,shape = circle,thick,draw=black,inner sep = 0pt}, edge/.style = {thick,line width=.4mm},every loop/.style={}}
    \def\y{0}
    \def\x{0}
     \draw  node[fill = black,circle,inner sep=0pt,minimum size=4pt] (m) at (\x,\y) {};
     \draw  node[draw=none] (a) at (\x-0.7,\y) {};
     \draw  node[draw=none] (b) at (\x+0.7,\y) {};
     \draw  node[draw=none] (c) at (\x,\y-0.7) {};
    \draw[edge] (a) -- (m);
    \draw[edge] (m) -- (b);
    \draw[edge] (m) -- (c);
\end{tikzpicture}\right)_{ijk} = \delta_{ijk} = \begin{cases}
      1 & \text{if } i=j=k\\
      0 & \text{otherwise}
    \end{cases}.
$$
借助复制张量，上述三重收缩可表示为
\begin{align*}
\begin{tikzpicture}[baseline=\baseline]
   \tikzset{tensor/.style = {minimum size = 0.5cm,shape = circle,thick,draw=black,inner sep = 0pt}, edge/.style = {thick,line width=.4mm},every loop/.style={}}
    \def\y{0}
    \def\x{0}
    \node[tensor,fill=blue!40!red!60!white] (a) at (\x,\y){\scalebox{0.85}{$\ab$}};
    \node[tensor,fill=blue!10!white] (b) at (\x+1.5,\y){\scalebox{0.85}{$\bb$}};
    \node[tensor,fill=blue!20!red!40!white] (c) at (\x+0.7,\y-0.75){\scalebox{0.85}{$\cbb$}};
    \draw  node[fill = black,circle,inner sep=0pt,minimum size=4pt] (m) at (\x+0.7,\y) {};
    \draw[edge] (a) -- (m) node [above,midway] {\scalebox{0.7}{\dimcolor{$d$}}};
    \draw[edge] (m) -- (b) node [above,midway] {\scalebox{0.7}{\dimcolor{$d$}}};
    \draw[edge] (m) -- (c) node [left,midway] {\scalebox{0.7}{\dimcolor{$d$}}};
    \end{tikzpicture} = \sum_{i,j,k=1}^d\delta_{ijk}\ab_i\bb_j\cbb_k =
\sum_{i=1}^d\ab_i\bb_i\cbb_i .
\end{align*}
下面给出任意阶 $N$ 的复制张量更严格的定义。
\begin{definition}
阶为 N 的复制张量
   \begin{tikzpicture}[baseline=\baseline]
    \tikzset{tensor/.style = {minimum size = 0.5cm,shape = circle,thick,draw=black,inner sep = 0pt}, edge/.style   = {thick,line width=.4mm},every loop/.style={}}
    \def\y{0}
    \def\x{0}
    \draw  node[fill = black,circle,inner sep=0pt,minimum size=4pt] (m) at (\x,\y) {};
    \draw[edge] (m) -- (\x-0.6,0) node [midway,above] {\scalebox{0.7}{\dimcolor{$d$}}};
    \draw[edge] (m) -- (\x+0.6,0) node [midway,above] {\scalebox{0.7}{\dimcolor{$d$}}};;
    \draw[edge] (m) -- (\x-0.5,-0.3) node [midway, below] {\scalebox{0.7}{\dimcolor{$d$}}};;
    \draw[dotted] (m) -- (\x,-0.6);
    \draw[dotted] (m) -- (\x,-0.6);
    \draw[dotted] (m) -- (\x+0.5,-0.5);
    \draw[dotted] (m) -- (\x+0.3,-0.6);
\end{tikzpicture}
是形状为 $\underbrace{d\times d\dots\times d}_{N\text{ times}}$ 的张量，定义为
\begin{align*}
       \begin{tikzpicture}[baseline=\baseline]
    \tikzset{tensor/.style = {minimum size = 0.5cm,shape = circle,thick,draw=black,inner sep = 0pt}, edge/.style   = {thick,line width=.4mm},every loop/.style={}}
    \def\y{0}
    \def\x{0}
    \draw  node[fill = black,circle,inner sep=0pt,minimum size=4pt] (m) at (\x,\y) {};
    \draw[edge] (m) -- (\x-0.6,0) node [midway,above] {\scalebox{0.7}{\dimcolor{$d$}}};
    \draw[edge] (m) -- (\x+0.6,0) node [midway,above] {\scalebox{0.7}{\dimcolor{$d$}}};;
    \draw[edge] (m) -- (\x-0.5,-0.3) node [midway, below] {\scalebox{0.7}{\dimcolor{$d$}}};;
    \draw[dotted] (m) -- (\x,-0.6);
    \draw[dotted] (m) -- (\x,-0.6);
    \draw[dotted] (m) -- (\x+0.5,-0.5);
    \draw[dotted] (m) -- (\x+0.3,-0.6);
\end{tikzpicture}
= \sum_{i=1}^d\underbrace{\eb_i\outprod\eb_i\outprod\dots\outprod\eb_i}_{N\text{ times}},%
\end{align*}
其中 $\eb_1,\eb_2,\dots,\eb_d$ 为 $\RR^{d}$ 的标准基向量。
\end{definition}
复制张量之名正由此而来：与复制张量收缩，便可``复制''标准基向量。事实上，设 $\eb_i\in\RR^{d}$ 为标准基向量，则容易
验证
$       
\begin{tikzpicture}[baseline=\baseline]
    \tikzset{tensor/.style = {minimum size = 0.5cm,shape = circle,thick,draw=black,inner sep = 0pt}, edge/.style   = {thick,line width=.4mm},every loop/.style={}}
    \def\y{0}
    \def\x{0}
    \node[tensor,fill=blue!40!red!60!white] (e) at (\x,\y){\scalebox{0.85}{$\eb_i$}};
    \draw  node[fill = black,circle,inner sep=0pt,minimum size=4pt] (m) at (\x+0.75,\y) {};
    \draw[edge] (e) -- (m) node [midway,above] {\scalebox{0.7}{\dimcolor{$d$}}};
    \draw[edge] (m) -- (\x+1.4,\y) node [midway,above] {\scalebox{0.7}{\dimcolor{$d$}}};
    \draw[edge] (m) -- (\x+0.75,\y-0.7) node [midway,right] {\scalebox{0.7}{\dimcolor{$d$}}};
    \node[draw=none] () at (\x+2,\y) {\scalebox{1}{$=$}};
    \node[tensor,fill=blue!40!red!60!white] (e1) at (\x+2.5,\y){\scalebox{0.85}{$\eb_i$}};
    \node[tensor,fill=blue!40!red!60!white] (e2) at (\x+3.5,\y){\scalebox{0.85}{$\eb_i$}};
    \draw[edge] (e2) -- (\x+3.5,\y-0.7) node [midway,right] {\scalebox{0.7}{\textcolor{gray}{$d$}}};
    \draw[edge] (e1) -- (\x+2.5,\y-0.7) node [midway,right] {\scalebox{0.7}{\textcolor{gray}{$d$}}};
\end{tikzpicture}.
$
简短证明如下：
\begin{align*}
    \left(
\begin{tikzpicture}[baseline=\baseline]
    \tikzset{tensor/.style = {minimum size = 0.5cm,shape = circle,thick,draw=black,inner sep = 0pt}, edge/.style   = {thick,line width=.4mm},every loop/.style={}}
    \def\y{0}
    \def\x{0}
    \node[tensor,fill=blue!40!red!60!white] (e) at (\x,\y){\scalebox{0.85}{$\eb_i$}};
    \draw  node[fill = black,circle,inner sep=0pt,minimum size=4pt] (m) at (\x+0.75,\y) {};
    \draw[edge] (e) -- (m) node [midway,above] {\scalebox{0.7}{\dimcolor{$d$}}};
    \draw[edge] (m) -- (\x+1.4,\y) node [midway,above] {\scalebox{0.7}{\dimcolor{$d$}}};
    \draw[edge] (m) -- (\x+0.75,\y-0.7)node [midway,right] {\scalebox{0.7}{\dimcolor{$d$}}};
\end{tikzpicture}\right)_{jk}
&= 
\sum_{s}(\eb_i)_{s}\delta_{sjk} = \sum_{s}\delta_{is}\delta_{sjk} = \delta_{ijk}\\ 
&= 
\begin{cases}
     1 & \text{if } i=j=k\\
      0 & \text{otherwise}.
    \end{cases} =
    (\eb_i\eb_i^\ts)_{j,k} =
\begin{tikzpicture}[baseline=\baseline]
    \tikzset{tensor/.style = {minimum size = 0.5cm,shape = circle,thick,draw=black,inner sep = 0pt}, edge/.style   = {thick,line width=.4mm},every loop/.style={}}
    \def\y{0}
    \def\x{0}
    \node[tensor,fill=blue!40!red!60!white] (e1) at (\x,\y){\scalebox{0.85}{$\eb_i$}};
    \node[tensor,fill=blue!40!red!60!white] (e2) at (\x+1,\y){\scalebox{0.85}{$\eb_i$}};
    \draw[edge] (e2) -- (\x+1,\y-0.7) node [below] {\scalebox{0.7}{\indcolor{$k$}}};
    \draw[edge] (e1) -- (\x,\y-0.7) node [below] {\scalebox{0.7}{\indcolor{$j$}}};
\end{tikzpicture}.
\end{align*}
务必注意，复制张量的``复制''性质只对标准基向量成立，对任意向量并不成立：一般而言
$       
\begin{tikzpicture}[baseline=\baseline]
    \tikzset{tensor/.style = {minimum size = 0.5cm,shape = circle,thick,draw=black,inner sep = 0pt}, edge/.style   = {thick,line width=.4mm},every loop/.style={}}
    \def\y{0}
    \def\x{0}
    \node[tensor,fill=blue!40!red!60!white] (e) at (\x,\y){\scalebox{0.85}{$\vb$}};
    \draw  node[fill = black,circle,inner sep=0pt,minimum size=4pt] (m) at (\x+0.75,\y) {};
    \draw[edge] (e) -- (m) node [midway,above] {\scalebox{0.7}{\dimcolor{$d$}}};
    \draw[edge] (m) -- (\x+1.4,\y) node [midway,above] {\scalebox{0.7}{\dimcolor{$d$}}};
    \draw[edge] (m) -- (\x+0.75,\y-0.7) node [midway,right] {\scalebox{0.7}{\dimcolor{$d$}}};
    \node[draw=none] () at (\x+2,\y) {\scalebox{1}{$\neq$}};
    \node[tensor,fill=blue!40!red!60!white] (e1) at (\x+2.5,\y){\scalebox{0.85}{$\vb$}};
    \node[tensor,fill=blue!40!red!60!white] (e2) at (\x+3.5,\y){\scalebox{0.85}{$\vb$}};
    \draw[edge] (e2) -- (\x+3.5,\y-0.7) node [midway,right] {\scalebox{0.7}{\dimcolor{$d$}}};
    \draw[edge] (e1) -- (\x+2.5,\y-0.7) node [midway,right] {\scalebox{0.7}{\dimcolor{$d$}}};
\end{tikzpicture}
$。复制张量之所以只``复制''基向量，是因为 $\delta_{ijk}$ 只挑出对角元。
\begin{proposition}\label{prop:copy.tensor} 下面列出复制张量的若干有用性质。
\begin{enumerate}
    \item 最简单的复制张量是元素全为 1 的向量，即
$$
  \begin{tikzpicture}[baseline=\baseline]
    \tikzset{tensor/.style = {minimum size = 0.5cm,shape = circle,thick,draw=black,inner sep = 0pt}, edge/.style   = {thick,line width=.4mm},every loop/.style={}}
    \def\y{0.5}
    \def\x{0}
    \draw  node[fill = black,circle,inner sep=0pt,minimum size=4pt] (m) at (\x,\y) {};
    \draw[edge] (m) -- (\x,\y-1) node [left,midway] {\scalebox{0.7}{\dimcolor{$n$}}};
 \end{tikzpicture}
 =\sum_{i=1}^n\eb_i =\overrightarrow{\mathbf{1}}.
$$\label{itm:one-copy-tensor}
\item 二阶复制张量就是单位矩阵，此处可直接省去中间的黑点，把它画成一条线，即
$
    \begin{tikzpicture}[baseline=\baseline]
    \tikzset{tensor/.style = {minimum size = 0.5cm,shape = circle,thick,draw=black,inner sep = 0pt}, edge/.style   = {thick,line width=.4mm},every loop/.style={}}
    \def\y{0}
    \def\x{0}
    \draw  node[fill = black,circle,inner sep=0pt,minimum size=4pt] (m) at (\x,\y) {};
    \draw  node[fill = black,circle,inner sep=0pt,minimum size=4pt] (m) at (\x,\y) {};
    \draw[edge] (m) -- (\x-0.6,0) node [very near end,left] {\scalebox{0.7}{\indcolor{$i$}}};
    \draw[edge] (m) -- (\x+0.6,0) node [very near end,right]
    {\scalebox{0.7}{\indcolor{$i$}}};
    \node[draw=none] at (\x+1,\y) {\scalebox{1}{\textcolor{black}{$=$}}};;
    \draw[edge] (\x+1.3,\y) -- (\x+2.5,\y) node [very near start,left=0.05cm] {\scalebox{0.7}{\indcolor{$i$}}} node [very near end,right=0.05cm] {\scalebox{0.7}{\indcolor{$i$}}};
    \end{tikzpicture} = \delta_{ij} = \begin{cases}
      1 & \text{if } i=j\\
      0 & \text{otherwise},
    \end{cases}
    =
    \Ib_{ij}.
$
\item 
有趣的是，把一个复制张量的任意若干条腿与另一个复制张量的腿收缩后，所得的张量网络本身仍是一个复制
张量（在图中把被收缩的腿与黑点合并即得）。例如：
\begin{align*}  
\sum_{l_1}\delta_{i_1j_1k_1l_1}\delta_{i_2j_2k_2l_1} 
&=
\begin{tikzpicture}[baseline=\baseline]
    \tikzset{tensor/.style = {minimum size = 0.5cm,shape = circle,thick,draw=black,inner sep = 0pt}, edge/.style   = {thick,line width=.4mm},every loop/.style={}}
    \def\y{0}
    \def\x{0}
    \draw  node[fill = black,circle,inner sep=0pt,minimum size=4pt] (m) at (\x,\y) {};
    \draw[edge] (m) -- (\x,\y-0.5)node [below] {\scalebox{0.6}{\indcolor{$i_1$}}};
    \draw[edge](m) -- (\x,\y+0.5)node [above] {\scalebox{0.6}{\indcolor{$
    k_1$}}};
    \draw[edge] (m) -- (\x-0.5,\y)node [left] {\scalebox{0.6}{\indcolor{$j_1$}}};
    \draw  node[fill = black,circle,inner sep=0pt,minimum size=4pt] (n) at (\x+1,\y) {};
    \draw[edge] (n) -- (\x+1.5,\y)node [right] {\scalebox{0.6}{\indcolor{$j_2$}}};
    \draw[edge] (n) --  (\x+1,\y+0.5)node [above] {\scalebox{0.6}{\indcolor{$k_2$}}};   
    \draw[edge] (n) -- (\x+1,\y-0.5)node [below] {\scalebox{0.6}{\indcolor{$i_2$}}};
    \draw[edge](m) -- (n)node [midway, above] {\scalebox{0.6}{\indcolor{$l_1$}}};
 \end{tikzpicture}
= 
\delta_{i_1j_1k_1i_2j_2k_2}\\
&=
\begin{tikzpicture}[baseline=\baseline]
    \tikzset{tensor/.style = {minimum size = 0.5cm,shape = circle,thick,draw=black,inner sep = 0pt}, edge/.style   = {thick,line width=.4mm},every loop/.style={}}
    \def\y{0}
    \def\x{0}
    \draw node[fill = black,circle,inner sep=0pt,minimum size=5pt] (m) at (\x,\y) {};
    \draw[edge] (m) -- (\x+0.75,\y)node [right] {\scalebox{0.6}{\indcolor{$j_2$}}};
    \draw[edge] (m) -- (\x-0.75,\y)node [left] {\scalebox{0.6}{\indcolor{$j_1$}}};
    \draw[edge] (m) -- (\x+0.5,\y-0.5)node [below] {\scalebox{0.6}{\indcolor{$i_2$}}};
    \draw[edge] (m) -- (\x-0.5,\y+0.5)node [above] {\scalebox{0.6}{\indcolor{$k_1$}}};
    \draw[edge] (m) -- (\x-0.5,\y-0.5)node [below] {\scalebox{0.6}{\indcolor{$i_1$}}};
    \draw[edge] (m) -- (\x+0.5,\y+0.5)node [above] {\scalebox{0.6}{\indcolor{$k_2$}}};
\end{tikzpicture}
=
 \begin{tikzpicture}[baseline=\baseline]
    \tikzset{tensor/.style = {minimum size = 0.5cm,shape = circle,thick,draw=black,inner sep = 0pt}, edge/.style   = {thick,line width=.4mm},every loop/.style={}}
    \def\y{0}
    \def\x{0}
    \draw  node[fill = black,circle,inner sep=0pt,minimum size=5pt] (m) at (\x,\y+0.4) {};
    \draw  node[fill = black,circle,inner sep=0pt,minimum size=5pt] (n) at (\x,\y-0.4) {};
    \draw[edge] (n) -- (m)node [midway, right] {\scalebox{0.6}{\indcolor{$i_1$}}};;
    \draw[edge](m) -- (\x,\y+1)node [above] {\scalebox{0.6}{\indcolor{$k_2$}}};
    \draw[edge](n) -- (\x,\y-1)node [below] {\scalebox{0.6}{\indcolor{$k_1$}}};
    \draw[edge](m) -- (\x+0.45,\y+0.65)node [above] {\scalebox{0.6}{\indcolor{$l_2$}}};
    \draw[edge](m) -- (\x-0.45,\y+0.65)node [above] {\scalebox{0.6}{\indcolor{$j_2$}}};
    \draw[edge](n) -- (\x+0.45,\y-0.65)node [above] {\scalebox{0.6}{\indcolor{$l_1$}}};;
    \draw[edge](n) -- (\x-0.45,\y-0.65)node [above] {\scalebox{0.6}{\indcolor{$j_1$}}};;
    \end{tikzpicture}
=
\sum_{i_1}\delta_{i_1j_1k_1l_1}\delta_{i_1j_2k_2l_2}.
\end{align*}
\item\label{prop:diagcopy}
对任意向量 $\vb\in\RR^n$，记 $\diag(\vb)\in\RR^{n\times n}$ 为以 $\vb$ 的各分量作对角元的对角矩阵；对任意矩阵 $\Ab\in\RR^{n\times n}$，记 $\diag(\Ab)\in\RR^{n}$ 为取出 $\Ab\in\RR^{n\times n}$ 全部对角元所得的向量。
于是
$$
    \diag(\vb) = 
     \begin{tikzpicture}[baseline=\baseline]
    \tikzset{tensor/.style = {minimum size = 0.5cm,shape = circle,thick,draw=black,inner sep = 0pt}, edge/.style   = {thick,line width=.4mm},every loop/.style={}}
    \def\y{0}
    \def\x{0}
    \draw  node[fill = black,circle,inner sep=0pt,minimum size=4pt] (m) at (\x,\y) {};
    \node[tensor,fill=blue!20!red!40!white] (v) at (\x,\y-0.85){\scalebox{0.85}{$\vb$}};
    \draw[edge] (m) -- (\x-0.6,0);
    \draw[edge] (m) -- (\x+0.6,0);
    \draw[edge] (m) -- (\x,\y-0.6); 
    \node[draw=none] at (\x+1,\y) {\scalebox{1}{\textcolor{black}{$=$}}};
    \node[tensor,fill=blue!20!red!40!white] (v) at (\x+2.5,\y){\scalebox{0.85}{}};
    \draw[edge] (v) -- (\x+3.5,\y);
    \draw[edge] (v) -- (\x+1.5,\y);
  \draw[edge=0.5] (v.south west) -- (v.north east);
\end{tikzpicture} 
~~\text{and}~~
\diag(\Ab) =\hspace{-0.7cm}
\begin{tikzpicture}[baseline=\baseline]
    \tikzset{tensor/.style = {minimum size = 0.5cm,shape = circle,thick,draw=black,inner sep = 0pt}, edge/.style = {thick,line width=.4mm},every loop/.style={}}
    \def\y{0}
    \def\x{0}
    \node[tensor,fill=red!50!white] (A) at (\x,\y+0.2){\scalebox{0.85}{$\Ab$}};
    \draw[edge] (A) to [out=-20,in=-160, distance=1.35cm] (A);
    \draw  node[fill = black,circle,inner sep=0pt,minimum size=4pt] (m) at (\x,\y-0.25) {};
    \draw[edge] (m) -- (\x,\y-0.65);
    \end{tikzpicture},
$$
其中
$
\begin{tikzpicture}[baseline=\baseline]
    \tikzset{tensor/.style = {minimum size = 0.5cm,shape = circle,thick,draw=black,inner sep = 0pt}, edge/.style = {thick,line width=.4mm},every loop/.style={}}
    \def\y{0}
    \def\x{0}
    \node[tensor,fill=blue!20!red!40!white] (v) at (\x+2.5,\y){\scalebox{0.85}{}};
    \draw[edge] (v) -- (\x+3.5,\y);
    \draw[edge] (v) -- (\x+1.5,\y);
  \draw[edge=0.5] (v.south west) -- (v.north east);
\end{tikzpicture} 
$ 
表示一个对角矩阵。
\begin{proof}
第一条：对任意 $i,j\in [n]$，有
\begin{align*}
    \begin{tikzpicture}[baseline=-15]
    \tikzset{tensor/.style = {minimum size = 0.5cm,shape = circle,thick,draw=black,inner sep = 0pt}, edge/.style   = {thick,line width=.4mm},every loop/.style={}}
    \def\y{0}
    \def\x{0}
    \draw  node[fill = black,circle,inner sep=0pt,minimum size=4pt] (m) at (\x,\y) {};
   \node[tensor,fill=blue!20!red!40!white] (v) at (\x,\y-0.85){\scalebox{0.85}{$\vb$}};
    \draw[edge] (m) -- (\x-0.6,0) node [very near end,left] {\scalebox{0.7}{\indcolor{$i$}}};
    \draw[edge] (m) -- (\x+0.6,0)node [very near end,right] {\scalebox{0.7}{\indcolor{$j$}}};;
    \draw[edge] (m) -- (\x,\y-0.6);
 \end{tikzpicture}
&=
 \sum_{k}
\left(\begin{tikzpicture}[baseline=-15]
    \tikzset{tensor/.style = {minimum size = 0.5cm,shape = circle,thick,draw=black,inner sep = 0pt}, edge/.style   = {thick,line width=.4mm},every loop/.style={}}
    \def\y{0}
    \def\x{0}
    \draw  node[fill = black,circle,inner sep=0pt,minimum size=4pt] (m) at (\x,\y) {};
    \draw[edge] (m) -- (\x-0.6,0)node [very near end, left] {\scalebox{0.7}{\indcolor{$i$}}};;
    \draw[edge] (m) -- (\x+0.6,0)node [very near end, right] {\scalebox{0.7}{\indcolor{$j$}}};
    \draw[edge] (m) -- (\x,\y-0.6)node [very near end, below]{\scalebox{0.7}{\indcolor{$k$}}};  \node[tensor,fill=blue!20!red!40!white] (v) at (\x+1.25,\y-0.8){\scalebox{0.85}{$\vb$}};
    \draw[edge] (v) -- (\x+1.25,0)node [very near end, above] {\scalebox{0.7}{\indcolor{$k$}}};;
 \end{tikzpicture}\right)\\
&= 
\sum_{k}\delta_{ijk}\vb_k = 
\delta_{ij}\vb_i
= 
\begin{cases}
\vb_i & \text{ if } i=j\\
0 &\text{ otherwise}
\end{cases}
= \text{diag}(\vb)_{ij}.
\end{align*}
第二条：对任意 $i\in [n]$，有
\begin{align*}
    \begin{tikzpicture}[baseline=\baseline]
    \tikzset{tensor/.style = {minimum size = 0.5cm,shape = circle,thick,draw=black,inner sep = 0pt}, edge/.style = {thick,line width=.4mm},every loop/.style={}}
    \def\y{0}
    \def\x{0}
    \node[tensor,fill=red!50!white] (A) at (\x,\y+0.2){\scalebox{0.85}{$\Ab$}};
    \draw[edge] (A) to [out=-20,in=-160, distance=1.35cm] (A);
    \draw  node[fill = black,circle,inner sep=0pt,minimum size=4pt] (m) at (\x,\y-0.25) {};
    \draw[edge] (m) -- (\x,\y-0.65)node [very near end, below] {\scalebox{0.7}{\indcolor{$i$}}};
\end{tikzpicture}\hspace{-0.7cm}
=
\sum_{j,k}\left(
\begin{tikzpicture}[baseline=\baseline]
    \tikzset{tensor/.style = {minimum size = 0.5cm,shape = circle,thick,draw=black,inner sep = 0pt}, edge/.style = {thick,line width=.4mm},every loop/.style={}}
    \def\y{0}
    \def\x{0}
    \node[tensor,fill=red!50!white] (A) at (\x,\y+0.4){\scalebox{0.85}{$\Ab$}};
    \draw  node[fill = black,circle,inner sep=0pt,minimum size=4pt] (m) at (\x,\y-0.25) {};
    \draw[edge] (m) -- (\x-0.7,\y-0.25)node [very near end,left] {\scalebox{0.7}{\indcolor{$k$}}};
    \draw[edge] (m) -- (\x+0.7,\y-0.25)node [very near end,right] {\scalebox{0.7}{\indcolor{$j$}}};
    \draw[edge] (m) -- (\x,\y-0.75)node [very near end, below] {\scalebox{0.7}{\indcolor{$i$}}};
    \draw[edge] (A) --(\x+0.7,\y+0.4)node [very near end,right] {\scalebox{0.7}{\indcolor{$j$}}};
    \draw[edge] (A) --(\x-0.7,\y+0.4)node [very near end,left] {\scalebox{0.7}{\indcolor{$k$}}};
    \end{tikzpicture}\right)
=
\sum_{j,k}\delta_{ijk}\Ab_{kj}
=
\Ab_{ii} = \diag(\Ab)
_{i}.
\end{align*}
\end{proof}
\item\label{prop:diag}
由此不难验证 $\diag(\vb)\diag(\ub)$ =
    \begin{tikzpicture}[baseline=\baseline]
    \tikzset{tensor/.style = {minimum size = 0.5cm,shape = circle,thick,draw=black,inner sep = 0pt}, edge/.style   = {thick,line width=.4mm},every loop/.style={}}
    \def\y{0}
    \def\x{0}
    \draw  node[fill = black,circle,inner sep=0pt,minimum size=4pt] (m) at (\x,\y) {};
    \node[tensor,fill=blue!20!red!40!white] (v) at (\x,\y-0.85){\scalebox{0.85}{$\vb$}};
    \draw[edge] (m) -- (\x-0.6,0);
    \draw[edge] (m) -- (\x,\y-0.6); 
    \draw  node[fill = black,circle,inner sep=0pt,minimum size=4pt] (m1) at (\x+1,\y) {};
    \node[tensor,fill=red!20!blue!40!white] (u) at (\x+1,\y-0.85){\scalebox{0.85}{$\ub$}};
    \draw[edge] (m1) -- (m);
    \draw[edge] (m1) -- (\x+1.6,0);
    \draw[edge] (m1) -- (\x+1,\y-0.6); 
\end{tikzpicture}.
\end{enumerate}
\end{proposition}

\subsection{张量网络中的矩阵分解}
与任何其他运算一样，张量分解也可以用张量网络图来描绘。本节将介绍矩阵分解的图形化图示。更高阶张量的一般分解将在第~\ref{chapter:tensor-decomposition}章中讨论。在张量网络中，分解是指把一个节点拆分为多个节点；我们称其逆运算为合并（merging），即把多个节点组合成一个节点。这一过程可以由下面的张量网络图清晰地说明：
\begin{align*}
\begin{tikzpicture}[baseline=\baseline]
    \tikzset{tensor/.style = {minimum size = 0.5cm,shape = circle,thick,draw=black,inner sep = 0pt}, edge/.style = {thick,line width=.4mm},every loop/.style={}}
    \def\y{0}
    \def\x{0}
    \node[tensor,fill=red!50!white] (A) at (\x,\y){\scalebox{0.85}{$\Ab$}};
    \node[tensor,fill=blue!50!white] (B) at (\x+2,\y){\scalebox{0.85}{$\Bb$}};
    \node[tensor,fill=blue!30!red!30!white] (C) at (\x+1,\y-1){\scalebox{0.85}{$\Cb$}};
    \node[tensor,fill=red!80!blue!60!white] (D) at (\x+3,\y-1){\scalebox{0.85}{$\Db$}};
    \draw[edge](A) -- (C) node [midway,above] {\scalebox{0.7}{\dimcolor{$r_1$}}};
    \draw[edge](C) -- (D) node [midway,above] {\scalebox{0.7}{\dimcolor{$r_2$}}};
    \draw[edge](B) -- (D) node [midway,above] {\scalebox{0.7}{\dimcolor{$r_3$}}};
    \draw[edge](B) -- (C) node [midway,above] {\scalebox{0.7}{\dimcolor{$r_3$}}};
    \draw[edge](A) -- (\x-0.8,\y) node [midway,above] {\scalebox{0.7}{\dimcolor{$d_1$}}};
    \draw[edge](B) -- (\x+2.8,\y) node [midway,above] {\scalebox{0.7}{\dimcolor{$d_2$}}};
    \node[draw=none]() at (\x+4.5,\y-0.15){\scalebox{1}{$\xrightarrow[]{\text{merging}}$}};
    \node[draw=none]() at 
    (\x+4.5,\y-0.75){\scalebox{1}{$\xleftarrow[\text{factorizing}]{}$}};
    \node[tensor,fill=red!30!blue!50!white] (M) at (\x+6.5,\y-0.5){\scalebox{0.85}{$\Mb$}};
    \draw[edge](M) -- (\x+7.5,\y-0.5) node [midway,above] {\scalebox{0.7}{\dimcolor{$d_2$}}};
    \draw[edge](M) -- (\x+5.5,\y-0.5) node [midway,above] {\scalebox{0.7}{\dimcolor{$d_1$}}};
\end{tikzpicture}
\end{align*}
因此，任何矩阵分解都可以用张量网络图来表示。后续章节将会看到，QR 分解和奇异值分解（SVD）是张量网络算法中常用于拆分节点的两种分解。这两种分解都把矩阵写成若干矩阵的乘积，其中一些因子满足正交性约束。我们将用视觉线索在张量网络图中表示这类正交性约束：
\paragraph{左、右标准正交矩阵与张量}
若把矩阵 $\Ub\in\RR^{m\times n}$ 的转置与它自身从左侧收缩得到单位矩阵~（$\Ub^\ts\Ub = \Ib_{n}$），则称该矩阵为\emph{左标准正交}矩阵。
类似地，若矩阵 $\Vb\in\RR^{m\times n}$ 与它的转置从右侧收缩得到单位矩阵
（$\Vb\Vb^\ts = \Ib_{m}$），则称其为\emph{右标准正交}矩阵。在张量网络图中，我们用半着色节点来表示这种正交结构，使得用未着色部分的边连接一个正交节点的两个副本时，结果就是单位矩阵。例如，左标准正交矩阵 $\Ub\in\RR^{m\times n}$ 将表示为
\begin{tikzpicture}[baseline=\baseline]
    \tikzset{tensor/.style = {minimum size = 0.5cm,shape = circle,thick,draw=black,inner sep = 0pt}, edge/.style   = {thick,line width=.4mm},every loop/.style={}}
    \def\y{0}
    \def\x{0};

    
    \node[tensor, path picture={
    \fill[fill=red!30!white](path picture bounding box.south east)
        -- 
    (path picture bounding box.north east) --
    (path picture bounding box.north west) -- cycle;
    }, shift={(\x,\y)}] (Ut) {$\Ub$};
    \draw[edge](Ut) -- (\x-1,\y) node [midway,above] {\scalebox{0.7}{\dimcolor{$m$}}};
    \draw[edge](Ut) -- (\x+1,\y) node [midway,above] {\scalebox{0.7}{\dimcolor{$n$}}};
\end{tikzpicture}，由此得到 
\begin{tikzpicture}[baseline=\baseline]
    \tikzset{tensor/.style = {minimum size = 0.5cm,shape = circle,thick,draw=black,inner sep = 0pt}, edge/.style   = {thick,line width=.4mm},every loop/.style={}}
    \def\y{0}
    \def\x{0};

    \node[tensor, path picture={
    \fill[fill=red!30!white](path picture bounding box.south west)
    -- 
    (path picture bounding box.north east) --
    (path picture bounding box.north west) -- cycle;
    }, shift={(\x,\y)}] (U) {$\Ub$};
    
    \node[tensor, path picture={
    \fill[fill=red!30!white](path picture bounding box.south east)
        -- 
    (path picture bounding box.north east) --
    (path picture bounding box.north west) -- cycle;
    }, shift={(\x+1,\y)}] (Ut) {$\Ub$};

    \draw[edge] (U) -- (Ut) node [midway,above] {\scalebox{0.7}{\dimcolor{$m$}}};
    \draw[edge](Ut) -- (\x+2,\y) node [midway,above] {\scalebox{0.7}{\dimcolor{$n$}}};
    \draw[edge](U) -- (\x-1,\y) node [midway,above] {\scalebox{0.7}{\dimcolor{$n$}}};
    \node[draw=none] () at (\x+2.3,\y) {\scalebox{0.8}{\dimcolor{$=$}}};
    \draw[edge](\x+2.5,\y) -- (\x+4,\y)node [midway,above] {\scalebox{0.7}{\dimcolor{$n$}}};
    \node[draw=none] () at (\x+4.3,\y) {\scalebox{0.8}{\dimcolor{$=$}}};
\end{tikzpicture} $\Ib_n$。 
类似地，右标准正交矩阵 $\Vb\in\RR^{m\times n}$ 将画成 
\begin{tikzpicture}[baseline=\baseline]
    \tikzset{tensor/.style = {minimum size = 0.5cm,shape = circle,thick,draw=black,inner sep = 0pt}, edge/.style   = {thick,line width=.4mm},every loop/.style={}}
    \def\y{0}
    \def\x{0};

    \node[tensor, path picture={
    \fill[fill=red!30!white](path picture bounding box.south west)
    -- 
    (path picture bounding box.north east) --
    (path picture bounding box.north west) -- cycle;
    }, shift={(\x,\y)}] (Vt) {$\Vb$};
    \draw[edge](Vt) -- (\x-1,\y) node [midway,above] {\scalebox{0.7}{\dimcolor{$m$}}};
    \draw[edge](Vt) -- (\x+1,\y) node [midway,above] {\scalebox{0.7}{\dimcolor{$n$}}};
\end{tikzpicture}，而我们有
\begin{tikzpicture}[baseline=\baseline]
    \tikzset{tensor/.style = {minimum size = 0.5cm,shape = circle,thick,draw=black,inner sep = 0pt}, edge/.style   = {thick,line width=.4mm},every loop/.style={}}
    \def\y{0}
    \def\x{0};

    \node[tensor, path picture={
    \fill[fill=red!30!white](path picture bounding box.south west)
    -- 
    (path picture bounding box.north east) --
    (path picture bounding box.north west) -- cycle;
    }, shift={(\x,\y)}] (U) {$\Vb$};
    
    \node[tensor, path picture={
    \fill[fill=red!30!white](path picture bounding box.south east)
        -- 
    (path picture bounding box.north west) --
    (path picture bounding box.north east) -- cycle;
    }, shift={(\x+1,\y)}] (Ut) {$\Vb$};

    \draw[edge] (U) -- (Ut) node [midway,above] {\scalebox{0.7}{\dimcolor{$n$}}};
    \draw[edge](Ut) -- (\x+2,\y) node [midway,above] {\scalebox{0.7}{\dimcolor{$m$}}};
    \draw[edge](U) -- (\x-1,\y) node [midway,above] {\scalebox{0.7}{\dimcolor{$m$}}};
    \node[draw=none] () at (\x+4.3,\y) {\scalebox{0.8}{\dimcolor{$=$}}};
    \node[draw=none] () at (\x+2.3,\y) {\scalebox{0.8}{\dimcolor{$=$}}};
    \draw[edge](\x+2.5,\y) -- (\x+4,\y)node [midway,above] {\scalebox{0.7}{\dimcolor{$m$}}};
\end{tikzpicture} $\Ib_m$。


注意，若矩阵 $\Ub$ 左标准正交，则 $\Ub\Ub^\ts$ 未必是单位矩阵，即 
$$
\begin{tikzpicture}[baseline=\baseline]
    \tikzset{tensor/.style = {minimum size = 0.5cm,shape = circle,thick,draw=black,inner sep = 0pt}, edge/.style   = {thick,line width=.4mm},every loop/.style={}}
    \def\y{0}
    \def\x{0};

    \node[tensor, path picture={
    \fill[fill=red!30!white](path picture bounding box.south east)
    -- 
    (path picture bounding box.north west) --
    (path picture bounding box.north east) -- cycle;
    }, shift={(\x,\y)}] (U) {$\Ub$};
    
    \node[tensor, path picture={
    \fill[fill=red!30!white](path picture bounding box.south west)
        -- 
    (path picture bounding box.north east) --
    (path picture bounding box.north west) -- cycle;
    }, shift={(\x+1,\y)}] (Ut) {$\Ub$};

    \draw[edge] (U) -- (Ut) node [midway,above] {\scalebox{0.7}{\dimcolor{$n$}}};
    \draw[edge](Ut) -- (\x+2,\y) node [midway,above] {\scalebox{0.7}{\dimcolor{$m$}}};
    \draw[edge](U) -- (\x-1,\y) node [midway,above] {\scalebox{0.7}{\dimcolor{$m$}}};
    \node[draw=none] () at (\x+4.3,\y) {\scalebox{0.8}{\dimcolor{$=$}}};
    \node[draw=none] () at (\x+2.3,\y) {\scalebox{0.8}{\dimcolor{$\neq$}}};
    \draw[edge](\x+2.5,\y) -- (\x+4,\y)node [midway,above] {\scalebox{0.7}{\dimcolor{$m$}}};
\end{tikzpicture}\Ib_m,$$ 只有当 $m=n$ 时才会如此。 



借助这些视觉线索，QR 分解与 SVD 对应的张量网络如下所示：
\begin{itemize}
    \item\textbf{QR 分解} 给定 $\Ab\in\RR^{d_1\times d_2}$，其~（秩揭示）QR 分解为 
    
$$
\begin{tikzpicture}[baseline=\baseline]
    \tikzset{tensor/.style = {minimum size = 0.5cm,shape = circle,thick,draw=black,inner sep = 0pt}, edge/.style   = {thick,line width=.4mm},every loop/.style={}}
    \def\y{-0.2}
    \def\x{0}
    \node[tensor,fill=red!30!white] (A) at (\x,\y+0.2){$\Ab$};
    \draw[edge] (A) -- (\x-0.8,\y+0.2)  node [midway,above] {\scalebox{0.7}{\dimcolor{$d_1$}}};
    \draw[edge] (A) -- (\x+0.8,\y+0.2)  node [midway,above] {\scalebox{0.7}{\dimcolor{$d_2$}}};
    \node[draw = none] at (\x+1.2,\y+0.2){$=$};
    \node[tensor, path picture={
    \fill[fill=red!30!white](path picture bounding box.south east)
    -- 
    (path picture bounding box.north west) --
    (path picture bounding box.north east) -- cycle;
    }, shift={(\x+2.5,\y+0.2)}] (Q) {$\Qb$};
    \node[tensor,fill=blue!30!white] (R) at (\x+3.5,\y+0.2){$\Rb$};
    \draw[edge] (Q) -- (R)  node [midway,above] {\scalebox{0.7}{\dimcolor{$R$}}};
    \draw[edge] (Q) -- (\x+1.7,\y+0.2)  node [midway,above] {\scalebox{0.7}{\dimcolor{$d_1$}}};
    \draw[edge] (R) -- (\x+4.3,\y+0.2)  node [midway,above] {\scalebox{0.7}{\dimcolor{$d_2$}}};
\end{tikzpicture},$$
其中 
$\begin{tikzpicture}[baseline=\baseline]
    \tikzset{tensor/.style = {minimum size = 0.5cm,shape = circle,thick,draw=black,inner sep = 0pt}, edge/.style   = {thick,line width=.4mm},every loop/.style={}}
    \def\y{0}
    \def\x{0}
    \node[tensor, path picture={
    \fill[fill=red!30!white](path picture bounding box.south east)
    -- 
    (path picture bounding box.north west) --
    (path picture bounding box.north east) -- cycle;
    }, shift={(\x,\y)}] (Q) {$\Qb$};
    \draw[edge] (Q) -- (\x+0.7,\y)  node [midway,above] {\scalebox{0.7}{\dimcolor{$R$}}};
    \draw[edge] (Q) -- (\x-0.7,\y)  node [midway,above] {\scalebox{0.7}{\dimcolor{$d_1$}}};
\end{tikzpicture}
\in\RR^{d_1\times R}$
是左标准正交矩阵，即 $\Qb^\ts\Qb = \Ib_R$；$\Rb\in\RR^{R\times d_2}$ 为上三角矩阵~（这一点在张量网络表示中并不明显），$R$ 是 $\Ab$ 的秩。 
\item \textbf{奇异值分解~（SVD）} 给定 $\Ab\in\RR^{d_1\times d_2}$，其~（紧凑）SVD 为 
\begin{align*}
\begin{tikzpicture}[baseline=\baseline]
    \tikzset{tensor/.style = {minimum size = 0.5cm,shape = circle,thick,draw=black,inner sep = 0pt}, edge/.style   = {thick,line width=.4mm},every loop/.style={}}
    \def\y{0}
    \def\x{0}
    \node[tensor,fill=red!30!white] (A) at (\x,\y){$\Ab$};
    \draw[edge] (A) -- (\x-0.8,\y)  node [midway,above] {\scalebox{0.7}{\dimcolor{$d_1$}}};
    \draw[edge] (A) -- (\x+0.8,\y)  node [midway,above] {\scalebox{0.7}{\dimcolor{$d_2$}}};
    \node[draw = none] at (\x+1.35,\y){$=$};
    \node[tensor, path picture={
    \fill[fill=red!30!white](path picture bounding box.south east)
    -- 
    (path picture bounding box.north west) --
    (path picture bounding box.north east) -- cycle;
    }, shift={(\x+3,\y)}] (U) {$\Ub$};
    \node[tensor,fill=blue!20!red!40!white] (sigma) at (\x+4.2,\y){\scalebox{0.85}{}};
    \node[tensor, path picture={
    \fill[fill=red!30!white](path picture bounding box.south west)
        -- 
    (path picture bounding box.north west) --
    (path picture bounding box.north east) -- cycle;
    }, shift={(\x+5.2,\y)}] (V) {$\Vb$};
    \draw[edge=0.5] (sigma.south west) -- (sigma.north        east);
    \draw[edge] (U) -- (sigma) node [midway,above] {\scalebox{0.7}{\dimcolor{$R$}}};
    \draw[edge] (sigma) -- (V)node [midway,above] {\scalebox{0.7}{\dimcolor{$R$}}};
    \draw[edge](U) -- (\x+2,\y) node [midway,above] {\scalebox{0.7}{\dimcolor{$d_1$}}};
    \draw[edge](V) -- (\x+6.2,\y) node [midway,above] {\scalebox{0.7}{\dimcolor{$d_2$}}};
\end{tikzpicture}~,
\end{align*}
其中 $\Ub\in\RR^{d_1\times R}$ 左标准正交~（$\Ub^\ts\Ub= \Ib_R$），$\Vb\in\RR^{R\times d_2}$ 右标准正交~（$\Vb\Vb^\ts = \Ib_{R}$），$\Sigmab\in\RR^{R\times R}
=
\begin{tikzpicture}[baseline=\baseline]
    \tikzset{tensor/.style = {minimum size = 0.5cm,shape = circle,thick,draw=black,inner sep = 0pt}, edge/.style   = {thick,line width=.4mm},every loop/.style={}}
    \def\y{0}
    \def\x{0}
    \node[tensor,fill=blue!20!red!40!white] (sigma) at (\x,\y){\scalebox{0.85}{}};
    \draw[edge=0.5] (sigma.south west) -- (sigma.north        east);
    \draw[edge](sigma) -- (\x-0.7,\y);
    \draw[edge](sigma) --(\x+0.7,\y);
\end{tikzpicture}$ 为对角矩阵，$R$ 是 $\Ab$ 的秩。 
\end{itemize}

设 $\Ab$ 具有上一图中所表示的 SVD，则我们可以用张量网络为恒等式 $\tr(\Ab^\ts\Ab) = \sum_{i}\sigma_i^2$ 给出一个简短的证明，其中 $\sigma_i$ 是矩阵 $\Ab\in\RR^
{m\times n}$：
$$
\tr(\Ab^\ts\Ab) 
\begin{tikzpicture}[baseline=\baseline]
   \tikzset{tensor/.style = {minimum size = 0.5cm,shape = circle,thick,draw=black,inner sep = 0pt}, edge/.style = {thick,line width=.4mm},every loop/.style={}}
    \def\y{0}
    \def\x{0}
    \node[draw=none] () at (\x-0.7,\y) {\scalebox{1}{$=$}};
    \node[tensor,fill=red!50!white] (A) at (\x,\y){\scalebox{0.85}{$\Ab$}};
    \node[tensor,fill=red!50!white] (B) at (\x+1,\y){\scalebox{0.85}{$\Ab$}};
    \draw[edge] (A) -- (B) node [midway,above]{\scalebox{0.7}{\dimcolor{$n$}}};
    \draw[edge] (A) to [out=155,in=25, distance=1cm] (B);
    \node[draw=none] () at (\x+0.5,\y+0.65) {\scalebox{0.6}{\dimcolor{$m$}}};
    \node[draw=none] () at (\x+1.7,\y) {\scalebox{1}{$=$}};
    \end{tikzpicture}
\begin{tikzpicture}[baseline=\baseline]
    \tikzset{tensor/.style = {minimum size = 0.5cm,shape = circle,thick,draw=black,inner sep = 0pt}, edge/.style   = {thick,line width=.4mm},every loop/.style={}}
    \def\y{0}
    \def\x{0}
   
    \node[tensor, path picture={
    \fill[fill=red!30!white](path picture bounding box.south east)
    -- 
    (path picture bounding box.north west) --
    (path picture bounding box.north east) -- cycle;
    }, shift={(\x,\y+0.5)}] (U) {$\Ub$};
    \node[tensor,fill=blue!20!red!40!white] (sigma) at (\x+1,\y+0.5){\scalebox{0.85}{}};
    \node[tensor, path picture={
    \fill[fill=red!30!white](path picture bounding box.south west)
        -- 
    (path picture bounding box.north west) --
    (path picture bounding box.north east) -- cycle;
    }, shift={(\x+2,\y+0.5)}] (V) {$\Vb$};
    \draw[edge=0.5] (sigma.south west) -- (sigma.north        east);
    \draw[edge] (U) -- (sigma) node [midway,above] {\scalebox{0.7}{\dimcolor{$R$}}};
    \draw[edge] (sigma) -- (V)node [midway,above] {\scalebox{0.7}{\dimcolor{$R$}}};
    
    \node[tensor, path picture={
    \fill[fill=red!30!white](path picture bounding box.south east)
    -- 
    (path picture bounding box.north west) --
    (path picture bounding box.north east) -- cycle;
    }, shift={(\x,\y-0.5)}] (U1) {$\Ub$};
    \node[tensor,fill=blue!20!red!40!white] (sigma1) at (\x+1,\y-0.5){\scalebox{0.85}{}};
    \node[tensor, path picture={
    \fill[fill=red!30!white](path picture bounding box.south west)
        -- 
    (path picture bounding box.north west) --
    (path picture bounding box.north east) -- cycle;
    }, shift={(\x+2,\y-0.5)}] (V1) {$\Vb$};
    \draw[edge=0.5] (sigma1.south west) -- (sigma1.north        east);
    \draw[edge] (U1) -- (sigma1) node [midway,above] {\scalebox{0.7}{\dimcolor{$R$}}};
    \draw[edge] (sigma1) -- (V1) node [midway,above] {\scalebox{0.7}{\dimcolor{$R$}}};
    \draw[edge] (U) to [out=-150,in=160, distance=0.5cm] (U1);
    \draw[edge] (V) to [out=-30,in=20, distance=0.5cm] (V1);
    \node[draw=none] () at (\x-0.8,\y) {\scalebox{0.7}{\dimcolor{$m$}}};
    \node[draw=none] () at (\x+2.8,\y) {\scalebox{0.7}{\dimcolor{$n$}}};
    \node[draw=none] () at (\x+3.5,\y) {\scalebox{1}{$=$}};
    \node[tensor,fill=blue!20!red!40!white] (sigma2) at (\x+4.5,\y-0.5){\scalebox{0.85}{}};
    \node[tensor,fill=blue!20!red!40!white] (sigma3) at (\x+4.5,\y+0.5){\scalebox{0.85}{}};
    \draw[edge=0.5] (sigma2.south west) -- (sigma2.north        east);
    \draw[edge=0.5] (sigma3.south west) -- (sigma3.north        east);
    \draw[edge] (sigma2) to [out=-150,in=160, distance=0.5cm] (sigma3);
    \draw[edge] (sigma2) to [out=-30,in=20, distance=0.5cm] (sigma3);
    \node[draw=none] () at (\x+4.2,\y) {\scalebox{0.7}{\dimcolor{$R$}}};
    \node[draw=none] () at (\x+4.9,\y) {\scalebox{0.7}{\dimcolor{$R$}}};
\end{tikzpicture}
= \tr(\Sigmab^2) = \sum_{i}\sigma_i^2,
$$
这里用到了 $\Ub$ 与 $\Vb$ 的左、右标准正交性质。

\section{张量运算}\label{chapter:operations}
回顾一下，张量 $\Tt \in \RR^{d_1 \times \cdots \times d_N}$ 可以看作一个 $N$ 阶多维数组：$N$ 阶（亦称 $N$ 路）张量有 $N$ 个模，每个模对应一个维度 \citep{kolda2009tensor}。
\paragraph{张量纤维} 对任意模 $i$（其中 $i = 1, \dots, N$），将除第 $i$ 个以外的指标全部固定，便得到张量的\textit{纤维}。例如，矩阵的一列对应模 1 纤维，而矩阵的一行对应模 2 纤维。在三阶张量 $\Tt \in \RR^{d_1 \times d_2 \times d_3}$ 中，模 1 纤维有 $d_2d_3$ 个，它们是尺寸为 $d_1$ 的向量，即对 $i_2 \in [d_2]$ 与 $i_3 \in [d_3]$，有 $\Tt_{:,i_2,i_3} \in \RR^{d_1}$。冒号表示第一个指标变化，而 $i_2$ 与 $i_3$ 保持固定。纤维是向量。
\paragraph{张量切片} 张量的切片是将除两个指标以外的所有指标固定后得到的二维数组（矩阵）。对三阶张量 $\Tt \in \RR^{d_1 \times d_2 \times d_3}$，分别有水平切片、侧面切片与正面切片，记为 $\Tt_{i_1,:,:} \in \RR^{d_2 \times d_3}$、$\Tt_{:,i_2,:} \in \RR^{d_1 \times d_3}$ 与 $\Tt_{:,:,i_3} \in \RR^{d_1 \times d_2}$。因此，切片是矩阵。
\paragraph{标量乘法} 
设 $\lambda \in \RR$ 为标量。标量乘积 $\lambda\Tt$ 
 仍是维度为 $d_1 \times \cdots \times d_N$ 的张量，逐元素定义为：
$
    (\lambda\Tt)_{i_1 \cdots i_N} = \lambda \cdot \Tt_{i_1 \cdots i_N}$
对所有可取的指标 $i_k\in[d_k]$ 与 $k\in[N]$。
\paragraph{张量求和} 设 $\St\in\RR^{d_1\times\cdots\times d_N}$。张量 $\Tt$ 与 $\St$ 之和仍是同维度的张量 $\Tt+\St\in\RR^{d_1\times\cdots\times d_N}$，逐元素定义为：$(\Tt+\St)_{i_1,\cdots,i_N} = \Tt_{i_1,\cdots,i_N} + \St_{i_1,\cdots,i_N}$，对所有可取的指标 $i_k\in[d_k]$ 与 $k\in[N]$。
\subsection{张量的置换与重排}
置换与重排是张量的两种基本运算。
\begin{itemize}
    \item \textbf{置换} 在不改变模的个数的情况下重新排列张量的指标。置换的一个常见例子是矩阵转置。
    \item  \textbf{重排} 把多个指标合并为较大的指标，或把一个较大的指标拆分为多个较小的指标，从而减少或增加指标的总数，同时保持张量的总体尺寸不变。
\end{itemize}
下面介绍矩阵化（matricization）与向量化（vectorization）这两种常见的张量重排运算。
\begin{definition}（矩阵化）\label{matricizitation:def}
设 $\Tt\in\RR^{d_1\times d_2\times \cdots\times d_N}$。对 $n\in[N]$，把 $\Tt$ 的所有模 $n$ 纤维取出并按列堆叠，即可展开为矩阵 $\Tt_{(n)}\in\RR^{d_n\times d_1d_2\cdots d_{n-1}d_{n+1}\cdots d_N}$，称为模 $n$ 矩阵化。例如，张量 $\Tt\in\RR^{d_1\times d_2\times d_3}$ 沿模 $1$ 的矩阵化记作 $\Tt_{(1)}$，它是如下所示的 $d_1\times d_2d_3$ 矩阵
$$
    \begin{bmatrix}
    \vline & \vline & \vline  & \vline \\
    \Tt_{:,1,1} & {\Tt}_{:,1,2}  & \dots  & {\Tt}_{:,d_2,d_3} \\
    \vline & \vline & \vline  & \vline 
    \end{bmatrix}
\in\RR^{d_1\times d_2d_3}.
$$

模 2 与模 3 矩阵化 $\Tt_{(2)}\in\RR^{d_2\times d_1d_3}$ 和 $\Tt_{(3)}\in\RR^{d_3\times d_1d_2}$ 的定义类似。
\end{definition}

注意，在模 1 矩阵化中，对应于 $\Tt$ 的第二和第三个模的两个指标被归并在一起，形成一个取值从 1 到 $d_2d_3$ 的新指标。在前面的定义中，我们选择按字典序排列这个新指标，例如当 $d_2=d_3=3$ 时为 $11,12,13,21,22,23,31,32,33$。矩阵化中各列的排列次序有时也采用别的约定~（例如 \citep{kolda2009tensor} 中的逆字典序）。一般说来，具体采用哪一种排列
并不重要，但在相关的计算中必须保持一致~（详见 \citep{kolda2006multilinear}）。



在张量网络图中，我们通过把相应的边合并到一起来表示这种指标归并：
$$\Tt_{(1)} = 
\left(\begin{tikzpicture}[baseline=\baseline]
    \tikzset{tensor/.style = {minimum size = 0.5cm,shape = circle,thick,draw=black,inner sep = 0pt}, edge/.style = {thick,line width=.4mm},every loop/.style={}}
    \def\y{0}
    \def\x{0}
    \node[tensor,fill=red!50!white] (T) at (\x,\y){\scalebox{0.85}{$\Tt$}};
    \draw[edge] (T) -- (\x+0.7,\y) node [midway,above]{\scalebox{0.7}{\dimcolor{$d_2$}}};
    \draw[edge] (T) -- (\x-0.7,\y) node [midway,above]{\scalebox{0.7}{\dimcolor{$d_1$}}};
    \draw[edge] (T) -- (\x,\y-0.7) node [midway,left]{\scalebox{0.7}{\dimcolor{$d_3$}}}; 
\end{tikzpicture}\right)_{(1)}
=
\begin{tikzpicture}[baseline=\baseline]
    \tikzset{tensor/.style = {minimum size = 0.5cm,shape = circle,thick,draw=black,inner sep = 0pt}, edge/.style = {thick,line width=.4mm},every loop/.style={}}
    \def\y{0}
    \def\x{0}
    \node[tensor,fill=red!50!white] (A) at (\x,\y){\scalebox{0.85}{$\Tt$}};
    \draw[edge] (\x+0.2,\y+0.15) -- (\x+0.5,\y+0.15)-- (\x+0.5,\y+0.05) -- (\x+1,\y+0.05) node [midway,above]
    {\scalebox{0.7}{\dimcolor{$d_2$}}};
    \draw[edge] (\x+0.2,\y-0.15) -- (\x+0.5,\y-0.15)-- (\x+0.5,\y-0.05)-- (\x+1,\y-0.05) node [midway,below]
    {\scalebox{0.7}{\dimcolor{$d_3$}}};
    \draw[edge] (A) -- (\x-0.7,\y)node [midway,above]
    {\scalebox{0.7}{\dimcolor{$d_1$}}};
\end{tikzpicture}
=
\begin{tikzpicture}[baseline=\baseline]
    \tikzset{tensor/.style = {minimum size = 0.5cm,shape = circle,thick,draw=black,inner sep = 0pt}, edge/.style = {thick,line width=.4mm},every loop/.style={}}
    \def\y{0}
    \def\x{0}
    \node[tensor,fill=red!50!white] (T) at (\x,\y){\scalebox{0.85}{$\Tt$}};
    \draw[edge](T) -- (\x+1,\y) node [midway, above]{\scalebox{0.7}{\dimcolor{$d_2d_3$}}};
    \draw[edge](T) -- (\x-1,\y) node [midway, above]{\scalebox{0.7}{\dimcolor{$d_1$}}};
\end{tikzpicture}
\in\RR^{d_1\times d_2d_3}.
$$
可以看到，形如
$\begin{tikzpicture}[baseline=\baseline]
    \tikzset{tensor/.style = {minimum size = 0.5cm,shape = circle,thick,draw=black,inner sep = 0pt}, edge/.style = {thick,line width=.4mm},every loop/.style={}}
    \def\y{0}
    \def\x{0}
    \draw[edge] (\x+0.15,\y+0.2) -- (\x+0.5,\y+0.2) -- (\x+0.5,\y+0.1) -- (\x+1,\y+0.1) node [midway,above]
    {\scalebox{0.7}{\dimcolor{$d_2$}}};
    \draw[edge] (\x+0.15,\y-0.1) -- (\x+0.5,\y-0.1)-- (\x+0.5,\y)-- (\x+1,\y) 
node [midway,below]
    {\scalebox{0.7}{\dimcolor{$d_3$}}};
    \end{tikzpicture}
$
的归并腿在张量网络图中表示一条大小为 $d_2d_3$ 的边。


一般地，矩阵化的概念可以推广到 $\Tt$ 的模的任一子集 $I\subset [N]$：把属于 $I$ 的模映为行，把不属于 $I$ 的模映为列，从而得到大小为 $\prod_{i\in I} d_i \times \prod_{j\in [N]\setminus I} d_j$ 的矩阵 $\Tt_{(I)}$。例如，对六阶张量 $\Tt\in\RR^{d_1\times\dots\times d_6}$，可以按如下方式归并指标：
\begin{align*}
\Tt_{i_1,\dots,i_6} 
&=
\begin{tikzpicture}[baseline=\baseline]
    \tikzset{tensor/.style = {minimum size = 0.5cm,shape = circle,thick,draw=black,inner sep = 0pt}, edge/.style = {thick,line width=.4mm},every loop/.style={}}
    \def\y{0}
    \def\x{0}
    \node[tensor,fill=red!20!blue!40!white] (T) at (\x,\y){\scalebox{0.85}{$\Tt$}};
    \draw[edge] (T) -- (\x+0.85,\y) node [very near end,right]
    {\scalebox{0.7}{\indcolor{$i_4$}}};
    \draw[edge] (T) -- (\x,\y-0.85) node [very near end, below]
    {\scalebox{0.7}{\indcolor{$i_6$}}};
    \draw[edge] (T) -- (\x-0.85,\y)node [very near end,left]
    {\scalebox{0.7}{\indcolor{$i_1$}}};
    \draw[edge] (T) -- (\x,\y+0.85) node [very near end,above]
    {\scalebox{0.7}{\indcolor{$i_3$}}};
    \draw[edge] (T) -- (\x-0.75,\y+0.5) node [pos=1.3]
    {\scalebox{0.7}{\indcolor{$i_2$}}};
    \draw[edge] (T) -- (\x+0.75,\y-0.5) node [pos=1.3]
    {\scalebox{0.7}{\indcolor{$i_5$}}};
    \node[draw=none] at (\x+2,\y+0.2){$\underrightarrow{I=\{1,2,6\}}$};
    \end{tikzpicture}
\begin{tikzpicture}[baseline=\baseline]
    \tikzset{tensor/.style = {minimum size = 0.5cm,shape = circle,thick,draw=black,inner sep = 0pt}, edge/.style = {thick,line width=.4mm},every loop/.style={}}
    \def\y{0}
    \def\x{0}
    \node[tensor,fill=red!20!blue!40!white] (T) at (\x,\y){\scalebox{0.85}{$\Tt$}};
    \draw[edge](T) -- (\x+0.85,\y)node [very near end, right]{\scalebox{0.7}{\indcolor{$i_5$}}};
    \draw[edge](\x+0.16,\y+0.2) -- (\x+0.5,\y+0.2) -- (\x+0.5,\y+0.1) -- (\x+0.85,\y+0.1)node [very near end, above]{\scalebox{0.7}{\indcolor{$i_3$}}};
    \draw[edge](\x+0.16,\y-0.2) -- (\x+0.5,\y-0.2) -- (\x+0.5,\y-0.1) -- (\x+0.85,\y-0.1)node [very near end, below]{\scalebox{0.7}{\indcolor{$i_4$}}};
    \draw[edge](T) -- (\x-0.85,\y)node [very near end, left]{\scalebox{0.7}{\indcolor{$i_2$}}};
    \draw[edge](\x-0.16,\y+0.2) -- (\x-0.5,\y+0.2) -- (\x-0.5,\y+0.1) -- (\x-0.85,\y+0.1)node [very near end, above]{\scalebox{0.7}{\indcolor{$i_6$}}};
    \draw[edge](\x-0.16,\y-0.2) -- (\x-0.5,\y-0.2) -- (\x-0.5,\y-0.1) -- (\x-0.85,\y-0.1)node [very near end, below]{\scalebox{0.7}{\indcolor{$i_1$}}};
\end{tikzpicture}\\
&=
(\Tt_{(\{1,2,6\})})_{i_1i_2i_6 , i_3i_4i_5}\in\RR^{d_1d_2d_6\times d_3d_4d_5}.
\end{align*}
这一概念虽然直观，写成形式化的记法却很繁琐：张量中属于 $I$ 的指标 $i_{r_1},i_{r_2},\cdots, i_{r_{|I|}}$（按降序排列）被映为行指标 $
j = 1+\sum_{l=1}^{|I|}
\left[(i_{r_l}-1)\prod_{m=1}^{l-1}I_{r_m}\right]
$，而其余指标 $i_{c_1},i_{c_2},\cdots, i_{c_{N-|I|}}$（同样按降序排列）被映为列指标 $k  = 1+\sum_{n=1}^{N - |I|}
\left[(i_{c_n}-1)\prod_{s=1}^{n-1}I_{r_s}\right]$。借助张量网络，我们只需把这类 
重排运算表示为边的归并，就能避开这种繁琐。 






与前面类似，张量既能转化为矩阵，也能转化为向量。这一过程称为向量化：
\begin{definition}（向量化）\label{def:vectorization}
设 $\Tt\in\RR^{d_1\times d_2\times\cdots\times d_N}$。$\Tt$ 的\emph{向量化}记为 $\vectorize(\Tt)\in\RR^{d_1d_2\cdots d_N}$，它是把自身的模 $N$ 纤维拼接起来所得到的向量\footnote{按字典序拼接。}。例如，张量 $\Tt\in\RR^{d_1\times d_2\times d_3}$ 的向量化为：
$$\vectorize(\Tt) =    
[(\Tt_{1,1,:})\ \  (\Tt_{1,2,:})\ \  \cdots \ \ (\Tt_{d_1,d_2,:})]^\ts
=
\begin{tikzpicture}[baseline=\baseline]
    \tikzset{tensor/.style = {minimum size = 0.5cm,shape = circle,thick,draw=black,inner sep = 0pt}, edge/.style = {thick,line width=.4mm},every loop/.style={}}
    \def\y{0}
    \def\x{0}
    \node[tensor,fill=blue!10!white] (A) at (\x,\y){\scalebox{0.85}{$\Tt$}};
    \draw[edge] (\x+0.15,\y+0.2) -- (\x+0.5,\y+0.2) -- (\x+0.5,\y+0.1) -- (\x+1,\y+0.1) node [midway,above]
    {\scalebox{0.7}{\dimcolor{$d_1$}}};
     \draw[edge] (A) -- (\x+1,\y)node [right]
    {\scalebox{0.7}{\dimcolor{$d_2$}}};

    \draw[edge] (\x+0.15,\y-0.2) -- (\x+0.5,\y-0.2)-- (\x+0.5,\y-0.1)-- (\x+1,\y-0.1) node [midway,below]
    {\scalebox{0.7}{\dimcolor{$d_3$}}};
    \end{tikzpicture}\in\RR^{d_1d_2d_3}.
$$
\end{definition}

向量化是一种展开操作，它把任意阶的张量化为向量。与矩阵化类似，在张量网络中，形如
$
\begin{tikzpicture}[baseline=\baseline]
    \tikzset{tensor/.style = {minimum size = 0.5cm,shape = circle,thick,draw=black,inner sep = 0pt}, edge/.style = {thick,line width=.4mm},every loop/.style={}}
    \def\y{0}
    \def\x{0}
    \draw[edge] (\x+0.15,\y+0.2) -- (\x+0.5,\y+0.2) -- (\x+0.5,\y+0.1) -- (\x+1,\y+0.1) node [midway,above]
    {\scalebox{0.7}{\textcolor{gray}{$d_1$}}};
    \draw[edge] (\x+0.15,\y) -- (\x+1,\y)node [right]
    {\scalebox{0.7}{\textcolor{gray}{$d_2$}}};
    \draw[edge] (\x+0.15,\y-0.2) -- (\x+0.5,\y-0.2)-- (\x+0.5,\y-0.1)-- (\x+1,\y-0.1) 
node [midway,below]
    {\scalebox{0.7}{\dimcolor{$d_3$}}};
    \end{tikzpicture}
$ 的边表示一条大小为 $d_1d_2d_3$ 的边。一般而言，全文中汇聚于一点的边表示一条边，其大小等于所有相关边大小之积。注意，张量的向量化等于该张量模 $1$ 矩阵化的向量化\footnote{为与本书所用的向量化定义保持一致，矩阵的向量化由其各行拼接而成，这有时称为\emph{行主序向量化}。注意，这正是 \textsf{Python} 中 \textsf{numpy} 包默认的向量化操作 \textsf{ravel()}。}，即对 $\Tt\in\RR^{d_1\times d_2\cdots \times d_N}$，有 $\vectorize(\Tt) = \vectorize(\Tt_{(1)})$。




\subsection{乘积}\label{sec:product}
张量可以用不同类型的运算组合起来，即进行各种\textit{乘法}。本节给出各类乘积运算的图示。
\paragraph{模 $n$ 张量乘积}\label{par:mode-n-tensor-product}
模 $n$ 乘积是张量沿某个特定的模与矩阵相乘，可视为矩阵乘法的推广。更确切地说，模 $n$ 乘积把矩阵与矩阵的乘法推广到某一个特定的模上，其中包括张量与矩阵、以及张量与向量之间的模 $n$ 乘积。
\begin{enumerate}
\item \textbf{模-n 乘积~（矩阵）。}\label{subsec:mode-n}
    张量 $\Xt\in\RR^{d_1\times\cdots\times d_N}$ 与矩阵 $\Mb\in\RR^{m\times d_n}$ 的模 $n$ 乘积记为
    $\Xt\ttm n \Mb \in \RR^{d_1\times \cdots\times d_{n-1}\times m\times d_{n+1}\times\cdots\times d_N}$，其定义为张量与矩阵沿张量的模 $n$、沿矩阵的第二模~（列指标）收缩，即 
    $$
  \Xt\ttm n \Mb
=
\begin{tikzpicture}[baseline=-0.5ex]
    \tikzset{tensor/.style = {minimum size = 0.5cm,shape = circle,thick,draw=black,inner sep = 0pt}, edge/.style = {thick,line width=.4mm},every loop/.style={}}
    \def\x{6}
    \node[tensor,fill=blue!20!white] (M) at (\x,-1.05) {$\Mb$};
    \node[tensor,fill=green!50!red!50!white] (At) at (\x,0) {$\Xt$};
    \draw[edge] (At) -- (\x-0.8,0) node [midway,above] {\scalebox{0.7}{\textcolor{gray}{$d_1$}}};
    \draw[edge] (At)-- (\x-0.6,-0.5) node [above]{\scalebox{0.7}{\dimcolor{$d_2$}}};
    \draw[edge] (At) -- (\x+0.8,0) node [midway,above] {\scalebox{0.7}{\dimcolor{$d_N$}}};;
    \draw[dotted] (At) -- (\x-0.3,-0.7);
    \draw[edge] (At) -- (\x,-0.8) node [midway,right] {\scalebox{0.7}{\dimcolor{$d_n$}}};
    \draw[dotted] (At) -- (\x+0.6,-0.7);
    \draw[dotted] (At) -- (\x+0.7,-0.5);
    \draw[edge] (M) -- (\x,-1.8) node [midway,right] {\scalebox{0.7}{\dimcolor{$m$}}};;
    \end{tikzpicture}
\in\RR^{d_1\times  \dots\times d_{n-1}\times m \times d_{n+1}\times \dots\times d_N}.
$$
该运算把张量的第 $n$ 模与矩阵的第二模~（列）收缩，原来的维数 $d_n$ 被新维数 $m$ 取代。若 $\Ab\in\RR^{d_n\times n},\Bb\in\RR^{d_m\times m}$ 与 $\Xt\in\RR^{d_1\times\cdots\times d_N}$ 所对应的模 $n\neq m$ 互不相同，则相乘的顺序无关紧要，即
$$
\Xt\ttm n \Ab \ttm m \Bb = 
\begin{tikzpicture}[baseline=-0.5ex]
    \tikzset{tensor/.style = {minimum size = 0.5cm,shape = circle,thick,draw=black,inner sep = 0pt}, edge/.style = {thick,line width=.4mm},every loop/.style={}}
    \def\x{0}
    \def\y{0}
    \node[tensor,fill=blue!20!red!30!white] (Xt) at (\x,\y) {$\Xt$};
    \node[tensor,fill=green!30!red!40!white] (A) at (\x,\y-1) {$\Ab$};
    \node[tensor,fill=green!50!red!50!white] (B) at (\x+1,\y-0.7) {$\Bb$};
    \draw[edge] (Xt) -- (\x-1,\y) node [midway,above] {\scalebox{0.7}{\dimcolor{$d_1$}}};
    \draw[dotted] (Xt) -- (\x-1,\y-0.5);
    \draw[edge] (Xt) -- (\x+1,\y) node [midway,above] {\scalebox{0.7}{\dimcolor{$d_N$}}};
    \node[draw=none] () at (\x+0.7,\y-0.27) {\scalebox{0.7}{\textcolor{gray}{$d_m$}}};
    \draw[edge] (Xt) -- (A) node [midway,left] {\scalebox{0.7}{\dimcolor{$d_n$}}};
    \draw[edge] (A) -- (\x,\y-1.7) node [midway,right] {\scalebox{0.7}{\textcolor{gray}{$n$}}};
    \draw[edge] (Xt) -- (B);
    \draw[edge] (B) -- (\x+1,\y-1.5) node [midway,right] {\scalebox{0.7}{\dimcolor{$m$}}};
    \draw[dotted] (Xt) -- (\x+1,\y-0.2);
    \draw[dotted] (Xt) -- (\x+0.5,\y-0.7);
    \end{tikzpicture}
= \Xt\ttm m\Bb\ttm n \Ab,~~\text{For}~~n\neq m.
$$
然而当 $m = n$ 时，乘积的顺序就至关重要。特别地，若 $\Ab\in\RR^{n\times d_n}$、$\Bb\in\RR^{p\times n}$，则 $(\Xt\ttm n \Ab) \ttm n \Bb = \Xt\ttm n(\Bb\Ab)$，而乘积 $(\Xt\ttm n \Bb) \ttm n \Ab$ 仅当 $n=d_n=p$ 时才有意义。
模 $n$ 张量乘积可以视为经典矩阵乘法的推广。特别地，
用二阶张量的模 $1$ 乘积与模 $2$ 乘积即可还原出矩阵乘积：
若 $\Ab\in\RR^{m\times n}$、$\Bb\in\RR^{p\times n}$、$\Cb\in\RR^{d\times m}$，则
\begin{align*}
\Ab\ttm2 \Bb
=
\begin{tikzpicture}[baseline=\baseline]
   \tikzset{tensor/.style = {minimum size = 0.5cm,shape = circle,thick,draw=black,inner sep = 0pt}, edge/.style = {thick,line width=.4mm},every loop/.style={}}
    \def\y{0}
    \def\x{0}
    \node[tensor,fill=red!50!white] (A) at (\x,\y){\scalebox{0.85}{$\Ab$}};
    \node[tensor,fill=blue!50!white] (B) at (\x+1,\y){\scalebox{0.85}{$\Bb$}};
    \draw[edge] (A) -- (B) node [above,midway] {\scalebox{0.7}{\dimcolor{$n$}}};
    \draw[edge] (\x-0.25,\y) -- (\x-0.6,\y)  node [midway,above] {\scalebox{0.7}{\dimcolor{$m$}}};
    \draw[edge] (\x+1.25,\y) -- (\x+1.6,\y)  node [midway,above] {\scalebox{0.7}{\dimcolor{$p$}}};
    \end{tikzpicture}
    = \Ab\Bb^\ts\in\RR^{m\times p},
\end{align*}
\begin{align*}
    \Ab\ttm1 \Cb =
    \begin{tikzpicture}[baseline=\baseline]
   \tikzset{tensor/.style = {minimum size = 0.5cm,shape = circle,thick,draw=black,inner sep = 0pt}, edge/.style = {thick,line width=.4mm},every loop/.style={}}
    \def\y{0}
    \def\x{0}
    \node[tensor,fill=red!20!white] (C) at (\x,\y){\scalebox{0.85}{$\Cb$}};
    \node[tensor,fill=red!50!white] (A) at (\x+1,\y){\scalebox{0.85}{$\Ab$}};
    \draw[edge] (C) -- (A) node [above,midway] {\scalebox{0.7}{\dimcolor{$m$}}};
    \draw[edge] (\x-0.25,\y) -- (\x-0.6,\y)  node [midway,above] {\scalebox{0.7}{\dimcolor{$d$}}};
    \draw[edge] (\x+1.25,\y) -- (\x+1.6,\y)  node [midway,above] {\scalebox{0.7}{\dimcolor{$n$}}};
    \end{tikzpicture}
    = \Cb\Ab\in\RR^{d\times n}.
\end{align*}

下面的命题说明如何借助矩阵化与矩阵乘积来表示模 $n$ 乘积。
\begin{proposition}
设 $\Xt\in\RR^{d_1\times\dots\times d_N}$ 且 $\Mb\in\RR^{m\times d_n}$，则 
$(\Xt\ttm n \Mb)_{(n)} = \Mb\Xt_{(n)}$.
\begin{proof}
我们用一张简单的张量网络图来说明特殊情形 $n=2$、$N=3$ 下的这一恒等式，一般情形的推广是直接的。于是
\begin{align*}
(\Xt\ttm2\Mb)_{(2)} 
&= 
\left(
    \begin{tikzpicture}[baseline=\baseline]
    \tikzset{tensor/.style = {minimum size = 0.5cm,shape = circle,thick,draw=black,inner sep = 0pt}, edge/.style = {thick,line width=.4mm},every loop/.style={}}
    \def\y{0}
    \def\x{0}
    \node[tensor,fill=green!50!red!50!white] (At) at (\x,\y){\scalebox{0.85}{$\Xt$}};
    \node[tensor,fill=blue!20!white] (M) at (\x-1,\y) {$\Mb$};
    \draw[edge] (At) -- (M) node [midway,above]
    {\scalebox{0.7}{\dimcolor{$d_2$}}};
    \draw[edge] (M) -- (\x-1.7,\y) node [midway,above]
    {\scalebox{0.7}{\dimcolor{$m$}}};);
    \draw[edge] (At)-- (\x+0.5,\y+0.5) node [below] {\scalebox{0.7}{\dimcolor{$d_1$}}};
    \draw[edge] (At)-- (\x+0.5,\y-0.5) node [above] {\scalebox{0.7}{\dimcolor{$d_3$}}};
    \end{tikzpicture}
\right)_{(2)}
=
    \begin{tikzpicture}[baseline=\baseline]
    \tikzset{tensor/.style = {minimum size = 0.5cm,shape = circle,thick,draw=black,inner sep = 0pt}, edge/.style = {thick,line width=.4mm},every loop/.style={}}
    \def\y{0}
    \def\x{0}
    \node[tensor,fill=green!50!red!50!white] (At) at (\x,\y){\scalebox{0.85}{$\Xt$}};
    \node[tensor,fill=blue!20!white] (M) at (\x-1,\y) {$\Mb$};
    \draw[edge] (At) -- (M) node [midway,above]
    {\scalebox{0.7}{\dimcolor{$d_2$}}};
    \draw[edge] (M) -- (\x-1.7,\y) node [midway,above]
    {\scalebox{0.7}{\dimcolor{$m$}}};);
    \draw[edge] (\x+0.15,\y+0.2)-- (\x+0.5,\y+0.2)-- (\x+0.5,\y+0.05) -- (\x+1,\y+0.05) node [above] {\scalebox{0.7}{\dimcolor{$d_1$}}};
    \draw[edge] (\x+0.15,\y-0.2) -- (\x+0.5,\y-0.2)--(\x+0.5,\y-0.04)--(\x+1,\y-0.04) node [below] {\scalebox{0.7}{\dimcolor{$d_3$}}};
    \end{tikzpicture}\\
&=
\Mb\Xt_{(2)} \in\RR^{m\times d_1d_3}.
\end{align*}
\end{proof}
\end{proposition}
\item\textbf{模-n 乘积~（向量）。} 张量 $\Xt\in\RR^{d_1\times\cdots\times d_N}$ 与向量 $\vb\in\RR^{d_n}$ 的模 $n$ 乘积记作
$\Xt \ttm n \vb \in\RR^{d_1\times\cdots \times d_{n-1}\times d_{n+1}\times \cdots\times d_N}$。注意结果是一个 $N-1$ 阶张量，用张量网络图可表示为 
    $$
\Xt\ttm n \vb
=
\begin{tikzpicture}[baseline=-0.5ex]
    \tikzset{tensor/.style = {minimum size = 0.5cm,shape = circle,thick,draw=black,inner sep = 0pt}, edge/.style = {thick,line width=.4mm},every loop/.style={}}
    \def\x{6}
    \node[tensor,fill=blue!20!white] (M) at (\x,-1.05) {$\vb$};
    \node[tensor,fill=green!50!red!50!white] (At) at 
    (\x,0) {$\Xt$};
    \draw[edge] (At) -- (\x-0.8,0) node [midway,above] {\scalebox{0.7}{\textcolor{gray}{$d_1$}}};
    \draw[edge] (At)-- (\x-0.6,-0.5) node [above]{\scalebox{0.7}{\textcolor{gray}{$d_2$}}};
    \draw[edge] (At) -- (\x+0.8,0) node [midway,above] {\scalebox{0.7}{\textcolor{gray}{$d_N$}}};;
    \draw[dotted] (At) -- (\x-0.3,-0.7);
    \draw[edge] (At) -- (\x,-0.8) node [midway,right] {\scalebox{0.7}{\textcolor{gray}{$d_n$}}};
    \draw[dotted] (At) -- (\x+0.6,-0.7);
    \draw[dotted] (At) -- (\x+0.7,-0.5);
    \end{tikzpicture}
\in\RR^{d_1\times  \dots\times d_{n-1}\times d_{n+1}\times \dots\times d_N}.
$$
注意，与模 $n$ 矩阵乘法不同，模 $n$ 向量乘法的相乘顺序会影响结果：向量收缩把张量的阶减 1，其后的模指标随之平移。设 $\Xt\in\RR^{d_1\times\cdots\times d_m\times\cdots\times d_n\times\cdots\times d_N}$ 且 $\ab\in\RR^{d_n},\bb\in\RR^{d_m}$，则
$
\Xt\ttm n \ab\ttm m\bb \neq
(\Xt \ttm m \bb)\ttm{n}\ab,
$
因为模 $n$ 向量乘法会去掉第 $m$ 维~\citep{bader2006algorithm}。
\end{enumerate}







下面介绍张量推导中常用的几种乘积：Kronecker 积、Khatri-Rao 积与 Hadamard 积。

\paragraph{Kronecker 积}\label{subsection-kron}
设 $\Ab\in\RR^{m\times n}$、$\Bb\in\RR^{p\times q}$。Kronecker 积 $\Ab\kron\Bb\in\RR^{mp\times nq}$ 定义为
\begin{align}
    \Ab\kron\Bb = 
    \begin{bmatrix}
    a_{11}\Bb & a_{12}\Bb & \dots & a_{1n}\Bb \\
    a_{21}\Bb & a_{22}\Bb & \dots & a_{2n}\Bb \\
    \vdots & \vdots & \ddots & \vdots \\
    a_{m1}\Bb & a_{m2}\Bb & \dots & a_{mn}\Bb 
    \end{bmatrix}.
\end{align}

在张量网络中，Kronecker 积由下图给出：
\begin{align}
    \Ab\kron\Bb = 
\begin{tikzpicture}[baseline=-0.5ex]
    \tikzset{tensor/.style = {minimum size = 0.5cm,shape = circle,thick,draw=black,inner sep = 0pt}, edge/.style   = {thick,line width=.4mm}}
    \def\x{0}
    \def\y{0}
    \node[tensor,fill=red!50!white] (A) at (\x,\y+0.3){$\Ab$};
    \node[tensor,fill=blue!50!white] (B) at (\x,\y-0.3){$\Bb$};
    \draw[edge] (\x-0.25,\y+0.3) -- (\x-0.6,\y+0.3) -- (\x-0.6,\y+0.05)   -- (\x-1.2,\y+0.05) node [midway,above] {\scalebox{0.7}{\dimcolor{$m$}}};
    \draw[edge] (\x-0.25,\y-0.3) -- (\x-0.6,\y-0.3) -- (\x-0.6,\y-0.05) -- (\x-1.2,\y-0.05) node [midway,below] {\scalebox{0.7}{\dimcolor{$p$}}};
    \draw[edge] (\x+0.25,\y+0.3) -- (\x+0.6,\y+0.3) -- (\x+0.6,\y+0.05) -- (\x+1.2,\y+0.05) node [midway,above] {\scalebox{0.7}{\dimcolor{$n$}}};
    \draw[edge] (\x+0.25,\y-0.3) -- (\x+0.6,\y-0.3) -- (\x+0.6,\y-0.05) -- (\x+1.2,\y-0.05) node [midway,below] {\scalebox{0.7}{\dimcolor{$q$}}};
\end{tikzpicture}.
\end{align}
从这张图可以看出，两个矩阵的 Kronecker 积不过是它们外积的一次重排。事实上，容易验证外积 $\Ab \outprod \Bb$ 的每个分量 $a_{i,j}b_{k,l}$ 都出现在 $\Ab \kron \Bb$ 中；外积是四阶张量，而 Kronecker 积是它的一种矩阵化。更精确地，可以证明 $\Ab \kron \Bb = (\Ab \outprod \Bb)_{(\{1,3\})}$，但我们不在此展开，而是相信张量网络图给出的简洁直观的看法：把 $\Ab$ 与 $\Bb$ 的行腿归到一侧、列腿归到另一侧，所得的矩阵就称为 $\Ab$ 与 $\Bb$ 的 Kronecker 积：
$$\Ab\outprod \Bb =    
\begin{tikzpicture}[baseline=-0.5ex]
    \tikzset{tensor/.style = {minimum size = 0.5cm,shape = circle,thick,draw=black,inner sep = 0pt}, edge/.style   = {thick,line width=.4mm}}
    \def\x{0}
    \def\y{0}
    \node[tensor,fill=red!50!white] (A) at (\x,\y+0.3){$\Ab$};
    \node[tensor,fill=blue!50!white] (B) at (\x,\y-0.3){$\Bb$};
    \draw[edge] (\x-0.25,\y+0.3) -- (\x-0.75,\y+0.3) node [midway,above] {\scalebox{0.7}{\dimcolor{$m$}}};
    \draw[edge] (\x-0.25,\y-0.3) -- (\x-0.75,\y-0.3) node [midway,below] {\scalebox{0.7}{\dimcolor{$p$}}};
    \draw[edge] (\x+0.25,\y+0.3) -- (\x+0.75,\y+0.3) node [midway,above] {\scalebox{0.7}{\dimcolor{$n$}}};
    \draw[edge] (\x+0.25,\y-0.3) -- (\x+0.75,\y-0.3) node [midway,below] {\scalebox{0.7}{\dimcolor{$q$}}};
\end{tikzpicture}
\longrightarrow
\begin{tikzpicture}[baseline=-0.5ex]
    \tikzset{tensor/.style = {minimum size = 0.5cm,shape = circle,thick,draw=black,inner sep = 0pt}, edge/.style   = {thick,line width=.4mm}}
    \def\x{0}
    \def\y{0}
    \node[tensor,fill=red!50!white] (A) at (\x,\y+0.3){$\Ab$};
    \node[tensor,fill=blue!50!white] (B) at (\x,\y-0.3){$\Bb$};
    \draw[edge] (\x-0.25,\y+0.3) -- (\x-0.6,\y+0.3) -- (\x-0.6,\y+0.05)   -- (\x-1.2,\y+0.05) node [midway,above] {\scalebox{0.7}{\dimcolor{$m$}}};
    \draw[edge] (\x-0.25,\y-0.3) -- (\x-0.6,\y-0.3) -- (\x-0.6,\y-0.05) -- (\x-1.2,\y-0.05) node [midway,below] {\scalebox{0.7}{\dimcolor{$p$}}};
    \draw[edge] (\x+0.25,\y+0.3) -- (\x+0.6,\y+0.3) -- (\x+0.6,\y+0.05) -- (\x+1.2,\y+0.05) node [midway,above] {\scalebox{0.7}{\dimcolor{$n$}}};
    \draw[edge] (\x+0.25,\y-0.3) -- (\x+0.6,\y-0.3) -- (\x+0.6,\y-0.05) -- (\x+1.2,\y-0.05) node [midway,below] {\scalebox{0.7}{\dimcolor{$q$}}};
\end{tikzpicture}
=\Ab \kron \Bb.
$$




更一般地，Kronecker 积可以定义在任意两个同阶张量之间。其做法是把两个张量的每一对腿
按次序两两归并到一起，例如，
$$
\At\kron\Bt = 
   \begin{tikzpicture}[baseline=-0.5ex]
    \tikzset{tensor/.style = {minimum size = 0.5cm,shape = circle,thick,draw=black,inner sep = 0pt}, edge/.style = {thick,line width=.4mm},every loop/.style={}}
    \def\x{6}
    \node[tensor,fill=green!50!red!50!white] (A) at (\x,0) {$\At$};
    \draw[edge] (A) -- (\x-0.6,0) node [midway,above] {\scalebox{0.7}{\dimcolor{$m_1$}}};
    \draw[edge] (A) -- (\x+0.6,0) node [midway,above] {\scalebox{0.7}{\dimcolor{$m_p$}}};;
    \draw[edge] (A) -- (\x-0.5,-0.3) node [midway, below] {\scalebox{0.7}{\dimcolor{$m_2$}}};;
    \draw[dotted] (A) -- (\x,-0.6);
    \draw[dotted] (A) -- (\x,-0.6);
    \draw[dotted] (A) -- (\x+0.5,-0.5);

    \draw[dotted] (A) -- (\x+0.3,-0.6);
    \end{tikzpicture}
\kron
    \begin{tikzpicture}[baseline=-0.5ex]
    \tikzset{tensor/.style = {minimum size = 0.5cm,shape = circle,thick,draw=black,inner sep = 0pt}, edge/.style = {thick,line width=.4mm},every loop/.style={}}
    \def\x{6}
    \def\y{0}
    \node[tensor,fill=green!70!red!20!white] (B) at (\x+1,\y) {$\Bt$};
    \draw[edge] (B) -- (\x+0.3,0) node [midway,above] {\scalebox{0.7}{\dimcolor{$n_1$}}};
    \draw[edge] (B) -- (\x+1.6,0) node [midway,above] {\scalebox{0.7}{\dimcolor{$n_p$}}};;
    \draw[edge] (B) -- (\x+0.4,-0.3) node [midway, below] {\scalebox{0.7}{\dimcolor{$n_2$}}};;
    \draw[dotted] (B) -- (\x+1,-0.6);
    \draw[dotted] (B) -- (\x+1.5,-0.5);
    \draw[dotted] (B) -- (\x+1.3,-0.6);
\end{tikzpicture}
=
    \begin{tikzpicture}[baseline=-0.5ex]
    \tikzset{tensor/.style = {minimum width = 3cm,minimum height = 0.6cm,shape = rounded rectangle,thick,draw=black,inner sep = 0pt}, edge/.style = {thick,line width=.4mm},every loop/.style={}}
    \def\x{0}
    \def\y{0}
    \node[tensor,fill=green!30!white] (At) at (\x,\y){$\At\kron\Bt$};
    \draw[edge] (\x-1.25,\y-0.3) -- (\x-1.25,\y-0.6) -- (\x-1.1,\y-0.6) -- (\x-1.1,\y-0.8);
    \draw[edge] (\x-0.85,\y-0.3) -- (\x-0.85,\y-0.6) -- (\x-1,\y-0.6) -- (\x-1,\y-0.8)node [below] {\scalebox{0.7}{\dimcolor{$m_1n_1$}}};
    \draw[edge] (\x-0.65,\y-0.3) -- (\x-0.65,\y-0.6) -- (\x-0.5,\y-0.6) -- (\x-0.5,\y-0.8);
    \draw[edge] (\x-0.25,\y-0.3) -- (\x-0.25,\y-0.6) -- (\x-0.4,\y-0.6) -- (\x-0.4,\y-0.8) node [below] {\scalebox{0.7}{\dimcolor{$m_2n_2$}}};
    \draw[edge] (\x+1.25,\y-0.3) -- (\x+1.25,\y-0.6) -- (\x+1.1,\y-0.6) -- (\x+1.1,\y-0.8);
    \draw[edge] (\x+0.85,\y-0.3) -- (\x+0.85,\y-0.6) -- (\x+1,\y-0.6) -- (\x+1,\y-0.8)  node [below] {\scalebox{0.7}{\dimcolor{$m_pn_p$}}};
    \draw[dotted] (\x+0.1,\y-0.3) -- (\x+0.1,\y-0.6) -- (\x+0.25,\y-0.6) -- (\x+0.25,\y-0.8);
    \draw[dotted] (\x+0.5,\y-0.3) -- (\x+0.5,\y-0.6) -- (\x+0.35,\y-0.6) -- (\x+0.35,\y-0.82) node [below] {\scalebox{0.7}{\dimcolor{$\dots$}}};;
    \end{tikzpicture}
    \in\RR^{m_1n_1\times\dots\times m_pn_p}.
$$

\begin{remark}
下面列出 Kronecker 积的一些有用性质。
\begin{enumerate}
    \item Kronecker 积不满足交换律，即一般地有 $\Ab\kron\Bb \neq \Bb\kron\Ab$%
    . 
    \item
    $
    \Ib\kron\Ab = 
    \begin{bmatrix}
    \Ab & 0 & \dots & 0\\
    0 & \Ab & \dots & 0\\
    \vdots & \vdots & \ddots & \vdots\\
    0 & 0 & \dots & \Ab
    \end{bmatrix}
    $
    是分块对角矩阵。
    \item 两个同阶张量的 Kronecker 积仍是同阶张量，而它们的外积则使阶数翻倍。
    \item 
    对 Kronecker 积 $\Ab\kron\Bb$ 作重排，即可得到外积 $\Ab\outprod\Bb$。
 \item
两个{单位矩阵}的 Kronecker 积是更大的单位矩阵，即 $\Ib_m \kron \Ib_n = \Ib_{mn}$，画出张量网络图便不难看出：
$$
\begin{tikzpicture}[baseline=-0.5ex]
    \tikzset{tensor/.style = {minimum width = 3cm,minimum height = 0.6cm,shape = rounded rectangle,thick,draw=black,inner sep = 0pt}, edge/.style = {thick,line width=.4mm},every loop/.style={}}
    \def\x{0}
    \def\y{0.1}
    \draw[edge](\x-0.1,\y) node [above] {\scalebox{0.7}{\dimcolor{$m$}}}  -- (\x+0.1,\y) -- (\x+0.1,\y+0.1)--(\x+1,\y+0.1) -- (\x+1,\y)-- (\x+1.2,\y) node [above] {\scalebox{0.7}{\dimcolor{$m$}}};
    \draw[edge](\x-0.1,\y-0.1)  node [below] {\scalebox{0.7}{\dimcolor{$n$}}}  -- (\x+0.1,\y-0.1) -- (\x+0.1,\y-0.2)--(\x+1,\y-0.2) -- (\x+1,\y-0.1) --(\x+1.2,\y-0.1)  node [below] {\scalebox{0.7}{\dimcolor{$n$}}} ;
    \node[draw=none] at (\x+1.4,\y-0.1) {\scalebox{1}{\textcolor{black}{$=$}}};
    \draw[edge](\x+1.7,\y-0.1) -- (\x+2.8,\y-0.1) node [midway, above] {\scalebox{0.7}{\dimcolor{$mn$}}} ;
\end{tikzpicture}.
$$
     \item
     两个向量的 Kronecker 积就是它们外积的向量化：
     $$\vectorize(\ub \outprod \vb) = \vectorize(\ub  \vb^\ts) =\ub \kron \vb =  \begin{tikzpicture}[baseline=-0.5ex]
    \tikzset{tensor/.style = {minimum size = 0.5cm,shape = circle,thick,draw=black,inner sep = 0pt}, edge/.style   = {thick,line width=.4mm}}
    \def\x{0}
    \def\y{0}
    \node[tensor,fill=red!50!white] (A) at (\x,\y+0.3){$\ub$};
    \node[tensor,fill=blue!50!white] (B) at (\x,\y-0.3){$\vb$};
    \draw[edge] (\x+0.25,\y+0.3) -- (\x+0.6,\y+0.3) -- (\x+0.6,\y+0.05) -- (\x+1.2,\y+0.05) node [midway,above] {\scalebox{0.7}{\dimcolor{$m$}}};
    \draw[edge] (\x+0.25,\y-0.3) -- (\x+0.6,\y-0.3) -- (\x+0.6,\y-0.05) -- (\x+1.2,\y-0.05) node [midway,below] {\scalebox{0.7}{\dimcolor{$n$}}};
\end{tikzpicture}
    .$$
    \item 
   Kronecker 积满足所谓的\textit{混合乘积}性质：$(\Ab\kron \Bb)(\Cb\kron\Db) = \Ab\Cb\kron\Bb\Db$，其中 $\Ab\in\RR^{m\times n},\Bb\in\RR^{p\times q},\Cb\in\RR^{n\times d}$、$\Db\in\RR^{q\times k}$。
\begin{proof}
 \begin{align*}
(\Ab\kron\Bb)(\Cb\kron\Db) 
&= 
\begin{tikzpicture}[baseline=\baseline]
    \tikzset{tensor/.style = {minimum size = 0.5cm,shape = circle,thick,draw=black,inner sep = 0pt}, edge/.style   = {thick,line width=.4mm}}
    \def\y{0.4}\def\x{0}
    \draw[edge] (\x-1,\y-0.3) -- (\x-0.7,\y-0.3) -- (\x-0.7,\y+0.1)   -- (\x+0.7,\y+0.1) -- (\x+0.7,\y-0.3) -- (\x+1.1,\y-0.3);
    \node[tensor,fill=red!50!white] (A) at (\x,\y){$\Ab$};
    \draw[edge] (\x+1,\y-0.3) -- (\x+1.3,\y-0.3) -- (\x+1.3,\y+0.1)   -- (\x+2.6,\y+0.1) -- (\x+2.6,\y-0.3) -- (\x+3,\y-0.3);
    \node[tensor,fill=orange!50!white] (C) at (\x+2,\y){$\Cb$};
    \def\y{-0.4}\def\x{0}
    \draw[edge] (\x-1,\y+0.4) -- (\x-0.7,\y+0.4) -- (\x-0.7,\y+0.1) -- (\x+0.7,\y+0.1) -- (\x+0.7,\y+0.4) -- (\x+0.7,\y+0.4) -- (\x+1.1,\y+0.4);
    \node[tensor,fill=blue!50!white] (B) at (\x,\y+0.1){$\Bb$};
    \draw[edge] (\x+1,\y+0.4) -- (\x+1.3,\y+0.4) -- (\x+1.3,\y+0.1) -- (\x+2.6,\y+0.1) -- (\x+2.6,\y+0.4) -- (\x+2.6,\y+0.4) -- (\x+3,\y+0.4);
    \node[draw=none] () at (\x+3.25,\y+0.45) {\scalebox{0.7}{\dimcolor{$dk$}}};
    \node[tensor,fill=green!50!red!40!white] (D) at (\x+2,\y+0.1){$\Db$};
    \node[draw=none] () at (\x+1,\y+0.7) {\scalebox{0.7}{\dimcolor{$nq$}}};
    \node[draw=none] () at (\x-1.25,\y+0.4) {\scalebox{0.7}{\dimcolor{$mp$}}};
\end{tikzpicture}\\
 &=
 \begin{tikzpicture}[baseline=\baseline]
    \tikzset{tensor/.style = {minimum size = 0.5cm,shape = circle,thick,draw=black,inner sep = 0pt}, edge/.style   = {thick,line width=.4mm}}
    \def\y{0.4}\def\x{0}
    \draw[edge] (\x-1,\y-0.35) -- (\x-0.7,\y-0.35) -- (\x-0.7,\y)   -- (\x,\y) -- (\x+0.5,\y) -- (\x+0.8,\y);
    \node[tensor,fill=red!50!white] (A) at (\x,\y){$\Ab$};
    \draw[edge] (\x+1.2,\y) -- (\x+1.7,\y) -- (\x+1.7,\y-0.35) -- (\x+2.1,\y-0.35);
    \node[tensor,fill=orange!50!white] (C) at (\x+1,\y){$\Cb$};
    \def\y{-0.4}\def\x{0}
    \draw[edge] (\x-1,\y+0.35) -- (\x-0.7,\y+0.35) -- (\x-0.7,\y) -- (\x,\y) -- (\x+0.8,\y);
    \node[tensor,fill=blue!50!white] (B) at (\x,\y){$\Bb$};
    \draw[edge] (\x+1,\y) -- (\x+1.7,\y) -- (\x+1.7,\y+0.35) -- (\x+2.1,\y+0.35);
    \node[draw=none] () at (\x+2.35,\y+0.4) {\scalebox{0.7}{\dimcolor{$dk$}}};
    \node[tensor,fill=green!50!red!40!white] (D) at (\x+1,\y){$\Db$};
    \node[draw=none] () at (\x-1.25,\y+0.35) {\scalebox{0.7}{\dimcolor{$mp$}}};
    \node[draw=none] () at (\x+0.5,\y+1) {\scalebox{0.7}{\dimcolor{$n$}}};
     \node[draw=none] () at (\x+0.5,\y+0.2) {\scalebox{0.7}{\dimcolor{$q$}}};
\end{tikzpicture}\\
&=
\Ab\Cb\kron\Bb\Db.
    \end{align*}
\end{proof}
特别地，$(\Ab\kron\Ib)(\Ib\kron\Bb)=\Ab\kron\Bb$。 
\item 由上面的推导可以看出，在处理含 Kronecker 积的恒等式时，常常需要把
$
    \begin{tikzpicture}[baseline=\baseline]
    \tikzset{tensor/.style = {minimum size = 0.5cm,shape = circle,thick,draw=black,inner sep = 0pt}, edge/.style   = {thick,line width=.4mm}}
    \def\y{0.6}\def\x{0}
    \draw[edge] (\x+3,\y-0.3)  node [above] {\scalebox{0.7}{\textcolor{gray}{$p$}}} -- (\x+4,\y-0.3)node [above] {\scalebox{0.7}{\textcolor{gray}{$n$}}};
    \def\y{-0.6}\def\x{0}
    \draw[edge] (\x+3,\y+0.3) node [above] {\scalebox{0.7}{\dimcolor{$m$}}}-- (\x+4,\y+0.3)node [above] {\scalebox{0.7}{\dimcolor{$q$}}};
    \end{tikzpicture}
$
重排成
$
 \begin{tikzpicture}[baseline=\baseline]
    \tikzset{tensor/.style = {minimum size = 0.5cm,shape = circle,thick,draw=black,inner sep = 0pt}, edge/.style   = {thick,line width=.4mm}}
    \def\y{0.4}\def\x{0}
    \draw[edge] (\x-1,\y-0.3) -- (\x-0.7,\y-0.3)-- (\x-0.7,\y+0.1) -- (\x+0.7,\y+0.1) -- (\x+0.7,\y-0.3) -- (\x+1.1,\y-0.3);
    \node[draw=none] () at (\x+1.35,\y-0.4) {\scalebox{0.7}{\dimcolor{$nq$}}};
    \node[draw=none] () at (\x-1.25,\y-0.4) {\scalebox{0.7}{\dimcolor{$pm$}}};
    \def\y{-0.4}\def\x{0}
    \draw[edge] (\x-1,\y+0.4) -- (\x-0.7,\y+0.4) -- (\x-0.7,\y+0.1) -- (\x+0.7,\y+0.1) -- (\x+0.7,\y+0.4) -- (\x+0.7,\y+0.4) -- (\x+1.1,\y+0.4);
    \end{tikzpicture}
$，反之亦然。
\item 对 $\Ab\in\RR^{m\times m}$ 与 $\Bb\in\RR^{n\times n},$，有等式 $\tr(\Ab\kron\Bb)=\tr(\Ab)\tr(\Bb)$，由下图不难证明：
$$
\tr(\Ab\kron\Bb) = 
\begin{tikzpicture}[baseline=-0.5ex]
    \tikzset{tensor/.style = {minimum size = 0.5cm,shape = circle,thick,draw=black,inner sep = 0pt}, edge/.style   = {thick,line width=.4mm}}
    \def\x{0}
    \def\y{0}
    \node[tensor,fill=red!50!white] (A) at (\x,\y+0.5){$\Ab$};
    \node[tensor,fill=blue!50!white] (B) at (\x,\y-0.2){$\Bb$};
    \draw[edge] (\x-0.25,\y+0.5) -- (\x-0.6,\y+0.5) -- (\x-0.6,\y+0.15)   -- (\x-1.2,\y+0.15);
    \draw[edge] (\x-0.25,\y-0.3) -- (\x-0.6,\y-0.3) -- (\x-0.6,\y+0.05) -- (\x-1.2,\y+0.05);
    \draw[edge] (\x+0.25,\y+0.5) -- (\x+0.6,\y+0.5) -- (\x+0.6,\y+0.15) -- (\x+1.2,\y+0.15);
    \draw[edge] (\x+0.25,\y-0.3) -- (\x+0.6,\y-0.3) -- (\x+0.6,\y+0.05) -- (\x+1.2,\y+0.05);
    \draw[edge](\x-1.2,\y+0.15) to [out=45,in=135, distance=1.5cm] (\x+1.2,\y+0.15);
    \draw[edge](\x-1.2,\y+0.05) to [out=-45,in=-135, distance=1.5cm] (\x+1.2,\y+0.05);
\end{tikzpicture}
=
\tr(\Ab)\tr(\Bb).
$$

\item \textbf{Sylvester 恒等式（Sylvester identity）}是一条联系向量化、矩阵乘法与 Kronecker 积的非常有用的恒等式：$$\vectorize(\Ab\Xb\Bb) = (\Ab\kron \Bb^\ts)\vectorize(\Xb),$$
其中 $\Ab\in\RR^{m\times n},\Xb\in\RR^{n\times p}$ 且 $\Bb\in\RR^{p\times q}.$
\begin{proof}
利用张量网络，证明相当直接： 
\begin{align*}
\vectorize(\Ab\Xb\Bb) &=
 \vectorize\left(\begin{tikzpicture}[baseline=-0.5ex]
    \tikzset{tensor/.style = {minimum size = 0.5cm,shape = circle,thick,draw=black,inner sep = 0pt}, edge/.style = {thick,line width=.4mm},every loop/.style={}}
    \def\x{6}
    \node[tensor,fill=blue!50!white] (A) at (\x-1,0) {$\Ab$};
    \node[tensor,fill=green!50!red!50!white] (X) at (\x,0) {$\Xb$};
    \node[tensor,fill=red!50!white] (B) at (\x+1,0) {$\Bb$};
    \draw[edge] (A) -- (X) node [midway,above] {\scalebox{0.7}{\dimcolor{$n$}}};
    \draw[edge] (X) -- (B) node [midway,above] {\scalebox{0.7}{\dimcolor{$p$}}};
    \draw[edge] (A) -- (\x-1.7,0) node [midway,above] {\scalebox{0.7}{\dimcolor{$m$}}};
    \draw[edge] (B) -- (\x+1.7,0) node [midway,above] {\scalebox{0.7}{\dimcolor{$q$}}};
    \end{tikzpicture}\right)
=
\begin{tikzpicture}[baseline=-0.5ex]
    \tikzset{tensor/.style = {minimum size = 0.5cm,shape = circle,thick,draw=black,inner sep = 0pt}, edge/.style = {thick,line width=.4mm},every loop/.style={}}
    \def\x{6}
    \node[tensor,fill=blue!50!white] (A) at (\x-1,0) {$\Ab$};
    \node[tensor,fill=green!50!red!50!white] (X) at (\x,0) {$\Xb$};
    \node[tensor,fill=red!50!white] (B) at (\x+1,0) {$\Bb$};
    \draw[edge] (A) -- (X) node [midway,below] {\scalebox{0.7}{\dimcolor{$n$}}};
    \draw[edge] (X) -- (B) node [midway,below] {\scalebox{0.7}{\dimcolor{$p$}}};
    \draw[edge] (A) -- (\x-1.5,0) -- (\x-1.5,-0.4) -- (\x-0.05,-0.4)--(\x-0.05,-0.8) node [midway,left] {\scalebox{0.7}{\dimcolor{$m$}}};
    \draw[edge] (B) -- (\x+1.5,0) -- (\x+1.5,-0.4) -- (\x+0.05,-0.4)--(\x+0.05,-0.8) node [midway,right] {\scalebox{0.7}{\dimcolor{$q$}}};
\end{tikzpicture}\\
&=
\begin{tikzpicture}[baseline=-2ex]
    \tikzset{tensor/.style = {minimum size = 0.5cm,shape = circle,thick,draw=black,inner sep = 0pt}, edge/.style = {thick,line width=.4mm},every loop/.style={}}
    \def\x{6}
    \draw[edge] (\x-0.05,0.25) -- (\x-0.05,0) node [midway, left] {\scalebox{0.7}{\dimcolor{$n$}}}   -- (\x-0.6,0) -- (\x-0.6,-0.25) ; 
    \draw[edge] (\x+0.05,0.25) -- (\x+0.05,0)   node [midway, right] {\scalebox{0.7}{\dimcolor{$p$}}}  -- (\x+0.6,0) -- (\x+0.6,-0.25);    
    \draw[edge] (\x-0.05,-1.4) -- (\x-0.05,-1)  node [midway, left] {\scalebox{0.7}{\dimcolor{$m$}}}  -- (\x-0.6,-1) -- (\x-0.6,-0.75) ; 
    \draw[edge] (\x+0.05,-1.4) -- (\x+0.05,-1)  node [midway, right] {\scalebox{0.7}{\dimcolor{$q$}}}   -- (\x+0.6,-1) -- (\x+0.6,-0.75); 
     \node[tensor,fill=green!50!red!50!white] (X) at (\x,0.5) {$\Xb$};
    \node[tensor,fill=blue!50!white] (A) at (\x-0.6,-0.5) {$\Ab$};
    \node[tensor,fill=red!50!white] (B) at (\x+0.6,-0.5) {$\Bb$};
\end{tikzpicture}
=
(\Ab\kron\Bb^\ts)\vectorize(\Xb).
\end{align*}
在最后一个图中，容易看出 $\Xb$ 的向量化以及 $\Ab$ 与 $\Bb$ 的 Kronecker 积，从而可以断定该张量网络表示二者之间的矩阵-向量乘法。然而，把张量网络翻译成经典数学记法时需要注意一些细节。首先，由于 $\vectorize(\Xb)$ 的维度为 $np$，Kronecker 积必须取在 $\Ab$ 与 $\Bb^\ts$ 之间，才能得到形状为 $mq\times np$ 的矩阵。其次，需要确定使用哪一个 Kronecker 积：$\Ab \kron \Bb^\ts$ 还是 $\Bb^\ts \kron \Ab$，二者的形状都正确。由于本手稿对 $\vectorize$ 这类重排运算采用字典序，正确的写法是  $\Ab \kron \Bb^\ts$。然而，若采用逆字典序，如~\citep{kolda2009tensor} 中的做法，正确的写法应为 $\Bb^\ts \kron \Ab$。
\end{proof}
\emph{Sylvester 恒等式}是求解形如 $\Ab\Xb+\Xb\Bb = \Cb$ 的方程组的有用恒等式。
\end{enumerate}
\end{remark}
模 $n$ 乘积、矩阵化与 Kronecker 积之间存在若干有用的关系： 
\begin{proposition}
设 $\Tt\in\RR^{d_1\times d_2\times\dots\times d_p}$，$\Ab_1\in\RR^{d_1\times n_1},
\Ab_2\in\RR^{d_2\times n_2},\dots,$
$\Ab_p\in\RR^{d_p\times n_p}$，则 
\begin{align*}(\Tt\ttm1\Ab_1\ttm2\Ab_2\ttm3\Ab_3\ttm4\dots\ttm p\Ab_p)_{(n)}=
\Ab_n\Tt_{(n)}(\Ab_1\kron\dots\kron\Ab_{n-1}\kron\Ab_{n+1}\kron\dots\kron\Ab_p)^\ts.
\end{align*}
\end{proposition}
\begin{proof}
当 $p =3$ 且 $n=1$ 时，
\begin{align*}
(\Tt\ttm1\Ab_1\ttm2\Ab_2\ttm3\Ab_3)_{(1)}
&= 
 \left(\begin{tikzpicture}[baseline=-0.5ex]
    \tikzset{tensor/.style = {minimum size = 0.5cm,shape = circle,thick,draw=black,inner sep = 0pt}, edge/.style = {thick,line width=.4mm},every loop/.style={}}
    \def\x{6}
    \node[tensor,fill=blue!20!white] (A1) at (\x,-1.05) {$\Ab_1$};
    \node[tensor,fill=green!50!red!50!white] (At) at (\x,0) {$\Tt$};
    \node[tensor,fill=red!50!white] (A2) at (\x+1,0) {$\Ab_2$};
    \node[tensor,fill=blue!50!white] (A3) at (\x-1,0) {$\Ab_3$};
    \draw[edge] (At) -- (A2) node [midway,above] {\scalebox{0.7}{\dimcolor{$d_2$}}};
    \draw[edge] (At) -- (A1) node [midway,left] {\scalebox{0.7}{\dimcolor{$d_1$}}};
    \draw[edge] (At) -- (A3) node [midway,above] {\scalebox{0.7}{\dimcolor{$d_3$}}};
    \draw[edge] (A1) -- (\x,-1.8) node [midway,right] {\scalebox{0.7}{\dimcolor{$n_1$}}};
    \draw[edge] (A2) -- (\x+1.8,0) node [midway,above] {\scalebox{0.7}{\dimcolor{$n_2$}}};
    \draw[edge] (A3) -- (\x-1.8,0) node [midway,above] {\scalebox{0.7}{\dimcolor{$n_3$}}};
    \end{tikzpicture}\right)_{(1)}
=
\begin{tikzpicture}[baseline=-0.5ex]
    \tikzset{tensor/.style = {minimum size = 0.5cm,shape = circle,thick,draw=black,inner sep = 0pt}, edge/.style = {thick,line width=.4mm},every loop/.style={}}
    \def\x{6}
    \def\y{0}
    \node[tensor,fill=blue!20!white] (A1) at (\x,-1.05) {$\Ab_1$};
    \node[tensor,fill=green!50!red!50!white] (At) at (\x,0) {$\Tt$};
    \node[tensor,fill=red!50!white] (A2) at (\x+1,\y+0.5) {$\Ab_2$};
    \node[tensor,fill=blue!50!white] (A3) at (\x+1,\y-0.5) {$\Ab_3$};
    \draw[edge] (At) -- (A1) node [midway,left] {\scalebox{0.7}{\dimcolor{$d_1$}}};
    \draw[edge] (At) -- (A3) node [midway,below] {\scalebox{0.7}{\dimcolor{$d_3$}}};
    \draw[edge] (At) -- (A2) node [midway,above] {\scalebox{0.7}{\dimcolor{$d_2$}}};
    \draw[edge] (A1) -- (\x,-1.8) node [midway,right] {\scalebox{0.7}{\dimcolor{$n_1$}}};
    \draw[edge] (A2) -- (\x+1.8,\y+0.5) -- (\x+1.8,\y)-- (\x+2.5,\y);
    \draw[edge] (A3) -- (\x+1.8,\y-0.5) -- (\x+1.8,\y-0.1) -- (\x+2.5,\y-0.1)node [right] {\scalebox{0.7}{\dimcolor{$n_2n_3$}}};
    \end{tikzpicture}\\
&=
\begin{tikzpicture}[baseline=-0.5ex]
    \tikzset{tensor/.style = {minimum size = 0.5cm,shape = circle,thick,draw=black,inner sep = 0pt}, edge/.style = {thick,line width=.4mm},every loop/.style={}}
    \def\x{6}
    \def\y{0}
    \node[tensor,fill=blue!20!white] (A1) at (\x,-1.05) {$\Ab_1$};
    \node[tensor,fill=green!50!red!50!white] (At) at (\x,0) {$\Tt$};
    \node[tensor,fill=red!50!white] (A2) at (\x+1,\y+0.5) {$\Ab_2$};
    \node[tensor,fill=blue!50!white] (A3) at (\x+1,\y-0.5) {$\Ab_3$};
    \draw[edge] (At) -- (A1) node [midway,left] {\scalebox{0.7}{\dimcolor{$d_1$}}};
    \draw[edge] (\x+0.7,\y-0.5) -- (\x+0.4,\y-0.5) -- (\x+0.4,\y) -- (\x+0.25,\y);
    \draw[edge] (\x+0.7,\y+0.5) -- (\x+0.4,\y+0.5) -- (\x+0.4,\y+0.1) -- (\x+0.25,\y+0.1);
    \draw[edge] (A1) -- (\x,-1.8) node [midway,right] {\scalebox{0.7}{\dimcolor{$n_1$}}};
    \draw[edge] (A2) -- (\x+1.6,\y+0.5) -- (\x+1.6,\y)-- (\x+2.35,\y);
    \draw[edge] (A3) -- (\x+1.6,\y-0.5) -- (\x+1.6,\y-0.1) -- (\x+2.35,\y-0.1)node [right] {\scalebox{0.7}{\dimcolor{$n_2n_3$}}};
    \node[draw=none] () at (\x+0.6,\y-0.3) {\scalebox{0.6}{\dimcolor{$d_3$}}};
    \node[draw=none] () at (\x+0.6,\y+0.3) {\scalebox{0.6}{\dimcolor{$d_2$}}};
\end{tikzpicture}
=
\Ab_1\Tt_{(1)}(\Ab_2\kron\Ab_3)^\ts.
\end{align*}
\end{proof}
\begin{proposition}
    设 $\Tt\in\RR^{d_1\times d_2\times\cdots\times d_p}$  且对 $i\in[p]$ 有 $\Ab_i\in\RR^{d_i\times n_i}$。设 $m \in [p]$，
记 $\Tt_{([m])}$ 表示把张量 $\Tt$ 重排为尺寸 $d_1d_2\cdots d_m\times d_{m+1}...d_p$ 的矩阵：前 $m$ 个指标映射到
行，其余指标映射到列。于是，
\begin{align*}
    (\Tt\ttm 1\Ab_1\ttm 2 \Ab_2\cdots\ttm p \Ab_p)_{([m])}=(\Ab_1\kron\Ab_2\kron\cdots\kron\Ab_{m})\Tb_{([m])}(\Ab_{m+1}\kron\Ab_{m+2}\kron\cdots\kron\Ab_{p})^\ts.
\end{align*}
此恒等式可推广到任意子集 $I\subset [p]$：
$$ (\Tt\ttm 1\Ab_1\ttm 2 \Ab_2\cdots\ttm p \Ab_p)_{(I)}    = \left(\bigotimes_{i\in I} \Ab_i\right)\Tt_{([m])}\left(\bigotimes_{i\in [p] \setminus I} \Ab_i\right)^\ts.$$
\begin{proof}
当 $p=6$ 且 $I = \{1,2,3\}$ 时，
\begin{align*}
\left(\begin{tikzpicture}[baseline=\baseline]
    \tikzset{tensor/.style = {minimum size = 0.5cm,shape = circle,thick,draw=black,inner sep = 0pt}, edge/.style = {thick,line width=.4mm},every loop/.style={}}
    \def\y{0}
    \def\x{0}
    \node[tensor,fill=red!20!blue!40!white] (T) at (\x,\y){\scalebox{0.85}{$\Tt$}};
\node[tensor,fill=red!20!white] (A) at (\x+1,\y){\scalebox{0.85}{$\Ab_1$}};
    \draw[edge](T) -- (A) node [midway,above] {\scalebox{0.7}{\dimcolor{$d_1$}}};
    \draw[edge](A) -- (\x+2,\y) node [midway,above] {\scalebox{0.7}{\dimcolor{$n_1$}}};
\node[tensor,fill=red!40!white] (B) at (\x,\y-1) {\scalebox{0.85}{$\Ab_2$}}; 
    \draw[edge](T) -- (B) node [midway,left] {\scalebox{0.7}{\dimcolor{$d_2$}}};
    \draw[edge](B) -- (\x,\y-2) node [midway,left] {\scalebox{0.7}{\dimcolor{$n_2$}}};
\node[tensor,fill=red!60!white] (C) at (\x-1,\y) {\scalebox{0.85}{$\Ab_3$}};
    \draw[edge](T) -- (C) node [midway,above] {\scalebox{0.7}{\dimcolor{$d_3$}}};
    \draw[edge](C) -- (\x-2,\y) node [midway,above] {\scalebox{0.7}{\dimcolor{$n_3$}}};
\node[tensor,fill=blue!20!white] (D) at (\x,\y+1){\scalebox{0.85}{$\Ab_4$}}; 
    \draw[edge](T) -- (D) node [midway,right] {\scalebox{0.7}{\dimcolor{$d_4$}}};
    \draw[edge](D) -- (\x,\y+2) node [midway,right] {\scalebox{0.7}{\dimcolor{$n_4$}}};
\node[tensor,fill=red!50!blue!10!white] (E) at (\x-0.85,\y+0.85){\scalebox{0.85}{$\Ab_5$}}; 
   \draw[edge](T) -- (E) node [midway,above] {\scalebox{0.7}{\dimcolor{$d_5$}}};
    \draw[edge](E) -- (\x-1.5,\y+1.5) node [midway,right] {\scalebox{0.7}{\dimcolor{$n_5$}}};
\node[tensor,fill=blue!50!white] (F) at (\x+0.85,\y-0.85){\scalebox{0.85}{$\Ab_6$}};
    \draw[edge](T) -- (F)node [midway,right] {\scalebox{0.7}{\dimcolor{$d_6$}}};
    \draw[edge](F) -- (\x+1.5,\y-1.5)node [midway,right] {\scalebox{0.7}{\dimcolor{$n_6$}}};
\end{tikzpicture}\right)_{(\{1,2,3\})}
\begin{tikzpicture}[baseline=\baseline]
    \tikzset{tensor/.style = {minimum size = 0.5cm,shape = circle,thick,draw=black,inner sep = 0pt}, edge/.style = {thick,line width=.4mm},every loop/.style={}}
    \def\y{0}
    \def\x{0}
    \node[draw=none] at (\x-3.5,\y){$=$};
    \node[tensor,fill=red!20!blue!40!white] (T) at (\x,\y){\scalebox{0.85}{$\Tb$}};
    \node[tensor,fill=red!20!white] (A) at (\x-1.5,\y-1){\scalebox{0.85}{$\Ab_1$}};
    \node[tensor,fill=red!40!white] (B) at (\x-1.5,\y) {\scalebox{0.85}{$\Ab_2$}}; 
    \node[tensor,fill=red!60!white] (C) at (\x-1.5,\y+1) {\scalebox{0.85}{$\Ab_3$}};
    \node[tensor,fill=blue!20!white] (D) at (\x+1.5,\y-1){\scalebox{0.85}{$\Ab_4$}};
\node[tensor,fill=red!50!blue!10!white] (E) at (\x+1.5,\y){\scalebox{0.85}{$\Ab_5$}}; 
\node[tensor,fill=blue!50!white] (F) at (\x+1.5,\y+1){\scalebox{0.85}{$\Ab_6$}};
    \draw[edge](T) -- (B);
    \node[draw=none] at (\x-1.15,\y+0.2) {\scalebox{0.7}{\dimcolor{$d_2$}}};;
    \node[draw=none] at (\x+1.15,\y+0.2) {\scalebox{0.7}{\dimcolor{$d_5$}}};;
    \draw[edge](B) -- (\x-2.5,\y);;
    \draw[edge](\x+0.16,\y+0.2) -- (\x+0.5,\y+0.2) -- (\x+0.5,\y+0.1) -- (\x+1,\y+0.1);
    \draw[edge](F) -- (\x+1,\y+1) -- (\x+1,\y+0.1)node [midway,left] {\scalebox{0.7}{\dimcolor{$d_6$}}};
    \draw[edge](F) -- (\x+2,\y+1)--(\x+2,\y+0.1)-- (\x+2.5,\y+0.1)node [near end,above] {\scalebox{0.7}{\dimcolor{$n_6$}}};
    \draw[edge](\x+0.16,\y-0.2) -- (\x+0.5,\y-0.2) -- (\x+0.5,\y-0.1) -- (\x+1,\y-0.1);
    \draw[edge](D) -- (\x+1,\y-1) -- (\x+1,\y-0.1)node [midway,left] {\scalebox{0.7}{\dimcolor{$d_4$}}};
    \draw[edge](D) -- (\x+2,\y-1)--(\x+2,\y-0.1)-- (\x+2.5,\y-0.1)node [near end,below] {\scalebox{0.7}{\dimcolor{$n_4$}}};
    \draw[edge](T) -- (E);
    \draw[edge](E) -- (\x+2.5,\y);
    \draw[edge](\x-0.16,\y+0.2) -- (\x-0.5,\y+0.2) -- (\x-0.5,\y+0.1) -- (\x-1,\y+0.1);
    \draw[edge](C) -- (\x-1,\y+1) -- (\x-1,\y+0.1)node [midway,right] {\scalebox{0.7}{\dimcolor{$d_3$}}};
    \draw[edge](C) -- (\x-2,\y+1)--(\x-2,\y+0.1)-- (\x-2.5,\y+0.1)node [near end,above] {\scalebox{0.7}{\dimcolor{$n_3$}}};;;
    \draw[edge](\x-0.16,\y-0.2) -- (\x-0.5,\y-0.2) -- (\x-0.5,\y-0.1) -- (\x-1,\y-0.1);
    \draw[edge](A) --  (\x-2,\y-1)--(\x-2,\y-0.1)-- (\x-2.5,\y-0.1)node [near end,below] {\scalebox{0.7}{\dimcolor{$n_1$}}};;
    \draw[edge](A) -- (\x-1,\y-1) -- (\x-1,\y-0.1)node [midway,right] {\scalebox{0.7}{\dimcolor{$d_1$}}};
    \node[draw=none] at (\x+2.8,\y){\scalebox{0.7}{\dimcolor{$n_4$}}};
    \node[draw=none] at (\x-2.8,\y){\scalebox{0.7}{\dimcolor{$n_2$}}};
\end{tikzpicture}.
\end{align*}
\end{proof}
\end{proposition}


\paragraph{Khatri-Rao 积}
Khatri-Rao 积是按列进行的 Kronecker 积~\citep{smilde2005multi,bro1998multi}。设 $\Ab\in\RR^{m\times R}$ 且 $\Bb\in\RR^{n\times R}$。$\Ab$ 与 $\Bb$ 的 Khatri-Rao 积记为 $\Ab\krao\Bb\in\RR^{mn\times R}$，定义为
\begin{align*}
    \def\baseline{-0.5ex}
    \Ab\krao\Bb =
    \begin{bmatrix}
    \vline & \vline  &   & \vline \\
    \ab_1\kron\bb_1 & \ab_2\kron\bb_2 & \dots & \ab_R\kron\bb_R\\
    \vline & \vline & & \vline 
    \end{bmatrix}
=
    \begin{tikzpicture}[baseline=\baseline]
    \tikzset{tensor/.style = {minimum size = 0.5cm,shape = circle,thick,draw=black,inner sep = 0pt}, edge/.style   = {thick,line width=.4mm},every loop/.style={}}
    \def\y{0}
    \def\x{0}
    \node[tensor,fill=red!50!white] (A) at (\x,\y){$\Ab$};
    \node[tensor,fill=blue!50!white] (B) at (\x+1,\y){$\Bb$};
    \draw[edge] (\x,\y+0.25) -- (\x,\y+0.5) -- (\x+0.5,\y+0.7);
    \draw[edge] (\x+1,\y+0.25) -- (\x+1,\y+0.5) -- (\x+0.5,\y+0.7);
    \draw[edge] (\x+0.5,\y+0.7) -- (\x+0.5,\y+1.2) node [left, midway] {\edgelabel{$R$}};
    \draw  node[fill = black,circle,inner sep=0pt,minimum size=4pt] (m) at (\x+0.5,\y+0.7) {};
    \draw[edge] (\x,\y-0.25) -- (\x,\y-0.7);
    \draw[edge] (\x+1,\y-0.25) -- (\x+1,\y-0.7);
    \draw[edge] (\x+1,\y-0.7) -- (\x+0.5,\y-0.7) -- (\x+0.5,\y-1);
    \draw[edge] (\x+1,\y-0.25) -- (\x+1,\y-0.7);
    \draw[edge] (\x,\y-0.7) -- (\x+0.4,\y-0.7) -- (\x+0.4,\y-1)node [below] {\edgelabel{$mn$}};
    \end{tikzpicture}\hspace{0.5cm},
\end{align*}
其中 $\ab_1,\dots,\ab_R\in\RR^{m}$ 为 $\Ab$ 的各列，$\bb_1,\dots,\bb_R\in\RR^{n}$ 为 $\Bb$ 的各列。  在对应的张量网络图中，复制张量强制所得矩阵只包含 $\Ab$ 与 $\Bb$ 中\emph{相互匹配}的各列的 Kronecker 积。 
\begin{remark}
Khatri-Rao 积的部分性质列举如下：
\begin{enumerate}
\item 与 Kronecker 积类似，Khatri-Rao 积不满足交换律，即一般情况下 $\Ab\krao\Bb\neq \Bb\krao\Ab$。
\item Khatri-Rao 积满足结合律，即 $\Ab\krao(\Bb\krao\Cb)=(\Ab\krao\Bb)\krao\Cb$，其中 $\Ab\in\RR^{m\times R},\Bb\in\RR^{n\times R}$ 且 $\Cb\in\RR^{s\times R}$。
\item $\Ab\krao\Bb$ 的各列构成 Kronecker 积 $\Ab\kron \Bb$ 各列的一个子集。
\end{enumerate}
\end{remark}

\paragraph{Hadamard 积}
设 $\Ab\in\RR^{m\times n}$ 与 $\Bb\in\RR^{m\times n}$ 为形状相同的矩阵。$\Ab$ 与 $\Bb$ 的 Hadamard 积（或称逐分量积）是矩阵  $\Ab\hadam\Bb\in\RR^{m\times n}$，按元素定义为 
$$(\Ab\hadam\Bb)_{ij} = \Ab_{ij}\Bb_{ij}.$$
在张量网络图中，Hadamard 积可以很方便地用复制张量表示：
\begin{equation*}
\Ab\hadam\Bb =
     \begin{tikzpicture}[baseline=\baseline]
    \tikzset{tensor/.style = {minimum size = 0.5cm,shape = circle,thick,draw=black,inner sep = 0pt}, edge/.style   = {thick,line width=.4mm},every loop/.style={}}
    \def\y{0}
    \def\x{0}
    \draw[edge] (\x,\y-0.7) -- (\x,\y+0.7)  node [very near end,left] {\edgelabel{$m$}} node [very near start,left] {\edgelabel{$n$}} ;
    \node[tensor,fill=red!50!white] (A) at (\x,\y){$\Ab$};
    \end{tikzpicture} 
    \ \hadam %
    \begin{tikzpicture}[baseline=\baseline]
    \tikzset{tensor/.style = {minimum size = 0.5cm,shape = circle,thick,draw=black,inner sep = 0pt}, edge/.style   = {thick,line width=.4mm},every loop/.style={}}
    \def\y{0}
    \def\x{0}
    \draw[edge] (\x,\y-0.7) -- (\x,\y+0.7)node [very near end,left] {\edgelabel{$m$}} node [very near start,left] {\edgelabel{$n$}} ;
    \node[tensor,fill=blue!50!white] (B) at (\x,\y){$\Bb$};
    \end{tikzpicture}
\ =
\begin{tikzpicture}[baseline=\baseline]
    \tikzset{tensor/.style = {minimum size = 0.5cm,shape = circle,thick,draw=black,inner sep = 0pt}, edge/.style   = {thick,line width=.4mm},every loop/.style={}}
    \def\y{0}
    \def\x{0}
    \node[tensor,fill=red!50!white] (A) at (\x,\y){$\Ab$};
    \node[tensor,fill=blue!50!white] (B) at (\x+1,\y){$\Bb$};
    \draw[edge] (\x,\y+0.25) -- (\x,\y+0.5) -- (\x+0.5,\y+0.7);
    \draw[edge] (\x+1,\y+0.25) -- (\x+1,\y+0.5) -- (\x+0.5,\y+0.7);
    \draw[edge] (\x+0.5,\y+0.7) -- (\x+0.5,\y+1) node [near start,left] {\edgelabel{$m$}};
    \draw  node[fill = black,circle,inner sep=0pt,minimum size=4pt] (m) at (\x+0.5,\y+0.7) {};
    \draw[edge] (\x,\y-0.25) -- (\x,\y-0.5) -- (\x+0.5,\y-0.7);
    \draw[edge] (\x+1,\y-0.25) -- (\x+1,\y-0.5) -- (\x+0.5,\y-0.7);
    \draw[edge] (\x+0.5,\y-0.7) -- (\x+0.5,\y-1) node [near start,left] {\edgelabel{$n$}};
    \draw  node[fill = black,circle,inner sep=0pt,minimum size=4pt] (m) at (\x+0.5,\y-0.7) {};
\end{tikzpicture}.
\end{equation*}
更一般地，对任意 $\Ab_1\in\RR^{m\times n},\dots,\Ab_d\in\RR^{m\times n}$，有 
\begin{align*}
\Ab_1\hadam\Ab_2\dots\hadam\Ab_{d} 
=
    \begin{tikzpicture}[baseline=\baseline]
    \tikzset{tensor/.style = {minimum size = 0.5cm,shape = circle,thick,draw=black,inner sep = 0pt}, edge/.style   = {thick,line width=.4mm},every loop/.style={}}
    \def\y{1}
    \def\x{0}
    \node[tensor,fill=red!30!white] (A_1) at (\x,\y){$\Ab_1$};
    \node[tensor,fill=red!30!white] (A_2) at (\x,\y-1){$\Ab_2$};
    \node[draw=none] at (\x,\y-1.5) {\scalebox{1}{\textcolor{black}{$\vdots$}}};
    \node[tensor,fill=red!30!white] (A_N) at (\x,\y-2.5){$\Ab_d$};
    \draw  node[fill = black,circle,inner sep=0pt,minimum size=4pt] (m) at (\x+1,\y-1) {};
    \draw  node[fill = black,circle,inner sep=0pt,minimum size=4pt] (n) at (\x-1,\y-1) {};
    \draw[edge] (A_1) -- (m);
    \draw[edge] (A_2) -- (m);
    \draw[edge] (A_N) -- (m);
    \draw[edge] (A_1) -- (n);
    \draw[edge] (A_2) -- (n);
    \draw[edge] (A_N) -- (n);
    \draw[edge] (m) -- (\x+1.5,\y-1)node [midway,above] {\edgelabel{$n$}};
    \draw[edge] (n) -- (\x-1.5,\y-1)node [midway,above] {\edgelabel{$m$}};
    \end{tikzpicture}.
\end{align*}
\begin{remark}
这里列出 Hadamard 积的若干有用性质。
\begin{itemize}
    \item 元素取值于 $\{0,1\}$ 的张量 $\Tt$ 与自身的 Hadamard 积仍是 $\Tt$。特别地，这对复制张量成立：
    
    设 $\Tt\in\RR^{d_1\times d_2\times d_3}$ 为一个三阶复制张量。则在张量网络中，Hadamard 积 $\Tt\hadam\Tt$ 为：
    $$\Tt\hadam\Tt=
    \begin{tikzpicture}[baseline=\baseline]
   \tikzset{tensor/.style = {minimum size = 0.5cm,shape = circle,thick,draw=black,inner sep = 0pt}, edge/.style = {thick,line width=.4mm},every loop/.style={}}
    \def\y{0.2}
    \def\x{0}
     \draw  node[fill = black,circle,inner sep=0pt,minimum size=4pt] (m) at (\x,\y-0.5) {};
     \draw  node[draw=none] (a) at (\x-0.7,\y-0.5) {};
     \draw  node[draw=none] (b) at (\x+0.7,\y-0.5) {};
     \draw  node[draw=none] (c) at (\x,\y-1.25) {};
    \draw[edge] (a) -- (m) node [midway,above] {\scalebox{0.7}{\dimcolor{$d_1$}}};
    \draw[edge] (m) -- (b)node [midway,above] {\scalebox{0.7}{\dimcolor{$d_2$}}};
    \draw[edge] (m) -- (c)node [midway,left] {\scalebox{0.7}{\dimcolor{$d_3$}}};
    \draw  node[fill = black,circle,inner sep=0pt,minimum size=4pt] (m1) at (\x,\y+0.3) {};
     \draw  node[draw=none] (d) at (\x-0.7,\y+0.3) {};
     \draw  node[draw=none] (e) at (\x+0.7,\y+0.3) {};
     \draw  node[draw=none] (f) at (\x,\y+1.05) {};
    \draw[edge] (d) -- (m1)node [midway,below] {\scalebox{0.7}{\dimcolor{$d_1$}}};
    \draw[edge] (m1) -- (e)node [midway,below] {\scalebox{0.7}{\dimcolor{$d_2$}}};
    \draw[edge] (m1) -- (f)node [midway,left] {\scalebox{0.7}{\dimcolor{$d_3$}}};
    \node[draw=none] at (\x,\y-0.15){$\hadam$};
\end{tikzpicture}
=
    \begin{tikzpicture}[baseline=\baseline]
   \tikzset{tensor/.style = {minimum size = 0.5cm,shape = circle,thick,draw=black,inner sep = 0pt}, edge/.style = {thick,line width=.4mm},every loop/.style={}}
    \def\y{0}
    \def\x{0}
    \draw[edge](\x-0.7,\y-0.5)--(\x+0.7,\y-0.5);
    \draw[edge](\x-0.7,\y+0.5)--(\x+0.7,\y+0.5);
    \draw[edge](\x-0.7,\y-0.5)--(\x-0.7,\y+0.5);
    \draw[edge](\x+0.7,\y-0.5)--(\x+0.7,\y+0.5);
    \draw[edge](\x,\y+0.5) -- (\x,\y-0.5);
     \draw  node[fill = black,circle,inner sep=0pt,minimum size=4pt] (m1) at (\x-0.7,\y) {};
    \draw  node[fill = black,circle,inner sep=0pt,minimum size=4pt] (m) at (\x,\y) {};
    \draw  node[fill = black,circle,inner sep=0pt,minimum size=4pt] (m2) at (\x+0.7,\y) {};
    \draw[edge](m1)--(\x-1,\y)node [midway,above] {\scalebox{0.7}{\dimcolor{$d_1$}}};
    \draw[edge](m2)--(\x+0.4,\y)node [midway,above] {\scalebox{0.7}{\dimcolor{$d_3$}}};
    \draw[edge](m)--(\x-0.3,\y)node [midway,above] {\scalebox{0.7}{\dimcolor{$d_2$}}};
     \end{tikzpicture}$$
    
    \item 利用命题~\ref{prop:copy.tensor}~（第 4、5 两条），可以看出 
    $\vb\hadam\ub=$
         \begin{tikzpicture}[baseline=\baseline]
    \tikzset{tensor/.style = {minimum size = 0.5cm,shape = circle,thick,draw=black,inner sep = 0pt}, edge/.style   = {thick,line width=.4mm},every loop/.style={}}
    \def\y{0}
    \def\x{0}
    \draw  node[fill = black,circle,inner sep=0pt,minimum size=4pt] (m) at (\x+0.5,\y+0.5) {};
    \node[tensor,fill=blue!20!red!40!white] (v) at (\x,\y-0.65){\scalebox{0.85}{$\vb$}};
    \node[tensor,fill=red!20!blue!40!white] (u) at (\x+1,\y-0.65){\scalebox{0.85}{$\ub$}};

    \draw[edge] (m) -- (\x,\y) -- (v); 
    
    \draw[edge] (m) -- (\x+1,\y) -- (u); 
    \draw[edge](m) -- (\x+0.5,\y+1);
    \end{tikzpicture}
    $=\text{diag}(\vb)\ub$
\end{itemize}
\end{remark}
作为本章的收尾，我们在下一节介绍若干有用的结果。

\subsection{张量网络的若干有用性质}\label{sec:useful-TNs}
本节列出关于张量网络的若干有用结论。 

\paragraph{矩阵化秩的上界。}
任何张量网络结构都会给出底层张量各矩阵化秩的一族上界。更准确地说，对任意张量网络，给定矩阵化的秩不超过图中分离行腿与列腿的任一割的权重。这里，割的权重指构成该割的各边维度之积。下面将其形式化为一条定理。 

\begin{theorem}(张量网络的矩阵化的秩)\label{thm: matricization-rank}
设 $\T \in \Rbb^{d_1\times d_2\times \dots \times d_p}$ 以张量网络给出，且 $I \subset [p]$。那么在该 TN 中，对任意把 $I$ 中指标对应的腿与 $[p] \setminus I$ 中的腿分离开的割 $\mathcal{C}$，
$$\rank(\T_{(I)}) \leq \prod_{e \in  \mathrm{cutset}(\mathcal C)} \mathrm{dim}(e),$$
其中 $\mathrm{cutset}(\mathcal{C})$ 为\footnote{严格地说，TN 是一张图 $G=(V,E)$。$\Tt$ 的每个维度与 $V$ 中的一个顶点相关联：即相应悬空腿所连接的顶点。该图的任一割 $\mathcal C$ 是 $V$ 的一个划分：$\mathcal{C} = (V_1, V_2)$。定理中考虑的是把 $I$ 中的指标与其余指标分离开的割，即满足 $I \subseteq V_1$~（由于 $(V_1,V_2)$ 是 $V$ 的划分，故 $([p]\setminus I) \subseteq V_2$）。$\mathcal{C}$ 的割集是一端在 $V_1$、另一端在 $V_2$ 的边的集合：$\mathrm{cutset}(\mathcal{C}) = \{(u,v) \in E \mid u\in V_1, v\in V_2 \}$。 } %
$\mathcal C$ 的割集，且 $\mathrm{dim}(e)$ 为张量网络中边 $e$ 的维度。
\end{theorem}
下面用几个例子说明这一定理。首先，矩阵乘积秩的经典上界可作为一个特例得到。设   $\Ab\in\Rbb^{m\times R}$ 且 $\Bb\in\RR^{R\times n}$。对矩阵~（亦即二阶张量）$\mat{AB} \in \Rbb^{m \times n}$，该定理给出上界

$$\rank(\Ab\Bb)=\rank
\left(\begin{tikzpicture}[baseline=-0.5ex]
    \tikzset{tensor/.style = {minimum size = 0.5cm,shape = circle,thick,draw=black,inner sep = 0pt}, edge/.style = {thick,line width=.4mm},every loop/.style={}}
    \def\x{6}
    \def\y{0} \node[tensor,fill=green!50!red!50!white] (A) at (\x+1,\y) {$\Ab$};
    \node[tensor,fill=green!50!red!50!white] (B) at (\x+2.5,\y) {$\Bb$};
    \draw[edge] (A) -- (\x+0.3,0) node [midway,above] {\scalebox{0.7}{\dimcolor{$m$}}};
    \draw[edge] (A) -- (B)node [above,midway] {\scalebox{0.7}{\dimcolor{$R$}}};;
    \draw[edge] (B) -- (\x+3.3,\y) node [midway,above] {\scalebox{0.7}{\dimcolor{$n$}}};
    \draw[dotted] (\x+1.6,\y+1) -- (\x+1.6,\y-1);
\end{tikzpicture}\right)\leq R.$$

再考虑如下 $d_1\times d_2$ 矩阵的张量网络表示：
$$
\begin{tikzpicture}[baseline=-0.5ex]
    \tikzset{tensor/.style = {minimum size = 0.5cm,shape = circle,thick,draw=black,inner sep = 0pt}, edge/.style = {thick,line width=.4mm},every loop/.style={}}
    \def\x{6}
    \def\y{0} \node[tensor,fill=green!50!red!50!white] (T) at (\x+1,\y) {$\Tt$};
    \node[tensor,fill=green!50!red!50!white] (A) at (\x+2.5,\y) {$\At$};
    \node[tensor,fill=green!50!red!50!white] (B) at (\x+3.5,\y) {$\Bb$};
      \node[tensor,fill=green!50!red!50!white] (S) at (\x+4.5,\y) {$\Sb$};
    \draw[edge] (T) -- (\x+0.3,\y)node [midway,above] {\scalebox{0.7}{\dimcolor{$d_1$}}};
     \draw[edge] (T) -- (A);
    \draw[edge] (A) -- (B)node [midway,above] {\scalebox{0.7}{\dimcolor{$R_4$}}};
    \draw[edge] (T) to [out=-45,in=-135] (A);
    \draw[edge] (T) to [out=45,in=135] (A);
    \node[draw=none] () at (\x+1.7,\y+0.15) {\scalebox{0.7}{\dimcolor{$R_1$}}};
    \node[draw=none] () at (\x+1.7,\y+0.6) {\scalebox{0.7}{\dimcolor{$R_2$}}};
    \node[draw=none] () at (\x+1.7,\y-0.6) {\scalebox{0.7}{\dimcolor{$R_3$}}};
    \draw[edge](B) -- (S) node [midway,above] {\scalebox{0.7}{\dimcolor{$R_5$}}};
    \draw[edge] (S) -- (\x+5.3,\y) node [midway,above] {\scalebox{0.7}{\dimcolor{$d_2$}}};
    \end{tikzpicture}$$
其中 $\Tt\in\RR^{d_1\times R_1\times R_2\times R_3}$，$\At\in\RR^{R_1\times R_2\times R_3\times R_4}$，$\Bb\in\RR^{R_4\times R_5}$ 且 $\Sb\in\RR^{R_5\times d_2}$。
例如，定理~\ref{thm: matricization-rank} 可用来给出上界
$$
\rank\left(\begin{tikzpicture}[baseline=-0.5ex]
    \tikzset{tensor/.style = {minimum size = 0.5cm,shape = circle,thick,draw=black,inner sep = 0pt}, edge/.style = {thick,line width=.4mm},every loop/.style={}}
    \def\x{6}
    \def\y{0} \node[tensor,fill=green!50!red!50!white] (T) at (\x+1,\y) {$\Tt$};
    \node[tensor,fill=green!50!red!50!white] (A) at (\x+2.5,\y) {$\At$};
    \node[tensor,fill=green!50!red!50!white] (B) at (\x+3.5,\y) {$\Bb$};
      \node[tensor,fill=green!50!red!50!white] (S) at (\x+4.5,\y) {$\Sb$};
    \draw[edge] (T) -- (\x+0.3,\y)node [midway,above] {\scalebox{0.7}{\dimcolor{$d_1$}}};
     \draw[edge] (T) -- (A);
    \draw[edge] (A) -- (B)node [midway,above] {\scalebox{0.7}{\dimcolor{$R_4$}}};
    \draw[edge] (T) to [out=-45,in=-135] (A);
    \draw[edge] (T) to [out=45,in=135] (A);
    \node[draw=none] () at (\x+1.7,\y+0.15) {\scalebox{0.7}{\dimcolor{$R_1$}}};
    \node[draw=none] () at (\x+1.7,\y+0.6) {\scalebox{0.7}{\dimcolor{$R_2$}}};
    \node[draw=none] () at (\x+1.7,\y-0.6) {\scalebox{0.7}{\dimcolor{$R_3$}}};
    \draw[edge](B) -- (S) node [midway,above] {\scalebox{0.7}{\dimcolor{$R_5$}}};
    \draw[edge] (S) -- (\x+5.3,\y) node [midway,above] {\scalebox{0.7}{\dimcolor{$d_2$}}};
    \draw[dotted] (\x+1.45,\y+1) -- (\x+1.45,\y-1);
    \end{tikzpicture}\right)\leq R_1R_2R_3$$
以及
$$
\rank
\left(\begin{tikzpicture}[baseline=-0.5ex]
    \tikzset{tensor/.style = {minimum size = 0.5cm,shape = circle,thick,draw=black,inner sep = 0pt}, edge/.style = {thick,line width=.4mm},every loop/.style={}}
    \def\x{6}
    \def\y{0} \node[tensor,fill=green!50!red!50!white] (T) at (\x+1,\y) {$\Tt$};
    \node[tensor,fill=green!50!red!50!white] (A) at (\x+2.5,\y) {$\At$};
    \node[tensor,fill=green!50!red!50!white] (B) at (\x+3.5,\y) {$\Bb$};
      \node[tensor,fill=green!50!red!50!white] (S) at (\x+5,\y) {$\Sb$};
    \draw[edge] (T) -- (\x+0.3,\y)node [midway,above] {\scalebox{0.7}{\dimcolor{$d_1$}}};
     \draw[edge] (T) -- (A);
    \draw[edge] (A) -- (B)node [midway,above] {\scalebox{0.7}{\dimcolor{$R_4$}}};
    \draw[edge] (T) to [out=-45,in=-135] (A);
    \draw[edge] (T) to [out=45,in=135] (A);
    \node[draw=none] () at (\x+1.7,\y+0.15) {\scalebox{0.7}{\dimcolor{$R_1$}}};
    \node[draw=none] () at (\x+1.7,\y+0.6) {\scalebox{0.7}{\dimcolor{$R_2$}}};
    \node[draw=none] () at (\x+1.7,\y-0.6) {\scalebox{0.7}{\dimcolor{$R_3$}}};
    \draw[edge](B) -- (S) node [midway,above] {\scalebox{0.7}{\dimcolor{$R_5$}}};
    \draw[edge] (S) -- (\x+5.75,\y) node [midway,above] {\scalebox{0.7}{\dimcolor{$d_2$}}};
    \draw[dotted] (\x+4.4,\y+1) -- (\x+4.4,\y-1);
    \end{tikzpicture}\right)\leq R_5.
$$
最后一个例子：设 $\Gt\in\RR^{R_1\times\cdots\times R_4}$，$\Ab_1\in\RR^{R_1\times d_1}, \Ab_2\in\RR^{R_2\times d_2},\Ab_3\in\RR^{R_3\times d_3}$ 且 $\Ab_4\in\RR^{R_4\times d_4}$，并考虑张量 $\T = \Gt \times_1 \Ab_1\times_2 \Ab_2\times_3 \Ab_3\times_4 \Ab_4$。定理~\ref{thm: matricization-rank} 例如可用来给出 $\Tt$ 的第一种矩阵化的秩的上界：
$$
\rank(\Tt_{(1)}) = 
\rank\left(
\begin{tikzpicture}[baseline=-0.5ex]
    \tikzset{tensor/.style = {minimum size = 0.5cm,shape = circle,thick,draw=black,inner sep = 0pt}, edge/.style = {thick,line width=.4mm},every loop/.style={}}
    \def\x{6}
    \def\y{0} 
    \node[tensor,fill=red!50!white] (G) at (\x+4.5,\y+1) {$\Gt$};
    \node[tensor,fill=red!20!white] (A_1) at (\x+3,\y){$\Ab_1$};
    \node[tensor,fill=red!40!white] (A_2) at (\x+4,\y){$\Ab_2$};
    \node[tensor, fill=blue!40!white] (A_3) at (\x+5,\y){$\Ab_3$};
    \node[tensor,fill=blue!20!white] (A_4) at (\x+6.5,\y){$\Ab_4$};
    \draw[edge] (G) -- (A_1) node [midway, above] {\scalebox{0.7}{\dimcolor{$R_1$}}};
    \draw[edge] (G) -- (A_2) node [midway, right] {\scalebox{0.7}{\dimcolor{$R_2$}}};
    \draw[edge] (G) -- (A_3) node [midway, right] {\scalebox{0.7}{\dimcolor{$R_3$}}};
    \draw[edge] (G) -- (A_4) node [midway, above] {\scalebox{0.7}{\dimcolor{$R_4$}}};
    \draw[edge] (A_1) -- (\x+3,\y-0.7) node [midway, left] {\scalebox{0.7}{\dimcolor{$d_1$}}};
    \draw[edge] (A_3) -- (\x+5,\y-0.7) node [midway, right=-0.11cm] {\scalebox{0.7}{\dimcolor{$d_3$}}};
    \draw[edge] (A_2) -- (\x+4,\y-0.35) -- (\x+4.85,\y-0.35) -- (\x+4.85,\y-0.7) node [midway, left] {\scalebox{0.7}{\dimcolor{$d_2$}}};
    \draw[edge] (A_4) -- (\x+6.5,\y-0.35) -- (\x+5.30,\y-0.35) -- (\x+5.30,\y-0.7) node [midway, right] {\scalebox{0.7}{\dimcolor{$d_4$}}};
    \draw[dotted](\x+3.5,\y+1)-- (\x+3.5,\y-0.5);
    \end{tikzpicture}\right)\leq R_1.
$$
\paragraph{TN 的边与和式张量。}设 $\T$ 为以 TN 给出的张量。TN 中的任意一条边，都对应一种把该张量表示为若干个与 $\T$ 同尺寸的张量之和的方式。这直接来自 TN 中边的定义。一条边表示一次收缩，而收缩就是一个求和。下面用例子说明这一概念。首先，考虑 $m\times n$ 矩阵的秩为 $R$ 的分解。TN 中的那条边对应把该矩阵表示为 $R$ 个（秩一）矩阵之和：  
$$
\begin{tikzpicture}[baseline=-0.5ex]
    \tikzset{tensor/.style = {minimum size = 0.5cm,shape = circle,thick,draw=black,inner sep = 0pt}, edge/.style = {thick,line width=.4mm},every loop/.style={}}
    \def\x{6}
    \def\y{0} \node[tensor,fill=green!50!red!50!white] (A) at (\x+1,\y) {$\Ab$};
    \node[tensor,fill=green!50!red!50!white] (B) at (\x+2.5,\y) {$\Bb$};
    \draw[edge] (A) -- (\x+0.3,0) node [midway,above] {\scalebox{0.7}{\dimcolor{$m$}}};
    \draw[edge] (A) -- (B)node [above,midway] {\scalebox{0.6}{\dimcolor{$R$}}};;
    \draw[edge] (B) -- (\x+3.3,\y) node [midway,above] {\scalebox{0.7}{\dimcolor{$n$}}};
    \draw[dotted] (\x+1.6,\y+1) -- (\x+1.6,\y-1);
\end{tikzpicture}
= 
\sum_{r=1}^R
\begin{tikzpicture}[baseline=-0.5ex]
    \tikzset{tensor/.style = {minimum size = 0.5cm,shape = circle,thick,draw=black,inner sep = 0pt}, edge/.style = {thick,line width=.4mm},every loop/.style={}}
    \def\x{6}
    \def\y{0} \node[tensor,fill=green!50!red!50!white] (A) at (\x+1,\y) {$\Ab$};
    \node[tensor,fill=green!50!red!50!white] (B) at (\x+2.5,\y) {$\Bb$};
    \draw[edge] (A) -- (\x+0.3,0) node [midway,above] {\scalebox{0.7}{\dimcolor{$m$}}};
    \draw[edge] (A) -- (\x+1.5,\y) node [pos=1.5] {\scalebox{0.7}{\indcolor{$r$}}};
    \draw[edge] (B) -- (\x+3.3,\y) node [midway,above] {\scalebox{0.7}{\dimcolor{$n$}}};
    \draw[edge](B) -- (\x+2,\y)node [pos=1.4] {\scalebox{0.7}{\indcolor{$r$}}};
\end{tikzpicture}
= \sum_{r=1}^R \Ab_{:,r} \Bb_{r,:} = \sum_{r=1}^R \ab_r \outprod \bb_r,
$$
其中 $\ab_r$ 为 $\Ab\in\RR^{m\times R}$ 的各列，$\bb_r$ 为 $\Bb\in\RR^{R\times n}$ 的各行。

类似地，表示矩阵的迹的那个 TN 中唯一的边，对应于把这个 TN（它是一个标量）写成若干标量（矩阵的对角元）之和，设 $\Ab\in\RR^{d\times d}$：
$$\begin{tikzpicture}[baseline=\baseline]
    \tikzset{tensor/.style = {minimum size = 0.5cm,shape = circle,thick,draw=black,inner sep = 0pt}, edge/.style = {thick,line width=.4mm},every loop/.style={}}
    \def\y{0}
    \def\x{0}
    \node[tensor,fill=red!50!white] (A) at (\x,\y+0.2){\scalebox{0.85}{$\Ab$}};
    \draw[edge] (A) to [out=-20,in=-160, distance=1.35cm] (A);
    \draw[dotted] (\x-0.15,\y-0.1) -- (\x-0.15,\y-0.8);
    \node[draw=none] () at (\x,\y-0.4) {\scalebox{0.7}{\dimcolor{$d$}}};
\end{tikzpicture}\hspace{-0.7cm}
=
\sum_{i=1}^d\left(
\begin{tikzpicture}[baseline=\baseline]
    \tikzset{tensor/.style = {minimum size = 0.5cm,shape = circle,thick,draw=black,inner sep = 0pt}, edge/.style = {thick,line width=.4mm},every loop/.style={}}
    \def\y{0}
    \def\x{0}
    \node[tensor,fill=red!50!white] (A) at (\x,\y){\scalebox{0.85}{$\Ab$}};
     \draw[edge] (\x-0.25,\y) to [out=110,in=60, distance=0.1cm] (\x-0.6,\y-0.2);    
    \draw[edge] (\x+0.25,\y) to [out=110,in=60, distance=0.1cm] (\x+0.6,\y-0.2); 
    \node[draw=none] () at (\x+0.7,\y-0.25) {\scalebox{0.7}{\indcolor{$i$}}};
    \node[draw=none] () at (\x-0.7,\y-0.25) {\scalebox{0.7}{\indcolor{$i$}}};
    \end{tikzpicture}\right)=\sum_{i=1}^d\Ab_{ii}.
    $$
现在考虑一个更有趣的例子：
$$\begin{tikzpicture}[baseline=\baseline]
    \tikzset{tensor/.style = {minimum size = 0.5cm,shape = circle,thick,draw=black,inner sep = 0pt}, edge/.style = {thick,line width=.4mm},every loop/.style={}}
    \def\y{0}
    \def\x{0}
    \node[tensor,fill=red!50!white] (A) at (\x,\y+0.5){\scalebox{0.85}{$\Ab$}};
    \node[tensor,fill=red!50!white] (B) at (\x,\y-0.5){\scalebox{0.85}{$\Bb$}};
    \draw[edge](A)--(B)node [midway, right] {\scalebox{0.7}{\dimcolor{$R$}}};
    \draw[edge](A) -- (\x+0.5,\y+0.5)--(\x+0.5,\y)--(\x+1,\y)node [midway, above] {\scalebox{0.7}{\dimcolor{$n$}}};
    \draw[edge](A) -- (\x-0.5,\y+0.5)--(\x-0.5,\y)--(\x-1,\y)node [midway, above] {\scalebox{0.7}{\dimcolor{$m$}}};
    \draw[edge](B) -- (\x+0.5,\y-0.5)--(\x+0.5,\y-0.1)--(\x+1,\y-0.1)node [midway, below] {\scalebox{0.7}{\dimcolor{$q$}}};
    \draw[edge](B) -- (\x-0.5,\y-0.5)--(\x-0.5,\y-0.1)--(\x-1,\y-0.1)node [midway, below] {\scalebox{0.7}{\dimcolor{$p$}}};
    \end{tikzpicture}.$$
这个 TN 中的边对应于一种表达矩阵的方式：把该 TN 所表示的矩阵写成若干矩阵之和：
$$\begin{tikzpicture}[baseline=\baseline]
    \tikzset{tensor/.style = {minimum size = 0.5cm,shape = circle,thick,draw=black,inner sep = 0pt}, edge/.style = {thick,line width=.4mm},every loop/.style={}}
    \def\y{0}
    \def\x{0}
    \node[tensor,fill=red!50!white] (A) at (\x,\y+0.5){\scalebox{0.85}{$\Ab$}};
    \node[tensor,fill=red!50!white] (B) at (\x,\y-0.5){\scalebox{0.85}{$\Bb$}};
    \draw[edge](A)--(B)node [midway, right] {\scalebox{0.7}{\dimcolor{$R$}}};
    \draw[edge](A) -- (\x+0.5,\y+0.5)--(\x+0.5,\y)--(\x+1,\y)node [midway, above] {\scalebox{0.7}{\dimcolor{$n$}}};
    \draw[edge](A) -- (\x-0.5,\y+0.5)--(\x-0.5,\y)--(\x-1,\y)node [midway, above] {\scalebox{0.7}{\dimcolor{$m$}}};
    \draw[edge](B) -- (\x+0.5,\y-0.5)--(\x+0.5,\y-0.1)--(\x+1,\y-0.1)node [midway, below] {\scalebox{0.7}{\dimcolor{$q$}}};
    \draw[edge](B) -- (\x-0.5,\y-0.5)--(\x-0.5,\y-0.1)--(\x-1,\y-0.1)node [midway, below] {\scalebox{0.7}{\dimcolor{$p$}}};
    \draw[dotted](\x+0.35,\y+0.15) -- (\x-0.35,\y+0.15);
    \end{tikzpicture}
    =
    \sum_{r=1}^R\begin{tikzpicture}[baseline=\baseline]
    \tikzset{tensor/.style = {minimum size = 0.5cm,shape = circle,thick,draw=black,inner sep = 0pt}, edge/.style = {thick,line width=.4mm},every loop/.style={}}
    \def\y{0}
    \def\x{0}
    \node[tensor,fill=red!50!white] (A) at (\x,\y+1){\scalebox{0.85}{$\Ab$}};
    \node[tensor,fill=red!50!white] (B) at (\x,\y-1){\scalebox{0.85}{$\Bb$}};
    \draw[edge](A) -- (\x,\y+0.45)node [pos=1.4]{\scalebox{0.7}{\indcolor{$r$}}};
    \draw[edge](B) -- (\x,\y-0.45)node [pos=1.3]{\scalebox{0.7}{\indcolor{$r$}}};
    \draw[edge](A) -- (\x+0.5,\y+1)--(\x+0.5,\y)--(\x+1,\y)node [midway, above] {\scalebox{0.7}{\dimcolor{$n$}}};
    \draw[edge](A) -- (\x-0.5,\y+1)--(\x-0.5,\y)--(\x-1,\y)node [midway, above] {\scalebox{0.7}{\dimcolor{$m$}}};
    \draw[edge](B) -- (\x+0.5,\y-1)--(\x+0.5,\y-0.1)--(\x+1,\y-0.1)node [midway, below] {\scalebox{0.7}{\dimcolor{$q$}}};
    \draw[edge](B) -- (\x-0.5,\y-1)--(\x-0.5,\y-0.1)--(\x-1,\y-0.1)node [midway, below] {\scalebox{0.7}{\dimcolor{$p$}}};
    \end{tikzpicture}
    =
    \sum_{r=1}^R\Ab_r\kron\Bb_r.$$
由此可见，这个 TN 对应的正是 Kronecker 积之和！ 
\paragraph{张量、多重线性映射与多项式。}
可以把张量看作\emph{多重线性映射的表示}，其含义与矩阵表示线性映射完全相同。回顾一下，矩阵 $\Ab\in\RR^{m\times n}$ 通过 $\xb\mapsto \Ab\xb$ 表示一个线性映射 $\RR^n\to\RR^m$。用张量网络记法，这就是
$$
\yb = \Ab\xb
=
\begin{tikzpicture}[baseline=-0.5ex]
    \tikzset{tensor/.style = {minimum size = 0.5cm,shape = circle,thick,draw=black,inner sep = 0pt},
             edge/.style   = {thick,line width=.4mm},every loop/.style={}}
    \def\x{0}\def\y{0}
    \node[tensor,fill=red!50!white] (A) at (\x,\y){$\Ab$};
    \node[tensor,fill=blue!20!white] (x) at (\x+1,\y){$\xb$};
    \draw[edge] (A) -- (x) node [midway,above]{\scalebox{0.7}{\dimcolor{$n$}}};
    \draw[edge] (A) -- (\x-1,\y) node [midway,above]{\scalebox{0.7}{\dimcolor{$m$}}};
\end{tikzpicture}
\in\RR^{m}.
$$
线性是指对一切 $\alpha,\beta\in\RR$ 与 $\xb,\zb\in\RR^n$ 都有 $\Ab(\alpha \xb+\beta \zb)=\alpha \Ab\xb+\beta \Ab\zb$。当映射\emph{对每个自变量分别线性}时，张量便是其自然推广。具体地说，一个 $N$ 阶张量 $\Tt\in\RR^{d_1\times\cdots\times d_N}$ 定义了一个映射：它接受 $N$ 个输入向量，并沿所有模收缩而给出一个标量：
$$
f(\xb^{(1)},\dots,\xb^{(N)})
:= \langle \Tt,\ \xb^{(1)}\outprod \cdots \outprod \xb^{(N)}\rangle
=
\sum_{i_1=1}^{d_1}\cdots\sum_{i_N=1}^{d_N}\Tt_{i_1\cdots i_N}\,\xb^{(1)}_{i_1}\cdots \xb^{(N)}_{i_N}.
$$
该映射是\emph{多重线性}的：固定其余向量时，它对每个 $\xb^{(k)}$ 都是线性的。在张量网络图中，这就是
\[
f(\ab,\bb,\dots,\zb)
=
\begin{tikzpicture}[baseline=2ex]
    \tikzset{tensor/.style = {minimum size = 0.5cm,shape = circle,thick,draw=black,inner sep = 0pt},
             edge/.style   = {thick,line width=.4mm},
             every loop/.style={}}
    \def\x{0}\def\y{0}

    \node[tensor,fill=green!50!red!50!white] (T) at (\x,\y+0.8){$\Tt$};

    \node[tensor,fill=blue!20!white] (a) at (\x-1.2,\y){$\ab$};
    \node[tensor,fill=blue!20!white] (b) at (\x-0.5,\y){$\bb$};
    \node[draw=none] (dots) at (\x+0.5,\y){$\cdots$};
    \node[tensor,fill=blue!20!white] (z) at (\x+1.2,\y){$\zb$};

    \draw[edge] (T) -- (a) node [near end,above]{\scalebox{0.7}{\dimcolor{$d_1\ $}}};
    \draw[edge] (T) -- (b) node [near end,right]{\scalebox{0.7}{\dimcolor{$d_2$}}};
    \draw[edge] (T) -- (z) node [near end,above]{\scalebox{0.7}{\dimcolor{$d_N$}}};
\end{tikzpicture}
\in\RR,
\]
其中 $\ab, \bb,\cdots,\zb\in\RR^d$。
若留出一条腿不作收缩，便得到一个\emph{向量值}多重线性映射。例如，若 $\Tt\in\RR^{d_1\times\cdots\times d_N \times p}$，则除最后一个模之外对所有的模收缩，就定义了一个输出为 $p$ 维向量的多重线性映射：
$$
f(\ab,\bb,\dots,\zb)
=
\begin{tikzpicture}[baseline=2ex]
    \tikzset{tensor/.style = {minimum size = 0.5cm,shape = circle,thick,draw=black,inner sep = 0pt},
             edge/.style   = {thick,line width=.4mm},
             every loop/.style={}}
    \def\x{0}\def\y{0}

    \node[tensor,fill=green!50!red!50!white] (T) at (\x,\y+0.8){$\Tt$};

    \node[tensor,fill=blue!20!white] (a) at (\x-1.2,\y){$\ab$};
    \node[tensor,fill=blue!20!white] (b) at (\x-0.5,\y){$\bb$};
    \node[draw=none] (dots) at (\x+0.5,\y){$\cdots$};
    \node[tensor,fill=blue!20!white] (z) at (\x+1.2,\y){$\zb$};

    \draw[edge] (T) -- (a) node [near end,above]{\scalebox{0.7}{\dimcolor{$d_1\ $}}};
    \draw[edge] (T) -- (b) node [near end,right]{\scalebox{0.7}{\dimcolor{$d_2$}}};
    \draw[edge] (T) -- (z) node [near end,above]{\scalebox{0.7}{\dimcolor{$d_N$}}};
    \draw[edge] (T) -- (\x,\y+1.6) node [midway,right]{\scalebox{0.7}{\dimcolor{$p$}}};
\end{tikzpicture}
\in\RR^p.
$$

\vspace{0.2cm}
\noindent
这就把张量直接与\emph{多元多项式}联系了起来。上面的标量映射
\[
f(\xb^{(1)},\dots,\xb^{(N)})=\sum_{i_1,\dots,i_N}\Tt_{i_1\cdots i_N}\,\xb^{(1)}_{i_1}\cdots \xb^{(N)}_{i_N},
\]
是关于各向量分量的多项式，并且是\emph{多重线性}的：每个变量出现的次数至多为一（就每一组变量 $\xb^{(k)}$ 而言如此）。在一个重要的特殊情形中，所有输入都是\emph{同一个}向量 $\xb\in\RR^{d}$，而 $\Tt\in\RR^{d\times\cdots\times d}$ 有 $N$ 个模，此时收缩
\[
\begin{tikzpicture}[baseline=2ex]
    \tikzset{tensor/.style = {minimum size = 0.5cm,shape = circle,thick,draw=black,inner sep = 0pt},
             edge/.style   = {thick,line width=.4mm},
             every loop/.style={}}
    \def\x{0}\def\y{0}

    \node[tensor,fill=green!50!red!50!white] (T) at (\x,\y+0.8){$\Tt$};

    \node[tensor,fill=blue!20!white] (a) at (\x-1.2,\y){$\xb$};
    \node[tensor,fill=blue!20!white] (b) at (\x-0.5,\y){$\xb$};
    \node[draw=none] (dots) at (\x+0.5,\y){$\cdots$};
    \node[tensor,fill=blue!20!white] (z) at (\x+1.2,\y){$\xb$};

    \draw[edge] (T) -- (a) node [midway,left]{\scalebox{0.7}{\dimcolor{$d\ $}}};
    \draw[edge] (T) -- (b) node [near end,right]{\scalebox{0.7}{\dimcolor{$d$}}};
    \draw[edge] (T) -- (z) node [midway,right]{\scalebox{0.7}{\dimcolor{$\ d$}}};
\end{tikzpicture}
= \Tt \ttm{1}\xb \ttm{2}\xb \cdots \ttm{N}\xb
=
\sum_{i_1,\dots,i_N}\Tt_{i_1\cdots i_N}\,\xb_{i_1}\cdots \xb_{i_N} \in \RR.
\]
是一个 $N$ 次\emph{齐次}多元多项式（每个单项式的总次数恰为 $N$）。张量的分量恰好是这一多项式在单项式基 $(\xb_{i_1}\cdots \xb_{i_N})_{i_1,\cdots,i_N}$ 下的系数。

反之，若把 $(N+1)$ 阶张量 $\Tt\in\RR^{ d\times\cdots\times d \times p}$ 与 $N$ 份 $\xb\in\RR^d$ 收缩，并让最后一条腿保持自由，就得到一个向量值的齐次多项式映射 $\RR^d\to\RR^m$：
\[
p(\xb) = \Tt \ttm{1}\xb \ttm{2}\xb\cdots \ttm{N}\xb \in\RR^{p}.
\]

\vspace{0.2cm}
\noindent
最后，\emph{非齐次}多项式可以用一个小技巧得到：\emph{在输入向量末尾添上一个常数 $1$}。设 $\xb\in\RR^d$，并定义增广向量 $\tilde{\xb} =\begin{bmatrix}\xb\\1\end{bmatrix}\in\RR^{d+1}$。若 $\widetilde{\Tt}\in\RR^{(d+1)\times\cdots\times(d+1)}$ 是一个 $N$ 阶张量，则
\[
\tilde{p}(\xb)
:=
\widetilde{\Tt} \ttm{1}\tilde{\xb}\ttm{2}\tilde{\xb}\cdots\ttm{N}\tilde{\xb},
\]
是 $\xb$ 的不超过 $N$ 次的多项式（某些因子会取到最后一个坐标，也就是常数 $1$，从而降低相应单项式的次数）。换句话说，增广变量 $(x_1,\dots,x_d,1)$ 的一个齐次张量多项式，就紧凑地表示了 $(x_1,\dots,x_d)$ 的一般多元多项式。这一观点在实践中尤为有用：只需给每个模添加一个额外的维度并把它固定为 $1$，就能用与齐次情形相同的收缩机制来处理偏置项与低次交互项。  

\section{张量分解}\label{chapter:tensor-decomposition}
处理高阶张量的计算代价很高，因为元素个数随张量的阶呈指数增长。为应对这一问题，张量分解应运而生，成为有力而高效的工具。
与矩阵分解类似，张量分解把一个高阶张量拆成阶数更低、元素更少的若干小分量，从而使计算更易处理。但矩阵秩的概念可以以多种方式推广到张量，每种推广都对应一种不同的因子化/分解模型。
本章介绍最广为人知的几种张量分解模型。
\subsection{并行因子/CP 分解（CANDECOMP/PARAFAC~Decomposition）}\label{sec:cp-decomposition}
$N$ 阶张量 $\At\in\RR^{d_1\times\cdots\times d_N}$ 的 CP 分解~\citep{hitchcock1927expression}把 $\At$ 分解为有限个秩一张量之和。等价地，它是 $R$ 个秩一张量的线性组合，其中 $R$ 称为该分解的秩：
$$\At = \sum_{r=1}^R \ab^{(r)}_1\outprod\dots\outprod\ab^{(r)}_N ,$$
其中对每个 $n\in [N]$，都有 $\ab^{(1)}_n,\dots,\ab^{(R)}_n\in\RR^{d_n}$。

把所有向量 $\ab^{(r)}_n$ 排成因子矩阵，
\[\Ab_1 = {\begin{pmatrix}
\ab^{(1)}_1 & \dots & \ab^{(R)}_1 \\
\end{pmatrix}}{\in\RR^{d_1\times R}}, \dots ,\Ab_N = {\begin{pmatrix}
\ab^{(1)}_N & \dots & \ab^{(R)}_N \\
\end{pmatrix}}{\in\RR^{d_N\times R}}, \]
CP 分解可简洁地记作 $\At = \CP{\Ab_1,\cdots,\Ab_N}$。

在张量网络中，CP 分解 $\At = \CP{\Ab,\Bb,\Cb}$ 表示为
$$
\begin{tikzpicture}[baseline=\baseline]
    \tikzset{tensor/.style = {minimum size = 0.5cm,shape = circle,thick,draw=black,inner sep = 0pt}, edge/.style   = {thick,line width=.4mm},every loop/.style={}}
    \def\y{0}
    \def\x{0}
    \node[tensor,fill=green!50!red!50!white] (A) at (\x,\y){\scalebox{0.85}{$\At$}};
    \draw[edge] (\x+0.25,\y) -- (\x+0.6,\y)  node [midway,above] {\scalebox{0.7}{\dimcolor{$d_2$}}};
    \draw[edge] (\x-0.25,\y) -- (\x-0.6,\y)  node [midway,above] {\scalebox{0.7}{\dimcolor{$d_1$}}};
    \draw[edge] (\x,\y-0.25) -- (\x,\y-0.7)  node [midway,left] {\scalebox{0.7}{\dimcolor{$d_3$}}};
    \end{tikzpicture}
=
\begin{tikzpicture}[baseline=\baseline]
    \tikzset{tensor/.style = {minimum size = 0.5cm,shape = circle,thick,draw=black,inner sep = 0pt}, edge/.style   = {thick,line width=.4mm},every loop/.style={}}
    \def\y{0}
    \def\x{0}
  \node[tensor,fill=red!50!white] (A) at (\x,\y){\scalebox{0.85}{$\Ab$}};
  \node[tensor,fill=blue!50!white] (B) at (\x+1,\y){$\Bb$};
   \node[tensor,fill=blue!20!white] (C) at (\x+2,\y){$\Cb$};
    \draw  node[fill = black,circle,inner sep=0pt,minimum size=4pt] (m) at (\x+1,\y+0.7) {};
    \draw[edge] (A) -- (m);
    \node[draw=none] () at (\x+0.4,\y+0.5) {\scalebox{0.7}{\dimcolor{$R$}}};
    \draw[edge] (B) -- (m) node [midway,left=-0.09cm] {\scalebox{0.7}{\dimcolor{$R$}}};
    \draw[edge] (C) -- (m);
    \node[draw=none] () at (\x+1.6,\y+0.5) {\scalebox{0.7}{\dimcolor{$R$}}};
    \draw[edge](A) -- (\x,\y-0.7)node [midway,left] {\edgelabel{$d_1$}};
    \draw[edge](B) -- (\x+1,\y-0.7)node [midway,left] {\edgelabel{$d_2$}};
    \draw[edge](C) -- (\x+2,\y-0.7)node [midway,left] {\edgelabel{$d_3$}};
\end{tikzpicture},
$$
其中黑点是第~\ref{sec:copy-tensor}~节引入的三阶复制张量。事实上，有
\begin{align*}
    \left(\begin{tikzpicture}[baseline=\baseline]
    \tikzset{tensor/.style = {minimum size = 0.5cm,shape = circle,thick,draw=black,inner sep = 0pt}, edge/.style   = {thick,line width=.4mm},every loop/.style={}}
    \def\y{0}
    \def\x{0}
  \node[tensor,fill=red!50!white] (A) at (\x,\y){\scalebox{0.85}{$\Ab$}};
  \node[tensor,fill=blue!50!white] (B) at (\x+1,\y){$\Bb$};
   \node[tensor,fill=blue!20!white] (C) at (\x+2,\y){$\Cb$};
    \draw  node[fill = black,circle,inner sep=0pt,minimum size=4pt] (m) at (\x+1,\y+0.7) {};
    \draw[edge] (A) -- (m);
    \node[draw=none] () at (\x+0.4,\y+0.5) {\scalebox{0.7}{\dimcolor{$R$}}};
    \draw[edge] (B) -- (m) node [midway,left=-0.09cm] {\scalebox{0.7}{\dimcolor{$R$}}};
    \draw[edge] (C) -- (m);
    \node[draw=none] () at (\x+1.6,\y+0.5) {\scalebox{0.7}{\dimcolor{$R$}}};
    \draw[edge](A) -- (\x,\y-0.7);
    \draw[edge](B) -- (\x+1,\y-0.7);
    \draw[edge](C) -- (\x+2,\y-0.7);
\end{tikzpicture}\right)_{ijk} 
&= \sum_{r_1=1}^R\sum_{r_2=1}^R\sum_{r_3=1}^R\delta_{r_1r_2r_3}\Ab_{ir_1}\Bb_{jr_2}\Cb_{kr_3}
=
\sum_{r=1}^R\Ab_{ir}\Bb_{jr}\Cb_{kr}\\
&=
\sum_{r=1}^R(\ab_r)_i(\bb_r)_j(\cbb_r)_k 
=
\sum_{r=1}^R(\ab_r\outprod\bb_r\outprod\cbb_r)_{ijk}\\
&=
\left(\sum_{r=1}^R\ab_r\outprod\bb_r\outprod\cbb_r\right)_{ijk}
\end{align*}
\paragraph{张量的 CP 秩}
张量秩最基本也是最早的概念是 CP 秩，最早由~\citep{hitchcock1927expression}提出。张量的 CP 秩定义为：把该张量表示为若干秩一张量之和所需的最少项数。\footnote{注意本文术语中「CP 分解的秩」与「张量的 CP 秩」有一细微差别：前者的秩指形如 $\At = \sum_{r=1}^R \ab^{(r)}_1\outprod\dots\outprod\ab^{(r)}_N$ 的分解中求和项的个数，而张量 $\At$ 的 CP 秩是使这样的分解存在的最小的 $R$。因此，若 $\At$ 有一个秩 $R$ 的 CP 分解，其 CP 秩至多为 $R$，也可能更小。}
尽管这一定义与矩阵秩的定义相似，张量秩的性质却与之大不相同。从计算的角度看，一个关键区别在于：与矩阵秩不同，不存在能在多项式时间内确定张量 CP 秩的算法。事实上，计算张量的秩是 NP 难问题~\citep{hillar2013most}。
张量秩还有若干变体，每个变体都与某种特定的张量分解相对应。随着本章的推进，我们会介绍其中几种不同的秩。


\begin{remark}
\newcommand*{\horzbar}{\rule[.5ex]{2.5ex}{0.5pt}}
这里列出 CP 秩与 CP 分解的一些有趣性质。
\begin{enumerate}
\item 
由定理~\ref{thm: matricization-rank}与命题~\ref{prop:copy.tensor}~（第 3 条）可知：若张量 $\At$ 存在秩 $R$ 的 CP 分解，则它的所有矩阵化的秩都不超过 $R$：
\begin{proposition}
    若 $\At = \CP{\Ab_1,\cdots, \Ab_N}$ 是 $\At$ 的一个秩 $R$ 的 CP 分解，则对任意 $I\subset [N]$ 都有 $\rank(\At_{(I)}) \leq R$。
\end{proposition}
下面的张量网络以四阶张量 $\At = \CP{\Ab_1,\Ab_2,\Ab_3, \Ab_N}$ 沿其前两个模的矩阵化为例，说明了这一结论。
\begin{align*}
    \left(
    \begin{tikzpicture}[baseline=-0.5ex]
    \tikzset{tensor/.style = {minimum size = 0.5cm,shape = circle,thick,draw=black,inner sep = 0pt}, edge/.style = {thick,line width=.4mm},every loop/.style={}}
    \def\x{6}
    \def\y{0}   \node[tensor,fill=green!50!red!50!white] (T) at (\x+1,\y) {$\At$};
    \draw[edge] (T) -- (\x+0.3,0) node [midway,above] {\scalebox{0.7}{\dimcolor{$d_1$}}};
    \draw[edge] (T) -- (\x+1.6,0) node [midway,above] {\scalebox{0.7}{\dimcolor{$d_4$}}};;
    \draw[edge] (T) -- (\x+0.4,-0.3) node [midway, below] {\scalebox{0.7}{\dimcolor{$d_2$}}};
    \draw[edge] (T) -- (\x+1.5,-0.4)node [midway, below] {\scalebox{0.7}{\dimcolor{$d_3$}}};
\end{tikzpicture}\right)_{(\{1,2\})}
&=
\left(\begin{tikzpicture}[baseline=\baseline]
    \tikzset{tensor/.style = {minimum size = 0.5cm,shape = circle,thick,draw=black,inner sep = 0pt}, edge/.style   = {thick,line width=.4mm},every loop/.style={}}
    \def\y{0}
    \def\x{0}
  \node[tensor,fill=red!50!white] (A) at (\x,\y){\scalebox{0.85}{$\Ab_1$}};
  \node[tensor,fill=blue!50!white] (B) at (\x+1,\y){$\Ab_2$};
   \node[tensor,fill=red!20!white] (C) at (\x+2,\y){$\Ab_3$};
    \node[tensor,fill=blue!20!white] (D) at (\x+3,\y){$\Ab_4$};
    \draw  node[fill = black,circle,inner sep=0pt,minimum size=4pt] (m) at (\x+1.5,\y+0.7) {};
    \draw[edge] (A) -- (m)node [midway, above] {\scalebox{0.7}{\dimcolor{$R$}}} ;
    \draw[edge] (B) -- (m);
    \draw[edge] (C) -- (m);
    \draw[edge] (D) -- (m)node [midway, above] {\scalebox{0.7}{\dimcolor{$R$}}};
    \draw[edge](A) -- (\x,\y-0.7)node [midway,left] {\scalebox{0.7}{\dimcolor{$d_1$}}};
    \draw[edge](B) -- (\x+1,\y-0.7)node [midway,left] {\scalebox{0.7}{\dimcolor{$d_2$}}};
    \draw[edge](C) -- (\x+2,\y-0.7)node [midway,left] {\scalebox{0.7}{\dimcolor{$d_3$}}};
    \draw[edge](D) -- (\x+3,\y-0.7) node [midway,left] {\scalebox{0.7}{\dimcolor{$d_4$}}};
\end{tikzpicture}\right)_{(\{1,2\})}\\
&=
\begin{tikzpicture}[baseline=\baseline]
    \tikzset{tensor/.style = {minimum size = 0.5cm,shape = circle,thick,draw=black,inner sep = 0pt}, edge/.style   = {thick,line width=.4mm},every loop/.style={}}
    \def\y{0}
    \def\x{0}
  \node[tensor,fill=red!50!white] (A) at (\x,\y){\scalebox{0.85}{$\Ab_1$}};
  \node[tensor,fill=blue!50!white] (B) at (\x+1,\y){$\Ab_2$};
   \node[tensor,fill=red!20!white] (C) at (\x+2,\y){$\Ab_3$};
    \node[tensor,fill=blue!20!white] (D) at (\x+3,\y){$\Ab_4$};
    \draw  node[fill = black,circle,inner sep=0pt,minimum size=4pt] (m1) at (\x+0.5,\y+0.7) {};
    \draw  node[fill = black,circle,inner sep=0pt,minimum size=4pt] (m2) at (\x+2.5,\y+0.7) {};
    \draw[edge](A) -- (\x,\y-0.65) -- (\x+0.5,\y-0.65) -- (\x+0.5,\y-1);
    \draw[edge](B) -- (\x+1,\y-0.65) -- (\x+0.6,\y-0.65) -- (\x+0.6,\y-1);
    \draw[edge](C) -- (\x+2,\y-0.65) -- (\x+2.5,\y-0.65) -- (\x+2.5,\y-1);
    \draw[edge](D) -- (\x+3,\y-0.65) -- (\x+2.6,\y-0.65) -- (\x+2.6,\y-1);
    \node[draw=none] at (\x+0.55,\y-1.15){\scalebox{0.7}{\dimcolor{$d_1d_2$}}};
    \node[draw=none] at (\x+2.55,\y-1.15){\scalebox{0.7}{\dimcolor{$d_3d_4$}}};
    \draw[dotted] (\x+1.5,\y+1.5) -- (\x+1.5,\y-0.7);
    \draw[edge](A) -- (m1)node [midway,left] {\scalebox{0.7}{\dimcolor{$R$}}};;
    \draw[edge](B) -- (m1)node [midway,right] {\scalebox{0.7}{\dimcolor{$R$}}};
    \draw[edge](C) -- (m2)node [midway,left] {\scalebox{0.7}{\dimcolor{$R$}}};
    \draw[edge](D) -- (m2)node [midway,right] {\scalebox{0.7}{\dimcolor{$R$}}};
    \draw[edge](m1)--(m2)node [near start, above] {\scalebox{0.7}{\dimcolor{$R$}}};
\end{tikzpicture}.
\end{align*}
\item  张量 $\At \in \RR^{d_1\times \cdots \times d_N}$ 的 CP 秩有一个容易得到的上界~\citep{kolda2009tensor}
$$\cprank(\At)\leq \min_{n\in [N]} \prod_{i\neq n} d_i.$$
\item 对二阶张量 $\Ab\in\RR^{d_1\times d_2}$，CP 秩恰好对应经典的矩阵秩概念。
\item CP 分解 $\At = \CP{\Ab,\Bb,\Cb}$ 可用克罗内克 $\delta$ 写为
$$\At_{i,j,k} = \sum_{r_1,r_2,r_3}\delta_{r_1,r_2,r_3}\Ab_{i,r_1}\Bb_{j,r_2}\Cb_{k,r_3},$$
也可借助模 $n$ 乘积写成（见~\ref{par:mode-n-tensor-product}）
$$\At = \It\ttm{1}\Ab\ttm{2}\Bb\ttm{3}\Cb$$
其中 $\It$ 为三阶复制张量。
\item 二阶张量（矩阵）的秩在域 $\RR$ 与 $\CC$ 上是相同的。但对高阶张量（阶数 $N\geq3$ 的 $N$ 阶张量），秩会随分解所用的域而改变~\citep{kruskal1989rank}。
\item 对一个 $N$ 阶、每维长度均为 $d$ 的张量（$d_1 = \dots = d_N = d$），其 CP 分解的参数个数为 $NdR$，而完整张量需要 $d^N$ 个参数。若 $R$ 很小，参数个数便可大幅减少。
\end{enumerate}
\end{remark}
下面的命题说明：张量 CP 分解的模 $n$ 矩阵化可以用 Khatri-Rao 积来表达。
\begin{proposition}
设 $\At\in\RR^{d_1\times d_2\times\dots\times d_N}$。若 $\At=\CP{\Ab_1,\Ab_2,\dots,\Ab_N}$，则
$$\At_{(n)}=\Ab_n\left(\Ab_1\krao\dots\krao\Ab_{(n-1)}\krao\Ab_{(n+1)}\krao\dots\krao\Ab_N\right)^\ts.$$
\begin{proof}
取 $N=3$ 且 $n=1$，
\begin{align*}
\At_{(1)}
&=
\left(\begin{tikzpicture}[baseline=\baseline]
    \tikzset{tensor/.style = {minimum size = 0.5cm,shape = circle,thick,draw=black,inner sep = 0pt}, edge/.style  = {thick,line width=.4mm},every loop/.style={}}
    \def\y{0}
    \def\x{0}
    \node[tensor,fill=green!50!red!50!white] (At) at (\x,\y){\scalebox{0.85}{$\At$}};
    \draw[edge] (\x+0.25,\y) -- (\x+0.6,\y)  node [midway,above] {\scalebox{0.7}{\dimcolor{$d_2$}}};
    \draw[edge] (\x-0.25,\y) -- (\x-0.6,\y)  node [midway,above] {\scalebox{0.7}{\dimcolor{$d_1$}}};
    \draw[edge] (\x,\y-0.25) -- (\x,\y-0.7)  node [midway,left] {\scalebox{0.7}{\dimcolor{$d_3$}}};
    \end{tikzpicture}\right)_{(1)}
=
\left(\begin{tikzpicture}[baseline=\baseline]
    \tikzset{tensor/.style = {minimum size = 0.5cm,shape = circle,thick,draw=black,inner sep = 0pt}, edge/.style   = {thick,line width=.4mm},every loop/.style={}}
    \def\y{0}
    \def\x{0}
  \node[tensor,fill=red!50!white] (A1) at (\x,\y){\scalebox{0.85}{$\Ab_1$}};
  \node[tensor,fill=blue!50!white] (A2) at (\x+1,\y){$\Ab_2$};
   \node[tensor,fill=blue!20!white] (A3) at (\x+2,\y){$\Ab_3$};
    \draw  node[fill = black,circle,inner sep=0pt,minimum size=4pt] (m) at (\x+1,\y+1) {};
    \draw[edge] (A1) -- (m) node [midway,left] {\scalebox{0.7}{\dimcolor{$R$}}};
    \draw[edge] (A2) -- (m)node [midway,left=-0.05cm] {\scalebox{0.7}{\dimcolor{$R$}}};
    \draw[edge] (A3) -- (m)node [midway,right] {\scalebox{0.7}{\dimcolor{$R$}}};
    \draw[edge](A1) -- (\x,\y-0.7)node [midway,left] {\edgelabel{$d_1$}};
    \draw[edge](A2) -- (\x+1,\y-0.7)node [midway,left] {\edgelabel{$d_2$}};
    \draw[edge](A3) -- (\x+2,\y-0.7)node [midway,left] {\edgelabel{$d_3$}};
\end{tikzpicture}\right)_{(1)}\\
&=
\begin{tikzpicture}[baseline=\baseline]
    \tikzset{tensor/.style = {minimum size = 0.5cm,shape = circle,thick,draw=black,inner sep = 0pt}, edge/.style   = {thick,line width=.4mm},every loop/.style={}}
    \def\y{0}
    \def\x{0}
    \node[tensor,fill=green!50!red!50!white] (At) at (\x,\y){\scalebox{0.85}{$\At$}};
    \draw[edge] (\x-0.25,\y) -- (\x-0.6,\y)  node [midway,above] {\scalebox{0.7}{\dimcolor{$d_1$}}};
    \draw[edge] (\x+0.2,\y+0.15) -- (\x+0.5,\y+0.15)-- (\x+0.5,\y+0.05) -- (\x+1,\y+0.05) node [midway,above]
    {\scalebox{0.7}{\dimcolor{$d_2$}}};
    \draw[edge] (\x+0.2,\y-0.15) -- (\x+0.5,\y-0.15)-- (\x+0.5,\y-0.05)-- (\x+1,\y-0.05) node [midway,below]
    {\scalebox{0.7}{\dimcolor{$d_3$}}};
    \end{tikzpicture}
=
\begin{tikzpicture}[baseline=\baseline]
    \tikzset{tensor/.style = {minimum size = 0.5cm,shape = circle,thick,draw=black,inner sep = 0pt}, edge/.style   = {thick,line width=.4mm},every loop/.style={}}
    \def\y{0}
    \def\x{0}
    \node[tensor,fill=red!50!white] (A1) at (\x,\y){\scalebox{0.85}{$\Ab_1$}};
    \node[tensor,fill=blue!50!white] (A2) at (\x+1.7,\y+0.7){$\Ab_2$};
    \node[tensor,fill=blue!20!white] (A3) at (\x+1.7,\y-0.7){$\Ab_3$};
    \draw  node[fill = black,circle,inner sep=0pt,minimum size=4pt] (m) at (\x+1,\y) {};
    \draw[edge] (A1) -- (m) node [very near end,above] {\scalebox{0.7}{\dimcolor{$R$}}};
    \draw[edge] (A2) -- (m);
    \draw[edge] (A3) -- (m);
    \draw[edge] (A1) -- (\x-0.7,\y) node [midway,above] {\scalebox{0.7}{\dimcolor{$d_1$}}};
    \draw[edge] (A2) -- (\x+2.5,\y+0.7) -- (\x+2.5,\y+0.05)-- (\x+3,\y+0.05) node [midway,above]
    {\scalebox{0.7}{\dimcolor{$d_2$}}};
    \draw[edge] (A3) -- (\x+2.5,\y-0.7) -- (\x+2.5,\y-0.05) -- (\x+3,\y-0.05) node [midway,below]
    {\scalebox{0.7}{\dimcolor{$d_3$}}};
    \end{tikzpicture}
=
\Ab_1(\Ab_2\krao\Ab_3)^\ts.
\end{align*}
\end{proof}
\end{proposition}


\subsection{Tucker 分解}
Tucker 分解 \cite{tucker1966some} 将一个 $N$ 阶张量
分解为较小的张量与因子矩阵。在这种分解中，较小的张量称为
\emph{核心张量}（core tensor）。Tucker 分解由核心张量与因子矩阵之间的
$N$ 个模 $n$ 乘积构成
（例如参见第~\ref{sec:product}~节）。

设 $\Tt \in \RR^{d_1 \times \dots \times d_N}$。$\Tt$ 的一个 Tucker 分解
由下式给出
\[
    \Tt
    = \Gt \times_1 \Ub_1 \times_2 \Ub_2\times_3 \dots \times_N \Ub_N,
\]
其中 $\Gt \in \RR^{R_1 \times \dots \times R_N}$ 为核心张量，
而 $\Ub_i \in \RR^{d_i \times R_i}$（$i \in [N]$）为因子矩阵。
该 $N$ 元组 $(R_1,\dots,R_N)$ 称为此 Tucker
分解的\emph{秩}。

四阶张量 $\Tt\in\RR^{d_1\times\cdots \times d_4}$ 的 Tucker 分解可用张量网络记法表示如下。
$$
\begin{tikzpicture}[baseline=-0.5ex]
    \tikzset{tensor/.style = {minimum size = 0.5cm,shape = circle,thick,draw=black,inner sep = 0pt}, edge/.style = {thick,line width=.4mm},every loop/.style={}}
    \def\x{6}
    \def\y{0}   \node[tensor,fill=green!50!red!50!white] (T) at (\x+1,\y) {$\Tt$};
    \draw[edge] (T) -- (\x+0.3,0) node [midway,above] {\scalebox{0.7}{\dimcolor{$d_1$}}};
    \draw[edge] (T) -- (\x+1.6,0) node [midway,above] {\scalebox{0.7}{\dimcolor{$d_4$}}};;
    \draw[edge] (T) -- (\x+0.4,-0.3) node [midway, below] {\scalebox{0.7}{\dimcolor{$d_2$}}};
    \draw[edge] (T) -- (\x+1.5,-0.4)node [midway, below] {\scalebox{0.7}{\dimcolor{$d_3$}}};
    \node[draw = none] at (\x+2,0){$=$};
\node[tensor,fill=red!50!white] (G) at (\x+4.5,\y+1) {$\Gt$};
    \node[tensor, path picture={
    \fill[fill=blue!60!white](path picture bounding box.south east)
    -- 
    (path picture bounding box.north west) --
    (path picture bounding box.north east) -- cycle;
    }, shift={(\x+3,\y)}] (U_1) {\scalebox{0.85}{$\Ub_1$}};
    \node[tensor, path picture={
    \fill[fill=blue!40!white](path picture bounding box.south east)
    -- 
    (path picture bounding box.north west) --
    (path picture bounding box.north east) -- cycle;
    }, shift={(\x+4,\y)}] (U_2) {\scalebox{0.85}{$\Ub_2$}};
    \node[tensor, path picture={
    \fill[fill=blue!30!white](path picture bounding box.south east)
    -- 
    (path picture bounding box.north west) --
    (path picture bounding box.north east) -- cycle;
    }, shift={(\x+5,\y)}] (U_3) {\scalebox{0.85}{$\Ub_3$}};
    \node[tensor, path picture={
    \fill[fill=blue!10!white](path picture bounding box.south east)
    -- 
    (path picture bounding box.north west) --
    (path picture bounding box.north east) -- cycle;
    }, shift={(\x+6,\y)}] (U_4) {\scalebox{0.85}{$\Ub_4$}};
    \draw[edge] (G) -- (U_1) node [midway, above] {\scalebox{0.7}{\dimcolor{$R_1$}}};
    \draw[edge] (G) -- (U_2) node [midway, right] {\scalebox{0.7}{\dimcolor{$R_2$}}};
    \draw[edge] (G) -- (U_3) node [midway, right] {\scalebox{0.7}{\dimcolor{$R_3$}}};
    \draw[edge] (G) -- (\x+6,\y+0.25) node [midway, above] {\scalebox{0.7}{\dimcolor{$R_4$}}};
    \draw[edge] (U_1) -- (\x+3,\y-0.7) node [midway, left] {\scalebox{0.7}{\dimcolor{$d_1$}}};
    \draw[edge] (U_2) -- (\x+4,\y-0.7) node [midway, left] {\scalebox{0.7}{\dimcolor{$d_2$}}};
    \draw[edge] (U_3) -- (\x+5,\y-0.7) node [midway, left] {\scalebox{0.7}{\dimcolor{$d_3$}}};
    \draw[edge] (U_4) -- (\x+6,\y-0.7) node [midway, left] {\scalebox{0.7}{\dimcolor{$d_4$}}};
\end{tikzpicture}.
$$


不难证明，因子矩阵
$\Ub_i \in \RR^{d_i \times R_i}$ 总可以取为酉矩阵（见下方注记~\ref{rem:tuckers}
中的第~\ref{item:Tucker_orth}~条）。


\paragraph{多重线性秩（Multi-linear rank）} 张量 $\Tt$ 的\emph{多重线性（或 Tucker）秩}\footnote{注意区分 Tucker \emph{分解}的秩 $(R_1,\dots,R_n)$ 与
张量 $\Tt$ 的（Tucker）多重线性秩：前者由
$\Tt = \Gt \times_1 \Ub_1 \times_2 \dots \times_N \Ub_N$ 中核心张量 $\Gt$ 的尺寸给出，
后者定义为满足条件的最小元组。} 是使得
$\Tt$ 存在秩 $(R_1,\dots,R_N)$ 的分解的所有元组中最小的
$(R_1,\dots,R_N)$。由于整数元组上只有
偏序关系，人们会问这一秩的概念是否良定义。
下面的命题说明它确实是良定义的。

\begin{proposition}
设 $\Tt = \Gt \times_1 \Ub_1 \times_2 \Ub_2\times_3\dots \times_N \Ub_N$ 与
$\Tt = \tilde\Gt \times_1 \tilde\Ub_1 \times_2 \tilde\Ub_2\times_3 \dots \times_N \tilde\Ub_N$ 是同一
张量 $\Tt$ 的两个 Tucker 分解，其秩分别为 $(R_1,R_2,\dots, R_n)$ 与
$(\tilde R_1,\tilde R_2,\dots,\tilde R_n)$。则 $\Tt$ 存在秩为 $(\min(R_1,\tilde R_1), \min(R_2,\tilde R_2), \dots, \min(R_N,\tilde R_N))$ 的 Tucker 分解。
\end{proposition}
\begin{proof}
设 $N=3$。对 $\Tt$ 的两个 Tucker 分解分别应用模 $1$ 展开与定理~\ref{thm: matricization-rank}，可得下列结果：
\begin{align*}
    \rank\left(\begin{tikzpicture}[baseline=\baseline]
    \tikzset{tensor/.style = {minimum size = 0.5cm,shape = circle,thick,draw=black,inner sep = 0pt}, edge/.style  = {thick,line width=.4mm},every loop/.style={}}
    \def\y{0}
    \def\x{0}
    \node[tensor,fill=green!50!red!50!white] (Tt) at (\x,\y){\scalebox{0.85}{$\Tt$}};
    \draw[edge] (\x+0.25,\y) -- (\x+0.6,\y)  node [midway,above] {\scalebox{0.7}{\dimcolor{$d_2$}}};
    \draw[edge] (\x-0.25,\y) -- (\x-0.6,\y)  node [midway,above] {\scalebox{0.7}{\dimcolor{$d_1$}}};
    \draw[edge] (\x,\y-0.25) -- (\x,\y-0.7)  node [midway,left] {\scalebox{0.7}{\dimcolor{$d_3$}}};
    \end{tikzpicture}\right)_{(1)}
=
\rank(\begin{tikzpicture}[baseline=\baseline]
    \tikzset{tensor/.style = {minimum size = 0.5cm,shape = circle,thick,draw=black,inner sep = 0pt}, edge/.style  = {thick,line width=.4mm},every loop/.style={}}
    \def\y{0}
    \def\x{0}
    \node[tensor,fill=green!50!red!50!white] (Tt) at (\x,\y){\scalebox{0.85}{$\Tt$}};
    \draw[edge] (Tt) -- (\x-0.6,\y)  node [midway,above] {\scalebox{0.7}{\dimcolor{$d_1$}}};
    \draw[edge] (Tt) -- (\x+0.6,\y)  node [midway,above] {\scalebox{0.7}{\dimcolor{$d_2$}}};    \draw[edge] (Tt) -- (\x+0.6,\y)  node [midway,above] {\scalebox{0.7}{\dimcolor{$d_2$}}};
    \draw[edge] (\x+0.25,\y-0.1) -- (\x+0.6,\y-0.1)  node [midway,below] {\scalebox{0.7}{\dimcolor{$d_3$}}};
\end{tikzpicture})
=
\rank\left(\begin{tikzpicture}[baseline=\baseline]
    \tikzset{tensor/.style = {minimum size = 0.5cm,shape = circle,thick,draw=black,inner sep = 0pt}, edge/.style   = {thick,line width=.4mm},every loop/.style={}}
    \def\y{0}
    \def\x{0}
    \node[tensor,fill=blue!60!white] (U_1) at (\x,\y){\scalebox{0.85}{$\Ub_1$}};
    \node[tensor,fill=blue!30!white] (U_2) at (\x+2.7,\y+0.7){$\Ub_2$};
    \node[tensor,fill=blue!10!white] (U_3) at (\x+2.7,\y-0.7){$\Ub_3$};
    \node[tensor,fill=red!60!white] (m) at (\x+1,\y) {$\Gt$}; 
     \draw[edge] (\x+1.2,\y+0.15) -- (\x+1.5,\y+0.15)-- (\x+1.5,\y+0.05) -- (\x+2,\y+0.05) node [midway,above]
    {\scalebox{0.7}{\dimcolor{$R_2$}}};
    \draw[edge] (\x+1.2,\y-0.15) -- (\x+1.5,\y-0.15)-- (\x+1.5,\y-0.05)-- (\x+2,\y-0.05) node [midway,below]
    {\scalebox{0.7}{\dimcolor{$R_3$}}};
    \draw[edge] (U_1) -- (m) node [midway,above] {\scalebox{0.7}{\dimcolor{$R_1$}}};
    \draw[edge] (U_1) -- (\x-0.7,\y) node [midway,above] {\scalebox{0.7}{\dimcolor{$d_1$}}};
    \draw[edge](\x+2,\y-0.05) -- (\x+2,\y-0.7) -- (U_3);
    \draw[edge](\x+2,\y+0.05) -- (\x+2,\y+0.7) -- (U_2);
    \draw[edge] (U_2) -- (\x+3.5,\y+0.7) -- (\x+3.5,\y+0.05)-- (\x+4,\y+0.05) node [midway,above]
    {\scalebox{0.7}{\dimcolor{$d_2$}}};
    \draw[edge] (U_3) -- (\x+3.5,\y-0.7) -- (\x+3.5,\y-0.05) -- (\x+4,\y-0.05) node [midway,below]
    {\scalebox{0.7}{\dimcolor{$d_3$}}};
    \draw[dotted] (\x+0.45,\y+1) -- (\x+0.45,\y-1);
    \end{tikzpicture}\right)\leq R_1,
\end{align*}
同时又有
\begin{align*}
    \rank\left(\begin{tikzpicture}[baseline=\baseline]
    \tikzset{tensor/.style = {minimum size = 0.5cm,shape = circle,thick,draw=black,inner sep = 0pt}, edge/.style  = {thick,line width=.4mm},every loop/.style={}}
    \def\y{0}
    \def\x{0}
    \node[tensor,fill=green!50!red!50!white] (Tt) at (\x,\y){\scalebox{0.85}{$\Tt$}};
    \draw[edge] (\x+0.25,\y) -- (\x+0.6,\y)  node [midway,above] {\scalebox{0.7}{\dimcolor{$d_2$}}};
    \draw[edge] (\x-0.25,\y) -- (\x-0.6,\y)  node [midway,above] {\scalebox{0.7}{\dimcolor{$d_1$}}};
    \draw[edge] (\x,\y-0.25) -- (\x,\y-0.7)  node [midway,left] {\scalebox{0.7}{\dimcolor{$d_3$}}};
    \end{tikzpicture}\right)_{(1)}
=
\rank(\begin{tikzpicture}[baseline=\baseline]
    \tikzset{tensor/.style = {minimum size = 0.5cm,shape = circle,thick,draw=black,inner sep = 0pt}, edge/.style  = {thick,line width=.4mm},every loop/.style={}}
    \def\y{0}
    \def\x{0}
    \node[tensor,fill=green!50!red!50!white] (Tt) at (\x,\y){\scalebox{0.85}{$\Tt$}};
    \draw[edge] (Tt) -- (\x-0.6,\y)  node [midway,above] {\scalebox{0.7}{\dimcolor{$d_1$}}};
    \draw[edge] (Tt) -- (\x+0.6,\y)  node [midway,above] {\scalebox{0.7}{\dimcolor{$d_2$}}};    \draw[edge] (Tt) -- (\x+0.6,\y)  node [midway,above] {\scalebox{0.7}{\dimcolor{$d_2$}}};
    \draw[edge] (\x+0.25,\y-0.1) -- (\x+0.6,\y-0.1)  node [midway,below] {\scalebox{0.7}{\dimcolor{$d_3$}}};
\end{tikzpicture})
=
\rank\left(\begin{tikzpicture}[baseline=\baseline]
    \tikzset{tensor/.style = {minimum size = 0.5cm,shape = circle,thick,draw=black,inner sep = 0pt}, edge/.style   = {thick,line width=.4mm},every loop/.style={}}
    \def\y{0}
    \def\x{0}
    \node[tensor,fill=red!60!blue!20!white] (U_1) at (\x,\y){\scalebox{0.85}{$\tilde{\Ub}_1$}};
    \node[tensor,fill=red!30!blue!50!white] (U_2) at (\x+2.7,\y+0.7){$\tilde{\Ub}_2$};
    \node[tensor,fill=red!20!blue!40!white] (U_3) at (\x+2.7,\y-0.7){$\tilde{\Ub}_3$};
    \node[tensor,fill=red!40!white] (m) at (\x+1,\y) {$\tilde{\Gt}$}; 
     \draw[edge] (\x+1.2,\y+0.15) -- (\x+1.5,\y+0.15)-- (\x+1.5,\y+0.05) -- (\x+2,\y+0.05) node [midway,above]
    {\scalebox{0.7}{\dimcolor{$\tilde{R}_2$}}};
    \draw[edge] (\x+1.2,\y-0.15) -- (\x+1.5,\y-0.15)-- (\x+1.5,\y-0.05)-- (\x+2,\y-0.05) node [midway,below]
    {\scalebox{0.7}{\dimcolor{$\tilde{R}_3$}}};
    \draw[edge] (U_1) -- (m) node [midway,above] {\scalebox{0.7}{\dimcolor{$\tilde{R_1}$}}};
    \draw[edge] (U_1) -- (\x-0.7,\y) node [midway,above] {\scalebox{0.7}{\dimcolor{$d_1$}}};
    \draw[edge](\x+2,\y-0.05) -- (\x+2,\y-0.7) -- (U_3);
    \draw[edge](\x+2,\y+0.05) -- (\x+2,\y+0.7) -- (U_2);
    \draw[edge] (U_2) -- (\x+3.5,\y+0.7) -- (\x+3.5,\y+0.05)-- (\x+4,\y+0.05) node [midway,above]
    {\scalebox{0.7}{\dimcolor{$d_2$}}};
    \draw[edge] (U_3) -- (\x+3.5,\y-0.7) -- (\x+3.5,\y-0.05) -- (\x+4,\y-0.05) node [midway,below]
    {\scalebox{0.7}{\dimcolor{$d_3$}}};
    \draw[dotted] (\x+0.45,\y+1) -- (\x+0.45,\y-1);
    \end{tikzpicture}\right)\leq \tilde{R}_1,
\end{align*}
因此，$\rank\left(\begin{tikzpicture}[baseline=\baseline]
    \tikzset{tensor/.style = {minimum size = 0.5cm,shape = circle,thick,draw=black,inner sep = 0pt}, edge/.style  = {thick,line width=.4mm},every loop/.style={}}
    \def\y{0}
    \def\x{0}
    \node[tensor,fill=green!50!red!50!white] (Tt) at (\x,\y){\scalebox{0.85}{$\Tt$}};
    \draw[edge] (\x+0.25,\y) -- (\x+0.6,\y)  node [midway,above] {\scalebox{0.7}{\dimcolor{$d_2$}}};
    \draw[edge] (\x-0.25,\y) -- (\x-0.6,\y)  node [midway,above] {\scalebox{0.7}{\dimcolor{$d_1$}}};
    \draw[edge] (\x,\y-0.25) -- (\x,\y-0.7)  node [midway,left] {\scalebox{0.7}{\dimcolor{$d_3$}}};
\end{tikzpicture}\right)_{(1)}\leq\min(R_1,\tilde{R}_1)$。 
对 $n=2,3$ 同样推理即知，$\Tt$ 存在秩为 $(\min(R_1,\tilde{R}_1),\min(R_2,\tilde{R}_2),\min(R_3,\tilde{R}_3))$ 的 Tucker 分解。
\end{proof}

\paragraph{计算 Tucker 分解。} 与计算起来是 NP 难的 CP 分解不同，求出一个精确的 Tucker 分解（特别是具有最小多重线性秩的那个）可以在多项式时间内通过逐次 SVD 完成。 

Tucker 分解算法的基本思路，是在模 $n$ 上独立于其余各模，求出
$R_n$ 个主导的左奇异向量。该算法称为 \textit{高阶奇异值分解（Higher-Order SVD, HOSVD）}~\citep{de2000multilinear}，其流程画在下面。 为简单起见，我们以 $N=4$ 为例演示这一过程：
\begin{align*}
    &\begin{tikzpicture}[baseline=\baseline]
    \tikzset{tensor/.style = {minimum size = 0.5cm,shape = circle,thick,draw=black,inner sep = 0pt}, edge/.style = {thick,line width=.4mm},every loop/.style={}}
    \def\x{0}
    \def\y{0}
    \node[tensor,fill=green!50!red!50!white] (A) at (\x-5.5,\y){\scalebox{0.85}{$\Tt$}};
    \draw[edge] (A) -- (\x-5.5,\y+0.75) node [midway,left]
    {\scalebox{0.7}{\dimcolor{$d_1$}}};
    \draw[edge] (A) -- (\x-5.5,\y-0.75) node [midway,left]
    {\scalebox{0.7}{\dimcolor{$d_3$}}};
    \draw[edge] (\x-5.75,\y+0.05) -- (\x-6.2,\y-0.3) node [midway,above]
    {\scalebox{0.7}{\dimcolor{$d_2$}}};
    \draw[edge] (\x-5.25,\y+0.05) -- (\x-4.7,\y-0.3) node [midway,above]
    {\scalebox{0.7}{\dimcolor{$d_4$}}};
    \node[draw=none] at (\x-3,\y){$\underrightarrow{\text{mode-1 matricization}}$};
    \node[tensor,fill=green!50!red!50!white] (A) at (\x-0.5,\y){\scalebox{0.85}{$\Tb$}};
    \draw[edge] (\x-0.35,\y+0.2) -- (\x,\y+0.2) --(\x,\y+0.1) -- (\x
    +0.35,\y+0.1) node [midway,above]
    {\scalebox{0.7}{\dimcolor{$d_2$}}};
    \draw[edge] (\x-0.35,\y-0.2) -- (\x,\y-0.2) -- (\x,\y-0.1) -- (\x+0.35,\y-0.1) node [midway,below]
    {\scalebox{0.7}{\dimcolor{$d_4$}}};
     \draw[edge] (A) -- (\x+0.35,\y)node [right]
    {\scalebox{0.7}{\dimcolor{$d_3$}}};
     \draw[edge] (A) -- (\x-1.25,\y) node [midway,above]
    {\scalebox{0.7}{\dimcolor{$d_1$}}};
    \node[draw=none] at (\x+1,\y){$\rightarrow$};
    \node[tensor, path picture={
    \fill[fill=blue!50!white](path picture bounding box.south east)
        -- 
    (path picture bounding box.north west) --
    (path picture bounding box.north east) -- cycle;
    }, shift={(\x+1.5,\y)}] (U_1) {\scalebox{0.8}{$\Ub_1$}};
    \node[tensor,fill=red!50!white] (Sigma_1) at (\x+2.5,\y){\scalebox{0.85}{$\Sigmab_1$}};
    \node[tensor, path picture={
    \fill[fill=red!60!white](path picture bounding box.south west)
        -- 
    (path picture bounding box.north east) --
    (path picture bounding box.north west) -- cycle;
    }, shift={(\x+3.5,\y)}] (V_1) {\scalebox{0.8}{$\Vb_1$}};
    \draw[edge](V_1) -- (\x+4.5,\y) node [right]{\scalebox{0.7}{\textcolor{gray}{$d_3$}}};
    \draw[edge] (U_1) -- (Sigma_1)node [midway,above]
    {\scalebox{0.7}{\dimcolor{$R_1$}}};
    \draw[edge] (U_1) -- (\x+1.5,\y-0.7)node [midway,left]
    {\scalebox{0.7}{\dimcolor{$d_1$}}};
    \draw[edge] (Sigma_1) -- (V_1) node[midway,above]    {\scalebox{0.7}{\dimcolor{$R_1$}}};
    \draw[edge] (\x+3.7,\y+0.15) -- (\x+4,\y+0.15) --(\x+4,\y+0.1) -- (\x+4.5,\y+0.1) node [midway,above]
    {\scalebox{0.7}{\dimcolor{$d_2$}}};
    \draw[edge] (\x+3.65,\y-0.2) -- (\x+4,\y-0.2) -- (\x+4,\y-0.1) -- (\x+4.5,\y-0.1) node [midway,below]
    {\scalebox{0.7}{\dimcolor{$d_4$}}};
    \node[draw=none] at (\x+5.75,\y){\text{Keep $\Ub_1$}.};
\end{tikzpicture}\\
&\begin{tikzpicture}[baseline=\baseline]
    \tikzset{tensor/.style = {minimum size = 0.5cm,shape = circle,thick,draw=black,inner sep = 0pt}, edge/.style = {thick,line width=.4mm},every loop/.style={}}
    \def\y{0}
    \def\x{0}
    \node[tensor,fill=green!50!red!50!white] (A) at (\x-2.5,\y){\scalebox{0.85}{$\Tt$}};
    \draw[edge] (A) -- (\x-2.5,\y+0.75) node [midway,left]
    {\scalebox{0.7}{\dimcolor{$d_1$}}};
    \draw[edge] (A) -- (\x-2.5,\y-0.75) node [midway,left]
    {\scalebox{0.7}{\dimcolor{$d_3$}}};
    \draw[edge] (\x-2.75,\y+0.05) -- (\x-3.2,\y-0.3) node [midway,above]
    {\scalebox{0.7}{\dimcolor{$d_2$}}};
    \draw[edge] (\x-2.25,\y+0.05) -- (\x-1.7,\y-0.3) node [midway,above]
    {\scalebox{0.7}{\dimcolor{$d_4$}}};
    \node[draw=none] at (\x,\y){$\underrightarrow{\text{mode-2 matricization}}$};
    \node[tensor,fill=green!50!red!50!white] (A) at (\x+2.5,\y){\scalebox{0.85}{$\Tb$}};
    \draw[edge] (A) -- (\x+1.75,\y) node [midway,above]
    {\scalebox{0.7}{\dimcolor{$d_2$}}};
    \draw[edge] (A) -- (\x+3.35,\y)node [right]
    {\scalebox{0.7}{\dimcolor{$d_3$}}};
    \draw[edge] (\x+2.65,\y+0.2) -- (\x+3,\y+0.2) -- (\x+3,\y+0.1) -- (\x+3.35,\y+0.1)node [midway,above]
    {\scalebox{0.7}{\dimcolor{$d_1$}}};
    \draw[edge] (\x+2.65,\y-0.2) -- (\x+3,\y-0.2) -- (\x+3,\y-0.1) -- (\x+3.35,\y-0.1)node [midway,below]
    {\scalebox{0.7}{\dimcolor{$d_4$}}};
    \node[draw=none] at (\x+4,\y){$\rightarrow$};
    \node[tensor, path picture={
    \fill[fill=blue!40!white](path picture bounding box.south east)
        -- 
    (path picture bounding box.north west) --
    (path picture bounding box.north east) -- cycle;
    }, shift={(\x+4.5,\y)}] (U_2) {\scalebox{0.8}{$\Ub_2$}};
    \node[tensor,fill=red!40!white] (Sigma_2) at (\x+5.5,\y){\scalebox{0.85}{$\Sigmab_2$}};
    \node[tensor, path picture={
    \fill[fill=red!50!white](path picture bounding box.south west)
        -- 
    (path picture bounding box.north east) --
    (path picture bounding box.north west) -- cycle;
    }, shift={(\x+6.5,\y)}] (V_2) {\scalebox{0.8}{$\Vb_2$}};
    \draw[edge](V_2) -- (\x+7.5,\y) node [right]{\scalebox{0.7}{\dimcolor{$d_3$}}};
    \draw[edge] (U_2) -- (Sigma_2)node [midway,above]
    {\scalebox{0.7}{\dimcolor{$R_2$}}};
    \draw[edge] (U_2) -- (\x+4.5,\y-0.7)node [midway,left]
    {\scalebox{0.7}{\dimcolor{$d_2$}}};
    \draw[edge] (Sigma_2) -- (V_2) node[midway,above] {\scalebox{0.7}{\dimcolor{$R_2$}}};
    \draw[edge] (\x+6.71,\y+0.15) -- (\x+7,\y+0.15) --(\x+7,\y+0.1) -- (\x+7.5,\y+0.1) node [midway,above]
    {\scalebox{0.7}{\dimcolor{$d_1$}}};
    \draw[edge] (\x+6.65,\y-0.2) -- (\x+7,\y-0.2) -- (\x+7,\y-0.1) -- (\x+7.5,\y-0.1) node [midway,below]
    {\scalebox{0.7}{\dimcolor{$d_4$}}};
    \node[draw=none] at (\x+8.75,\y){\text{Keep $\Ub_2$}.};
\end{tikzpicture}\\
&\begin{tikzpicture}[baseline=\baseline]
    \tikzset{tensor/.style = {minimum size = 0.5cm,shape = circle,thick,draw=black,inner sep = 0pt}, edge/.style = {thick,line width=.4mm},every loop/.style={}}
    \def\y{0}
    \def\x{0}
    \node[tensor,fill=green!50!red!50!white] (A) at (\x-2.5,\y){\scalebox{0.85}{$\Tt$}};
    \draw[edge] (A) -- (\x-2.5,\y+0.75) node [midway,left]
    {\scalebox{0.7}{\dimcolor{$d_1$}}};
    \draw[edge] (A) -- (\x-2.5,\y-0.75) node [midway,left]
    {\scalebox{0.7}{\dimcolor{$d_3$}}};
    \draw[edge] (\x-2.75,\y+0.05) -- (\x-3.2,\y-0.3) node [midway,above]
    {\scalebox{0.7}{\dimcolor{$d_2$}}};
    \draw[edge] (\x-2.25,\y+0.05) -- (\x-1.7,\y-0.3) node [midway,above]
    {\scalebox{0.7}{\dimcolor{$d_4$}}};
    \node[draw=none] at (\x,\y){$\underrightarrow{\text{mode-3 matricization}}$};
    \node[tensor,fill=green!50!red!50!white] (A) at (\x+2.5,\y){\scalebox{0.85}{$\Tb$}};
    \draw[edge] (A) -- (\x+1.75,\y) node [midway,above]
    {\scalebox{0.7}{\dimcolor{$d_3$}}};
    \draw[edge] (A) -- (\x+3.35,\y)node [right]
    {\scalebox{0.7}{\dimcolor{$d_2$}}};
    \draw[edge] (\x+2.65,\y+0.2) -- (\x+3,\y+0.2) -- (\x+3,\y+0.1) -- (\x+3.35,\y+0.1)node [midway,above]
    {\scalebox{0.7}{\dimcolor{$d_1$}}};
    \draw[edge] (\x+2.65,\y-0.2) -- (\x+3,\y-0.2) -- (\x+3,\y-0.1) -- (\x+3.35,\y-0.1)node [midway,below]
    {\scalebox{0.7}{\dimcolor{$d_4$}}};
    \node[draw=none] at (\x+4,\y){$\rightarrow$};
    \node[tensor, path picture={
    \fill[fill=blue!30!white](path picture bounding box.south east)
        -- 
    (path picture bounding box.north west) --
    (path picture bounding box.north east) -- cycle;
    }, shift={(\x+4.5,\y)}] (U_3) {\scalebox{0.8}{$\Ub_3$}};
    \node[tensor,fill=red!30!white] (Sigma_3) at (\x+5.5,\y){\scalebox{0.85}{$\Sigmab_3$}};
    \node[tensor, path picture={
    \fill[fill=red!40!white](path picture bounding box.south west)
        -- 
    (path picture bounding box.north east) --
    (path picture bounding box.north west) -- cycle;
    }, shift={(\x+6.5,\y)}] (V_3) {\scalebox{0.8}{$\Vb_3$}};
    \draw[edge](V_3) -- (\x+7.5,\y) node [right]{\scalebox{0.7}{\dimcolor{$d_2$}}};
    \draw[edge] (U_3) -- (Sigma_3)node [midway,above]
    {\scalebox{0.7}{\dimcolor{$R_3$}}};
    \draw[edge] (U_3) -- (\x+4.5,\y-0.7)node [midway,left]
    {\scalebox{0.7}{\dimcolor{$d_3$}}};
    \draw[edge] (Sigma_3) -- (V_3) node[midway,above] {\scalebox{0.7}{\dimcolor{$R_3$}}};
    \draw[edge] (\x+6.71,\y+0.15) -- (\x+7,\y+0.15) --(\x+7,\y+0.1) -- (\x+7.5,\y+0.1) node [midway,above]
    {\scalebox{0.7}{\dimcolor{$d_1$}}};
    \draw[edge] (\x+6.65,\y-0.2) -- (\x+7,\y-0.2) -- (\x+7,\y-0.1) -- (\x+7.5,\y-0.1) node [midway,below]
    {\scalebox{0.7}{\dimcolor{$d_4$}}};
    \node[draw=none] at (\x+8.75,\y){\text{Keep $\Ub_3$.}};
\end{tikzpicture}\\
&\begin{tikzpicture}[baseline=\baseline]
    \tikzset{tensor/.style = {minimum size = 0.5cm,shape = circle,thick,draw=black,inner sep = 0pt}, edge/.style = {thick,line width=.4mm},every loop/.style={}}
    \def\y{0}
    \def\x{0}
    \node[tensor,fill=green!50!red!50!white] (A) at (\x-2.5,\y){\scalebox{0.85}{$\Tt$}};
    \draw[edge] (A) -- (\x-2.5,\y+0.75) node [midway,left]
    {\scalebox{0.7}{\dimcolor{$d_1$}}};
    \draw[edge] (A) -- (\x-2.5,\y-0.75) node [midway,left]
    {\scalebox{0.7}{\dimcolor{$d_3$}}};
    \draw[edge] (\x-2.75,\y+0.05) -- (\x-3.2,\y-0.3) node [midway,above]
    {\scalebox{0.7}{\dimcolor{$d_2$}}};
    \draw[edge] (\x-2.25,\y+0.05) -- (\x-1.7,\y-0.3) node [midway,above]
    {\scalebox{0.7}{\dimcolor{$d_4$}}};
    \node[draw=none] at (\x,\y){$\underrightarrow{\text{mode-4 matricization}}$};
    \node[tensor,fill=green!50!red!50!white] (A) at (\x+2.5,\y){\scalebox{0.85}{$\Tb$}};
    \draw[edge] (A) -- (\x+1.75,\y) node [midway,above]
    {\scalebox{0.7}{\dimcolor{$d_4$}}};
    \draw[edge] (A) -- (\x+3.35,\y)node [right]
    {\scalebox{0.7}{\dimcolor{$d_2$}}};
    \draw[edge] (\x+2.65,\y+0.2) -- (\x+3,\y+0.2) -- (\x+3,\y+0.1) -- (\x+3.35,\y+0.1)node [midway,above]
    {\scalebox{0.7}{\dimcolor{$d_1$}}};
    \draw[edge] (\x+2.65,\y-0.2) -- (\x+3,\y-0.2) -- (\x+3,\y-0.1) -- (\x+3.35,\y-0.1)node [midway,below]
    {\scalebox{0.7}{\dimcolor{$d_3$}}};
    \node[draw=none] at (\x+4,\y){$\rightarrow$};
    \node[tensor, path picture={
    \fill[fill=blue!20!white](path picture bounding box.south east)
        -- 
    (path picture bounding box.north west) --
    (path picture bounding box.north east) -- cycle;
    }, shift={(\x+4.5,\y)}] (U_4) {\scalebox{0.8}{$\Ub_4$}};
    \node[tensor,fill=red!20!white] (Sigma_4) at (\x+5.5,\y){\scalebox{0.85}{$\Sigmab_4$}};
    \node[tensor, path picture={
    \fill[fill=red!30!white](path picture bounding box.south west)
        -- 
    (path picture bounding box.north east) --
    (path picture bounding box.north west) -- cycle;
    }, shift={(\x+6.5,\y)}] (V_4) {\scalebox{0.8}{$\Vb_4$}};
    \draw[edge](V_4) -- (\x+7.5,\y) node [right]{\scalebox{0.7}{\dimcolor{$d_2$}}};
    \draw[edge] (U_4) -- (Sigma_4)node [midway,above]
    {\scalebox{0.7}{\dimcolor{$R_4$}}};
    \draw[edge] (U_4) -- (\x+4.5,\y-0.7)node [midway,left]
    {\scalebox{0.7}{\dimcolor{$d_4$}}};
    \draw[edge] (Sigma_4) -- (V_4) node[midway,above] {\scalebox{0.7}{\dimcolor{$R_4$}}};
    \draw[edge] (\x+6.71,\y+0.15) -- (\x+7,\y+0.15) --(\x+7,\y+0.1) -- (\x+7.5,\y+0.1) node [midway,above]
    {\scalebox{0.7}{\dimcolor{$d_1$}}};
    \draw[edge] (\x+6.65,\y-0.2) -- (\x+7,\y-0.2) -- (\x+7,\y-0.1) -- (\x+7.5,\y-0.1) node [midway,below]
    {\scalebox{0.7}{\dimcolor{$d_3$}}};
    \node[draw=none] at (\x+8.75,\y){\text{Keep $\Ub_4$.}};
\end{tikzpicture}
\end{align*}
若上述所有 SVD 都是精确的，即对每个 $n$ 有 $R_n = \rank(\Tt_{(n)})$，则对所有 $n$ 均有 $\Tt \times_n \Ub_n\Ub_n^\ts = \Tt$，从而 $\Tt = \Tt\times_1\Ub_1\Ub_1^\ts\times_2\cdots\times_N\Ub_N\Ub_N^\ts$。举例来说，当 $N=4$ 时有：
$$
\begin{tikzpicture}[baseline=-0.5ex]
    \tikzset{tensor/.style = {minimum size = 0.5cm,shape = circle,thick,draw=black,inner sep = 0pt}, edge/.style = {thick,line width=.4mm},every loop/.style={}}
    \def\x{6}
    \def\y{0}   \node[tensor,fill=green!50!red!50!white] (T) at (\x+1,\y) {$\Tt$};
    \draw[edge] (T) -- (\x+0.3,0) node [midway,above] {\scalebox{0.7}{\dimcolor{$d_1$}}};
    \draw[edge] (T) -- (\x+1.6,0) node [midway,above] {\scalebox{0.7}{\dimcolor{$d_4$}}};;
    \draw[edge] (T) -- (\x+0.4,-0.3) node [midway, below] {\scalebox{0.7}{\dimcolor{$d_2$}}};
    \draw[edge] (T) -- (\x+1.5,-0.4)node [midway, below] {\scalebox{0.7}{\dimcolor{$d_3$}}};
    \node[draw = none] at (\x+2.2,0){$=$};
\end{tikzpicture}
\begin{tikzpicture}[baseline=-0.5ex]
    \tikzset{tensor/.style = {minimum size = 0.5cm,shape = circle,thick,draw=black,inner sep = 0pt}, edge/.style = {thick,line width=.4mm},every loop/.style={}}
    \def\x{6}
    \def\y{0} 
    \node[tensor,fill=green!50!red!50!white] (T) at (\x+4.5,\y+1) {$\Tt$};
    \node[tensor, path picture={
    \fill[fill=blue!60!white](path picture bounding box.north east)
        -- 
    (path picture bounding box.south east) --
    (path picture bounding box.south west) -- cycle;
    }, shift={(\x+3,\y)}] (U_1) {\scalebox{0.85}{$\Ub_1$}};
    \node[tensor, path picture={
    \fill[fill=blue!60!white](path picture bounding box.south east)
        -- 
    (path picture bounding box.north east) --
    (path picture bounding box.north west) -- cycle;
    }, shift={(\x+3,\y-0.95)}] (Ut_1) {\scalebox{0.85}{$\Ub_1$}};
    \node[tensor, path picture={
    \fill[fill=blue!40!white](path picture bounding box.north east)
        -- 
    (path picture bounding box.south east) --
    (path picture bounding box.south west) -- cycle;
    }, shift={(\x+4,\y)}] (U_2) {\scalebox{0.85}{$\Ub_2$}};
    \node[tensor, path picture={
    \fill[fill=blue!40!white](path picture bounding box.south east)
        -- 
    (path picture bounding box.north east) --
    (path picture bounding box.north west) -- cycle;
    }, shift={(\x+4,\y-0.95)}] (Ut_2) {\scalebox{0.85}{$\Ub_2$}};
    \node[tensor, path picture={
    \fill[fill=blue!30!white](path picture bounding box.north east)
    -- 
    (path picture bounding box.south west) --
    (path picture bounding box.south east) -- cycle;
    }, shift={(\x+5,\y)}] (U_3) {\scalebox{0.85}{$\Ub_3$}};
    \node[tensor, path picture={
    \fill[fill=blue!30!white](path picture bounding box.south east)
        -- 
    (path picture bounding box.north east) --
    (path picture bounding box.north west) -- cycle;
    }, shift={(\x+5,\y-0.95)}] (Ut_3) {\scalebox{0.85}{$\Ub_3$}};
    \node[tensor, path picture={
    \fill[fill=blue!10!white](path picture bounding box.north east)
    -- 
    (path picture bounding box.south west) --
    (path picture bounding box.south east) -- cycle;
    }, shift={(\x+6,\y)}] (U_4) {\scalebox{0.85}{$\Ub_4$}};
    \node[tensor, path picture={
    \fill[fill=blue!10!white](path picture bounding box.south east)
        -- 
    (path picture bounding box.north east) --
    (path picture bounding box.north west) -- cycle;
    }, shift={(\x+6,\y-0.95)}] (Ut_4) {\scalebox{0.85}{$\Ub_4$}};
    \draw[edge] (T) -- (\x+2.95,\y+0.25) node [midway, above] {\scalebox{0.7}{\dimcolor{$d_1$}}};
    \draw[edge] (T) -- (U_2) node [midway, right] {\scalebox{0.7}{\dimcolor{$d_2$}}};
    \draw[edge] (T) -- (U_3) node [midway, right] {\scalebox{0.7}{\dimcolor{$d_3$}}};
    \draw[edge] (T) -- (\x+6,\y+0.25) node [midway, above] {\scalebox{0.7}{\dimcolor{$d_4$}}};
    \draw[edge] (U_1) -- (\x+3,\y-0.7) node [midway, left] {\scalebox{0.7}{\dimcolor{$R_1$}}};
    \draw[edge] (U_2) -- (\x+4,\y-0.7) node [midway, left] {\scalebox{0.7}{\dimcolor{$R_2$}}};
    \draw[edge] (U_3) -- (\x+5,\y-0.7) node [midway, left] {\scalebox{0.7}{\dimcolor{$R_3$}}};
    \draw[edge] (U_4) -- (\x+6,\y-0.7) node [midway, left] {\scalebox{0.7}{\dimcolor{$R_4$}}};
    \draw[edge] (Ut_1) -- (\x+3,\y-1.7) node [midway, left] {\scalebox{0.7}{\dimcolor{$d_1$}}};
    \draw[edge] (Ut_2) -- (\x+4,\y-1.7) node [midway, left] {\scalebox{0.7}{\dimcolor{$d_2$}}};
    \draw[edge] (Ut_3) -- (\x+5,\y-1.7) node [midway, left] {\scalebox{0.7}{\dimcolor{$d_3$}}};
    \draw[edge] (Ut_4) -- (\x+6,\y-1.7) node [midway, left] {\scalebox{0.7}{\dimcolor{$d_4$}}};
    \draw[dash pattern=on 1pt off 2pt] plot[smooth cycle] coordinates { 
    (\x+2.5,\y) 
    (\x+2.5, \y+1.2) (\x+6.3,\y+1.2) (\x+6.3, \y-0.55)  (\x+2.5,\y-0.55)};
    \node[draw = none] at (\x+7,0){$=$};
\end{tikzpicture}
\begin{tikzpicture}[baseline=-0.5ex]
    \tikzset{tensor/.style = {minimum size = 0.5cm,shape = circle,thick,draw=black,inner sep = 0pt}, edge/.style = {thick,line width=.4mm},every loop/.style={}}
    \def\x{6}
    \def\y{0} 
    \node[tensor,fill=red!50!white] (G) at (\x+13.5,\y+1) {$\Gt$};
    \node[tensor, path picture={
    \fill[fill=blue!60!white](path picture bounding box.south east)
    -- 
    (path picture bounding box.north west) --
    (path picture bounding box.north east) -- cycle;
    }, shift={(\x+12.2,\y)}] (U_1)  {\scalebox{0.85}{$\Ub_1$}};  \node[tensor, path picture={
    \fill[fill=blue!40!white](path picture bounding box.south east)
    -- 
    (path picture bounding box.north west) --
    (path picture bounding box.north east) -- cycle;
    }, shift={(\x+13.2,\y)}] (U_2) {\scalebox{0.85}{$\Ub_2$}};
    \node[tensor, path picture={
    \fill[fill=blue!30!white](path picture bounding box.south east)
    -- 
    (path picture bounding box.north west) --
    (path picture bounding box.north east) -- cycle;
    }, shift={(\x+14.2,\y)}] (U_3) {\scalebox{0.85}{$\Ub_3$}};
        \node[tensor, path picture={
    \fill[fill=blue!10!white](path picture bounding box.south east)
    -- 
    (path picture bounding box.north west) --
    (path picture bounding box.north east) -- cycle;
    }, shift={(\x+15.2,\y)}] (U_4) {\scalebox{0.85}{$\Ub_4$}};
     \draw[edge] (U_1) -- (\x+12.2,\y-0.7) node [midway, left] {\scalebox{0.7}{\dimcolor{$d_1$}}};
    \draw[edge] (U_2) -- (\x+13.2,\y-0.7) node [midway, left] {\scalebox{0.7}{\dimcolor{$d_2$}}};
    \draw[edge] (U_3) -- (\x+14.2,\y-0.7) node [midway, left] {\scalebox{0.7}{\dimcolor{$d_3$}}};
    \draw[edge] (U_4) -- (\x+15.2,\y-0.7) node [midway, left] {\scalebox{0.7}{\dimcolor{$d_4$}}};
    \draw[edge](G) -- (U_1)node[midway,above]{\scalebox{0.7}{\dimcolor{$R_1$}}};
    \draw[edge](G) -- (U_2)node[midway,right]{\scalebox{0.7}{\dimcolor{$R_2$}}};
    \draw[edge](G) -- (U_3)node[midway,right]{\scalebox{0.7}{\dimcolor{$R_3$}}};
    \draw[edge](G) -- (\x+15.2,\y+0.25)node[midway,above]{\scalebox{0.7}{\dimcolor{$R_4$}}};
\end{tikzpicture}.
$$
其中张量 $\Gt = \Tt\ttm{1}\Ub_1^\ts\ttm{2}\dots\ttm{4}\Ub_4^\ts$。 

核心张量  $\Gt\in\RR^{R_1\times\dots\times R_4}$ 与因子矩阵 $\Ub_1,\dots,\Ub_4$ 一起构成张量 $\Tt$ 的一个 Tucker 分解。若 $\Tt$ 各矩阵化上的 SVD 全部精确，则所得的 Tucker 分解也是精确的。使用截断 SVD 时，Tucker 分解的整体逼近误差可表示为各截断 SVD 所引入误差的函数~(见~\citep{de2000multilinear})。

\paragraph{多重线性秩与矩阵化的秩}
上述算法表明：若对所有 $n$ 有 $\tilde R_n = \rank(\Tt_{(n)})$，则存在秩为 $(\tilde R_1,\dots, \tilde R_N)$ 的 Tucker 分解，这说明 $\Tt$ 的多重线性秩 $( R_1,\dots, R_N)$ 以矩阵化的秩为上界，即对所有 $n$ 有 $ R_n \leq \rank(\Tt_{(n)} )$。也容易证明，多重线性秩以矩阵化的秩为下界。事实上，设 $\T$ 的多重线性秩为 $(R_1,\dots, R_N)$，则存在秩 $(R_1,\dots, R_N)$ 的 Tucker 分解 $\Tt = \Gt\times_1\Ub_1\times_2 \dots \times_N \Ub_N$。由定理~\ref{thm: matricization-rank} 即可看出，这意味着对所有 $n$ 有 $\rank(\Tt_{(n)} )\leq  R_n$。
这就证明了如下定理。
\begin{theorem}
张量 $\Tt$ 的多重线性秩 $(R_1,R_2,\dots, R_N)$ 由其矩阵化的秩给出，即对 $n\in[N]$ 有 $R_n = \rank(\Tt_{(n)})$。
\end{theorem}


\begin{remark}\label{rem:tuckers}
这里列出 Tucker 分解以及用于计算它的高阶奇异值分解~(HOSVD) 算法的一些有趣性质。
\begin{enumerate}

\item  若 $\Tt = \Gt\ttm{1}\Ub_1\ttm{2}\dots\ttm{N-1}\Ub_{N-1}\ttm{N}\Ub_N$ 且所有 $\Ub_i$ 均正交，则 $\norm{\Tt}_F = \norm{\Gt}_F$，且 $\vectorize(\Gt) = (\Ub_1\kron\cdots\kron\Ub_N)\vectorize(\Tt)$。
事实上，当 $N=3$ 时有
\begin{align*}
(\Ub_1\kron\Ub_2\kron\Ub_3)\vectorize(\Tt) &= 
\begin{tikzpicture}[baseline=-0.5ex]
    \tikzset{tensor/.style = {minimum size = 0.5cm,shape = circle,thick,draw=black,inner sep = 0pt}, edge/.style = {thick,line width=.4mm},every loop/.style={}}
    \def\x{0}
    \def\y{0.5}   
    \node[tensor, path picture={
    \fill[fill=blue!60!white](path picture bounding box.south east)
        -- 
    (path picture bounding box.north west) --
    (path picture bounding box.north east) -- cycle;
    }, shift={(\x-0.85,\y-0.65)}] (U_1) {\scalebox{0.85}{$\Ub_1$}};
    \node[tensor, path picture={
    \fill[fill=blue!40!white](path picture bounding box.south east)
        -- 
    (path picture bounding box.north west) --
    (path picture bounding box.north east) -- cycle;
    }, shift={(\x,\y-0.65)}] (U_2) {\scalebox{0.85}{$\Ub_2$}};
    \node[tensor, path picture={
    \fill[fill=blue!30!white](path picture bounding box.south east)
        -- 
    (path picture bounding box.north west) --
    (path picture bounding box.north east) -- cycle;
    }, shift={(\x+1,\y-0.65)}] (U_3) {\scalebox{0.85}{$\Ub_3$}};
    \draw[edge] (U_1) -- (\x-0.85,\y-1.35) -- (\x-0.1,\y-1.35) -- (\x-0.1,\y-1.85) node [midway, left] {\scalebox{0.7}{\dimcolor{$d_1$}}};
    \draw[edge] (U_2) -- (\x,\y+0.55) node [very near start, right=-0.1cm] {\scalebox{0.7}{\dimcolor{$R_2$}}};
    \draw[edge] (U_3) -- (\x+1,\y-1.35) -- (\x+0.1,\y-1.35) -- (\x+0.1, \y-1.85) node [midway, right] {\scalebox{0.7}{\dimcolor{$d_3$}}};
    \draw[edge](U_1) -- (\x-0.85,\y) -- (\x-0.1,\y) -- (\x-0.1,\y+0.55)node [midway, left] {\scalebox{0.7}{\dimcolor{$R_1$}}};
    \draw[edge](U_3) -- (\x+1,\y) -- (\x+0.1,\y) -- (\x+0.1,\y+0.55)node [midway, right] {\scalebox{0.7}{\dimcolor{$R_3$}}};
    \node[tensor,fill=green!50!red!50!white](T) at (\x,\y-1.85) {$\Tt$};
    \draw[edge](U_2)--(T)node [near start, right=-0.1cm] {\scalebox{0.7}{\dimcolor{$d_2$}}};
\end{tikzpicture}
=
\begin{tikzpicture}[baseline=-0.5ex]
    \tikzset{tensor/.style = {minimum size = 0.5cm,shape = circle,thick,draw=black,inner sep = 0pt}, edge/.style = {thick,line width=.4mm},every loop/.style={}}
    \def\x{0}
    \def\y{1.25}   
    \node[tensor, path picture={
    \fill[fill=blue!60!white](path picture bounding box.south east)
        -- 
    (path picture bounding box.north west) --
    (path picture bounding box.north east) -- cycle;
    }, shift={(\x-0.85,\y-0.65)}] (U_1) {\scalebox{0.85}{$\Ub_1$}};
    \node[tensor, path picture={
    \fill[fill=blue!60!white](path picture bounding box.north east)
        -- 
    (path picture bounding box.south west) --
    (path picture bounding box.south east) -- cycle;
    }, shift={(\x-0.85,\y-1.5)}] (Ut_1) {\scalebox{0.85}{$\Ub_1$}};
    \node[tensor, path picture={
    \fill[fill=blue!40!white](path picture bounding box.south east)
        -- 
    (path picture bounding box.north west) --
    (path picture bounding box.north east) -- cycle;
    }, shift={(\x,\y-0.65)}] (U_2) {\scalebox{0.85}{$\Ub_2$}};
     \node[tensor, path picture={
    \fill[fill=blue!40!white](path picture bounding box.north east)
        -- 
    (path picture bounding box.south west) --
    (path picture bounding box.south east) -- cycle;
    }, shift={(\x,\y-1.5)}] (Ut_2) {\scalebox{0.85}{$\Ub_2$}};
    \node[tensor, path picture={
    \fill[fill=blue!30!white](path picture bounding box.south east)
        -- 
    (path picture bounding box.north west) --
    (path picture bounding box.north east) -- cycle;
    }, shift={(\x+1,\y-0.65)}] (U_3) {\scalebox{0.85}{$\Ub_3$}};
    \node[tensor, path picture={
    \fill[fill=blue!30!white](path picture bounding box.north east)
        -- 
    (path picture bounding box.south west) --
    (path picture bounding box.south east) -- cycle;
    }, shift={(\x+1,\y-1.5)}] (Ut_3) {\scalebox{0.85}{$\Ub_3$}};
    \draw[edge] (U_2) -- (\x,\y+0.55) node [very near start, right=-0.1cm] {\scalebox{0.7}{\dimcolor{$R_2$}}};
    \draw[edge](U_1) -- (\x-0.85,\y) -- (\x-0.1,\y) -- (\x-0.1,\y+0.55)node [midway, left] {\scalebox{0.7}{\dimcolor{$R_1$}}};
    \draw[edge](U_3) -- (\x+1,\y) -- (\x+0.1,\y) -- (\x+0.1,\y+0.55)node [midway, right] {\scalebox{0.7}{\dimcolor{$R_3$}}};
    \draw[edge](U_1)--(Ut_1)node [midway, right] {\scalebox{0.7}{\dimcolor{$d_1$}}};
    \draw[edge](U_2)--(Ut_2)node [midway, right] {\scalebox{0.7}{\dimcolor{$d_2$}}};
    \draw[edge](U_3)--(Ut_3)node [midway, right] {\scalebox{0.7}{\dimcolor{$d_3$}}};
    \node[tensor,fill=red!50!white] (G) at (\x,\y-2.5) {$\Gt$};
    \draw[edge](G)--(Ut_1)node [midway, left] {\scalebox{0.7}{\dimcolor{$R_1$}}};
    \draw[edge](G)--(Ut_2)node [midway, right=-0.1cm] {\scalebox{0.7}{\dimcolor{$R_2$}}};
    \draw[edge](G)--(Ut_3)node [midway, right] {\scalebox{0.7}{\dimcolor{$R_3$}}};
\end{tikzpicture}
=
\begin{tikzpicture}[baseline=-0.5ex]
    \tikzset{tensor/.style = {minimum size = 0.5cm,shape = circle,thick,draw=black,inner sep = 0pt}, edge/.style = {thick,line width=.4mm},every loop/.style={}}
    \def\x{6}
    \def\y{0}
    \node[tensor,fill=red!50!white] (G) at (\x+4.5,\y) {$\Gt$};
    \draw[edge](G) -- (\x+5.5,\y) node [right] {\scalebox{0.7}{\dimcolor{$R_2$}}};
    \draw[edge](G) -- (\x+4.5,\y+0.25) -- (\x+5,\y+0.25) -- (\x+5,\y+0.1) -- (\x+5.5, \y+0.1) node [midway,above] {\scalebox{0.7}{\dimcolor{$R_1$}}};
    \draw[edge](G) -- (\x+4.5,\y-0.25) -- (\x+5,\y-0.25) -- (\x+5,\y-0.1) -- (\x+5.5, \y-0.1) node [midway,below] {\scalebox{0.7}{\dimcolor{$R_3$}}};
    \end{tikzpicture}
    =\vectorize(\Gt).
\end{align*}

注意到 $\Ub_1\kron\Ub_2\kron\Ub_3$ 是左正交矩阵，于是 
 $\norm{\Tt} = \norm{\vectorize(\Tt)} = \norm{(\Ub_1\kron\Ub_2\kron\Ub_3)\vectorize(\Tt)}= \norm{\vectorize(\Gt)} = \norm{\Gt}$.
    


    \item 若不对因子矩阵施加任何正交约束，并将核心张量固定为复制张量，
$
\begin{tikzpicture}[baseline=-0.5ex]
    \tikzset{tensor/.style = {minimum size = 0.5cm,shape = circle,thick,draw=black,inner sep = 0pt}, edge/.style = {thick,line width=.4mm},every loop/.style={}}
    \def\x{6}
    \def\y{0}
    \node[tensor,fill=red!50!white] (G) at (\x+4.5,\y) {$\Gt$};
    \draw[edge] (G) -- (\x+5.2,\y) node [midway, above] {\scalebox{0.7}{\textcolor{gray}{$R_4$}}};
    \draw[edge] (G) -- (\x+3.7,\y-0.3) node [below, near end] {\scalebox{0.7}{\dimcolor{$R_2$}}};
    \draw[edge] (G) -- (\x+5.2,\y-0.3) node [below, near end] {\scalebox{0.7}{\textcolor{gray}{$R_3$}}};
    \draw[edge] (G) -- (\x+3.7,\y) node [midway, above] {\scalebox{0.7}{\textcolor{gray}{$R_1$}}};
    \node[draw = none] at (\x+5.5,0){$=$};
    \draw  node[fill = black,circle,inner sep=0pt,minimum size=4pt] (m) at (\x+6.5,\y) {};
    \draw[edge] (m) -- (\x+7.1,\y) node [midway, above] {\scalebox{0.7}{\textcolor{gray}{$R_4$}}};
    \draw[edge] (m) -- (\x+5.9,\y-0.3) node [below, near end] {\scalebox{0.7}{\textcolor{gray}{$R_2$}}};
    \draw[edge] (m) -- (\x+7.1,\y-0.3) node [below, near end] {\scalebox{0.7}{\textcolor{gray}{$R_3$}}};
    \draw[edge] (m) -- (\x+5.9,\y) node [midway, above] {\scalebox{0.7}{\textcolor{gray}{$R_1$}}};
\end{tikzpicture},
$
便得到 CP 分解。 %
    \item \label{item:Tucker_orth} 若 $\Tt = \Gt\ttm{1}\Ab_1\ttm{2}\Ab_2\ttm{3}\Ab_3\ttm{4}\Ab_4$，其中 $\Ab_1,\Ab_2,\Ab_3$ 与 $\Ab_4$ 未必正交，则存在与 $\Gt$ 同尺寸的 $\Tilde{\Gt}$ 以及正交的 $\Qb_1,\Qb_2,\Qb_3$ 与 $\Qb_4$，使得 $\Tt = \Tilde{\Gt}\ttm{1}\Qb_1\ttm{2}\Qb_2\ttm{3}\Qb_3\ttm{4}\Qb_4$。
    \begin{proof}
对每个因子矩阵作 QR 分解，得到
\begin{align*}
\begin{tikzpicture}[baseline=-0.5ex]
    \tikzset{tensor/.style = {minimum size = 0.5cm,shape = circle,thick,draw=black,inner sep = 0pt}, edge/.style = {thick,line width=.4mm},every loop/.style={}}
    \def\x{6}
    \def\y{0}   \node[tensor,fill=green!50!red!50!white] (T) at (\x+1,\y) {$\Tt$};
    \draw[edge] (T) -- (\x+0.3,0) node [midway,above] {\scalebox{0.7}{\dimcolor{$d_1$}}};
    \draw[edge] (T) -- (\x+1.6,0) node [midway,above] {\scalebox{0.7}{\dimcolor{$d_4$}}};;
    \draw[edge] (T) -- (\x+0.4,-0.3) node [midway, below] {\scalebox{0.7}{\dimcolor{$d_2$}}};
    \draw[edge] (T) -- (\x+1.5,-0.4)node [midway, below] {\scalebox{0.7}{\dimcolor{$d_3$}}};
\end{tikzpicture}
&=
\begin{tikzpicture}[baseline=-0.5ex]
    \tikzset{tensor/.style = {minimum size = 0.5cm,shape = circle,thick,draw=black,inner sep = 0pt}, edge/.style = {thick,line width=.4mm},every loop/.style={}}
    \def\x{6}
    \def\y{0}   
    \node[tensor,fill=red!50!white] (G) at (\x+4.5,\y+1) {$\Gt$};
    \node[tensor,fill=red!20!white] (A_1) at (\x+3,\y){$\Ab_1$};
    \node[tensor,fill=red!40!white] (A_2) at (\x+4,\y){$\Ab_2$};
    \node[tensor, fill=blue!40!white] (A_3) at (\x+5,\y){$\Ab_3$};
    \node[tensor,fill=blue!20!white] (A_4) at (\x+6,\y){$\Ab_4$};
    \draw[edge] (G) -- (A_1) node [midway, above] {\scalebox{0.7}{\dimcolor{$R_1$}}};
    \draw[edge] (G) -- (A_2) node [midway, right=-0.1cm] {\scalebox{0.7}{\dimcolor{$R_2$}}};
    \draw[edge] (G) -- (A_3) node [midway, right=-0.05] {\scalebox{0.7}{\dimcolor{$R_3$}}};
    \draw[edge] (G) -- (A_4) node [midway, above] {\scalebox{0.7}{\dimcolor{$R_4$}}};
    \draw[edge] (A_1) -- (\x+3,\y-0.7) node [midway, left] {\scalebox{0.7}{\dimcolor{$d_1$}}};
    \draw[edge] (A_2) -- (\x+4,\y-0.7) node [midway, left] {\scalebox{0.7}{\dimcolor{$d_2$}}};
    \draw[edge] (A_3) -- (\x+5,\y-0.7) node [midway, left] {\scalebox{0.7}{\dimcolor{$d_3$}}};
    \draw[edge] (A_4) -- (\x+6,\y-0.7) node [midway, left] {\scalebox{0.7}{\dimcolor{$d_4$}}};
\end{tikzpicture}\\
&=
\begin{tikzpicture}[baseline=-0.5ex]
    \tikzset{tensor/.style = {minimum size = 0.5cm,shape = circle,thick,draw=black,inner sep = 0pt}, edge/.style = {thick,line width=.4mm},every loop/.style={}}
    \def\x{6}
    \def\y{0}   
   \node[tensor,fill=red!50!white] (G) at (\x+9,\y+1) {$\Gt$};
    \node[tensor,fill=red!20!white] (R_1) at (\x+7.5,\y){$\Rb_1$};
    \node[tensor,fill=red!40!white] (R_2) at (\x+8.5,\y){$\Rb_2$};
    \node[tensor, fill=blue!40!white] (R_3) at (\x+9.5,\y){$\Rb_3$};
    \node[tensor,fill=blue!20!white] (R_4) at (\x+10.5,\y){$\Rb_4$};
    \node[tensor, path picture={
    \fill[fill=red!20!white](path picture bounding box.south east)
    -- 
    (path picture bounding box.north west) --
    (path picture bounding box.north east) -- cycle;
    }, shift={(\x+7.5,\y-1)}] (Q_1)  {\scalebox{0.85}{$\Qb_1$}};  \node[tensor, path picture={
    \fill[fill=red!40!white](path picture bounding box.south east)
    -- 
    (path picture bounding box.north west) --
    (path picture bounding box.north east) -- cycle;
    }, shift={(\x+8.5,\y-1)}] (Q_2) {\scalebox{0.85}{$\Qb_2$}};
    \node[tensor, path picture={
    \fill[fill=blue!40!white](path picture bounding box.south east)
    -- 
    (path picture bounding box.north west) --
    (path picture bounding box.north east) -- cycle;
    }, shift={(\x+9.5,\y-1)}] (Q_3) {\scalebox{0.85}{$\Qb_3$}};
        \node[tensor, path picture={
    \fill[fill=blue!20!white](path picture bounding box.south east)
    -- 
    (path picture bounding box.north west) --
    (path picture bounding box.north east) -- cycle;
    }, shift={(\x+10.5,\y-1)}] (Q_4) {\scalebox{0.85}{$\Qb_4$}};
    \draw[edge] (G) -- (R_1) node [midway, above] {\scalebox{0.7}{\dimcolor{$R_1$}}};
    \draw[edge] (G) -- (R_2) node [midway, right=-0.05cm] {\scalebox{0.7}{\dimcolor{$R_2$}}};
    \draw[edge] (G) -- (R_3) node [midway, right=-0.05cm] {\scalebox{0.7}{\dimcolor{$R_3$}}};
    \draw[edge] (G) -- (R_4) node [midway, above] {\scalebox{0.7}{\dimcolor{$R_4$}}};
    \draw[edge] (R_1) -- (Q_1) node [midway, right] {\scalebox{0.7}{\dimcolor{$R_1$}}};
    \draw[edge] (R_2) -- (Q_2) node [midway, right] {\scalebox{0.7}{\dimcolor{$R_2$}}};
    \draw[edge] (R_3) -- (Q_3) node [midway, right=-0.1cm] {\scalebox{0.7}{\dimcolor{$R_3$}}};
    \draw[edge] (R_4) -- (Q_4) node [midway, right] {\scalebox{0.7}{\dimcolor{$R_4$}}};
    \draw[edge] (Q_1) -- (\x+7.5,\y-1.7) node [midway, left] {\scalebox{0.7}{\dimcolor{$d_1$}}};
    \draw[edge] (Q_2) -- (\x+8.5,\y-1.7) node [midway, left] {\scalebox{0.7}{\dimcolor{$d_2$}}};
    \draw[edge] (Q_3) -- (\x+9.5,\y-1.7) node [midway, left] {\scalebox{0.7}{\dimcolor{$d_3$}}};
    \draw[edge] (Q_4) -- (\x+10.5,\y-1.7) node [midway, left] {\scalebox{0.7}{\dimcolor{$d_4$}}};
    \draw[dash pattern=on 1pt off 2pt] plot[smooth cycle] coordinates { 
    (\x+7,\y) 
    (\x+7.3, \y+1.2) (\x+10.7,\y+1.2) (\x+10.7, \y-0.55)  (\x+7.3,\y-0.55)};
\end{tikzpicture}
=
\begin{tikzpicture}[baseline=-0.5ex]
    \tikzset{tensor/.style = {minimum size = 0.5cm,shape = circle,thick,draw=black,inner sep = 0pt}, edge/.style = {thick,line width=.4mm},every loop/.style={}}
    \def\x{6}
    \def\y{0}   
    \node[tensor,fill=blue!50!white] (G) at (\x+13.5,\y+1) {$\Tilde{\Gt}$};
    \node[tensor, path picture={
    \fill[fill=red!20!white](path picture bounding box.south east)
    -- 
    (path picture bounding box.north west) --
    (path picture bounding box.north east) -- cycle;
    }, shift={(\x+12.2,\y)}] (Q_1)  {\scalebox{0.85}{$\Qb_1$}};  \node[tensor, path picture={
    \fill[fill=red!40!white](path picture bounding box.south east)
    -- 
    (path picture bounding box.north west) --
    (path picture bounding box.north east) -- cycle;
    }, shift={(\x+13.2,\y)}] (Q_2) {\scalebox{0.85}{$\Qb_2$}};
    \node[tensor, path picture={
    \fill[fill=blue!40!white](path picture bounding box.south east)
    -- 
    (path picture bounding box.north west) --
    (path picture bounding box.north east) -- cycle;
    }, shift={(\x+14.2,\y)}] (Q_3) {\scalebox{0.85}{$\Qb_3$}};
        \node[tensor, path picture={
    \fill[fill=blue!20!white](path picture bounding box.south east)
    -- 
    (path picture bounding box.north west) --
    (path picture bounding box.north east) -- cycle;
    }, shift={(\x+15.2,\y)}] (Q_4) {\scalebox{0.85}{$\Qb_4$}};
     \draw[edge] (Q_1) -- (\x+12.2,\y-0.7) node [midway, left] {\scalebox{0.7}{\dimcolor{$d_1$}}};
    \draw[edge] (Q_2) -- (\x+13.2,\y-0.7) node [midway, left] {\scalebox{0.7}{\dimcolor{$d_2$}}};
    \draw[edge] (Q_3) -- (\x+14.2,\y-0.7) node [midway, left] {\scalebox{0.7}{\dimcolor{$d_3$}}};
    \draw[edge] (Q_4) -- (\x+15.2,\y-0.7) node [midway, left] {\scalebox{0.7}{\dimcolor{$d_4$}}};
    \draw[edge](G) -- (Q_1)node[midway,above]{\scalebox{0.7}{\dimcolor{$R_1$}}};
    \draw[edge](G) -- (Q_2)node[midway,right]{\scalebox{0.7}{\dimcolor{$R_2$}}};
    \draw[edge](G) -- (Q_3)node[midway,right]{\scalebox{0.7}{\dimcolor{$R_3$}}};
    \draw[edge](G) -- (\x+15.2,\y+0.25)node[midway,above]{\scalebox{0.7}{\dimcolor{$R_4$}}};
\end{tikzpicture}.
\end{align*}
\end{proof}
\item 对于一个 $N$ 阶 $d$ 维张量，在假设 $R_1=\dots=R_N=R$ 之下，其 Tucker 分解的参数个数为 $\bigo(R^N+NdR)$。
\item
本节最后说明，张量的 Tucker 分解的模 $n$ 矩阵化如何用 Kronecker 积来表达。
\begin{proposition}\label{decomp:tucker-kron}
设 $\At\in\RR^{d_1\times d_2\times\dots\times d_N}$。若 $\At=\Gt\ttm{1}\Ub_1\ttm{2}\dots\ttm{N}\Ub_N$，则
$$\At_{(i)}=\Ub_i\Gt_{(i)}\left(\Ub_1\kron\dots\kron\Ub_{(i-1)}\kron\Ub_{(i+1)}\kron\dots\kron\Ub_N\right)^\ts.$$
\begin{proof}
设 $N=3$ 且 $n=1$，则有
\begin{align*}
\At_{(1)}
&=
\left(\begin{tikzpicture}[baseline=\baseline]
    \tikzset{tensor/.style = {minimum size = 0.5cm,shape = circle,thick,draw=black,inner sep = 0pt}, edge/.style  = {thick,line width=.4mm},every loop/.style={}}
    \def\y{0}
    \def\x{0}
    \node[tensor,fill=green!50!red!50!white] (At) at (\x,\y){\scalebox{0.85}{$\At$}};
    \draw[edge] (\x+0.25,\y) -- (\x+0.6,\y)  node [midway,above] {\scalebox{0.7}{\dimcolor{$d_2$}}};
    \draw[edge] (\x-0.25,\y) -- (\x-0.6,\y)  node [midway,above] {\scalebox{0.7}{\dimcolor{$d_1$}}};
    \draw[edge] (\x,\y-0.25) -- (\x,\y-0.7)  node [midway,left] {\scalebox{0.7}{\dimcolor{$d_3$}}};
    \end{tikzpicture}\right)_{(1)}
=
\left(\begin{tikzpicture}[baseline=\baseline]
    \tikzset{tensor/.style = {minimum size = 0.5cm,shape = circle,thick,draw=black,inner sep = 0pt}, edge/.style   = {thick,line width=.4mm},every loop/.style={}}
    \def\y{0}
    \def\x{0}
  \node[tensor,fill=red!50!white] (A1) at (\x,\y){\scalebox{0.85}{$\Ub_1$}};
  \node[tensor,fill=blue!50!white] (A2) at (\x+1,\y){$\Ub_2$};
   \node[tensor,fill=blue!20!white] (A3) at (\x+2,\y){$\Ub_3$};
   \node[tensor,fill=red!60!blue!20!white] (m) at (\x+1,\y+1) {$\Gt$};
    \draw[edge] (A1) -- (m)node [midway,left] {\scalebox{0.7}{\dimcolor{$R_1$}}};
    \draw[edge] (A2) -- (m)node [midway,right=-0.05cm] {\scalebox{0.7}{\dimcolor{$R_2$}}};
    \draw[edge] (A3) -- (m)node [midway,right] {\scalebox{0.7}{\dimcolor{$R_3$}}};
    \draw[edge](A1) -- (\x,\y-0.7)node [midway,left] {\edgelabel{$d_1$}};
    \draw[edge](A2) -- (\x+1,\y-0.7)node [midway,left] {\edgelabel{$d_2$}};
    \draw[edge](A3) -- (\x+2,\y-0.7)node [midway,left] {\edgelabel{$d_3$}};
\end{tikzpicture}\right)_{(1)}\\
&=
\begin{tikzpicture}[baseline=\baseline]
    \tikzset{tensor/.style = {minimum size = 0.5cm,shape = circle,thick,draw=black,inner sep = 0pt}, edge/.style   = {thick,line width=.4mm},every loop/.style={}}
    \def\y{0}
    \def\x{0}
    \node[tensor,fill=green!50!red!50!white] (At) at (\x,\y){\scalebox{0.85}{$\At$}};
    \draw[edge] (\x-0.25,\y) -- (\x-0.6,\y)  node [midway,above] {\scalebox{0.7}{\dimcolor{$d_1$}}};
    \draw[edge] (\x+0.2,\y+0.15) -- (\x+0.5,\y+0.15)-- (\x+0.5,\y+0.05) -- (\x+1,\y+0.05) node [midway,above]
    {\scalebox{0.7}{\dimcolor{$d_2$}}};
    \draw[edge] (\x+0.2,\y-0.15) -- (\x+0.5,\y-0.15)-- (\x+0.5,\y-0.05)-- (\x+1,\y-0.05) node [midway,below]
    {\scalebox{0.7}{\dimcolor{$d_3$}}};
    \end{tikzpicture}
=
\begin{tikzpicture}[baseline=\baseline]
    \tikzset{tensor/.style = {minimum size = 0.5cm,shape = circle,thick,draw=black,inner sep = 0pt}, edge/.style   = {thick,line width=.4mm},every loop/.style={}}
    \def\y{0}
    \def\x{0}
    \node[tensor,fill=red!50!white] (A1) at (\x,\y){\scalebox{0.85}{$\Ub_1$}};
    \node[tensor,fill=blue!50!white] (A2) at (\x+2.7,\y+0.7){$\Ub_2$};
    \node[tensor,fill=blue!20!white] (A3) at (\x+2.7,\y-0.7){$\Ub_3$};
    \node[tensor,fill=red!60!blue!20!white] (m) at (\x+1,\y) {$\Gt$}; 
     \draw[edge] (\x+1.2,\y+0.15) -- (\x+1.5,\y+0.15)-- (\x+1.5,\y+0.05) -- (\x+2,\y+0.05) node [midway,above]
    {\scalebox{0.7}{\dimcolor{$R_2$}}};
    \draw[edge] (\x+1.2,\y-0.15) -- (\x+1.5,\y-0.15)-- (\x+1.5,\y-0.05)-- (\x+2,\y-0.05) node [midway,below]
    {\scalebox{0.7}{\dimcolor{$R_3$}}};
    \draw[edge] (A1) -- (m) node [midway,above] {\scalebox{0.7}{\dimcolor{$R_1$}}};
    \draw[edge] (A1) -- (\x-0.7,\y) node [midway,above] {\scalebox{0.7}{\dimcolor{$d_1$}}};
    \draw[edge](\x+2,\y-0.05) -- (\x+2,\y-0.7) -- (A3);
    \draw[edge](\x+2,\y+0.05) -- (\x+2,\y+0.7) -- (A2);
    \draw[edge] (A2) -- (\x+3.5,\y+0.7) -- (\x+3.5,\y+0.05)-- (\x+4,\y+0.05) node [midway,above]
    {\scalebox{0.7}{\dimcolor{$d_2$}}};
    \draw[edge] (A3) -- (\x+3.5,\y-0.7) -- (\x+3.5,\y-0.05) -- (\x+4,\y-0.05) node [midway,below]
    {\scalebox{0.7}{\dimcolor{$d_3$}}};
    \end{tikzpicture}
=
\Ub_1\Gt_{(1)}(\Ub_2\kron\Ub_3)^\ts.
\end{align*}
\end{proof}
\end{proposition}
\end{enumerate}
\end{remark}
\subsection{张量列~(TT) 分解}\label{sec:tt-decomposition}
张量列~(TT) 分解~\citep{oseledets2010approximation} 是另一种流行的张量分解模型。它将一个 $N$ 阶张量分解为 $N$ 个较小的三阶张量。 
张量 $\St\in\RR^{d_1\times\cdots\times d_N}$ 的秩 $(R_1,\dots,R_{N-1})$ \textit{张量列分解}将其分解为 $N$ 个三阶张量之积\footnote{首尾两个核心可视为矩阵而非三阶张量，因为它们各含一个维数为 1 的模。不过，把它们看作三阶张量能简化一些记法。}，其中 $\Gt_n\in\RR^{R_{n-1}\times d_n\times R_n}$，$n\in [N]$，且 $R_0=R_N=1$：
$$\St_{i_1,\cdots,i_N}= \sum_{r_0=1}^{R_0}\sum_{r_1=1}^{R_1} \cdots \sum_{r_{N-1}=1}^{R_{N-1}}\sum_{r_N=1}^{R_N} \prod_{n=1}^N\Gt_n(r_{n-1},i_n,r_n) := \TT(({\Gt_n})_{n=1}^N),$$ 
对所有 $i_1\in[d_1],\cdots,i_N\in[d_N]$ 成立。  
$\St$ 的\emph{TT 秩}是满足如下条件的最小
$(R_1,R_2,\dots,R_{N-1})$：$\St$ 存在精确的秩 $(R_1,R_2,\dots,R_{N-1})$ 的 TT 分解 $\St=\TT(({\Gt_n})_{n=1}^N)$。
与 CP 分解和 Tucker 分解类似，这里的秩可视为控制 TT 分解表达能力的超参数。也就是说，只要秩足够大，TT 分解就能表示任意张量。下面的张量网络展示了四阶张量的秩 $(R_1,R_2,R_3)$ TT 分解：

$$
\begin{tikzpicture}[baseline=-0.5ex]
    \tikzset{tensor/.style = {minimum size = 0.5cm,shape = circle,thick,draw=black,inner sep = 0pt}, edge/.style = {thick,line width=.4mm},every loop/.style={}}
    \def\x{6}
    \def\y{0}
     \node[tensor,fill=green!50!red!50!white] (B) at (\x+1,\y) {$\St$};
    \draw[edge] (B) -- (\x+0.3,0) node [midway,above] {\scalebox{0.7}{\textcolor{gray}{$d_1$}}};
    \draw[edge] (B) -- (\x+1.6,0) node [midway,above] {\scalebox{0.7}{\textcolor{gray}{$d_4$}}};;
    \draw[edge] (B) -- (\x+0.4,-0.3) node [midway, below] {\scalebox{0.7}{\textcolor{gray}{$d_2$}}};;
    \draw[edge] (B) -- (\x+1.5,-0.4)node [midway, below] {\scalebox{0.7}{\textcolor{gray}{$d_3$}}};;
\end{tikzpicture}=
\begin{tikzpicture}[baseline=\baseline]
    \tikzset{tensor/.style = {minimum size = 0.5cm,shape = circle,thick,draw=black,inner sep = 0pt}, edge/.style   = {thick,line width=.4mm},every loop/.style={}}
    \def\y{0}
    \def\x{0}
  \node[tensor,fill=red!30!white] (D) at (\x-1,\y){$\Gt_1$};
  \node[tensor,fill=red!50!white] (A) at (\x,\y){\scalebox{0.85}{$\Gt_2$}};
  \node[tensor,fill=blue!50!white] (B) at (\x+1,\y){$\Gt_3$};
   \node[tensor,fill=blue!20!white] (C) at (\x+2,\y){$\Gt_4$};
    \draw[edge](A) -- (\x,\y-0.7)node [midway,left] {\edgelabel{$d_2$}};
    \draw[edge](B) -- (\x+1,\y-0.7)node [midway,left] {\edgelabel{$d_3$}};
    \draw[edge](C) -- (\x+2,\y-0.7)node [midway,left] {\edgelabel{$d_4$}};
     \draw[edge](D) -- (\x-1,\y-0.7)node [midway,left] {\edgelabel{$d_1$}};
    \draw[edge](D) -- (A) node [midway,above] {\edgelabel{$R_1$}};
    \draw[edge](A) -- (B) node [midway,above] {\edgelabel{$R_2$}};
    \draw[edge](B) -- (C) node [midway,above] {\edgelabel{$R_3$}};
\end{tikzpicture}.
$$
TT 分解最早在量子物理与凝聚态物理社区中以 \emph{矩阵乘积态}（Matrix Product States, MPS）之名提出。在这一语境下，内部边的维数被称为 \textit{键维数}，自由腿则称为 \textit{物理维数}。
接下来我们说明如何在多项式时间内，利用逐次奇异值分解（SVD）计算张量 $\St\in\RR^{d_1\times\dots\times d_N}$ 的精确 TT 分解。该算法被称为 TT-SVD 算法~\citep{oseledets2010approximation}，但更早由量子物理社区引入，并被称作 \textit{逐次 Schmidt 分解}~\citep{vidal2004efficient,mcculloch2007density}。
在给出 TT-SVD 算法之前，先介绍三阶张量的左正交与右正交概念。
\begin{definition}\label{def:left-right-ortho}若 $(\At_n)_{(3)}^\ts(\At_n)_{(3)}=\Ib_{R_n}$，则称 $n\in[N]$ 对应的核心张量 $\At_n\in\RR^{R_{n-1}\times d_n\times R_n}$ 左正交；若 $(\At_n)_{(1)}(\At_n)_{(1)}^\ts=\Ib_{R_{n-1}}$，则称其右正交。这可以用张量网络
图示如下：
\begin{align*}
\begin{tikzpicture}[baseline=\baseline]
    \tikzset{tensor/.style = {minimum size = 0.5cm,shape = circle,thick,draw=black,inner sep = 0pt}, edge/.style = {thick,line width=.4mm},every loop/.style={}}
    \def\y{0}
    \def\x{0}
    \node[tensor, path picture={
    \fill[fill=red!30!white](path picture bounding box.south east)
        -- 
    (path picture bounding box.north west) --
    (path picture bounding box.north east) -- cycle;
    }, shift={(\x,\y)}] (A) {\scalebox{0.85}{$\At_n$}};
    \draw[edge](A) -- (\x,\y-0.7)node [midway,left] {\scalebox{0.7}{\dimcolor{$d_n$}}};
    \draw[edge](A) -- (\x-0.85,\y)node [midway,above] {\scalebox{0.7}{\dimcolor{$R_{n-1}$}}};
    \draw[edge](A) -- (\x+0.85,\y)node [midway,above] {\scalebox{0.7}{\dimcolor{$R_n$}}};
\end{tikzpicture}\ \  \text{is left-orthogonal}
\iff
\begin{tikzpicture}[baseline=\baseline]
    \tikzset{tensor/.style = {minimum size = 0.5cm,shape = circle,thick,draw=black,inner sep = 0pt}, edge/.style = {thick,line width=.4mm},every loop/.style={}}
    \def\y{0}
    \def\x{0}
    \node[tensor, path picture={
    \fill[fill=red!30!white](path picture bounding box.south west)
        -- 
    (path picture bounding box.north east) --
    (path picture bounding box.north west) -- cycle;
    }, shift={(\x,\y)}] (A1) {\scalebox{0.85}{$\At_n$}};
    \node[tensor, path picture={
    \fill[fill=red!30!white](path picture bounding box.south east)
        -- 
    (path picture bounding box.north west) --
    (path picture bounding box.north east) -- cycle;
    }, shift={(\x+1,\y)}] (A2) {\scalebox{0.85}{$\At_n$}};
    \draw[edge](A1)--(A2)node [midway,above] {\scalebox{0.7}{\dimcolor{$R_{n-1}$}}};
    \draw[edge](\x+0.23,\y-0.1) -- (\x+0.77,\y-0.1) node [midway,below] {\scalebox{0.7}{\dimcolor{$d_n$}}};
    \draw[edge](A1) -- (\x-0.85,\y)node [midway,above] {\scalebox{0.7}{\dimcolor{$R_{n}$}}};
    \draw[edge](A2) -- (\x+1.85,\y)node [midway,above] {\scalebox{0.7}{\dimcolor{$R_{n}$}}};
\end{tikzpicture}
=
\begin{tikzpicture}[baseline=\baseline]
    \tikzset{tensor/.style = {minimum size = 0.5cm,shape = circle,thick,draw=black,inner sep = 0pt}, edge/.style = {thick,line width=.4mm},every loop/.style={}}
    \def\y{0}
    \def\x{0}
    \draw[edge](\x,\y) -- (\x+1.5,\y)node [midway,above] {\scalebox{0.7}{\dimcolor{$R_n$}}};
\end{tikzpicture}=\Ib_{R_n},
\end{align*}
类似地，
\begin{align*}
\begin{tikzpicture}[baseline=\baseline]
    \tikzset{tensor/.style = {minimum size = 0.5cm,shape = circle,thick,draw=black,inner sep = 0pt}, edge/.style = {thick,line width=.4mm},every loop/.style={}}
    \def\y{0}
    \def\x{0}
    \node[tensor, path picture={
    \fill[fill=blue!30!white](path picture bounding box.south west)
        -- 
    (path picture bounding box.north east) --
    (path picture bounding box.north west) -- cycle;
    }, shift={(\x,\y)}] (A) {\scalebox{0.85}{$\At_n$}};
    \draw[edge](A) -- (\x,\y-0.7)node [midway,left] {\scalebox{0.7}{\dimcolor{$d_n$}}};
    \draw[edge](A) -- (\x-0.85,\y)node [midway,above] {\scalebox{0.7}{\dimcolor{$R_{n-1}$}}};
    \draw[edge](A) -- (\x+0.85,\y)node [midway,above] {\scalebox{0.7}{\dimcolor{$R_n$}}};
\end{tikzpicture}\ \  \text{is right-orthogonal}
\iff
\begin{tikzpicture}[baseline=\baseline]
    \tikzset{tensor/.style = {minimum size = 0.5cm,shape = circle,thick,draw=black,inner sep = 0pt}, edge/.style = {thick,line width=.4mm},every loop/.style={}}
    \def\y{0}
    \def\x{0}
    \node[tensor, path picture={
    \fill[fill=blue!30!white](path picture bounding box.south west)
        -- 
    (path picture bounding box.north east) --
    (path picture bounding box.north west) -- cycle;
    }, shift={(\x,\y)}] (A1) {\scalebox{0.85}{$\At_n$}};
    \node[tensor, path picture={
    \fill[fill=blue!30!white](path picture bounding box.south east)
        -- 
    (path picture bounding box.north west) --
    (path picture bounding box.north east) -- cycle;
    }, shift={(\x+1,\y)}] (A2) {\scalebox{0.85}{$\At_n$}};
    \draw[edge](A1)--(A2)node [midway,above] {\scalebox{0.7}{\dimcolor{$R_{n}$}}};
    \draw[edge](A1) -- (\x-0.85,\y)node [midway,above] {\scalebox{0.7}{\dimcolor{$R_{n-1}$}}};
    \draw[edge](A2) -- (\x+1.85,\y)node [midway,above] {\scalebox{0.7}{\dimcolor{$R_{n-1}$}}};
    \draw[edge](\x+0.23,\y-0.1) -- (\x+0.77,\y-0.1) node [midway,below] {\scalebox{0.7}{\dimcolor{$d_n$}}};
\end{tikzpicture}
=
\begin{tikzpicture}[baseline=\baseline]
    \tikzset{tensor/.style = {minimum size = 0.5cm,shape = circle,thick,draw=black,inner sep = 0pt}, edge/.style = {thick,line width=.4mm},every loop/.style={}}
    \def\y{0}
    \def\x{0}
    \draw[edge](\x,\y) -- (\x+1.5,\y)node [midway,above] {\scalebox{0.7}{\dimcolor{$R_{n-1}$}}};
\end{tikzpicture}=\Ib_{R_{n-1}}.
\end{align*}
\end{definition}
下面的定理说明任意张量都存在最小的张量列（TT）分解，其证明是构造性的，给出了 TT-SVD 算法的流程。
\begin{theorem}(\textbf{计算（正交）张量列分解。})\label{thm:tt-decomposition}
对任意 $\St\in\RR^{d_1\times\dots\times d_N}$，设 $\St_{([n])}\in\RR^{d_1\dots d_n\times d_{n+1}\dots d_N}$ 为把 $\St$ 的前 $n$  个模映射到 $\St_{[n]}$ 的行上所得到的矩阵化。则 $\St$ 的 TT 秩 $(R_1,\cdots,R_{N-1})$ 由
    下式给出
    $R_n = \rank~(\St_{([n])})$ 对所有 $n\in[N-1]$ 成立。
\end{theorem}
\begin{proof}
    设 $R_n = \rank~(\St_{([n])})$。我们先说明如何用逐次 QR 分解构造出 $\St\in\RR^{d_1\times\cdots\times d_N}$ 的一个秩为 $(R_1,R_2,\cdots,R_{N-1})$ 的 TT 分解。为简单起见，只考虑 $N=4$ 的情形，但上述论证很容易推广。

我们先对 $\ten S$ 的第一模矩阵化做 QR 分解，该矩阵化的秩为 $R_1$：
\begin{align*}
    \begin{tikzpicture}[baseline=\baseline]
    \tikzset{tensor/.style = {minimum size = 0.5cm,shape = circle,thick,draw=black,inner sep = 0pt}, edge/.style = {thick,line width=.4mm},every loop/.style={}}
    \def\y{0}
    \def\x{0}
    \node[tensor,fill=green!50!red!50!white] (A) at (\x-5.5,\y){\scalebox{0.85}{$\St$}};
    \draw[edge] (A) -- (\x-5.5,\y+0.75) node [midway,left]
    {\scalebox{0.7}{\dimcolor{$d_1$}}};
    \draw[edge] (A) -- (\x-5.5,\y-0.75) node [midway,left]
    {\scalebox{0.7}{\dimcolor{$d_3$}}};
    \draw[edge] (\x-5.75,\y+0.05) -- (\x-6.2,\y-0.3) node [midway,above]
    {\scalebox{0.7}{\dimcolor{$d_2$}}};
    \draw[edge] (\x-5.25,\y+0.05) -- (\x-4.7,\y-0.3) node [midway,above]
    {\scalebox{0.7}{\dimcolor{$d_4$}}};
    \node[draw=none] at (\x-3,\y){$\underrightarrow{\text{mode-1 matricization}}$};
    %
    %
    %
    %
    \node[tensor,fill=green!50!red!50!white] (A) at (\x-0.5,\y){\scalebox{0.85}{$\St$}};
    \draw[edge] (\x-0.35,\y+0.2) -- (\x,\y+0.2) --(\x,\y+0.1) -- (\x
    +0.35,\y+0.1) node [midway,above]
    {\scalebox{0.7}{\dimcolor{$d_2$}}};
    \draw[edge] (\x-0.35,\y-0.2) -- (\x,\y-0.2) -- (\x,\y-0.1) -- (\x+0.35,\y-0.1) node [midway,below]
    {\scalebox{0.7}{\dimcolor{$d_4$}}};
     \draw[edge] (A) -- (\x+0.35,\y)node [right]
    {\scalebox{0.7}{\dimcolor{$d_3$}}};
     \draw[edge] (A) -- (\x-1.25,\y) node [midway,above]
    {\scalebox{0.7}{\dimcolor{$d_1$}}};
    \node[draw=none] at (\x+1.25,\y){$\underrightarrow{\text{QR}}$};
    %
    %
    %
    \node[tensor, path picture={
    \fill[fill=blue!50!white](path picture bounding box.south east)
        -- 
    (path picture bounding box.north west) --
    (path picture bounding box.north east) -- cycle;
    }, shift={(\x+2,\y)}] (Q_1) {\scalebox{0.8}{$\Qb_1$}};
    \node[tensor,fill=red!50!white] (R_1) at (\x+3,\y){\scalebox{0.85}{$\Rb_1$}};
    \draw[edge] (Q_1) -- (\x+2,\y-0.75) node [midway,left]
    {\scalebox{0.7}{\dimcolor{$d_1$}}};
    \draw[edge] (Q_1) -- (R_1) node [midway,above]
    {\scalebox{0.7}{\dimcolor{$R_1$}}};
    \draw[edge] (\x+3.15,\y+0.2) -- (\x+3.5,\y+0.2) --(\x+3.5,\y+0.1) -- (\x+4,\y+0.1) node [midway,above]
    {\scalebox{0.7}{\dimcolor{$d_2$}}};
    \draw[edge] (\x+3.15,\y-0.2) -- (\x+3.5,\y-0.2) -- (\x+3.5,\y-0.1) -- (\x+4,\y-0.1) node [midway,below]
    {\scalebox{0.7}{\dimcolor{$d_4$}}};
    \draw[edge] (R_1) -- (\x+4,\y) node [right]
    {\scalebox{0.7}{\dimcolor{$d_3$}}};
    \node[draw=none] at (\x+6.5,\y){\text{Keep $\Qb_1$, resume with $\Rb_1$.}};
    \end{tikzpicture}
\end{align*}
然后，我们对上一步得到的 $\mat R_1$ 再做一次 QR 分解：
\begin{align*}
    \begin{tikzpicture}[baseline=\baseline]
    \tikzset{tensor/.style = {minimum size = 0.5cm,shape = circle,thick,draw=black,inner sep = 0pt}, edge/.style = {thick,line width=.4mm},every loop/.style={}}
    \def\y{0}
    \def\x{0}
    \node[tensor,fill=red!50!white] (R_1) at (\x-5.5,\y){\scalebox{0.85}{$\Rb_1$}};
    \draw[edge] (R_1) -- (\x-6.5,\y) node [midway,above]
    {\scalebox{0.7}{\dimcolor{$R_1$}}};
    \draw[edge] (\x-5.35,\y+0.2) -- (\x-5,\y+0.2) -- (\x-5,\y+0.1) -- (\x-4.5,\y+0.1) node [midway,above]
    {\scalebox{0.7}{\dimcolor{$d_2$}}};
    \draw[edge] (\x-5.35,\y-0.2) -- (\x-5,\y-0.2) -- (\x-5,\y-0.1) -- (\x-4.5,\y-0.1) node [midway,below]
    {\scalebox{0.7}{\dimcolor{$d_4$}}};
    \draw[edge] (R_1) -- (\x-4.5,\y) node [right]
    {\scalebox{0.7}{\dimcolor{$d_3$}}};
    \node[draw=none] at (\x-3.25,\y){$\underrightarrow{\text{reshape}}$};
    \node[tensor,fill=red!50!white] (R_1) at (\x-1.25,\y){\scalebox{0.85}{$\Rb_1$}};
    \draw[edge] (R_1) -- (\x-2.25,\y) node [midway,above]
    {\scalebox{0.7}{\dimcolor{$R_1$}}};
    \draw[edge] (\x-1.5,\y-0.1) -- (\x-2.25,\y-0.1) node [midway,below]
    {\scalebox{0.7}{\dimcolor{$d_2$}}};
    \draw[edge] (\x-1.1,\y-0.2) -- (\x-0.7,\y-0.2) -- (\x-0.7,\y-0.1) -- (\x-0.35,\y-0.1) node [midway,below]
    {\scalebox{0.7}{\dimcolor{$d_4$}}};
    \draw[edge] (R_1) -- (\x-0.35,\y) node [right]
    {\scalebox{0.7}{\dimcolor{$d_3$}}};
    \node[draw=none] at (\x+0.5,\y){$\underrightarrow{\text{QR}}$};
    \node[tensor, path picture={
    \fill[fill=blue!30!white](path picture bounding box.south east)
        -- 
    (path picture bounding box.north west) --
    (path picture bounding box.north east) -- cycle;
    }, shift={(\x+1.5,\y)}] (Q_2) {\scalebox{0.8}{$\Qb_2$}};
    \node[tensor,fill=red!50!white] (R_2) at (\x+2.5,\y){\scalebox{0.85}{$\Rb_2$}};
    \draw[edge] (Q_2) -- (\x+1.5,\y-0.75) node [midway,left=0.1cm]
    {\scalebox{0.7}{\dimcolor{$d_2$}}};
    \draw[edge] (\x+1.38,\y-0.22) -- (\x+1.38,\y-0.75) node [midway,right=0.1cm]
    {\scalebox{0.7}{\dimcolor{$R_1$}}};
    \draw[edge] (Q_2) -- (R_2) node [midway,above]
    {\scalebox{0.7}{\dimcolor{$R_2$}}};
    \draw[edge] (\x+2.65,\y-0.2) -- (\x+3,\y-0.2) -- (\x+3,\y-0.1) -- (\x+3.45,\y-0.1) node [midway,below]
    {\scalebox{0.7}{\dimcolor{$d_4$}}};
    \draw[edge] (R_2) -- (\x+3.45,\y) node [right]
    {\scalebox{0.7}{\dimcolor{$d_3$}}};
    \node[draw=none] at (\x+6,\y){\text{Keep $\Qb_2$}, resume with $\Rb_2$.};
    \end{tikzpicture}
\end{align*}\\
要使这一精确 QR 分解的秩为 $R_2$，我们首先需要证明：
将 $\Rb_1$ 重排成 $R_1d_2 \times d_3d_4$ 矩阵后，其秩至多为 $R_2$。注意到，由于 $\rank(\St_{([2])}) = R_2$，
存在 $\Ab \in \RR^{d_1 d_2 \times R_2}$ 与
$\Bb \in \RR^{R_2 \times d_3 d_4}$ 使得 $\St_{([2])} = \Ab\Bb$，即
$$
\begin{tikzpicture}[baseline=\baseline]
\tikzset{
  tensor/.style = {minimum size = 0.5cm,shape = circle,thick,draw=black,inner sep = 0pt},
  edge/.style   = {thick,line width=.4mm},
  every loop/.style={}
}
\def\y{0}
\def\x{0}
\node[tensor,fill=green!50!red!50!white] (A) at (\x-5.5,\y){\scalebox{0.85}{$\St$}};
\draw[edge] (A) -- (\x-5.5,\y+0.75) node [midway,left]
{\scalebox{0.7}{\dimcolor{$d_1$}}};
\draw[edge] (A) -- (\x-5.5,\y-0.75) node [midway,left]
{\scalebox{0.7}{\dimcolor{$d_3$}}};
\draw[edge] (\x-5.75,\y+0.05) -- (\x-6.2,\y-0.3) node [midway,above]
{\scalebox{0.7}{\dimcolor{$d_2$}}};
\draw[edge] (\x-5.25,\y+0.05) -- (\x-4.7,\y-0.3) node [midway,above]
{\scalebox{0.7}{\dimcolor{$d_4$}}};
\end{tikzpicture}
\ \ \ \underrightarrow{\text{reshape}}\ \ \ 
\begin{tikzpicture}[baseline=\baseline]
\tikzset{
  tensor/.style = {minimum size = 0.5cm,shape = circle,thick,draw=black,inner sep = 0pt},
  edge/.style   = {thick,line width=.4mm},
  every loop/.style={}
}
\def\y{0}
\def\x{0}
\node[tensor,fill=green!50!red!50!white] (A) at (\x+0.5,\y){\scalebox{0.85}{$\St$}};
\draw[edge] (\x+0.25,\y+0.1) -- (\x-0.35,\y+0.1) node [midway,above]
{\scalebox{0.7}{\dimcolor{$d_1$}}};
\draw[edge] (\x+0.75,\y-0.1) -- (\x+1.35,\y-0.1) node [midway,below]
{\scalebox{0.7}{\dimcolor{$d_4$}}};
\draw[edge] (A) -- (\x+1.35,\y) node [midway,above]
{\scalebox{0.7}{\dimcolor{$d_3$}}};
\draw[edge] (A) -- (\x-0.35,\y) node [midway,below]
{\scalebox{0.7}{\dimcolor{$d_2$}}};
\end{tikzpicture}
\ \ \ =\ \ \ 
\begin{tikzpicture}[baseline=\baseline]
\tikzset{
  tensor/.style = {minimum size = 0.5cm,shape = circle,thick,draw=black,inner sep = 0pt},
  edge/.style   = {thick,line width=.4mm},
  every loop/.style={}
}
\def\y{-0.05}
\def\x{0}
\node[tensor,fill=blue!50!red!50!white] (B) at (\x+6.5,\y){\scalebox{0.85}{$\Ab$}};
\node[tensor,fill=blue!50!red!20!white] (C) at (\x+7.5,\y){\scalebox{0.85}{$\Bb$}};
\draw[edge] (B) -- (C) node [midway,above]
{\scalebox{0.7}{\dimcolor{$R_2$}}};
\draw[edge] (B) -- (\x+5.5,\y) node [midway,below]
{\scalebox{0.7}{\dimcolor{$d_2$}}};
\draw[edge] (\x+6.25,\y+0.1) -- (\x+5.5,\y+0.1) node [midway,above]
{\scalebox{0.7}{\dimcolor{$d_1$}}};
\draw[edge] (C) -- (\x+8.5,\y) node [midway,above]
{\scalebox{0.7}{\dimcolor{$d_3$}}};
\draw[edge] (\x+7.75,\y-0.1) -- (\x+8.5,\y-0.1) node [midway,below]
{\scalebox{0.7}{\dimcolor{$d_4$}}};
\end{tikzpicture}.
$$

将 $\Qb_1$ 与第一模矩阵化 $\Sb_{(1)}$ 收缩，并利用上面的分解式，可得
\begin{align}\label{eq:tt-qr-fac}
    \begin{tikzpicture}[baseline=\baseline]
    \tikzset{tensor/.style = {minimum size = 0.5cm,shape = circle,thick,draw=black,inner sep = 0pt}, edge/.style = {thick,line width=.4mm},every loop/.style={}}
    \def\y{0}
    \def\x{0}
    \node[tensor, path picture={
    \fill[fill=blue!50!white](path picture bounding box.south west)
        -- 
    (path picture bounding box.north east) --
    (path picture bounding box.north west) -- cycle;
    }, shift={(\x-5.5,\y)}] (Q_1) {\scalebox{0.8}{$\Qb_1$}};
    \node[tensor,fill=green!50!red!50!white] (S) at (\x-4.5,\y){\scalebox{0.85}{$\St$}};
    \draw[edge](Q_1) -- (S)node [midway,above]
    {\scalebox{0.7}{\dimcolor{$d_1$}}};
    \draw[edge](Q_1) -- (\x-6.5,\y)node [midway,above]
    {\scalebox{0.7}{\dimcolor{$R_1$}}};
    \draw[edge](S) -- (\x-3.5,\y)node [right]
    {\scalebox{0.7}{\dimcolor{$d_3$}}};
    \draw[edge](\x-4.25,\y-0.1) -- (\x-3.5,\y-0.1)node [midway,below]
    {\scalebox{0.7}{\dimcolor{$d_4$}}};
    \draw[edge](\x-4.25,\y+0.1) -- (\x-3.5,\y+0.1)node [midway,above]
    {\scalebox{0.7}{\dimcolor{$d_2$}}};
\end{tikzpicture}
=
\begin{tikzpicture}[baseline=\baseline]
    \tikzset{tensor/.style = {minimum size = 0.5cm,shape = circle,thick,draw=black,inner sep = 0pt}, edge/.style = {thick,line width=.4mm},every loop/.style={}}
    \def\y{0}
    \def\x{0}
    \node[tensor, path picture={
    \fill[fill=blue!50!white](path picture bounding box.south west)
        -- 
    (path picture bounding box.north east) --
    (path picture bounding box.north west) -- cycle;
    }, shift={(\x-5.5,\y)}] (Q_1) {\scalebox{0.8}{$\Qb_1$}};
    \node[tensor,fill=green!50!red!50!white] (S) at (\x-4.5,\y){\scalebox{0.85}{$\St$}};
    \draw[edge](Q_1) -- (S)node [midway,above]
    {\scalebox{0.7}{\dimcolor{$d_1$}}};
    \draw[edge](Q_1) -- (\x-6.5,\y)node [midway,above]
    {\scalebox{0.7}{\dimcolor{$R_1$}}};
    \draw[edge](S) -- (\x-3.5,\y)node [right]
    {\scalebox{0.7}{\dimcolor{$d_3$}}};
    \draw[edge](\x-4.25,\y-0.1) -- (\x-3.5,\y-0.1)node [midway,below]
    {\scalebox{0.7}{\dimcolor{$d_4$}}};
    \draw[edge](\x-4.5,\y-0.25) -- (\x-4.5,\y-0.45) -- (\x-5.8,\y-0.45)node [midway,below]
    {\scalebox{0.7}{\dimcolor{$d_2$}}} -- (\x-5.8,\y-0.1)-- (\x-6.5,\y-0.1);
\end{tikzpicture}
=
\begin{tikzpicture}
    [baseline=\baseline]
    \tikzset{tensor/.style = {minimum size = 0.5cm,shape = circle,thick,draw=black,inner sep = 0pt}, edge/.style = {thick,line width=.4mm},every loop/.style={}}
    \def\y{0}
    \def\x{0}    
    \node[tensor, path picture={
    \fill[fill=blue!50!white](path picture bounding box.south west)
        -- 
    (path picture bounding box.north east) --
    (path picture bounding box.north west) -- cycle;
    }, shift={(\x,\y)}] (Q_1) {\scalebox{0.8}{$\Qb_1$}};
    \node[tensor,fill=blue!50!red!50!white] (A) at (\x+1,\y){\scalebox{0.85}{$\Ab$}};
    \node[tensor,fill=blue!50!red!20!white] (B) at (\x+2.5,\y){\scalebox{0.85}{$\Bb$}};
    \draw[edge] (A) -- (B) node [midway,above]
    {\scalebox{0.7}{\dimcolor{$R_2$}}};
    \draw[edge] (Q_1) -- (A) node [midway,above]
    {\scalebox{0.7}{\dimcolor{$d_1$}}};
    \draw[edge] (Q_1) -- (\x-1,\y) node [midway,above]
    {\scalebox{0.7}{\dimcolor{$R_1$}}};
    \draw[edge] (\x+2.74,\y-0.1) -- (\x+3.5,\y-0.1) node [midway,below]
    {\scalebox{0.7}{\dimcolor{$d_4$}}};
    \draw[edge] (B) -- (\x+3.5,\y) node [midway,above]
    {\scalebox{0.7}{\dimcolor{$d_3$}}};
 \draw[edge](\x+1,\y-0.25) -- (\x+1,\y-0.45) -- (\x-0.3,\y-0.45)node [midway,below]
    {\scalebox{0.7}{\dimcolor{$d_2$}}} -- (\x-0.3,\y-0.1)-- (\x-1,\y-0.1);
\end{tikzpicture}.
\end{align}
此外，由于 $\Qb_1$ 正交且来自 $\Sb_{(1)}$ 的 QR 分解，故 $\Qb_1\Sb_{(1)}$ 等于 $\Rb_1$，于是：
\begin{align}\label{eq:tt-equal-r1}
    \begin{tikzpicture}[baseline=\baseline]
    \tikzset{tensor/.style = {minimum size = 0.5cm,shape = circle,thick,draw=black,inner sep = 0pt}, edge/.style = {thick,line width=.4mm},every loop/.style={}}
    \def\y{0}
    \def\x{0}
    \node[tensor, path picture={
    \fill[fill=blue!50!white](path picture bounding box.south west)
        -- 
    (path picture bounding box.north east) --
    (path picture bounding box.north west) -- cycle;
    }, shift={(\x-5.5,\y)}] (Q_1) {\scalebox{0.8}{$\Qb_1$}};
    \node[tensor,fill=green!50!red!50!white] (S) at (\x-4.5,\y){\scalebox{0.85}{$\St$}};
    \draw[edge](Q_1) -- (S)node [midway,above]
    {\scalebox{0.7}{\dimcolor{$d_1$}}};
    \draw[edge](Q_1) -- (\x-6.5,\y)node [midway,above]
    {\scalebox{0.7}{\dimcolor{$R_1$}}};
    \draw[edge](S) -- (\x-3.5,\y)node [right]
    {\scalebox{0.7}{\dimcolor{$d_3$}}};
    \draw[edge](\x-4.25,\y-0.1) -- (\x-3.5,\y-0.1)node [midway,below]
    {\scalebox{0.7}{\dimcolor{$d_4$}}};
    \draw[edge](\x-4.5,\y-0.25) -- (\x-4.5,\y-0.45) -- (\x-5.8,\y-0.45)node [midway,below]
    {\scalebox{0.7}{\dimcolor{$d_2$}}} -- (\x-5.8,\y-0.1)-- (\x-6.5,\y-0.1);
   \def\x{-3}
     \node[draw=none] at (\x+0.2,\y){$=$};
   \node[tensor, path picture={
    \fill[fill=blue!50!white](path picture bounding box.south west)
        -- 
    (path picture bounding box.north east) --
    (path picture bounding box.north west) -- cycle;
    }, shift={(\x+1.5,\y)}] (Q_1) {\scalebox{0.8}{$\Qb_1$}};
    \node[tensor, path picture={
    \fill[fill=blue!50!white](path picture bounding box.south east)
        -- 
    (path picture bounding box.north west) --
    (path picture bounding box.north east) -- cycle;
    }, shift={(\x+2.5,\y)}] (Q_1t) {\scalebox{0.8}{$\Qb_1$}};
    \node[tensor,fill=red!50!white] (R_1) at (\x+3.5,\y){\scalebox{0.85}{$\Rb_1$}};
    \draw[edge](Q_1) -- (\x+0.5,\y) node [midway,above]
    {\scalebox{0.7}{\dimcolor{$R_1$}}};
    \draw[edge](Q_1) -- (Q_1t) node [midway,above]
    {\scalebox{0.7}{\dimcolor{$d_1$}}};
 \draw[edge](\x+3.5,\y-0.25) -- (\x+3.5,\y-0.45) -- (\x+1.1,\y-0.45)node [midway,below]
    {\scalebox{0.7}{\dimcolor{$d_2$}}} -- (\x+1.1,\y-0.1)-- (\x+0.5,\y-0.1);
    \draw[edge] (\x+3.65,\y-0.2) -- (\x+4,\y-0.2) -- (\x+4,\y-0.1) -- (\x+4.5,\y-0.1) node [midway,below]
    {\scalebox{0.7}{\dimcolor{$d_4$}}};
    \draw[edge] (R_1) -- (\x+4.5,\y) node [right]
    {\scalebox{0.7}{\dimcolor{$d_3$}}};
    \draw[edge](Q_1t) -- (R_1) node [midway,above]
    {\scalebox{0.7}{\dimcolor{$R_1$}}};
    \node[draw=none] at (\x+5.25,\y){$=$};
    \node[tensor,fill=red!50!white] (R_1) at (\x+6.5,\y){\scalebox{0.85}{$\Rb_1$}};
    \draw[edge] (\x+6.28,\y-0.1) -- (\x+5.5, \y-0.1) node [midway,below]
    {\scalebox{0.7}{\dimcolor{$d_2$}}};
    \draw[edge] (\x+6.65,\y-0.2) -- (\x+7,\y-0.2) -- (\x+7,\y-0.1) -- (\x+7.5,\y-0.1) node [midway,below]
    {\scalebox{0.7}{\dimcolor{$d_4$}}};
    \draw[edge] (R_1) -- (\x+7.5,\y) node [midway,above]
    {\scalebox{0.7}{\dimcolor{$d_3$}}};
    \draw[edge] (R_1) -- (\x+5.5,\y) node [midway, above]
    {\scalebox{0.7}{\dimcolor{$R_1$}}};
\end{tikzpicture}.
\end{align}
结合恒等式~\eqref{eq:tt-qr-fac} 与~\eqref{eq:tt-equal-r1}，可得
$$
\begin{tikzpicture}[baseline=\baseline]
    \tikzset{tensor/.style = {minimum size = 0.5cm,shape = circle,thick,draw=black,inner sep = 0pt}, edge/.style = {thick,line width=.4mm},every loop/.style={}}
    \def\x{-6}
    \def\y{0}
    \node[tensor,fill=red!50!white] (R_1) at (\x+6.5,\y){\scalebox{0.85}{$\Rb_1$}};
    \draw[edge] (\x+6.28,\y-0.1) -- (\x+5.5, \y-0.1) node [midway,below]
    {\scalebox{0.7}{\dimcolor{$d_2$}}};
    \draw[edge] (\x+6.65,\y-0.2) -- (\x+7,\y-0.2) -- (\x+7,\y-0.1) -- (\x+7.5,\y-0.1) node [midway,below]
    {\scalebox{0.7}{\dimcolor{$d_4$}}};
    \draw[edge] (R_1) -- (\x+7.5,\y) node [midway,above]
    {\scalebox{0.7}{\dimcolor{$d_3$}}};
    \draw[edge] (R_1) -- (\x+5.5,\y) node [midway, above]
    {\scalebox{0.7}{\dimcolor{$R_1$}}};
\end{tikzpicture}
=
\begin{tikzpicture}
    [baseline=\baseline]
    \tikzset{tensor/.style = {minimum size = 0.5cm,shape = circle,thick,draw=black,inner sep = 0pt}, edge/.style = {thick,line width=.4mm},every loop/.style={}}
    \def\y{0}
    \def\x{0}    
    \node[tensor, path picture={
    \fill[fill=blue!50!white](path picture bounding box.south west)
        -- 
    (path picture bounding box.north east) --
    (path picture bounding box.north west) -- cycle;
    }, shift={(\x,\y)}] (Q_1) {\scalebox{0.8}{$\Qb_1$}};
    \node[tensor,fill=blue!50!red!50!white] (A) at (\x+1,\y){\scalebox{0.85}{$\Ab$}};
    \node[tensor,fill=blue!50!red!20!white] (B) at (\x+2.5,\y){\scalebox{0.85}{$\Bb$}};
    \draw[edge] (A) -- (B) node [midway,above]
    {\scalebox{0.7}{\dimcolor{$R_2$}}};
    \draw[edge] (Q_1) -- (A) node [midway,above]
    {\scalebox{0.7}{\dimcolor{$d_1$}}};
    \draw[edge] (Q_1) -- (\x-1,\y) node [midway,above]
    {\scalebox{0.7}{\dimcolor{$R_1$}}};
    \draw[edge] (\x+2.74,\y-0.1) -- (\x+3.5,\y-0.1) node [midway,below]
    {\scalebox{0.7}{\dimcolor{$d_4$}}};
    \draw[edge] (B) -- (\x+3.5,\y) node [midway,above]
    {\scalebox{0.7}{\dimcolor{$d_3$}}};
 \draw[edge](\x+1,\y-0.25) -- (\x+1,\y-0.45) -- (\x-0.3,\y-0.45)node [midway,below]
    {\scalebox{0.7}{\dimcolor{$d_2$}}} -- (\x-0.3,\y-0.1)-- (\x-1,\y-0.1);
    \draw[dotted](\x+1.5,\y+0.75) -- (\x+1.5,\y-0.75);
\end{tikzpicture}.
$$
利用定理~\ref{thm: matricization-rank}，上面的张量网络图表明：把 $\Rb_1$ 重排为 $R_1d_2\times d_3d_4$
矩阵后，所得矩阵的秩至多为 $R_2$。
对 $\Rb_2$ 重复上述步骤，最终得到秩为 $(R_1,R_2,R_3)$ 的 TT 分解，其中 $R_n = \rank{(\St_{([n])})}$ 对 $n\in [3]$ 成立，即
\begin{align}\label{eq:tt-format}
    \begin{tikzpicture}[baseline=\baseline]
    \tikzset{tensor/.style = {minimum size = 0.5cm,shape = circle,thick,draw=black,inner sep = 0pt}, edge/.style = {thick,line width=.4mm},every loop/.style={}}
    \def\y{0}
    \def\x{0}
    \node[tensor,fill=green!50!red!50!white] (S) at (\x,\y){\scalebox{0.85}{$\St$}};
    \draw[edge] (S) -- (\x,\y+0.75) node [midway,left]
    {\scalebox{0.7}{\dimcolor{$d_1$}}};
    \draw[edge] (S) -- (\x,\y-0.75) node [midway,left]
    {\scalebox{0.7}{\dimcolor{$d_3$}}};
    \draw[edge] (\x+0.25,\y+0.05) -- (\x+0.8,\y-0.3) node [midway,above]
    {\scalebox{0.7}{\dimcolor{$d_4$}}};
    \draw[edge] (\x-0.25,\y+0.05) -- (\x-0.8,\y-0.3) node [midway,above]
    {\scalebox{0.7}{\dimcolor{$d_2$}}};
\end{tikzpicture}
=
\begin{tikzpicture}[baseline=\baseline]
    \tikzset{tensor/.style = {minimum size = 0.5cm,shape = circle,thick,draw=black,inner sep = 0pt}, edge/.style = {thick,line width=.4mm},every loop/.style={}}
    \def\y{0}
    \def\x{0}
    \node[tensor, path picture={
    \fill[fill=blue!50!white](path picture bounding box.south east)
        -- 
    (path picture bounding box.north west) --
    (path picture bounding box.north east) -- cycle;
    }, shift={(\x,\y)}] (Q_1) {\scalebox{0.8}{$\Qt_1$}};
    \node[tensor, path picture={
    \fill[fill=blue!35!white](path picture bounding box.south east)
        -- 
    (path picture bounding box.north west) --
    (path picture bounding box.north east) -- cycle;
    }, shift={(\x+1,\y)}] (Q_2) {\scalebox{0.8}{$\Qt_2$}};
    \node[tensor, path picture={
    \fill[fill=blue!20!white](path picture bounding box.south east)
        -- 
    (path picture bounding box.north west) --
    (path picture bounding box.north east) -- cycle;
    }, shift={(\x+2,\y)}] (Q_3) {\scalebox{0.8}{$\Qt_3$}};
    \draw[edge](Q_1) -- (Q_2)node [midway,above]  {\scalebox{0.7}{\dimcolor{$R_1$}}};
    \draw[edge](Q_1) -- (\x,\y-0.85)node [midway,left]  {\scalebox{0.7}{\dimcolor{$d_1$}}};
    \draw[edge](Q_2) -- (\x+1,\y-0.85)node [midway,left]  {\scalebox{0.7}{\dimcolor{$d_2$}}};
    \draw[edge](Q_2) -- (Q_3)node [midway,above]  {\scalebox{0.7}{\dimcolor{$R_2$}}};
    \draw[edge](Q_3) -- (\x+2,\y-0.85)node [midway,left]  {\scalebox{0.7}{\dimcolor{$d_3$}}};
    \node[tensor,fill=red!50!white] (R_1) at (\x+3,\y){\scalebox{0.85}{$\Rb_3$}};
    \draw[edge](R_1) -- (\x+3,\y-0.85)node [midway,left]  {\scalebox{0.7}{\dimcolor{$d_4$}}};
    \draw[edge](Q_3) -- (R_1) node [midway, above]{\scalebox{0.7}{\dimcolor{$R_3$}}};
    \end{tikzpicture}.
    \end{align}
由此得到的 TT 分解中，除最后一个核心张量外，其余核心张量均是正交的。

上述算法表明，若对所有 $n$ 都有 $\tilde R_n = \rank(\St_{([n])})$，则存在一个秩为 $(\tilde R_1,\dots, \tilde R_{N-1})$ 的 TT 分解，这说明 $\St$ 的 TT 秩 $(R_1,\dots, R_{N-1})$ 以各矩阵化的秩为上界：对所有 $n$ 有 $ R_n \leq \rank(\St_{([n])} )$。也容易证明 TT 秩以各矩阵化的秩为下界。事实上，设 $\St$ 的 TT 秩为 $(R_1,\dots, R_{N-1})$，则存在一个秩为 $(R_1,\dots, R_{N-1})$ 的 TT 分解。利用定理~\ref{thm: matricization-rank}，容易看出这意味着对所有 $n$ 有 $\rank(\St_{([n])} )\leq  R_n$。
\end{proof}
尽管上面的证明用逐次 QR 分解构造出精确的 TT 分解，但通常被称为 TT-SVD 的算法是对逐次展开施行奇异值分解。在精确情形下，保留全部非零奇异值即得 TT 秩
$
R_n = \rank(\St_{([n])} )
$.
奇异值分解这一形式的优点在于，它直接给出了一个低秩近似算法：在每一步都可使用截断奇异值分解，或截到给定的 TT 秩，或按被舍弃奇异值不超过给定容差来截断~\citep{oseledets2010approximation,holtz2012manifolds}。

另一种被广泛使用的、把张量分解为张量列（tensor train, TT）格式的方法是交替最小二乘（Alternating Least Squares, ALS）算法。在介绍这一方法之前，我们先描述 TT 分解的标准形。
\begin{definition}\label{def:canonical-form}
   若所有 $n<j$ 的核心张量都是左正交的、所有 $n>j$ 的核心张量都是右正交的，则称 TT 分解 $\TT(({\At_n})_{n=1}^N)\in\RR^{d_1\times\dots\times d_N}$ 关于固定指标 $j\in [N]$ 处于标准形；例如，下面的 TT 分解就关于其第三个核心张量处于标准形：
$$
\begin{tikzpicture}[baseline=-0.5ex]
    \tikzset{tensor/.style = {minimum size = 0.5cm,shape = circle,thick,draw=black,inner sep = 0pt}, edge/.style = {thick,line width=.4mm},every loop/.style={}}
    \def\x{0}
    \def\y{0}
 \node[tensor, path picture={
        \fill[fill=red!30!white](path picture bounding box.south east)
        -- 
        (path picture bounding box.north west) --
        (path picture bounding box.north east) -- cycle;
        }, shift={(\x,\y)}] (A) {\scalebox{0.85}{$\At_1$}};
        \draw[edge](A) -- (\x,\y-0.7) node [midway,left] {\scalebox{0.7}{\dimcolor{$d_1$}}};
       \node[tensor, path picture={
        \fill[fill=red!50!white](path picture bounding box.south east)
        -- 
        (path picture bounding box.north west) --
        (path picture bounding box.north east) -- cycle;
        }, shift={(\x+1,\y)}] (B) {\scalebox{0.85}{$\At_2$}};
         \draw[edge](B) -- (\x+1,\y-0.7)node [midway,left] {\scalebox{0.7}{\dimcolor{$d_2$}}};
         \node[tensor,fill=blue!50!white] (C) at (\x+2,\y){$\At_3$};
          \draw[edge](C) -- (\x+2,\y-0.7)node [midway,left] {\scalebox{0.7}{\dimcolor{$d_3$}}};
          \node[tensor, path picture={
        \fill[fill=blue!60!red!50!](path picture bounding box.south west)
        -- 
        (path picture bounding box.north east) --
        (path picture bounding box.north west) -- cycle;
        }, shift={(\x+3,\y)}] (D) {\scalebox{0.85}{$\At_4$}};
         \draw[edge](D) -- (\x+3,\y-0.7)node [midway,left] {\scalebox{0.7}{\dimcolor{$d_4$}}};
           \node[tensor, path picture={
        \fill[fill=blue!20!white](path picture bounding box.south west)
        -- 
        (path picture bounding box.north east) --
        (path picture bounding box.north west) -- cycle;
        }, shift={(\x+4,\y)}] (E) {\scalebox{0.85}{$\At_5$}};
         \draw[edge](E) -- (\x+4,\y-0.7)node [midway,left] {\scalebox{0.7}{\dimcolor{$d_5$}}};
    \draw[edge](A) -- (B) node [midway,above] {\scalebox{0.7}{\dimcolor{$R_1$}}};
    \draw[edge](B) -- (C) node [midway,above] {\scalebox{0.7}{\dimcolor{$R_2$}}};
    \draw[edge](C) -- (D) node [midway,above] {\scalebox{0.7}{\dimcolor{$R_3$}}};
    \draw[edge](D) -- (E) node [midway,above] {\scalebox{0.7}{\dimcolor{$R_4$}}};
\end{tikzpicture}.
$$
$\At_3$ 左侧的核心张量都是左正交的，右侧的核心张量都是右正交的。核心张量 $\At_3$ 称为正交中心。注意，对核心张量施加一系列 QR 分解，就能把任意 TT 分解高效地转化为关于任意指标 $j\in [N]$ 的标准形~\citep{holtz2012alternating,evenbly2018gauge}。
\end{definition}
\paragraph{用交替最小二乘~(ALS) 计算 TT 分解} 单站点张量列交替最小二乘（TT-ALS）方法~\citep{holtz2012alternating} 从一个处于标准形的 TT 分解出发，该分解由一个粗略近似初始化~(例如用随机核心张量，或用 TT-SVD 得到)。
在第一步中，分解的核心张量 
$\At_1$ 是非正交的。在自左向右（相应地，自右向左）的扫描过程中，算法固定除第 $n$ 个核心张量 $\At_n$ 之外的所有核心张量，并对其加以优化，使其与目标张量之间的距离最小。由此得到的极小化问题是一个简单的最小二乘问题，故得名 ALS。  
更新 $\At_n$ 之后，施加一次 QR 分解以保证数值稳定性并维持标准形。非标准正交的因子随后被合并到左侧（或右侧，取决于扫描方向），这一步称为 \textit{核心张量正交化}。
\begin{figure}[h!]
    \centering
\begin{tikzpicture}[baseline=-0.5ex]
    \tikzset{tensor/.style = {minimum size = 0.5cm,shape = circle,thick,draw=black,inner sep = 0pt}, edge/.style = {thick,line width=.4mm},every loop/.style={}}
    \def\x{0}
    \def\y{0}
    \node[tensor, 
        fill=red!30!white, shift = {(\x-4,\y)}] (A) {\scalebox{0.85}{$\At_1$}};
        \draw[edge](A) -- (\x-4,\y-0.6) node [midway,left] {\scalebox{0.8}{\dimcolor{$d_1$}}};
       \node[tensor, path picture={
        \fill[fill=red!50!white](path picture bounding box.south west)
        -- 
        (path picture bounding box.north east) --
        (path picture bounding box.north west) -- cycle;
        }, shift={(\x-3,\y)}] (B) {\scalebox{0.85}{$\At_2$}};
       \draw[edge](B) -- (\x-3,\y-0.6)node [midway,left] {\scalebox{0.7}{\dimcolor{$d_2$}}};
        \node[tensor, path picture={
        \fill[fill=blue!50!white](path picture bounding box.south west)
        -- 
        (path picture bounding box.north east) --
        (path picture bounding box.north west) -- cycle;
        }, shift={(\x-2,\y)}] (C) {\scalebox{0.85}{$\At_3$}};
        \draw[edge](C) -- (\x-2,\y-0.6)node [midway,left] {\scalebox{0.7}{\dimcolor{$d_3$}}};
        \node[tensor, path picture={
       \fill[fill=blue!60!red!50!](path             picture bounding box.south west)
        -- 
        (path picture bounding box.north east) --
       (path picture bounding box.north west) -- cycle;}, shift={(\x-1,\y)}] (D) {\scalebox{0.85}{$\At_4$}};
    \draw[edge](D) -- (\x-1,\y-0.6)node [midway,left] {\scalebox{0.7}{\dimcolor{$d_4$}}};
    \node[tensor, path picture={
    \fill[fill= blue!20!white](path             picture bounding box.south west)
    -- 
    (path picture bounding box.north east) --
    (path picture bounding box.north west) -- cycle;}, shift={(\x,\y)}] (E) {\scalebox{0.85}{$\At_5$}};
    \draw[edge](E) -- (\x,\y-0.6)node [midway,left] {\scalebox{0.7}{\dimcolor{$d_5$}}};
    \draw[edge](A) -- (B) node [midway,above] {\scalebox{0.7}{\dimcolor{$R_1$}}};
    \draw[edge](B) -- (C) node [midway,above] {\scalebox{0.7}{\dimcolor{$R_2$}}};
    \draw[edge](C) -- (D) node [midway,above] {\scalebox{0.7}{\dimcolor{$R_3$}}};
    \draw[edge](D) -- (E) node [midway,above] {\scalebox{0.7}{\dimcolor{$R_4$}}};
\end{tikzpicture}~~~\text{step:~1}\\
\begin{tikzpicture}[baseline=-0.5ex]
    \tikzset{tensor/.style = {minimum size = 0.5cm,shape = circle,thick,draw=black,inner sep = 0pt}, edge/.style = {thick,line width=.4mm},every loop/.style={}}
    \def\x{0}
    \def\y{0}
    \node[tensor, path picture = {
        \fill[fill=red!30!white] (path picture bounding box.south east)
        -- 
        (path picture bounding box.north west) --
        (path picture bounding box.north east) -- cycle;
        }, shift = {(\x-4,\y)}] (A) {\scalebox{0.85}{$\At_1$}};
        \draw[edge](A) -- (\x-4,\y-0.6) node [midway,left] {\scalebox{0.8}{\dimcolor{$d_1$}}};
         \draw  node[fill = black,circle,inner sep=0pt,minimum size=6pt] (m) at (\x-3.5,\y) {};
       \node[tensor, path picture={
        \fill[fill=red!50!white](path picture bounding box.south west)
        -- 
        (path picture bounding box.north east) --
        (path picture bounding box.north west) -- cycle;
        }, shift={(\x-3,\y)}] (B) {\scalebox{0.85}{$\At_2$}};
       \draw[edge](B) -- (\x-3,\y-0.6)node [midway,left] {\scalebox{0.7}{\dimcolor{$d_2$}}};
        \node[tensor, path picture={
        \fill[fill=blue!50!white](path picture bounding box.south west)
        -- 
        (path picture bounding box.north east) --
        (path picture bounding box.north west) -- cycle;
        }, shift={(\x-2,\y)}] (C) {\scalebox{0.85}{$\At_3$}};
        \draw[edge](C) -- (\x-2,\y-0.6)node [midway,left] {\scalebox{0.7}{\dimcolor{$d_3$}}};
        \node[tensor, path picture={
       \fill[fill=blue!60!red!50!](path             picture bounding box.south west)
        -- 
        (path picture bounding box.north east) --
       (path picture bounding box.north west) -- cycle;}, shift={(\x-1,\y)}] (D) {\scalebox{0.85}{$\At_4$}};
    \draw[edge](D) -- (\x-1,\y-0.6)node [midway,left] {\scalebox{0.7}{\dimcolor{$d_4$}}};
    \node[tensor, path picture={
    \fill[fill= blue!20!white](path             picture bounding box.south west)
    -- 
    (path picture bounding box.north east) --
    (path picture bounding box.north west) -- cycle;}, shift={(\x,\y)}] (E) {\scalebox{0.85}{$\At_5$}};
    \draw[edge](E) -- (\x,\y-0.6)node [midway,left] {\scalebox{0.7}{\dimcolor{$d_5$}}};
    \draw[edge](A) -- (B) node [midway,above] {\scalebox{0.7}{\dimcolor{$R_1$}}};
    \draw[edge](B) -- (C) node [midway,above] {\scalebox{0.7}{\dimcolor{$R_2$}}};
    \draw[edge](C) -- (D) node [midway,above] {\scalebox{0.7}{\dimcolor{$R_3$}}};
    \draw[edge](D) -- (E) node [midway,above] {\scalebox{0.7}{\dimcolor{$R_4$}}};
\end{tikzpicture}~~~~~~\text{QR}\\
\begin{tikzpicture}[baseline=-0.5ex]
    \tikzset{tensor/.style = {minimum size = 0.5cm,shape = circle,thick,draw=black,inner sep = 0pt}, edge/.style = {thick,line width=.4mm},every loop/.style={}}
    \def\x{0}
    \def\y{0}
    \node[tensor, path picture = {
        \fill[fill=red!30!white] (path picture bounding box.south east)
        -- 
        (path picture bounding box.north west) --
        (path picture bounding box.north east) -- cycle;
        }, shift = {(\x-4,\y)}] (A) {\scalebox{0.85}{$\At_1$}};
        \draw[edge](A) -- (\x-4,\y-0.6) node [midway,left] {\scalebox{0.8}{\dimcolor{$d_1$}}};
       \node[tensor, 
        fill=red!50!white, shift = {(\x-3,\y)}] (B) {\scalebox{0.85}{$\At_2$}};
       \draw[edge](B) -- (\x-3,\y-0.6)node [midway,left] {\scalebox{0.7}{\dimcolor{$d_2$}}};
        \node[tensor, path picture={
        \fill[fill=blue!50!white](path picture bounding box.south west)
        -- 
        (path picture bounding box.north east) --
        (path picture bounding box.north west) -- cycle;
        }, shift={(\x-2,\y)}] (C) {\scalebox{0.85}{$\At_3$}};
        \draw[edge](C) -- (\x-2,\y-0.6)node [midway,left] {\scalebox{0.7}{\dimcolor{$d_3$}}};
        \node[tensor, path picture={
       \fill[fill=blue!60!red!50!](path             picture bounding box.south west) -- (path picture bounding box.north east) -- (path picture bounding box.north west) -- cycle;}, shift={(\x-1,\y)}] (D) {\scalebox{0.85}{$\At_4$}};
    \draw[edge](D) -- (\x-1,\y-0.6)node [midway,left] {\scalebox{0.7}{\dimcolor{$d_4$}}};
    \node[tensor, path picture={
    \fill[fill= blue!20!white](path picture bounding box.south west)
    -- 
    (path picture bounding box.north east) --
    (path picture bounding box.north west) -- cycle;}, shift={(\x,\y)}] (E) {\scalebox{0.85}{$\At_5$}};
    \draw[edge](E) -- (\x,\y-0.6)node [midway,left] {\scalebox{0.7}{\dimcolor{$d_5$}}};
    \draw[edge](A) -- (B) node [midway,above] {\scalebox{0.7}{\dimcolor{$R_1$}}};
    \draw[edge](B) -- (C) node [midway,above] {\scalebox{0.7}{\dimcolor{$R_2$}}};
    \draw[edge](C) -- (D) node [midway,above] {\scalebox{0.7}{\dimcolor{$R_3$}}};
    \draw[edge](D) -- (E) node [midway,above] {\scalebox{0.7}{\dimcolor{$R_4$}}};
\end{tikzpicture}~~~\text{step:~2}\\
\begin{tikzpicture}[baseline=-0.5ex]
    \tikzset{tensor/.style = {minimum size = 0.5cm,shape = circle,thick,draw=black,inner sep = 0pt}, edge/.style = {thick,line width=.4mm},every loop/.style={}}
    \def\x{0}
    \def\y{0}
    \node[tensor, path picture = {
        \fill[fill=red!30!white] (path picture bounding box.south east)
        -- 
        (path picture bounding box.north west) --
        (path picture bounding box.north east) -- cycle;
        }, shift = {(\x-4,\y)}] (A) {\scalebox{0.85}{$\At_1$}};
        \draw[edge](A) -- (\x-4,\y-0.6) node [midway,left] {\scalebox{0.8}{\dimcolor{$d_1$}}};
       \node[tensor, path picture={
        \fill[fill=red!50!white](path picture bounding box.south east)
        -- 
        (path picture bounding box.north west) --
        (path picture bounding box.north east) -- cycle;
        }, shift={(\x-3,\y)}] (B) {\scalebox{0.85}{$\At_2$}};
       \draw[edge](B) -- (\x-3,\y-0.6)node [midway,left] {\scalebox{0.7}{\dimcolor{$d_2$}}};
        \draw  node[fill = black,circle,inner sep=0pt,minimum size=6pt] (m) at (\x-2.5,\y) {};
        \node[tensor, path picture={
        \fill[fill=blue!50!white](path picture bounding box.south west)
        -- 
        (path picture bounding box.north east) --
        (path picture bounding box.north west) -- cycle;
        }, shift={(\x-2,\y)}] (C) {\scalebox{0.85}{$\At_3$}};
        \draw[edge](C) -- (\x-2,\y-0.6)node [midway,left] {\scalebox{0.7}{\dimcolor{$d_3$}}};
        \node[tensor, path picture={
       \fill[fill=blue!60!red!50!](path             picture bounding box.south west)
        -- 
        (path picture bounding box.north east) --
       (path picture bounding box.north west) -- cycle;}, shift={(\x-1,\y)}] (D) {\scalebox{0.85}{$\At_4$}};
    \draw[edge](D) -- (\x-1,\y-0.6)node [midway,left] {\scalebox{0.7}{\dimcolor{$d_4$}}};
    \node[tensor, path picture={
    \fill[fill= blue!20!white](path             picture bounding box.south west)
    -- 
    (path picture bounding box.north east) --
    (path picture bounding box.north west) -- cycle;}, shift={(\x,\y)}] (E) {\scalebox{0.85}{$\At_5$}};
    \draw[edge](E) -- (\x,\y-0.6)node [midway,left] {\scalebox{0.7}{\dimcolor{$d_5$}}};
    \draw[edge](A) -- (B) node [midway,above] {\scalebox{0.7}{\textcolor{gray}{$R_1$}}};
    \draw[edge](B) -- (C) node [midway,above] {\scalebox{0.7}{\textcolor{gray}{$R_2$}}};
    \draw[edge](C) -- (D) node [midway,above] {\scalebox{0.7}{\textcolor{gray}{$R_3$}}};
    \draw[edge](D) -- (E) node [midway,above] {\scalebox{0.7}{\textcolor{gray}{$R_4$}}};
\end{tikzpicture}~~~~~~\text{QR}\\
\begin{tikzpicture}[baseline=-0.5ex]
    \tikzset{tensor/.style = {minimum size = 0.5cm,shape = circle,thick,draw=black,inner sep = 0pt}, edge/.style = {thick,line width=.4mm},every loop/.style={}}
    \def\x{0}
    \def\y{0}
 \node[tensor, path picture={
        \fill[fill=red!30!white](path picture bounding box.south east)
        -- 
        (path picture bounding box.north west) --
        (path picture bounding box.north east) -- cycle;
        }, shift={(\x,\y)}] (A) {\scalebox{0.85}{$\At_1$}};
        \draw[edge](A) -- (\x,\y-0.7) node [midway,left] {\scalebox{0.7}{\dimcolor{$d_1$}}};
       \node[tensor, path picture={
        \fill[fill=red!50!white](path picture bounding box.south east)
        -- 
        (path picture bounding box.north west) --
        (path picture bounding box.north east) -- cycle;
        }, shift={(\x+1,\y)}] (B) {\scalebox{0.85}{$\At_2$}};
         \draw[edge](B) -- (\x+1,\y-0.7)node [midway,left] {\scalebox{0.7}{\dimcolor{$d_2$}}};
         \node[tensor,fill=blue!50!white] (C) at (\x+2,\y){$\At_3$};
          \draw[edge](C) -- (\x+2,\y-0.7)node [midway,left] {\scalebox{0.7}{\dimcolor{$d_3$}}};
          \node[tensor, path picture={
        \fill[fill=blue!60!red!50!](path picture bounding box.south west)
        -- 
        (path picture bounding box.north east) --
        (path picture bounding box.north west) -- cycle;
        }, shift={(\x+3,\y)}] (D) {\scalebox{0.85}{$\At_4$}};
         \draw[edge](D) -- (\x+3,\y-0.7)node [midway,left] {\scalebox{0.7}{\dimcolor{$d_4$}}};
        \node[tensor, path picture={
        \fill[fill=blue!20!white](path picture bounding box.south west)
        -- 
        (path picture bounding box.north east) --
        (path picture bounding box.north west) -- cycle;
        }, shift={(\x+4,\y)}] (E) {\scalebox{0.85}{$\At_5$}};
         \draw[edge](E) -- (\x+4,\y-0.7)node [midway,left] {\scalebox{0.7}{\dimcolor{$d_5$}}};
    \draw[edge](A) -- (B) node [midway,above] {\scalebox{0.7}{\dimcolor{$R_1$}}};
    \draw[edge](B) -- (C) node [midway,above] {\scalebox{0.7}{\dimcolor{$R_2$}}};
    \draw[edge](C) -- (D) node [midway,above] {\scalebox{0.7}{\dimcolor{$R_3$}}};
    \draw[edge](D) -- (E) node [midway,above] {\scalebox{0.7}{\dimcolor{$R_4$}}};
\end{tikzpicture}~~~\text{step:~3}
\caption{TT-ALS 的一次半扫描}\label{fig:ortho-tt-als-alg}
\end{figure}
图~\ref{fig:ortho-tt-als-alg} 给出了 TT-ALS 的一次半扫描。在每一个非 QR 步骤中，完全着色的核心张量被优化；在每一个 QR 步骤中，非正交分量（用黑色圆圈表示）~被吸收进下一个核心张量。这一过程不断重复直到抵达分解的右端，然后再从右端重复同样的过程直到抵达左端（图中未展示）。整个过程反复进行直至收敛。 
\begin{remark}
    除重构误差之外，该算法还可配合其他损失函数使用。此时，损失关于某个核心张量的优化未必归结为最小二乘问题，因而不再称为 ALS，而只是 \emph{交替最小化}。该算法还有一个双站点变体：在每次迭代中，把相邻的两个核心张量合并后联合优化，再用截断奇异值分解将其重新一分为二~\citep{holtz2012alternating}。
\end{remark}

在下一节中，我们先引入 TT 矩阵，也称矩阵乘积算符（Matrix Product Operator, MPO），然后演示如何在 TT 格式下高效地完成各类运算。
\subsection{TT 格式下的高效运算}\label{sec:efficient-operations}
张量列（TT）格式的一个关键优势在于，它不仅给出高维张量的低秩表示，还使得关键代数运算能够高效执行。
本节给出可以在 TT 格式下高效执行的主要运算。我们先说明如何计算两个张量列（TT）张量之间的内积，接着引入大矩阵的矩阵乘积算符（MPO）表示，并讨论 TT 格式下的矩阵-向量乘积~(mat–vec)。然后讨论 TT 舍入，这是控制 TT 秩增长的关键步骤。最后概述 TT 框架中其他可高效执行的运算及其计算复杂度。
\paragraph{内积}
如第~\ref{subsec:contraction} 节所述，无论对张量还是对向量，内积都由连接所有对应指标得到。设 TT 格式下有两个四阶张量 $\Tt=\TT(({\Gt_n})_{n=1}^N)$、$\widetilde{\Tt}=\TT(({\Tilde{\Gt_n}})_{n=1}^N)\in\RR^{d_1\times...\times d_4}$，则该内积可以高效地算出，例如，
\begin{align}\label{eq:tt-inner}
\langle\Tt,\widetilde{\Tt}\rangle 
&=
    \begin{tikzpicture}[baseline=\baseline]
    \tikzset{tensor/.style = {minimum size = 0.5cm,shape = circle,thick,draw=black,inner sep = 0pt}, edge/.style   = {thick,line width=.4mm},every loop/.style={}}
    \def\y{0}
    \def\x{0}
  \node[tensor,fill=red!30!white] (G1) at (\x,\y+0.5){$\Gt_1$};
  \node[tensor,fill=red!50!white] (G2) at (\x+1,\y+0.5){$\Gt_2$};
  \node[tensor,fill=blue!50!white] (G3) at (\x+2,\y+0.5){$\Gt_3$};
   \node[tensor,fill=blue!20!white] (G4) at (\x+3,\y+0.5){$\Gt_4$};
\node[tensor,fill=red!30!white] (Gtilde1) at (\x,\y-0.5){$\Tilde{\Gt_1}$};
  \node[tensor,fill=red!50!white] (Gtilde2) at (\x+1,\y-0.5){$\Tilde{\Gt_2}$};
  \node[tensor,fill=blue!50!white] (Gtilde3) at (\x+2,\y-0.5){$\Tilde{\Gt_3}$};
   \node[tensor,fill=blue!20!white] (Gtilde4) at (\x+3,\y-0.5){$\Tilde{\Gt_4}$};
    \draw[edge](G1) -- (Gtilde1)node [midway,left] {\edgelabel{$d_1$}};
    \draw[edge](G2) -- (Gtilde2)node [midway,left] {\edgelabel{$d_2$}};
    \draw[edge](G3) -- (Gtilde3)node [midway,left] {\edgelabel{$d_4$}};
     \draw[edge](G4) -- (Gtilde4)node [midway,left] {\edgelabel{$d_1$}};
    \draw[edge](G1) -- (G2) node [midway,above] {\edgelabel{$R$}};
    \draw[edge](G2) -- (G3) node [midway,above] {\edgelabel{$R$}};
    \draw[edge](G3) -- (G4) node [midway,above] {\edgelabel{$R$}};
    \draw[edge](Gtilde1) -- (Gtilde2) node [midway,below] {\edgelabel{$R$}};
    \draw[edge](Gtilde2) -- (Gtilde3) node [midway,below] {\edgelabel{$R$}};
    \draw[edge](Gtilde3) -- (Gtilde4) node [midway,below] {\edgelabel{$R$}};
\end{tikzpicture}\nonumber\\
&=
\begin{tikzpicture}[baseline=\baseline]
    \tikzset{tensor/.style = {minimum size = 0.5cm,shape = circle,thick,draw=black,inner sep = 0pt}, edge/.style   = {thick,line width=.4mm},every loop/.style={}}
    \def\y{0}
    \def\x{0}
      \node[tensor,fill=red!30!white] (At_1) at (\x+4.2,\y){$\At_1$};
  \node[tensor,fill=red!50!white] (At_2) at (\x+5.5,\y){$\At_2$};
  \node[tensor,fill=blue!50!white] (At_3) at (\x+6.8,\y){$\At_3$};
   \node[tensor,fill=blue!20!white] (At_4) at (\x+8,\y){$\At_4$};
    \draw[edge](At_1) -- (At_2) node [midway,above] {\edgelabel{$R$}};
    \draw[edge](\x+4.45,\y-0.1) --  (\x+5.25,\y-0.1) node [midway,below] {\edgelabel{$R$}};
    \draw[edge](At_2) -- (At_3) node [midway,above] {\edgelabel{$R$}};
    \draw[edge](\x+5.75,\y-0.1) --  (\x+6.55,\y-0.1) node [midway,below] {\edgelabel{$R$}};
    \draw[edge](At_3) -- (At_4) node [midway,above] {\edgelabel{$R$}};
    \draw[edge](\x+7.05,\y-0.1) --  (\x+7.75,\y-0.1) node [midway,below] {\edgelabel{$R$}};
    \node[draw=none] at (\x+8.6,\y) {$=$};
    \node[tensor,fill=red!30!white] (At_1) at (\x+9.2,\y){$\At_1$};
  \node[tensor,fill=red!50!white] (At_2) at (\x+10.2,\y){$\At_2$};
  \node[tensor,fill=blue!50!white] (At_3) at (\x+11.2,\y){$\At_3$};
   \node[tensor,fill=blue!20!white] (At_4) at (\x+12.2,\y){$\At_4$};
    \draw[edge](At_1) -- (At_2) node [midway,above] {\edgelabel{$R^2$}};
    \draw[edge](At_2) -- (At_3) node [midway,above] {\edgelabel{$R^2$}};
    \draw[edge](At_3) -- (At_4) node [midway,above] {\edgelabel{$R^2$}};
    \end{tikzpicture},
\end{align}
其中，为简单起见，我们假设所有秩都等于 $R$。
按式~\eqref{eq:tt-inner} 所描述的过程计算 TT 格式下两个 $N$ 阶张量的内积，其复杂度为 $\bigo(NdR^4)$，这里假设 $d_1=\cdots=d_N=d$ 且 $R_1=\cdots=R_{N-1} = R$。与复杂度为 $\bigo(d^N)$ 的标准方法相比，这是一个显著的改进。因此，TT 格式为在高维且 TT 秩较低的张量上进行运算提供了高效而实用的途径。此外，核心张量还能以更高效的方式收缩：
\begin{align*}
\begin{tikzpicture}[baseline=\baseline]
    \tikzset{tensor/.style = {minimum size = 0.5cm,shape = circle,thick,draw=black,inner sep = 0pt}, edge/.style   = {thick,line width=.4mm},every loop/.style={}}
    \def\y{0}
    \def\x{0}
    \node[tensor,fill=red!30!white] (D) at (\x-1,\y){$\Gt_1$};
    \node[tensor,fill=red!50!white] (A) at (\x,\y){$\Gt_2$};
    \node[tensor,fill=blue!50!white] (B) at (\x+1,\y){$\Gt_3$};
    \node[tensor,fill=blue!20!white] (C) at (\x+2,\y){$\Gt_4$};
    \draw[edge](A) -- (\x,\y-0.7)node [midway,left] {\edgelabel{$d_2$}};
    \draw[edge](B) -- (\x+1,\y-0.7)node [midway,left] {\edgelabel{$d_3$}};
    \draw[edge](C) -- (\x+2,\y-0.7)node [midway,left] {\edgelabel{$d_4$}};
    \draw[edge](D) -- (\x-1,\y-0.7)node [midway,left] {\edgelabel{$d_1$}};
    \draw[edge](D) -- (A) node [midway,above] {\edgelabel{$R$}};
    \draw[edge](A) -- (B) node [midway,above] {\edgelabel{$R$}};
    \draw[edge](B) -- (C) node [midway,above] {\edgelabel{$R$}};
    \node[tensor,fill=red!30!white] (Dtild) at (\x-1,\y-1){$\Tilde{\Gt_1}$};
    \node[tensor,fill=red!50!white] (Atild) at (\x,\y-1){$\Tilde{\Gt_2}$};
    \node[tensor,fill=blue!50!white] (Btild) at (\x+1,\y-1){$\Tilde{\Gt_3}$};
    \node[tensor,fill=blue!20!white] (Ctild) at (\x+2,\y-1){$\Tilde{\Gt_4}$};
    \draw[edge](Dtild) -- (Atild) node [midway,below] {\edgelabel{$R$}};
    \draw[edge](Atild) -- (Btild) node [midway,below] {\edgelabel{$R$}};
    \draw[edge](Btild) -- (Ctild) node [midway,below] {\edgelabel{$R$}};
    \draw[thick, dotted,rotate=90] (\x-0.5,\y+1) ellipse (1 cm and 0.5 cm);
    \node[tensor,fill=red!30!white] (Gt_1) at (\x+3.5,\y){$\Gt_1$};
    \node[tensor,fill=red!50!white] (Gt_2) at (\x+4.5,\y){$\Gt_2$};
    \node[tensor,fill=blue!50!white] (Gt_3) at (\x+5.5,\y){$\Gt_3$};
    \node[tensor,fill=blue!20!white] (Gt_4) at (\x+6.5,\y){$\Gt_4$};
    \draw[edge](Gt_1) -- (Gt_2) node [midway,above]{\scalebox{0.7}{\dimcolor{$R$}}};
    \draw[edge](Gt_2) -- (Gt_3) node [midway,above,very near start]{\scalebox{0.7}{\dimcolor{$R$}}};
    \draw[edge](Gt_3) -- (Gt_4) node [midway,above]{\scalebox{0.7}{\dimcolor{$R$}}};
    \node[tensor,fill=red!30!white] (Gtild_1) at (\x+3.5,\y-1){$\Tilde{\Gt_1}$};
    \node[tensor,fill=red!50!white] (Gtild_2) at (\x+4.5,\y-1){$\Tilde{\Gt_2}$};
    \node[tensor,fill=blue!50!white] (Gtild_3) at (\x+5.5,\y-1){$\Tilde{\Gt_3}$};
    \node[tensor,fill=blue!20!white] (Gtild_4) at (\x+6.5,\y-1){$\Tilde{\Gt_4}$};
    \draw[edge](Gtild_1) -- (Gtild_2) node [midway,below,near end] {\scalebox{0.7}{\dimcolor{$R$}}};
    \draw[edge](Gtild_2) -- (Gtild_3) node [midway,below] {\scalebox{0.7}{\dimcolor{$R$}}};
    \draw[edge](Gtild_3) -- (Gtild_4) node [midway,below]{\scalebox{0.7}{\dimcolor{$R$}}};
    \draw[edge](Gt_1) -- (Gtild_1) node [midway,left] {\scalebox{0.7}{\dimcolor{$d_1$}}};
    \draw[edge](Gt_2) -- (Gtild_2) node [midway,left]{\scalebox{0.7}{\dimcolor{$d_2$}}};
    \draw[edge](Gt_3) -- (Gtild_3) node [midway,left]{\scalebox{0.7}{\dimcolor{$d_3$}}};
    \draw[edge](Gt_4) -- (Gtild_4) node [midway,left]{\scalebox{0.7}{\dimcolor{$d_4$}}};
    \node[draw=none] at (\x+2.6,\y-0.5){$\rightarrow$};
    \draw[dash pattern=on 1pt off 2pt] plot[smooth cycle] coordinates { 
    (\x+3,\y+0.35) 
    (\x+3.9, \y+0.35) (\x+4.9,\y+0.35) (\x+4.9, \y-0.35)  (\x+3.9,\y-0.35) (\x+3.9,\y-1.4) (\x+3, \y-1.4) };
    \node[tensor,fill=red!30!white] (Gt_1) at (\x+8,\y){$\Gt_1$};
    \node[tensor,fill=red!50!white] (Gt_2) at (\x+9,\y){$\Gt_2$};
    \node[tensor,fill=blue!50!white] (Gt_3) at (\x+10,\y){$\Gt_3$};
    \node[tensor,fill=blue!20!white] (Gt_4) at (\x+11,\y){$\Gt_4$};
    \draw[edge](Gt_1) -- (Gt_2) node [midway,above]{\scalebox{0.7}{\dimcolor{$R$}}};
    \draw[edge](Gt_2) -- (Gt_3) node [midway,above]{\scalebox{0.7}{\dimcolor{$R$}}};
    \draw[edge](Gt_3) -- (Gt_4) node [midway,above]{\scalebox{0.7}{\dimcolor{$R$}}};
    \node[tensor,fill=red!30!white] (Gtild_1) at (\x+8,\y-1){$\Tilde{\Gt_1}$};
    \node[tensor,fill=red!50!white] (Gtild_2) at (\x+9,\y-1){$\Tilde{\Gt_2}$};
    \node[tensor,fill=blue!50!white] (Gtild_3) at (\x+10,\y-1){$\Tilde{\Gt_3}$};
    \node[tensor,fill=blue!20!white] (Gtild_4) at (\x+11,\y-1){$\Tilde{\Gt_4}$};
    \draw[edge](Gtild_1) -- (Gtild_2) node [midway,below,] {\scalebox{0.7}{\dimcolor{$R$}}};
    \draw[edge](Gtild_2) -- (Gtild_3) node [midway,below] {\scalebox{0.7}{\dimcolor{$R$}}};
    \draw[edge](Gtild_3) -- (Gtild_4) node [midway,below]{\scalebox{0.7}{\dimcolor{$R$}}};
    \draw[edge](Gt_1) -- (Gtild_1) node [midway,left] {\scalebox{0.7}{\dimcolor{$d_1$}}};
    \draw[edge](Gt_2) -- (Gtild_2) node [midway,left]{\scalebox{0.7}{\dimcolor{$d_2$}}};
    \draw[edge](Gt_3) -- (Gtild_3) node [midway,left]{\scalebox{0.6}{\dimcolor{$d_3$}}};
    \draw[edge](Gt_4) -- (Gtild_4) node [midway,left]{\scalebox{0.7}{\dimcolor{$d_4$}}};
    \node[draw=none] at (\x+7.2,\y-0.5){$\rightarrow$};
    \draw[dash pattern=on 1pt off 2pt] plot[smooth cycle] coordinates { 
    (\x+7.7,\y+0.3) 
    (\x+9.25, \y+0.35) (\x+9.25,\y-1.4) (\x+7.6, \y-1.4)};
    \node[draw=none] at (\x+11.7,\y-0.5){$\rightarrow$};
\end{tikzpicture}
\end{align*}
\begin{align*}
    \begin{tikzpicture}[baseline=\baseline]
    \tikzset{tensor/.style = {minimum size = 0.5cm,shape = circle,thick,draw=black,inner sep = 0pt}, edge/.style   = {thick,line width=.4mm},every loop/.style={}}
    \def\y{0}
    \def\x{0}
    \node[tensor,fill=red!30!white] (Gt_1) at (\x,\y){$\Gt_1$};
    \node[tensor,fill=red!50!white] (Gt_2) at (\x+1,\y){$\Gt_2$};
    \node[tensor,fill=blue!50!white] (Gt_3) at (\x+2,\y){$\Gt_3$};
    \node[tensor,fill=blue!20!white] (Gt_4) at (\x+3,\y){$\Gt_4$};
    \node[tensor,fill=red!30!white] (Gtild_1) at (\x,\y-1){$\Tilde{\Gt_1}$};
    \node[tensor,fill=red!50!white] (Gtild_2) at (\x+1,\y-1){$\Tilde{\Gt_2}$};
    \node[tensor,fill=blue!50!white] (Gtild_3) at (\x+2,\y-1){$\Tilde{\Gt_3}$};
    \node[tensor,fill=blue!20!white] (Gtild_4) at (\x+3,\y-1){$\Tilde{\Gt_4}$};
    \draw[edge](Gt_1) -- (Gtild_1)node [midway,left] {\edgelabel{$d_1$}};
    \draw[edge](Gt_2) -- (Gtild_2)node [midway,left] {\edgelabel{$d_2$}};
    \draw[edge](Gt_3) -- (Gtild_3)node [midway,left] {\edgelabel{$d_3$}};
    \draw[edge](Gt_4) -- (Gtild_4)node [midway,left] {\edgelabel{$d_4$}};
    \draw[edge](Gt_1) -- (Gt_2) node [midway,above] {\edgelabel{$R$}};
    \draw[edge](Gt_2) -- (Gt_3) node [midway,above] {\edgelabel{$R$}};
    \draw[edge](Gt_3) -- (Gt_4) node [midway,above,near end] {\scalebox{0.7}{\dimcolor{$R$}}};
    \draw[edge](Gtild_1) -- (Gtild_2) node [midway,below,near end] {\scalebox{0.7}{\dimcolor{$R$}}};
    \draw[edge](Gtild_2) -- (Gtild_3) node [midway,below] {\edgelabel{$R$}};
    \draw[edge](Gtild_3) -- (Gtild_4) node [midway,below] {\edgelabel{$R$}};
    \draw[dash pattern=on 1pt off 2pt] plot[smooth cycle] coordinates { 
    (\x-0.5,\y+0.5) (\x,\y+0.5)
    (\x+2.3, \y+0.5) (\x+2.3,\y-0.4) (\x+0.5, \y-0.4)  (\x+0.5, \y-1.34) (\x-0.5, \y-1.34) (\x-0.5, \y+0.1)};
    \node[draw=none] at (\x+3.7,\y-0.5){$\rightarrow$};
    \node[draw=none] at (\x+4.5,\y-0.5){$\cdots$};
    \node[draw=none] at (\x+5.5,\y-0.5){$\rightarrow$};
\end{tikzpicture}
\begin{tikzpicture}[baseline=\baseline]
    \tikzset{tensor/.style = {minimum size = 0.5cm,shape = circle,thick,draw=black,inner sep = 0pt}, edge/.style   = {thick,line width=.4mm},every loop/.style={}}
    \def\y{0}
    \def\x{0}
    \node[tensor,fill=red!30!white] (Gt_1) at (\x+6,\y){$\Gt_1$};
    \node[tensor,fill=red!50!white] (Gt_2) at (\x+7,\y){$\Gt_2$};
    \node[tensor,fill=blue!50!white] (Gt_3) at (\x+8,\y){$\Gt_3$};
    \node[tensor,fill=blue!20!white] (Gt_4) at (\x+9,\y){$\Gt_4$};
    \node[tensor,fill=red!30!white] (Gtild_1) at (\x+6,\y-1){$\Tilde{\Gt_1}$};
    \node[tensor,fill=red!50!white] (Gtild_2) at (\x+7,\y-1){$\Tilde{\Gt_2}$};
    \node[tensor,fill=blue!50!white] (Gtild_3) at (\x+8,\y-1){$\Tilde{\Gt_3}$};
    \node[tensor,fill=blue!20!white] (Gtild_4) at (\x+9,\y-1){$\Tilde{\Gt_4}$};
    \draw[edge](Gt_1) -- (Gtild_1)node [midway,left] {\edgelabel{$d_1$}};
    \draw[edge](Gt_2) -- (Gtild_2)node [midway,left] {\edgelabel{$d_2$}};
    \draw[edge](Gt_3) -- (Gtild_3)node [midway,left] {\edgelabel{$d_3$}};
    \draw[edge](Gt_4) -- (Gtild_4)node [midway,left] {\edgelabel{$d_4$}};
    \draw[edge](Gt_1) -- (Gt_2) node [midway,above] {\edgelabel{$R$}};
    \draw[edge](Gt_2) -- (Gt_3) node [midway,above] {\edgelabel{$R$}};
    \draw[edge](Gt_3) -- (Gt_4) node [midway,above] {\scalebox{0.65}{\textcolor{gray}{$R$}}};
    \draw[edge](Gtild_1) -- (Gtild_2) node [midway,below] {\scalebox{0.65}{\textcolor{gray}{$R$}}};
    \draw[edge](Gtild_2) -- (Gtild_3) node [midway,below] {\edgelabel{$R$}};
    \draw[edge](Gtild_3) -- (Gtild_4) node [midway,below,near end] {\edgelabel{$R$}};
    \draw[dash pattern=on 1pt off 2pt] plot[smooth cycle] coordinates { 
    (\x+5.5,\y+0.5) (\x+6,\y+0.5)
    (\x+9.5, \y+0.5) (\x+9.5,\y-0.4) (\x+8.3, \y-0.4)  (\x+8.3, \y-1.34) (\x+5.5, \y-1.34) (\x+5.5, \y+0.1)};
\end{tikzpicture}
\end{align*}
由上可见，计算两个长度为 $\bigo(d^N)$ 的向量的内积，其复杂度可降至 $\bigo(NdR^3)$。一般而言，为任意张量网络确定最优收缩顺序是一个 NP 难问题~\citep{chi1997optimizing}。
\paragraph{矩阵乘积算符~(MPO)}\label{par:mpo}
作为 TT 分解的推广，我们引入矩阵乘积算符（MPO）~\citep{oseledets2010approximation}。MPO 又称 TT 矩阵，是一条由四阶张量组成的链，用来表示一个矩阵。它最初是为描述作用于多体量子系统的算符而发展起来的。简单地说，MPO 就是一种用张量表示矩阵的方法。设有矩阵 $\Ab\in\RR^{I_1I_2\dots I_N\times J_1J_2\dots J_N}$。对 $n\in[N]$，令  $\At_n\in\RR^{R_{n-1}\times I_n\times J_n\times R_n}$ ，其中 $R_0=R_N=1$。$\Ab$ 的秩为 $(R_1,\cdots,R_N)$ 的 MPO 分解定义为 
$$\Ab_{i_1i_2\dotsi_N,j_1j_2\dots j_N} = (\At_1)_{i_1,j_1,:}(\At_2)_{:,i_2,j_2,:}\dots(\At_{N-1})_{:,i_{N-1},j_{N-1},:}(\At_N)_{:,i_N,j_N},$$ 
对所有指标 $i_1\in[I_1],\cdots,i_N\in[I_N]$ 与 $j_1\in[J_1],\dots, j_N\in[J_N]$ 成立。
下面用张量网络记法给出矩阵 $\Ab\in\RR^{I_1I_2I_3\times J_1J_2J_3}$ 的一个秩~$(R_1,R_2)$ MPO 分解的例子：
$$
\begin{tikzpicture}[baseline=-0.5ex]
    \tikzset{tensor/.style = {minimum size = 0.5cm,shape = circle,thick,draw=black,inner sep = 0pt}, edge/.style = {thick,line width=.4mm},every loop/.style={}}
    \def\x{6}
    \def\y{0}
     \node[tensor,fill=green!50!red!50!white] (B) at (\x+1,\y) {$\Ab$};
    \draw[edge] (\x+1.15,\y+0.2) -- (\x+1.5,\y+0.2) -- (\x+1.5,\y+0.1) -- (\x+2,\y+0.1) node [midway,above]
    {\scalebox{0.7}{\dimcolor{$J_3$}}};
     \draw[edge] (B) -- (\x+2,\y)node [right]
    {\scalebox{0.7}{\dimcolor{$J_2$}}};
    \draw[edge] (\x+1.15,\y-0.2) -- (\x+1.5,\y-0.2)-- (\x+1.5,\y-0.1)-- (\x+2,\y-0.1) node [midway,below]
    {\scalebox{0.7}{\dimcolor{$J_1$}}};
    
     \draw[edge] (\x+0.85,\y+0.2) -- (\x+0.5,\y+0.2) -- (\x+0.5,\y+0.1) -- (\x,\y+0.1) node [midway,above]
    {\scalebox{0.7}{\dimcolor{$I_3$}}};
     \draw[edge] (B) -- (\x,\y)node [left]
    {\scalebox{0.7}{\dimcolor{$I_2$}}};

    \draw[edge] (\x+0.85,\y-0.2) -- (\x+0.5,\y-0.2)-- (\x+0.5,\y-0.1)-- (\x,\y-0.1) node [midway,below]
    {\scalebox{0.7}{\dimcolor{$I_1$}}};
\end{tikzpicture}
=
\begin{tikzpicture}[baseline=-0.5ex]
    \tikzset{tensor/.style = {minimum size = 0.5cm,shape = circle,thick,draw=black,inner sep = 0pt}, edge/.style = {thick,line width=.4mm},every loop/.style={}}
    \def\x{6}
    \def\y{0}
     \node[tensor,fill=green!50!red!50!white] (B) at (\x+1,\y) {$\At$};
    \draw[edge] (B) -- (\x-0.15,\y) node [midway,above=-0.07cm] {\scalebox{0.7}{\dimcolor{$I_2$}}};
    \draw[edge] (B) -- (\x,\y-0.5) node [midway,below] {\scalebox{0.7}{\dimcolor{$I_1$}}};
    \draw[edge] (B) -- (\x,\y+0.5) node [midway,above] {\scalebox{0.7}{\dimcolor{$I_3$}}};
    \draw[edge] (B) -- (\x+2.15,\y) node [midway,above=-0.07] {\scalebox{0.7}{\dimcolor{$J_2$}}};
    \draw[edge] (B) -- (\x+2,\y-0.5) node [midway,below] {\scalebox{0.7}{\dimcolor{$J_1$}}};
    \draw[edge] (B) -- (\x+2,\y+0.5) node [midway,above] {\scalebox{0.7}{\dimcolor{$J_3$}}};
\end{tikzpicture}
=
\begin{tikzpicture}[baseline=\baseline]
    \tikzset{tensor/.style = {minimum size = 0.5cm,shape = circle,thick,draw=black,inner sep = 0pt}, edge/.style   = {thick,line width=.4mm},every loop/.style={}}
    \def\y{0}
    \def\x{0}
  \node[tensor,fill=red!30!white] (D) at (\x-1.5,\y){$\At_1$};
  \node[tensor,fill=red!50!white] (A) at (\x,\y){$\At_2$};
  \node[tensor,fill=blue!50!white] (B) at (\x+1.5,\y){$\At_3$};
    \draw[edge](A) -- (\x,\y-0.75)node [midway, right] {\edgelabel{$I_2$}};
    \draw[edge](B) -- (\x+1.5,\y-0.75)node [midway,right] {\edgelabel{$I_3$}};
     \draw[edge](D) -- (\x-1.5,\y-0.75)node [midway,right] {\edgelabel{$I_1$}};
    \draw[edge](D) -- (A) node [midway,above] {\edgelabel{$R_1$}};
    \draw[edge](A) -- (B) node [midway,above] {\edgelabel{$R_2$}};
    \draw[edge](A) -- (\x,\y+0.75)node [midway, right] {\edgelabel{$J_2$}};
    \draw[edge](B) -- (\x+1.5,\y+0.75) node [midway,right] {\edgelabel{$J_3$}};
    \draw[edge](D) -- (\x-1.5,\y+0.75)node [midway,right] {\edgelabel{$J_1$}};
\end{tikzpicture}
=
\begin{tikzpicture}[baseline=\baseline]
    \tikzset{tensor/.style = {minimum size = 0.5cm,shape = circle,thick,draw=black,inner sep = 0pt}, edge/.style   = {thick,line width=.4mm},every loop/.style={}}
    \def\y{0}
    \def\x{0}
  \node[tensor,fill=red!30!white] (D) at (\x-1.5,\y){$\At_1$};
  \node[tensor,fill=red!50!white] (A) at (\x,\y){$\At_2$};
  \node[tensor,fill=blue!50!white] (B) at (\x+1.5,\y){$\At_3$};
    \draw[edge](A) -- (\x,\y-1.25)node [near start=-0.05, right=-0.07cm] {\edgelabel{$I_2$}};
    \draw[edge](B) -- (\x+1.5,\y-0.75)-- (\x+0.15,\y-0.75)node [midway,below] {\edgelabel{$I_3$}}--(\x+0.15,\y-1.25);
     \draw[edge](D) -- (\x-1.5,\y-0.75)-- (\x-0.15,\y-0.75)node [midway,below] {\edgelabel{$I_1$}}--(\x-0.15,\y-1.25);
    \draw[edge](D) -- (A) node [midway,above] {\edgelabel{$R_1$}};
    \draw[edge](A) -- (B) node [midway,above] {\edgelabel{$R_2$}};
    \draw[edge](A) -- (\x,\y+1.25)node [near start=-0.05, right=-0.07cm] {\edgelabel{$J_2$}};
    \draw[edge](B) -- (\x+1.5,\y+0.75)-- (\x+0.15,\y+0.75)node [midway,above] {\edgelabel{$J_3$}}--(\x+0.15,\y+1.25);
    \draw[edge](D) -- (\x-1.5,\y+0.75)-- (\x-0.15,\y+0.75)node [midway,above] {\edgelabel{$J_1$}}--(\x-0.15,\y+1.25);
\end{tikzpicture}.
$$
\paragraph{MPO 中的矩阵-向量乘积} 矩阵 

$\Ab\in\RR^{I_1I_2I_3\times J_1J_2J_3}$
（以 TT 矩阵形式给出）与向量 $\vb\in\RR^{I_1I_2I_3}$
（以 TT 格式给出）的乘积可以高效计算：
\begin{align}\label{eq:tt-mat-vec}
\begin{tikzpicture}[baseline=-0.5ex]
    \tikzset{tensor/.style = {minimum size = 0.5cm,shape = circle,thick,draw=black,inner sep = 0pt}, edge/.style = {thick,line width=.4mm},every loop/.style={}}
    \def\x{0}
    \def\y{0}
     \node[tensor,fill=green!50!red!50!white] (B) at (\x+1,\y) {$\Ab$};
    \node[tensor,fill=blue!50!red!50!white] (a) at (\x-0.25,\y) {$\vb$};
    \draw[edge] (B) -- (\x,\y) node [midway,above] {\scalebox{0.7}{\dimcolor{$I_1I_2I_3$}}};
    \draw[edge] (B) -- (\x+2,\y) node [midway,above] {\scalebox{0.7}{\dimcolor{$J_1J_2J_3$}}};
\end{tikzpicture}
=
\begin{tikzpicture}[baseline=\baseline]
    \tikzset{tensor/.style = {minimum size = 0.5cm,shape = circle,thick,draw=black,inner sep = 0pt}, edge/.style   = {thick,line width=.4mm},every loop/.style={}}
    \def\y{0}
    \def\x{0}
    \node[tensor,fill=red!30!white] (D) at (\x-1.5,\y){$\At_1$};
    \node[tensor,fill=red!50!white] (A) at (\x,\y){$\At_2$};
  \node[tensor,fill=blue!50!white] (B) at (\x+1.5,\y){$\At_3$};
    \draw[edge](A) -- (\x,\y-0.7)node [midway,left] {\edgelabel{$I_2$}};
    \draw[edge](B) -- (\x+1.5,\y-0.7)node [midway,left] {\edgelabel{$I_3$}};
     \draw[edge](D) -- (\x-1.5,\y-0.7)node [midway,left] {\edgelabel{$I_1$}};
    \draw[edge](D) -- (A) node [midway,above] {\edgelabel{$R_1$}};
    \draw[edge](A) -- (B) node [midway,above] {\edgelabel{$R_2$}};
    \draw[edge](A) -- (\x,\y+0.7)node [midway,left] {\edgelabel{$J_2$}};
    \draw[edge](B) -- (\x+1.5,\y+0.7)node [midway,left] {\edgelabel{$J_3$}};
    \draw[edge](D) -- (\x-1.5,\y+0.7)node [midway,left] {\edgelabel{$J_1$}};
    \node[tensor,fill=red!70!blue!40!white] (G1) at (\x-1.5,\y-0.9){$\Gt_1$};
    \node[tensor,fill=red!30!blue!30!white] (G2) at (\x,\y-0.9){$\Gt_2$};
  \node[tensor,fill=red!50!blue!50!white] (G3) at (\x+1.5,\y-0.9){$\Gt_3$};
  \draw[edge] (G1) -- (G2) node [midway,above] {\edgelabel{$S_1$}};
    \draw[edge] (G2) -- (G3) node [midway,above] {\edgelabel{$S_2$}};
\end{tikzpicture}
=~
\begin{tikzpicture}[baseline=\baseline]
    \tikzset{tensor/.style = {minimum size = 0.5cm,shape = circle,thick,draw=black,inner sep = 0pt}, edge/.style   = {thick,line width=.4mm},every loop/.style={}}
    \def\y{0}
    \def\x{0}
    \node[tensor,fill=cyan] (H1) at (\x-1.5,\y){$\Ht_1$};
    \node[tensor,fill=cyan!50!] (H2) at (\x,\y){$\Ht_2$};
   \node[tensor,fill=cyan!20!] (H3) at (\x+1.5,\y){$\Ht_3$};
   \draw[edge] (H1) -- (H2) node [midway,above]{\edgelabel{$R_1S_1$}};
    \draw[edge] (H2) -- (H3) node [midway,above]{\edgelabel{$R_2S_2$}};
    \draw[edge](H1) -- (\x-1.5,\y-0.7) node [midway, left]{\edgelabel{$J_1$}};
    \draw[edge](H2) -- (\x,\y-0.7)node [midway, left]{\edgelabel{$J_2$}};
    \draw[edge](H3) -- (\x+1.5,\y-0.7)node [midway, left]{\edgelabel{$J_3$}};
\end{tikzpicture}.
\end{align}
最后一个等式中，每一对核心张量 $\At_i,\Gt_i$ 被合并为核心 $\Ht_i$~（$i = 1,2,3$），由此得到秩为 $(R_1 S_1, R_2 S_2)$ 的 TT 向量。更一般地，对于 $N$ 阶张量，若 $I_1 = \cdots = I_N = I$，$J_1 = \cdots = J_N = J$，$R_1 = \cdots = R_{N-1} = R$，且 $S_1 = \cdots = S_{N-1} = S$，则矩阵-向量乘积的总计算复杂度为 $\mathcal{O}(N R^2 S^2 I J)$。

上面的矩阵-向量乘积已经表明，在 TT 格式下进行运算会使 TT 秩增大。基础线性代数运算也是如此，
例如求和与张量的逐元素乘积。所得张量仍是 TT 格式，但秩通常因这些运算而增长~\citep{oseledets2010approximation}。这些运算及其对 TT 秩的影响汇总于表~\ref{tab:tt-operations}~\citep{novikov2014putting}。
\begin{table}
\centering
\begin{tabular}{lll}
\toprule
\textsc{运算} & \textsc{输出秩} & \textsc{复杂度} \\
\midrule
$\Ct = \At \cdot \text{const}$ & $R(\Ct) = R(\At)$ & $\mathcal{O}(d\,R(\At))$ \\
$\Ct = \At + \text{const}$ & $R(\Ct) \leq R(\At) + 1$ & $\mathcal{O}(nd\,R(\At)^2)$ \\
$\Ct = \At + \Bt$ & $R(\Ct) \leq R(\At) + R(\Bt)$ & $\mathcal{O}(nd\,R(\Ct)^2)$ \\
$\Ct = \At \krao \Bt$ & $R(\Ct) \leq R(\At)\,R(\Bt)$ & $\mathcal{O}(nd\,R(\At)^2\,R(\Bt)^2)$ \\
$\mathrm{sum} \At$ & -- & $\mathcal{O}(nd\,R(\At)^2)$ \\
$\norm{\At}_F$ & -- & $\mathcal{O}(nd\,R(\mathbf{A})^3)$ \\
\bottomrule
\end{tabular}
\caption{ TT 格式张量上的高效运算~（表格取自~\citep{novikov2014putting}。）}
\label{tab:tt-operations}
\end{table}
表中 $R$ 指最大的 TT 秩，即 $R(\At) = \max_{i=0,\dots,N} R_i$，$n$ 为张量的阶，$d$ 为其最大维数。

需要时，可以用 TT 舍入（TT-rounding）算法~\citep{oseledets2010approximation} 来降低秩~（以一定的近似误差为代价）。

\paragraph{TT 舍入} 如前所述，TT 格式下的张量运算，例如加法或乘法，一般会使 TT 秩增大。为了在保持精度的同时避免秩不受控制地增长，就必须进行秩缩减。这可以通过施加一个 \emph{TT 舍入}流程实现：给定一个 $N$ 阶 TT 张量，其秩为 $R^\prime_k$（$k\in[N-1]$），目标是求出一个 TT 秩 $R_k$ 更低而质量良好的近似。为此，一种朴素而低效的做法是重构整个张量，再执行 TT-SVD，得到秩为 $(R_1,\cdots,R_{N-1})$ 的 TT 近似。事实上，无需重构整个张量就能完成此事！我们将说明，下面这一流程与对整个张量执行 TT-SVD 完全等价。

首先对 TT 核心张量依次做 QR 分解，按自左向右（或自右向左）的扫描完成正交化。随后对每个核心张量施以截断，也就是用截断 SVD 按给定的秩 $R_k$ 舍弃奇异值。被截断的因子再重新分配到相邻的 TT 核心张量中，于是得到一个秩更小、且近似误差受到控制的新 TT 表示。作为示例，考虑一个以 TT 格式给出、秩为 $R^\prime_k$ 的四阶张量。第一步，我们自右向左执行 QR 分解：
\begin{align*}
&\begin{tikzpicture}[baseline=\baseline]
    \tikzset{tensor/.style = {minimum size = 0.5cm,shape = circle,thick,draw=black,inner sep = 0pt}, edge/.style   = {thick,line width=.4mm},every loop/.style={}}
    \def\y{0}
    \def\x{0}
    \node[tensor,fill=cyan] (G1) at (\x-3,\y){$\Gt_1$};
    \node[tensor,fill=cyan!60!] (G2) at (\x-2,\y){$\Gt_2$};
    \node[tensor,fill=cyan!40!] (G3) at (\x-1,\y){$\Gt_3$};
   \node[tensor,fill=cyan!20!] (G4) at (\x,\y){$\Gt_4$};
   \draw[edge] (G1) -- (G2) node [midway,above]{\edgelabel{$R^\prime_1$}};
    \draw[edge] (G2) -- (G3) node [midway,above]{\edgelabel{$R^\prime_2$}};
    \draw[edge] (G3) -- (G4) node [midway,above]{\edgelabel{$R^\prime_3$}};
    \draw[edge](G1) -- (\x-3,\y-0.7) node [midway, left]{\edgelabel{$d_1$}};
    \draw[edge](G2) -- (\x-2,\y-0.7)node [midway, left]{\edgelabel{$d_2$}};
    \draw[edge](G3) -- (\x-1,\y-0.7)node [midway, left]{\edgelabel{$d_3$}};
    \draw[edge](G4) -- (\x,\y-0.7)node [midway, left]{\edgelabel{$d_4$}};
    \node[draw=none] at (\x+3.5,\y){$\underrightarrow{\text{QR decomposition from right to left}}$};
\end{tikzpicture}
\begin{tikzpicture}[baseline=\baseline]
    \tikzset{tensor/.style = {minimum size = 0.5cm,shape = circle,thick,draw=black,inner sep = 0pt}, edge/.style = {thick,line width=.4mm},every loop/.style={}}
    \def\y{0}
    \def\x{0}
        \node[tensor,fill=cyan] (G1) at (\x-3,\y){$\Gt_1$};
    \node[tensor,fill=cyan!60!] (G2) at (\x-2,\y){$\Gt_2$};
    \node[tensor,fill=cyan!40!] (G3) at (\x-1,\y){$\Gt_3$};
    \node[tensor, path picture={
    \fill[fill=cyan!20!](path picture bounding box.south west)
    -- 
    (path picture bounding box.north east) --
    (path picture bounding box.north west) -- cycle;}, shift={(\x+1,\y)}] (G4) {$\Qt_4$};
    \node[tensor,fill=cyan!20!] (R4) at (\x,\y){$\Rb_4$};
   \draw[edge] (G1) -- (G2) node [midway,above]{\edgelabel{$R^\prime_1$}};
    \draw[edge] (G2) -- (G3) node [midway,above]{\edgelabel{$R^\prime_2$}};
    \draw[edge] (G3) -- (R4) node [midway,above]{\edgelabel{$R^\prime_3$}};
    \draw[edge] (R4) -- (G4) node [midway,above]{\edgelabel{$R^\prime_3$}};
    \draw[edge](G1) -- (\x-3,\y-0.7) node [midway, left]{\edgelabel{$d_1$}};
    \draw[edge](G2) -- (\x-2,\y-0.7)node [midway, left]{\edgelabel{$d_2$}};
    \draw[edge](G3) -- (\x-1,\y-0.7)node [midway, left]{\edgelabel{$d_3$}};
    \draw[edge](G4) -- (\x+1,\y-0.7)node [midway,left] {\scalebox{0.7}{\dimcolor{$d_4$}}};
    \node[draw=none] at (\x-0.5,\y+0.7){$\overbrace{\hspace{1cm}}^{\text{merge \& apply QR}}$};
    \end{tikzpicture}\\
&\begin{tikzpicture}[baseline=\baseline]
    \tikzset{tensor/.style = {minimum size = 0.5cm,shape = circle,thick,draw=black,inner sep = 0pt}, edge/.style = {thick,line width=.4mm},every loop/.style={}}
    \def\y{0}
    \def\x{0}
        \node[tensor,fill=cyan] (G1) at (\x-3,\y){$\Gt_1$};
    \node[tensor,fill=cyan!60!] (G2) at (\x-2,\y){$\Gt_2$};
    \node[tensor, path picture={
    \fill[fill=cyan!40!](path picture bounding box.south west)
    -- 
    (path picture bounding box.north east) --
    (path picture bounding box.north west) -- cycle;}, shift={(\x,\y)}] (G3) {$\Qt_3$};
    \node[tensor, path picture={
    \fill[fill=cyan!20!](path picture bounding box.south west)
    -- 
    (path picture bounding box.north east) --
    (path picture bounding box.north west) -- cycle;}, shift={(\x+1,\y)}] (G4) {$\Qt_4$};
    \node[tensor,fill=cyan!40!] (R3) at (\x-1,\y){$\Rb_3$};
   \draw[edge] (G1) -- (G2) node [midway,above]{\edgelabel{$R^\prime_1$}};
    \draw[edge] (G2) -- (R3) node [midway,above]{\edgelabel{$R^\prime_2$}};
    \draw[edge] (R3) -- (G3) node [midway,above]{\edgelabel{$R^\prime_3$}};
    \draw[edge] (G3) -- (G4) node [midway,above]{\edgelabel{$R^\prime_3$}};
    \draw[edge](G1) -- (\x-3,\y-0.7) node [midway, left]{\edgelabel{$d_1$}};
    \draw[edge](G2) -- (\x-2,\y-0.7)node [midway, left]{\edgelabel{$d_2$}};
    \draw[edge](G3) -- (\x,\y-0.7)node [midway, left]{\edgelabel{$d_3$}};
    \draw[edge](G4) -- (\x+1,\y-0.7)node [midway,left] {\scalebox{0.7}{\dimcolor{$d_4$}}};
    \node[draw=none] at (\x-1.5,\y+0.7){$\overbrace{\hspace{1cm}}^{\text{merge \& apply QR}}$};
    \node[draw=none] at (\x+2,\y){$\cdots$};
    \end{tikzpicture}
\begin{tikzpicture}[baseline=\baseline]
    \tikzset{tensor/.style = {minimum size = 0.5cm,shape = circle,thick,draw=black,inner sep = 0pt}, edge/.style = {thick,line width=.4mm},every loop/.style={}}
    \def\y{0}
    \def\x{0}
    \node[tensor,fill=cyan] (G1) at (\x-3,\y){$\Tilde{\Gt_1}$};
    \node[tensor, path picture={
    \fill[fill=cyan!60!](path picture bounding box.south west)
    -- 
    (path picture bounding box.north east) --
    (path picture bounding box.north west) -- cycle;}, shift={(\x-2,\y)}] (G2) {$\Qt_2$};
    \node[tensor, path picture={
    \fill[fill=cyan!40!](path picture bounding box.south west)
    -- 
    (path picture bounding box.north east) --
    (path picture bounding box.north west) -- cycle;}, shift={(\x-1,\y)}] (G3) {$\Qt_3$};
    \node[tensor, path picture={
    \fill[fill=cyan!20!](path picture bounding box.south west)
    -- 
    (path picture bounding box.north east) --
    (path picture bounding box.north west) -- cycle;}, shift={(\x,\y)}] (G4) {$\Qt_4$};
   \draw[edge] (G1) -- (G2) node [midway,above]{\edgelabel{$R^\prime_1$}};
      \draw[edge] (G2) -- (G3) node [midway,above]{\edgelabel{$R^\prime_2$}};
    \draw[edge] (G3) -- (G4) node [midway,above]{\edgelabel{$R^\prime_3$}};
    \draw[edge](G1) -- (\x-3,\y-0.7) node [midway, left]{\edgelabel{$d_1$}};
    \draw[edge](G2) -- (\x-2,\y-0.7)node [midway, left]{\edgelabel{$d_2$}};
    \draw[edge](G3) -- (\x-1,\y-0.7)node [midway, left]{\edgelabel{$d_3$}};
    \draw[edge](G4) -- (\x,\y-0.7)node [midway,left] {\scalebox{0.7}{\dimcolor{$d_4$}}};
    \end{tikzpicture}.
\end{align*}
随后自左向右继续施加截断 SVD：
\begin{align*}
&\begin{tikzpicture}[baseline=\baseline]
    \tikzset{tensor/.style = {minimum size = 0.5cm,shape = circle,thick,draw=black,inner sep = 0pt}, edge/.style = {thick,line width=.4mm},every loop/.style={}}
    \def\x{0}
    \def\y{0}
    \node[tensor,fill=cyan] (G1) at (\x-3,\y){$\Tilde{\Gt_1}$};
    \node[tensor, path picture={
    \fill[fill=cyan!60!](path picture bounding box.south west)
    -- 
    (path picture bounding box.north east) --
    (path picture bounding box.north west) -- cycle;}, shift={(\x-2,\y)}] (G2) {$\Qt_2$};
    \node[tensor, path picture={
    \fill[fill=cyan!40!](path picture bounding box.south west)
    -- 
    (path picture bounding box.north east) --
    (path picture bounding box.north west) -- cycle;}, shift={(\x-1,\y)}] (G3) {$\Qt_3$};
    \node[tensor, path picture={
    \fill[fill=cyan!20!](path picture bounding box.south west)
    -- 
    (path picture bounding box.north east) --
    (path picture bounding box.north west) -- cycle;}, shift={(\x,\y)}] (G4) {$\Qt_4$};
   \draw[edge] (G1) -- (G2) node [midway,above]{\edgelabel{$R^\prime_1$}};
      \draw[edge] (G2) -- (G3) node [midway,above]{\edgelabel{$R^\prime_2$}};
    \draw[edge] (G3) -- (G4) node [midway,above]{\edgelabel{$R^\prime_3$}};
    \draw[edge](G1) -- (\x-3,\y-0.7) node [midway, left]{\edgelabel{$d_1$}};
    \draw[edge](G2) -- (\x-2,\y-0.7)node [midway, left]{\edgelabel{$d_2$}};
    \draw[edge](G3) -- (\x-1,\y-0.7)node [midway, left]{\edgelabel{$d_3$}};
    \draw[edge](G4) -- (\x,\y-0.7)node [midway,left] {\scalebox{0.7}{\dimcolor{$d_4$}}};
    \node[draw=none] at (\x+3.25,\y){$\underrightarrow{\text{rank-$R_1$ truncated SVD of $\tilde \Gt_1$}}$};
    \end{tikzpicture}
\begin{tikzpicture}[baseline=\baseline]
    \tikzset{tensor/.style = {minimum size = 0.5cm,shape = circle,thick,draw=black,inner sep = 0pt}, edge/.style = {thick,line width=.4mm},every loop/.style={}}
    \def\y{0}
    \def\x{0}
     \node[tensor, path picture={
    \fill[fill=cyan](path picture bounding box.south east)
    -- 
    (path picture bounding box.north west) --
    (path picture bounding box.north east) -- cycle;}, shift={(\x-3,\y)}] (U1) {$\Gt_1$};
    \node[tensor,fill=cyan] (sigma) at (\x-2,\y){{}};
    \draw[edge=0.5] (sigma.south west) -- (sigma.north        east);
    \node[tensor, path picture={
    \fill[fill=cyan](path picture bounding box.south west)
    -- 
    (path picture bounding box.north east) --
    (path picture bounding box.north west) -- cycle;}, shift={(\x-1,\y)}] (V1) {$\Vb_1$};
    \node[tensor, path picture={
    \fill[fill=cyan!60!](path picture bounding box.south west)
    -- 
    (path picture bounding box.north east) --
    (path picture bounding box.north west) -- cycle;}, shift={(\x,\y)}] (G2) {$\Qt_2$};
    \node[tensor, path picture={
    \fill[fill=cyan!40!](path picture bounding box.south west)
    -- 
    (path picture bounding box.north east) --
    (path picture bounding box.north west) -- cycle;}, shift={(\x+1,\y)}] (G3) {$\Qt_3$};
    \node[tensor, path picture={
    \fill[fill=cyan!20!](path picture bounding box.south west)
    -- 
    (path picture bounding box.north east) --
    (path picture bounding box.north west) -- cycle;}, shift={(\x+2,\y)}] (G4) {$\Qt_4$};
   \draw[edge] (U1) -- (sigma) node [midway,above]{\edgelabel{$R_1$}};
      \draw[edge] (sigma) -- (V1) node [midway,above]{\edgelabel{$R_1$}};
    \draw[edge] (V1) -- (G2) node [midway,above]{\edgelabel{$R_1$}};
    \draw[edge] (G2) -- (G3) node [midway,above]{\edgelabel{$R^\prime_2$}};
    \draw[edge] (G3) -- (G4) node [midway,above]{\edgelabel{$R^\prime_3$}};
    \draw[edge](U1) -- (\x-3,\y-0.7) node [midway, left]{\edgelabel{$d_1$}};
    \draw[edge](G2) -- (\x,\y-0.7)node [midway, left]{\edgelabel{$d_2$}};
    \draw[edge](G3) -- (\x+1,\y-0.7)node [midway, left]{\edgelabel{$d_3$}};
    \draw[edge](G4) -- (\x+2,\y-0.7)node [midway,left] {\scalebox{0.7}{\dimcolor{$d_4$}}};
    \node[draw=none] at (\x-1,\y+0.7){$\overbrace{\hspace{2cm}}^{\text{merge \& apply rank-$R_2$ SVD}}$};
    \end{tikzpicture}\\
&\begin{tikzpicture}[baseline=\baseline]
    \tikzset{tensor/.style = {minimum size = 0.5cm,shape = circle,thick,draw=black,inner sep = 0pt}, edge/.style = {thick,line width=.4mm},every loop/.style={}}
    \def\y{0}
    \def\x{0}
     \node[tensor, path picture={
    \fill[fill=cyan](path picture bounding box.south east)
    -- 
    (path picture bounding box.north west) --
    (path picture bounding box.north east) -- cycle;}, shift={(\x-3,\y)}] (U1) {$\Gt_1$};
    \node[tensor, path picture={
    \fill[fill=cyan!60!](path picture bounding box.south east)
    -- 
    (path picture bounding box.north west) --
    (path picture bounding box.north east) -- cycle;}, shift={(\x-2,\y)}] (G2) {$\Gt_2$};
    \node[tensor,fill=cyan!60!] (sigma) at (\x-1,\y){{}};
    \draw[edge=0.5] (sigma.south west) -- (sigma.north        east);
    \node[tensor, path picture={
    \fill[fill=cyan!60!](path picture bounding box.south west)
    -- 
    (path picture bounding box.north east) --
    (path picture bounding box.north west) -- cycle;}, shift={(\x,\y)}] (V2) {$\Vb_2$};
    \node[tensor, path picture={
    \fill[fill=cyan!40!](path picture bounding box.south west)
    -- 
    (path picture bounding box.north east) --
    (path picture bounding box.north west) -- cycle;}, shift={(\x+1,\y)}] (G3) {$\Qt_3$};
    \node[tensor, path picture={
    \fill[fill=cyan!20!](path picture bounding box.south west)
    -- 
    (path picture bounding box.north east) --
    (path picture bounding box.north west) -- cycle;}, shift={(\x+2,\y)}] (G4) {$\Qt_4$};
   \draw[edge] (U1) -- (G2) node [midway,above]{\edgelabel{$R_1$}};
      \draw[edge] (G2) -- (sigma) node [midway,above]{\edgelabel{$R_2$}};
    \draw[edge] (sigma) -- (V2) node [midway,above]{\edgelabel{$R_2$}};
    \draw[edge] (V2) -- (G3) node [midway,above]{\edgelabel{$R^\prime_2$}};
    \draw[edge] (G3) -- (G4) node [midway,above]{\edgelabel{$R^\prime_3$}};
    \draw[edge](U1) -- (\x-3,\y-0.7) node [midway, left]{\edgelabel{$d_1$}};
    \draw[edge](G2) -- (\x-2,\y-0.7)node [midway, left]{\edgelabel{$d_2$}};
    \draw[edge](G3) -- (\x+1,\y-0.7)node [midway, left]{\edgelabel{$d_3$}};
    \draw[edge](G4) -- (\x+2,\y-0.7)node [midway,left] {\scalebox{0.7}{\dimcolor{$d_4$}}};
    \node[draw=none] at (\x,\y+0.7){$\overbrace{\hspace{2cm}}^{\text{merge \& apply rank-$R_3$ SVD}}$};
    \node[draw=none] at (\x+3,\y){$\longrightarrow$};
    \end{tikzpicture}
\begin{tikzpicture}[baseline=\baseline]
    \tikzset{tensor/.style = {minimum size = 0.5cm,shape = circle,thick,draw=black,inner sep = 0pt}, edge/.style = {thick,line width=.4mm},every loop/.style={}}
    \def\y{0}
    \def\x{0}
     \node[tensor, path picture={
    \fill[fill=cyan](path picture bounding box.south east)
    -- 
    (path picture bounding box.north west) --
    (path picture bounding box.north east) -- cycle;}, shift={(\x-3,\y)}] (U1) {$\Gt_1$};
    \node[tensor, path picture={
    \fill[fill=cyan!60!](path picture bounding box.south east)
    -- 
    (path picture bounding box.north west) --
    (path picture bounding box.north east) -- cycle;}, shift={(\x-2,\y)}] (G2) {$\Gt_2$};
    \node[tensor, path picture={
    \fill[fill=cyan!40!](path picture bounding box.south east)
    -- 
    (path picture bounding box.north west) --
    (path picture bounding box.north east) -- cycle;}, shift={(\x-1,\y)}] (G3) {$\Gt_3$};
    \node[tensor,fill=cyan!40!] (sigma) at (\x,\y){{}};
    \draw[edge=0.5] (sigma.south west) -- (sigma.north        east);
    \node[tensor, path picture={
    \fill[fill=cyan!40!](path picture bounding box.south west)
    -- 
    (path picture bounding box.north east) --
    (path picture bounding box.north west) -- cycle;}, shift={(\x+1,\y)}] (V3) {$\Vb_3$};
    \node[tensor, path picture={
    \fill[fill=cyan!20!](path picture bounding box.south west)
    -- 
    (path picture bounding box.north east) --
    (path picture bounding box.north west) -- cycle;}, shift={(\x+2,\y)}] (G4) {$\Qt_4$};
   \draw[edge] (U1) -- (G2) node [midway,above]{\edgelabel{$R_1$}};
      \draw[edge] (G2) -- (G3) node [midway,above]{\edgelabel{$R_2$}};
    \draw[edge] (G3) -- (sigma) node [midway,above]{\edgelabel{$R_3$}};
    \draw[edge] (sigma) -- (V3) node [midway,above]{\edgelabel{$R_3$}};
    \draw[edge] (V3) -- (G4) node [midway,above]{\edgelabel{$R^\prime_3$}};
    \draw[edge](U1) -- (\x-3,\y-0.7) node [midway, left]{\edgelabel{$d_1$}};
    \draw[edge](G2) -- (\x-2,\y-0.7)node [midway, left]{\edgelabel{$d_2$}};
    \draw[edge](G3) -- (\x-1,\y-0.7)node [midway, left]{\edgelabel{$d_3$}};
    \draw[edge](G4) -- (\x+2,\y-0.7)node [midway,left] {\scalebox{0.7}{\dimcolor{$d_4$}}};
    \node[draw=none] at (\x+1,\y+0.7){$\overbrace{\hspace{2cm}}^{\text{merge}}$};
    \node[draw=none] at (\x+3,\y){$\longrightarrow$};
    \end{tikzpicture}\\
&\begin{tikzpicture}[baseline=\baseline]
    \tikzset{tensor/.style = {minimum size = 0.5cm,shape = circle,thick,draw=black,inner sep = 0pt}, edge/.style = {thick,line width=.4mm},every loop/.style={}}
    \def\y{0}
    \def\x{0}
     \node[tensor, path picture={
    \fill[fill=cyan](path picture bounding box.south east)
    -- 
    (path picture bounding box.north west) --
    (path picture bounding box.north east) -- cycle;}, shift={(\x-3,\y)}] (U1) {$\Gt_1$};
    \node[tensor, path picture={
    \fill[fill=cyan!60!](path picture bounding box.south east)
    -- 
    (path picture bounding box.north west) --
    (path picture bounding box.north east) -- cycle;}, shift={(\x-2,\y)}] (G2) {$\Gt_2$};
    \node[tensor, path picture={
    \fill[fill=cyan!40!](path picture bounding box.south east)
    -- 
    (path picture bounding box.north west) --
    (path picture bounding box.north east) -- cycle;}, shift={(\x-1,\y)}] (G3) {$\Gt_3$};
    \node[tensor,fill=cyan!20!] (G4) at (\x,\y){$\Gt_4$};
   \draw[edge] (U1) -- (G2) node [midway,above]{\edgelabel{$R_1$}};
      \draw[edge] (G2) -- (G3) node [midway,above]{\edgelabel{$R_2$}};
    \draw[edge] (G3) -- (G4) node [midway,above]{\edgelabel{$R_3$}};
    \draw[edge](U1) -- (\x-3,\y-0.7) node [midway, left]{\edgelabel{$d_1$}};
    \draw[edge](G2) -- (\x-2,\y-0.7)node [midway, left]{\edgelabel{$d_2$}};
    \draw[edge](G3) -- (\x-1,\y-0.7)node [midway, left]{\edgelabel{$d_3$}};
    \draw[edge](G4) -- (\x,\y-0.7)node [midway,left] {\scalebox{0.7}{\dimcolor{$d_4$}}};
    \end{tikzpicture},   
\end{align*}
其中最后一步把剩余的所有张量收缩并合并为最终的核心张量 $\Gt_4$。可以证明，这一流程与执行 TT-SVD \emph{完全}等价。

如人所料，本章并未涵盖其他一些张量网络。下一节简要介绍其中几种，它们都是 Tucker 分解与 TT 分解的重要推广。
\subsection{其他分解：张量环、投影纠缠对态（PEPS）\& 层次 Tucker}

张量列（TT）的一种常见推广是张量环（Tensor Ring, TR）分解。
\paragraph{张量环分解}
张量环分解是 TT 分解的推广~\citep{zhao2016tensor}。它最初出现于量子物理，近年来在机器学习社区中逐渐流行起来~\citep{wang2017efficient,wang2018wide,yuan2018higher}。尽管 TR 分解存在若干数值稳定性方面的问题~\citep{batselier2018trouble,chen2020tensor}，实践中它一般所需参数更少，压缩率也高于 TT 分解~\citep{mickelin2020algorithms}。设 $\Xt\in\RR^{d_1\times\dots\times d_N}$
    为 $N$ 阶张量。再设
    $\Gt_n\in\RR^{R_{n-1}\times d_n\times R_n}$（$n\in[N]$）为核心张量。
    张量 $\Xt$ 的秩 $R$ 的\emph{张量环分解}由下式给出
\begin{align*}
    \Xt_{i_1,\dots,i_N}=\hspace{0.25cm}\sum_{r_1=1}^{R_1}\dots\sum_{r_{N}=1}^{R_N}
    (\Gt_1)_{r_N,i_1,r_1} (\Gt_2)_{ r_1,i_2,r_2}\dots  (\Gt_{N-1})_{r_{N-2},i_{N-1},r_{N-1}} (\Gt_{N})_{r_{N-1},i_{N},r_N},
\end{align*}
对所有指标 $i_1\in[d_1],\cdots,i_N\in[d_N]$ 均成立。
四阶张量的 TR 分解可用张量网络记法表示如下：
$$
\begin{tikzpicture}[baseline=-0.5ex]
    \tikzset{tensor/.style = {minimum size = 0.5cm,shape = circle,thick,draw=black,inner sep = 0pt}, edge/.style = {thick,line width=.4mm},every loop/.style={}}
    \def\x{6}
    \def\y{0}
     \node[tensor,fill=green!50!red!50!white] (B) at (\x+1,\y) {$\Xt$};
    \draw[edge] (B) -- (\x+0.3,0) node [midway,above] {\scalebox{0.7}{\dimcolor{$d_1$}}};
    \draw[edge] (B) -- (\x+1.6,0) node [midway,above] {\scalebox{0.7}{\dimcolor{$d_4$}}};;
    \draw[edge] (B) -- (\x+0.4,-0.3) node [midway, below] {\scalebox{0.7}{\textcolor{gray}{$d_2$}}};;
    \draw[edge] (B) -- (\x+1.5,-0.4)node [midway, below] {\scalebox{0.7}{\dimcolor{$d_3$}}};;
\end{tikzpicture}=
\begin{tikzpicture}[baseline=\baseline]
    \tikzset{tensor/.style = {minimum size = 0.5cm,shape = circle,thick,draw=black,inner sep = 0pt}, edge/.style   = {thick,line width=.4mm},every loop/.style={}}
    \def\y{0}
    \def\x{0}
  \node[tensor,fill=red!30!white] (D) at (\x-1,\y){$\Gt_1$};
  \node[tensor,fill=red!50!white] (A) at (\x,\y){$\Gt_2$};
  \node[tensor,fill=blue!50!white] (B) at (\x+1,\y){$\Gt_3$};
   \node[tensor,fill=blue!20!white] (C) at (\x+2,\y){$\Gt_4$};
    \draw[edge](A) -- (\x,\y-0.7)node [midway,left] {\edgelabel{$d_2$}};
    \draw[edge](B) -- (\x+1,\y-0.7)node [midway,left] {\edgelabel{$d_3$}};
    \draw[edge](C) -- (\x+2,\y-0.7)node [midway,left] {\edgelabel{$d_4$}};
     \draw[edge](D) -- (\x-1,\y-0.7)node [midway,left] {\edgelabel{$d_1$}};
    \draw[edge](D) -- (A) node [midway,above] {\edgelabel{$R$}};
    \draw[edge](A) -- (B) node [midway,above] {\edgelabel{$R$}};
    \draw[edge](B) -- (C) node [midway,above] {\edgelabel{$R$}};
    \draw[edge] (D) to [out=145,in=35, distance=1cm] (C);
     \node[draw=none] () at (\x+0.5,\y+0.8) {\scalebox{0.7}{\dimcolor{$R$}}};
\end{tikzpicture}.
$$
作为张量环分解的特例，令所有秩都等于 1，便可得到一个秩一分解：
$$
\begin{tikzpicture}[baseline=-0.5ex]
    \tikzset{tensor/.style = {minimum size = 0.5cm,shape = circle,thick,draw=black,inner sep = 0pt}, edge/.style = {thick,line width=.4mm},every loop/.style={}}
    \def\x{6}
    \def\y{0}
     \node[tensor,fill=green!50!red!50!white] (B) at (\x+1,\y) {$\Xt$};
    \draw[edge] (B) -- (\x+0.3,0) node [midway,above] {\scalebox{0.7}{\dimcolor{$d_1$}}};
    \draw[edge] (B) -- (\x+1.6,0) node [midway,above] {\scalebox{0.7}{\dimcolor{$d_4$}}};
    \draw[edge] (B) -- (\x+0.4,-0.3) node [midway, below] {\scalebox{0.7}{\dimcolor{$d_2$}}};
    \draw[edge] (B) -- (\x+1.5,-0.4)node [midway, below] {\scalebox{0.7}{\dimcolor{$d_3$}}};
\end{tikzpicture}=
\begin{tikzpicture}[baseline=\baseline]
    \tikzset{tensor/.style = {minimum size = 0.5cm,shape = circle,thick,draw=black,inner sep = 0pt}, edge/.style   = {thick,line width=.4mm},every loop/.style={}}
    \def\y{0}
    \def\x{0}
  \node[tensor,fill=red!30!white] (D) at (\x-1,\y){$\gb_1$};
  \node[tensor,fill=red!50!white] (A) at (\x,\y){$\gb_2$};
  \node[tensor,fill=blue!50!white] (B) at (\x+1,\y){$\gb_3$};
   \node[tensor,fill=blue!20!white] (C) at (\x+2,\y){$\gb_4$};
    \draw[edge](A) -- (\x,\y-0.7)node [midway,left] {\edgelabel{$d_2$}};
    \draw[edge](B) -- (\x+1,\y-0.7)node [midway,left] {\edgelabel{$d_3$}};
    \draw[edge](C) -- (\x+2,\y-0.7)node [midway,left] {\edgelabel{$d_4$}};
     \draw[edge](D) -- (\x-1,\y-0.7)node [midway,left] {\edgelabel{$d_1$}};
\end{tikzpicture}
=
\xb^1\outprod\xb^2\outprod\xb^3\outprod\xb^4.
$$
\paragraph{投影纠缠对态（PEPS）}
投影纠缠对态~（PEPS）一族由此成为把 TT 分解推广到量子物理中高阶格点结构的代表之一~\citep{verstraete2004renormalization,orus2014practical}。
PEPS 是由五阶核心张量按二维阵列排布而成的张量网络。对应于 $4\times 4$ 方格子的 PEPS 张量图如下：
$$
\begin{tikzpicture}[baseline=-0.5ex]
    \tikzset{tensor/.style = {minimum size = 0.5cm,shape = circle,thick,draw=black,inner sep = 0pt}, edge/.style = {thick,line width=.4mm},every loop/.style={}}
    \def\x{6}
    \def\y{0}
     \node[tensor,fill=red!30!blue!50!white] (A) at (\x,\y) {};
     \node[tensor,fill=red!30!blue!50!white] (B) at (\x+1,\y) {};
     \node[tensor,fill=red!30!blue!50!white] (C) at (\x+2,\y) {};
     \node[tensor,fill=red!30!blue!50!white] (D) at (\x+0.5,\y-1) {};
     \node[tensor,fill=red!30!blue!50!white] (E) at (\x+1.5,\y-1) {};
     \node[tensor,fill=red!30!blue!50!white] (F) at (\x+2.5,\y-1) {};
     \node[tensor,fill=red!30!blue!50!white] (G) at (\x+1.25,\y-2) {};
     \node[tensor,fill=red!30!blue!50!white] (H) at (\x+2.25,\y-2) {};
    \node[tensor,fill=red!30!blue!50!white] (I) at (\x+3.25,\y-2) {};
     \node[tensor,fill=red!30!blue!50!white] (J) at (\x+3,\y) {};
     \node[tensor,fill=red!30!blue!50!white] (K) at (\x+3.5,\y-1) {};
     \node[tensor,fill=red!30!blue!50!white] (L) at (\x+4.25,\y-2) {};
     \node[tensor,fill=red!30!blue!50!white] (M) at (\x+1.75,\y-3) {};

     \node[tensor,fill=red!30!blue!50!white] (N) at (\x+2.75,\y-3) {};
     \node[tensor,fill=red!30!blue!50!white] (O) at (\x+3.75,\y-3) {};
     \node[tensor,fill=red!30!blue!50!white] (P) at (\x+4.75,\y-3) {};
     
     \draw[edge](A) --(B);
    \draw[edge](A) --(\x,\y-0.65);
    \draw[edge](A) --(D);
    \draw[edge](B) --(C);
    \draw[edge](B) --(\x+1,\y-0.65);
    \draw[edge](B) --(E);
    \draw[edge](D) --(E);
    \draw[edge](C) --(F);
    \draw[edge](C) -- (\x+2,\y-0.65);
    \draw[edge](E) --(F);
    \draw[edge](F) -- (\x+2.5,\y-1.65);
    \draw[edge](D) --(G);
    \draw[edge](D) -- (\x+0.5,\y-1.65);
    \draw[edge](G) --(H);
    \draw[edge](E) --(H);
    \draw[edge](E) -- (\x+1.5,\y-1.65);
    \draw[edge](H) --(I);
    \draw[edge](F) --(I);
    \draw[edge](G) -- (\x+1.25,\y-2.65);
    \draw[edge](H) -- (\x+2.25,\y-2.65);
    \draw[edge](I) -- (\x+3.25,\y-2.65);
    \draw[edge](C)--(J);
     \draw[edge](K)--(F);
    \draw[edge](I)--(L);
    \draw[edge](J)--(K);
    \draw[edge](K)--(L);
    \draw[edge](J) -- (\x+3,\y-0.65);
    \draw[edge](K) -- (\x+3.5,\y-1.65);
    \draw[edge](L) -- (\x+4.25,\y-2.65);
    \draw[edge](M)--(N);
    \draw[edge](N)--(O);
    \draw[edge](O)--(P);
    \draw[edge](G)--(M);
    \draw[edge](H)--(N);
    \draw[edge](I)--(O);
    \draw[edge](L)--(P);
    \draw[edge](M)--(\x+1.75,\y-3.65);
    \draw[edge](N)--(\x+2.75,\y-3.65);
    \draw[edge](O)--(\x+3.75,\y-3.65);
    \draw[edge](P)--(\x+4.75,\y-3.65);
\end{tikzpicture}
$$
\paragraph{树形张量网络/~层次 Tucker 分解}
另一种颇受关注的张量网络是树形张量网络~（Tree Tensor Networks, TTN）或层次 Tucker~（Hierarchical Tucker, HT）分解，最早见于~\citep{hackbusch2009new, grasedyck2010hierarchical}。这一张量网络的关键优势在于，与 TT 分解这类较简单的模型相比，它的表达能力更强。四阶张量的 HT 分解可用张量网络图表示如下：
$$
\begin{tikzpicture}[baseline=\baseline]
    \tikzset{tensor/.style = {minimum size = 0.5cm,shape = circle,thick,draw=black,inner sep = 0pt}, edge/.style   = {thick,line width=.4mm},every loop/.style={}}
    \def\y{0}
    \def\x{0}
    \node[tensor,fill=blue!30!red!50!white] (A) at (\x+0.5,\y+1) {};
    \node[tensor,fill=blue!30!red!50!white] (B) at (\x,\y) {};
  \node[tensor,fill=blue!30!red!50!white] (C) at (\x+1,\y){};
   \node[tensor,fill=blue!30!red!50!white] (D) at (\x-1,\y-1){};
   \node[tensor,fill=blue!30!red!50!white] (E) at (\x,\y-1){};
   \node[tensor,fill=blue!30!red!50!white] (F) at (\x+1,\y-1){};
   \node[tensor,fill=blue!30!red!50!white] (G) at (\x+2,\y-1){};
   \draw[edge](A) -- (B);
   \draw[edge](A) -- (C);
   \draw[edge](B) -- (D);
   \draw[edge](B) -- (E);
    \draw[edge](C) -- (F);
   \draw[edge](C) -- (G);
    \draw[edge](D) -- (\x-1,\y-1.65);
   \draw[edge](E) -- (\x,\y-1.65);
    \draw[edge](F) -- (\x+1,\y-1.65);
   \draw[edge](G) -- (\x+2,\y-1.65);
\end{tikzpicture}
$$








\section{计算张量网络的梯度}\label{chapter:gradients}
在一般情形下优化张量网络，是许多研究领域共同面临的重大挑战。针对矩阵的优化技术虽已发展良多，但把这些方法推广到张量网络仍然不易~\citep{liao2019differentiable}。其困难源于张量收缩的高昂计算代价，也源于高维情形缺少高效的优化算法。此外，只要张量网络的结构略微复杂，手工用链式法则求梯度便令人却步。本章介绍一种优雅而直观的方法：借助张量网络的图示记法来推导高阶导数。该方法表明，张量网络函数的导数往往继承初始张量网络的低秩结构。

下一节先针对以向量为输入的函数，给出梯度与雅可比矩阵（Jacobian）的经典概念，再把它们推广到以张量为输入、以张量为输出的函数。
\subsection{雅可比矩阵}
我们先回顾定义在 $\mathbb{R}^n$ 上的函数的梯度与雅可比矩阵这两个经典概念。

\begin{definition}
  设 $f:\RR^n\to\RR$ 与 $g:\RR^n\to\RR^p$。$f$ 的梯度与 $g$ 在 $\theta\in\mathbb{R}^n$ 处的雅可比矩阵定义为
    $$\nabla_\theta f = \left[\frac{\partial f(\theta)}
{\partial\theta_1},\frac{\partial f(\theta)}{\partial\theta_2},\dots,\frac{\partial f(\theta)} 
    {\partial\theta_n}\right]^\ts
        =
    \begin{tikzpicture}[baseline=\baseline]
    \tikzset{tensor/.style = {minimum size = 0.5cm,shape = circle,thick,draw=black,inner sep = 0pt}, edge/.style   = {thick,line width=.4mm}}
    \def\y{0}
     \def\x{0}
    \node[tensor,fill=blue!40!red!60!white] (a1) at (\x,\y){};
    \draw[edge](a1) -- (\x+0.75,\y)node [midway, above]{\edgelabel{$n$}};
    \end{tikzpicture} \in\mathbb{R}^n,$$
    以及
    $$\frac{\partial g(\theta)}{\partial\theta} = \left(\frac{\partial g(\theta)_i}{\partial\theta_j}\right)_{i=1,\cdots,p\atop j=1,\cdots, n}
    =
    \begin{tikzpicture}[baseline=\baseline]
    \tikzset{tensor/.style = {minimum size = 0.5cm,shape = circle,thick,draw=black,inner sep = 0pt}, edge/.style   = {thick,line width=.4mm}}
    \def\y{0}
     \def\x{0}
    \node[tensor,fill=blue!20!red!50!white] (A) at (\x,\y){};
    \draw[edge](A) -- (\x+0.75,\y)node [midway, above]{\scalebox{0.7}{\dimcolor{$n$}}};
    \draw[edge](A) -- (\x-0.75,\y)node [midway, above]{\scalebox{0.7}{\dimcolor{$p$}}};
    \end{tikzpicture}\in\mathbb{R}^{p\times n}.$$  
\end{definition}
雅可比矩阵的概念可以自然推广到以张量为输入、输出张量（阶数任意）的函数。

\begin{definition}
 设 $f:\RR^{n_1\times\dots\times n_N}\to\RR^{m_1\times\dots\times m_M}$，。$f$ 在 $\theta\in\RR^{n_1\times\dots\times n_N}$ 处的\emph{雅可比张量}是张量 $\frac{\partial f}{\partial\theta}\in\RR^{m_1\times\dots\times m_M\times n_1\times\dots\times n_N}$
 定义为
\begin{align*}
   \frac{\partial f(\theta)}{\partial\theta} 
   &= 
   \left(\frac{\partial f(\theta)_{i_1,\dots,i_M}}{\partial\theta_{j_1,\dots,j_N}}\right)_{i_1\in[m_1],\cdots,i_M\in[m_M]\atop j_1\in[n_1],\cdots,j_N
\in[n_N]}\\
&=
\begin{tikzpicture}[baseline=-0.5ex]
    \tikzset{tensor/.style = {minimum size = 0.5cm,shape = circle,thick,draw=black,inner sep = 0pt}, edge/.style = {thick,line width=.4mm},every loop/.style={}}
    \def\x{6}
    \node[tensor,fill=green!50!red!50!white] (T) at (\x,0) {};
    \draw[edge] (T) -- (\x-0.85,0.2)node [midway, above]{\scalebox{0.7}{\dimcolor{$m_1$}}};
    \draw[edge] (T) -- (\x-0.85,-0.2)node [left]{\scalebox{0.7}{\dimcolor{$m_2$}}};
    \draw[edge] (T) -- (\x-0.1,-0.75)node [near end, left]{\scalebox{0.7}{\dimcolor{$m_M$}}};
    \draw[dotted] (\x-0.55,-0.3) -- (\x-0.2,-0.45);
    \draw[edge] (T) -- (\x+0.85,0.2)node [midway, above]{\scalebox{0.7}{\dimcolor{$n_1$}}};
    \draw[edge] (T) -- (\x+0.85,-0.2)node [right]{\scalebox{0.7}{\dimcolor{$n_2$}}};
    \draw[edge] (T) -- (\x+0.1,-0.75)node [near end, right]{\scalebox{0.7}{\dimcolor{$n_N$}}};
    \draw[dotted] (\x+0.55,-0.3) -- (\x+0.2,-0.45);
    \end{tikzpicture}\in\RR^{m_1\times\cdots\times m_M\times n_1\times\cdots\times n_N}. 
\end{align*}
\end{definition}




下面的定理表明，张量网络关于某个核心张量的一阶导数，用张量网络本身就能轻松算出。

\begin{theorem}\label{thm:gradients}
    设 $\Tt$ 是由张量网络给出的张量，其中 $\Gt$ 是在该张量网络中\emph{只出现一次}的核心张量。那么，把 $\Gt$ 从 $\Tt$ 的张量网络中移除，所得即 $\frac{\partial \Tt}{\partial \Gt}$。 
\end{theorem}
下面的例子说明如何将定理~\ref{thm:gradients}用于张量网络。
\paragraph{示例}
设 
$\begin{tikzpicture}[baseline=\baseline]
    \tikzset{tensor/.style = {minimum size = 0.5cm,shape = circle,thick,draw=black,inner sep = 0pt}, edge/.style = {thick,line width=.4mm},every loop/.style={}}
    \def\x{0}
    \def\y{0}
     \node[tensor,fill=green!50!red!50!white] (B) at (\x+1,\y) {$\Xt$};
    \draw[edge] (B) -- (\x+0.25,0) node [midway,above] {\scalebox{0.7}{\dimcolor{$d_1$}}};
    \draw[edge] (B) -- (\x+1.85,\y) node [midway,above] {\scalebox{0.7}{\dimcolor{$d_3$}}};
    \draw[edge] (B) -- (\x+1,\y-0.85) node [midway,left] {\scalebox{0.7}{\dimcolor{$d_2$}}};
\end{tikzpicture}
=
\begin{tikzpicture}[baseline=\baseline]
    \tikzset{tensor/.style = {minimum size = 0.5cm,shape = circle,thick,draw=black,inner sep = 0pt}, edge/.style = {thick,line width=.4mm},every loop/.style={}}
    \def\x{0}
    \def\y{0}
    \node[tensor,fill=red!30!white] (A) at (\x+2,\y) {$\Gt_1$};
    \node[tensor,fill=red!50!white] (B) at (\x+3.2,\y) {$\Gt_2$};
    \node[tensor,fill=blue!30!white] (C) at (\x+2.6,\y-0.7) {$\Gt_3$};
    \node[tensor,fill=blue!50!white] (D) at (\x+4.3,\y) {$\Gt_4$};
    \draw[edge] (A) -- (B) node [midway,above] {\scalebox{0.7}{\dimcolor{$R_1$}}};
    \draw[edge] (A) -- (C);
    \node[draw=none] () at (\x+2,\y-0.5) {\scalebox{0.7}{\dimcolor{$R_2$}}};
    \draw[edge] (B) -- (C);
    \node[draw=none] () at (\x+3.2,\y-0.5) {\scalebox{0.7}{\dimcolor{$R_3$}}};
    \draw[edge] (B) -- (D) node [midway,above] {\scalebox{0.7}{\dimcolor{$R_4$}}};
    \draw[edge] (A) -- (\x+2,\y+0.7) node [midway,left] {\scalebox{0.7}{\dimcolor{$d_1$}}};
    \draw[edge] (D) -- (\x+4.3,\y+0.7) node [midway,right] {\scalebox{0.7}{\dimcolor{$d_3$}}};
    \draw[edge] (C) -- (\x+2.6,\y-1.5) node [midway,left] {\scalebox{0.7}{\dimcolor{$d_2$}}};
\end{tikzpicture}
$ 为张量网络。注意，可以把 $\Xt$ 看作其张量网络表示中任一核心张量的函数，于是自然会问：$\Xt$ 关于例如 $\Gt_2$ 或 $\Gt_4$ 的雅可比矩阵是什么。
上述定理给出了计算这些导数的简便途径：
\begin{itemize}
\item 根据定理~\ref{thm:gradients}，要求 $\Xt$ 关于 $\Gt_2$ 的梯度，需从张量网络中移除核心张量 $\Gt_2$，即
$$
\frac{\partial}{\partial\Gt_2}\left(\begin{tikzpicture}[baseline=\baseline]
    \tikzset{tensor/.style = {minimum size = 0.5cm,shape = circle,thick,draw=black,inner sep = 0pt}, edge/.style = {thick,line width=.4mm},every loop/.style={}}
    \def\y{0}
    \def\x{0}
    \node[tensor,fill=red!30!white] (A) at (\x+2,\y) {$\Gt_1$};
    \node[tensor,fill=red!50!white] (B) at (\x+3.2,\y) {$\Gt_2$};
    \node[tensor,fill=blue!30!white] (C) at (\x+2.6,\y-0.7) {$\Gt_3$};
    \node[tensor,fill=blue!50!white] (D) at (\x+4.3,\y) {$\Gt_4$};
    \draw[edge] (A) -- (B) node [midway,above] {\scalebox{0.7}{\dimcolor{$R_1$}}};
    \draw[edge] (A) -- (C);
    \node[draw=none] () at (\x+2,\y-0.5) {\scalebox{0.7}{\dimcolor{$R_2$}}};
    \draw[edge] (B) -- (C);
    \node[draw=none] () at (\x+3.2,\y-0.5) {\scalebox{0.7}{\dimcolor{$R_3$}}};
    \draw[edge] (B) -- (D) node [midway,above] {\scalebox{0.7}{\dimcolor{$R_4$}}};
    \draw[edge] (A) -- (\x+2,\y+0.7) node [midway,left] {\scalebox{0.7}{\dimcolor{$d_1$}}};
    \draw[edge] (D) -- (\x+4.3,\y+0.7) node [midway,right] {\scalebox{0.7}{\dimcolor{$d_3$}}};
    \draw[edge] (C) -- (\x+2.6,\y-1.5) node [midway,left] {\scalebox{0.7}{\dimcolor{$d_2$}}};
    \end{tikzpicture}\right)
=
    \begin{tikzpicture}[baseline=\baseline]
    \tikzset{tensor/.style = {minimum size = 0.5cm,shape = circle,thick,draw=black,inner sep = 0pt}, edge/.style = {thick,line width=.4mm},every loop/.style={}}
    \def\y{-0.5}
    \def\x{0}
    \node[tensor,fill=red!30!white] (A) at (\x,\y+1) {$\Gt_1$};
    \node[tensor,fill=blue!30!white] (C) at (\x+0.6,\y+0.2) {$\Gt_3$};
    \node[tensor,fill=blue!50!white] (D) at (\x+2,\y+0.5) {$\Gt_4$};
    \draw[edge] (A) -- (C);
    \draw[edge] (A) -- (\x,\y+1.7) node [midway,left] {\scalebox{0.7}{\dimcolor{$d_1$}}};
    \draw[edge] (A) -- (\x+0.7,\y+1) node [midway,above] 
{\scalebox{0.7}{\dimcolor{$R_1$}}};
    \draw[edge] (D) -- (\x+3,\y+0.5) node [midway,above] {\scalebox{0.7}{\dimcolor{$R_4$}}};
    \draw[edge] (C) -- (\x+0.6,\y-0.7) node [midway,left] {\scalebox{0.7}{\dimcolor{$d_2$}}};
    \draw[edge] (C) -- (\x+1.3,\y+0.2) node [midway,above] {\scalebox{0.7}{\dimcolor{$R_3$}}};
    \draw[edge] (D) -- (\x+2,\y+1.3) node [midway,left] {\scalebox{0.7}{\dimcolor{$d_3$}}};
    \node[draw=none] () at (\x,\y+0.4) {\scalebox{0.7}{\dimcolor{$R_2$}}};
\end{tikzpicture}.
$$ 
要理解这一做法为何成立，可把 $\Xt$ 的张量分解看作矩阵-向量乘积：
$$
\Mb\vb 
=
\begin{tikzpicture}[baseline=\baseline]
    \tikzset{tensor/.style = {minimum size = 0.5cm,shape = circle,thick,draw=black,inner sep = 0pt}, edge/.style = {thick,line width=.4mm},every loop/.style={}}
    \def\y{1}
    \def\x{0}
     \node[tensor,fill=red!30!white](G_1) at (\x,\y) {$\Gt_1$};
    \node[tensor,fill=blue!30!white] (G_3) at (\x,\y-1) {$\Gt_3$};
     \node[tensor,fill=red!50!white](G_2) at (\x+2,\y-1) {$\Gt_2$};
    \node[tensor,fill=blue!50!white] (G_4) at (\x,\y-2) {$\Gt_4$};
    \node[tensor,fill=red!50!white] (G_2) at (\x+2,\y-1) {$\Gt_2$};
    \draw[edge](G_1) -- (\x-1,\y) -- (\x-1,\y-0.85) -- (\x-1.75,\y-0.85);
    \draw[edge](G_3)-- (\x-1.75,\y-1)node [left] {\scalebox{0.7}{\dimcolor{$d_1d_2d_3$}}};
    \draw[edge](G_2)--(G_1)node [midway,above] {\scalebox{0.7}{\dimcolor{$R_1$}}};
    \draw[edge](G_2)--(G_3)node [midway,above] {\scalebox{0.7}{\dimcolor{$R_3$}}};
    \draw[edge](G_1)--(G_3)node [midway,left] {\scalebox{0.7}{\dimcolor{$R_2$}}};
    \draw[edge](G_2)--(G_4)node [midway,above] {\scalebox{0.7}{\dimcolor{$R_4$}}};
    \draw[edge] (G_4) -- (\x-1,\y-2) -- (\x-1,\y-1.15) -- (\x-1.75,\y-1.15);
    \draw[dotted] (\x+0.45,\y+0.65) -- (\x+0.45,\y-2.35);
\end{tikzpicture},
$$
其中 $\vb = \vectorize(\Gt_2) $，而 $\Mb$ 是把张量网络余下部分矩阵化后得到的矩阵。 
利用线性映射的雅可比矩阵这一经典恒等式，便得到期望的结果：
$$\frac{\partial\Mb\vb}{\partial\vb} = \Mb = 
\begin{tikzpicture}[baseline=\baseline]
    \tikzset{tensor/.style = {minimum size = 0.5cm,shape = circle,thick,draw=black,inner sep = 0pt}, edge/.style = {thick,line width=.4mm},every loop/.style={}}
    \def\y{1}
    \def\x{0}
     \node[tensor,fill=red!30!white](G_1) at (\x,\y) {$\Gt_1$};
    \node[tensor,fill=blue!30!white] (G_3) at (\x,\y-1) {$\Gt_3$};
    \node[tensor,fill=blue!50!white] (G_4) at (\x,\y-2) {$\Gt_4$};
    \draw[edge](G_1) -- (\x-0.85,\y) -- (\x-0.85,\y-0.9) -- (\x-1.5,\y-0.9);
    \draw[edge](G_3)-- (\x-1.5,\y-1)node [left] {\scalebox{0.7}{\dimcolor{$d_1d_2d_3$}}};
    \draw[edge](G_1)-- (\x+0.85,\y) -- (\x+0.85,\y-0.9) -- (\x+1.5,\y-0.9);
    \draw[edge](G_3) -- (\x+1.5,\y-1) node [right] {\scalebox{0.7}{\dimcolor{$R_1R_3R_4$}}};
    \draw[edge](G_1)--(G_3)node [midway,left] {\scalebox{0.7}{\dimcolor{$R_2$}}};
    \draw[edge] (G_4) -- (\x+0.85,\y-2) -- (\x+0.85,\y-1.1) -- (\x+1.5,\y-1.1);
    \draw[edge] (G_4) -- (\x-0.85,\y-2) -- (\x-0.85,\y-1.1) -- (\x-1.5,\y-1.1);
\end{tikzpicture}\in\RR^{d_1d_2d_3\times R_1R_3R_4}.
$$
\item 根据定理~\ref{thm:gradients}，要求 $\Xt$ 关于 $\Gt_4$ 的梯度，需从相应的张量网络中移除核心张量 $\Gt_4$。即 
$$
    \frac{\partial}{\partial\Gt_4}\left(\begin{tikzpicture}[baseline=\baseline]
    \tikzset{tensor/.style = {minimum size = 0.5cm,shape = circle,thick,draw=black,inner sep = 0pt}, edge/.style = {thick,line width=.4mm},every loop/.style={}}
    \def\x{0}
    \def\y{0}
    \node[tensor,fill=red!30!white] (A) at (\x+2,\y) {$\Gt_1$};
    \node[tensor,fill=red!50!white] (B) at (\x+3.2,\y) {$\Gt_2$};
    \node[tensor,fill=blue!30!white] (C) at (\x+2.6,\y-0.7) {$\Gt_3$};
    \node[tensor,fill=blue!50!white] (D) at (\x+4.3,\y) {$\Gt_4$};
    \draw[edge] (A) -- (B) node [midway,above] {\scalebox{0.7}{\dimcolor{$R_1$}}};
    \draw[edge] (A) -- (C);
    \node[draw=none] () at (\x+2,\y-0.5) {\scalebox{0.7}{\dimcolor{$R_2$}}};
    \draw[edge] (B) -- (C);
    \node[draw=none] () at (\x+3.2,\y-0.5) {\scalebox{0.7}{\dimcolor{$R_3$}}};
    \draw[edge] (B) -- (D) node [midway,above] {\scalebox{0.7}{\dimcolor{$R_4$}}};
    \draw[edge] (A) -- (\x+2,\y+0.7) node [midway,left] {\scalebox{0.7}{\dimcolor{$d_1$}}};
    \draw[edge] (D) -- (\x+4.3,\y+0.7) node [midway,right] {\scalebox{0.7}{\dimcolor{$d_3$}}};
    \draw[edge] (C) -- (\x+2.6,\y-1.5) node [midway,left] {\scalebox{0.7}{\dimcolor{$d_2$}}};
    \end{tikzpicture}\right)
=
    \begin{tikzpicture}[baseline=\baseline]
    \tikzset{tensor/.style = {minimum size = 0.5cm,shape = circle,thick,draw=black,inner sep = 0pt}, edge/.style = {thick,line width=.4mm},every loop/.style={}}
    \def\y{-0.5}
    \def\x{0}
    \node[tensor,fill=red!30!white] (A) at (\x+2,\y+1) {$\Gt_1$};
    \node[tensor,fill=red!50!white] (B) at (\x+3.2,\y+1) {$\Gt_2$};
    \node[tensor,fill=blue!30!white] (C) at (\x+2.6,\y) {$\Gt_3$};
    \draw[edge] (A) -- (B) node [midway,above] {\scalebox{0.7}{\dimcolor{$R_1$}}};
    \draw[edge] (A) -- (C);
    \node[draw=none] () at (\x+2,\y+0.5) {\scalebox{0.7}{\dimcolor{$R_2$}}};
    \draw[edge] (B) -- (C);
    \node[draw=none] () at (\x+3.2,\y+0.5) {\scalebox{0.7}{\dimcolor{$R_3$}}};
    \draw[edge] (B) -- (\x+4,\y+1) node [midway,above] {\scalebox{0.7}{\dimcolor{$R_4$}}};
    \draw[edge] (A) -- (\x+2,\y+1.7) node [midway,left] {\scalebox{0.7}{\dimcolor{$d_1$}}};
    \draw[edge] (C) -- (\x+2.6,\y-0.7) node [midway,left] {\scalebox{0.7}{\dimcolor{$d_2$}}};
    \draw[edge] (\x+4.5,\y+1) -- (\x+4.5,\y+1.65) node [midway,right] {\scalebox{0.7}{\dimcolor{$d_3$}}};
    \end{tikzpicture}.
$$
注意 $\Gt_4$ 的悬空腿（dangling leg）仍留在雅可比张量中！为理解这一点，可再次把 $\Xt$ 的张量分解看作矩阵-向量乘积 $\Mb\vb$，使得 
$$
\Mb\vb = \begin{tikzpicture}[baseline=\baseline]
    \tikzset{tensor/.style = {minimum size = 0.5cm,shape = circle,thick,draw=black,inner sep = 0pt}, edge/.style = {thick,line width=.4mm},every loop/.style={}}
    \def\y{0.7}
    \def\x{0}
     \node[tensor,fill=red!30!white](G_1) at (\x,\y) {$\Gt_1$};
     \node[tensor,fill=red!50!white](G_2) at (\x+1,\y) {$\Gt_2$};
    \node[tensor,fill=blue!30!white] (G_3) at (\x+0.5,\y-1) {$\Gt_3$};
    \node[tensor,fill=blue!50!white] (G_4) at (\x+2.25,\y-1.5) {$\Gt_4$};
    \draw[edge](G_2)--(G_3);
    \draw[edge](G_1)--(G_2)node [midway,above] {\scalebox{0.7}{\dimcolor{$R_1$}}};
    \draw[edge](G_1)--(G_3);
    \draw[edge](G_2)--(G_4)node [midway,above=0.05cm] {\scalebox{0.7}{\dimcolor{$R_4$}}};
    \draw[edge](G_1) -- (\x,\y-1.35) -- (\x+0.35,\y-1.35) -- (\x+0.35,\y-2);
    \draw[edge](G_3)--(\x+0.5,\y-2)node [below] {\scalebox{0.7}{\dimcolor{$d_1d_2d_3$}}};
    \draw[edge](G_4) -- (\x+0.75,\y-1.5) -- (\x+0.75,\y-2);
    \draw[dotted] (\x+2.25,\y-0.75) -- (\x+1,\y-2);
\end{tikzpicture},
$$
此时 $\vb  =\vectorize(\Gt_4)$。
对该张量网络关于 $\vb = \vectorize(\Gb_4)$ 求导，便得到期望的结果：  $$\frac{\partial\Mb\vb}{\partial\vb} = \Mb
=
\begin{tikzpicture}[baseline=\baseline]
\tikzset{tensor/.style = {minimum size = 0.5cm,shape = circle,thick,draw=black,inner sep = 0pt}, edge/.style = {thick,line width=.4mm},every loop/.style={}}
    \def\y{1}
    \def\x{0}
     \node[tensor,fill=red!30!white](G_1) at (\x,\y) {$\Gt_1$};
     \node[tensor,fill=red!50!white](G_2) at (\x+1,\y) {$\Gt_2$};
    \node[tensor,fill=blue!30!white] (G_3) at (\x+0.5,\y-1) {$\Gt_3$};
    \draw[edge](G_2)--(G_3);
    \draw[edge](G_1)--(G_2)node [midway,above] {\scalebox{0.7}{\dimcolor{$R_1$}}};
    \draw[edge](G_1)--(G_3);
    \draw[edge](G_1) -- (\x,\y-1.35) -- (\x+0.35,\y-1.35) -- (\x+0.35,\y-2);
    \draw[edge](G_3)--(\x+0.5,\y-2)node [below] {\scalebox{0.7}{\dimcolor{$d_1d_2d_3$}}};
     \draw[edge](G_2) -- (\x+1.85,\y)node [midway, above] {\scalebox{0.7}{\dimcolor{$R_4$}}};
    \draw[edge](\x+1.85,\y-1.5) -- (\x+0.75,\y-1.5) -- (\x+0.75,\y-2);
\end{tikzpicture}\in\RR^{d_1d_2d_3\times R_4}.$$
注意，对某个张量求导只是从网络中移除相应的节点，而与它相连的边（腿）保持原样。
\end{itemize}
下面说明如何用定理~\ref{thm:gradients}不费力气地推出经典的导数与梯度计算式。 
\paragraph{用张量网络重新推导经典恒等式}
\begin{enumerate}
\item 
    $\frac{\partial\langle\ub,\vb\rangle}{\partial \ub}
=
    \frac{\partial\left(\begin{tikzpicture}[baseline=-0.1ex]
    \tikzset{tensor/.style = {minimum size = 0.5cm,shape = circle,thick,draw=black,inner sep = 0pt}, edge/.style = {thick,line width=.4mm},every loop/.style={}}
    \def\x{0}
    \def\y{0}
    \node[tensor,fill=red!50!white] (u) at (\x+1,\y) {$\ub$};
    \node[tensor,fill=red!20!white] (v) at (\x+2,\y) {$\vb$};
    \draw[edge] (u) -- (v);
    \end{tikzpicture}\right)}{\partial\left(\begin{tikzpicture}[baseline=-0.1ex]
    \tikzset{tensor/.style = {minimum size = 0.5cm,shape = circle,thick,draw=black,inner sep = 0pt}, edge/.style = {thick,line width=.4mm},every loop/.style={}}
    \def\x{0}
    \def\y{0}
    \node[tensor,fill=red!50!white] (u) at (\x+1,\y) {$\ub$};
    \draw[edge] (u) -- (\x+1.7,\y);
    \end{tikzpicture}\right)}
=
    \begin{tikzpicture}[baseline=-0.1ex]
    \tikzset{tensor/.style = {minimum size = 0.5cm,shape = circle,thick,draw=black,inner sep = 0pt}, edge/.style = {thick,line width=.4mm},every loop/.style={}}
    \def\x{6}
    \def\y{0}
    \node[tensor,fill=red!20!white] (v) at (\x+2,\y) {$\vb$};
    \draw[edge] (v) -- (\x+2.7,\y);
    \end{tikzpicture}
    =
    \vb
    $
\item
    $\frac{\partial\Ab\xb}{\partial\xb}
=
    \frac{\partial\left(\begin{tikzpicture}[baseline=-0.1ex]
    \tikzset{tensor/.style = {minimum size = 0.5cm,shape = circle,thick,draw=black,inner sep = 0pt}, edge/.style = {thick,line width=.4mm},every loop/.style={}}
    \def\x{6}
    \def\y{0}
    \node[tensor,fill=blue!50!white] (A) at (\x+1,\y) {$\Ab$};
    \node[tensor,fill=blue!20!white] (x) at (\x+2,\y) {$\xb$};
    \draw[edge] (A) -- (x);
    \draw[edge] (A) -- (\x+0.2,\y);
    \end{tikzpicture}\right)}{\partial\left(\begin{tikzpicture}[baseline=-0.1ex]
    \tikzset{tensor/.style = {minimum size = 0.5cm,shape = circle,thick,draw=black,inner sep = 0pt}, edge/.style = {thick,line width=.4mm},every loop/.style={}}
    \def\x{6}
    \def\y{0}
    \node[tensor,fill=blue!20!white] (x) at (\x+1,\y) {$\xb$};
    \draw[edge] (x) -- (\x+0.2,\y);
    \end{tikzpicture}\right)}
=
    \begin{tikzpicture}[baseline=-0.1ex]
    \tikzset{tensor/.style = {minimum size = 0.5cm,shape = circle,thick,draw=black,inner sep = 0pt}, edge/.style = {thick,line width=.4mm},every loop/.style={}}
    \def\x{6}
    \def\y{0}
    \node[tensor,fill=blue!50!white] (A) at (\x+1,\y) {$\Ab$};
    \draw[edge] (A) -- (\x+0.3,\y);
    \draw[edge] (A) -- (\x+1.7,\y);
    \end{tikzpicture}
=
    \Ab
    $
    
\item 
    $\frac{\partial \xb^\ts\Ab\xb}{\partial\Ab}
=
    \frac{\partial\left(\begin{tikzpicture}[baseline=-0.1ex]
    \tikzset{tensor/.style = {minimum size = 0.5cm,shape = circle,thick,draw=black,inner sep = 0pt}, edge/.style = {thick,line width=.4mm},every loop/.style={}}
    \def\x{6}
    \def\y{0}
    \node[tensor,fill=blue!50!white] (A) at (\x+1,\y) {$\Ab$};
    \node[tensor,fill=blue!20!white] (x) at (\x+2,\y) {$\xb$};
    \node[tensor,fill=blue!20!white] (xt) at (\x,\y) {$\xb$};
    \draw[edge] (A) -- (x);
    \draw[edge] (A) -- (xt);
    \end{tikzpicture}\right)}{\partial\left(\begin{tikzpicture}[baseline=-0.1ex]
    \tikzset{tensor/.style = {minimum size = 0.5cm,shape = circle,thick,draw=black,inner sep = 0pt}, edge/.style = {thick,line width=.4mm},every loop/.style={}}
    \def\x{6}
    \def\y{0}
    \node[tensor,fill=blue!50!white] (A) at (\x+1,\y) {$\Ab$};
    \draw[edge] (A) -- (\x+1.7,\y);
    \draw[edge] (A) -- (\x+0.3,\y);
    \end{tikzpicture}\right)}
=
    \begin{tikzpicture}[baseline=-0.1ex]
    \tikzset{tensor/.style = {minimum size = 0.5cm,shape = circle,thick,draw=black,inner sep = 0pt}, edge/.style = {thick,line width=.4mm},every loop/.style={}}
    \def\x{6}
    \def\y{0}
    \node[tensor,fill=blue!20!white] (x) at (\x+1,\y) {$\xb$};
    \node[tensor,fill=blue!20!white] (xt) at (\x-0.5,\y) {$\xb$};
    \draw[edge] (x) -- (\x+0.3,\y);
    \draw[edge] (xt) -- (\x+0.1,\y);
    \end{tikzpicture}
=
    \xb\outprod\xb
$
\item 
    $\frac{\partial\tr(\Ab)}{\partial\Ab}
    =
    \frac{\partial\left(\begin{tikzpicture}[baseline=-0.1ex]
    \tikzset{tensor/.style = {minimum size = 0.5cm,shape = circle,thick,draw=black,inner sep = 0pt}, edge/.style = {thick,line width=.4mm},every loop/.style={}}
    \def\x{6}
    \def\y{0}
    \node[tensor,fill=blue!50!white] (A) at (\x+1,\y) {$\Ab$};
    \draw[edge] (A) to [out=20,in=160, distance=1cm] (A);
    \end{tikzpicture}\right)}
    {\partial\left(\begin{tikzpicture}[baseline=-0.1ex]
    \tikzset{tensor/.style = {minimum size = 0.5cm,shape = circle,thick,draw=black,inner sep = 0pt}, edge/.style = {thick,line width=.4mm},every loop/.style={}}
    \def\x{6}
    \def\y{0}
    \node[tensor,fill=blue!50!white] (A) at (\x+1,\y) {$\Ab$};
    \draw[edge] (A) -- (\x+1.7,\y);
    \draw[edge] (A) -- (\x+0.3,\y);
    \end{tikzpicture}\right)}
=
\begin{tikzpicture}[baseline=-0.1ex]
    \tikzset{tensor/.style = {minimum size = 0.5cm,shape = circle,thick,draw=black,inner sep = 0pt}, edge/.style = {thick,line width=.4mm},every loop/.style={}}
    \def\x{6}
    \def\y{0}
    \node[draw=none] (A) at (\x+1,\y) {};
    \draw[edge] (A) to [out=20,in=160, distance=1cm] (A);
    \end{tikzpicture}
=
    \Ib
    $
    \item
    $\frac{\partial\Ab\xb}{\partial\Ab}
=
    \frac{\partial\left(\begin{tikzpicture}[baseline=-0.1ex]
    \tikzset{tensor/.style = {minimum size = 0.5cm,shape = circle,thick,draw=black,inner sep = 0pt}, edge/.style = {thick,line width=.4mm},every loop/.style={}}
    \def\x{6}
    \def\y{0}
    \node[tensor,fill=blue!50!white] (A) at (\x+1,\y) {$\Ab$};
    \node[tensor,fill=blue!20!white] (x) at (\x+2,\y) {$\xb$};
    \draw[edge] (A) -- (x);
    \draw[edge] (A) -- (\x+0.2,\y);
    \end{tikzpicture}\right)}{\partial\left(\begin{tikzpicture}[baseline=-0.1ex]
    \tikzset{tensor/.style = {minimum size = 0.5cm,shape = circle,thick,draw=black,inner sep = 0pt}, edge/.style = {thick,line width=.4mm},every loop/.style={}}
    \def\x{6}
    \def\y{0}
    \node[tensor,fill=blue!50!white] (A) at (\x+1,\y) {$\Ab$};
    \draw[edge] (A) -- (\x+0.3,\y);
    \draw[edge] (A) -- (\x+1.7,\y);
    \end{tikzpicture}\right)}
=
    \begin{tikzpicture}[baseline=-0.1ex]
    \tikzset{tensor/.style = {minimum size = 0.5cm,shape = circle,thick,draw=black,inner sep = 0pt}, edge/.style = {thick,line width=.4mm},every loop/.style={}}
    \def\x{6}
    \def\y{0}
    \draw[edge] (\x+2,\y) -- (\x+3,\y);
    \node[tensor,fill=blue!20!white] (x) at (\x+4.5,\y) {$\xb$};
    \draw[edge] (x) -- (\x+3.5,\y);
    \node[draw=none] at (\x+3.25,\y){$\outprod$};
    \end{tikzpicture}
=
    \Ib\outprod\xb
$,
\end{enumerate}
其中最后一个恒等式的结果是三阶张量。
定理~\ref{thm:gradients}适用于被求导的张量在网络中只出现一次的情形。下面的定理说明如何处理同一张量出现多次的情形。
\begin{theorem}\label{thm:multiple-tensors-gradients}
    设 $\Tt$ 为张量网络，$\Gt$ 是其中的核心张量。若 $\Gt$ 在 $\Tt$ 的张量网络中出现 $k$ 次，则把 $\Tt$ 的张量网络的 $k$ 份拷贝相加，每份拷贝中移除 $\Gt$ 的一次不同出现，所得即 $\frac{\partial\Tt}{\partial\Gt}$。
\end{theorem}
下面的例子说明如何应用定理~\ref{thm:multiple-tensors-gradients}。
\paragraph{示例}
\begin{itemize}
    \item 设 $\xb\in\RR^{d}$，$\Ab\in\RR^{d\times d}$
    \begin{align*}
    \frac{\partial \xb^\ts\Ab\xb}{\partial\xb}
=
    \frac{\partial\left(\begin{tikzpicture}[baseline=-0.1ex]
    \tikzset{tensor/.style = {minimum size = 0.5cm,shape = circle,thick,draw=black,inner sep = 0pt}, edge/.style = {thick,line width=.4mm},every loop/.style={}}
    \def\x{6}
    \def\y{0}
    \node[tensor,fill=blue!50!white] (A) at (\x+1,\y) {$\Ab$};
    \node[tensor,fill=blue!20!white] (x) at (\x+2,\y) {$\xb$};
    \node[tensor,fill=blue!20!white] (xprime) at 
    (\x,\y){$\xb$};
    \draw[edge] (A) -- (x)node [midway, above] {\scalebox{0.7}{\dimcolor{$d$}}};
    \draw[edge] (A) -- (xprime)node [midway, above] {\scalebox{0.7}{\dimcolor{$d$}}};
    \end{tikzpicture}\right)}{\partial\left(\begin{tikzpicture}[baseline=-0.1ex]
    \tikzset{tensor/.style = {minimum size = 0.5cm,shape = circle,thick,draw=black,inner sep = 0pt}, edge/.style = {thick,line width=.4mm},every loop/.style={}}
    \def\x{6}
    \def\y{0}
    \node[tensor,fill=blue!20!white] (x) at (\x,\y) {$\xb$};
    \draw[edge] (x) -- (\x-0.75,\y)node [midway, above] {\scalebox{0.7}{\dimcolor{$d$}}};
    \end{tikzpicture}\right)}
&=
    \begin{tikzpicture}[baseline=-0.1ex]
    \tikzset{tensor/.style = {minimum size = 0.5cm,shape = circle,thick,draw=black,inner sep = 0pt}, edge/.style = {thick,line width=.4mm},every loop/.style={}}
    \def\x{6}
    \def\y{0}
    \node[tensor,fill=blue!50!white] (A) at (\x+1,\y) {$\Ab$};
    \node[tensor,fill=blue!20!white] (x) at (\x+2,\y) {$\xb$};
    \draw[edge] (A) -- (x)node [midway, above] {\scalebox{0.7}{\dimcolor{$d$}}};
    \draw[edge] (A) -- (\x+0.25,\y)node [midway, above] {\scalebox{0.7}{\dimcolor{$d$}}};
    \end{tikzpicture}  
    +
    \begin{tikzpicture}[baseline=-0.1ex]
    \tikzset{tensor/.style = {minimum size = 0.5cm,shape = circle,thick,draw=black,inner sep = 0pt}, edge/.style = {thick,line width=.4mm},every loop/.style={}}
    \def\x{6}
    \def\y{0}
    \node[tensor,fill=blue!20!white] (x) at (\x+1,\y) {$\xb$};
    \node[tensor,fill=blue!50!white] (A) at (\x+2,\y) {$\Ab$};
    \draw[edge] (A) -- (x)node [midway, above] {\scalebox{0.7}{\dimcolor{$d$}}};
    \draw[edge] (A) -- (\x+2.75,\y)node [midway, above] {\scalebox{0.7}{\dimcolor{$d$}}};
    \end{tikzpicture}\\  
&=
    \Ab\xb + \Ab^\ts\xb = (\Ab\xb+\Ab^\ts)\xb.
\end{align*}
    \item 
设 $\Xb\in\RR^{n\times m},~\Wb\in\RR^{m\times n}$
    \begin{align*}
    \frac{\partial\norm{\Xb\Wb - \Yb}^2_F}{\partial \Wb}
&= 
    \frac{\partial\tr(\Wb^\ts\Xb^\ts\Xb\Wb)}{\partial\Wb}\\
&= 
    \frac{\partial\left( \begin{tikzpicture}[baseline=-0.1ex]
    \tikzset{tensor/.style = {minimum size = 0.5cm,shape = circle,thick,draw=black,inner sep = 0pt}, edge/.style = {thick,line width=.4mm},every loop/.style={}}
    \def\x{6}
    \def\y{0}
    \node[tensor,fill=red!20!white] (Wt) at (\x+1,\y) {$\Wb$};
    \node[tensor,fill=red!50!white] (Xt) at (\x+2,\y) {$\Xb$};
    \node[tensor,fill=red!50!white] (X) at (\x+3,\y) {$\Xb$};
    \node[tensor,fill=red!20!white] (W) at (\x+4,\y) {$\Wb$};
    \draw[edge] (Wt) -- (Xt)node [midway, above] {\scalebox{0.7}{\dimcolor{$m$}}};
    \draw[edge] (Xt) -- (X)node [midway, above] {\scalebox{0.7}{\dimcolor{$n$}}};
    \draw[edge] (X) -- (W)node [midway, above] {\scalebox{0.7}{\dimcolor{$n$}}};
    \draw[edge] (Wt) to [out=145,in=35, distance=1cm] (W);
    \end{tikzpicture}\right)}{\partial\left(\begin{tikzpicture}[baseline=-0.1ex]
    \tikzset{tensor/.style = {minimum size = 0.5cm,shape = circle,thick,draw=black,inner sep = 0pt}, edge/.style = {thick,line width=.4mm},every loop/.style={}}
    \def\x{6}
    \def\y{0}   
    \node[tensor,fill=red!20!white] (W) at (\x+4,\y) {$\Wb$};
    \draw[edge] (W) -- (\x+4.75,\y)node [midway, above] {\scalebox{0.7}{\dimcolor{$m$}}};
    \draw[edge] (W) -- (\x+3.25,\y)node [midway, above] {\scalebox{0.7}{\dimcolor{$n$}}};
    \end{tikzpicture}\right)}\\
&=
    \begin{tikzpicture}[baseline=-0.1ex]
    \tikzset{tensor/.style = {minimum size = 0.5cm,shape = circle,thick,draw=black,inner sep = 0pt}, edge/.style = {thick,line width=.4mm},every loop/.style={}}
    \def\x{6}
    \def\y{0}
    \node[tensor,fill=red!50!white] (Xt) at (\x+2,\y) {$\Xb$};
    \node[tensor,fill=red!50!white] (X) at (\x+3,\y) {$\Xb$};
    \node[tensor,fill=red!20!white] (W) at (\x+4,\y) {$\Wb$};
    \draw[edge] (Xt) -- (\x+1.25,\y)node [midway, above] {\scalebox{0.7}{\dimcolor{$m$}}};
    \draw[edge] (Xt) -- (X)node [midway, above] {\scalebox{0.7}{\dimcolor{$n$}}};
    \draw[edge] (X) -- (W)node [midway, above] {\scalebox{0.7}{\dimcolor{$m$}}};
    \draw[edge] (W) -- (\x+4.75,\y)node [midway, above] {\scalebox{0.7}{\dimcolor{$n$}}};
    \end{tikzpicture}
+
    \begin{tikzpicture}[baseline=-0.1ex]
    \tikzset{tensor/.style = {minimum size = 0.5cm,shape = circle,thick,draw=black,inner sep = 0pt}, edge/.style = {thick,line width=.4mm},every loop/.style={}}
    \def\x{6}
    \def\y{0}
    \node[tensor,fill=red!20!white] (Wt) at (\x+1,\y) {$\Wb$};
    \node[tensor,fill=red!50!white] (Xt) at (\x+2,\y) {$\Xb$};
    \node[tensor,fill=red!50!white] (X) at (\x+3,\y) {$\Xb$};
    \draw[edge] (Wt) -- (Xt)node [midway, above] {\scalebox{0.7}{\dimcolor{$m$}}};
    \draw[edge] (Xt) -- (X)node [midway, above] {\scalebox{0.7}{\dimcolor{$n$}}};
    \draw[edge] (Wt) -- (\x+0.25,\y)node [midway, above] {\scalebox{0.7}{\dimcolor{$n$}}};
    \draw[edge] (X) -- (\x+3.75,\y)node [midway, above] {\scalebox{0.7}{\dimcolor{$m$}}};    
    \end{tikzpicture}\\
&=
2\Xb^\ts\Xb\Wb.
\end{align*}
\end{itemize}
本章最后给出张量网络梯度的若干应用。
\subsection{应用}
张量网络梯度的应用十分广泛，这里着重介绍两个典型例子。 

我们先简要介绍损失函数的概念，并说明链式法则如何方便地用于涉及张量与雅可比张量的表达式。损失函数 $L:S\times S\to \RR^{\geq 0}$ 用于度量计算模型、算法或决策过程在给出输出 $\hat{y}\in S$ 作为真值 $y\in S$ 的预测时所产生的误差或代价。许多学习与优化算法的目标就是最小化这一损失。一种做法是用 Frobenius 范数来度量预测输出与真值之差。当输入与输出都是张量时，可由雅可比张量算出所需梯度，从而最小化该 Frobenius 范数。设 $\At\in\RR^{d_1\times d_2\times d_3}$ 为真值，$\Tt\in\RR^{d_1\times d_2\times d_3}$ 为预测输出。利用定理~\ref{thm:gradients}，损失 $L = \norm{\At-\Tt}_F^2$ 的梯度可写为：
    \begin{align*}
        \frac{\partial\norm{\At-\Tt}_F^2}{\partial\Tt} 
        &=
        \frac{\partial}{\partial\Tt}\left(\norm{\At}_F^2 + \norm{\Tt}_F^2 - 2\langle\At,\Tt\rangle\right)\\
        &= 
    \frac{\partial}{\partial\Tt}\left(\begin{tikzpicture}[baseline=\baseline]
    \tikzset{tensor/.style = {minimum size = 0.5cm,shape = circle,thick,draw=black,inner sep = 0pt}, edge/.style = {thick,line width=.4mm},every loop/.style={}}
    \def\y{0}
    \def\x{0}
    \node[tensor,fill=green!50!red!50!white] (St) at (\x,\y){\scalebox{0.85}{$\Tt$}};
    \node[tensor,fill=green!50!red!50!white] (Tt) at (\x+1,\y){\scalebox{0.85}{$\Tt$}};
    \draw[edge] (St) -- (Tt);
    \draw[edge] (St) to [out=45,in=135] (Tt);
    \draw[edge] (St) to [out=-45,in=-135] (Tt);
    \node[draw=none] () at (\x+0.5,\y+0.5) {\scalebox{0.7}{\dimcolor{$d_1$}}};
    \node[draw=none] () at (\x+0.5,\y+0.15) {\scalebox{0.7}{\dimcolor{$d_2$}}};
    \node[draw=none] () at (\x+0.5,\y-0.5) {\scalebox{0.7}{\dimcolor{$d_3$}}};
\end{tikzpicture}
-2
~
\begin{tikzpicture}[baseline=\baseline]
    \tikzset{tensor/.style = {minimum size = 0.5cm,shape = circle,thick,draw=black,inner sep = 0pt}, edge/.style = {thick,line width=.4mm},every loop/.style={}}
    \def\y{0}
    \def\x{0}
    \node[tensor,fill=blue!40!red!20!white] (St) at (\x,\y){\scalebox{0.85}{$\At$}};
    \node[tensor,fill=green!50!red!50!white] (Tt) at (\x+1,\y){\scalebox{0.85}{$\Tt$}};
    \draw[edge] (St) -- (Tt);
    \draw[edge] (St) to [out=45,in=135] (Tt);
    \draw[edge] (St) to [out=-45,in=-135] (Tt);
    \node[draw=none] () at (\x+0.5,\y+0.5) {\scalebox{0.7}{\dimcolor{$d_1$}}};
    \node[draw=none] () at (\x+0.5,\y+0.15) {\scalebox{0.7}{\dimcolor{$d_2$}}};
    \node[draw=none] () at (\x+0.5,\y-0.5) {\scalebox{0.7}{\dimcolor{$d_3$}}};
\end{tikzpicture}
\right)\\
&=
\left(\begin{tikzpicture}[baseline=\baseline]
    \tikzset{tensor/.style = {minimum size = 0.5cm,shape = circle,thick,draw=black,inner sep = 0pt}, edge/.style = {thick,line width=.4mm},every loop/.style={}}
    \def\y{0}
    \def\x{0}
    \node[tensor,fill=green!50!red!50!white] (Tt) at (\x,\y){\scalebox{0.85}{$\Tt$}};
    \draw[edge](Tt) -- (\x+0.85,\y)node[midway, above]{\edgelabel{$d_2$}};
    \draw[edge](Tt) -- (\x-0.85,\y)node[midway, above]{\edgelabel{$d_1$}};
    \draw[edge](Tt) -- (\x,\y-0.85)node[midway, left]{\edgelabel{$d_3$}};
    \end{tikzpicture}
+
\begin{tikzpicture}[baseline=\baseline]
    \tikzset{tensor/.style = {minimum size = 0.5cm,shape = circle,thick,draw=black,inner sep = 0pt}, edge/.style = {thick,line width=.4mm},every loop/.style={}}
    \def\y{0}
    \def\x{0}
    \node[tensor,fill=green!50!red!50!white] (Tt) at (\x,\y){\scalebox{0.85}{$\Tt$}};
    \draw[edge](Tt) -- (\x+0.85,\y)node[midway, above]{\edgelabel{$d_2$}};
    \draw[edge](Tt) -- (\x-0.85,\y)node[midway, above]{\edgelabel{$d_1$}};
    \draw[edge](Tt) -- (\x,\y-0.85)node[midway, left]{\edgelabel{$d_3$}};
    \end{tikzpicture}
-2~~
\begin{tikzpicture}[baseline=\baseline]
    \tikzset{tensor/.style = {minimum size = 0.5cm,shape = circle,thick,draw=black,inner sep = 0pt}, edge/.style = {thick,line width=.4mm},every loop/.style={}}
    \def\y{0}
    \def\x{0}
    \node[tensor,fill=blue!40!red!20!white] (Tt) at (\x,\y){\scalebox{0.85}{$\At$}};
    \draw[edge](Tt) -- (\x+0.85,\y)node[midway, above]{\edgelabel{$d_2$}};
    \draw[edge](Tt) -- (\x-0.85,\y)node[midway, above]{\edgelabel{$d_1$}};
    \draw[edge](Tt) -- (\x,\y-0.85)node[midway, left]{\edgelabel{$d_3$}};
    \end{tikzpicture}\right)\\
&=
2\left(\begin{tikzpicture}[baseline=\baseline]
    \tikzset{tensor/.style = {minimum size = 0.5cm,shape = circle,thick,draw=black,inner sep = 0pt}, edge/.style = {thick,line width=.4mm},every loop/.style={}}
    \def\y{0}
    \def\x{0}
    \node[tensor,fill=green!50!red!50!white] (Tt) at (\x,\y){\scalebox{0.85}{$\Tt$}};
    \draw[edge](Tt) -- (\x+0.85,\y)node[midway, above]{\edgelabel{$d_2$}};
    \draw[edge](Tt) -- (\x-0.85,\y)node[midway, above]{\edgelabel{$d_1$}};
    \draw[edge](Tt) -- (\x,\y-0.85)node[midway, left]{\edgelabel{$d_3$}};
    \end{tikzpicture}
-
\begin{tikzpicture}[baseline=\baseline]
    \tikzset{tensor/.style = {minimum size = 0.5cm,shape = circle,thick,draw=black,inner sep = 0pt}, edge/.style = {thick,line width=.4mm},every loop/.style={}}
    \def\y{0}
    \def\x{0}
    \node[tensor,fill=blue!40!red!20!white] (Tt) at (\x,\y){\scalebox{0.85}{$\At$}};
    \draw[edge](Tt) -- (\x+0.85,\y)node[midway, above]{\edgelabel{$d_2$}};
    \draw[edge](Tt) -- (\x-0.85,\y)node[midway, above]{\edgelabel{$d_1$}};
    \draw[edge](Tt) -- (\x,\y-0.85)node[midway, left]{\edgelabel{$d_3$}};
    \end{tikzpicture}\right).
\end{align*}
链式法则很容易推广到涉及张量变量的表达式。例如，考虑任意损失函数 $\mathcal{ L}: \mathbb{R}^{d_1\times \cdots \times d_N} \to \mathbb{R}$，我们希望它在 $N$ 阶张量输入 $\Tt$ 上被最小化。再设该张量输入以张量网络的形式参数化，其中 $\Gt \in \mathbb{R}^{m_1\times \cdots \times m_M}$ 是一个 $M$ 阶核心张量。要计算 $\mathcal L$ 关于 $\Gt$ 的雅可比张量，可如下使用链式法则： 
$$\frac{\partial  \mathcal L}{\partial\Gt} 
=
\frac{\partial \mathcal L}{\partial \Tt }\frac{\partial\Tt}{\partial\Gt}  $$
其中 $\frac{\partial  \mathcal L}{\partial\Gt}$ 的尺寸为 $m_1\times\cdots \times m_M$，$\frac{\partial  \mathcal L}{\partial \Tt }$ 的尺寸为%
\footnote{严格来说，$\frac{\partial  \mathcal L}{\partial\Gt}$ 的形状为 $1\times m_1 \cdots \times m_M$，$\frac{\partial  \mathcal L}{\partial \Tt }$ 的形状为 $1\times d_1\times \cdots \times d_N$，只是维度为 $1$ 的第一个平凡模自然被省略了。}
$d_1\times \cdots \times d_N$，而 $\frac{\partial\Tt}{\partial\Gt} $ 的尺寸为 $d_1\times \cdots \times d_N \times m_1 \cdots \times m_M$。$\frac{\partial  \mathcal L}{\partial \Tt }$ 与 $\frac{\partial\Tt}{\partial\Gt}  $ 之间的乘积默认理解为：将 $\frac{\partial  \mathcal L}{\partial \Tt }$ 的 $M$ 个模与 $\frac{\partial\Tt}{\partial\Gt}  $ 的前 $M$ 个模相收缩。

有了这些工具，我们就可以用两个应用来结束本章，它们都涉及张量网络的梯度计算。 
\begin{itemize}
\item\textbf{用梯度下降求 CP 分解} 要求张量 $\Tt\in\RR^{d_1\times\cdots\times d_N}$ 的近似 CP 分解~(见第~\ref{sec:cp-decomposition}~节)，可用梯度下降最小化损失 $L = \norm{\Tt-\CP{\Gb_1,\cdots,\Gb_N}}^2_F$。这里 $\CP{\Gb_1,\cdots,\Gb_N}$ 表示 CP 分解，其中 $\Gb_1,\cdots,\Gb_N$ 为因子矩阵。为对因子矩阵作优化，可对损失函数施加梯度下降。例如在 $N=3$ 时，为求 $\Gb_1$，应用链式法则可写 $\frac{\partial L}{\partial\Gb_1} = \frac{\partial L}{\partial\CP{\Gb_1,\Gb_2,\Gb_3}}\frac{\partial\CP{\Gb_1,\Gb_2,\Gb_3}}{\partial\Gb_1}$，其中 $\frac{\partial L}{\partial\CP{\Gb_1,\Gb_2,\Gb_3}}= 2(\Tt-\CP{\Gb_1,\Gb_2,\Gb_3})\in\RR^{d_1\times d_2\times d_3}$，且 
\begin{align*}
\frac{\partial\CP{\Gb_1,\Gb_2,\Gb_3}}{\partial\Gb_1}
=
\begin{tikzpicture}[baseline=\baseline]
    \tikzset{tensor/.style = {minimum size = 0.5cm,shape = circle,thick,draw=black,inner sep = 0pt}, edge/.style   = {thick,line width=.4mm},every loop/.style={}}
    \def\y{0}
    \def\x{0}
 \node[tensor,fill=blue!50!white] (G2) at (\x+1,\y){$\Gb_2$};
   \node[tensor,fill=blue!20!white] (G3) at (\x+2,\y){$\Gb_3$};
    \draw  node[fill = black,circle,inner sep=0pt,minimum size=4pt] (m) at (\x+1,\y+1) {};
    \draw[edge] (m) -- (\x,\y+0.1) node [midway,left] {\edgelabel{$R$}};
    \draw[edge] (G2) -- (m)node [midway,left] {\edgelabel{$R$}};
    \draw[edge] (G3) -- (m)node [midway,right] {\edgelabel{$R$}};
    \draw[edge](\x,\y-0.1) -- (\x,\y-0.7) node [midway,right] {\edgelabel{$d_1$}};
    \draw[edge](G2) -- (\x+1,\y-0.7)node [midway,left] {\edgelabel{$d_2$}};
    \draw[edge](G3) -- (\x+2,\y-0.7) node [midway,left] {\edgelabel{$d_3$}};
\end{tikzpicture}\in\RR^{d_1\times d_2\times d_3\times R}.
\end{align*}
观察所得的张量网络可知，它能用矩阵化与 Khatri-Rao 积表示：
\begin{align*}
  \frac{\partial L}{\partial\Gb_1} 
&=
\frac{\partial L}{\partial\CP{\Gb_1,\Gb_2,\Gb_3}}\frac{\partial\CP{\Gb_1,\Gb_2,\Gb_3}}{\partial\Gb_1}\\
&=
\begin{tikzpicture}[baseline=-0.5ex]
    \tikzset{tensor1/.style = {minimum width = 3cm,minimum height = 0.6cm,shape = rounded rectangle,thick,draw=black,inner sep = 0pt}, edge/.style = {thick,line width=.4mm},every loop/.style={}}
     \tikzset{tensor2/.style = {minimum size = 0.5cm,shape = circle,thick,draw=black,inner sep = 0pt}, edge/.style   = {thick,line width=.4mm},every loop/.style={}}
    \def\y{1}
    \def\x{0}
    \node[tensor1,fill=blue!45!red!20!white] (Tt) at (\x,\y){$2(\Tt-\CP{\Gb_1,\Gb_2,\Gb_3})$};
    \node[tensor2,fill=blue!50!white] (G2) at (\x,\y-1){$\Gb_2$};
   \node[tensor2,fill=blue!20!white] (G3) at (\x+1,\y-1){$\Gb_3$};
    \draw  node[fill = black,circle,inner sep=0pt,minimum size=4pt] (m) at (\x+0.5,\y-1.85) {};
    \draw[edge](G2)--(m)node [midway,left] {\edgelabel{$R$}};
    \draw[edge](G3)--(m)node [midway,right] {\edgelabel{$R$}};    \draw[edge](Tt)--(G2)node [midway,left] {\edgelabel{$d_2$}};
    \draw[edge](\x+1,\y-0.3)--(G3)node [midway,left] {\edgelabel{$d_3$}};
    \draw[edge](m)--(\x+0.5,\y-2.25)node [midway,left] {\edgelabel{$R$}};
    \draw[edge](\x-0.75,\y-0.3) -- (\x-0.75,\y-2.25)node [midway,left] {\edgelabel{$d_1$}};
\end{tikzpicture}\\
&=\begin{tikzpicture}[baseline=-0.5ex]
    \tikzset{tensor/.style = {minimum width = 3cm,minimum height = 0.6cm,shape = rounded rectangle,thick,draw=black,inner sep = 0pt}, edge/.style = {thick,line width=.4mm},every loop/.style={}}
    \def\y{1}
    \def\x{0}
    \node[tensor,fill=blue!45!red!20!white] (Tt) at (\x,\y){$2(\Tt-\CP{\Gb_1,\Gb_2,\Gb_3})$};
    \node[tensor,fill=blue!25!red!45!white] (Gt) at (\x+1,\y-1){$\Gb_2\krao\Gb_3$};
    \draw[edge](\x+0.5,\y-0.3) -- (\x+0.5,\y-0.7);
    \draw[edge](\x+0.6,\y-0.3) -- (\x+0.6,\y-0.7)node [midway,right] {\edgelabel{$d_2d_3$}};
    \draw[edge](\x-0.75,\y-0.3) -- (\x-0.75,\y-2)node [midway,left] {\edgelabel{$d_1$}};
    \draw[edge](\x+0.85,\y-1.3) -- (\x+0.85,\y-2)node [midway,left] {\edgelabel{$R$}};
\end{tikzpicture}\\
&=
2(\Tt-\CP{\Gb_1,\Gb_2,\Gb_3})_{(1)}(\Gb_2\krao\Gb_3)
\in\RR^{d_1\times R}.
\end{align*}
 \item\textbf{机器学习} 为压缩神经网络，可用张量分解来参数化全连接层的稠密权重矩阵，从而减少内存占用~\citep{novikov2015tensorizing}。设 $\Wb\in\RR^{P\times D}$ 为权重矩阵，$\xb\in\RR^{D}$ 为神经网络某一层的输入向量。输出 $\yb\in\RR^{P}$ 由权重矩阵 $\Wb$ 与输入向量 $\xb$ 收缩而得。 
    设 $P = \Pi_{i=1}^N{p_i}$ 且 $D =\Pi_{i=1}^Nd_i$。将 $\Wb$ 与 $\xb$ 重排为张量后，便有
\begin{align*}
    \begin{tikzpicture}[baseline=-0.1ex]
    \tikzset{tensor/.style = {minimum size = 0.5cm,shape = circle,thick,draw=black,inner sep = 0pt}, edge/.style = {thick,line width=.4mm},every loop/.style={}}
    \def\x{6}
    \def\y{0}
    \node[tensor,fill=blue!40!white] (y) at (\x,\y) {$\Yt$};
    \draw[edge](y)--(\x,\y+0.85)node [midway,right] {\scalebox{0.7}{\dimcolor{$p_1$}}};
    \draw[edge](y)--(\x-0.75,\y+0.7)node [midway,right=-0.01cm] {\scalebox{0.7}{\dimcolor{$p_2$}}};
    \draw[edge](y)--(\x-0.85,\y-0.7)node [very near end,right] {\scalebox{0.7}{\dimcolor{$p_{N-1}$}}};
    \draw[edge](y)--(\x,\y-0.85)node [midway,right] {\scalebox{0.7}{\dimcolor{$p_N$}}};
    \node[draw=none] at (\x-0.5,\y+0.1){$\vdots$};
    \end{tikzpicture}
=\sum_{d_1,\cdots,d_N}\Wt_{p_1,\cdots, p_N,d_1,\cdots, d_N}\Xt_{d_1,\cdots, d_N} =
\begin{tikzpicture}[baseline=-0.1ex]
    \tikzset{tensor/.style = {minimum size = 0.5cm,shape = circle,thick,draw=black,inner sep = 0pt}, edge/.style = {thick,line width=.4mm},every loop/.style={}}
    \def\x{6}
    \def\y{0}
    \node[tensor,fill=red!30!blue!20!white] (w) at (\x,\y) {$\Wt$};
    \node[tensor,fill=blue!30!red!20!white] (x) at (\x+1.5,\y) {$\Xt$};
    \draw[edge](w)--(x);
    \draw[edge](w)--(\x,\y+1)node [midway,left] {\scalebox{0.7}{\dimcolor{$p_1$}}};
    \draw[edge](w)--(\x-0.75,\y+0.8)node [midway,left] {\scalebox{0.7}{\dimcolor{$p_2$}}};
    \draw[edge](w)--(\x-0.75,\y-0.8)node [midway,left] {\scalebox{0.7}{\dimcolor{$p_{N-1}$}}};
    \draw[edge](w)--(\x,\y-1)node [midway,left=-0.1cm] {\scalebox{0.7}{\dimcolor{$p_N$}}};
    \node[draw=none] at (\x-0.5,\y+0.1){$\vdots$};
   \draw[edge] (w) to [out=-75,in=-100] (x);
    \draw[edge] (w) to [out=75,in=100] (x);
    \draw[dotted] (w) to [out=35,in=145] (x);
    \node[draw=none] at (\x+0.85,\y-0.15){$\vdots$};
    \node[draw=none] () at (\x+0.85,\y+0.85) {\scalebox{0.7}{\dimcolor{$d_1$}}};
    (\x+1.2,\y-0.2){$\vdots$};
    \node[draw=none] () at (\x+0.85,\y-0.85) {\scalebox{0.7}{\dimcolor{$d_N$}}};
\end{tikzpicture}.
\end{align*}
为简单起见，下面只考虑 $N=4$ 的情形。文献~\citep{novikov2015tensorizing} 的作者提出以 MPO 格式参数化 $\Wt$~(见第~\ref{sec:efficient-operations}~节)，即 
$\Wt = 
\begin{tikzpicture}[baseline=\baseline]
    \tikzset{tensor/.style = {minimum size = 0.5cm,shape = circle,thick,draw=black,inner sep = 0pt}, edge/.style   = {thick,line width=.4mm},every loop/.style={}}
    \def\y{0}
    \def\x{0}
  \node[tensor,fill=red!30!white] (G1) at (\x,\y){$\Gt_1$};
  \node[tensor,fill=red!50!white] (G2) at (\x+1,\y){\scalebox{0.85}{$\Gt_2$}};
  \node[tensor,fill=blue!30!white] (G3) at (\x+2,\y){$\Gt_3$};
  \node[tensor,fill=blue!50!white] (G4) at (\x+3,\y){$\Gt_4$};
    \draw[edge](G1) -- (\x,\y-0.85)node [midway,left] {\edgelabel{$d_1$}};
    \draw[edge](G2) -- (\x+1,\y-0.85)node [midway,left] {\edgelabel{$d_2$}};
     \draw[edge](G3) -- (\x+2,\y-0.85)node [midway,left] {\edgelabel{$d_3$}};
    \draw[edge](G4) -- (\x+3,\y-0.85)node [midway,left] {\edgelabel{$d_4$}};
    \draw[edge](G1) -- (G2) node [midway,above] {\edgelabel{$R_1$}};
    \draw[edge](G2) -- (G3) node [midway,above] {\edgelabel{$R_2$}};
    \draw[edge](G3) -- (G4) node [midway,above] {\edgelabel{$R_3$}};
    \draw[edge](G1) -- (\x,\y+0.85)node [midway,left] {\edgelabel{$p_1$}};
    \draw[edge](G2) -- (\x+1,\y+0.85)node [midway,left] {\edgelabel{$p_2$}};
    \draw[edge](G3) -- (\x+2,\y+0.85)node [midway,left] {\edgelabel{$p_3$}};
    \draw[edge](G4) -- (\x+3,\y+0.85)node [midway,left] {\edgelabel{$p_4$}};
\end{tikzpicture}.$
在反向传播步骤中，由于要更新的是权重矩阵，就需要对全部核心张量作优化。按链式法则，对 $i\in[4]$ 有 $\frac{\partial L}{\partial \Gt_i} = \frac{\partial L}{\partial \Yt}\frac{\partial\Yt}{\partial\Gt_i}$，其中 $L$ 表示损失函数。例如 $i=2$ 时，需求 $\Yt$ 关于 $\Gt_2$ 的导数。利用定理~\ref{thm:gradients}，只需从张量网络中删去 $\Gt_2$ 便可得到结果：
\begin{align*}
\frac{\partial\Yt}{\partial\Gt_2} 
&=
\frac{\partial\left(\begin{tikzpicture}[baseline=\baseline]
    \tikzset{tensor/.style = {minimum size = 0.5cm,shape = circle,thick,draw=black,inner sep = 0pt}, edge/.style   = {thick,line width=.4mm},every loop/.style={}}
    \def\y{0}
    \def\x{0}
  \node[tensor,fill=red!30!white] (G1) at (\x,\y){$\Gt_1$};
  \node[tensor,fill=red!50!white] (G2) at (\x+1,\y){\scalebox{0.85}{$\Gt_2$}};
  \node[tensor,fill=blue!30!white] (G3) at (\x+2,\y){$\Gt_3$};
  \node[tensor,fill=blue!50!white] (G4) at (\x+3,\y){$\Gt_4$};
  \node[tensor,fill=blue!30!red!20!white] (x) at (\x+1.5,\y-1.5) {$\Xt$};
    \draw[edge](G1) -- (\x,\y-0.85)node [midway,left] {\edgelabel{$d_1$}};
    \draw[edge](x)--(\x,\y-0.85);
    \draw[edge](G2) -- (\x+1,\y-0.85)node [midway,left] {\edgelabel{$d_2$}};
    \draw[edge](x)--(\x+1,\y-0.85);
     \draw[edge](G3) -- (\x+2,\y-0.85)node [midway,left] {\edgelabel{$d_3$}};
    \draw[edge](x)--(\x+2,\y-0.85);
    \draw[edge](G4) -- (\x+3,\y-0.85)node [midway,left] {\edgelabel{$d_4$}};
    \draw[edge](x)--(\x+3,\y-0.85);
    \draw[edge](G1) -- (G2) node [midway,above] {\edgelabel{$R_1$}};
    \draw[edge](G2) -- (G3) node [midway,above] {\edgelabel{$R_2$}};
    \draw[edge](G3) -- (G4) node [midway,above] {\edgelabel{$R_3$}};
    \draw[edge](G1) -- (\x,\y+0.85)node [midway,left] {\edgelabel{$p_1$}};
    \draw[edge](G2) -- (\x+1,\y+0.85)node [midway,left] {\edgelabel{$p_2$}};
    \draw[edge](G3) -- (\x+2,\y+0.85)node [midway,left] {\edgelabel{$p_3$}};
    \draw[edge](G4) -- (\x+3,\y+0.85)node [midway,left] {\edgelabel{$p_4$}};
\end{tikzpicture}\right)}{\partial\Gt_2}\\
&=
\begin{tikzpicture}[baseline=\baseline]
    \tikzset{tensor/.style = {minimum size = 0.5cm,shape = circle,thick,draw=black,inner sep = 0pt}, edge/.style   = {thick,line width=.4mm},every loop/.style={}}
    \def\y{-1}
    \def\x{0}
  \node[tensor,fill=red!30!white] (G1) at (\x,\y+1.5){$\Gt_1$};
    \node[tensor,fill=blue!30!white] (G3) at (\x+2,\y+1.5){$\Gt_3$};
  \node[tensor,fill=blue!50!white] (G4) at (\x+3,\y+1.5){$\Gt_4$};
  \node[tensor,fill=blue!30!red!20!white] (x) at (\x+1.5,\y) {$\Xt$};
    \draw[edge](G1) -- (\x,\y+0.85)node [midway,left] {\edgelabel{$d_1$}};
    \draw[edge](x)--(\x,\y+0.85);
    \draw[edge](\x+1,\y+1.35) -- (\x+1,\y+0.85)node [midway,left] {\edgelabel{$d_2$}};
    \draw[edge](x)--(\x+1,\y+0.85);
     \draw[edge](G3) -- (\x+2,\y+0.85)node [midway,left] {\edgelabel{$d_3$}};
    \draw[edge](x)--(\x+2,\y+0.85);
    \draw[edge](G4) -- (\x+3,\y+0.85)node [midway,left] {\edgelabel{$d_4$}};
    \draw[edge](x)--(\x+3,\y+0.85);
    \draw[edge](G3) -- (G4) node [midway,above] {\edgelabel{$R_3$}};
    \draw[edge](G1) -- (\x,\y+2.35)node [midway,left] {\edgelabel{$p_1$}};
    \draw[edge](G3) -- (\x+2,\y+2.35)node [midway,left] {\edgelabel{$p_3$}};
    \draw[edge](G4) -- (\x+3,\y+2.35)node [midway,left] {\edgelabel{$p_4$}};
    \draw[edge](G1)--(\x+0.85,\y+1.5) node [midway,above] {\edgelabel{$R_1$}};
    \draw[edge](G3)--(\x+1.25,\y+1.5) node [midway,above] {\edgelabel{$R_2$}};
    \draw[edge](\x+1,\y+1.6)--(\x+1,\y+2.35) node [midway,left] {\edgelabel{$p_2$}};
\end{tikzpicture}\in\RR^{p_1\times R_1\times p_2\times R_2\times p_3\times p_4\times d_2}.
\end{align*}
于是损失的梯度为
\begin{align*}
    \frac{\partial L}{\partial\Gt_2} = \frac{\partial L}{\partial\Yt}\frac{\partial \Yt}{\partial\Gt_2}
    =
    \begin{tikzpicture}[baseline=\baseline]
    \tikzset{tensor/.style = {minimum size = 0.5cm,shape = circle,thick,draw=black,inner sep = 0pt}, edge/.style   = {thick,line width=.4mm},every loop/.style={}}
    \def\y{-1.5}
    \def\x{0}
    \node[tensor,fill=green!30!white] (L) at (\x+1.5,\y+3.25){$\frac{\partial L}{\partial\Yt}$};
  \node[tensor,fill=red!30!white] (G1) at (\x,\y+1.5){$\Gt_1$};
    \node[tensor,fill=blue!30!white] (G3) at (\x+2,\y+1.5){$\Gt_3$};
  \node[tensor,fill=blue!50!white] (G4) at (\x+3,\y+1.5){$\Gt_4$};
  \node[tensor,fill=blue!30!red!20!white] (x) at (\x+1.5,\y) {$\Xt$};
    \draw[edge](G1) -- (\x,\y+0.85)node [midway,left] {\edgelabel{$d_1$}};
    \draw[edge](x)--(\x,\y+0.85);
    \draw[edge](\x+1,\y+1.35) -- (\x+1,\y+0.85)node [midway,left] {\edgelabel{$d_2$}};
    \draw[edge](x)--(\x+1,\y+0.85);
     \draw[edge](G3) -- (\x+2,\y+0.85)node [midway,left] {\edgelabel{$d_3$}};
    \draw[edge](x)--(\x+2,\y+0.85);
    \draw[edge](G4) -- (\x+3,\y+0.85)node [midway,left] {\edgelabel{$d_4$}};
    \draw[edge](x)--(\x+3,\y+0.85);
    \draw[edge](G3) -- (G4) node [midway,above] {\edgelabel{$R_3$}};
    \draw[edge](G1) -- (\x,\y+2.35)node [midway,left] {\edgelabel{$p_1$}};
    \draw[edge](L)--(\x,\y+2.35);
    \draw[edge](G3) -- (\x+2,\y+2.35)node [midway,left] {\edgelabel{$p_3$}};
    \draw[edge](L)--(\x+2,\y+2.35);
    \draw[edge](G4) -- (\x+3,\y+2.35)node [midway,left] {\edgelabel{$p_4$}};
    \draw[edge](L)--(\x+3,\y+2.35);
    \draw[edge](G1)--(\x+0.85,\y+1.5) node [midway,above] {\edgelabel{$R_1$}};
    \draw[edge](G3)--(\x+1.25,\y+1.5) node [midway,above] {\edgelabel{$R_2$}};
    \draw[edge](\x+1,\y+1.6)--(\x+1,\y+2.35) 
    node [midway,left] {\edgelabel{$p_2$}};
    \draw[edge](L)--(\x+1,\y+2.35);
\end{tikzpicture}\in\RR^{R_1\times p_2\times R_2\times d_2}.
\end{align*}

\end{itemize}







\section{概率分布与随机张量网络}\label{chapter:prob}
本章讨论张量网络的一类非常有用的应用：以直观且计算高效的方式表示高维概率分布。例如，TT 等张量网络分解能够高效编码概率分布及其边缘分布~\citep{gardiner2024tensor, glasser2019expressive}。我们先说明如何用张量网络高效地建模概率分布，随后把重心转向研究随机张量彼此收缩所得张量的统计性质。本章两部分都将表明，张量网络能帮助我们轻松地得到涉及概率分布与高阶张量的非平凡结论。 
\subsection{概率分布与张量分解}
本节讨论如何用张量网络参数化概率分布。
\subsubsection{作为张量的概率分布}\label{sub:probs}
 用张量网络建模概率分布的一种常见做法，是直接表示概率本身。考虑 $N$ 个离散随机变量的多元概率质量函数 $\PP(X_1,\cdots,X_N)$。这一联合概率分布可以看作一个元素非负、且元素之和为一的 $N$ 阶张量： 
\begin{definition}\label{def:probability}
设 $\PP$ 是 $N$ 个离散随机变量 $X_1,\cdots,X_N$ 上的联合分布，它们分别取值于 $[d_1], [d_2], \cdots,[d_N]$。 
称由 $\Tt_{i_1\cdots i_N} = \PP(X_1=i_1,\cdots,X_N=i_N)$ 定义的张量 $\Tt\in\RR^{d_1\times\cdots\times d_N}$ \emph{表示}概率 $\PP$，并记为 $\Tt \simeq  \PP(X_1,\cdots,X_N)$。由构造可知，$\Tt$ 的元素非负且总和为一： 
$$
\sum_{i_1,\cdots,i_N}\Tt_{i_1,\cdots,i_N} 
=
\begin{tikzpicture}[baseline=\baseline]
    \tikzset{tensor/.style = {minimum size = 0.5cm,shape = circle,thick,draw=black,inner sep = 0pt}, edge/.style = {thick,line width=.4mm},every loop/.style={}}
    \def\x{0}
    \def\y{0}
    \node[tensor,fill=blue!20!white] (T) at (\x+1,\y) {$\Tt$};
    \draw  node[fill = black,circle,inner sep=0pt,minimum size=4pt] (m_1) at (\x,\y-0.65) {};
    \draw  node[fill = black,circle,inner sep=0pt,minimum size=4pt] (m_2) at (\x+0.5,\y-0.65) {};
    \draw  node[fill = black,circle,inner sep=0pt,minimum size=4pt] (m_3) at (\x+2,\y-0.65) {};
    \node[draw=none] at (\x+1.25,\y-0.65){$\dots$};
    \draw[edge] (T) -- (m_1)node [very near end,above]
    {\scalebox{0.7}{\dimcolor{$d_1$}}};
    \draw[edge] (T) -- (m_2)node [very near end,right]
    {\scalebox{0.7}{\dimcolor{$d_2$}}};
    \draw[edge] (T) -- (m_3)node [very near end,above]
    {\scalebox{0.7}{\dimcolor{$d_N$}}};
\end{tikzpicture}
= 1,
$$
其中 $
\begin{tikzpicture}[baseline=\baseline]
    \tikzset{tensor/.style = {minimum size = 0.5cm,shape = circle,thick,draw=black,inner sep = 0pt}, edge/.style = {thick,line width=.4mm},every loop/.style={}}
    \def\x{0}
    \def\y{0}
    \draw  node[fill = black,circle,inner sep=0pt,minimum size=4pt] (m_1) at (\x,\y) {};
    \draw[edge](m_1)--(\x+1,\y);
    \end{tikzpicture}
$
是全 1 向量~(见命题~\ref{prop:copy.tensor} 中的条目~\ref{itm:one-copy-tensor})。
\end{definition}
下面说明如何用张量网络记法表示边缘概率与条件概率。
\paragraph{边缘概率}
设 $\Tt\in\RR^{d_1\times d_2\times d_3}$ 为表示概率分布 $\PP(X_1, X_2, X_3)$ 的张量，则 
\begin{enumerate}
    \item 由全概率公式，$X_2$ 与 $X_3$ 的边缘分布可表示为
    
\begin{align*}
    \PP(X_2=i_2,X_3=i_3) 
    &=
    \sum_{i_1=1}^{d_1}\PP(X_1= i_1,X_2=i_2,X_3=i_3) \\
    &= 
\begin{tikzpicture}[baseline=\baseline]
    \tikzset{tensor/.style = {minimum size = 0.5cm,shape = circle,thick,draw=black,inner sep = 0pt}, edge/.style = {thick,line width=.4mm},every loop/.style={}}
    \def\y{0.25}
    \def\x{0}
    \node[tensor,fill=blue!20!white] (T) at (\x,\y) {$\Tt$};
    \draw  node[fill = black,circle,inner sep=0pt,minimum size=4pt] (m_1) at (\x-0.5,\y-0.65) {};
    \draw[edge] (T) -- (m_1);
    \draw[edge](T) -- (\x,\y-0.65)node [very near end,below]
    {\scalebox{0.7}{\indcolor{$i_2$}}};
    \draw[edge](T) -- (\x+0.5,\y-0.65)node [very near end,below]
    {\scalebox{0.7}{\indcolor{$i_3$}}};
\end{tikzpicture}
 = \sum_{i_1=1}^{d_1}\Tt_{i_1,i_2,i_3},
\end{align*}
对所有 $i_2,i_3$ 成立。这表明，表示边缘分布的张量可由表示联合分布的张量经一次简单运算得到，即
\begin{align*}
    \text{If}~~~ \begin{tikzpicture}[baseline=\baseline]
    \tikzset{tensor/.style = {minimum size = 0.5cm,shape = circle,thick,draw=black,inner sep = 0pt}, edge/.style = {thick,line width=.4mm},every loop/.style={}}
    \def\y{0.25}
    \def\x{0}
    \node[tensor,fill=blue!20!white] (T) at (\x,\y) {$\Tt$};
    \draw[edge] (T) -- (\x-0.5,\y-0.65);
    \draw[edge] (T) -- (m_1);
    \draw[edge](T) -- (\x,\y-0.65);
    \draw[edge](T) -- (\x+0.5,\y-0.65);
\end{tikzpicture}\simeq\PP(X_1, X_2, X_3),~~~\text{then}~~~
\begin{tikzpicture}[baseline=\baseline]
    \tikzset{tensor/.style = {minimum size = 0.5cm,shape = circle,thick,draw=black,inner sep = 0pt}, edge/.style = {thick,line width=.4mm},every loop/.style={}}
    \def\y{0.25}
    \def\x{0}
    \node[tensor,fill=blue!20!white] (T) at (\x,\y) {$\Tt$};
    \draw  node[fill = black,circle,inner sep=0pt,minimum size=4pt] (m_1) at (\x-0.5,\y-0.65) {};
    \draw[edge] (T) -- (m_1);
    \draw[edge](T) -- (\x,\y-0.65);
    \draw[edge](T) -- (\x+0.5,\y-0.65);
\end{tikzpicture}\simeq\PP(X_2, X_3).
\end{align*}
\item
类似地，$X_2$ 的边缘分布可表示为
\begin{align*}
\PP(X_2=i_2) 
&= \sum_{i_1,i_3=1}^{d_1,d_3}\PP(X_1=i_1,X_2= i_2, X_3=i_3)\\
&=
\begin{tikzpicture}[baseline=\baseline]
    \tikzset{tensor/.style = {minimum size = 0.5cm,shape = circle,thick,draw=black,inner sep = 0pt}, edge/.style = {thick,line width=.4mm},every loop/.style={}}
    \def\y{0.25}
    \def\x{0}
    \node[tensor,fill=blue!20!white] (T) at (\x,\y) {$\Tt$};
    \draw  node[fill = black,circle,inner sep=0pt,minimum size=4pt] (m_1) at (\x-0.5,\y-0.65) {};        \draw  node[fill = black,circle,inner sep=0pt,minimum size=4pt] (m_2) at (\x+0.5,\y-0.65) {};
    \draw[edge] (T) -- (m_1);
    \draw[edge](T) -- (m_2);
    \draw[edge](T) -- (\x,\y-0.65)node [very near end,below]
    {\scalebox{0.7}{\indcolor{$i_2$}}};
\end{tikzpicture}=\sum_{i_1,i_3=1}^{d_1,d_3}\Tt_{i_1,i_2,i_3},
\end{align*}
对所有 $i_2$ 成立。因此在张量网络记法下，$X_2$ 的边缘分布可表示为：
\begin{align*}
    \text{If}~~~ \begin{tikzpicture}[baseline=\baseline]
    \tikzset{tensor/.style = {minimum size = 0.5cm,shape = circle,thick,draw=black,inner sep = 0pt}, edge/.style = {thick,line width=.4mm},every loop/.style={}}
    \def\y{0.25}
    \def\x{0}
    \node[tensor,fill=blue!20!white] (T) at (\x,\y) {$\Tt$};
    \draw[edge] (T) -- (\x-0.5,\y-0.65);
    \draw[edge] (T) -- (m_1);
    \draw[edge](T) -- (\x,\y-0.65);
    \draw[edge](T) -- (\x+0.5,\y-0.65);
\end{tikzpicture}\simeq\PP(X_1, X_2, X_3),~~~\text{then}~~~
\begin{tikzpicture}[baseline=\baseline]
    \tikzset{tensor/.style = {minimum size = 0.5cm,shape = circle,thick,draw=black,inner sep = 0pt}, edge/.style = {thick,line width=.4mm},every loop/.style={}}
    \def\y{0.25}
    \def\x{0}
    \node[tensor,fill=blue!20!white] (T) at (\x,\y) {$\Tt$};
    \draw  node[fill = black,circle,inner sep=0pt,minimum size=4pt] (m_1) at (\x-0.5,\y-0.65) {};
    \draw  node[fill = black,circle,inner sep=0pt,minimum size=4pt] (m_2) at (\x+0.5,\y-0.65) {};
    \draw[edge] (T) -- (m_1);
    \draw[edge](T) -- (\x,\y-0.65);
    \draw[edge](T) -- (\x+0.5,\y-0.65);
\end{tikzpicture}\simeq\PP(X_2).
\end{align*}
\end{enumerate}
下文为简便起见，用 $\PP(i_1,\cdots,i_N)$ 表示概率 $\PP(X_1=i_1,\cdots,X_N=i_N)$。 
可见，要在张量网络记法中表示边缘分布，只需用全 1 向量与相应的边收缩。
类似地，我们也能表示条件概率分布。
\paragraph{条件概率}
设 
$\PP(i_n \mid i_1, \dots, i_{n-1})
$
表示在前 $n-1$ 个随机变量取值 $i_1, \dots, i_{n-1}$ 的条件下 $X_n = i_n$ 的条件概率。  
对三阶概率张量 $\Tt\in \RR^{d_1 \times d_2 \times d_3}$，这些条件概率可用张量网络记法紧凑地表示：
\begin{enumerate}
    \item 
有 $
\PP(i_1|i_2) = \frac{\PP(i_1,i_2)}{\PP(i_2)} = \frac{\sum_{i_3}\PP(i_1,i_2,i_3)}{\sum_{i_1,i_3}\PP(i_1,i_2,i_3)} = 
\frac{\begin{tikzpicture}[baseline=\baseline]
    \tikzset{tensor/.style = {minimum size = 0.5cm,shape = circle,thick,draw=black,inner sep = 0pt}, edge/.style = {thick,line width=.4mm},every loop/.style={}}
    \def\x{0}
    \def\y{0}
    \node[tensor,fill=blue!20!white] (T) at (\x,\y) {$\Tt$};
    \draw  node[fill = black,circle,inner sep=0pt,minimum size=4pt] (m_1) at (\x+0.5,\y-0.65) {};
    \draw[edge] (T) -- (\x-0.5,\y-0.65)node [very near end,below]
    {\scalebox{0.7}{\indcolor{$i_1$}}};
    \draw[edge](T) -- (\x,\y-0.75)node [very near end,below]
    {\scalebox{0.7}{\indcolor{$i_2$}}};
    \draw[edge](T) -- (m_1);
\end{tikzpicture}}{\begin{tikzpicture}[baseline=\baseline]
    \tikzset{tensor/.style = {minimum size = 0.5cm,shape = circle,thick,draw=black,inner sep = 0pt}, edge/.style = {thick,line width=.4mm},every loop/.style={}}
    \def\x{0}
    \def\y{0}
    \node[tensor,fill=blue!20!white] (T) at (\x,\y) {$\Tt$};
    \draw  node[fill = black,circle,inner sep=0pt,minimum size=4pt] (m_1) at (\x-0.5,\y-0.65) {}; 
    \draw  node[fill = black,circle,inner sep=0pt,minimum size=4pt] (m_2) at (\x+0.5,\y-0.65) {};
    \draw[edge] (T) -- (m_1);
    \draw[edge](T) -- (m_2);
    \draw[edge](T) -- (\x,\y-0.75)node [very near end,below]
    {\scalebox{0.7}{\indcolor{$i_2$}}};
\end{tikzpicture}},
$
对所有 $i_1, i_2$ 成立。 

因此，对于固定的 $i_2$，条件分布 $\PP(X_1|X_2=i_2)\propto\begin{tikzpicture}[baseline=\baseline]
    \tikzset{tensor/.style = {minimum size = 0.5cm,shape = circle,thick,draw=black,inner sep = 0pt}, edge/.style = {thick,line width=.4mm},every loop/.style={}}
    \def\y{0.25}
    \def\x{0}
    \node[tensor,fill=blue!20!white] (T) at (\x,\y) {$\Tt$};
    \draw  node[fill = black,circle,inner sep=0pt,minimum size=4pt] (m_1) at (\x+0.5,\y-0.65) {};
    \draw[edge] (T) -- (\x-0.5,\y-0.65);
    \draw[edge](T) -- (\x,\y-0.75)node [very near end,below]
    {\scalebox{0.7}{\indcolor{$i_2$}}};
    \draw[edge](T) -- (m_1);
\end{tikzpicture}.$
\item 我们有 $\PP(i_3|i_1,i_2) = \frac{\PP(i_1,i_2,i_3)}{\PP(i_1,i_2)}=\frac{\PP(i_1,i_2,i_3)}{\sum_{i_3}\PP(i_1,i_2,i_3)}
=
\frac{\begin{tikzpicture}[baseline=\baseline]
    \tikzset{tensor/.style = {minimum size = 0.5cm,shape = circle,thick,draw=black,inner sep = 0pt}, edge/.style = {thick,line width=.4mm},every loop/.style={}}
    \def\x{0}
    \def\y{0}
    \node[tensor,fill=blue!20!white] (T) at (\x,\y) {$\Tt$};
    \draw[edge] (T) -- (\x-0.5,\y-0.65)node [very near end,below]
    {\scalebox{0.7}{\indcolor{$i_1$}}};
    \draw[edge](T) -- (\x,\y-0.75)node [very near end,below]
    {\scalebox{0.7}{\indcolor{$i_2$}}};
    \draw[edge](T) -- (\x+0.5,\y-0.65)node [very near end,below]
    {\scalebox{0.7}{\indcolor{$i_3$}}};
\end{tikzpicture}}{\begin{tikzpicture}[baseline=\baseline]
    \tikzset{tensor/.style = {minimum size = 0.5cm,shape = circle,thick,draw=black,inner sep = 0pt}, edge/.style = {thick,line width=.4mm},every loop/.style={}}
    \def\x{0}
    \def\y{0}
    \node[tensor,fill=blue!20!white] (T) at (\x,\y) {$\Tt$};
    \draw  node[fill = black,circle,inner sep=0pt,minimum size=4pt] (m_2) at (\x+0.5,\y-0.65) {};
    \draw[edge] (T) -- (\x-0.5,\y-0.65)node [very near end,below]
    {\scalebox{0.7}{\indcolor{$i_1$}}};
    \draw[edge](T) -- (m_2);
    \draw[edge](T) -- (\x,\y-0.75)node [very near end,below]
    {\scalebox{0.7}{\indcolor{$i_2$}}};
\end{tikzpicture}},
$
对任意 $i_1, i_2$ 和 $i_3$ 成立。

因此，对于固定的 $i_1$ 和 $ i_2$，条件分布 $\PP(X_3|X_1=i_1, X_2 = i_2)\propto\begin{tikzpicture}[baseline=\baseline]
    \tikzset{tensor/.style = {minimum size = 0.5cm,shape = circle,thick,draw=black,inner sep = 0pt}, edge/.style = {thick,line width=.4mm},every loop/.style={}}
    \def\y{0.25}
    \def\x{0}
    \node[tensor,fill=blue!20!white] (T) at (\x,\y) {$\Tt$};
    \draw[edge] (T) -- (\x-0.5,\y-0.65)node [very near end,below]
    {\scalebox{0.7}{\indcolor{$i_1$}}};
    \draw[edge](T) -- (\x,\y-0.75)node [very near end,below]
    {\scalebox{0.7}{\indcolor{$i_2$}}};
    \draw[edge](T) -- (\x+0.5,\y-0.65);
\end{tikzpicture}.$ 
\end{enumerate}
\subsubsection{概率分布的张量列参数化}
如第~\ref{chapter:tensor-decomposition} 章所见，处理高阶张量的计算代价很高，因为元素个数随张量的阶指数增长。表示联合概率分布的张量也是如此：张量表示的规模随随机变量的个数指数增长。解决这一问题的一个途径，是把概率张量参数化为低秩张量网络。

这里我们说明如何用张量列分解做到这一点。其他分解模型同样可行，但张量列格式尤为有趣，并已在量子物理中非常成功地用于建模多体系统~\citep{TensorNetworkOrg}。

设 $\Tt\in\RR^{d_1\times\cdots\times d_N}$ 为以张量列格式参数化的概率张量： 
\begin{center}
\begin{tikzpicture}[baseline=-0.5ex]
    \tikzset{tensor/.style = {minimum size = 0.5cm,shape = circle,thick,draw=black,inner sep = 0pt}, edge/.style = {thick,line width=.4mm},every loop/.style={}}
    \def\x{6}
    \def\y{0}
     \node[tensor,fill=green!50!red!50!white] (B) at (\x+1,\y) {$\Tt$};
    \draw[edge] (B) -- (\x+0.3,0) node [midway,above] {\scalebox{0.7}{\dimcolor{$d_1$}}};
    \draw[edge] (B) -- (\x+1.6,0) node [midway,above] {\scalebox{0.7}{\dimcolor{$d_4$}}};
    \draw[edge] (B) -- (\x+0.4,-0.3) node [midway, below] {\scalebox{0.7}{\dimcolor{$d_2$}}};
    \draw[edge] (B) -- (\x+1.5,-0.4)node [midway, below] {\scalebox{0.7}{\dimcolor{$d_3$}}};
\end{tikzpicture}
=
\begin{tikzpicture}[baseline=\baseline]
    \tikzset{tensor/.style = {minimum size = 0.5cm,shape = circle,thick,draw=black,inner sep = 0pt}, edge/.style   = {thick,line width=.4mm},every loop/.style={}}
    \def\y{0}
    \def\x{0}
  \node[tensor,fill=red!30!white] (D) at (\x-1,\y){$\Gt_1$};
  \node[tensor,fill=red!50!white] (A) at (\x,\y){\scalebox{0.85}{$\Gt_2$}};
  \node[tensor,fill=blue!50!white] (B) at (\x+1,\y){$\Gt_3$};
   \node[tensor,fill=blue!20!white] (C) at (\x+2,\y){$\Gt_4$};
    \draw[edge](A) -- (\x,\y-0.7)node [midway,left] {\edgelabel{$d_2$}};
    \draw[edge](B) -- (\x+1,\y-0.7)node [midway,left] {\edgelabel{$d_3$}};
    \draw[edge](C) -- (\x+2,\y-0.7)node [midway,left] {\edgelabel{$d_4$}};
     \draw[edge](D) -- (\x-1,\y-0.7)node [midway,left] {\edgelabel{$d_1$}};
    \draw[edge](D) -- (A) node [midway,above] {\edgelabel{$R_1$}};
    \draw[edge](A) -- (B) node [midway,above] {\edgelabel{$R_2$}};
    \draw[edge](B) -- (C) node [midway,above] {\edgelabel{$R_3$}};
\end{tikzpicture}.
\end{center}

由于 $\Tt$ 表示一个概率分布，它的元素必须非负，且总和为一。
在机器学习场景中，我们要针对某个损失函数~（如对数似然）来优化核心张量，此时可以对核心张量 $\Gt_1,\cdots,\Gt_N$ 施加若干约束以保证这两条性质。 
要保证 $\Tt$ 的元素非负，一个简单的办法是约束核心张量的元素非负，或者对所有核心张量的元素逐分量地施加一个取值非负的“激活”函数。另一种做法源于量子物理：假设 $\Tt$ 的元素是概率的平方根，而非概率本身： $\Tt_{i_1,\cdots,i_N} = \sqrt{\PP(X_1=i_1,\cdots,X_N=i_N)}$。当 $\Tt$ 处于 TT 格式时，这对应于如下张量网络：
$$
\PP(X_1=i_1,\cdots,X_4=i_4) = (\Tt_{i_1,\cdots,i_4})^2 = 
\begin{tikzpicture}[baseline=\baseline]
    \tikzset{tensor/.style = {minimum size = 0.5cm,shape = circle,thick,draw=black,inner sep = 0pt}, edge/.style   = {thick,line width=.4mm},every loop/.style={}}
    \def\y{-0.35}
    \def\x{0}
  \node[tensor,fill=red!30!white] (D) at (\x-1,\y){$\Gt_1$};
  \node[tensor,fill=red!50!white] (A) at (\x,\y){\scalebox{0.85}{$\Gt_2$}};
  \node[tensor,fill=blue!50!white] (B) at (\x+1,\y){$\Gt_3$};
   \node[tensor,fill=blue!20!white] (C) at (\x+2,\y){$\Gt_4$};
   \node[tensor,fill=red!30!white] (G1) at (\x-1,\y+0.85){$\Gt_1$};
  \node[tensor,fill=red!50!white] (G2) at (\x,\y+0.85){\scalebox{0.85}{$\Gt_2$}};
  \node[tensor,fill=blue!50!white] (G3) at (\x+1,\y+0.85){$\Gt_3$};
   \node[tensor,fill=blue!20!white] (G4) at (\x+2,\y+0.85){$\Gt_4$};
    \draw[edge](G1) -- (\x-1,\y+1.65)node [very near end, above]{\scalebox{0.7} {\indcolor{$i_1$}}};
    \draw[edge](G2) -- (\x,\y+1.65)node [very near end, above]{\scalebox{0.7} {\indcolor{$i_2$}}};
    \draw[edge](G3) -- (\x+1,\y+1.65)node [very near end, above]{\scalebox{0.7}{\indcolor{$i_3$}}};
    \draw[edge](G4) -- (\x+2,\y+1.65)node [very near end, above]{\scalebox{0.7}{\indcolor{$i_4$}}};
    \draw[edge](G1) -- (G2)node [midway,above] {\edgelabel{$R_1$}};
    \draw[edge](G2) -- (G3)node [midway,above] {\edgelabel{$R_2$}};
    \draw[edge](G3) -- (G4)node [midway,above] {\edgelabel{$R_3$}};
    \draw[edge](A) -- (\x,\y-0.7)node [very near end, below] {\scalebox{0.7}{\indcolor{$i_2$}}};
    \draw[edge](B) -- (\x+1,\y-0.7)node [very near end, below]{\scalebox{0.7} {\indcolor{$i_3$}}};
    \draw[edge](C) -- (\x+2,\y-0.7)node [very near end, below] {\scalebox{0.7}{\indcolor{$i_4$}}};
     \draw[edge](D) -- (\x-1,\y-0.7)node [very near end, below] {\scalebox{0.7}{\indcolor{$i_1$}}};
    \draw[edge](D) -- (A) node [midway,above] {\edgelabel{$R_1$}};
    \draw[edge](A) -- (B) node [midway,above] {\edgelabel{$R_2$}};
    \draw[edge](B) -- (C) node [midway,above] {\edgelabel{$R_3$}};
\end{tikzpicture}
.
$$

利用复制张量，我们可以将其写为：
\begin{align}\label{eqn:parametrized-prob}
    \PP(X_1,\cdots,X_4) \simeq
    \begin{tikzpicture}[baseline=\baseline]
    \tikzset{tensor/.style = {minimum size = 0.5cm,shape = circle,thick,draw=black,inner sep = 0pt}, edge/.style   = {thick,line width=.4mm},every loop/.style={}}
    \def\y{-0.35}
    \def\x{0}
  \node[tensor,fill=red!30!white] (D) at (\x-1,\y-0.25){$\Gt_1$};
  \node[tensor,fill=red!50!white] (A) at (\x,\y-0.25){\scalebox{0.85}{$\Gt_2$}};
  \node[tensor,fill=blue!50!white] (B) at (\x+1,\y-0.25){$\Gt_3$};
   \node[tensor,fill=blue!20!white] (C) at (\x+2,\y-0.25){$\Gt_4$};
   \node[tensor,fill=red!30!white] (G1) at (\x-1,\y+1){$\Gt_1$};
  \node[tensor,fill=red!50!white] (G2) at (\x,\y+1){\scalebox{0.85}{$\Gt_2$}};
  \node[tensor,fill=blue!50!white] (G3) at (\x+1,\y+1){$\Gt_3$};
   \node[tensor,fill=blue!20!white] (G4) at (\x+2,\y+1){$\Gt_4$};
    \draw[edge](G1) -- (G2)node [midway,above] {\edgelabel{$R_1$}};
    \draw[edge](G2) -- (G3)node [midway,above] {\edgelabel{$R_2$}};
    \draw[edge](G3) -- (G4)node [midway,above] {\edgelabel{$R_3$}};
    \draw[edge](D) -- (A) node [midway,above] {\edgelabel{$R_1$}};
    \draw[edge](A) -- (B) node [midway,above] {\edgelabel{$R_2$}};
    \draw[edge](B) -- (C) node [midway,above] {\edgelabel{$R_3$}};
    \draw  node[fill = black,circle,inner sep=0pt,minimum size=4pt] (m_1) at (\x-1,\y+0.35) {};
    \draw  node[fill = black,circle,inner sep=0pt,minimum size=4pt] (m_2) at (\x,\y+0.35) {};
    \draw  node[fill = black,circle,inner sep=0pt,minimum size=4pt] (m_3) at (\x+1,\y+0.35) {};
    \draw  node[fill = black,circle,inner sep=0pt,minimum size=4pt] (m_4) at (\x+2,\y+0.35) {};
    \draw[edge](D)--(G1);
    \draw[edge](A)--(G2);
    \draw[edge](B)--(G3);    \draw[edge](C)--(G4);
    \draw[edge](m_1) -- (\x-1.35,\y+0.35);
    \draw[edge](m_4) -- (\x+2.35,\y+0.35);
    \draw[edge](m_2) -- (\x-0.35,\y+0.35);
    \draw[edge](m_3) -- (\x+1.35,\y+0.35);
\end{tikzpicture}
\end{align}

为保证概率之和为一，我们可以选择建模未归一化的平方根概率，因为正如后文所见，当概率张量处于 TT 格式时，归一化常数的计算非常容易。于是我们假设
$$\PP(X_1=i_1,\cdots,X_N=i_N) = \frac{(\Tt_{i_1,\cdots,i_N})^2}{\zeta},$$
其中 $\zeta = \sum_{i_1,\cdots,i_N}(\Tt_{i_1,\cdots,i_N})^2$ 为归一化常数。与第~\ref{sec:efficient-operations} 节介绍的 TT 格式高效运算类似，该归一化因子同样可以在 TT 格式下高效计算： 
$$\zeta= \sum_{i_1,\cdots,i_4}\Tt_{i_1,\cdots,i_4}^2
=
\begin{tikzpicture}[baseline=\baseline]
    \tikzset{tensor/.style = {minimum size = 0.5cm,shape = circle,thick,draw=black,inner sep = 0pt}, edge/.style   = {thick,line width=.4mm},every loop/.style={}}
    \def\y{-0.35}
    \def\x{0}
  \node[tensor,fill=red!30!white] (D) at (\x-1,\y-0.25){$\Gt_1$};
  \node[tensor,fill=red!50!white] (A) at (\x,\y-0.25){\scalebox{0.85}{$\Gt_2$}};
    \node[tensor,fill=blue!50!white] (B) at (\x+1,\y-0.25){$\Gt_3$};
    \node[tensor,fill=blue!20!white] (C) at (\x+2,\y-0.25){$\Gt_4$};
     \node[tensor,fill=red!30!white] (G1) at (\x-1,\y+1){$\Gt_1$};
    \node[tensor,fill=red!50!white] (G2) at (\x,\y+1){\scalebox{0.85}{$\Gt_2$}};
    \node[tensor,fill=blue!50!white] (G3) at (\x+1,\y+1){$\Gt_3$};
    \node[tensor,fill=blue!20!white] (G4) at (\x+2,\y+1){$\Gt_4$};

    \draw[edge](G1) -- (G2)node [midway,above] {\edgelabel{$R_1$}};
    \draw[edge](G2) -- (G3)node [midway,above] {\edgelabel{$R_2$}};
    \draw[edge](G3) -- (G4)node [midway,above] {\edgelabel{$R_3$}};
    \draw[edge](D) -- (A) node [midway,above] {\edgelabel{$R_1$}};
    \draw[edge](A) -- (B) node [midway,above] {\edgelabel{$R_2$}};
    \draw[edge](B) -- (C) node [midway,above] {\edgelabel{$R_3$}};
    \draw[edge](G1)--(D);
    \draw[edge](G2)--(A);
    \draw[edge](G3)--(B);
    \draw[edge](G4)--(C);
\end{tikzpicture}.
$$
上述思想也被称为\emph{Born 机器}~\citep{glasser2019expressive,han2018unsupervised}。
与第~\ref{sub:probs} 节所示类似，TT 参数化的概率分布的边缘分布同样可以用张量网络来计算~（见式~\eqref{eqn:parametrized-prob}）：
\begin{enumerate}
    \item %
    利用全概率公式，$X_2,X_3$ 和 $X_4$ 的边缘分布可以表示为 
    \begin{align*}
    \PP(X_2=i_2, X_3=i_3, X_4=i_4) 
    &= \sum_{i_1=1}^{d_1}\PP(X_1=i_1,\cdots,X_4=i_4)\\ 
    &=
    \begin{tikzpicture}[baseline=\baseline]
    \tikzset{tensor/.style = {minimum size = 0.5cm,shape = circle,thick,draw=black,inner sep = 0pt}, edge/.style   = {thick,line width=.4mm},every loop/.style={}}
    \def\y{-0.35}
    \def\x{0}
  \node[tensor,fill=red!30!white] (D) at (\x-1,\y-0.25){$\Gt_1$};
  \node[tensor,fill=red!50!white] (A) at (\x,\y-0.25){\scalebox{0.85}{$\Gt_2$}};
  \node[tensor,fill=blue!50!white] (B) at (\x+1,\y-0.25){$\Gt_3$};
   \node[tensor,fill=blue!20!white] (C) at (\x+2,\y-0.25){$\Gt_4$};
   \node[tensor,fill=red!30!white] (G1) at (\x-1,\y+1){$\Gt_1$};
  \node[tensor,fill=red!50!white] (G2) at (\x,\y+1){\scalebox{0.85}{$\Gt_2$}};
  \node[tensor,fill=blue!50!white] (G3) at (\x+1,\y+1){$\Gt_3$};
   \node[tensor,fill=blue!20!white] (G4) at (\x+2,\y+1){$\Gt_4$};
    \draw[edge](G1) -- (G2)node [midway,above] {\edgelabel{$R_1$}};
    \draw[edge](G2) -- (G3)node [midway,above] {\edgelabel{$R_2$}};
    \draw[edge](G3) -- (G4)node [midway,above] {\edgelabel{$R_3$}};
    \draw[edge](D) -- (A) node [midway,above] {\edgelabel{$R_1$}};
    \draw[edge](A) -- (B) node [midway,above] {\edgelabel{$R_2$}};
    \draw[edge](B) -- (C) node [midway,above] {\edgelabel{$R_3$}};
    \draw  node[fill = black,circle,inner sep=0pt,minimum size=4pt] (m_1) at (\x-1,\y+0.35) {};
    \draw  node[fill = black,circle,inner sep=0pt,minimum size=4pt] (m_2) at (\x,\y+0.35) {};
    \draw  node[fill = black,circle,inner sep=0pt,minimum size=4pt] (m_3) at (\x+1,\y+0.35) {};
    \draw  node[fill = black,circle,inner sep=0pt,minimum size=4pt] (m_4) at (\x+2,\y+0.35) {};
    \draw[edge](D)--(G1);
    \draw[edge](A)--(G2);
    \draw[edge](B)--(G3);    \draw[edge](C)--(G4);
    \draw[edge](m_1) -- (\x-1.35,\y+0.35);
    \draw  node[fill = black,circle,inner sep=0pt,minimum size=4pt] (m) at (\x-1.35,\y+0.35) {};
    \draw[edge](m_4) -- (\x+2.35,\y+0.35) node [right] {\scalebox{0.7}{\indcolor{$i_4$}}};;
    \draw[edge](m_2) -- (\x-0.35,\y+0.35) node [left] {\scalebox{0.7}{\indcolor{$i_2$}}};;
    \draw[edge](m_3) -- (\x+1.35,\y+0.35) node [right] {\scalebox{0.7}{\indcolor{$i_3$}}};;
\end{tikzpicture}
=
\sum_{i_1}^{d_1}\Tt_{i_1,\cdots,i_4},
    \end{align*}
对任意 $i_2\in[d_1],i_3\in[d_3]$ 和 $i_4\in[d_4]$ 成立。 
\item 其余的边缘分布可以类似地表示。例如，$X_2$ 的边缘分布可以写为
\begin{align*}
    \PP(X_2=i_2) 
    &= \sum_{i_1,i_3,i_4=1}^{d_1,d_3,d_4}\PP(X_1=i_1,\cdots,X_4=i_4)\\ 
    &=
    \begin{tikzpicture}[baseline=\baseline]
    \tikzset{tensor/.style = {minimum size = 0.5cm,shape = circle,thick,draw=black,inner sep = 0pt}, edge/.style   = {thick,line width=.4mm},every loop/.style={}}
    \def\y{-0.35}
    \def\x{0}
  \node[tensor,fill=red!30!white] (D) at (\x-1,\y-0.25){$\Gt_1$};
  \node[tensor,fill=red!50!white] (A) at (\x,\y-0.25){\scalebox{0.85}{$\Gt_2$}};
  \node[tensor,fill=blue!50!white] (B) at (\x+1,\y-0.25){$\Gt_3$};
   \node[tensor,fill=blue!20!white] (C) at (\x+2,\y-0.25){$\Gt_4$};
   \node[tensor,fill=red!30!white] (G1) at (\x-1,\y+1){$\Gt_1$};
  \node[tensor,fill=red!50!white] (G2) at (\x,\y+1){\scalebox{0.85}{$\Gt_2$}};
  \node[tensor,fill=blue!50!white] (G3) at (\x+1,\y+1){$\Gt_3$};
   \node[tensor,fill=blue!20!white] (G4) at (\x+2,\y+1){$\Gt_4$};
    \draw[edge](G1) -- (G2)node [midway,above] {\edgelabel{$R_1$}};
    \draw[edge](G2) -- (G3)node [midway,above] {\edgelabel{$R_2$}};
    \draw[edge](G3) -- (G4)node [midway,above] {\edgelabel{$R_3$}};
    \draw[edge](D) -- (A) node [midway,above] {\edgelabel{$R_1$}};
    \draw[edge](A) -- (B) node [midway,above] {\edgelabel{$R_2$}};
    \draw[edge](B) -- (C) node [midway,above] {\edgelabel{$R_3$}};
    \draw  node[fill = black,circle,inner sep=0pt,minimum size=4pt] (m_1) at (\x-1,\y+0.35) {};
    \draw  node[fill = black,circle,inner sep=0pt,minimum size=4pt] (m_2) at (\x,\y+0.35) {};
    \draw  node[fill = black,circle,inner sep=0pt,minimum size=4pt] (m_3) at (\x+1,\y+0.35) {};
    \draw  node[fill = black,circle,inner sep=0pt,minimum size=4pt] (m_4) at (\x+2,\y+0.35) {};
    \draw[edge](D)--(G1);
    \draw[edge](A)--(G2);
    \draw[edge](B)--(G3);    
    \draw[edge](C)--(G4);
    \draw[edge](m_1) -- (\x-1.35,\y+0.35);
    \draw  node[fill = black,circle,inner sep=0pt,minimum size=4pt] (m_5) at (\x-1.35,\y+0.35) {};
    \draw[edge](m_4) -- (\x+2.35,\y+0.35);
    \draw  node[fill = black,circle,inner sep=0pt,minimum size=4pt] (m_6) at (\x+2.35,\y+0.35) {};
    \draw[edge](m_2) -- (\x-0.35,\y+0.35) node [left] {\scalebox{0.7}{\indcolor{$i_2$}}};;
    \draw[edge](m_3) -- (\x+1.35,\y+0.35);
    \draw  node[fill = black,circle,inner sep=0pt,minimum size=4pt] (m_7) at (\x+1.35,\y+0.35) {};
\end{tikzpicture}
=
\sum_{i_1,i_3,i_4=1}^{d_1,d_3,d_4}\Tt_{i_1,\cdots,i_4},
    \end{align*}
\end{enumerate}
\subsection{计算随机张量网络的期望} 
下面说明张量网络形式体系如何帮助我们推导随机张量的复杂性质。 
本节先给出两个非常有用的命题。
第一条只是期望的线性性的直接推论，但在本节中十分有用。
\begin{proposition}\label{prop:inner-expectation}
    设 $\At$ 与 $\Bt$ 为两个独立的随机张量网络。张量网络。则二者内积的期望为 
$$\EE\left(\begin{tikzpicture}[baseline=\baseline]
    \tikzset{tensor/.style = {minimum size = 0.5cm,shape = circle,thick,draw=black,inner sep = 0pt}, edge/.style = {thick,line width=.4mm},every loop/.style={}}
    \def\y{0}
    \def\x{0}
    \node[tensor,fill=red!50!blue!50!white] (At) at (\x,\y){\scalebox{0.85}{$\At$}};
    \node[tensor,fill=red!30!blue!50!white] (Bt) at (\x+1,\y){\scalebox{0.85}{$\Bt$}};
    \draw[edge] (At) -- (Bt);
    \draw[edge] (At) to [out=45,in=135] (Bt);
    \draw[edge] (At) to [out=-45,in=-135] (Bt);
    \draw[edge](At)--(\x-0.85,\y);
    \draw[edge](At)--(\x-0.75,\y+0.5);    \draw[edge](At)--(\x-0.75,\y-0.5);
    \draw[edge](Bt)--(\x+1.85,\y);
    \draw[edge](Bt)--(\x+1.75,\y+0.5);    \draw[edge](Bt)--(\x+1.75,\y-0.5);
\end{tikzpicture} \right)
=
\begin{tikzpicture}[baseline=\baseline]
    \tikzset{tensor/.style = {minimum size = 0.8cm,shape = circle,thick,draw=black,inner sep = 0pt}, edge/.style = {thick,line width=.4mm},every loop/.style={}}
    \def\y{0}
    \def\x{0}
    \node[tensor,fill=red!50!blue!50!white] (At) at (\x,\y){\scalebox{0.85}{$\EE(\At)$}};
    \node[tensor,fill=red!30!blue!50!white] (Bt) at (\x+1.5,\y){\scalebox{0.85}{$\EE(\Bt)$}};
    \draw[edge] (At) -- (Bt);
    \draw[edge] (At) to [out=45,in=135] (Bt);
    \draw[edge] (At) to [out=-45,in=-135] (Bt);
    \draw[edge](At)--(\x-0.85,\y);
    \draw[edge](At)--(\x-0.75,\y+0.5);    \draw[edge](At)--(\x-0.75,\y-0.5);
    \draw[edge](Bt)--(\x+2.35,\y);
    \draw[edge](Bt)--(\x+2.25,\y+0.5);    \draw[edge](Bt)--(\x+2.25,\y-0.5);
\end{tikzpicture}.
$$
\begin{proof}
结论直接由期望的线性性以及 $\At$ 与 $\Bt$ 的独立性得到。
\end{proof}
\end{proposition}
第二条给出一个有趣的性质：对于元素取自高斯分布的矩阵，其与自身的外积的期望归结为一个非常简单的张量网络。 
\begin{proposition}\label{prop:outer-prod}设 $\Ab\in\RR^{m\times n}$ 为随机矩阵，其元素独立同分布于标准正态分布。$\Ab$ 与自身的外积的期望为 
$$\EE\left(
\begin{tikzpicture}[baseline=\baseline]
    \tikzset{tensor/.style = {minimum size = 0.5cm,shape = circle,thick,draw=black,inner sep = 0pt}, edge/.style = {thick,line width=.4mm},every loop/.style={}}
    \def\y{0.15}
    \def\x{0}
    \node[tensor,fill=green!50!blue!20!white] (A) at (\x,\y){\scalebox{0.85}{$\Ab$}};
    \node[tensor,fill=green!50!blue!20!white] (B) at (\x+1.5,\y){\scalebox{0.85}{$\Ab$}};
    \draw[edge](A) -- (\x-0.55,\y-0.55);
    \node[draw=none] () at (\x-0.5,\y-0.3) {\scalebox{0.7}{\dimcolor{$m$}}};
    \draw[edge](A) -- (\x+0.55,\y-0.55);
    \node[draw=none] () at (\x+0.5,\y-0.3) {\scalebox{0.7}{\dimcolor{$n$}}};
    \draw[edge](B) -- (\x+1,\y-0.55);
    \node[draw=none] () at (\x+1,\y-0.3) {\scalebox{0.7}{\dimcolor{$m$}}};
    \draw[edge](B) -- (\x+2,\y-0.55);
    \node[draw=none] () at (\x+2,\y-0.3) {\scalebox{0.7}{\dimcolor{$n$}}};
\end{tikzpicture}
\right)
=
\begin{tikzpicture}[baseline=\baseline]
    \tikzset{tensor/.style = {minimum size = 0.5cm,shape = circle,thick,draw=black,inner sep = 0pt}, edge/.style = {thick,line width=.4mm},every loop/.style={}}
    \def\y{0}
    \def\x{0}
    \draw[edge] (\x+0.15,\y-0.1) to [out=100,in=80] (\x+1.85,\y-0.1);
    \draw[edge] (\x+1.45,\y-0.1) to [out=100,in=80] (\x+3,\y-0.1);
    \node[draw=none] () at (\x+2.25,\y+0.5) {\scalebox{0.6}{\dimcolor{$n$}}};
    \node[draw=none] () at (\x+1,\y+0.5) {\scalebox{0.6}{\dimcolor{$m$}}};
\end{tikzpicture},
$$
其中右端
是两个克罗内克 $\delta$ 的外积~（即 $
\begin{tikzpicture}[baseline=\baseline]
    \tikzset{tensor/.style = {minimum size = 0.5cm,shape = circle,thick,draw=black,inner sep = 0pt}, edge/.style = {thick,line width=.4mm},every loop/.style={}}
    \def\y{0}
    \def\x{0}
    \draw[edge] (\x+0.15,\y-0.1) to [out=100,in=80] (\x+1.85,\y-0.1);
    \draw[edge] (\x+1.45,\y-0.1) to [out=100,in=80] (\x+3,\y-0.1);
    \node[draw=none] () at (\x+0.15,\y-0.25) {\scalebox{0.7}{\indcolor{$i_1$}}};
    \node[draw=none] () at (\x+3,\y-0.25) {\scalebox{0.7}{\indcolor{$i_4$}}};
    \node[draw=none] () at (\x+1.85,\y-0.25) {\scalebox{0.7}{\indcolor{$i_3$}}};
    \node[draw=none] () at (\x+1.45,\y-0.25) {\scalebox{0.7}{\indcolor{$i_2$}}};
\end{tikzpicture}
=\delta_{i_1i_3}\delta_{i_2i_4}$，对任意 $i_1,i_3\in[m]$ 和 $i_2,i_4\in[n]$ 成立。）
\begin{proof}
由于 $\Ab$ 的元素独立取自 $\mathcal{N}(0,1)$，其向量化服从多元正态分布，$\vectorize(\Ab) \sim \mathcal{N}(\mathbf{0},\Ib_{mn})$，因而 $\EE(\vectorize(\Ab)\vectorize(\Ab)^\ts) = \Ib_{mn}$。这意味着 
$$ \EE(\Ab_{i_1i_2}\Ab_{i_3i_4})=
\begin{cases}
        1 & \text{ if } i_1=i_3\text{ and } i_2=i_4\\
        0 & \text{ otherwise, }
    \end{cases} 
    $$
于是 $\EE(\Ab\outprod\Ab)_{i_1i_2i_3i_4}=\EE(\Ab_{i_1i_2}\Ab_{i_3i_4}) = \delta_{i_1i_3}\delta_{i_2i_4}=
\begin{tikzpicture}[baseline=\baseline]
    \tikzset{tensor/.style = {minimum size = 0.5cm,shape = circle,thick,draw=black,inner sep = 0pt}, edge/.style = {thick,line width=.4mm},every loop/.style={}}
    \def\y{0}
    \def\x{0}
    \draw[edge] (\x+0.15,\y-0.1) to [out=100,in=80] (\x+1.85,\y-0.1);
    \draw[edge] (\x+1.45,\y-0.1) to [out=100,in=80] (\x+3,\y-0.1);
    \node[draw=none] () at (\x+0.15,\y-0.25) {\scalebox{0.7}{\indcolor{$i_1$}}};
    \node[draw=none] () at (\x+3,\y-0.25) {\scalebox{0.7}{\indcolor{$i_4$}}};
    \node[draw=none] () at (\x+1.85,\y-0.25) {\scalebox{0.7}{\indcolor{$i_3$}}};
    \node[draw=none] () at (\x+1.45,\y-0.25) {\scalebox{0.7}{\indcolor{$i_2$}}};
\end{tikzpicture}$，对任意 $i_1,i_3\in[m]$ 和 $i_2,i_4\in[n]$ 成立。


\end{proof}
\end{proposition}
\begin{remark}
    注意，$\EE(\Ab\outprod\Ab)$ 的指标不同排列顺序对应如下张量网络：  
    \begin{align*}
            \EE(\Ab\outprod\Ab)_{i_1i_2i_4i_3}
            &= 
            \EE(\Ab_{i_1i_2}\Ab_{i_4i_3}) 
    = \delta_{i_1i_4}\delta_{i_2i_3}\\
    &=
    \begin{cases}
        1 & \text{ if } i_1=i_4\text{ and } i_2=i_3\\
        0 & \text{ otherwise, }
    \end{cases} 
    =
    \begin{tikzpicture}[baseline=\baseline]
    \tikzset{tensor/.style = {minimum size = 0.5cm,shape = circle,thick,draw=black,inner sep = 0pt}, edge/.style = {thick,line width=.4mm},every loop/.style={}}
    \def\y{0}
    \def\x{0}
    \draw[edge] (\x,\y-0.1) to [out=100,in=80] (\x+2,\y-0.1);
    \draw[edge] (\x+0.15,\y-0.1) to [out=100,in=80] (\x+1.85,\y-0.1);
    \node[draw=none] () at (\x,\y-0.25) {\scalebox{0.7}{\indcolor{$i_1$}}};
    \node[draw=none] () at (\x+2.15,\y-0.25) {\scalebox{0.7}{\indcolor{$i_4$}}};
    \node[draw=none] () at (\x+0.35,\y-0.25) {\scalebox{0.7}{\indcolor{$i_2$}}};
    \node[draw=none] () at (\x+1.7,\y-0.25) {\scalebox{0.7}{\indcolor{$i_3$}}};
\end{tikzpicture},
    \end{align*}
以及
\begin{align*}
\EE(\Ab\outprod\Ab)_{i_1i_4i_2i_3} 
&= 
\EE(\Ab_{i_1i_4}\Ab_{i_2i_3})= \delta_{i_1i_2}\delta_{i_4i_3}\\
&=
\begin{cases}
        1 & \text{ if } i_1=i_2\text{ and } i_3=i_4\\
        0 & \text{ otherwise, }
\end{cases}
=
\begin{tikzpicture}[baseline=\baseline]
    \tikzset{tensor/.style = {minimum size = 0.5cm,shape = circle,thick,draw=black,inner sep = 0pt}, edge/.style = {thick,line width=.4mm},every loop/.style={}}
    \def\y{0}
    \def\x{0}
    \draw[edge] (\x+0.15,\y-0.1) to [out=100,in=80] (\x+1.85,\y-0.1);
    \draw[edge] (\x+2,\y-0.1) to [out=100,in=80] (\x+3.65,\y-0.1);
    \node[draw=none] () at (\x+0.15,\y-0.25) {\scalebox{0.7}{\indcolor{$i_1$}}};
    \node[draw=none] () at (\x+1.85,\y-0.25) {\scalebox{0.7}{\indcolor{$i_2$}}};
    \node[draw=none] () at (\x+2.15,\y-0.25) {\scalebox{0.7}{\indcolor{$i_3$}}};
    \node[draw=none] () at (\x+3.65,\y-0.25) {\scalebox{0.7}{\indcolor{$i_4$}}};
\end{tikzpicture}.
\end{align*}
\end{remark}
接下来说明如何借助命题~\ref{prop:inner-expectation} 与~\ref{prop:outer-prod}，为高斯随机矩阵与其自身乘积的期望给出一个极为简洁的张量网络推导。 
\begin{proposition}\label{prop:wishart-mean}
若随机矩阵 $\Ab\in\RR^{m\times n}$ 的元素独立同分布（i.i.d.）地取自标准正态分布，则 $\EE(\Ab^\ts\Ab) = m\Ib_n$。 
\begin{proof}
        借助命题~\ref{prop:inner-expectation} 与~\ref{prop:outer-prod}，可写为 
\begin{align*}
\EE(\Ab^\ts\Ab) 
&= 
\EE\left(\begin{tikzpicture}[baseline=\baseline]
    \tikzset{tensor/.style = {minimum size = 0.5cm,shape = circle,thick,draw=black,inner sep = 0pt}, edge/.style = {thick,line width=.4mm},every loop/.style={}}
    \def\y{0}
    \def\x{0}
    \node[tensor,fill=red!30!blue!30!white] (A1) at (\x,\y){\scalebox{0.85}{$\Ab$}};
    \node[tensor,fill=red!30!blue!30!white] (A2) at (\x+1,\y){\scalebox{0.85}{$\Ab$}};
    \draw[edge](A1) -- (A2) node [midway, above]{\scalebox{0.7}{\dimcolor{$m$}}};
    \draw[edge](A1) -- (\x-0.85,\y) node [midway, above]{\scalebox{0.7}{\dimcolor{$n$}}};
    \draw[edge](A2) -- (\x+1.85,\y) node [midway, above]{\scalebox{0.7}{\dimcolor{$n$}}};
\end{tikzpicture}\right)
=
\EE\left(\begin{tikzpicture}[baseline=\baseline]
    \tikzset{tensor/.style = {minimum size = 0.5cm,shape = circle,thick,draw=black,inner sep = 0pt}, edge/.style = {thick,line width=.4mm},every loop/.style={}}
    \def\y{0}
    \def\x{0}
    \node[tensor,fill=red!30!blue!30!white] (A2) at (\x,\y-0.5){\scalebox{0.85}{$\Ab$}};
    \node[tensor,fill=red!30!blue!30!white] (A1) at (\x,\y+0.5){\scalebox{0.85}{$\Ab$}};
    \draw[edge](\x+0.25,\y+0.4) -- (\x+1.25,\y+0.4) -- (\x+1.25,\y-0.4) -- (\x+0.25,\y-0.4);
    \draw[edge](\x+0.2,\y+0.65)--(\x+1.25,\y+0.65)node [midway, above]{\scalebox{0.7}{\dimcolor{$n$}}};
    \draw[edge](\x+0.2,\y-0.65)--(\x+1.25,\y-0.65)node [midway, below]{\scalebox{0.7}{\dimcolor{$n$}}};
    \node[draw=none] () at (\x+1.45,\y) {\scalebox{0.7}{\dimcolor{$m$}}};
    \draw[dotted](\x+0.55,\y+1)--(\x+0.55,\y-1);
\end{tikzpicture}\right)\\
\ \
&\textequal{Prop.~\ref{prop:inner-expectation}}
\ \
\begin{tikzpicture}[baseline=\baseline]
   \tikzset{tensor/.style = {minimum size = 0.5cm,shape = circle,thick,draw=black,inner sep = 0pt}, edge/.style = {thick,line width=.4mm},every loop/.style={}}
    \def\y{0}
    \def\x{0}
    \node[tensor,fill=red!30!blue!30!white] (A1) at (\x,\y+0.5){\scalebox{0.85}{$\Ab$}};
    \node[tensor,fill=red!30!blue!30!white] (A2) at (\x,\y-0.5){\scalebox{0.85}{$\Ab$}};
    \draw[edge,draw=blue](\x+0.25,\y+0.4) -- (\x+1.65,\y+0.4);
    \draw[edge, draw=red](\x+1.65,\y+0.4)-- (\x+1.95,\y+0.4) -- (\x+1.95,\y-0.4)--(\x+1.65,\y+0-0.4);
    \draw[edge, draw=blue](\x+1.65,\y-0.4) -- (\x+0.25,\y-0.4);
    \draw[edge, draw=blue](\x+0.2,\y+0.65)--(\x+1.95,\y+0.65)node [midway, above]{\scalebox{0.7}{\dimcolor{$n$}}};
    \draw[edge,draw=blue](\x+0.2,\y-0.65)--(\x+1.95,\y-0.65)node [midway, below]{\scalebox{0.7}{\dimcolor{$n$}}};
    \node[draw=none] () at (\x+2.15,\y) {\scalebox{0.7}{\dimcolor{$m$}}};
    \node at (\x-0.55,\y) {$\EE\left( \vphantom{\rule{0pt}{1cm}} \right.$};
    \node at (\x+0.55,\y) {$\left. \vphantom{\rule{0pt}{1cm}} \right)$};
    \node at (\x+1.15,\y) {$\EE\left( \vphantom{\rule{0pt}{1cm}} \right.$};
    \node at (\x+2.35,\y) {$\left. \vphantom{\rule{0pt}{1cm}} \right)$};
\end{tikzpicture}
~~\textequal{Prop.~\ref{prop:outer-prod}}
\ \ 
\begin{tikzpicture}[baseline=\baseline]
   \tikzset{tensor/.style = {minimum size = 0.5cm,shape = circle,thick,draw=black,inner sep = 0pt}, edge/.style = {thick,line width=.4mm},every loop/.style={}}
    \def\y{0}
    \def\x{0}
    \draw[edge,draw=blue](\x+1.35,\y+0.4) -- (\x+1.65,\y+0.4);
    \draw[edge, draw=red](\x+1.65,\y+0.4)-- (\x+1.95,\y+0.4) -- (\x+1.95,\y-0.4)--(\x+1.65,\y+0-0.4);
    \draw[edge, draw=blue](\x+1.35,\y-0.4) -- (\x+1.65,\y-0.4);
    \draw[edge, draw=blue](\x+1.35,\y-0.4) -- (\x+1.35,\y-0.4);
    \draw[edge, draw=blue](\x+1.35,\y+0.65)--(\x+1.95,\y+0.65);
    \draw[edge,draw=blue](\x+1.35,\y-0.65)--(\x+1.95,\y-0.65);
    \node[draw=none] () at (\x+2.15,\y) {\scalebox{0.7}{\dimcolor{$m$}}};
    \node[draw=none] () at (\x+1.35,\y) {\scalebox{0.7}{\dimcolor{$m$}}};
    \node[draw=none] () at (\x+0.85,\y) {\scalebox{0.7}{\dimcolor{$n$}}};
    \draw[edge, draw=blue](\x+1.35,\y+0.65)to[out=170,in=190](\x+1.35,\y-0.65);
    \draw[edge, draw=blue](\x+1.35,\y+0.4)to[out=190,in=170](\x+1.35,\y-0.4);
\end{tikzpicture}
= m\Ib_n.
\end{align*}
最后一个图由收缩长度为 $m$ 的腿得到。可以看到，一旦完成收缩，所得表达式就是单位矩阵的迹，对应上面等式中长度为 $m$ 的圆环。
\end{proof}
\end{proposition}
上述命题可解释为关于 Wishart 分布均值的陈述。  
回顾一下，从 Wishart 分布 $W_p(\mat \Sigma, m)$（其中 $\mat \Sigma \in \RR^{n\times n}$ 为协方差矩阵）采样矩阵 $\mat X$ 的做法是：把矩阵 $\mat A\in\mathbb R^{m\times n}$ 的每一列独立地取自多元正态 $\mathcal N(\vec 0 , \mat \Sigma)$，再令 $\mat X = \mat A^\top\mat A$。因此，命题~\ref{prop:inner-expectation} 中的矩阵 $\Ab ^\top \Ab$ 服从 Wishart 分布  $W_p(\mat I,m)$。
\begin{remark} 由命题~\ref{prop:inner-expectation} 可知，若 $\Ab \in \RR^{m \times n}$ 的元素取自标准正态分布，则 $\EE(\norm{\Ab}_F^2) = mn$。 
该结论可推广到元素独立同分布地取自标准正态分布的高阶张量，例如若 $\Tt\in\RR^{d_1\times d_2\times\cdots\times d_N}$，则 $\EE(\norm{\Tt}_F^2) = d_1d_2\cdots d_N$。
\end{remark}
利用前述命题给出的高斯随机矩阵性质，可以计算其乘积的期望范数。同样地，用张量网络表述后证明极为简洁，而若用指标与求和显式处理，则要繁琐得多。
\begin{proposition}\label{prop:expected-norm}
    设 $\Ab\in\RR^{m\times r}$ 与 $\Bb\in\RR^{r\times n}$ 为随机矩阵，其元素独立同分布（i.i.d.）于均值为零、方差为一的标准正态分布。则其乘积的 Frobenius 范数平方的期望为 $\EE({\norm{\Ab\Bb}^2_F}) = mnr$。
\begin{proof}
借助命题~\ref{prop:inner-expectation}，可写为
\begin{align*}
 \EE({\norm{\Ab\Bb}^2_F})
 &= 
 \EE\left(\begin{tikzpicture}[baseline=\baseline]
   \tikzset{tensor/.style = {minimum size = 0.5cm,shape = circle,thick,draw=black,inner sep = 0pt}, edge/.style = {thick,line width=.4mm},every loop/.style={}}
    \def\y{0}
    \def\x{0}
    \node[tensor,fill=red!50!white] (A) at (\x,\y+0.5){\scalebox{0.85}{$\Ab$}};
    \node[tensor,fill=blue!50!white] (B) at (\x+1,\y+0.5){\scalebox{0.85}{$\Bb$}};
    \draw[edge] (A) -- (B) node [midway,above]{\scalebox{0.7}{\dimcolor{$r$}}};
    \node[tensor,fill=red!50!white] (A1) at (\x,\y-0.5){\scalebox{0.85}{$\Ab$}};
    \node[tensor,fill=blue!50!white] (B1) at (\x+1,\y-0.5){\scalebox{0.85}{$\Bb$}};
    \draw[edge] (A1) -- (B1) node [midway,above]{\scalebox{0.7}{\dimcolor{$r$}}};
    \draw[edge] (A) to [out=-150,in=160, distance=0.5cm] (A1);
    \draw[edge] (B) to [out=-30,in=20, distance=0.5cm] (B1);
    \node[draw=none] () at (\x-0.75,\y) {\scalebox{0.7}{\dimcolor{$m$}}};
    \node[draw=none] () at (\x+1.75,\y) {\scalebox{0.7}{\dimcolor{$n$}}};
    \end{tikzpicture}\right)\\
&=
\EE\left(\begin{tikzpicture}[baseline=\baseline]
   \tikzset{tensor/.style = {minimum size = 0.5cm,shape = circle,thick,draw=black,inner sep = 0pt}, edge/.style = {thick,line width=.4mm},every loop/.style={}}
    \def\y{0.25}
    \def\x{0}
    \node[tensor,fill=red!50!white] (A) at (\x,\y){\scalebox{0.85}{$\Ab$}}; 
    \node[tensor,fill=red!50!white] (A1) at (\x+1,\y){\scalebox{0.85}{$\Ab$}};
    \draw[edge](A)--(A1)node[midway,above] {\scalebox{0.7}{\dimcolor{$m$}}};
    \draw[edge](A)--(\x,\y-0.85)node[midway,left] {\scalebox{0.7}{\dimcolor{$r$}}};
    \draw[edge](A1)--(\x+1,\y-0.85)node[midway,right] {\scalebox{0.7}{\dimcolor{$r$}}};
    \node[tensor,fill=blue!50!white] (B) at (\x+2,\y){\scalebox{0.85}{$\Bb$}}; 
    \node[tensor,fill=blue!50!white] (B1) at (\x+3,\y){\scalebox{0.85}{$\Bb$}};
    \draw[edge](B)--(B1)node[midway,above] {\scalebox{0.7}{\dimcolor{$n$}}};
    \draw[edge](B)--(\x+2,\y-0.85)node[midway,left] {\scalebox{0.7}{\dimcolor{$r$}}};
    \draw[edge](B1)--(\x+3,\y-0.85)node[midway,right] {\scalebox{0.7}{\dimcolor{$r$}}};
    \draw[edge](\x+2,\y-0.85) -- (\x+2,\y-1) -- (\x+1,\y-1) -- (\x+1,\y-0.85);
    \draw[edge](\x+3,\y-0.85) -- (\x+3,\y-1.25) -- (\x,\y-1.25) -- (\x,\y-0.85);
    \draw[dotted](\x+1.5,\y+0.5) -- (\x+1.5,\y-1.65);
\end{tikzpicture}\right)
=
\begin{tikzpicture}[baseline=\baseline]
   \tikzset{tensor/.style = {minimum size = 0.5cm,shape = circle,thick,draw=black,inner sep = 0pt}, edge/.style = {thick,line width=.4mm},every loop/.style={}}
    \def\y{0}
    \def\x{0}
    \node[tensor,fill=red!50!white] (A) at (\x,\y){\scalebox{0.85}{$\Ab$}}; 
    \node[tensor,fill=red!50!white] (A1) at (\x+1,\y){\scalebox{0.85}{$\Ab$}};
    \draw[edge, draw=red](A)--(A1)node[midway,above] {\scalebox{0.7}{\dimcolor{$m$}}};
   \node[tensor,fill=blue!50!white] (B) at (\x+2,\y){\scalebox{0.85}{$\Bb$}}; 
    \node[tensor,fill=blue!50!white] (B1) at (\x+3,\y){\scalebox{0.85}{$\Bb$}};
    \draw[edge, draw=blue](B)--(B1)node[midway,above] {\scalebox{0.7}{\dimcolor{$n$}}};
    \node at (\x-0.55,\y) {$\EE\Biggl ( $};
    \node (eq) at (\x+1.39,\y) {$\Biggr )\EE$};
    \node (eq) at (\x+1.73,\y) {$\Biggl($};
    \draw[edge, draw=blue](B) -- (\x+2,\y-0.85);
    \draw[edge, draw=blue](B1) -- (\x+3,\y-0.85);
    \draw[edge, draw=red](A) -- (\x,\y-0.85);
    \draw[edge, draw=red](A1) -- (\x+1,\y-0.85);
    \draw[edge](\x+1,\y-0.85) -- (\x+1,\y-1) -- (\x+2,\y-1)--(\x+2,\y-0.85);
    \draw[edge](\x+3,\y-0.85) -- (\x+3,\y-1.2) -- (\x,\y-1.2) -- (\x,\y-0.85);
    \node[draw=none] () at (\x+1.65,\y-0.85) {\scalebox{0.7}{\dimcolor{$r$}}};
    \node[draw=none] () at (\x+0.5,\y-0.95) {\scalebox{0.7}{\dimcolor{$r$}}};
\end{tikzpicture}\Biggr).
\end{align*}
由命题~\ref{prop:wishart-mean} 有 $\EE\left(\begin{tikzpicture}[baseline=\baseline]
   \tikzset{tensor/.style = {minimum size = 0.5cm,shape = circle,thick,draw=black,inner sep = 0pt}, edge/.style = {thick,line width=.4mm},every loop/.style={}}
    \def\y{0.15}
    \def\x{0}
    \node[tensor,fill=red!50!white] (A) at (\x,\y){\scalebox{0.85}{$\Ab$}}; 
    \node[tensor,fill=red!50!white] (A1) at (\x+1,\y){\scalebox{0.85}{$\Ab$}};
    \draw[edge, draw=red](A)--(A1)node[midway,above] {\scalebox{0.7}{\dimcolor{$m$}}};
    \draw[edge, draw=red](A) -- (\x,\y-0.85)node[midway,left] {\scalebox{0.7}{\dimcolor{$r$}}};
    \draw[edge, draw=red](A1) -- (\x+1,\y-0.85)node[midway,right] {\scalebox{0.7}{\dimcolor{$r$}}};
\end{tikzpicture}\right)
=
m\begin{tikzpicture}[baseline=\baseline]
   \tikzset{tensor/.style = {minimum size = 0.5cm,shape = circle,thick,draw=black,inner sep = 0pt}, edge/.style = {thick,line width=.4mm},every loop/.style={}}
    \def\y{0}
    \def\x{0}
     \draw[edge, draw=red] (\x,\y-0.1) to [out=100,in=80] (\x+2,\y-0.1);
    \node[draw=none] () at (\x+1,\y+0.65) {\scalebox{0.7}{\dimcolor{$r$}}};
\end{tikzpicture}$
以及 
$\EE\left(\begin{tikzpicture}[baseline=\baseline]
   \tikzset{tensor/.style = {minimum size = 0.5cm,shape = circle,thick,draw=black,inner sep = 0pt}, edge/.style = {thick,line width=.4mm},every loop/.style={}}
    \def\y{0.15}
    \def\x{0}
    \node[tensor,fill=blue!50!white] (B) at (\x,\y){\scalebox{0.85}{$\Bb$}}; 
    \node[tensor,fill=blue!50!white] (B1) at (\x+1,\y){\scalebox{0.85}{$\Bb$}};
    \draw[edge, draw=blue](B)--(B1)node[midway,above] {\scalebox{0.7}{\dimcolor{$n$}}};
    \draw[edge, draw=blue](B) -- (\x,\y-0.85)node[midway,left] {\scalebox{0.7}{\dimcolor{$r$}}};
    \draw[edge, draw=blue](B1) -- (\x+1,\y-0.85)node[midway,right] {\scalebox{0.7}{\dimcolor{$r$}}};
\end{tikzpicture}\right)
= 
n\begin{tikzpicture}[baseline=\baseline]
   \tikzset{tensor/.style = {minimum size = 0.5cm,shape = circle,thick,draw=black,inner sep = 0pt}, edge/.style = {thick,line width=.4mm},every loop/.style={}}
    \def\y{0}
    \def\x{0}
     \draw[edge, draw=blue] (\x,\y-0.1) to [out=100,in=80] (\x+2,\y-0.1);
    \node[draw=none] () at (\x+1,\y+0.65) {\scalebox{0.7}{\dimcolor{$r$}}};
\end{tikzpicture}.
$ 
由于长度为 $r$ 的腿之间存在收缩，最终得到
\begin{align*}
 \EE({\norm{\Ab\Bb}^2_F})
 = 
 \EE\left(\begin{tikzpicture}[baseline=\baseline]
   \tikzset{tensor/.style = {minimum size = 0.5cm,shape = circle,thick,draw=black,inner sep = 0pt}, edge/.style = {thick,line width=.4mm},every loop/.style={}}
    \def\y{0}
    \def\x{0}
    \node[tensor,fill=red!50!white] (A) at (\x,\y+0.5){\scalebox{0.85}{$\Ab$}};
    \node[tensor,fill=blue!50!white] (B) at (\x+1,\y+0.5){\scalebox{0.85}{$\Bb$}};
    \draw[edge] (A) -- (B) node [midway,above]{\scalebox{0.7}{\dimcolor{$r$}}};
    \node[tensor,fill=red!50!white] (A1) at (\x,\y-0.5){\scalebox{0.85}{$\Ab$}};
    \node[tensor,fill=blue!50!white] (B1) at (\x+1,\y-0.5){\scalebox{0.85}{$\Bb$}};
    \draw[edge] (A1) -- (B1) node [midway,above]{\scalebox{0.7}{\dimcolor{$r$}}};
    \draw[edge] (A) to [out=-150,in=160, distance=0.5cm] (A1);
    \draw[edge] (B) to [out=-30,in=20, distance=0.5cm] (B1);
    \node[draw=none] () at (\x-0.75,\y) {\scalebox{0.7}{\dimcolor{$m$}}};
    \node[draw=none] () at (\x+1.75,\y) {\scalebox{0.7}{\dimcolor{$n$}}};
    \end{tikzpicture}\right)
=
mn\  \begin{tikzpicture}     [baseline=\baseline]
 \tikzset{tensor/.style = {minimum size = 0.5cm,shape = circle,thick,draw=black,inner sep = 0pt}, edge/.style = {thick,line 
width=.4mm},every loop/.style={}}
    \def\y{0}
    \def\x{1}
    \draw[edge](\x+2,\y) -- (\x+2,\y-0.35) -- (\x+3,\y-0.35)--(\x+3,\y);
    \draw[edge](\x+1,\y) -- (\x+1,\y-0.5) -- (\x+4,\y-0.5) -- (\x+4,\y);
    \draw[edge, draw=red] (\x+2,\y) to [out=100,in=80](\x+1,\y);
    \draw[edge, draw=blue] (\x+3,\y) to [out=100,in=80](\x+4,\y);
    \node[draw=none] () at (\x+2.5,\y-0.65) {\scalebox{0.7}{\dimcolor{$r$}}};
    \node[draw=none] () at (\x+3.5,\y+0.5) {\scalebox{0.7}{\dimcolor{$r$}}};
    \node[draw=none] () at (\x+1.5,\y+0.5) {\scalebox{0.7}{\dimcolor{$r$}}};
\end{tikzpicture}
= mnr.
\end{align*}
\end{proof}
\end{proposition}
掌握了高斯随机矩阵乘积期望范数的计算方法后，便可把这一方法推广到更复杂的张量网络的期望。

\paragraph{例} 设有一个任意的张量网络，例如
$$\Tt = 
\begin{tikzpicture}[baseline=\baseline]
   \tikzset{tensor/.style = {minimum size = 0.5cm,shape = circle,thick,draw=black,inner sep = 0pt}, edge/.style = {thick,line width=.4mm},every loop/.style={}}
    \def\y{0}
    \def\x{0}
    \node[tensor,fill=red!20!blue!20!white] (A) at (\x,\y){\scalebox{0.85}{$\At$}}; 
    \node[tensor,fill=red!40!blue!20!white] (B) at (\x+1,\y){\scalebox{0.85}{$\Bt$}};
    \node[tensor,fill=red!40!blue!30!white] (C) at (\x+2,\y){\scalebox{0.85}{$\Ct$}};
    \node[tensor,fill=red!60!blue!40!white] (G) at (\x+1,\y+1){\scalebox{0.85}{$\Gt$}};
    \draw[edge](A)--(B)node[midway,above] {\scalebox{0.7}{\dimcolor{$r_1$}}};
    \draw[edge](B)--(C)node[midway,above] {\scalebox{0.7}{\dimcolor{$r_2$}}};
    \draw[edge](G)--(A)node[midway,left] {\scalebox{0.7}{\dimcolor{$d_1$}}};
    \draw[edge](G)--(B)node[midway,left] {\scalebox{0.7}{\dimcolor{$d_2$}}};
    \draw[edge](G)--(C)node[midway,right] {\scalebox{0.7}{\dimcolor{$d_3$}}};
    \draw[edge](A)--(\x,\y-0.85)node[midway,left] {\scalebox{0.7}{\dimcolor{$m_1$}}};
    \draw[edge](B)--(\x+1,\y-0.85)node[midway,left] {\scalebox{0.7}{\dimcolor{$m_2$}}};
    \draw[edge](C)--(\x+2,\y-0.85)node[midway,left] {\scalebox{0.7}{\dimcolor{$m_3$}}};
\end{tikzpicture},
$$
其中 $\At\in\RR^{m_1\times d_1\times r_1}, \Bt\in\RR^{m_2\times r_1\times d_2\times r_2}, \Ct\in\RR^{m_3\times r_2\times d_3}$ 与 $\Gt\in\RR^{d_1\times d_2\times d_3}$ 的元素均独立同分布地取自标准正态分布，我们要计算 $\norm{\Tt}^2_F$ 的期望。利用张量网络图，有
\begin{align*}
    \EE\left(\norm{\begin{tikzpicture}[baseline=\baseline]
   \tikzset{tensor/.style = {minimum size = 0.5cm,shape = circle,thick,draw=black,inner sep = 0pt}, edge/.style = {thick,line width=.4mm},every loop/.style={}}
    \def\y{0}
    \def\x{0}
    \node[tensor,fill=red!20!blue!20!white] (A) at (\x,\y){\scalebox{0.85}{$\At$}}; 
    \node[tensor,fill=red!40!blue!20!white] (B) at (\x+1,\y){\scalebox{0.85}{$\Bt$}};
    \node[tensor,fill=red!40!blue!30!white] (C) at (\x+2,\y){\scalebox{0.85}{$\Ct$}};
    \node[tensor,fill=red!60!blue!40!white] (G) at (\x+1,\y+1){\scalebox{0.85}{$\Gt$}};
    \draw[edge](A)--(B)node[midway,above] {\scalebox{0.7}{\dimcolor{$r_1$}}};
    \draw[edge](B)--(C)node[midway,above] {\scalebox{0.7}{\dimcolor{$r_2$}}};
    \draw[edge](G)--(A)node[midway,left] {\scalebox{0.7}{\dimcolor{$d_1$}}};
    \draw[edge](G)--(B)node[midway,left] {\scalebox{0.7}{\dimcolor{$d_2$}}};
    \draw[edge](G)--(C)node[midway,right] {\scalebox{0.7}{\dimcolor{$d_3$}}};
    \draw[edge](A)--(\x,\y-0.85)node[midway,left] {\scalebox{0.7}{\dimcolor{$m_1$}}};
    \draw[edge](B)--(\x+1,\y-0.85)node[midway,left] {\scalebox{0.7}{\dimcolor{$m_2$}}};
    \draw[edge](C)--(\x+2,\y-0.85)node[midway,left] {\scalebox{0.7}{\dimcolor{$m_3$}}};
\end{tikzpicture}}_F^2\right)
=
\EE\left(\begin{tikzpicture}[baseline=\baseline]
   \tikzset{tensor/.style = {minimum size = 0.5cm,shape = circle,thick,draw=black,inner sep = 0pt}, edge/.style = {thick,line width=.4mm},every loop/.style={}}
    \def\y{0.5}
    \def\x{0}
    \node[tensor,fill=red!20!blue!20!white] (A_1) at (\x,\y){\scalebox{0.85}{$\At$}}; 
    \node[tensor,fill=red!40!blue!20!white] (B_1) at (\x+1,\y){\scalebox{0.85}{$\Bt$}};
    \node[tensor,fill=red!40!blue!30!white] (C_1) at (\x+2,\y){\scalebox{0.85}{$\Ct$}};
    \node[tensor,fill=red!60!blue!40!white] (G_1) at (\x+1,\y+1){\scalebox{0.85}{$\Gt$}};
    \draw[edge](A_1)--(B_1)node[midway,above] {\scalebox{0.7}{\dimcolor{$r_1$}}};
    \draw[edge](B_1)--(C_1)node[midway,above] {\scalebox{0.7}{\dimcolor{$r_2$}}};
    \draw[edge](G_1)--(A_1)node[midway,left] {\scalebox{0.7}{\dimcolor{$d_1$}}};
    \draw[edge](G_1)--(B_1)node[midway,left] {\scalebox{0.7}{\dimcolor{$d_2$}}};
    \draw[edge](G_1)--(C_1)node[midway,right] {\scalebox{0.7}{\dimcolor{$d_3$}}};
    \node[tensor,fill=red!20!blue!20!white] (A) at (\x,\y-1){\scalebox{0.85}{$\At$}}; 
    \node[tensor,fill=red!40!blue!20!white] (B) at (\x+1,\y-1){\scalebox{0.85}{$\Bt$}};
    \node[tensor,fill=red!40!blue!30!white] (C) at (\x+2,\y-1){\scalebox{0.85}{$\Ct$}};
    \node[tensor,fill=red!60!blue!40!white] (G) at (\x+1,\y-2){\scalebox{0.85}{$\Gt$}};
    \draw[edge](A)--(B)node[midway,above] {\scalebox{0.7}{\dimcolor{$r_1$}}};
    \draw[edge](B)--(C)node[midway,above] {\scalebox{0.7}{\dimcolor{$r_2$}}};
    \draw[edge](G)--(A)node[midway,left] {\scalebox{0.7}{\dimcolor{$d_1$}}};
    \draw[edge](G)--(B)node[midway,left] {\scalebox{0.7}{\dimcolor{$d_2$}}};
    \draw[edge](G)--(C)node[midway,right] {\scalebox{0.7}{\dimcolor{$d_3$}}};
    \draw[edge](A_1)--(A)node[midway,left] {\scalebox{0.7}{\dimcolor{$m_1$}}};
    \draw[edge](B_1)--(B)node[midway,left] {\scalebox{0.7}{\dimcolor{$m_2$}}};
    \draw[edge](C_1)--(C)node[midway,left] {\scalebox{0.7}{\dimcolor{$m_3$}}};
\end{tikzpicture}\right).
\end{align*}
由命题~\ref{prop:wishart-mean} 可知 

$$
\EE\left(\begin{tikzpicture}[baseline=\baseline]
   \tikzset{tensor/.style = {minimum size = 0.5cm,shape = circle,thick,draw=black,inner sep = 0pt}, edge/.style = {thick,line width=.4mm},every loop/.style={}}
    \def\y{0.5}
    \def\x{0}
    \node[tensor,fill=red!20!blue!20!white] (A1) at (\x,\y){\scalebox{0.85}{$\At$}}; 
    \node[tensor,fill=red!20!blue!20!white] (A2) at (\x,\y-1){\scalebox{0.85}{$\At$}}; 
    \draw[edge, draw=blue](A1) -- (A2)node [midway,right]{\scalebox{0.7}{\dimcolor{$m_1$}}};
    \draw[edge, draw=blue](A1) -- (\x+0.65,\y)node [midway,above]{\scalebox{0.7}{\dimcolor{$r_1$}}};
    \draw[edge, draw=blue](A2) -- (\x+0.65,\y-1)node [midway,above]{\scalebox{0.7}{\dimcolor{$r_1$}}};
    \draw[edge,draw=blue](A1) -- (\x,\y+0.65)node [midway,right]{\scalebox{0.7}{\dimcolor{$d_1$}}};
    \draw[edge,draw=blue](A2) -- (\x,\y-1.65)node [midway,right]{\scalebox{0.7}{\dimcolor{$d_1$}}};
\end{tikzpicture}\right) = m_1
\begin{tikzpicture}[baseline=\baseline]
   \tikzset{tensor/.style = {minimum size = 0.5cm,shape = circle,thick,draw=black,inner sep = 0pt}, edge/.style = {thick,line width=.4mm},every loop/.style={}}
    \def\y{0}
    \def\x{0}
    \draw[edge, draw=blue](\x-0.65,\y+0.5) -- (\x-0.65,\y-0.5);
    \draw[edge, draw=blue](\x,\y+0.5) to [out=190,in=170] (\x,\y-0.5);
    \node[draw=none] () at (\x-0.8,\y) {\scalebox{0.7}{\dimcolor{$d_1$}}};
    \node[draw=none] () at (\x-0.1,\y) {\scalebox{0.7}{\dimcolor{$r_1$}}};
\end{tikzpicture},\  
\EE\left(\begin{tikzpicture}[baseline=\baseline]
   \tikzset{tensor/.style = {minimum size = 0.5cm,shape = circle,thick,draw=black,inner sep = 0pt}, edge/.style = {thick,line width=.4mm},every loop/.style={}}
    \def\y{0.5}
    \def\x{0}
    \node[tensor,fill=red!40!blue!20!white] (B1) at (\x,\y){\scalebox{0.85}{$\Bt$}};
    \node[tensor,fill=red!40!blue!20!white] (B2) at (\x,\y-1){\scalebox{0.85}{$\Bt$}};
    \draw[edge, draw=red](B1) -- (B2)node [midway,right]{\scalebox{0.7}{\dimcolor{$m_2$}}};
    \draw[edge, draw=red](B1) -- (\x,\y+0.65)node [midway,right]{\scalebox{0.7}{\dimcolor{$d_2$}}};
    \draw[edge, draw=red](B1) -- (\x+0.65,\y)node [midway,above]{\scalebox{0.7}{\dimcolor{$r_2$}}};
    \draw[edge, draw=red](B2) -- (\x+0.65,\y-1)node [midway,above]{\scalebox{0.7}{\dimcolor{$r_2$}}};
    \draw[edge, draw=red](B1) -- (\x-0.65,\y)node [midway,above]{\scalebox{0.7}{\dimcolor{$r_1$}}};
    \draw[edge, draw=red](B2) -- (\x-0.65,\y-1)node [midway,above]{\scalebox{0.7}{\dimcolor{$r_1$}}};
     \draw[edge, draw=red](B2) --(\x,\y-1.65)node [midway,right]{\scalebox{0.7}{\dimcolor{$d_2$}}};
    \end{tikzpicture}\right)
    =
m_2
\begin{tikzpicture}[baseline=\baseline]
   \tikzset{tensor/.style = {minimum size = 0.5cm,shape = circle,thick,draw=black,inner sep = 0pt}, edge/.style = {thick,line width=.4mm},every loop/.style={}}
    \def\y{0}
    \def\x{0}
    \draw[edge, draw=red](\x-0.65,\y+0.5) -- (\x-0.65,\y-0.5);
    \draw[edge, draw=red](\x,\y+0.5) to [out=190,in=170] (\x,\y-0.5);
    \draw[edge, draw=red](\x-1.3,\y+0.5) to [out=-20,in=-5] (\x-1.3,\y-0.5);
    \node[draw=none] () at (\x-0.8,\y) {\scalebox{0.7}{\dimcolor{$d_2$}}};
    0.5);
    \node[draw=none] () at (\x-1.2,\y) {\scalebox{0.7}{\dimcolor{$r_1$}}};
    \node[draw=none] () at (\x-0.1,\y) {\scalebox{0.7}{\dimcolor{$r_2$}}};
\end{tikzpicture}
\text{ and }
    \EE\left(\begin{tikzpicture}[baseline=\baseline]
   \tikzset{tensor/.style = {minimum size = 0.5cm,shape = circle,thick,draw=black,inner sep = 0pt}, edge/.style = {thick,line width=.4mm},every loop/.style={}}
    \def\y{0.5}
    \def\x{0}
    \node[tensor,fill=red!40!blue!30!white] (C1) at (\x,\y){\scalebox{0.85}{$\Ct$}};
    \node[tensor,fill=red!40!blue!30!white] (C2) at (\x,\y-1){\scalebox{0.85}{$\Ct$}};
    \draw[edge, draw=orange](C1) -- (C2)node [midway,right]{\scalebox{0.7}{\dimcolor{$m_3$}}};
    \draw[edge, draw=orange](C1) -- (\x,\y+0.65)node [midway,right]{\scalebox{0.7}{\dimcolor{$d_3$}}};
    \draw[edge, draw=orange](C1) -- (\x-0.65,\y)node [midway,above]{\scalebox{0.7}{\dimcolor{$r_2$}}};
    \draw[edge, draw=orange](C2) -- (\x-0.65,\y-1)node [midway,above]{\scalebox{0.7}{\dimcolor{$r_2$}}};
     \draw[edge, draw=orange](C2) --(\x,\y-1.65)node [midway,right]{\scalebox{0.7}{\dimcolor{$d_3$}}};
    \end{tikzpicture}\right)
=
m_3
\begin{tikzpicture}[baseline=\baseline]
   \tikzset{tensor/.style = {minimum size = 0.5cm,shape = circle,thick,draw=black,inner sep = 0pt}, edge/.style = {thick,line width=.4mm},every loop/.style={}}
    \def\y{0}
    \def\x{0}
    \draw[edge, draw=orange](\x-0.65,\y+0.5) -- (\x-0.65,\y-0.5);
    \draw[edge, draw=orange](\x-1.3,\y+0.5) to [out=-20,in=-5] (\x-1.3,\y-0.5);
    \node[draw=none] () at (\x-0.8,\y) {\scalebox{0.7}{\dimcolor{$d_3$}}};
    0.5);
    \node[draw=none] () at (\x-1.2,\y) {\scalebox{0.7}{\dimcolor{$r_2$}}};
\end{tikzpicture},
$$
于是由命题~\ref{prop:inner-expectation} 可写为：
\begin{align*}
\EE\left(\begin{tikzpicture}[baseline=\baseline]
   \tikzset{tensor/.style = {minimum size = 0.5cm,shape = circle,thick,draw=black,inner sep = 0pt}, edge/.style = {thick,line width=.4mm},every loop/.style={}}
    \def\y{0.5}
    \def\x{0}
    \node[tensor,fill=red!20!blue!20!white] (A_1) at (\x,\y){\scalebox{0.85}{$\At$}}; 
    \node[tensor,fill=red!40!blue!20!white] (B_1) at (\x+1,\y){\scalebox{0.85}{$\Bt$}};
    \node[tensor,fill=red!40!blue!30!white] (C_1) at (\x+2,\y){\scalebox{0.85}{$\Ct$}};
    \draw[edge](A_1)--(B_1)node[midway,above] {\scalebox{0.7}{\dimcolor{$r_1$}}};
    \draw[edge](B_1)--(C_1)node[midway,above] {\scalebox{0.7}{\dimcolor{$r_2$}}};
    \draw[edge](A_1) -- (\x+0.65,\y+0.65) node[midway,left] {\scalebox{0.7}{\dimcolor{$d_1$}}};
    \draw[edge](B_1) -- (\x+1,\y+0.65) node[midway,right] {\scalebox{0.7}{\dimcolor{$d_2$}}};
    \draw[edge](C_1) -- (\x+1.65,\y+0.65) node[midway,right] {\scalebox{0.7}{\dimcolor{$d_3$}}};  \node[tensor,fill=red!20!blue!20!white] (A) at (\x,\y-1){\scalebox{0.85}{$\At$}}; 
    \node[tensor,fill=red!40!blue!20!white] (B) at (\x+1,\y-1){\scalebox{0.85}{$\Bt$}};
    \node[tensor,fill=red!40!blue!30!white] (C) at (\x+2,\y-1){\scalebox{0.85}{$\Ct$}};
    \draw[edge](A)--(B) node[midway,above] {\scalebox{0.7}{\dimcolor{$r_1$}}};
    \draw[edge](B)--(C)node[midway,above] {\scalebox{0.7}{\dimcolor{$r_2$}}};
    \draw[edge](A_1)--(A) node[midway,left] {\scalebox{0.7}{\dimcolor{$m_1$}}};
    \draw[edge](B_1)--(B) node[midway,left] {\scalebox{0.7}{\dimcolor{$m_2$}}};
    \draw[edge](C_1)--(C) node[midway,left] {\scalebox{0.7}{\dimcolor{$m_3$}}};
    \draw[edge](A) -- (\x+0.65,\y-1.65)node[midway,left] {\scalebox{0.7}{\dimcolor{$d_1$}}};
    \draw[edge](B) -- (\x+1,\y-1.65)node[midway,right] {\scalebox{0.7}{\dimcolor{$d_2$}}};
    \draw[edge](C) -- (\x+1.65,\y-1.65)node[midway,right] {\scalebox{0.7}{\dimcolor{$d_3$}}};
\end{tikzpicture}\right)
&=
\begin{tikzpicture}[baseline=\baseline]
   \tikzset{tensor/.style = {minimum size = 0.5cm,shape = circle,thick,draw=black,inner sep = 0pt}, edge/.style = {thick,line width=.4mm},every loop/.style={}}
    \def\y{0.5}
    \def\x{0}
    \node[tensor,fill=red!20!blue!20!white] (A1) at (\x,\y){\scalebox{0.85}{$\At$}}; 
    \node[tensor,fill=red!20!blue!20!white] (A2) at (\x,\y-1){\scalebox{0.85}{$\At$}}; 
    \draw[edge, draw=blue](A1) -- (A2)node[midway,right] {\scalebox{0.7}{\dimcolor{$m_1$}}};
    \draw[edge, draw=blue](A1) -- (\x+0.65,\y)node[midway,above] {\scalebox{0.7}{\dimcolor{$r_1$}}};
    \draw[edge, draw=blue](A2) -- (\x+0.65,\y-1)node[midway,above] {\scalebox{0.7}{\dimcolor{$r_1$}}};
    \draw[edge,draw=blue](A1) -- (\x,\y+0.65)node[midway,right] {\scalebox{0.7}{\dimcolor{$d_1$}}};
    \draw[edge,draw=blue](A2) -- (\x,\y-1.65)node[midway,right] {\scalebox{0.7}{\dimcolor{$d_1$}}};
    \node at (\x-0.55,\y-0.5) {$\EE\left( \vphantom{\rule{0pt}{1cm}} \right.$};
    \node at (\x+1,\y-0.5) {$\left.\vphantom{\rule{0pt}{1cm}} \right)$};
    \node at (\x+1.45,\y-0.5){$\EE\left( \vphantom{\rule{0pt}{1cm}} \right.$};
    \node[tensor,fill=red!40!blue!20!white] (B1) at (\x+2,\y){\scalebox{0.85}{$\Bt$}};
    \node[tensor,fill=red!40!blue!20!white] (B2) at (\x+2,\y-1){\scalebox{0.85}{$\Bt$}};
    \draw[edge, draw=red](B1) -- (B2)node[midway,right] {\scalebox{0.7}{\dimcolor{$m_2$}}};
    \draw[edge, draw=red](B1) -- (\x+2.65,\y)node[midway,above] {\scalebox{0.7}{\dimcolor{$r_2$}}};
    \draw[edge, draw=red](B2) -- (\x+2.65,\y-1)node[midway,above] {\scalebox{0.7}{\dimcolor{$r_2$}}};
     \draw[edge, draw=red](B1) -- (\x+2,\y+0.65)node[midway,right] {\scalebox{0.7}{\dimcolor{$d_1$}}};
    \draw[edge, draw=red](B2) -- (\x+2,\y-1.65)node[midway,right] {\scalebox{0.7}{\dimcolor{$d_2$}}};
    \draw[edge, draw=red](B1) -- (\x+1.35,\y);
     \draw[edge, draw=red](B2) -- (\x+1.35,\y-1);
     \draw[edge](\x+1.35,\y) -- (\x+0.65,\y);
    \draw[edge](\x+1.35,\y-1) -- (\x+0.65,\y-1);
    \node at (\x+3,\y-0.5) {$\left.\vphantom{\rule{0pt}{1cm}} \right)$};
    \node at (\x+3.45,\y-0.5){$\EE\left( \vphantom{\rule{0pt}{1cm}} \right.$};
    \node[tensor,fill=red!40!blue!30!white] (C1) at (\x+4,\y){\scalebox{0.85}{$\Ct$}};
    \node[tensor,fill=red!40!blue!30!white] (C2) at (\x+4,\y-1){\scalebox{0.85}{$\Ct$}};
    \draw[edge, draw=orange](C1) -- (C2)node[midway,right] {\scalebox{0.7}{\dimcolor{$m_3$}}};
    \draw[edge, draw=orange](C1) -- (\x+3.25,\y);
     \draw[edge, draw=orange](C2) -- (\x+3.25,\y-1);
    \draw[edge, draw=orange](C1) --(\x+4,\y+0.65)node[midway,right] {\scalebox{0.7}{\dimcolor{$d_3$}}};
    \draw[edge, draw=orange](C2) -- (\x+4,\y-1.65)node[midway,right] {\scalebox{0.7}{\dimcolor{$d_3$}}};
    \draw[edge](\x+3.25,\y)--(\x+2.65,\y);
    \draw[edge](\x+3.25,\y-1)--(\x+2.65,\y-1);
    \node at (\x+4.5,\y-0.5) {$\left.\vphantom{\rule{0pt}{1cm}} \right)$};
\end{tikzpicture}\\
&= 
m_1m_2m_3
\begin{tikzpicture}[baseline=\baseline]
   \tikzset{tensor/.style = {minimum size = 0.5cm,shape = circle,thick,draw=black,inner sep = 0pt}, edge/.style = {thick,line width=.4mm},every loop/.style={}}
    \def\y{0}
    \def\x{0}
    \draw[edge, draw=blue](\x,\y+0.5) -- (\x,\y-0.5);
    \draw[edge, draw=blue](\x+0.5,\y+0.5) to [out=190,in=170] (\x+0.5,\y-0.5);
    \node[draw=none] () at (\x-0.15,\y) {\scalebox{0.7}{\dimcolor{$d_1$}}};
     \draw[edge, draw=red](\x+0.5,\y+0.5) to [out=-20,in=-5] (\x+0.5,\y-0.5);
    \node[draw=none] () at (\x+0.5,\y+0.65) {\scalebox{0.7}{\dimcolor{$r_1$}}};
    \node[draw=none] () at (\x+1.1,\y) {\scalebox{0.7}{\dimcolor{$d_2$}}};
    \draw[edge, draw=red](\x+1.25,\y+0.5) -- (\x+1.25,\y-0.5);
    \draw[edge, draw=red](\x+1.75,\y+0.5) to [out=190,in=170] (\x+1.75,\y-0.5);
    \draw[edge, draw=orange](\x+1.75,\y+0.5) to [out=-20,in=-5] (\x+1.75,\y-0.5);
    \node[draw=none] () at (\x+1.75,\y+0.65) {\scalebox{0.7}{\dimcolor{$r_2$}}};
    \draw[edge, draw=orange](\x+2.5,\y+0.5) -- (\x+2.5,\y-0.5);
    \node[draw=none] () at (\x+2.35,\y) {\scalebox{0.7}{\dimcolor{$d_3$}}};
  \end{tikzpicture}\\
&=
m_1m_2m_3r_1r_2
\begin{tikzpicture}[baseline=\baseline]
   \tikzset{tensor/.style = {minimum size = 0.5cm,shape = circle,thick,draw=black,inner sep = 0pt}, edge/.style = {thick,line width=.4mm},every loop/.style={}}
    \def\y{0}
    \def\x{0}
    \draw[edge, draw=blue](\x,\y+0.5) -- (\x,\y-0.5);
    \node[draw=none] () at (\x-0.15,\y) {\scalebox{0.7}{\dimcolor{$d_1$}}};
    \draw[edge, draw=red](\x+0.5,\y+0.5) -- (\x+0.5,\y-0.5);
    \node[draw=none] () at (\x+0.35,\y) {\scalebox{0.7}{\dimcolor{$d_2$}}};
    \draw[edge, draw=orange](\x+1,\y+0.5) -- (\x+1,\y-0.5);
    \node[draw=none] () at (\x+0.85,\y) {\scalebox{0.7}{\dimcolor{$d_3$}}};
\end{tikzpicture}.
\end{align*}
因此，由于 $\Gt$ 与 $\At,\Bt$ 及 $\Ct$ 相互独立，可再次使用命题~\ref{prop:inner-expectation}，得到 
\begin{align*}
    \EE\left(\begin{tikzpicture}[baseline=\baseline]
   \tikzset{tensor/.style = {minimum size = 0.5cm,shape = circle,thick,draw=black,inner sep = 0pt}, edge/.style = {thick,line width=.4mm},every loop/.style={}}
    \def\y{0.5}
    \def\x{0}
    \node[tensor,fill=red!20!blue!20!white] (A_1) at (\x,\y){\scalebox{0.85}{$\At$}}; 
    \node[tensor,fill=red!40!blue!20!white] (B_1) at (\x+1,\y){\scalebox{0.85}{$\Bt$}};
    \node[tensor,fill=red!40!blue!30!white] (C_1) at (\x+2,\y){\scalebox{0.85}{$\Ct$}};
    \node[tensor,fill=red!60!blue!40!white] (G_1) at (\x+1,\y+1){\scalebox{0.85}{$\Gt$}};
    \draw[edge](A_1)--(B_1)node[midway,above] {\scalebox{0.7}{\dimcolor{$r_1$}}};
    \draw[edge](B_1)--(C_1)node[midway,above] {\scalebox{0.7}{\dimcolor{$r_2$}}};
    \draw[edge](G_1)--(A_1)node[midway,left] {\scalebox{0.7}{\dimcolor{$d_1$}}};
    \draw[edge](G_1)--(B_1)node[midway,left] {\scalebox{0.7}{\dimcolor{$d_2$}}};
    \draw[edge](G_1)--(C_1)node[midway,right] {\scalebox{0.7}{\dimcolor{$d_3$}}};
    \node[tensor,fill=red!20!blue!20!white] (A) at (\x,\y-1){\scalebox{0.85}{$\At$}}; 
    \node[tensor,fill=red!40!blue!20!white] (B) at (\x+1,\y-1){\scalebox{0.85}{$\Bt$}};
    \node[tensor,fill=red!40!blue!30!white] (C) at (\x+2,\y-1){\scalebox{0.85}{$\Ct$}};
    \node[tensor,fill=red!60!blue!40!white] (G) at (\x+1,\y-2){\scalebox{0.85}{$\Gt$}};
    \draw[edge](A)--(B)node[midway,above] {\scalebox{0.7}{\dimcolor{$r_1$}}};
    \draw[edge](B)--(C)node[midway,above] {\scalebox{0.7}{\dimcolor{$r_2$}}};
    \draw[edge](G)--(A)node[midway,left] {\scalebox{0.7}{\dimcolor{$d_1$}}};
    \draw[edge](G)--(B)node[midway,left] {\scalebox{0.7}{\dimcolor{$d_2$}}};
    \draw[edge](G)--(C)node[midway,right] {\scalebox{0.7}{\dimcolor{$d_3$}}};
    \draw[edge](A_1)--(A)node[midway,left] {\scalebox{0.7}{\dimcolor{$m_1$}}};
    \draw[edge](B_1)--(B)node[midway,left] {\scalebox{0.7}{\dimcolor{$m_2$}}};
    \draw[edge](C_1)--(C)node[midway,left] {\scalebox{0.7}{\dimcolor{$m_3$}}};
\end{tikzpicture}\right)
&=
m_1m_2m_3r_1r_2
\EE\left(\begin{tikzpicture}[baseline=\baseline]
   \tikzset{tensor/.style = {minimum size = 0.5cm,shape = circle,thick,draw=black,inner sep = 0pt}, edge/.style = {thick,line width=.4mm},every loop/.style={}}
    \def\y{0.5}
    \def\x{0}
    \node[tensor,fill=red!60!blue!40!white] (G_1) at (\x+1,\y+1){\scalebox{0.85}{$\Gt$}};
    \node[tensor,fill=red!60!blue!40!white] (G) at (\x+1,\y-2){\scalebox{0.85}{$\Gt$}};
    \draw[edge](G_1)--(\x+1,\y+0.25);
    \draw[edge](G_1) -- (\x+0.5,\y+0.5);
    \draw[edge](G_1) -- (\x+1.5,\y+0.5);
    \draw[edge](G)--(\x+1,\y-1.25);
    \draw[edge](G)--(\x+0.5,\y-1.5);
    \draw[edge](G)--(\x+1.5,\y-1.5);
    \draw[edge, draw=blue](\x+0.5,\y+0.5)--(\x+0.5,\y-1.5)node[midway,left] {\scalebox{0.7}{\dimcolor{$d_1$}}};
    \draw[edge, draw=red](\x+1,\y+0.25)--(\x+1,\y-1.25)node[midway,right] {\scalebox{0.7}{\dimcolor{$d_2$}}};
     \draw[edge, draw=orange](\x+1.5,\y+0.5)--(\x+1.5,\y-1.5)node[midway,right] {\scalebox{0.7}{\dimcolor{$d_3$}}};
\end{tikzpicture}\right)
\ \ \ \textequal{Prop.~\ref{prop:wishart-mean}}\ \  m_1m_2m_3r_1r_2d_1d_2d_3.
\end{align*}
由此可见，借助张量网络可以高效地完成复杂计算。下面介绍涉及高斯随机矩阵的其他恒等式。
\begin{enumerate}
\item  第一个恒等式是 \emph{Isserlis} 定理~\citep{isserlis1918formula} 的一个有用推论。
\begin{theorem}\label{thm:isserlis}
设随机向量 $\ab\in\RR^{n}$ 的元素独立同分布地取自均值为零、方差为一的正态分布。则
\begin{align*}
\EE(\ab^{\outprod 4}) 
&= 
\EE\left(\begin{tikzpicture}[baseline=\baseline]
   \tikzset{tensor/.style = {minimum size = 0.5cm,shape = circle,thick,draw=black,inner sep = 0pt}, edge/.style = {thick,line width=.4mm},every loop/.style={}}
    \def\y{0.25}
    \def\x{0}
    \node[tensor,fill=red!20!blue!20!white] (a1) at (\x,\y){\scalebox{0.85}{$\ab$}}; 
    \node[tensor,fill=red!20!blue!20!white] (a2) at (\x+1,\y){\scalebox{0.85}{$\ab$}}; 
    \node[tensor,fill=red!20!blue!20!white] (a3) at (\x+2,\y){\scalebox{0.85}{$\ab$}}; 
    \node[tensor,fill=red!20!blue!20!white] (a4) at (\x+3,\y){\scalebox{0.85}{$\ab$}}; 
    \draw[edge](a1) -- (\x,\y-0.85);
    \draw[edge](a2) -- (\x+1,\y-0.85);
    \draw[edge](a3) -- (\x+2,\y-0.85);
    \draw[edge](a4) -- (\x+3,\y-0.85);
\end{tikzpicture}
\right)\\
&=
\begin{tikzpicture}[baseline=\baseline]
    \tikzset{tensor/.style = {minimum size = 0.5cm,shape = circle,thick,draw=black,inner sep = 0pt}, edge/.style = {thick,line width=.4mm},every loop/.style={}}
    \def\y{0}
    \def\x{0}
    \draw[edge] (\x+0.15,\y-0.1) to [out=100,in=80] (\x+1.85,\y-0.1);
    \draw[edge] (\x+2,\y-0.1) to [out=100,in=80] (\x+3.65,\y-0.1);
\end{tikzpicture}
+
\begin{tikzpicture}[baseline=\baseline]
    \tikzset{tensor/.style = {minimum size = 0.5cm,shape = circle,thick,draw=black,inner sep = 0pt}, edge/.style = {thick,line width=.4mm},every loop/.style={}}
    \def\y{0}
    \def\x{0}
    \draw[edge] (\x+0.15,\y-0.1) to [out=100,in=80] (\x+1.85,\y-0.1);
    \draw[edge] (\x+1.45,\y-0.1) to [out=100,in=80] (\x+3,\y-0.1);
\end{tikzpicture}
+
    \begin{tikzpicture}[baseline=\baseline]
    \tikzset{tensor/.style = {minimum size = 0.5cm,shape = circle,thick,draw=black,inner sep = 0pt}, edge/.style = {thick,line width=.4mm},every loop/.style={}}
    \def\y{0}
    \def\x{0}
    \draw[edge] (\x,\y-0.1) to [out=100,in=80] (\x+2,\y-0.1);
    \draw[edge] (\x+0.15,\y-0.1) to [out=100,in=80] (\x+1.85,\y-0.1);
\end{tikzpicture}
\end{align*}
其中 $\ab^{\outprod 4}$ 表示 $\ab$ 与自身作 4 次外积。最后一个等式中，求和的各项对应单位矩阵的不同重排。

该结果可推广到任意协方差结构：若 $\ab$ 取自多元正态 $\mathcal N(\mathbf 0, \Sigmab )$，则
\begin{align*}
\EE(\ab^{\outprod 4})=
\begin{tikzpicture}[baseline=\baseline]
    \tikzset{tensor/.style = {minimum size = 0.4cm,shape = circle,thick,draw=black,inner sep = 0pt}, edge/.style = {thick,line width=.4mm},every loop/.style={}}
    \def\y{0}`
    \def\x{0}
    \draw[edge] (\x+0.15,\y-0.1) to [out=100,in=80] (\x+1.85,\y-0.1);
    \draw[edge] (\x+2,\y-0.1) to [out=100,in=80] (\x+3.65,\y-0.1);
    \node[tensor,fill=red!50!white] (A) at (\x+0.9,\y+0.35){\scalebox{0.85}{$\Sigmab$}};
    \node[tensor,fill=red!50!white] (A) at (\x+2.8,\y+0.35){\scalebox{0.85}{$\Sigmab$}};
\end{tikzpicture}
+
\begin{tikzpicture}[baseline=\baseline]
    \tikzset{tensor/.style = {minimum size = 0.4cm,shape = circle,thick,draw=black,inner sep = 0pt}, edge/.style = {thick,line width=.4mm},every loop/.style={}}
    \def\y{0}`
    \def\x{0}
    \draw[edge] (\x+0.15,\y-0.1) to [out=100,in=80] (\x+1.85,\y-0.1);
    \draw[edge] (\x+1.45,\y-0.1) to [out=100,in=80] (\x+3,\y-0.1);
    \node[tensor,fill=red!50!white] (A) at (\x+0.9,\y+0.35){\scalebox{0.85}{$\Sigmab$}};
    \node[tensor,fill=red!50!white] (A) at (\x+2.25,\y+0.35){\scalebox{0.85}{$\Sigmab$}};
\end{tikzpicture}
+
\begin{tikzpicture}[baseline=\baseline]
    \tikzset{tensor/.style = {minimum size = 0.4cm,shape = circle,thick,draw=black,inner sep = 0pt}, edge/.style = {thick,line width=.4mm},every loop/.style={}}
    \def\y{0.3}`
    \def\x{0}
    \draw[edge] (\x,\y-0.1) to [out=100,in=80] (\x+2,\y-0.1);
    \draw[edge] (\x+0.15,\y-0.45) to [out=100,in=80] (\x+1.85,\y-0.45);
    \node[tensor,fill=red!50!white] (A) at (\x+0.9,\y+0.45){\scalebox{0.85}{$\Sigmab$}};
    \node[tensor,fill=red!50!white] (A) at (\x+0.95,\y){\scalebox{0.85}{$\Sigmab$}};
\end{tikzpicture} .
\end{align*}


\begin{proof}
    我们证明协方差为单位矩阵的情形。对任意协方差矩阵，证明可直接推广。 
    
    逐元素地，利用 Isserlis
    定理，并结合 $\ab\sim\mathcal{N}(\mathbf{0},\Ib)$ 这一事实，对 $i_1,\cdots,i_4\in[n]$ 有
\begin{align*}
(\EE(&\ab^{\outprod 4}))_{i_1,i_2,i_3,i_4} 
=
    \EE(\ab_{i_1}\ab_{i_2}\ab_{i_3}\ab_{i_4})\\
&= 
    \EE(\ab_{i_1}\ab_{i_2})\EE(\ab_{i_3}\ab_{i_4}) + \EE(\ab_{i_1}\ab_{i_3})\EE(\ab_{i_2}\ab_{i_4})
    +\EE(\ab_{i_1}\ab_{i_4})\EE(\ab_{i_2}\ab_{i_3})\\
&=
    \delta_{i_1i_2}\delta_{i_3i_4} + \delta_{i_1i_3}\delta_{i_2i_4} + \delta_{i_1i_4}\delta_{i_2i_3},
\end{align*}
用张量网络记法写作：
\begin{align*}
\EE(\ab^{\outprod 4})_{i_1i_2i_3i_4}
=
\begin{tikzpicture}[baseline=\baseline]
    \tikzset{tensor/.style = {minimum size = 0.5cm,shape = circle,thick,draw=black,inner sep = 0pt}, edge/.style = {thick,line width=.4mm},every loop/.style={}}
    \def\y{0}
    \def\x{0}
    \draw[edge] (\x+0.15,\y-0.1) to [out=100,in=80] (\x+1.85,\y-0.1);
    \draw[edge] (\x+2,\y-0.1) to [out=100,in=80] (\x+3.65,\y-0.1);
    \node[draw=none] () at (\x+0.15,\y-0.25) {\scalebox{0.7}{\indcolor{$i_1$}}};
    \node[draw=none] () at (\x+1.85,\y-0.25) {\scalebox{0.7}{\indcolor{$i_2$}}};
    \node[draw=none] () at (\x+2.15,\y-0.25) {\scalebox{0.7}{\indcolor{$i_3$}}};
    \node[draw=none] () at (\x+3.65,\y-0.25) {\scalebox{0.7}{\indcolor{$i_4$}}};
\end{tikzpicture}
+
\begin{tikzpicture}[baseline=\baseline]
    \tikzset{tensor/.style = {minimum size = 0.5cm,shape = circle,thick,draw=black,inner sep = 0pt}, edge/.style = {thick,line width=.4mm},every loop/.style={}}
    \def\y{0}
    \def\x{0}
    \draw[edge] (\x+0.15,\y-0.1) to [out=100,in=80] (\x+1.85,\y-0.1);
    \draw[edge] (\x+1.45,\y-0.1) to [out=100,in=80] (\x+3,\y-0.1);
    \node[draw=none] () at (\x+0.15,\y-0.25) {\scalebox{0.7}{\indcolor{$i_1$}}};
    \node[draw=none] () at (\x+3,\y-0.25) {\scalebox{0.7}{\indcolor{$i_4$}}};
    \node[draw=none] () at (\x+1.85,\y-0.25) {\scalebox{0.7}{\indcolor{$i_3$}}};
    \node[draw=none] () at (\x+1.45,\y-0.25) {\scalebox{0.7}{\indcolor{$i_2$}}};
\end{tikzpicture}
+
    \begin{tikzpicture}[baseline=\baseline]
    \tikzset{tensor/.style = {minimum size = 0.5cm,shape = circle,thick,draw=black,inner sep = 0pt}, edge/.style = {thick,line width=.4mm},every loop/.style={}}
    \def\y{0}
    \def\x{0}
    \draw[edge] (\x,\y-0.1) to [out=100,in=80] (\x+2,\y-0.1);
    \draw[edge] (\x+0.15,\y-0.1) to [out=100,in=80] (\x+1.85,\y-0.1);
    \node[draw=none] () at (\x,\y-0.25) {\scalebox{0.7}{\indcolor{$i_1$}}};
    \node[draw=none] () at (\x+2.15,\y-0.25) {\scalebox{0.7}{\indcolor{$i_4$}}};
    \node[draw=none] () at (\x+0.35,\y-0.25) {\scalebox{0.7}{\indcolor{$i_2$}}};
    \node[draw=none] () at (\x+1.7,\y-0.25) {\scalebox{0.7}{\indcolor{$i_3$}}};
\end{tikzpicture}.
\end{align*}
\end{proof}
\end{theorem}
也可以把 \emph{Isserlis 定理} 的结论推广到随机变量本身为高阶张量的情形。例如，对三阶张量 $\At\in\RR^{d_1\times d_2\times d_3}$，可写为
\begin{align*}
    (\EE(\At^{\outprod 4}))_{i_1j_1r_1i_2j_2r_2i_3j_3r_3i_4j_4r_4}
&=\\
&\hspace{-3cm}\begin{tikzpicture}[baseline=\baseline]
    \tikzset{tensor/.style = {minimum size = 0.5cm,shape = circle,thick,draw=black,inner sep = 0pt}, edge/.style = {thick,line width=.4mm},every loop/.style={}}
    \def\x{0}
    \def\y{0}
    \draw[edge] (\x+0.15,\y-0.1) to [out=100,in=80] (\x+1.85,\y-0.1);
    \draw[edge] (\x+2,\y-0.1) to [out=100,in=80] (\x+3.65,\y-0.1);
    \node[draw=none] () at (\x+0.15,\y-0.25) {\scalebox{0.7}{\indcolor{$j_1$}}};
    \node[draw=none] () at (\x+1.85,\y-0.25) {\scalebox{0.7}{\indcolor{$j_2$}}};
    \node[draw=none] () at (\x+2.15,\y-0.25) {\scalebox{0.7}{\indcolor{$j_3$}}};
    \node[draw=none] () at (\x+3.65,\y-0.25) {\scalebox{0.7}{\indcolor{$j_4$}}};
    \draw[edge] (\x+0.15,\y+1) to [out=100,in=80] (\x+1.85,\y+1);
    \draw[edge] (\x+2,\y+1) to [out=100,in=80] (\x+3.65,\y+1);
    \node[draw=none] () at (\x+0.15,\y+0.9) {\scalebox{0.7}{\indcolor{$i_1$}}};
    \node[draw=none] () at (\x+1.85,\y+0.9) {\scalebox{0.7}{\indcolor{$i_2$}}};
    \node[draw=none] () at (\x+2.15,\y+0.9) {\scalebox{0.7}{\indcolor{$i_3$}}};
    \node[draw=none] () at (\x+3.65,\y+0.9) {\scalebox{0.7}{\indcolor{$i_4$}}};
    \draw[edge] (\x+0.15,\y-1) to [out=100,in=80] (\x+1.85,\y-1);
    \draw[edge] (\x+2,\y-1) to [out=100,in=80] (\x+3.65,\y-1);
    \node[draw=none] () at (\x+0.15,\y-1.15) {\scalebox{0.7}{\indcolor{$r_1$}}};
    \node[draw=none] () at (\x+1.85,\y-1.15) {\scalebox{0.7}{\indcolor{$r_2$}}};
    \node[draw=none] () at (\x+2.15,\y-1.15) {\scalebox{0.7}{\indcolor{$r_3$}}};
    \node[draw=none] () at (\x+3.65,\y-1.15) {\scalebox{0.7}{\indcolor{$r_4$}}};
\end{tikzpicture}
+
\begin{tikzpicture}[baseline=\baseline]
    \tikzset{tensor/.style = {minimum size = 0.5cm,shape = circle,thick,draw=black,inner sep = 0pt}, edge/.style = {thick,line width=.4mm},every loop/.style={}}
    \def\x{0}
    \def\y{0}
    \draw[edge] (\x+0.15,\y-0.1) to [out=100,in=80] (\x+1.85,\y-0.1);
    \draw[edge] (\x+1.45,\y-0.1) to [out=100,in=80] (\x+3,\y-0.1);
    \node[draw=none] () at (\x+0.15,\y-0.25) {\scalebox{0.7}{\indcolor{$j_1$}}};
    \node[draw=none] () at (\x+3,\y-0.25) {\scalebox{0.7}{\indcolor{$j_4$}}};
    \node[draw=none] () at (\x+1.85,\y-0.25) {\scalebox{0.7}{\indcolor{$j_3$}}};
    \node[draw=none] () at (\x+1.45,\y-0.25) {\scalebox{0.7}{\indcolor{$j_2$}}};
    \draw[edge] (\x+0.15,\y+1) to [out=100,in=80] (\x+1.85,\y+1);
    \draw[edge] (\x+1.45,\y+1) to [out=100,in=80] (\x+3,\y+1);
    \node[draw=none] () at (\x+0.15,\y+0.9) {\scalebox{0.7}{\indcolor{$i_1$}}};
    \node[draw=none] () at (\x+3,\y+0.9) {\scalebox{0.7}{\indcolor{$i_4$}}};
    \node[draw=none] () at (\x+1.85,\y+0.9) {\scalebox{0.7}{\indcolor{$i_3$}}};
    \node[draw=none] () at (\x+1.45,\y+0.9) {\scalebox{0.7}{\indcolor{$i_2$}}};
    \draw[edge] (\x+0.15,\y-1) to [out=100,in=80] (\x+1.85,\y-1);
    \draw[edge] (\x+1.45,\y-1) to [out=100,in=80] (\x+3,\y-1);
    \node[draw=none] () at (\x+0.15,\y-1.15) {\scalebox{0.7}{\indcolor{$r_1$}}};
    \node[draw=none] () at (\x+3,\y-1.15) {\scalebox{0.7}{\indcolor{$r_4$}}};
    \node[draw=none] () at (\x+1.85,\y-1.15) {\scalebox{0.7}{\indcolor{$r_3$}}};
    \node[draw=none] () at (\x+1.45,\y-1.15) {\scalebox{0.7}{\indcolor{$r_2$}}};
\end{tikzpicture}
+
    \begin{tikzpicture}[baseline=\baseline]
    \tikzset{tensor/.style = {minimum size = 0.5cm,shape = circle,thick,draw=black,inner sep = 0pt}, edge/.style = {thick,line width=.4mm},every loop/.style={}}
    \def\x{0}
    \def\y{0}
    \draw[edge] (\x,\y+1) to [out=100,in=80] (\x+2,\y+1);
    \draw[edge] (\x+0.15,\y+1) to [out=100,in=80] (\x+1.85,\y+1);
    \node[draw=none] () at (\x,\y+0.9) {\scalebox{0.7}{\indcolor{$i_1$}}};
    \node[draw=none] () at (\x+2.15,\y+0.9) {\scalebox{0.7}{\indcolor{$i_4$}}};
    \node[draw=none] () at (\x+0.35,\y+0.9) {\scalebox{0.7}{\indcolor{$i_2$}}};
    \node[draw=none] () at (\x+1.7,\y+0.9) {\scalebox{0.7}{\indcolor{$i_3$}}};
    \draw[edge] (\x,\y-0.1) to [out=100,in=80] (\x+2,\y-0.1);
    \draw[edge] (\x+0.15,\y-0.1) to [out=100,in=80] (\x+1.85,\y-0.1);
    \node[draw=none] () at (\x,\y-0.25) {\scalebox{0.7}{\indcolor{$j_1$}}};
    \node[draw=none] () at (\x+2.15,\y-0.25) {\scalebox{0.7}{\indcolor{$j_4$}}};
    \node[draw=none] () at (\x+0.35,\y-0.25) {\scalebox{0.7}{\indcolor{$j_2$}}};
    \node[draw=none] () at (\x+1.7,\y-0.25) {\scalebox{0.7}{\indcolor{$j_3$}}};
    \draw[edge] (\x,\y-1) to [out=100,in=80] (\x+2,\y-1);
    \draw[edge] (\x+0.15,\y-1) to [out=100,in=80] (\x+1.85,\y-1);
    \node[draw=none] () at (\x,\y-1.25) {\scalebox{0.7}{\indcolor{$r_1$}}};
    \node[draw=none] () at (\x+2.15,\y-1.25) {\scalebox{0.7}{\indcolor{$r_4$}}};
    \node[draw=none] () at (\x+0.35,\y-1.25) {\scalebox{0.7}{\indcolor{$r_2$}}};
    \node[draw=none] () at (\x+1.7,\y-1.25) {\scalebox{0.7}{\indcolor{$r_3$}}};
\end{tikzpicture},
\end{align*}
对 $i_1,\cdots,i_4\in[d_1], j_1,\cdots,j_4\in [d_2]$ 以及 $r_1,\cdots,r_4\in[d_3]$ 成立。
\item 利用上述结果，当 $\Ab$ 为随机矩阵时，可轻松导出  $(\Ab^\ts\Ab)^{\outprod 2}$ 的期望：
\begin{proposition}\label{prop:kronecker-exp}
    设随机矩阵 $\Ab\in\RR^{m\times n}$ 的元素独立同分布地取自标准正态分布。则
\begin{align*}
\EE((\Ab^\ts\Ab)^{\outprod 2})) 
=
m^2\begin{tikzpicture}[baseline=\baseline]
    \tikzset{tensor/.style = {minimum size = 0.5cm,shape = circle,thick,draw=black,inner sep = 0pt}, edge/.style = {thick,line width=.4mm},every loop/.style={}}
    \def\y{0}
    \def\x{0}
    \draw[edge] (\x+0.15,\y-0.1) to [out=100,in=80] (\x+1.85,\y-0.1);
    \draw[edge] (\x+2,\y-0.1) to [out=100,in=80] (\x+3.65,\y-0.1);
    \node[draw=none] () at (\x+1,\y+0.5) {\scalebox{0.7}{\dimcolor{$n$}}};
    \node[draw=none] () at (\x+2.85,\y+0.5) {\scalebox{0.7}{\dimcolor{$n$}}};
\end{tikzpicture}
+m\begin{tikzpicture}[baseline=\baseline]
    \tikzset{tensor/.style = {minimum size = 0.5cm,shape = circle,thick,draw=black,inner sep = 0pt}, edge/.style = {thick,line width=.4mm},every loop/.style={}}
    \def\y{0}
    \def\x{0}
    \draw[edge] (\x+0.15,\y-0.1) to [out=100,in=80] (\x+1.85,\y-0.1);
    \draw[edge] (\x+1.45,\y-0.1) to [out=100,in=80] (\x+3,\y-0.1);
    \node[draw=none] () at (\x+1,\y+0.5) {\scalebox{0.7}{\dimcolor{$n$}}};
    \node[draw=none] () at (\x+2.25,\y+0.5) {\scalebox{0.7}{\dimcolor{$n$}}};
\end{tikzpicture}
+
m\begin{tikzpicture}[baseline=\baseline]
    \tikzset{tensor/.style = {minimum size = 0.5cm,shape = circle,thick,draw=black,inner sep = 0pt}, edge/.style = {thick,line width=.4mm},every loop/.style={}}
    \def\y{0}
    \def\x{0}
    \draw[edge] (\x,\y-0.1) to [out=100,in=80] (\x+2,\y-0.1);
    \draw[edge] (\x+0.15,\y-0.1) to [out=100,in=80] (\x+1.85,\y-0.1);
    \node[draw=none] () at (\x+1,\y+0.65) {\scalebox{0.7}{\dimcolor{$n$}}};
    \node[draw=none] () at (\x+1,\y+0.25) {\scalebox{0.7}{\dimcolor{$n$}}};
\end{tikzpicture}.
\end{align*}
\begin{proof}
将定理~\ref{thm:isserlis} 用于矩阵，并借助命题~\ref{prop:wishart-mean}，即可得到所需表达式。首先，由定理~\ref{thm:isserlis} 有
\begin{align*}
\EE(\Ab^{\outprod4})_{i_1j_1i_2j_2i_3j_3i_4j_4}
&= 
\EE\left(\begin{tikzpicture}[baseline=\baseline]
    \tikzset{tensor/.style = {minimum size = 0.5cm,shape = circle,thick,draw=black,inner sep = 0pt}, edge/.style = {thick,line width=.4mm},every loop/.style={}}
    \def\y{0.25}
    \def\x{0}
    \node[tensor,fill=green!50!blue!20!white] (A1) at (\x,\y){\scalebox{0.85}{$\Ab$}};
    \node[tensor,fill=green!50!blue!20!white] (A2) at (\x+1,\y){\scalebox{0.85}{$\Ab$}};
    \node[tensor,fill=green!50!blue!20!white] (A3) at (\x+2,\y){\scalebox{0.85}{$\Ab$}};
    \node[tensor,fill=green!50!blue!20!white] (A4) at (\x+3,\y){\scalebox{0.85}{$\Ab$}};
    \draw[edge](\x-0.15,\y-0.18) -- (\x-0.15,\y-0.65)node [below]{\scalebox{0.7}{\indcolor{$i_1$}}};
    \draw[edge](\x+0.15,\y-0.18) -- (\x+0.15,\y-0.65)node [below]{\scalebox{0.7}{\indcolor{$j_1$}}};
    \draw[edge](\x+0.85,\y-0.18) -- (\x+0.85,\y-0.65)node [below]{\scalebox{0.7}{\indcolor{$i_2$}}};
    \draw[edge](\x+1.15,\y-0.18) -- (\x+1.15,\y-0.65)node [below]{\scalebox{0.7}{\indcolor{$j_2$}}};
    \draw[edge](\x+1.85,\y-0.18) -- (\x+1.85,\y-0.65)node [below]{\scalebox{0.7}{\indcolor{$i_3$}}};
    \draw[edge](\x+2.15,\y-0.18) -- (\x+2.15,\y-0.65)node [below]{\scalebox{0.7}{\indcolor{$j_3$}}};
    \draw[edge](\x+2.85,\y-0.18) -- (\x+2.85,\y-0.65)node [below]{\scalebox{0.7}{\indcolor{$i_4$}}};
    \draw[edge](\x+3.15,\y-0.18) -- (\x+3.15,\y-0.65)node [below]{\scalebox{0.7}{\indcolor{$j_4$}}};
     \node[draw=none] () at (\x-0.35,\y-0.4) {\scalebox{0.7}{\dimcolor{$n$}}};
    \node[draw=none] () at (\x+0.35,\y-0.4) {\scalebox{0.7}{\dimcolor{$m$}}};
    \node[draw=none] () at (\x+0.65,\y-0.4) {\scalebox{0.7}{\dimcolor{$n$}}};
    \node[draw=none] () at (\x+1.35,\y-0.4) {\scalebox{0.7}{\dimcolor{$m$}}};
    \node[draw=none] () at (\x+1.65,\y-0.4) {\scalebox{0.7}{\dimcolor{$n$}}};
    \node[draw=none] () at (\x+2.35,\y-0.4) {\scalebox{0.7}{\dimcolor{$m$}}};
    \node[draw=none] () at (\x+2.65,\y-0.4) {\scalebox{0.7}{\dimcolor{$n$}}};
    \node[draw=none] () at (\x+3.35,\y-0.4) {\scalebox{0.7}{\dimcolor{$m$}}};
\end{tikzpicture}\right)\\
&\hspace{-1.5cm}=  
\begin{tikzpicture}[baseline=\baseline]
    \tikzset{tensor/.style = {minimum size = 0.5cm,shape = circle,thick,draw=black,inner sep = 0pt}, edge/.style = {thick,line width=.4mm},every loop/.style={}}
    \def\y{0.5}
    \def\x{0}
    \draw[edge] (\x+0.15,\y-0.1) to [out=100,in=80] (\x+1.85,\y-0.1);
    \draw[edge] (\x+2,\y-0.1) to [out=100,in=80] (\x+3.65,\y-0.1);
    \node[draw=none] () at (\x+0.15,\y-0.25) {\scalebox{0.7}{\indcolor{$i_1$}}};
    \node[draw=none] () at (\x+1.85,\y-0.25) {\scalebox{0.7}{\indcolor{$i_2$}}};
    \node[draw=none] () at (\x+2.15,\y-0.25) {\scalebox{0.7}{\indcolor{$i_3$}}};
    \node[draw=none] () at (\x+3.85,\y-0.25) {\scalebox{0.7}{\indcolor{$i_4$}}};
    \draw[edge](\x+0.15,\y-0.85) to [out=100,in=80] (\x+1.85,\y-0.85);
    \draw[edge] (\x+2,\y-0.85) to [out=100,in=80] (\x+3.65,\y-0.85);
    \node[draw=none] () at (\x+0.1,\y-0.95) {\scalebox{0.7}{\indcolor{$j_1$}}};
    \node[draw=none] () at (\x+1.8, \y-0.95) {\scalebox{0.7}{\indcolor{$j_2$}}};
    \node[draw=none] () at (\x+2.15,\y-0.95) {\scalebox{0.7}{\indcolor{$j_3$}}};
    \node[draw=none] () at (\x+3.85,\y-0.95) {\scalebox{0.7}{\indcolor{$j_4$}}};
\end{tikzpicture}
+
\begin{tikzpicture}[baseline=\baseline]
    \tikzset{tensor/.style = {minimum size = 0.5cm,shape = circle,thick,draw=black,inner sep = 0pt}, edge/.style = {thick,line width=.4mm},every loop/.style={}}
    \def\y{0.5}
    \def\x{0}
    \draw[edge] (\x+0.15,\y-0.1) to [out=100,in=80] (\x+1.85,\y-0.1);
    \draw[edge] (\x+1.45,\y-0.1) to [out=100,in=80] (\x+3,\y-0.1);
    \node[draw=none] () at (\x+0.15,\y-0.25) {\scalebox{0.7}{\indcolor{$i_1$}}};
    \node[draw=none] () at (\x+3,\y-0.25) {\scalebox{0.7}{\indcolor{$i_4$}}};
    \node[draw=none] () at (\x+1.85,\y-0.25) {\scalebox{0.7}{\indcolor{$i_3$}}};
    \node[draw=none] () at (\x+1.45,\y-0.25) {\scalebox{0.7}{\indcolor{$i_2$}}};
    \draw[edge] (\x+0.15,\y-0.85) to [out=100,in=80] (\x+1.85,\y-0.85);
    \draw[edge] (\x+1.45,\y-0.85) to [out=100,in=80] (\x+3,\y-0.85);
    \node[draw=none] () at (\x+0.1,\y-0.95) {\scalebox{0.7}{\indcolor{$j_1$}}};
    \node[draw=none] () at (\x+2.95,\y-0.95) {\scalebox{0.7}{\indcolor{$j_4$}}};
    \node[draw=none] () at (\x+1.8,\y-0.95) {\scalebox{0.7}{\indcolor{$j_3$}}};
    \node[draw=none] () at (\x+1.4,\y-0.95) {\scalebox{0.7}{\indcolor{$j_2$}}};
\end{tikzpicture}
+
 \begin{tikzpicture}[baseline=\baseline]
    \tikzset{tensor/.style = {minimum size = 0.5cm,shape = circle,thick,draw=black,inner sep = 0pt}, edge/.style = {thick,line width=.4mm},every loop/.style={}}
    \def\y{0.5}
    \def\x{0}
    \draw[edge] (\x,\y-0.1) to [out=100,in=80] (\x+2,\y-0.1);
    \draw[edge] (\x+0.15,\y-0.1) to [out=100,in=80] (\x+1.85,\y-0.1);
    \node[draw=none] () at (\x,\y-0.25) {\scalebox{0.7}{\indcolor{$i_1$}}};
    \node[draw=none] () at (\x+2.15,\y-0.25) {\scalebox{0.7}{\indcolor{$i_4$}}};
    \node[draw=none] () at (\x+0.35,\y-0.23) {\scalebox{0.7}{\indcolor{$i_2$}}};
    \node[draw=none] () at (\x+1.7,\y-0.25) {\scalebox{0.7}{\indcolor{$i_3$}}};
    \draw[edge] (\x,\y-0.85) to [out=100,in=80] (\x+2,\y-0.85);
    \draw[edge] (\x+0.15,\y-0.85) to [out=100,in=80] (\x+1.85,\y-0.85);
    \node[draw=none] () at (\x-0.03,\y-0.95) {\scalebox{0.7}{\indcolor{$j_1$}}};
    \node[draw=none] () at (\x+2.15,\y-0.95) {\scalebox{0.7}{\indcolor{$j_4$}}};
    \node[draw=none] () at (\x+0.35,\y-0.95) {\scalebox{0.7}{\indcolor{$j_2$}}};
    \node[draw=none] () at (\x+1.7,\y-0.95) {\scalebox{0.7}{\indcolor{$j_3$}}};
\end{tikzpicture},
\end{align*}
对 $i_1,\cdots, i_4\in[n]$ 与 $j_1,\cdots,j_4\in [m]$ 均成立。由此可得
\begin{align*}
\EE((\Ab^\ts\Ab)^{\outprod 2})_{i_1i_2i_3i_4}
&=
\EE\left(\begin{tikzpicture}[baseline=\baseline]
    \tikzset{tensor/.style = {minimum size = 0.5cm,shape = circle,thick,draw=black,inner sep = 0pt}, edge/.style = {thick,line width=.4mm},every loop/.style={}}
    \def\y{0.25}
    \def\x{0}
    \node[tensor,fill=green!50!blue!20!white] (A1) at (\x,\y){\scalebox{0.85}{$\Ab$}};
    \node[tensor,fill=green!50!blue!20!white] (A2) at (\x+1,\y){\scalebox{0.85}{$\Ab$}};
    \node[tensor,fill=green!50!blue!20!white] (A3) at (\x+2,\y){\scalebox{0.85}{$\Ab$}};
    \node[tensor,fill=green!50!blue!20!white] (A4) at (\x+3,\y){\scalebox{0.85}{$\Ab$}};
    \draw[edge](A1) -- (A2)node [midway, above]{\scalebox{0.7}{\dimcolor{$m$}}};
    \draw[edge](A3) -- (A4)node [midway, above]{\scalebox{0.7}{\dimcolor{$m$}}};
    \draw[edge](A1) -- (\x,\y-0.85)node [below]{\scalebox{0.7}{\indcolor{$i_1$}}};
    \draw[edge](A2) -- (\x+1,\y-0.85)node [below]{\scalebox{0.7}{\indcolor{$i_2$}}};
    \draw[edge](A3) -- (\x+2,\y-0.85)node [below]{\scalebox{0.7}{\indcolor{$i_3$}}};
    \draw[edge](A4) -- (\x+3,\y-0.85)node [below]{\scalebox{0.7}{\indcolor{$i_4$}}};
    \node[draw=none] () at (\x-0.25,\y-0.5) {\scalebox{0.7}{\dimcolor{$n$}}};
    \node[draw=none] () at (\x+1.25,\y-0.5) {\scalebox{0.7}{\dimcolor{$n$}}};
    \node[draw=none] () at (\x+1.75,\y-0.5) {\scalebox{0.7}{\dimcolor{$n$}}};
    \node[draw=none] () at (\x+3.25,\y-0.5) {\scalebox{0.7}{\dimcolor{$n$}}};
\end{tikzpicture}\right)\\
&=
\EE\left(\begin{tikzpicture}[baseline=\baseline]
    \tikzset{tensor/.style = {minimum size = 0.5cm,shape = circle,thick,draw=black,inner sep = 0pt}, edge/.style = {thick,line width=.4mm},every loop/.style={}}
    \def\y{0.25}
    \def\x{0}
    \node[tensor,fill=green!50!blue!20!white] (A1) at (\x,\y){\scalebox{0.85}{$\Ab$}};
    \node[tensor,fill=green!50!blue!20!white] (A2) at (\x+1,\y){\scalebox{0.85}{$\Ab$}};
    \node[tensor,fill=green!50!blue!20!white] (A3) at (\x+2,\y){\scalebox{0.85}{$\Ab$}};
    \node[tensor,fill=green!50!blue!20!white] (A4) at (\x+3,\y){\scalebox{0.85}{$\Ab$}};
    \draw[edge](\x-0.15,\y-0.18) -- (\x-0.15,\y-0.65)node [below]{\scalebox{0.7}{\indcolor{$i_1$}}};
    \draw[edge](\x+0.15,\y-0.18) -- (\x+0.15,\y-0.65);
    \draw[edge,draw=red](\x+0.15,\y-0.65) --(\x+0.15,\y-0.75) -- (\x+1.15,\y-0.75)--(\x+1.15,\y-0.65);
    \node[draw=none] () at (\x+0.4,\y-0.6) {\scalebox{0.7}{\dimcolor{$m$}}};
    \draw[edge](\x+0.85,\y-0.18) -- (\x+0.85,\y-0.65)node [below]{\scalebox{0.7}{\indcolor{$i_2$}}};
    \draw[edge](\x+1.15,\y-0.18) -- (\x+1.15,\y-0.65);
    \draw[edge](\x+1.85,\y-0.18) -- (\x+1.85,\y-0.65)node [below]{\scalebox{0.7}{\indcolor{$i_3$}}};
    \draw[edge](\x+2.15,\y-0.18) -- (\x+2.15,\y-0.65);
    \draw[edge,draw=red](\x+2.15,\y-0.65) --(\x+2.15,\y-0.75) -- (\x+3.15,\y-0.75)--(\x+3.15,\y-0.65);
    \node[draw=none] () at (\x+2.4,\y-0.6) {\scalebox{0.7}{\dimcolor{$m$}}};
    \draw[edge](\x+2.85,\y-0.18) -- (\x+2.85,\y-0.65)node [below]{\scalebox{0.7}{\indcolor{$i_4$}}};
    \draw[edge](\x+3.15,\y-0.18) -- (\x+3.15,\y-0.65);
    \node[draw=none] () at (\x-0.35,\y-0.45) {\scalebox{0.7}{\dimcolor{$n$}}};
    \node[draw=none] () at (\x+0.7,\y-0.45) {\scalebox{0.7}{\dimcolor{$n$}}};
    \node[draw=none] () at (\x+1.7,\y-0.45) {\scalebox{0.7}{\dimcolor{$n$}}};
    \node[draw=none] () at (\x+2.7,\y-0.45) {\scalebox{0.7}{\dimcolor{$n$}}};
\end{tikzpicture}\right)\\
&=
\begin{tikzpicture}[baseline=\baseline]
    \tikzset{tensor/.style = {minimum size = 0.5cm,shape = circle,thick,draw=black,inner sep = 0pt}, edge/.style = {thick,line width=.4mm},every loop/.style={}}
    \def\y{0.5}
    \def\x{0}
    \draw[edge] (\x+0.15,\y-0.1) to [out=100,in=80] (\x+1.85,\y-0.1);
    \draw[edge] (\x+2,\y-0.1) to [out=100,in=80] (\x+3.65,\y-0.1);
    \node[draw=none] () at (\x+0.15,\y-0.25) {\scalebox{0.7}{\indcolor{$i_1$}}};
    \node[draw=none] () at (\x+1.85,\y-0.25) {\scalebox{0.7}{\indcolor{$i_2$}}};
    \node[draw=none] () at (\x+2.15,\y-0.25) {\scalebox{0.7}{\indcolor{$i_3$}}};
    \node[draw=none] () at (\x+3.85,\y-0.25) {\scalebox{0.7}{\indcolor{$i_4$}}};
    \draw[edge](\x+0.15,\y-0.85) to [out=100,in=80] (\x+1.85,\y-0.85);
    \draw[edge, draw=red](\x+0.15,\y-0.85) to [out=-100,in=-80] (\x+1.85,\y-0.85);
    \draw[edge] (\x+2,\y-0.85) to [out=100,in=80] (\x+3.65,\y-0.85);
    \draw[edge, draw=red] (\x+2,\y-0.85) to [out=-100,in=-80] (\x+3.65,\y-0.85);
\end{tikzpicture}
+
\begin{tikzpicture}[baseline=\baseline]
    \tikzset{tensor/.style = {minimum size = 0.5cm,shape = circle,thick,draw=black,inner sep = 0pt}, edge/.style = {thick,line width=.4mm},every loop/.style={}}
    \def\y{0.5}
    \def\x{0}
    \draw[edge] (\x+0.15,\y-0.1) to [out=100,in=80] (\x+1.85,\y-0.1);
    \draw[edge] (\x+1.45,\y-0.1) to [out=100,in=80] (\x+3,\y-0.1);
    \node[draw=none] () at (\x+0.15,\y-0.25) {\scalebox{0.7}{\indcolor{$i_1$}}};
    \node[draw=none] () at (\x+3,\y-0.25) {\scalebox{0.7}{\indcolor{$i_4$}}};
    \node[draw=none] () at (\x+1.85,\y-0.25) {\scalebox{0.7}{\indcolor{$i_3$}}};
    \node[draw=none] () at (\x+1.45,\y-0.25) {\scalebox{0.7}{\indcolor{$i_2$}}};
    \draw[edge] (\x+0.15,\y-0.85) to [out=100,in=80] (\x+1.85,\y-0.85);
    \draw[edge, draw=red](\x+0.15,\y-0.85) to [out=-100,in=-80](\x+1.45,\y-0.85);
    \draw[edge] (\x+1.45,\y-0.85) to [out=100,in=80] (\x+3,\y-0.85);
    \draw[edge, draw=red](\x+1.85,\y-0.85) to [out=-100,in=-80](\x+3,\y-0.85);
\end{tikzpicture}
+
 \begin{tikzpicture}[baseline=\baseline]
    \tikzset{tensor/.style = {minimum size = 0.5cm,shape = circle,thick,draw=black,inner sep = 0pt}, edge/.style = {thick,line width=.4mm},every loop/.style={}}
    \def\y{0.5}
    \def\x{0}
    \draw[edge] (\x,\y-0.1) to [out=100,in=80] (\x+2,\y-0.1);
    \draw[edge] (\x+0.15,\y-0.1) to [out=100,in=80] (\x+1.85,\y-0.1);
    \node[draw=none] () at (\x,\y-0.25) {\scalebox{0.7}{\indcolor{$i_1$}}};
    \node[draw=none] () at (\x+2.15,\y-0.25) {\scalebox{0.7}{\indcolor{$i_4$}}};
    \node[draw=none] () at (\x+0.35,\y-0.23) {\scalebox{0.7}{\indcolor{$i_2$}}};
    \node[draw=none] () at (\x+1.7,\y-0.25) {\scalebox{0.7}{\indcolor{$i_3$}}};
    \draw[edge] (\x,\y-0.85) to [out=100,in=80] (\x+2,\y-0.85);
    \draw[edge, draw=red](\x,\y-0.85)to [out=-100,in=-80](\x+0.15,\y-0.85);
    \draw[edge] (\x+0.15,\y-0.85) to [out=100,in=80] (\x+1.85,\y-0.85);
    \draw[edge, draw=red](\x+1.85,\y-0.85)to [out=-100,in=-80](\x+2,\y-0.85);
\end{tikzpicture}\\
&=
m^2\begin{tikzpicture}[baseline=\baseline]
    \tikzset{tensor/.style = {minimum size = 0.5cm,shape = circle,thick,draw=black,inner sep = 0pt}, edge/.style = {thick,line width=.4mm},every loop/.style={}}
    \def\y{0}
    \def\x{0}
    \draw[edge] (\x+0.15,\y-0.1) to [out=100,in=80] (\x+1.85,\y-0.1);
    \draw[edge] (\x+2,\y-0.1) to [out=100,in=80] (\x+3.65,\y-0.1);
    \node[draw=none] () at (\x+0.15,\y-0.25) {\scalebox{0.7}{\indcolor{$i_1$}}};
    \node[draw=none] () at (\x+1.85,\y-0.25) {\scalebox{0.7}{\indcolor{$i_2$}}};
    \node[draw=none] () at (\x+2.15,\y-0.25) {\scalebox{0.7}{\indcolor{$i_3$}}};
    \node[draw=none] () at (\x+3.85,\y-0.25) {\scalebox{0.7}{\indcolor{$i_4$}}};
    \end{tikzpicture}
+m\begin{tikzpicture}[baseline=\baseline]
    \tikzset{tensor/.style = {minimum size = 0.5cm,shape = circle,thick,draw=black,inner sep = 0pt}, edge/.style = {thick,line width=.4mm},every loop/.style={}}
    \def\y{0}
    \def\x{0}
    \draw[edge] (\x+0.15,\y-0.1) to [out=100,in=80] (\x+1.85,\y-0.1);
    \draw[edge] (\x+1.45,\y-0.1) to [out=100,in=80] (\x+3,\y-0.1);
    \node[draw=none] () at (\x+0.15,\y-0.25) {\scalebox{0.7}{\indcolor{$i_1$}}};
    \node[draw=none] () at (\x+3,\y-0.25) {\scalebox{0.7}{\indcolor{$i_4$}}};
    \node[draw=none] () at (\x+1.85,\y-0.25) {\scalebox{0.7}{\indcolor{$i_3$}}};
    \node[draw=none] () at (\x+1.45,\y-0.25) {\scalebox{0.7}{\indcolor{$i_2$}}};
    \end{tikzpicture}
+
m \begin{tikzpicture}[baseline=\baseline]
    \tikzset{tensor/.style = {minimum size = 0.5cm,shape = circle,thick,draw=black,inner sep = 0pt}, edge/.style = {thick,line width=.4mm},every loop/.style={}}
    \def\y{0}
    \def\x{0}
    \draw[edge] (\x,\y-0.1) to [out=100,in=80] (\x+2,\y-0.1);
    \draw[edge] (\x+0.15,\y-0.1) to [out=100,in=80] (\x+1.85,\y-0.1);
    \node[draw=none] () at (\x,\y-0.25) {\scalebox{0.7}{\indcolor{$i_1$}}};
    \node[draw=none] () at (\x+2.15,\y-0.25) {\scalebox{0.7}{\indcolor{$i_4$}}};
    \node[draw=none] () at (\x+0.35,\y-0.23) {\scalebox{0.7}{\indcolor{$i_2$}}};
    \node[draw=none] () at (\x+1.7,\y-0.25) {\scalebox{0.7}{\indcolor{$i_3$}}};
    \end{tikzpicture}.
\end{align*}
\end{proof}
\end{proposition}
\item 更一般地，若 $\Ab\in\RR^{m\times n}$ 是元素独立同分布、服从标准正态分布的随机矩阵，而 $\Tt\in\RR^{n\times n\times n\times n}$ 是与 $\Ab$ 相互独立的任意随机张量，则把 $\Tt$ 与 $\Ab$ 收缩后得到
\begin{align*}
&\hspace*{-2cm}\EE\left(\begin{tikzpicture}[baseline=\baseline]
    \tikzset{tensor/.style = {minimum size = 0.5cm,shape = circle,thick,draw=black,inner sep = 0pt}, edge/.style = {thick,line width=.4mm},every loop/.style={}}
    \def\y{0.25}
    \def\x{0}
    \node[tensor,fill=green!50!blue!20!white] (A1) at (\x,\y){\scalebox{0.85}{$\Ab$}};
    \node[tensor,fill=green!50!blue!20!white] (A2) at (\x+1,\y){\scalebox{0.85}{$\Ab$}};
    \node[tensor,fill=green!50!blue!20!white] (A3) at (\x+2,\y){\scalebox{0.85}{$\Ab$}};
    \node[tensor,fill=green!50!blue!20!white] (A4) at (\x+3,\y){\scalebox{0.85}{$\Ab$}};
    \draw[edge](A1) -- (A2)node [midway, above]{\scalebox{0.7}{\dimcolor{$m$}}};
    \draw[edge](A3) -- (A4)node [midway, above]{\scalebox{0.7}{\dimcolor{$m$}}};
    \draw[edge](A1) -- (\x,\y-0.85)node[midway,left]{\scalebox{0.7}{\dimcolor{$n$}}};
    \draw[edge](A2) -- (\x+1,\y-0.85)node[midway,left]{\scalebox{0.7}{\dimcolor{$n$}}};
    \draw[edge](A3) -- (\x+2,\y-0.85)node[midway,left]{\scalebox{0.7}{\dimcolor{$n$}}};
    \draw[edge](A4) -- (\x+3,\y-0.85)node[midway,left]{\scalebox{0.7}{\dimcolor{$n$}}};
    \tikzset{tensor/.style = {minimum width = 3.5cm,minimum height = 0.6cm,shape = rounded rectangle,thick,draw=black,inner sep = 0pt}, edge/.style = {thick,line width=.4mm},every loop/.style={}}
    \def\x{0}
    \def\y{0}
    \node[tensor,fill=green!30!white] (At) at (\x+1.5,\y-0.85){$\Tt$};
    \end{tikzpicture}\right)= \sum_{i_1i_2i_3i_4}\EE(\Tt_{i_1i_2i_3i_4})(m^2\delta_{i_1i_2}\delta_{i_3i_4}+m\delta_{i_1i_3}\delta_{i_2i_4}+m\delta_{i_1i_4}\delta_{i_2i_3}).
\end{align*}
\item 设 $\Ab\in\RR^{m\times n}$ 与 $\Bb\in\RR^{n\times p}$ 为两个随机矩阵，其元素从标准正态分布中独立同分布地抽取，则
$$\EE(\Ab\Bb)^{\outprod 2} = 
\EE\left(\begin{tikzpicture}[baseline=\baseline]
    \tikzset{tensor/.style = {minimum size = 0.5cm,shape = circle,thick,draw=black,inner sep = 0pt}, edge/.style = {thick,line width=.4mm},every loop/.style={}}
    \def\y{0.25}
    \def\x{0}
    \node[tensor,fill=red!20!white] (A1) at (\x,\y){\scalebox{0.85}{$\Ab$}};
    \node[tensor,fill=blue!20!white] (B1) at (\x+1,\y){\scalebox{0.85}{$\Bb$}};
    \node[tensor,fill=red!20!white] (A2) at (\x+2,\y){\scalebox{0.85}{$\Ab$}};
    \node[tensor,fill=blue!20!white] (B2) at (\x+3,\y){\scalebox{0.85}{$\Bb$}};
    \draw[edge](A1) -- (B1)node [midway, above]{\scalebox{0.7}{\dimcolor{$n$}}};
    \draw[edge](A2) -- (B2)node [midway, above]{\scalebox{0.7}{\dimcolor{$n$}}};
    \draw[edge](A1) -- (\x,\y-0.85)node [midway, left]{\scalebox{0.7}{\dimcolor{$m$}}};
    \draw[edge](B1) -- (\x+1,\y-0.85)node [midway, right]{\scalebox{0.7}{\dimcolor{$p$}}};
    \draw[edge](A2) -- (\x+2,\y-0.85)node [midway, left]{\scalebox{0.7}{\dimcolor{$m$}}};
    \draw[edge](B2) -- (\x+3,\y-0.85)node [midway, right]{\scalebox{0.7}{\dimcolor{$p$}}};
\end{tikzpicture}\right)
=
n\begin{tikzpicture}[baseline=\baseline]
    \tikzset{tensor/.style = {minimum size = 0.5cm,shape = circle,thick,draw=black,inner sep = 0pt}, edge/.style = {thick,line width=.4mm},every loop/.style={}}
    \def\y{0}
    \def\x{0}
    \draw[edge] (\x+0.15,\y-0.1) to [out=100,in=80] (\x+1.85,\y-0.1);
    \draw[edge] (\x+1.45,\y-0.1) to [out=100,in=80] (\x+3,\y-0.1);
    \node[draw=none] () at (\x+1,\y+0.5) {\scalebox{0.7}{\dimcolor{$m$}}};
    \node[draw=none] () at (\x+2.25,\y+0.5) {\scalebox{0.7}{\dimcolor{$p$}}};
    \end{tikzpicture}.
$$
\begin{proof}
    利用命题~\ref{prop:inner-expectation} 与命题~\ref{prop:outer-prod} 可以写出：
\begin{align*}
\EE\left(\begin{tikzpicture}[baseline=\baseline]
    \tikzset{tensor/.style = {minimum size = 0.5cm,shape = circle,thick,draw=black,inner sep = 0pt}, edge/.style = {thick,line width=.4mm},every loop/.style={}}
    \def\y{0.25}
    \def\x{0}
    \node[tensor,fill=red!20!white] (A1) at (\x,\y){\scalebox{0.85}{$\Ab$}};
    \node[tensor,fill=blue!20!white] (B1) at (\x+1,\y){\scalebox{0.85}{$\Bb$}};
    \node[tensor,fill=red!20!white] (A2) at (\x+2,\y){\scalebox{0.85}{$\Ab$}};
    \node[tensor,fill=blue!20!white] (B2) at (\x+3,\y){\scalebox{0.85}{$\Bb$}};
    \draw[edge](A1) -- (B1)node [midway, above]{\scalebox{0.7}{\dimcolor{$n$}}};
    \draw[edge](A2) -- (B2)node [midway, above]{\scalebox{0.7}{\dimcolor{$n$}}};
    \draw[edge](A1) -- (\x,\y-0.85)node [midway,left]{\scalebox{0.7}{\dimcolor{$m$}}};
    \draw[edge](B1) -- (\x+1,\y-0.85)node [midway,right]{\scalebox{0.7}{\dimcolor{$p$}}};
    \draw[edge](A2) -- (\x+2,\y-0.85)node [midway,left]
    {\scalebox{0.7}{\dimcolor{$m$}}};
    \draw[edge](B2) -- (\x+3,\y-0.85)node [midway,right]{\scalebox{0.7}{\dimcolor{$p$}}};
\end{tikzpicture}\right)
&=
\EE\left(
\begin{tikzpicture}[baseline=\baseline]
    \tikzset{tensor/.style = {minimum size = 0.5cm,shape = circle,thick,draw=black,inner sep = 0pt}, edge/.style = {thick,line width=.4mm},every loop/.style={}}
    \def\y{0.25}
    \def\x{0}
    \node[tensor,fill=red!20!white] (A1) at (\x,\y){\scalebox{0.85}{$\Ab$}};
    \node[tensor,fill=red!20!white] (A2) at (\x+1,\y){\scalebox{0.85}{$\Ab$}};
    \node[tensor,fill=blue!20!white] (B1) at (\x+2,\y){\scalebox{0.85}{$\Bb$}};
    \node[tensor,fill=blue!20!white] (B2) at (\x+3,\y){\scalebox{0.85}{$\Bb$}};
    \draw[edge](\x-0.15,\y-0.2) -- (\x-0.15,\y-0.65)node [midway,left]
    {\scalebox{0.7}{\dimcolor{$m$}}};
    \draw[edge](\x+0.15,\y-0.2) -- (\x+0.15,\y-0.65)node [midway,right]
    {\scalebox{0.7}{\dimcolor{$n$}}};
    \draw[edge](\x+0.85,\y-0.2) -- (\x+0.85,\y-0.65)node [midway,left]
    {\scalebox{0.7}{\dimcolor{$m$}}};
    \draw[edge](\x+1.15,\y-0.2) -- (\x+1.15,\y-0.65)node [midway,right]
    {\scalebox{0.7}{\dimcolor{$n$}}};
    \draw[edge](\x+1.85,\y-0.2) -- (\x+1.85,\y-0.65)node [midway,left]
    {\scalebox{0.7}{\dimcolor{$n$}}};
    \draw[edge](\x+2.15,\y-0.2) -- (\x+2.15,\y-0.65)node [midway,right]
    {\scalebox{0.7}{\dimcolor{$p$}}};
    \draw[edge](\x+2.85,\y-0.2) -- (\x+2.85,\y-0.65)node [midway,left]
    {\scalebox{0.7}{\dimcolor{$n$}}};
    \draw[edge](\x+3.15,\y-0.2) -- (\x+3.15,\y-0.65)node [midway,right]
    {\scalebox{0.7}{\dimcolor{$p$}}};
    \draw[edge,draw=red](\x+0.15,\y-0.65)--(\x+0.15,\y-0.75) -- (\x+1.85,\y-0.75)--(\x+1.85,\y-0.65);
    \draw[edge,draw=red](\x+1.15,\y-0.65)--(\x+1.15,\y-0.85) -- (\x+2.85,\y-0.85)--(\x+2.85,\y-0.65);
\end{tikzpicture}\right)\\
&\textequal{Prop.~\ref{prop:inner-expectation}}\ \
  \begin{tikzpicture}[baseline=\baseline]
    \tikzset{tensor/.style = {minimum size = 0.5cm,shape = circle,thick,draw=black,inner sep = 0pt}, edge/.style = {thick,line width=.4mm},every loop/.style={}}
    \def\y{0.15}
    \def\x{0}
    \node[tensor,fill=red!20!white] (A1) at (\x,\y){\scalebox{0.85}{$\Ab$}};
    \node[tensor,fill=red!20!white] (A2) at (\x+1,\y){\scalebox{0.85}{$\Ab$}};
    \node[tensor,fill=blue!20!white] (B1) at (\x+2.5,\y){\scalebox{0.85}{$\Bb$}};
    \node[tensor,fill=blue!20!white] (B2) at (\x+3.5,\y){\scalebox{0.85}{$\Bb$}};
    \draw[edge](\x-0.15,\y-0.2) -- (\x-0.15,\y-0.65)node [midway,left]
    {\scalebox{0.7}{\dimcolor{$m$}}};
    \draw[edge](\x+0.15,\y-0.2) -- (\x+0.15,\y-0.65)node [midway,right]
    {\scalebox{0.7}{\dimcolor{$n$}}};
    \draw[edge](\x+0.85,\y-0.2) -- (\x+0.85,\y-0.65)node [midway,left]
    {\scalebox{0.7}{\dimcolor{$m$}}};
    \draw[edge](\x+1.15,\y-0.2) -- (\x+1.15,\y-0.65)node [midway,right]
    {\scalebox{0.7}{\dimcolor{$n$}}};
    \draw[edge](\x+2.35,\y-0.2) -- (\x+2.35,\y-0.65)node [midway,left]
    {\scalebox{0.7}{\dimcolor{$n$}}};
    \draw[edge](\x+2.65,\y-0.2) -- (\x+2.65,\y-0.65)node [midway,right]
    {\scalebox{0.7}{\dimcolor{$p$}}};
    \draw[edge](\x+3.35,\y-0.2) -- (\x+3.35,\y-0.65)node [midway,left]
    {\scalebox{0.7}{\dimcolor{$n$}}};
    \draw[edge](\x+3.65,\y-0.2) -- (\x+3.65,\y-0.65)node [midway,right]
    {\scalebox{0.7}{\dimcolor{$p$}}};
    \draw[edge,draw=red](\x+0.15,\y-0.65)--(\x+0.15,\y-0.75) -- (\x+2.35,\y-0.75)--(\x+2.35,\y-0.65);
    \draw[edge,draw=red](\x+1.15,\y-0.65)--(\x+1.15,\y-0.85) -- (\x+3.35,\y-0.85)--(\x+3.35,\y-0.65);
    \node at (\x-0.65,\y-0.1) {$\EE\left( \vphantom{\rule{0pt}{1cm}} \right.$};
    \node at (\x+1.45,\y-0.1) {$ \left.\vphantom{\rule{0pt}{1cm}} \right)$};
    \node at (\x+1.95,\y-0.15){$\EE\left(\vphantom{\rule{0pt}{1cm}}\right.$};
    \node at (\x+4,\y-0.15) {$ \left.\vphantom{\rule{0pt}{1cm}} \right)$};
\end{tikzpicture}\\
&\textequal{Prop.~\ref{prop:outer-prod}}\ \
\begin{tikzpicture}[baseline=\baseline]
    \tikzset{tensor/.style = {minimum size = 0.5cm,shape = circle,thick,draw=black,inner sep = 0pt}, edge/.style = {thick,line width=.4mm},every loop/.style={}}
    \def\y{0}
    \def\x{0}
    \draw[edge] (\x+0.15,\y-0.1) to [out=100,in=80] (\x+1.85,\y-0.1);
    \draw[edge] (\x+1.45,\y-0.1) to [out=100,in=80] (\x+3,\y-0.1);
    \node[draw=none] () at (\x+1,\y+0.5) {\scalebox{0.7}{\dimcolor{$m$}}};
    \node[draw=none] () at (\x+2.25,\y+0.5) {\scalebox{0.7}{\dimcolor{$n$}}};
    \draw[edge] (\x+3.45,\y-0.1) to [out=100,in=80] (\x+5,\y-0.1);
    \draw[edge] (\x+4.65,\y-0.1) to [out=100,in=80] (\x+6.25,\y-0.1);
    \node[draw=none] () at (\x+4,\y+0.5) {\scalebox{0.7}{\dimcolor{$n$}}};
    \node[draw=none] () at (\x+5.5,\y+0.5) {\scalebox{0.7}{\dimcolor{$p$}}};
    \draw[edge, draw=red] (\x+1.45,\y-0.1) to [out=-100,in=-80] (\x+3.45,\y-0.1);
    \draw[edge, draw=red] (\x+3,\y-0.1) to [out=-100,in=-80] (\x+5,\y-0.1);
\end{tikzpicture}\\
&= 
n\begin{tikzpicture}[baseline=\baseline]
    \tikzset{tensor/.style = {minimum size = 0.5cm,shape = circle,thick,draw=black,inner sep = 0pt}, edge/.style = {thick,line width=.4mm},every loop/.style={}}
    \def\y{0}
    \def\x{0}
    \draw[edge] (\x+0.15,\y-0.1) to [out=100,in=80] (\x+1.85,\y-0.1);
    \draw[edge] (\x+1.45,\y-0.1) to [out=100,in=80] (\x+3,\y-0.1);
    \node[draw=none] () at (\x+1,\y+0.5) {\scalebox{0.7}{\dimcolor{$m$}}};
    \node[draw=none] () at (\x+2.25,\y+0.5) {\scalebox{0.7}{\dimcolor{$p$}}};
    \end{tikzpicture}.
\end{align*}
\end{proof}
\item 下一条命题是文献~\citep{gupta2018matrix} 中定理 3.3.15 第 (iv) 条的一个有用的特例。
\begin{proposition}
    设 $\Ab\in\RR^{m\times n}$ 是元素独立同分布、服从标准正态分布的随机矩阵，$\Xb\in\RR^{d\times m}$ 为常矩阵，则 
$$\EE(\tr(\Xb\Ab\Ab^\ts\Xb^\ts)^2) = n^2\norm{\Xb}_F^4 + 2n\tr((\Xb^\ts\Xb)^2).$$
\begin{proof}
    我们借助张量网络记法来证明这一恒等式：
\begin{align*}
\EE(\tr(\Xb\Ab\Ab^\ts\Xb^\ts)^2) 
&=
\EE\left(\begin{tikzpicture}[baseline=\baseline]
    \tikzset{tensor/.style = {minimum size = 0.5cm,shape = circle,thick,draw=black,inner sep = 0pt}, edge/.style = {thick,line width=.4mm},every loop/.style={}}
    \def\y{0.25}
    \def\x{0}
    \node[tensor,fill=red!40!white] (X1) at (\x,\y){\scalebox{0.85}{$\Xb$}};
    \node[tensor,fill=blue!40!white] (A1) at (\x+1,\y){\scalebox{0.85}{$\Ab$}};
    \node[tensor,fill=blue!40!white] (A2) at (\x+2,\y){\scalebox{0.85}{$\Ab$}};
    \node[tensor,fill=red!40!white] (X2) at (\x+3,\y){\scalebox{0.85}{$\Xb$}};
    \draw[edge](X1) -- (A1)node [midway, above]{\scalebox{0.7}{\dimcolor{$m$}}};
    \draw[edge](A1) -- (A2)node [midway, above]{\scalebox{0.7}{\dimcolor{$n$}}};
    \draw[edge](A2) -- (X2)node [midway, above]{\scalebox{0.7}{\dimcolor{$m$}}};
    \draw[edge](X1) to [out=-100,in=-80] (X2);
    \node[tensor,fill=red!40!white](X_1) at (\x+4,\y){\scalebox{0.85}{$\Xb$}};
    \node[tensor,fill=blue!40!white](A_1) at (\x+5,\y){\scalebox{0.85}{$\Ab$}};
    \node[tensor,fill=blue!40!white] (A_2) at (\x+6,\y){\scalebox{0.85}{$\Ab$}};
    \node[tensor,fill=red!40!white] (X_2) at (\x+7,\y){\scalebox{0.85}{$\Xb$}};
    \draw[edge](X_1) -- (A_1)node [midway, above]{\scalebox{0.7}{\dimcolor{$m$}}};
    \draw[edge](A_1) -- (A_2)node [midway, above]{\scalebox{0.7}{\dimcolor{$n$}}};
    \draw[edge](A_2) -- (X_2)node [midway, above]{\scalebox{0.7}{\dimcolor{$m$}}};
    \draw[edge](X_1) to [out=-100,in=-80] (X_2);
    \node[draw=none] () at (\x+1.5,\y-1) {\scalebox{0.7}{\dimcolor{$d$}}};
    \node[draw=none] () at (\x+5.5,\y-1) {\scalebox{0.7}{\dimcolor{$d$}}};
\end{tikzpicture}\right)\\
&\textequal{Prop.~\ref{prop:inner-expectation}}\ \
\begin{tikzpicture}[baseline=\baseline]
    \tikzset{tensor/.style = {minimum size = 0.5cm,shape = circle,thick,draw=black,inner sep = 0pt}, edge/.style = {thick,line width=.4mm},every loop/.style={}}
    \def\y{0.25}
    \def\x{0}
    \node[tensor,fill=blue!40!white] (A1) at (\x,\y){\scalebox{0.85}{$\Ab$}};
    \node[tensor,fill=blue!40!white] (A2) at (\x+1,\y){\scalebox{0.85}{$\Ab$}};
    \node[tensor,fill=blue!40!white] (A_1) at (\x+2,\y){\scalebox{0.85}{$\Ab$}};
    \node[tensor,fill=blue!40!white] (A_2) at (\x+3,\y){\scalebox{0.85}{$\Ab$}};
    \draw[edge](A1)--(A2)node [midway, above]{\scalebox{0.7}{\dimcolor{$n$}}};
    \draw[edge](A_1)--(A_2)node [midway, above]{\scalebox{0.7}{\dimcolor{$n$}}};
    \node[tensor,fill=red!40!white] (X1) at (\x,\y-1){\scalebox{0.85}{$\Xb$}};
    \node[tensor,fill=red!40!white] (X2) at (\x+1,\y-1){\scalebox{0.85}{$\Xb$}};
    \node[tensor,fill=red!40!white] (X_1) at (\x+2,\y-1){\scalebox{0.85}{$\Xb$}};
    \node[tensor,fill=red!40!white] (X_2) at (\x+3,\y-1){\scalebox{0.85}{$\Xb$}};
    \draw[edge](X1)--(X2)node [midway, above]{\scalebox{0.7}{\dimcolor{$d$}}};
    \draw[edge](X_1)--(X_2)node [midway, above]{\scalebox{0.7}{\dimcolor{$d$}}};
    \node at (\x-0.65,\y) {$\EE\left( \vphantom{\rule{0pt}{0.5cm}} \right.$};
    \node at (\x+3.45,\y) {$ \left.\vphantom{\rule{0pt}{0.5cm}} \right)$};
    \node at (\x-0.65,\y-1){$\EE\left(\vphantom{\rule{0pt}{0.5cm}}\right.$};
    \node at (\x+3.45,\y-1) {$ \left.\vphantom{\rule{0pt}{0.5cm}} \right)$};
    \draw[edge, draw=red](A1)--(X1)node [midway, left]{\scalebox{0.7}{\dimcolor{$m$}}};
    \draw[edge, draw=red](A2)--(X2)node [midway, right]{\scalebox{0.7}{\dimcolor{$m$}}};
    \draw[edge, draw=red](A_1)--(X_1)node [midway, left=-0.05cm]{\scalebox{0.7}{\dimcolor{$m$}}};
    \draw[edge, draw=red](A_2)--(X_2)node [midway, right=-0.05cm]{\scalebox{0.7}{\dimcolor{$m$}}};
\end{tikzpicture}\\
&=
\begin{tikzpicture}[baseline=\baseline]
    \tikzset{tensor/.style = {minimum size = 0.5cm,shape = circle,thick,draw=black,inner sep = 0pt}, edge/.style = {thick,line width=.4mm},every loop/.style={}}
    \def\y{0.15}
    \def\x{0}
    \node[tensor,fill=blue!40!white] (A1) at (\x,\y){\scalebox{0.85}{$\Ab$}};
    \node[tensor,fill=blue!40!white] (A2) at (\x+1,\y){\scalebox{0.85}{$\Ab$}};
    \node[tensor,fill=blue!40!white] (A_1) at (\x+2,\y){\scalebox{0.85}{$\Ab$}};
    \node[tensor,fill=blue!40!white] (A_2) at (\x+3,\y){\scalebox{0.85}{$\Ab$}};
    \draw[edge](A1)--(A2)node [midway, above]{\scalebox{0.7}{\dimcolor{$n$}}};
    \draw[edge](A_1)--(A_2)node [midway, above]{\scalebox{0.7}{\dimcolor{$n$}}};
    \node[tensor,fill=red!40!white] (X1) at (\x,\y-1){\scalebox{0.85}{$\Xb$}};
    \node[tensor,fill=red!40!white] (X2) at (\x+1,\y-1){\scalebox{0.85}{$\Xb$}};
    \node[tensor,fill=red!40!white] (X_1) at (\x+2,\y-1){\scalebox{0.85}{$\Xb$}};
    \node[tensor,fill=red!40!white] (X_2) at (\x+3,\y-1){\scalebox{0.85}{$\Xb$}};
    \draw[edge](X1)--(X2)node [midway, above]{\scalebox{0.7}{\dimcolor{$d$}}};
    \draw[edge](X_1)--(X_2)node [midway, above]{\scalebox{0.7}{\dimcolor{$d$}}};
    \node at (\x-0.65,\y) {$\EE\left( \vphantom{\rule{0pt}{0.5cm}} \right.$};
    \node at (\x+3.45,\y) {$ \left.\vphantom{\rule{0pt}{0.5cm}} \right)$};
    \draw[edge, draw=red](A1)--(X1)node [midway, left]{\scalebox{0.7}{\dimcolor{$m$}}};
    \draw[edge, draw=red](A2)--(X2)node [midway, right]{\scalebox{0.7}{\dimcolor{$m$}}};
    \draw[edge, draw=red](A_1)--(X_1)node [midway, left]{\scalebox{0.7}{\dimcolor{$m$}}};
    \draw[edge, draw=red](A_2)--(X_2)node [midway, right=-0.05cm]{\scalebox{0.7}{\dimcolor{$m$}}};
\end{tikzpicture},
\end{align*}
最后一个等号成立是因为矩阵 $\Xb$ 是常矩阵。由命题~\ref{prop:kronecker-exp} 可知：
\begin{align*}
\EE(\Ab^\ts\Ab)^{\outprod 2}
&=
n^2\begin{tikzpicture}[baseline=\baseline]
    \tikzset{tensor/.style = {minimum size = 0.5cm,shape = circle,thick,draw=black,inner sep = 0pt}, edge/.style = {thick,line width=.4mm},every loop/.style={}}
    \def\y{0}
    \def\x{0}
    \draw[edge] (\x+0.15,\y-0.1) to [out=100,in=80] (\x+1.85,\y-0.1);
    \draw[edge] (\x+2,\y-0.1) to [out=100,in=80] (\x+3.65,\y-0.1);
    \node[draw=none] () at (\x+1,\y+0.5) {\scalebox{0.7}{\dimcolor{$m$}}};
    \node[draw=none] () at (\x+2.85,\y+0.5) {\scalebox{0.7}{\dimcolor{$m$}}};
\end{tikzpicture}
+n\begin{tikzpicture}[baseline=\baseline]
    \tikzset{tensor/.style = {minimum size = 0.5cm,shape = circle,thick,draw=black,inner sep = 0pt}, edge/.style = {thick,line width=.4mm},every loop/.style={}}
    \def\y{0}
    \def\x{0}
    \draw[edge] (\x+0.15,\y-0.1) to [out=100,in=80] (\x+1.85,\y-0.1);
    \draw[edge] (\x+1.45,\y-0.1) to [out=100,in=80] (\x+3,\y-0.1);
    \node[draw=none] () at (\x+1,\y+0.5) {\scalebox{0.7}{\dimcolor{$m$}}};
    \node[draw=none] () at (\x+2.25,\y+0.5) {\scalebox{0.7}{\dimcolor{$m$}}};
\end{tikzpicture}
+
n\begin{tikzpicture}[baseline=\baseline]
    \tikzset{tensor/.style = {minimum size = 0.5cm,shape = circle,thick,draw=black,inner sep = 0pt}, edge/.style = {thick,line width=.4mm},every loop/.style={}}
    \def\y{0}
    \def\x{0}
    \draw[edge] (\x,\y-0.1) to [out=100,in=80] (\x+2,\y-0.1);
    \draw[edge] (\x+0.15,\y-0.1) to [out=100,in=80] (\x+1.85,\y-0.1);
    \node[draw=none] () at (\x+1,\y+0.65) {\scalebox{0.7}{\dimcolor{$m$}}};
    \node[draw=none] () at (\x+1,\y+0.25) {\scalebox{0.7}{\dimcolor{$m$}}};
\end{tikzpicture}.
\end{align*}
因此， 
\begin{align*}
    &\begin{tikzpicture}[baseline=\baseline]
    \tikzset{tensor/.style = {minimum size = 0.5cm,shape = circle,thick,draw=black,inner sep = 0pt}, edge/.style = {thick,line width=.4mm},every loop/.style={}}
    \def\y{0.15}
    \def\x{0}
    \node[tensor,fill=blue!40!white] (A1) at (\x,\y){\scalebox{0.85}{$\Ab$}};
    \node[tensor,fill=blue!40!white] (A2) at (\x+1,\y){\scalebox{0.85}{$\Ab$}};
    \node[tensor,fill=blue!40!white] (A_1) at (\x+2,\y){\scalebox{0.85}{$\Ab$}};
    \node[tensor,fill=blue!40!white] (A_2) at (\x+3,\y){\scalebox{0.85}{$\Ab$}};
    \draw[edge](A1)--(A2)node [midway, above]{\scalebox{0.7}{\dimcolor{$n$}}};
    \draw[edge](A_1)--(A_2)node [midway, above]{\scalebox{0.7}{\dimcolor{$n$}}};
    \node[tensor,fill=red!40!white] (X1) at (\x,\y-1){\scalebox{0.85}{$\Xb$}};
    \node[tensor,fill=red!40!white] (X2) at (\x+1,\y-1){\scalebox{0.85}{$\Xb$}};
    \node[tensor,fill=red!40!white] (X_1) at (\x+2,\y-1){\scalebox{0.85}{$\Xb$}};
    \node[tensor,fill=red!40!white] (X_2) at (\x+3,\y-1){\scalebox{0.85}{$\Xb$}};
    \draw[edge](X1)--(X2)node [midway, above]{\scalebox{0.7}{\dimcolor{$d$}}};
    \draw[edge](X_1)--(X_2)node [midway, above]{\scalebox{0.7}{\dimcolor{$d$}}};
    \node at (\x-0.65,\y) {$\EE\left( \vphantom{\rule{0pt}{0.5cm}} \right.$};
    \node at (\x+3.45,\y) {$ \left.\vphantom{\rule{0pt}{0.5cm}} \right)$};
    \draw[edge, draw=red](A1)--(X1)node [midway, left]{\scalebox{0.7}{\dimcolor{$m$}}};
    \draw[edge, draw=red](A2)--(X2)node [midway, right]{\scalebox{0.7}{\dimcolor{$m$}}};
    \draw[edge, draw=red](A_1)--(X_1)node [midway, left]{\scalebox{0.7}{\dimcolor{$m$}}};
    \draw[edge, draw=red](A_2)--(X_2)node [midway, right]{\scalebox{0.7}{\dimcolor{$m$}}};
\end{tikzpicture}\\
&=
\left(\begin{tikzpicture}[baseline=\baseline]
    \tikzset{tensor/.style = {minimum size = 0.5cm,shape = circle,thick,draw=black,inner sep = 0pt}, edge/.style = {thick,line width=.4mm},every loop/.style={}}
    \def\y{0}
    \def\x{0}
    \node[tensor,fill=red!40!white] (X1) at (\x,\y){\scalebox{0.85}{$\Xb$}};
    \node[tensor,fill=red!40!white] (X2) at (\x+1,\y){\scalebox{0.85}{$\Xb$}};
    \node[tensor,fill=red!40!white] (X_1) at (\x+2,\y){\scalebox{0.85}{$\Xb$}};
    \node[tensor,fill=red!40!white] (X_2) at (\x+3,\y){\scalebox{0.85}{$\Xb$}};
    \draw[edge](X1)--(X2)node [midway, above]{\scalebox{0.7}{\dimcolor{$d$}}};
    \draw[edge](X_1)--(X_2)node [midway, above]{\scalebox{0.7}{\dimcolor{$d$}}};
    \draw[edge,draw=red](X1) -- (\x,\y+0.85)node [midway, right]{\scalebox{0.7}{\dimcolor{$m$}}};
    \draw[edge,draw=red](X2) -- (\x+1,\y+0.85)node [midway, right]{\scalebox{0.7}{\dimcolor{$m$}}};
    \draw[edge,draw=red](X_1) -- (\x+2,\y+0.85)node [midway, right]{\scalebox{0.7}{\dimcolor{$m$}}};
    \draw[edge,draw=red](X_2) -- (\x+3,\y+0.85)node [midway, right]{\scalebox{0.7}{\dimcolor{$m$}}};
\end{tikzpicture}\right)\left(n^2\begin{tikzpicture}[baseline=\baseline]
    \tikzset{tensor/.style = {minimum size = 0.5cm,shape = circle,thick,draw=black,inner sep = 0pt}, edge/.style = {thick,line width=.4mm},every loop/.style={}}
    \def\y{0}
    \def\x{0}
    \draw[edge] (\x+0.15,\y-0.1) to [out=100,in=80] (\x+1.85,\y-0.1);
    \draw[edge] (\x+2,\y-0.1) to [out=100,in=80] (\x+3.65,\y-0.1);
    \node[draw=none] () at (\x+1,\y+0.5) {\scalebox{0.7}{\dimcolor{$m$}}};
    \node[draw=none] () at (\x+2.85,\y+0.5) {\scalebox{0.7}{\dimcolor{$m$}}};
\end{tikzpicture}
+n\begin{tikzpicture}[baseline=\baseline]
    \tikzset{tensor/.style = {minimum size = 0.5cm,shape = circle,thick,draw=black,inner sep = 0pt}, edge/.style = {thick,line width=.4mm},every loop/.style={}}
    \def\y{0}
    \def\x{0}
    \draw[edge] (\x+0.15,\y-0.1) to [out=100,in=80] (\x+1.85,\y-0.1);
    \draw[edge] (\x+1.45,\y-0.1) to [out=100,in=80] (\x+3,\y-0.1);
    \node[draw=none] () at (\x+1,\y+0.5) {\scalebox{0.7}{\dimcolor{$m$}}};
    \node[draw=none] () at (\x+2.25,\y+0.5) {\scalebox{0.7}{\dimcolor{$m$}}};
\end{tikzpicture}
+
n\begin{tikzpicture}[baseline=\baseline]
    \tikzset{tensor/.style = {minimum size = 0.5cm,shape = circle,thick,draw=black,inner sep = 0pt}, edge/.style = {thick,line width=.4mm},every loop/.style={}}
    \def\y{0}
    \def\x{0}
    \draw[edge] (\x,\y-0.1) to [out=100,in=80] (\x+2,\y-0.1);
    \draw[edge] (\x+0.15,\y-0.1) to [out=100,in=80] (\x+1.85,\y-0.1);
    \node[draw=none] () at (\x+1,\y+0.65) {\scalebox{0.7}{\dimcolor{$m$}}};
    \node[draw=none] () at (\x+1,\y+0.25) {\scalebox{0.7}{\dimcolor{$m$}}};
\end{tikzpicture}\right)\\
&=
n^2\begin{tikzpicture}[baseline=\baseline]
    \tikzset{tensor/.style = {minimum size = 0.5cm,shape = circle,thick,draw=black,inner sep = 0pt}, edge/.style = {thick,line width=.4mm},every loop/.style={}}
    \def\y{0}
    \def\x{0}
    \draw[edge] (\x+0.15,\y-0.1) to [out=90,in=90] (\x+1.85,\y-0.1);
    \draw[edge] (\x+2.5,\y-0.1) to [out=90,in=90] (\x+4.15,\y-0.1);
    \node[draw=none] () at (\x+1,\y+0.5) {\scalebox{0.7}{\dimcolor{$m$}}};
    \node[draw=none] () at (\x+3.35,\y+0.5) {\scalebox{0.7}{\dimcolor{$m$}}};
    \node[tensor,fill=red!40!white] (X1) at (\x+0.15,\y-0.85){\scalebox{0.85}{$\Xb$}};
    \node[tensor,fill=red!40!white] (X2) at (\x+1.85,\y-0.85){\scalebox{0.85}{$\Xb$}};
    \draw[edge, draw=red](X1)--(\x+0.15,\y-0.1);
    \draw[edge](X1)--(X2) node[midway,above]{\scalebox{0.7}{\dimcolor{$d$}}};
    \draw[edge, draw=red](X2)--(\x+1.85,\y-0.1);
    \node[tensor,fill=red!40!white] (X_1) at (\x+2.5,\y-0.85){\scalebox{0.85}{$\Xb$}};
    \node[tensor,fill=red!40!white] (X_2) at (\x+4.15,\y-0.85){\scalebox{0.85}{$\Xb$}};
    \draw[edge, draw=red](X_1)--(\x+2.5,\y-0.1);
    \draw[edge](X_1)--(X_2) node[midway,above]{\scalebox{0.7}{\dimcolor{$d$}}};
    \draw[edge, draw=red](X_2)--(\x+4.15,\y-0.1);
\end{tikzpicture}
+n\begin{tikzpicture}[baseline=\baseline]
    \tikzset{tensor/.style = {minimum size = 0.5cm,shape = circle,thick,draw=black,inner sep = 0pt}, edge/.style = {thick,line width=.4mm},every loop/.style={}}
    \def\y{0}
    \def\x{0}
    \draw[edge] (\x+0.15,\y-0.1) to [out=60,in=120] (\x+2.5,\y-0.1);
    \draw[edge] (\x+1.85,\y-0.1) to [out=60,in=120] (\x+4.15,\y-0.1);
    \node[draw=none] () at (\x+1.15,\y+0.6) {\scalebox{0.7}{\dimcolor{$m$}}};
    \node[draw=none] () at (\x+3.25,\y+0.6) {\scalebox{0.7}{\dimcolor{$m$}}};
    \node[tensor,fill=red!40!white] (X1) at (\x+0.15,\y-0.85){\scalebox{0.85}{$\Xb$}};
    \node[tensor,fill=red!40!white] (X2) at (\x+1.85,\y-0.85){\scalebox{0.85}{$\Xb$}};
    \draw[edge, draw=red](X1)--(\x+0.15,\y-0.1);
    \draw[edge](X1)--(X2) node[midway,above]{\scalebox{0.7}{\dimcolor{$d$}}};
    \draw[edge, draw=red](X2)--(\x+1.85,\y-0.1);
    \node[tensor,fill=red!40!white] (X_1) at (\x+2.5,\y-0.85){\scalebox{0.85}{$\Xb$}};
    \node[tensor,fill=red!40!white] (X_2) at (\x+4.15,\y-0.85){\scalebox{0.85}{$\Xb$}};
    \draw[edge, draw=red](X_1)--(\x+2.5,\y-0.1);
    \draw[edge](X_1)--(X_2) node[midway,above]{\scalebox{0.7}{\dimcolor{$d$}}};
    \draw[edge, draw=red](X_2)--(\x+4.15,\y-0.1);
\end{tikzpicture}
+n \begin{tikzpicture}[baseline=\baseline]
    \tikzset{tensor/.style = {minimum size = 0.5cm,shape = circle,thick,draw=black,inner sep = 0pt}, edge/.style = {thick,line width=.4mm},every loop/.style={}}
    \def\y{0}
    \def\x{0}
    \draw[edge] (\x+0.15,\y-0.1) to [out=30,in=150] (\x+4.15,\y-0.1);
    \draw[edge] (\x+1.85,\y-0.1) to [out=30,in=150] (\x+2.5,\y-0.1);
    \node[draw=none] () at (\x+2.15,\y+0.6) {\scalebox{0.7}{\dimcolor{$m$}}};
    \node[draw=none] () at (\x+2.15,\y+0.15) {\scalebox{0.6}{\dimcolor{$m$}}};\node[tensor,fill=red!40!white] (X1) at (\x+0.15,\y-0.85){\scalebox{0.85}{$\Xb$}};
    \node[tensor,fill=red!40!white] (X2) at (\x+1.85,\y-0.85){\scalebox{0.85}{$\Xb$}};
    \draw[edge, draw=red](X1)--(\x+0.15,\y-0.1);
    \draw[edge](X1)--(X2) node[midway,above]{\scalebox{0.7}{\dimcolor{$d$}}};
    \draw[edge, draw=red](X2)--(\x+1.85,\y-0.1);
    \node[tensor,fill=red!40!white] (X_1) at (\x+2.5,\y-0.85){\scalebox{0.85}{$\Xb$}};
    \node[tensor,fill=red!40!white] (X_2) at (\x+4.15,\y-0.85){\scalebox{0.85}{$\Xb$}};
    \draw[edge, draw=red](X_1)--(\x+2.5,\y-0.1);
    \draw[edge](X_1)--(X_2) node[midway,above]{\scalebox{0.7}{\dimcolor{$d$}}};
    \draw[edge, draw=red](X_2)--(\x+4.15,\y-0.1);
\end{tikzpicture}\\
&=
n^2\begin{tikzpicture}[baseline=\baseline]
    \tikzset{tensor/.style = {minimum size = 0.5cm,shape = circle,thick,draw=black,inner sep = 0pt}, edge/.style = {thick,line width=.4mm},every loop/.style={}}
    \def\y{0}
    \def\x{0}
    \draw[edge] (\x+0.15,\y-0.1) to [out=100,in=80] (\x+1.85,\y-0.1);
    \draw[edge] (\x+2.5,\y-0.1) to [out=100,in=80] (\x+4.15,\y-0.1);
    \node[draw=none] () at (\x+1,\y+0.5) {\scalebox{0.7}{\dimcolor{$m$}}};
    \node[draw=none] () at (\x+3.35,\y+0.5) {\scalebox{0.7}{\dimcolor{$m$}}};
    \node[tensor,fill=red!40!white] (X1) at (\x+0.15,\y-0.85){\scalebox{0.85}{$\Xb$}};
    \node[tensor,fill=red!40!white] (X2) at (\x+1.85,\y-0.85){\scalebox{0.85}{$\Xb$}};
    \draw[edge](X1)--(\x+0.15,\y-0.1);
    \draw[edge](X1)--(X2) node[midway,above]{\scalebox{0.7}{\dimcolor{$d$}}};
    \draw[edge](X2)--(\x+1.85,\y-0.1);
    \node[tensor,fill=red!40!white] (X_1) at (\x+2.5,\y-0.85){\scalebox{0.85}{$\Xb$}};
    \node[tensor,fill=red!40!white] (X_2) at (\x+4.15,\y-0.85){\scalebox{0.85}{$\Xb$}};
    \draw[edge](X_1)--(\x+2.5,\y-0.1);
    \draw[edge](X_1)--(X_2) node[midway,above]{\scalebox{0.7}{\dimcolor{$d$}}};
    \draw[edge](X_2)--(\x+4.15,\y-0.1);
\end{tikzpicture}
+2n\begin{tikzpicture}[baseline=\baseline]
    \tikzset{tensor/.style = {minimum size = 0.5cm,shape = circle,thick,draw=black,inner sep = 0pt}, edge/.style = {thick,line width=.4mm},every loop/.style={}}
    \def\y{0}
    \def\x{0}
    \node[tensor,fill=red!40!white] (X1) at (\x,\y){\scalebox{0.85}{$\Xb$}};
    \node[tensor,fill=red!40!white] (X2) at (\x+1,\y){\scalebox{0.85}{$\Xb$}};
    \node[tensor,fill=red!40!white] (X_1) at (\x,\y-1){\scalebox{0.85}{$\Xb$}};
    \node[tensor,fill=red!40!white] (X_2) at (\x+1,\y-1){\scalebox{0.85}{$\Xb$}};
    \draw[edge](X_1) -- (X_2)node [midway, above]{\scalebox{0.7}{\dimcolor{$d$}}};
    \draw[edge](X1) -- (X2)node [midway, above]{\scalebox{0.7}{\dimcolor{$d$}}};
    \draw[edge](X1)--(X_1) node [midway, left]{\scalebox{0.7}{\dimcolor{$m$}}};
    \draw[edge](X2)--(X_2) node [midway, right]{\scalebox{0.7}{\dimcolor{$m$}}};
\end{tikzpicture}\\
&= 
n^2\norm{\Xb}_F^4+2n\tr((\Xb^\ts\Xb)^2).
\end{align*}
\end{proof}
\end{proposition}
\end{enumerate}

本章给出的这些推导表明，对张量网络表达式加以运算，能够帮助我们以极为直观的方式得到非平凡的结果。 

    

\section{结论}\label{chapter:conclusion}

本手稿提出了一套自成体系的张量网络图形语言，从基础记法与基本概念出发，逐步展开到张量分解、梯度计算与概率分布。其主要目标，是为在高维对象上执行从简单到复杂的各类运算提供可用的框架。
贯穿全文的一个核心观点是：张量网络的图形语言并不只是记法上的便利。
恰恰相反，它提供了一套清晰而精确的数学框架，把许多经典结果收拢进同一个统一的结构之中。
若干例子可以说明这一点。在张量图中，把行腿与列腿分开的任一割的权重决定了矩阵化秩的界（定理~\ref{thm: matricization-rank}）。第~\ref{chapter:gradients} 章的梯度法则（定理~\ref{thm:gradients} 与定理~\ref{thm:multiple-tensors-gradients}）把
复杂张量网络表达式的雅可比矩阵计算归结为一个单一的图示操作，
即删去节点；经典恒等式
$\partial(\Ab\xb)/\partial \Ab = \xb \outprod \Ib$ 与 $\partial \tr(\Ab)/\partial \Ab = \Ib$
便成为立即可得的推论。第~\ref{chapter:prob} 章的概率计算同样表明，
一系列非平凡结果都能用张量网络记法给出简洁的推导，其中包括 Wishart 分布的均值（命题~\ref{prop:wishart-mean}）、高斯张量的 Isserlis 定理（定理~\ref{thm:isserlis}），以及高斯随机矩阵乘积范数的期望（命题~\ref{prop:expected-norm}）。

综合来看，这些例子凸显出一条普遍的原则：在处理高维张量时，
把计算组织成张量网络的收缩、以图示的方式进行推理，
所得到的表示便既简洁，
又在概念上一目了然。

















\section*{致谢}
作者感谢 Jacob Miller 富有启发的讨论，他也是促使我们撰写本手稿的最初动力之一。
我们衷心感谢加拿大高等研究院（CIFAR AI chair program）、加拿大自然科学工程
研究理事会（Discovery program, RGPIN-2019-05949）以及 IVADO 为本研究提供的资助。

\\end{document}
